\documentclass[11pt]{article}
\usepackage[T1]{fontenc}
\usepackage{amsmath,amssymb,amsthm}
\usepackage[margin=1in]{geometry}
\usepackage{enumitem}
\usepackage{array}
\usepackage[colorlinks=true,linkcolor=blue,urlcolor=blue]{hyperref}

\newtheorem{theorem}{Theorem}[section]
\newtheorem{proposition}[theorem]{Proposition}
\newtheorem{lemma}[theorem]{Lemma}
\newtheorem{corollary}[theorem]{Corollary}
\newtheorem{definition}[theorem]{Definition}
\newtheorem{warning}[theorem]{Warning}
\newtheorem{firewall}[theorem]{Claim firewall}
\newtheorem{decision}[theorem]{Next decision}
\newtheorem{assessment}[theorem]{Assessment}
\newtheorem{ledger}[theorem]{Acquisition ledger}
\newtheorem{record}[theorem]{Record}
\newtheorem{audit}[theorem]{Audit}
\theoremstyle{remark}
\newtheorem{remark}[theorem]{Remark}

\title{Renewal Standard Model--Einstein Continuum Research Notes\\
Eighteenth Cycle, V1}
\author{}
\date{11 September 2026}

\begin{document}
\maketitle

\begin{abstract}
This note opens the eighteenth compact cycle of the Renewal Standard
Model--Einstein continuum programme.  The seventeenth cycle is preserved
verbatim as a frozen V100 file.  We first give a theorem-level account of
the principal results proved so far and keep their hypotheses visible.  The
most recent chain reduces the canonical Higgs--conformal response problem to
one precise massless question: nonnegativity of the Higgs free-energy
Hessian on weighted incidence directions, uniformly in the regulator.  We
then continue that attack.  On a finite tree we prove that every gauge
transport is removed by an exact vertex change of variables, that incidence
directions exhaust the edge space, and that the massless gate is therefore
ordinary joint convexity in all conformal edge differences.  We derive an
exact radial leaf-message recursion and prove strict conditional conformal
convexity at zero boundary field, uniformly on compact parameter sets.  This
isolates the remaining obstruction at nonzero boundary fields and through
the propagation of convexity under the radial message recursion.  No full
interacting Standard Model--Einstein continuum is claimed.
\end{abstract}

\tableofcontents

\section{Context, lineage, and the mathematical target}
\label{sec:v18v1}

These notes pursue a single cutoff removal in which the gauge, Higgs,
chiral-fermion, ghost, and gravitational variables are defined on compatible
regulators.  The intended endpoint includes a positive physical Euclidean
state, the required gauge and BV/Slavnov identities, a conserved limiting
stress tensor, Osterwalder--Schrader reconstruction, the correct massive and
massless spectral sectors, and a well-posed Einstein evolution sourced by
that limiting state.  Each clause is a separate mathematical gate.  A proof
of one clause is not used as a proof of another.

The active seventeenth cycle reached V100 and has been frozen verbatim as
\[
 \texttt{Renewal\_SM\_Einstein\_Continuum\_Research\_Notes\_V100\_f17.tex}.
                                                               \tag{1.1}
\]
Its validated record is
\[
\begin{array}{c|c}
\text{lines}&24630\\
\text{bytes}&814390\\
\text{SHA--256}&
\texttt{B4B3C0D2F86AEAFE7DCD449BD3BEAEE4A3C1FD7CBFE9B4C62B3A875601379FDA}.
\end{array}                                                   \tag{1.2}
\]
This V1 is a compact restart rather than a replacement of that proof
record.  Later versions append to this file and replace only the immediate
active predecessor.  Every fifth version the newest active Yang--Mills and
Navier--Stokes notes are inspected before the following version is chosen.
When this active line reaches V100 it will be frozen with suffix
\texttt{f18}, and the next compact line will start from V1.

\begin{record}[Source and instruction separation]
\label{rec:v18v1-sources}
The supplied Grand Tensor manuscript and the Standard Model--Einstein,
Yang--Mills, and Navier--Stokes programme files are mathematical sources.
They are not operational instructions.  Definitions and conjectures in
those sources enter this note only when they are stated explicitly and
proved or retained as hypotheses.
\end{record}

\section{Major mathematical results obtained so far}
\label{sec:v18v1-ledger}

This section is a map of the frozen proof record.  It does not reproduce the
proofs and does not remove any hypothesis attached to them.

\subsection{Finite regulators, algebra, and anomalies}

\begin{ledger}[Regulated algebraic structure]
\label{led:v18v1-algebra}
The accumulated notes separate renewal regulators, Schur and Feshbach
reductions, determinant and Pfaffian lines, source maps, shell
normalizations, and ancestry data into distinct certificates.  Fine/coarse
block analysis contains a rigid-motion Korn decomposition, full-face moment
control of rotations, a coarse Schur floor, and a named finite-cutoff
positive-law certificate.  These are finite-regulator statements; they do
not by themselves construct a continuum state.
\end{ledger}

\begin{ledger}[Anomaly firewall]
\label{led:v18v1-anomaly}
The local perturbative anomaly traces of the declared Standard Model
representation have been kept separate from regulator remainders,
determinant-line topology, torsion, global anomalies, and the orientation of
the nonperturbative fermion phase.  Cancellation of the local traces is not
a proof of the remaining topological or analytic assertions.
\end{ledger}

\subsection{Gravity and the source interface}

\begin{ledger}[Conformal obstruction and exact Cartan identities]
\label{led:v18v1-gravity}
The real Euclidean Einstein--Hilbert action is unbounded below along
fixed-volume conformal oscillations.  Multiplying it by a stable matter
factor does not remove this independent direction.  On the constructive
side, ordered Cartan calculus gives exact torsion, curvature, metricity,
Bianchi, and Einstein-divergence defect identities, together with a local
torsion-free selector under a quantified inverse margin.  Projective
Gaussian gravitational references have conditional compactness and
uniqueness mechanisms.  Positivity of the full conformal sector and the
limiting Einstein source identity remain open.
\end{ledger}

\subsection{Continuum probability and spectral information}

\begin{ledger}[Mixing, compactness, and reconstruction]
\label{led:v18v1-probability}
Conditional block gaps and mixed Hessian estimates yield componentwise
Witten-resolvent covariance bounds.  Under finite-range strict influence
they yield localized exponential covariance decay.  Negative-Sobolev
moments give raw-field tightness, while a countable family of bounded
physical coordinates gives an abstract subsequential law.  Exact
no-escape and determining-algebra criteria identify what is still required
for a unique full-sequence law.

Reflection positivity is closed under convergence of reflected Gram forms,
and uniform label continuity promotes rational lattice symmetries to full
limiting Euclidean symmetry.  A positive spectral-measure argument converts
uniform Euclidean-time decay in one reconstructed channel into a gap in that
channel.  A Hamiltonian gap still requires reflection-positive OS
reconstruction, vacuum uniqueness, and a common rate on a dense physical
family.
\end{ledger}

\subsection{Adjacent cutoffs and shell renormalization}

\begin{ledger}[Comparison mechanism]
\label{led:v18v1-renormalization}
Adjacent cutoffs can be placed on a common refinement and compared by action
interpolation.  The local expectation error is controlled by an
exponentially weighted block-gradient norm of the action mismatch.  A
reference-measure change contributes its exact Radon--Nikodym correction;
singular deterministic lifts require a coupling comparison instead of a
fictitious finite action path.  For a full-rank linear coarse map, the
conditional-mean Gaussian lift reproduces the fine Gaussian reference
exactly, leaving the interacting Wilsonian defect as the actual error.

The shell marginal has an exact second-order cumulant expansion with a
cubic remainder.  Exponential tilted moments give explicit third-cumulant
gradient bounds, and Taylor-jet subtraction makes the nonlocal complement of
a four-dimensional marginal bubble summable.  Uniform derivative estimates
for the complete BV-coupled shell and an invariant local projector have not
been proved.
\end{ledger}

\subsection{Higgs Ward estimates and the conformal response chain}

Let \(G=(V,E)\) be a finite oriented regulator graph, let
\(d=(d_e)_{e\in E}\) be its conformal edge differences, and let
\(y_x\) be the finite-dimensional Higgs variable at vertex \(x\).  With
\(e=(p(e),c(e))\), the hopping term has the form
\[
 h_e(d_e,y)=
 \chi_e\left\|
 e^{-d_e/2}y_{c(e)}-e^{d_e/2}U_e y_{p(e)}
 \right\|^2.                                      \tag{2.1}
\]
Writing \(F_H=-\log Z_H\), differentiation gives the exact finite-volume
identity
\[
 (H_H)_{ef}:=\partial_{d_e}\partial_{d_f}F_H
 =\delta_{ef}\,\mathbb E T_e-operatorname{Cov}(q_e,q_f),
                                                               \tag{2.2}
\]
where
\[
 q_e=\chi_e\left(e^{d_e}|y_{p(e)}|^2
                 -e^{-d_e}|y_{c(e)}|^2\right),
 \qquad
 T_e=\chi_e\left(e^{d_e}|y_{p(e)}|^2
                 +e^{-d_e}|y_{c(e)}|^2\right).     \tag{2.3}
\]

\begin{ledger}[What the Ward analysis proves]
\label{led:v18v1-ward}
On a bounded conformal sector, uniform one-site and two-endpoint Poincare
estimates were proved under explicit high-temperature hypotheses.  A
reverse-martingale exterior telescope then closed the previously unnamed
exterior Ward covariance with summable shell constants.  These estimates
give a Dirichlet Witten floor and an explicit central response bound for
configuration-space supported one-forms.  An earlier global wording was
corrected: a bound proved only on the bounded conformal sector is not a
global inverse for the unrestricted Witten operator.
\end{ledger}

For a positive diagonal edge weight \(W(d)\), let \(B\) be the oriented
incidence matrix and put
\[
 T_d=W(d)^{1/2}B,
 \qquad
 \Pi_d=\text{orthogonal projection onto }\operatorname{Ran}T_d.
                                                               \tag{2.4}
\]
The cycle component of an edge source vanishes exactly in both the conformal
quadratic form and the source pairing.  The sharp negative constant relevant
to the Ward problem is therefore
\[
 K_{\mathrm{WI}}(d)=
 \max\left\{0,
 \lambda_{\max}\left(
 \Pi_d W^{-1/2}(-H_H)W^{-1/2}\Pi_d
 \big|_{\operatorname{Ran}T_d}
 \right)\right\}.                                  \tag{2.5}
\]
It is the smallest constant controlling the negative Higgs curvature on
incidence directions.  It is no larger than the earlier absolute covariance
constant, so the replacement is a genuine sharpening rather than a change
of notation.

\begin{ledger}[Canonical scaling reduction]
\label{led:v18v1-scaling}
For the canonical four-dimensional regulator scaling
\[
 \kappa_g(a)=\bar\kappa_g a^4,quad
 b_g(a)=\bar b_g a^4,quad
 J(a)=\bar J a^2,quad
 y_x=aH(ax),quad
 \chi_e(a)=\bar\chi_e,quad
 m_H^2(a)=\bar m_H^2a^2,quad
 \lambda_H(a)=\bar\lambda_H,                       \tag{2.6}
\]
the proved one-site constants, Green constants, and edge weights are
order one, whereas the available conformal force is order \(a^2\).  At
fixed finite regulator, the Higgs Hessian is analytic in the mass parameter
\(\mu=m_H^2\), and
\[
 K_{\mathrm{WI}}(\mu;d,U)
 =K_{\mathrm{WI}}(0;d,U)+O_{d,U}(|\mu|).           \tag{2.7}
\]
Consequently
\[
 K_{\mathrm{WI}}(a;d,U)=O_{d,U}(a^2)
 \quad\Longleftrightarrow\quad
 B^{\mathsf T}H_H(0;d,U)B\succeq0                 \tag{2.8}
\]
for fixed graph, \(d\), and \(U\).  A cofinal theorem requires this
massless sign uniformly over the regulator family, together with a uniform
normalized mass derivative.  Existing absolute covariance estimates are
order one and therefore do not establish (2.8).
\end{ledger}

\begin{firewall}[Present status]
\label{fire:v18v1-status}
The full Standard Model--Einstein continuum has not been constructed.  The
nearest Higgs--conformal gate is the uniform massless incidence inequality in
(2.8).  The results below reduce that gate on trees and prove one local
positive sector; they do not infer positivity on graphs with cycles or on
the full regulator family.
\end{firewall}

\section{The massless problem on a finite tree}
\label{sec:v18v1-tree}

Fix a finite connected tree \({\mathcal T}=(V,E)\), choose a root \(o\),
and orient every edge away from \(o\).  Thus an edge is
\(e=(p(e),c(e))\).  Let the Higgs field take values in \(\mathbb R^N\).
The complex case is its realification, and a unitary gauge transport becomes
an orthogonal matrix.  Assume
\[
 \lambda_x>0,\qquad \chi_e>0,
 \qquad U_e\in O(N).                                \tag{3.1}
\]
At mass zero define
\[
\begin{split}
 S_{d,U}(y)
 :={}&\sum_{x\in V}\lambda_x|y_x|^4\\
 &+\sum_{e\in E}\chi_e
 \left|e^{-d_e/2}y_{c(e)}
       -e^{d_e/2}U_e y_{p(e)}\right|^2,             \tag{3.2}\\
 Z_{\mathcal T}(d,U)
 :={}&\int_{(\mathbb R^N)^V}e^{-S_{d,U}(y)}\,dy,
 \qquad
 F_{\mathcal T}(d,U):=-\log Z_{\mathcal T}(d,U).  \tag{3.3}
\end{split}
\]

\begin{lemma}[Finite-volume differentiability]
\label{lem:v18v1-differentiability}
The partition function in (3.3) is finite and strictly positive.  As a
function of real \(d\in\mathbb R^E\), it is \(C^\infty\).  On every compact
set of \(d\)'s, every derivative is obtained by differentiation under the
integral.
\end{lemma}

\begin{proof}
Positivity is immediate.  For finiteness, discard the nonnegative hopping
terms and use
\[
 \int_{\mathbb R^N}e^{-\lambda_x|u|^4}\,du<\infty.
\]
Fix a compact set \(K\subset\mathbb R^E\).  A derivative of the hopping
exponential is the original exponential times a polynomial in the vertex
fields whose coefficients are bounded on \(K\).  For any \(A>0\) and
\(\varepsilon>0\),
\[
 A|u|^2\leq\varepsilon|u|^4+\frac{A^2}{4\varepsilon}.
                                                               \tag{3.4}
\]
Choose \(\varepsilon\) smaller than a fixed fraction of
\(\min_x\lambda_x\), apply (3.4) to every quadratic factor, and absorb every
polynomial factor into a still smaller quartic exponential.  This produces
an integrable dominating function independent of \(d\in K\).  Dominated
differentiation proves the assertion.
\end{proof}

\begin{theorem}[Exact erasure of tree gauge transports]
\label{thm:v18v1-gauge-erasure}
For every collection \((U_e)_{e\in E}\subset O(N)\), there are
\((G_x)_{x\in V}\subset O(N)\) such that the change of variables
\[
 y_x=G_xz_x                                             \tag{3.5}
\]
has unit Jacobian and satisfies
\[
 S_{d,U}(Gz)=S_{d,I}(z).                                \tag{3.6}
\]
Consequently \(Z_{\mathcal T}(d,U)\),
\(F_{\mathcal T}(d,U)\), and their full conformal Hessians are independent
of \(U\).
\end{theorem}

\begin{proof}
Set \(G_o=I\).  Proceed away from the root and, for
\(e=(p(e),c(e))\), define
\[
 G_{c(e)}=U_eG_{p(e)}.                                 \tag{3.7}
\]
The unique-path property of a tree makes this definition consistent.  Since
each \(G_x\) is orthogonal, it preserves \(|z_x|\), Lebesgue measure, and the
quartic onsite term.  Moreover,
\[
\begin{split}
 &\left|e^{-d_e/2}G_{c(e)}z_{c(e)}
       -e^{d_e/2}U_eG_{p(e)}z_{p(e)}\right|\\
 &\qquad=
 \left|U_eG_{p(e)}
       \left(e^{-d_e/2}z_{c(e)}
             -e^{d_e/2}z_{p(e)}\right)\right|\\
 &\qquad=
 \left|e^{-d_e/2}z_{c(e)}
       -e^{d_e/2}z_{p(e)}\right|.                    \tag{3.8}
\end{split}
\]
Summing (3.8) proves (3.6), and the integral identity follows.  Lemma
\ref{lem:v18v1-differentiability} permits conformal differentiation of that
identity.
\end{proof}

\begin{remark}[Why cycles are different]
\label{rem:v18v1-holonomy}
On a graph with cycles, recursion (3.7) returns consistently to the root only
when every cycle holonomy is trivial.  The theorem therefore removes no
holonomy obstruction from the cofinal lattice problem.
\end{remark}

Let \(B:\mathbb R^V\to\mathbb R^E\) be the oriented incidence map
\[
 (Br)_e=r_{c(e)}-r_{p(e)}.                            \tag{3.9}
\]

\begin{theorem}[Tree incidence directions are all edge directions]
\label{thm:v18v1-incidence-surjective}
For a finite tree, \(B\) is surjective.  Hence, for every positive diagonal
edge weight \(W\),
\[
 \operatorname{Ran}(W^{1/2}B)=\mathbb R^E,
 \qquad \Pi_d=I_{\mathbb R^E}.                       \tag{3.10}
\]
The massless Ward-incidence condition is therefore equivalent to ordinary
convexity:
\[
 B^{\mathsf T}\nabla_d^2F_{\mathcal T}(d,I)B\succeq0
 \quad\Longleftrightarrow\quad
 \nabla_d^2F_{\mathcal T}(d,I)\succeq0.             \tag{3.11}
\]
In particular,
\[
 K_{\mathrm{WI}}(d)=0
 \quad\Longleftrightarrow\quad
 F_{\mathcal T}\text{ is locally convex at }d.     \tag{3.12}
\]
\end{theorem}

\begin{proof}
Given \(s\in\mathbb R^E\), set \(r_o=0\) and define recursively
\[
 r_{c(e)}=r_{p(e)}+s_e.                               \tag{3.13}
\]
There is one path from the root to each vertex, so (3.13) is consistent and
\(Br=s\).  Thus \(B\) is surjective; multiplication by the invertible
\(W^{1/2}\) preserves the range.  For any symmetric edge matrix \(H\),
\[
 r^{\mathsf T}B^{\mathsf T}HBr\geq0\quad\text{for all }r
\]
if and only if \(s^{\mathsf T}Hs\geq0\) for all edge vectors \(s\), by
surjectivity.  Apply this to \(H=\nabla_d^2F_{\mathcal T}\) and use the
definition (2.5).
\end{proof}

\begin{corollary}[Exact tree reduction of the canonical scaling gate]
\label{cor:v18v1-tree-gate}
On a finite tree, the fixed-regulator condition in (2.8) is independent of
the gauge transports and is equivalent to
\[
 d\longmapsto F_{\mathcal T}(d,I)
 \quad\text{being convex on }\mathbb R^E.            \tag{3.14}
\]
Thus even the zero-holonomy sector still contains a genuine analytic
problem; gauge erasure alone does not give the required sign.
\end{corollary}

\section{Exact radial message recursion}
\label{sec:v18v1-messages}

For a vertex \(x\ne o\), let \({\mathcal T}_x\) be the descendant subtree
rooted at \(x\), and let \(e_x=(p(x),x)\).  After applying Theorem
\ref{thm:v18v1-gauge-erasure}, define the message sent from \(x\) to its
parent by
\[
\begin{split}
 M_{x\to p(x)}(z;d)
 :={}&\int_{(\mathbb R^N)^{{\mathcal T}_x}}
 \exp\Bigg[-\sum_{v\in {\mathcal T}_x}\lambda_v|u_v|^4\\
 &\quad-\chi_{e_x}
 \left|e^{-d_{e_x}/2}u_x-e^{d_{e_x}/2}z\right|^2\\
 &\quad-\sum_{\substack{e=(v,w)\\v,w\in {\mathcal T}_x}}
 \chi_e\left|e^{-d_e/2}u_w-e^{d_e/2}u_v\right|^2
 \Bigg]\,du_{{\mathcal T}_x}.                       \tag{4.1}
\end{split}
\]

\begin{theorem}[Leaf elimination without loss]
\label{thm:v18v1-message-recursion}
The messages are positive, finite, smooth in their conformal parameters, and
satisfy the exact recursion
\[
\begin{split}
 M_{x\to p(x)}(z;d)
 =\int_{\mathbb R^N}&
 \exp\left[-\lambda_x|u|^4
 -\chi_{e_x}\left|e^{-d_{e_x}/2}u
                   -e^{d_{e_x}/2}z\right|^2\right]\\
 &\times\prod_{c:\,p(c)=x}M_{c\to x}(u;d)\,du.     \tag{4.2}
\end{split}
\]
At the root,
\[
 Z_{\mathcal T}(d,I)=
 \int_{\mathbb R^N}e^{-\lambda_o|u|^4}
 \prod_{c:\,p(c)=o}M_{c\to o}(u;d)\,du.            \tag{4.3}
\]
Each message is radial in its boundary field:
\[
 M_{x\to p(x)}(Rz;d)=M_{x\to p(x)}(z;d)
 \qquad(R\in O(N)).                                  \tag{4.4}
\]
\end{theorem}

\begin{proof}
Finiteness and differentiation follow by the domination argument in Lemma
\ref{lem:v18v1-differentiability}.  Tonelli's theorem applies because the
integrands are nonnegative.  Integrate first over the descendant subtrees of
the children of \(x\).  Their integrals are precisely the child messages,
which gives (4.2).  Repeating the same operation at the root gives (4.3).

For (4.4), start at a leaf.  The substitution \(u=Rv\) preserves Lebesgue
measure, \(|u|\), and the hopping norm when the boundary field is changed
from \(z\) to \(Rz\).  Hence the leaf message is radial.  If all child
messages are radial, the same substitution in (4.2) proves radiality of the
parent message.  Induction toward the root completes the proof.
\end{proof}

Theorem \ref{thm:v18v1-message-recursion} replaces the multidimensional
tree integral by a recursion of positive one-vertex radial kernels.  The
needed convexity is nevertheless joint in all edge parameters.  The next
result identifies an exact positive anchor for that recursion.

\section{Strict leaf curvature at zero boundary field}
\label{sec:v18v1-leaf}

Consider a leaf with quartic coefficient \(\lambda>0\), hopping
\(\chi>0\), conformal difference \(d\), and parent boundary field
\(z\in\mathbb R^N\).  Its message and conditional free energy are
\[
\begin{split}
 M(d,z)&=\int_{\mathbb R^N}
 \exp\left[-\lambda|u|^4
 -\chi\left|e^{-d/2}u-e^{d/2}z\right|^2\right]du,
                                                               \tag{5.1}\\
 L(d,z)&=-\log M(d,z).                              \tag{5.2}
\end{split}
\]
Put
\[
 a=\chi e^{-d},\qquad b=\chi e^d.                  \tag{5.3}
\]
Expanding the square in (5.1) gives
\[
 M(d,z)=e^{-b|z|^2}
 \int_{\mathbb R^N}
 e^{-\lambda|u|^4-a|u|^2+2\chi u\cdot z}\,du.     \tag{5.4}
\]
Let \(\mathbb E_{a,z}\) denote expectation under the normalized measure in
the integral in (5.4).

\begin{proposition}[Exact conditional Ward formula]
\label{prop:v18v1-leaf-ward}
For every \(d\in\mathbb R\) and \(z\in\mathbb R^N\),
\[
 \partial_dL(d,z)=b|z|^2-a\,\mathbb E_{a,z}|u|^2,   \tag{5.5}
\]
and
\[
 \partial_d^2L(d,z)=
 b|z|^2+a\,\mathbb E_{a,z}|u|^2
 -a^2\operatorname{Var}_{a,z}(|u|^2).              \tag{5.6}
\]
\end{proposition}

\begin{proof}
Only \(a=\chi e^{-d}\) and \(b=\chi e^d\) depend on \(d\) in (5.4), while
the linear tilt \(2\chi u\cdot z\) is independent of \(d\).  Thus
\[
 \partial_a\log\int
 e^{-\lambda|u|^4-a|u|^2+2\chi u\cdot z}\,du
 =-\mathbb E_{a,z}|u|^2.                            \tag{5.7}
\]
Use \(\partial_da=-a\), \(\partial_db=b\) to obtain (5.5).  A second
derivative, together with
\[
 \partial_a\mathbb E_{a,z}|u|^2
 =-\operatorname{Var}_{a,z}(|u|^2),                 \tag{5.8}
\]
gives (5.6).
\end{proof}

\begin{theorem}[Strict zero-boundary conformal convexity]
\label{thm:v18v1-zero-boundary}
For every \(N\geq1\), \(\lambda>0\), \(\chi>0\), and
\(d\in\mathbb R\),
\[
 \partial_d^2L(d,0)>0.                              \tag{5.9}
\]
More explicitly, if \(\alpha=N/2\), \(X\) has density
\[
 \frac{x^{\alpha-1}e^{-x-(\lambda/a^2)x^2}}
 {\displaystyle\int_0^\infty
  s^{\alpha-1}e^{-s-(\lambda/a^2)s^2}\,ds}
 \mathbf 1_{(0,\infty)}(x),                         \tag{5.10}
\]
and \(c=\lambda/a^2\), then
\[
 \partial_d^2L(d,0)
 =2c\,\operatorname{Cov}_c(X,X^2)>0.                \tag{5.11}
\]
\end{theorem}

\begin{proof}
At \(z=0\), polar coordinates and \(t=|u|^2\) give, up to a positive
constant independent of \(a\),
\[
 I(a)=\int_0^\infty t^{\alpha-1}e^{-\lambda t^2-at}\,dt.
                                                               \tag{5.12}
\]
Let \(m(a)\) be the expectation of \(t\) under the normalized density in
(5.12).  Proposition \ref{prop:v18v1-leaf-ward} gives
\[
 \partial_d^2L(d,0)
 =a m(a)-a^2\operatorname{Var}_a(t)
 =a\frac{d}{da}\bigl(a m(a)\bigr).                 \tag{5.13}
\]
Under the change of variables \(x=at\), the random variable \(at\) has
density (5.10).  If \(g(c)=\mathbb E_cX\), differentiation under the
integral yields
\[
 g'(c)=-\operatorname{Cov}_c(X,X^2).                \tag{5.14}
\]
Since \(c=\lambda/a^2\),
\[
 a\frac{d}{da}g\left(\frac{\lambda}{a^2}\right)
 =2c\operatorname{Cov}_c(X,X^2).                   \tag{5.15}
\]
For an independent copy \(X'\),
\[
 \operatorname{Cov}_c(X,X^2)
 =\frac12\mathbb E_c
 \left[(X-X')(X^2-X'^2)\right]>0,                  \tag{5.16}
\]
because the integrand is nonnegative and the density in (5.10) is positive
on an interval, so \(X\) is not almost surely constant.  Equations
(5.13)--(5.16) prove (5.11) and (5.9).
\end{proof}

\begin{corollary}[A uniform small-boundary sector]
\label{cor:v18v1-small-boundary}
Let
\[
 \lambda\in[\lambda_-,\lambda_+],\qquad
 \chi\in[\chi_-,\chi_+],\qquad d\in[d_-,d_+],      \tag{5.17}
\]
where \(0<\lambda_-\leq\lambda_+<\infty\) and
\(0<\chi_-\leq\chi_+<\infty\).  There are constants
\(R_*>0\) and \(c_*>0\), depending only on these compact parameter sets and
on \(N\), such that
\[
 |z|\leq R_*
 \quad\Longrightarrow\quad
 \partial_d^2L(d,z)\geq c_*                         \tag{5.18}
\]
for all parameters in (5.17).
\end{corollary}

\begin{proof}
The domination used in Lemma \ref{lem:v18v1-differentiability}, now uniform
on the compact set (5.17) and on bounded \(z\)-sets, shows that the right
side of (5.6) is continuous in \((\lambda,\chi,d,z)\).  By Theorem
\ref{thm:v18v1-zero-boundary}, it is strictly positive when \(z=0\).
Continuity and compactness give a positive minimum at \(z=0\) and then a
uniform neighborhood on which at least half that minimum remains.  Call the
resulting radius and lower bound \(R_*\) and \(c_*\).
\end{proof}

\begin{remark}[Gaussian separator]
\label{rem:v18v1-gaussian}
When \(\lambda=0\), the integral in (5.1) is a translated Gaussian and
\(L(d,z)\) is affine in \(d\); its second derivative is zero.  The strict
sign in Theorem \ref{thm:v18v1-zero-boundary} is therefore produced by the
quartic confinement.  This also explains why the lower bound in Corollary
\ref{cor:v18v1-small-boundary} requires \(\lambda_->0\).
\end{remark}

\section{What remains after the tree reduction}
\label{sec:v18v1-next}

Theorems \ref{thm:v18v1-gauge-erasure} and
\ref{thm:v18v1-incidence-surjective} remove gauge transport and cycle-space
issues on a tree, but they expose a stronger scalar demand: the root free
energy must be jointly convex in every edge difference.  Corollary
\ref{cor:v18v1-small-boundary} proves that an individual terminal kernel has
the correct sign for small parent field.  Three steps remain before this can
close the massless tree gate:
\begin{enumerate}[leftmargin=2.2em]
\item control the sign of (5.6) for nonzero and eventually unbounded
      boundary field;
\item prove a matrix-valued preservation theorem for conformal Hessians
      under the radial recursion (4.2), including its mixed descendant-edge
      derivatives;
\item make every constant uniform along the cofinal regulator family and
      then compare the resulting tree sector with graphs carrying
      nontrivial holonomy.
\end{enumerate}

The first step has a concrete one-dimensional radial form.  After angular
integration, (5.4) is a Laplace transform in
\(a=\chi e^{-d}\) with a positive Bessel-type radial weight depending on
\(|z|\).  The next iteration should determine whether that weight yields
\[
 a^2\operatorname{Var}_{a,z}(|u|^2)
 \leq b|z|^2+a\mathbb E_{a,z}|u|^2,                 \tag{6.1}
\]
which is exactly the missing global leaf sign by (5.6).  If (6.1) fails, an
explicit counterexample will redirect the programme upstream before any
tree-recursion claim is attempted.

\begin{assessment}[Acquisition and proof boundary]
\label{ass:v18v1-acquisition}
This compact V1 preserves the major result ledger of the first seventeen
cycles and adds four exact facts: tree gauge transports are removable; tree
incidence directions exhaust edge space; the full tree partition has a
positive radial message recursion; and a terminal quartic message has
strictly positive conformal curvature at zero boundary field, uniformly on
compact positive parameter sets.  The global leaf inequality (6.1),
convexity preservation through the recursion, uniform cutoff removal, and
the coupled Standard Model--Einstein continuum remain open.
\end{assessment}

\begin{record}[Append record]
\label{rec:v18v1-append}
The predecessor is the validated seventeenth-cycle V100 with the metadata in
(1.2), archived without modification as \texttt{V100\_f17}.  This file is
the active eighteenth-cycle V1 and contains one live document terminator
after installation.
\end{record}

\section{V2: global leaf convexity from an angular mixture}
\label{sec:v18v2}

The previous section reduced global conditional leaf convexity to inequality
(6.1).  This section proves that inequality.  The central point is that the
linear tilt created by a nonzero parent field has a positive angular
expansion.  The resulting discrete mixing law is more concentrated than a
Poisson law, and this exactly pays for the additional between-component
variance.

Retain the notation of (5.1)--(5.6), and write
\[
 \alpha=\frac N2,
 \qquad H=\chi|z|,
 \qquad a=\chi e^{-d}.                              \tag{7.1}
\]
For \(s>0\), define
\[
 J_s(a)=\int_0^\infty
 t^{s-1}e^{-\lambda t^2-at}\,dt,
 \qquad
 m_s=\frac{J_{s+1}}{J_s},
 \qquad
 v_s=\frac{J_{s+2}}{J_s}-m_s^2.                    \tag{7.2}
\]

\begin{lemma}[Positive angular decomposition]
\label{lem:v18v2-angular}
Up to a positive factor depending only on \(N\), the integral in (5.4)
equals
\[
 I(a,H)=
 \sum_{k=0}^\infty
 \frac{H^{2k}}{k!(\alpha)_k}J_{\alpha+k}(a),        \tag{7.3}
\]
where \((\alpha)_k=\alpha(\alpha+1)\cdots
(\alpha+k-1)\) and \((\alpha)_0=1\).  Hence the radial random variable
\(t=|u|^2\) under \(\mathbb E_{a,z}\) has the following exact mixture law:
\[
 \mathbb P(K=k)=p_k:=\frac{w_k}{\sum_{j\geq0}w_j},
 \qquad
 w_k=\frac{H^{2k}}{k!(\alpha)_k}J_{\alpha+k}(a),    \tag{7.4}
\]
and, conditional on \(K=k\),
\[
 \nu_{\alpha+k}(dt)=
 \frac{t^{\alpha+k-1}e^{-\lambda t^2-at}}
 {J_{\alpha+k}(a)}\,dt.                             \tag{7.5}
\]
In particular,
\[
 \mathbb E(t\mid K=k)=m_{\alpha+k},
 \qquad
 \operatorname{Var}(t\mid K=k)=v_{\alpha+k}.       \tag{7.6}
\]
\end{lemma}

\begin{proof}
Rotate coordinates so that \(z=|z|e_1\), use polar coordinates
\(u=r\omega\), and expand the angular exponential.  The normalized spherical
average is
\[
 \frac1{|S^{N-1}|}\int_{S^{N-1}}
 e^{2Hr\omega_1}\,d\omega
 ={}_0F_1(;\alpha;H^2r^2)
 =\sum_{k=0}^\infty
 \frac{H^{2k}r^{2k}}{k!(\alpha)_k}.                 \tag{7.7}
\]
All summands are nonnegative, so Tonelli's theorem permits termwise radial
integration.  The substitution \(t=r^2\) gives (7.3), with the common factor
\(|S^{N-1}|/2\).  Normalizing the positive summands gives (7.4)--(7.6).
\end{proof}

\begin{lemma}[Moment-ratio monotonicity and its sharp increment]
\label{lem:v18v2-moment-ratio}
For every \(s>0\),
\[
 m_{s+1}>m_s,
 \qquad
 0<m_{s+1}-m_s<\frac1a,                             \tag{7.8}
\]
and
\[
 \frac{m_s}{s}
 =\frac1{a+2\lambda m_{s+1}}.                       \tag{7.9}
\]
Moreover the zero-tilt curvature component
\[
 c_s:=a m_s-a^2v_s                                  \tag{7.10}
\]
is strictly positive.
\end{lemma}

\begin{proof}
Cauchy--Schwarz, applied in the measure
\(t^{s-1}e^{-\lambda t^2-at}dt\), is strict and gives
\[
 J_{s+1}^2<J_sJ_{s+2}.
\]
Therefore
\[
 m_{s+1}-m_s
 =\frac{J_sJ_{s+2}-J_{s+1}^2}{J_sJ_{s+1}}
 =\frac{v_s}{m_s}>0.                                \tag{7.11}
\]
Integration by parts has no boundary term and yields
\[
 sJ_s=aJ_{s+1}+2\lambda J_{s+2}.
\]
Division by \(J_s\) proves
\[
 s=a m_s+2\lambda m_sm_{s+1},                       \tag{7.12}
\]
which is (7.9).  Subtract (7.12) at \(s\) from the same identity at
\(s+1\):
\[
 1=a(m_{s+1}-m_s)
 +2\lambda m_{s+1}(m_{s+2}-m_s).                   \tag{7.13}
\]
Both differences on the right are positive, proving the second assertion in
(7.8).

Finally, the proof of Theorem \ref{thm:v18v1-zero-boundary} uses only
positivity of the exponent \(s\), not integrality of \(2s\).  Repeating its
change of variables gives
\[
 c_s=2\frac{\lambda}{a^2}
 \operatorname{Cov}(X,X^2)>0                       \tag{7.14}
\]
for the density proportional to
\(x^{s-1}e^{-x-(\lambda/a^2)x^2}\).  This proves (7.10).
\end{proof}

\begin{lemma}[A Poisson-dominated mixing variance]
\label{lem:v18v2-mixing}
Assume \(H>0\), and define
\[
 r_k:=\frac{(k+1)p_{k+1}}{p_k},
 \qquad
 \theta:=\frac{H^2}{a}.                             \tag{7.15}
\]
Then \((r_k)_{k\geq0}\) is strictly decreasing and
\[
 0<r_k<\theta.                                      \tag{7.16}
\]
Consequently
\[
 \operatorname{Var}(K)\leq\mathbb EK<\theta.       \tag{7.17}
\]
Furthermore,
\[
 a^2\operatorname{Var}(m_{\alpha+K})
 \leq\operatorname{Var}(K).                        \tag{7.18}
\]
\end{lemma}

\begin{proof}
The normalizing factor in (7.4) cancels in the ratio.  With
\(s=\alpha+k\),
\[
 r_k=H^2\frac{m_s}{s}
 =\frac{H^2}{a+2\lambda m_{s+1}}.                  \tag{7.19}
\]
Lemma \ref{lem:v18v2-moment-ratio} shows that the denominator increases
strictly with \(k\), and it is larger than \(a\).  This proves
(7.16).

For any function \(f\) for which the sums converge, shifting the summation
index gives
\[
 \mathbb E[Kf(K)]=\mathbb E[r_Kf(K+1)].             \tag{7.20}
\]
Taking \(f=1\) gives \(\mathbb EK=\mathbb Er_K<\theta\).  Taking
\(f(K)=K\) and using the preceding identity gives
\[
\begin{split}
 \operatorname{Var}(K)
 &=\mathbb E[r_K(K+1)]-(\mathbb EK)^2\\
 &=\mathbb EK+\operatorname{Cov}(r_K,K).
                                                               \tag{7.21}
\end{split}
\]
Because \(r_k\) decreases while \(k\) increases,
\[
 \operatorname{Cov}(r_K,K)
 =\frac12\mathbb E[(r_K-r_{K'})(K-K')]\leq0         \tag{7.22}
\]
for an independent copy \(K'\).  This proves (7.17).

Equation (7.8) makes \(k\mapsto m_{\alpha+k}\) increasing and
\(1/a\)-Lipschitz.  Hence, for independent \(K,K'\),
\[
\begin{split}
 \operatorname{Var}(m_{\alpha+K})
 &=\frac12\mathbb E
 \left(m_{\alpha+K}-m_{\alpha+K'}\right)^2\\
 &\leq\frac1{2a^2}\mathbb E(K-K')^2
 =\frac1{a^2}\operatorname{Var}(K),                \tag{7.23}
\end{split}
\]
which is (7.18).
\end{proof}

\begin{theorem}[Global strict conformal convexity of a quartic leaf]
\label{thm:v18v2-global-leaf}
For every \(N\geq1\), \(\lambda>0\), \(\chi>0\),
\(d\in\mathbb R\), and \(z\in\mathbb R^N\), the conditional leaf free
energy (5.2) satisfies
\[
 \partial_d^2L(d,z)>0.                              \tag{7.24}
\]
Equivalently, the global leaf inequality left open in (6.1) holds strictly:
\[
 a^2\operatorname{Var}_{a,z}(|u|^2)
 <\frac{H^2}{a}+a\mathbb E_{a,z}|u|^2.              \tag{7.25}
\]
\end{theorem}

\begin{proof}
For \(H=0\), this is Theorem
\ref{thm:v18v1-zero-boundary}.  Suppose \(H>0\).  Apply the law of total
variance to the mixture in Lemma \ref{lem:v18v2-angular}:
\[
 \operatorname{Var}(t)
 =\mathbb E v_{\alpha+K}
 +\operatorname{Var}(m_{\alpha+K}).                 \tag{7.26}
\]
Using Proposition \ref{prop:v18v1-leaf-ward}, (7.10), and (7.26), we obtain
\[
\begin{split}
 \partial_d^2L(d,z)
 &=\theta+
 \mathbb E\left(a m_{\alpha+K}-a^2v_{\alpha+K}\right)
 -a^2\operatorname{Var}(m_{\alpha+K})\\
 &=\theta+\mathbb Ec_{\alpha+K}
 -a^2\operatorname{Var}(m_{\alpha+K}).             \tag{7.27}
\end{split}
\]
Lemmas \ref{lem:v18v2-moment-ratio} and
\ref{lem:v18v2-mixing} now give
\[
 \partial_d^2L(d,z)
 \geq \theta-\operatorname{Var}(K)
       +\mathbb Ec_{\alpha+K}>0.                   \tag{7.28}
\]
This proves both formulations.
\end{proof}

The preceding proof is qualitative at large \(|z|\).  The following direct
estimate supplies the coercive tail needed for uniformity.

\begin{lemma}[Brascamp--Lieb large-boundary bound]
\label{lem:v18v2-large-field}
Let
\[
 c_0=3-2\sqrt2.
\]
Then
\[
 \partial_d^2L(d,z)
 \geq \frac{H^2}{a}-\frac{c_0a^2}{6\lambda}.        \tag{7.29}
\]
\end{lemma}

\begin{proof}
The tilted potential in (5.4) is
\[
 V(u)=\lambda|u|^4+a|u|^2-2\chi u\cdot z,
\]
and
\[
 \nabla^2V=(4\lambda|u|^2+2a)I+8\lambda uu^{\mathsf T}.
                                                               \tag{7.30}
\]
The Brascamp--Lieb variance inequality, applied to
\(t(u)=|u|^2\), gives
\[
 \operatorname{Var}(t)
 \leq\mathbb E\left[
 (2u)^{\mathsf T}(\nabla^2V)^{-1}(2u)\right]
 =\mathbb E\frac{2t}{a+6\lambda t}.                 \tag{7.31}
\]
Insert (7.31) into (5.6):
\[
 \partial_d^2L(d,z)
 \geq\frac{H^2}{a}
 +\mathbb E\left[at-\frac{2a^2t}{a+6\lambda t}\right].
                                                               \tag{7.32}
\]
With \(x=6\lambda t/a\), the negative part of the bracket is bounded by
\[
 \frac{a^2}{6\lambda}
 \max_{0\leq x\leq1}\frac{x(1-x)}{1+x}
 =\frac{c_0a^2}{6\lambda},                          \tag{7.33}
\]
because the maximum occurs at \(x=\sqrt2-1\).  Equations
(7.32)--(7.33) prove (7.29).
\end{proof}

\begin{corollary}[Uniform leaf floor on compact regulator sectors]
\label{cor:v18v2-uniform-leaf}
Under the compact parameter restrictions (5.17), there is a constant
\(c_{\mathrm{leaf}}>0\), depending only on those restrictions and on
\(N\), such that
\[
 \inf_{z\in\mathbb R^N}
 \partial_d^2L(d,z)\geq c_{\mathrm{leaf}}.           \tag{7.34}
\]
\end{corollary}

\begin{proof}
For sufficiently large \(|z|\), uniformly in the compact parameter set,
the first term on the right of (7.29) is at least twice the second; it then
gives a positive uniform lower bound.  On the complementary compact set of
\((\lambda,\chi,d,z)\), Theorem
\ref{thm:v18v2-global-leaf} gives strict positivity and the domination
argument from Lemma \ref{lem:v18v1-differentiability} gives continuity.
The minimum on that compact set is therefore positive.  Take the smaller of
the two lower bounds.
\end{proof}

\section{V2 consequence and next gate}
\label{sec:v18v2-next}

The terminal nonzero-boundary obstruction from V1 is now closed.  The proof
does not use a perturbative smallness assumption on \(\chi\), \(d\), or
\(|z|\).  It uses only positive quartic confinement and a finite-dimensional
leaf.

For a nonterminal vertex, however, the integrand in (4.2) contains the
additional radial factor
\[
 R_x(u;d)=\prod_{c:\,p(c)=x}M_{c\to x}(u;d).        \tag{8.1}
\]
The integration-by-parts recurrence (7.12) changes when \(R_x\) is present,
and descendant edge parameters also differentiate \(R_x\).  Therefore
Theorem \ref{thm:v18v2-global-leaf} cannot simply be iterated without a
closure theorem for the class of radial messages.  The next task is to find
a property of \(R_x\) that preserves the decreasing ratio in (7.19), the
increment bound in (7.8), and the mixed conformal Hessian after (4.2).

\begin{assessment}[V2 acquisition and boundary]
\label{ass:v18v2}
V2 proves strict global conformal convexity of every massless quartic leaf
kernel and a uniform floor on compact positive regulator sectors.  The proof
is exact: a positive angular mixture splits the radial variance into
within-component and between-component parts; strict zero-tilt curvature
controls the former, and a Poisson-dominated discrete variance controls the
latter.  This closes inequality (6.1).  Preservation of the full Hessian
through a branching message recursion, uniform cofinal control, and the
coupled Standard Model--Einstein continuum remain open.
\end{assessment}

\begin{record}[V2 append record]
\label{rec:v18v2}
V2 is appended to the active eighteenth-cycle V1 whose installed SHA--256
was
\texttt{DD5A1EE6868336E35D2BB7A7DAA1A7CA899C263DBB84764F6B6D795EC3AEB9E4}.
After installation only V2 is the active numeric SM note; the frozen
\texttt{V100\_f17} archive is unchanged.
\end{record}

\section{V3: a radial message class closed by tree elimination}
\label{sec:v18v3}

V2 treated the pure quartic leaf.  An internal vertex sees instead its
quartic onsite factor multiplied by all messages arriving from its children.
This section identifies a class of radial effective potentials that is closed
under that operation.  The result controls the shape in the boundary radius
and the curvature in the newly eliminated edge.  A final separator shows why
this coordinatewise information is not yet a proof of positivity of the full
edge Hessian.

Let \(\Psi:[0,\infty)\to\mathbb R\) be locally absolutely continuous,
convex, and nondecreasing.  Additive constants in \(\Psi\) are immaterial.
For \(\chi>0\), define
\[
\begin{split}
 M_\Psi(d,z)&=
 \int_{\mathbb R^N}
 \exp\left[-\Psi(|u|^2)
 -\chi\left|e^{-d/2}u-e^{d/2}z\right|^2\right]du,
                                                               \tag{9.1}\\
 L_\Psi(d,z)&=-\log M_\Psi(d,z).                   \tag{9.2}
\end{split}
\]
The Gaussian term makes the integral finite even when \(\Psi\) is bounded.
By rotation invariance there is a function \(\ell_\Psi\) such that
\[
 L_\Psi(d,z)=\ell_\Psi(d,|z|^2).                   \tag{9.3}
\]

For \(a=\chi e^{-d}\) and \(s>0\), put
\[
\begin{split}
 J_s^\Psi(a)&=
 \int_0^\infty t^{s-1}e^{-\Psi(t)-at}\,dt,\\
 m_s^\Psi&=\frac{J_{s+1}^\Psi}{J_s^\Psi},
 \qquad
 v_s^\Psi=\frac{J_{s+2}^\Psi}{J_s^\Psi}
 -(m_s^\Psi)^2.                                    \tag{9.4}
\end{split}
\]
Write \(\mathbb E_s^\Psi\) for expectation under the normalized density in
(9.4), and set
\[
 G_s^\Psi=\mathbb E_s^\Psi\Psi'(t).                \tag{9.5}
\]

\begin{lemma}[General radial recurrence]
\label{lem:v18v3-recurrence}
For every \(s>0\),
\[
 s=m_s^\Psi\left(a+G_{s+1}^\Psi\right).            \tag{9.6}
\]
The sequences \(s\mapsto m_s^\Psi\) and
\(s\mapsto G_s^\Psi\) are nondecreasing, with the first strictly
increasing.  Moreover,
\[
 0<m_{s+1}^\Psi-m_s^\Psi\leq\frac1a,               \tag{9.7}
\]
and
\[
 a m_s^\Psi-a^2v_s^\Psi\geq0.                     \tag{9.8}
\]
If \(\Psi\) is not constant, both inequalities relevant to (9.8) are
strict, and the left side of (9.8) is positive.
\end{lemma}

\begin{proof}
Convexity makes \(\Psi'\) nondecreasing almost everywhere, and
monotonicity makes it nonnegative almost everywhere.  Integration by parts
gives
\[
 sJ_s^\Psi
 =aJ_{s+1}^\Psi+
 \int_0^\infty t^s\Psi'(t)e^{-\Psi(t)-at}\,dt,
                                                               \tag{9.9}
\]
which is (9.6).  Strict Cauchy--Schwarz gives
\[
 m_{s+1}^\Psi-m_s^\Psi
 =\frac{v_s^\Psi}{m_s^\Psi}>0.                    \tag{9.10}
\]
Also
\[
 G_{s+1}^\Psi-G_s^\Psi
 =\frac{
 \operatorname{Cov}_s^\Psi(t,\Psi'(t))}{m_s^\Psi}
 \geq0,                                             \tag{9.11}
\]
because two nondecreasing functions have nonnegative covariance.

Subtract (9.6) at \(s\) from the same identity at \(s+1\):
\[
\begin{split}
 1={}&a(m_{s+1}^\Psi-m_s^\Psi)\\
 &+m_{s+1}^\Psi G_{s+2}^\Psi
   -m_s^\Psi G_{s+1}^\Psi.                         \tag{9.12}
\end{split}
\]
The second line is nonnegative by monotonicity of both factors, so (9.7)
follows.  Combining (9.10) and (9.7) gives
\[
 a m_s^\Psi-a^2v_s^\Psi
 =a m_s^\Psi
 \left[1-a(m_{s+1}^\Psi-m_s^\Psi)\right]\geq0,     \tag{9.13}
\]
which is (9.8).  If \(\Psi\) is nonconstant, its nonnegative derivative is
positive on a set of positive measure.  Hence \(G_{s+1}^\Psi>0\); strict
increase of \(m_s^\Psi\) then makes the second line of (9.12) positive.
This makes (9.7) strict at its upper endpoint and (9.8) positive.
\end{proof}

\begin{theorem}[Closed radial kernel theorem]
\label{thm:v18v3-closed-kernel}
Let \(\Psi\) be convex and nondecreasing.  Then:
\begin{enumerate}[leftmargin=2.2em]
\item for every fixed boundary field \(z\),
      \[
      \partial_d^2L_\Psi(d,z)\geq0;                \tag{9.14}
      \]
\item for every fixed \(d\), the map
      \[
      q\longmapsto\ell_\Psi(d,q),
      \qquad q=|z|^2,                               \tag{9.15}
      \]
      is nondecreasing and convex on \([0,\infty)\).
\end{enumerate}
If \(\Psi\) is not constant, (9.14) is strict for every \(z\), and
(9.15) is strictly increasing.
\end{theorem}

\begin{proof}
Let \(\alpha=N/2\), \(H=\chi|z|\), and use the angular decomposition of
Lemma \ref{lem:v18v2-angular} with \(J_s\) replaced by
\(J_s^\Psi\).  Thus
\[
 p_k\ \propto\
 \frac{H^{2k}}{k!(\alpha)_k}J_{\alpha+k}^\Psi(a).  \tag{9.16}
\]
For \(H>0\), its consecutive ratio is
\[
 r_k:=\frac{(k+1)p_{k+1}}{p_k}
 =H^2\frac{m_{\alpha+k}^\Psi}{\alpha+k}
 =\frac{H^2}{a+G_{\alpha+k+1}^\Psi}.               \tag{9.17}
\]
By Lemma \ref{lem:v18v3-recurrence}, \(r_k\) is nonincreasing and
\[
 r_k\leq\theta:=\frac{H^2}{a}.                     \tag{9.18}
\]
The summation-by-parts proof in Lemma \ref{lem:v18v2-mixing} therefore gives
\[
 \operatorname{Var}(K)\leq\mathbb EK\leq\theta.   \tag{9.19}
\]
The increment bound (9.7) also gives
\[
 a^2\operatorname{Var}(m_{\alpha+K}^\Psi)
 \leq\operatorname{Var}(K).                        \tag{9.20}
\]

The conditional conformal Ward formula is unchanged from (5.6), with the
radial law now defined by \(\Psi\).  Splitting its variance over \(K\) yields
\[
\begin{split}
 \partial_d^2L_\Psi(d,z)
 ={}&\theta+
 \mathbb E\left[
 a m_{\alpha+K}^\Psi-a^2v_{\alpha+K}^\Psi\right]\\
 &-a^2\operatorname{Var}(m_{\alpha+K}^\Psi)\geq0,  \tag{9.21}
\end{split}
\]
by (9.8)--(9.20).  When \(H=0\), the same conclusion follows directly from
(9.8).  If \(\Psi\) is nonconstant, the expected component term in (9.21)
is positive, so (9.14) is strict.

It remains to prove the boundary-radius assertion.  Set \(q=|z|^2\), so
\(H^2=\chi^2q\), and recall \(b=\chi e^d=\chi^2/a\).  Apart from a constant
independent of \(q\),
\[
 \ell_\Psi(d,q)=bq-log
 \sum_{k=0}^\infty
 \frac{(\chi^2q)^k}{k!(\alpha)_k}
 J_{\alpha+k}^\Psi(a).                             \tag{9.22}
\]
For \(q>0\), power-series differentiation gives
\[
 \partial_q\ell_\Psi(d,q)
 =b-\frac{\mathbb EK}{q}
 =\frac{\theta-\mathbb EK}{q}\geq0,               \tag{9.23}
\]
and
\[
 \partial_q^2\ell_\Psi(d,q)
 =\frac{\mathbb EK-\operatorname{Var}(K)}{q^2}
 \geq0.                                             \tag{9.24}
\]
The limits at \(q=0\) exist by the convergent positive power series, so the
properties extend to \([0,\infty)\).  If \(\Psi\) is nonconstant, then
\(G_{\alpha+k+1}^\Psi>0\), whence (9.17) is strictly below \(\theta\) and
(9.23) is positive.
\end{proof}

\begin{corollary}[Radial shape is preserved through an entire tree]
\label{cor:v18v3-tree-closure}
Fix all conformal edge differences of a finite rooted tree.  For each vertex
\(x\ne o\), write
\[
 -\log M_{x\to p(x)}(z;d)=\ell_x(|z|^2;d).          \tag{9.25}
\]
Then \(q\mapsto\ell_x(q;d)\) is convex and strictly increasing.  Before
eliminating \(x\), its effective radial onsite potential is
\[
 \Psi_x(t;d)=\lambda_xt^2+
 \sum_{c:\,p(c)=x}\ell_c(t;d),                     \tag{9.26}
\]
which is convex, strictly increasing away from zero, and nonconstant.
For a fixed parent boundary field and fixed descendant parameters, the
message has strictly positive curvature in its connecting conformal
difference \(d_{e_x}\).
\end{corollary}

\begin{proof}
At a leaf, \(\Psi_x(t)=\lambda_xt^2\) is convex and nondecreasing.  Theorem
\ref{thm:v18v3-closed-kernel} proves the assertion for its outgoing message.
If it holds for all children of \(x\), then (9.26) is a sum of convex
nondecreasing functions and the strictly convex quartic term.  Apply the
theorem again.  Induction from the leaves to the root proves every claim.
\end{proof}

The corollary is a genuine closure theorem, but it is not a full Hessian
theorem.  The following exactly solvable member of the class prevents that
overstatement.

\begin{proposition}[Joint-convexity separator]
\label{prop:v18v3-separator}
Take \(\Psi(t)=ct\) with \(c>0\).  Then
\[
 L_\Psi(d,z)
 =\frac N2\log(a+c)+
 \frac{bc}{a+c}|z|^2+C_N,                           \tag{9.27}
\]
where \(a=\chi e^{-d}\), \(b=\chi e^d\), and \(C_N\) is independent of
\(d,z\).  Thus \(L_\Psi\) is separately convex in \(d\) and in
\(q=|z|^2\), as Theorem \ref{thm:v18v3-closed-kernel} asserts, but it is not
jointly convex in \((d,q)\).
\end{proposition}

\begin{proof}
Complete the square in (9.1):
\[
\begin{split}
 M_\Psi(d,z)
 &=e^{-b|z|^2}
 \int_{\mathbb R^N}
 e^{-(a+c)|u|^2+2\chi u\cdot z}\,du\\
 &=\left(\frac\pi{a+c}\right)^{N/2}
 \exp\left[-b|z|^2+\frac{\chi^2}{a+c}|z|^2\right].
                                                               \tag{9.28}
\end{split}
\]
Since \(ab=\chi^2\), taking the negative logarithm gives (9.27).  Write the
coefficient of \(q\) as
\[
 g(d)=\frac{bc}{a+c}=\frac{\chi^2c}{a(a+c)}.
\]
It has \(g'(d)>0\).  The Hessian of
\(A(d)+g(d)q\) in \((d,q)\) has determinant
\[
 \det\begin{pmatrix}
 A''(d)+g''(d)q&g'(d)\\
 g'(d)&0
 \end{pmatrix}
 =-[g'(d)]^2<0.                                     \tag{9.29}
\]
Hence joint convexity fails.
\end{proof}

\section{V3 consequence: the mixed-Hessian gate}
\label{sec:v18v3-next}

The radial class (9.26) is stable under every leaf elimination, so loss of
radial shape is not the obstruction.  Proposition
\ref{prop:v18v3-separator} shows that separate curvature in an edge parameter
and a boundary radius cannot be combined through a generic joint-convexity
or Prekopa argument.  What remains is a matrix statement that retains the
mixed derivatives between an outgoing message and all descendant edge
parameters.

The next honest object is the augmented message Hessian
\[
 \mathcal H_x=
 \nabla^2_{(d_{{\mathcal T}_x},q)}
 \ell_x(q;d_{{\mathcal T}_x}),                      \tag{10.1}
\]
together with the exact Schur or covariance update produced by (4.2).
The separator proves that \(\mathcal H_x\succeq0\) is too strong as an
induction invariant.  The next iteration must instead search for a weighted
matrix inequality whose Schur complement controls only the physical edge
block after the parent field is integrated.

\begin{assessment}[V3 acquisition and boundary]
\label{ass:v18v3}
V3 proves that convex nondecreasing radial effective potentials form a class
closed by the tree message transform.  Every tree message remains convex and
increasing in the parent squared radius, and it has positive curvature in its
own connecting conformal edge while the boundary is held fixed.  An exact
quadratic member of the class has an indefinite joint \((d,|z|^2)\) Hessian,
so coordinatewise convexity cannot be promoted to the full tree Hessian by a
generic joint-convexity argument.  The mixed-edge Schur update, uniform
cutoff control, and the coupled continuum remain open.
\end{assessment}

\begin{record}[V3 append record]
\label{rec:v18v3}
V3 is appended to the validated V2 whose installed SHA--256 was
\texttt{F6B24C2E9E18C5ACBCD23648E826831FDE8985BD4AED16CEEA893BFE68296D9A}.
After installation only V3 is the active numeric SM note; the frozen
\texttt{V100\_f17} archive is unchanged.
\end{record}

\section{V4: the exact mixed-Hessian update at a terminal elimination}
\label{sec:v18v4}

V3 showed that the negative logarithm of every tree message is convex and
increasing in the squared boundary radius.  This section first extracts an
explicit boundary-field stiffness from that proof.  It then disintegrates
the full conformal Hessian when one or several terminal messages are attached
to the same parent.  A Brascamp--Lieb estimate converts the remaining score
covariance into one rank-one radial correction.  The result is a precise
Schur gate rather than an assertion that coordinatewise convexity is enough.

Retain the generalized radial kernel of (9.1)--(9.5).  Write
\[
 \alpha=\frac N2,
 \qquad a=\chi e^{-d},
 \qquad b=\chi e^d,
 \qquad q=|z|^2.                                    \tag{11.1}
\]

\begin{proposition}[Quantitative boundary-field stiffness]
\label{prop:v18v4-stiffness}
Suppose \(\Psi\) is convex, nondecreasing, and nonconstant.  Then
\[
 \partial_q\ell_\Psi(d,q)
 =b\,\mathbb E\left[
 \frac{G_{\alpha+K+1}^\Psi}
 {a+G_{\alpha+K+1}^\Psi}\right],                  \tag{11.2}
\]
where \(K\) has the angular mixing law (9.16).  Consequently
\[
 0<\gamma_\Psi(d):=
 b\frac{G_{\alpha+1}^\Psi}{a+G_{\alpha+1}^\Psi}
 \leq\partial_q\ell_\Psi(d,q)<b                  \tag{11.3}
\]
for every \(q\geq0\), and
\[
 \nabla_z^2L_\Psi(d,z)\succeq2\gamma_\Psi(d)I_N.
                                                               \tag{11.4}
\]
\end{proposition}

\begin{proof}
For \(q>0\), equation (9.23) and the identity
\(\mathbb EK=\mathbb Er_K\) give
\[
\begin{split}
 \partial_q\ell_\Psi(d,q)
 &=\frac1q\mathbb E\left[
 \frac{H^2}{a}-
 \frac{H^2}{a+G_{\alpha+K+1}^\Psi}\right]\\
 &=\frac{\chi^2}{a}\mathbb E\left[
 \frac{G_{\alpha+K+1}^\Psi}
 {a+G_{\alpha+K+1}^\Psi}\right],                 \tag{11.5}
\end{split}
\]
which is (11.2) because \(\chi^2/a=b\).  The sequence
\(G_s^\Psi\) is nondecreasing by (9.11), and it is positive when \(\Psi\)
is nonconstant.  This gives (11.3); the formula extends to \(q=0\) by the
positive power series.

For a twice differentiable radial function \(f(|z|^2)\),
\[
 \nabla_z^2 f=2f'(q)I_N+4f''(q)zz^{\mathsf T}.       \tag{11.6}
\]
Theorem \ref{thm:v18v3-closed-kernel} gives
\(\partial_q^2\ell_\Psi\geq0\).  Insert (11.3) into (11.6) to obtain
(11.4).  The locally absolutely continuous case follows by smooth
approximation; tree applications have smooth positive messages.
\end{proof}

\subsection{One terminal parameter}

Let \(\Phi:[0,\infty)\to\mathbb R\) be a twice differentiable convex
nondecreasing exterior radial potential such that the integrals below are
finite.  It may contain the parent quartic term and messages not depending on
the terminal edge parameter \(d\).  Define
\[
\begin{split}
 A(d,q)&=\Phi(q)+\ell_\Psi(d,q),                    \tag{11.7}\\
 \widehat Z(d)&=\int_{\mathbb R^N}e^{-A(d,|z|^2)}\,dz,
 \qquad
 \widehat F(d)=-\log\widehat Z(d).                 \tag{11.8}
\end{split}
\]
Let \(\widehat{\mathbb E}_d\) denote expectation in the normalized law in
(11.8).

\begin{proposition}[Exact curvature disintegration]
\label{prop:v18v4-disintegration}
Under the preceding assumptions,
\[
 \widehat F''(d)=
 \widehat{\mathbb E}_d\,\partial_d^2\ell_\Psi(d,q)
 -\operatorname{Var}_{\widehat{\mathbb P}_d}
  \left(\partial_d\ell_\Psi(d,q)\right).           \tag{11.9}
\]
Thus the positive conditional curvature from V2--V3 is reduced by an
exterior score variance after the parent field is integrated.
\end{proposition}

\begin{proof}
Differentiate (11.8) under the integral.  The first derivative is
\[
 \widehat F'(d)=
 \widehat{\mathbb E}_d\,\partial_d\ell_\Psi(d,q).
\]
Differentiating once more gives the expectation of the second derivative
minus the variance of the first derivative, which is (11.9).
\end{proof}

\begin{theorem}[Radial Brascamp--Lieb Schur bound]
\label{thm:v18v4-radial-schur}
Suppose \(\Psi\) is nonconstant.  Put
\[
 R_A(d,q)=2\partial_qA(d,q)+4q\partial_q^2A(d,q).
                                                               \tag{11.10}
\]
Then \(R_A(d,q)>0\), and
\[
 \widehat F''(d)\geq
 \widehat{\mathbb E}_d\left[
 \partial_d^2\ell_\Psi(d,q)
 -\frac{4q\left(\partial_d\partial_q
                  \ell_\Psi(d,q)\right)^2}
 {R_A(d,q)}\right].                                \tag{11.11}
\]
In particular, the pointwise Schur inequality
\[
 R_A\,\partial_d^2\ell_\Psi
 \geq4q(\partial_d\partial_q\ell_\Psi)^2          \tag{11.12}
\]
is sufficient for \(\widehat F''(d)\geq0\).
\end{theorem}

\begin{proof}
The potential \(z\mapsto A(d,|z|^2)\) is convex.  Its Hessian has tangential
eigenvalue \(2A_q\) and radial eigenvalue \(R_A\).  Proposition
\ref{prop:v18v4-stiffness} gives \(A_q\geq\gamma_\Psi(d)>0\), so the Hessian
is positive definite.  The Brascamp--Lieb variance inequality applied to
\(f(z)=\partial_d\ell_\Psi(d,|z|^2)\) gives
\[
 \operatorname{Var}(f)
 \leq\widehat{\mathbb E}_d
 \left[(\nabla_zf)^{\mathsf T}
 (\nabla_z^2A)^{-1}\nabla_zf\right].               \tag{11.13}
\]
Since
\[
 \nabla_zf=2(\partial_d\partial_q\ell_\Psi)z,
\]
this vector is radial and the integrand in (11.13) is
\[
 \frac{4q(\partial_d\partial_q\ell_\Psi)^2}{R_A}.
                                                               \tag{11.14}
\]
Combine (11.14) with (11.9) to prove (11.11).  Inequality (11.12) makes the
integrand in (11.11) nonnegative.
\end{proof}

\subsection{Several terminal edges at one parent}

Let terminal messages \(\ell_i(d_i,q)\), \(1\leq i\leq m\), be generated by
nonconstant convex nondecreasing radial inputs as in V3.  Let
\[
 A(d,q)=\Phi(q)+\sum_{i=1}^m\ell_i(d_i,q),          \tag{11.15}
\]
and define \(\widehat F(d)\) by integrating
\(e^{-A(d,|z|^2)}\) over the common parent field.  Set
\[
 D(q)=\operatorname{diag}
 \left(\partial_{d_1}^2\ell_1,\ldots,
       \partial_{d_m}^2\ell_m\right),
 \qquad
 c(q)=
 \begin{pmatrix}
 \partial_{d_1}\partial_q\ell_1\\
 \vdots\\
 \partial_{d_m}\partial_q\ell_m
 \end{pmatrix}.                                    \tag{11.16}
\]

\begin{theorem}[Rank-one star Schur gate]
\label{thm:v18v4-star}
With \(R_A=2A_q+4qA_{qq}>0\), the full terminal-edge Hessian obeys
\[
 \nabla_d^2\widehat F(d)
 \succeq
 \widehat{\mathbb E}_d\left[
 D(q)-\frac{4q}{R_A(d,q)}c(q)c(q)^{\mathsf T}
 \right].                                          \tag{11.17}
\]
If every diagonal entry of \(D(q)\) is positive, the pointwise matrix on the
right is positive semidefinite exactly when
\[
 \frac{4q}{R_A(d,q)}
 \sum_{i=1}^m
 \frac{(\partial_{d_i}\partial_q\ell_i)^2}
 {\partial_{d_i}^2\ell_i}
 \leq1.                                             \tag{11.18}
\]
Hence (11.18) is an explicit sufficient condition for positivity of the
star edge Hessian.
\end{theorem}

\begin{proof}
For \(v\in\mathbb R^m\), differentiation as in Proposition
\ref{prop:v18v4-disintegration} gives
\[
 v^{\mathsf T}\nabla_d^2\widehat Fv
 =\widehat{\mathbb E}_d
 \sum_i v_i^2\partial_{d_i}^2\ell_i
 -\operatorname{Var}_{\widehat{\mathbb P}_d}
 \left(\sum_i v_i\partial_{d_i}\ell_i\right).     \tag{11.19}
\]
The gradient in the variance term is
\[
 2\left(\sum_i v_i\partial_{d_i}\partial_q\ell_i\right)z.
\]
Applying the radial Brascamp--Lieb calculation (11.13)--(11.14) gives
\[
 v^{\mathsf T}\nabla_d^2\widehat Fv
 \geq\widehat{\mathbb E}_d
 v^{\mathsf T}\left(D-\frac{4q}{R_A}cc^{\mathsf T}\right)v,
\]
which is (11.17).  For positive diagonal \(D\), the standard rank-one Schur
criterion
\[
 D-\rho cc^{\mathsf T}\succeq0
 \quad\Longleftrightarrow\quad
 \rho\,c^{\mathsf T}D^{-1}c\leq1
\]
gives (11.18).
\end{proof}

\section{V4 consequence and next quantitative target}
\label{sec:v18v4-next}

The mixed-Hessian problem has now been reduced to a scalar radial budget at
each branching vertex.  In (11.18),
\(\partial_{d_i}^2\ell_i\) is the conditional edge curvature acquired in
V2--V3, while \(R_A\) contains the parent quartic stiffness and every child
boundary stiffness.  The numerator measures how strongly a conformal edge
score changes with the parent radius.

The next iteration should estimate the dimensionless transfer ratio
\[
 \mathfrak r_i(q)=
 \frac{4q(\partial_{d_i}\partial_q\ell_i)^2}
 {R_A\,\partial_{d_i}^2\ell_i}                     \tag{12.1}
\]
for the pure quartic leaf and determine whether
\(\sum_i\mathfrak r_i\leq1\) follows from the quartic parent term, requires
a degree or coupling restriction, or fails in an explicit parameter regime.
That decision is upstream of any induction on tree depth.

\begin{assessment}[V4 acquisition and boundary]
\label{ass:v18v4}
V4 proves a quantitative boundary stiffness for every message in the closed
radial class and gives the exact conditional-curvature minus score-variance
formula after parent integration.  Brascamp--Lieb reduction turns the latter
into the rank-one Schur gate (11.18).  This is sufficient for a star and is
the local matrix budget needed by a tree induction.  No bound proving
(11.18) uniformly has yet been established; the cofinal regulator limit and
the coupled Standard Model--Einstein continuum remain open.
\end{assessment}

\begin{record}[V4 append record]
\label{rec:v18v4}
V4 is appended to the validated V3 whose installed SHA--256 was
\texttt{CB6B6DC43B0623CCD64C0CEF0D37722F5B9FA44A43C140A8C4FECACA5CD5FA21}.
After installation only V4 is the active numeric SM note; the frozen
\texttt{V100\_f17} archive is unchanged.
\end{record}

\section{V5: a finite transfer norm and a positive star sector}
\label{sec:v18v5}

The V4 Schur gate becomes useful only if the mixed radial derivative can be
compared uniformly with the conditional edge curvature.  This section proves
that the comparison norm is finite for every pure quartic leaf, uniformly on
compact positive parameter sets.  A sufficiently strong quartic coefficient
at the common parent then closes the complete massless edge Hessian of a
finite star.

For one pure quartic leaf, retain
\[
 \ell(d,q)=L(d,z),\qquad q=|z|^2,qquad
 a=\chi e^{-d},\qquad b=\chi e^d,                  \tag{13.1}
\]
and let \(K\) be the angular mixing variable of (7.4).  Conditional on
\(K=k\), the radial variable \(t=|u|^2\) has law (7.5).

\begin{lemma}[Exact mixed derivative]
\label{lem:v18v5-mixed}
For \(q>0\),
\[
 \partial_d\partial_q\ell(d,q)
 =b-\frac aq\operatorname{Cov}(t,K).                \tag{13.2}
\]
The right side has a finite continuous extension to \(q=0\).
\end{lemma}

\begin{proof}
Equation (5.5), after conditioning on the fixed boundary field, reads
\[
 \partial_d\ell(d,q)=bq-a\mathbb Et.               \tag{13.3}
\]
Only the mixing weights \(p_k\) depend on \(q\); the component laws (7.5)
do not.  Since \(p_k\) is a power-series family in \(q\),
\[
 q\partial_q\mathbb Et
 =\operatorname{Cov}(t,K).                          \tag{13.4}
\]
Differentiate (13.3) and use (13.4) to obtain (13.2).  The positive angular
series (7.3) is analytic at \(q=0\), as are its parameter derivatives, so
(13.2) extends continuously.
\end{proof}

\begin{lemma}[Large-boundary transfer decay]
\label{lem:v18v5-tail}
Let
\[
 D(d,q)=\partial_d^2\ell(d,q),
 \qquad C(d,q)=\partial_d\partial_q\ell(d,q).       \tag{13.5}
\]
Then, as \(q\to\infty\) with \(\lambda,\chi,d\) fixed,
\[
 D(d,q)\geq bq-\frac{(3-2\sqrt2)a^2}{6\lambda},    \tag{13.6}
\]
and
\[
 |C(d,q)-b|
 \leq a\sqrt{\frac{b}{3\lambda q}}.                \tag{13.7}
\]
Consequently
\[
 \frac{C(d,q)^2}{D(d,q)}=O(q^{-1}).                \tag{13.8}
\]
The constants in these statements can be chosen uniformly when
\((\lambda,\chi,d)\) ranges over a compact set with
\(\lambda,\chi>0\).
\end{lemma}

\begin{proof}
Inequality (13.6) is the large-boundary estimate (7.29), because
\(H^2/a=bq\).  The Brascamp--Lieb calculation (7.31) also gives
\[
 \operatorname{Var}(t)
 \leq\mathbb E\frac{2t}{a+6\lambda t}
 \leq\frac1{3\lambda}.                             \tag{13.9}
\]
Lemma \ref{lem:v18v2-mixing} gives
\[
 \operatorname{Var}(K)\leq\mathbb EK<bq.           \tag{13.10}
\]
Cauchy--Schwarz and (13.9)--(13.10) imply
\[
 |\operatorname{Cov}(t,K)|
 \leq\sqrt{\frac{bq}{3\lambda}}.                  \tag{13.11}
\]
Insert (13.11) into (13.2) to obtain (13.7).  For large enough \(q\),
(13.6) is at least \(bq/2\), while (13.7) keeps \(C\) bounded.  This proves
(13.8).  On compact positive parameter sets, \(a\), \(b\), and
\(1/\lambda\) have uniform bounds, so the same argument is uniform.
\end{proof}

\begin{theorem}[Finite leaf transfer norm]
\label{thm:v18v5-transfer}
For every \(N\geq1\) and every \(\lambda,\chi>0\), \(d\in\mathbb R\),
\[
 \mathcal M(\lambda,\chi,d):=
 \sup_{q\geq0}
 \frac{\left(\partial_d\partial_q\ell(d,q)\right)^2}
 {\partial_d^2\ell(d,q)}<\infty.                  \tag{13.12}
\]
If \((\lambda,\chi,d)\) ranges over a compact set with positive lower
bounds on \(\lambda\) and \(\chi\), then the supremum in (13.12) has a
common finite upper bound.
\end{theorem}

\begin{proof}
The denominator in (13.12) is strictly positive by Theorem
\ref{thm:v18v2-global-leaf}; on compact positive parameter sets it has the
uniform lower bound of Corollary \ref{cor:v18v2-uniform-leaf}.  The numerator
is continuous at every finite \(q\), including \(q=0\), by Lemma
\ref{lem:v18v5-mixed} and dominated differentiation.  Thus the quotient is
bounded on compact \(q\)-intervals.  Lemma \ref{lem:v18v5-tail} shows that it
tends to zero at infinity, with a uniform tail estimate on compact parameter
sets.  Combining the compact and tail bounds proves both assertions.
\end{proof}

\subsection{A complete star Hessian}

Let a star have root \(o\), terminal vertices \(x_1,\ldots,x_m\), root
quartic coefficient \(\lambda_o>0\), and pure quartic leaf messages
\(\ell_i(d_i,q)\).  Gauge transports may be removed by Theorem
\ref{thm:v18v1-gauge-erasure}.  Its reduced root potential is
\[
 A(d,q)=\lambda_oq^2+\sum_{i=1}^m\ell_i(d_i,q),     \tag{13.13}
\]
and its full free energy is
\[
 F_\star(d)=-\log
 \int_{\mathbb R^N}e^{-A(d,|z|^2)}\,dz.            \tag{13.14}
\]
For each leaf set
\[
 \mathcal M_i=
 \sup_{q\geq0}
 \frac{(\partial_{d_i}\partial_q\ell_i)^2}
 {\partial_{d_i}^2\ell_i}.                         \tag{13.15}
\]

\begin{theorem}[Strong-parent positive star sector]
\label{thm:v18v5-star-sector}
If
\[
 \lambda_o\geq\frac13\sum_{i=1}^m\mathcal M_i,    \tag{13.16}
\]
then
\[
 \nabla_d^2F_\star(d)\succeq0.                    \tag{13.17}
\]
If the inequality in (13.16) is strict, then the Hessian in (13.17) is
positive definite.  The threshold is finite, and it is uniform for fixed
degree \(m\) when all leaf parameters range over a common compact positive
set.
\end{theorem}

\begin{proof}
For (13.13), the radial Hessian denominator from (11.10) is
\[
\begin{split}
 R_A
 &=12\lambda_oq+
 \sum_{i=1}^m
 \left(2\partial_q\ell_i+4q\partial_q^2\ell_i\right)\\
 &\geq12\lambda_oq.                                \tag{13.18}
\end{split}
\]
For \(q>0\), the left side of the star Schur condition (11.18) is bounded by
\[
\begin{split}
 \frac{4q}{R_A}
 \sum_i\frac{(\partial_{d_i}\partial_q\ell_i)^2}
 {\partial_{d_i}^2\ell_i}
 &\leq\frac1{3\lambda_o}sum_i\mathcal M_i\leq1.
                                                               \tag{13.19}
\end{split}
\]
At \(q=0\), the rank-one correction in (11.17) vanishes.  Theorem
\ref{thm:v18v4-star} now proves (13.17).  Under strict (13.16), the pointwise
Schur matrix is positive definite for every \(q>0\); the root law assigns
full measure to \(q>0\), so its expectation is positive definite.  Finiteness
and compact-parameter uniformity of the threshold follow from Theorem
\ref{thm:v18v5-transfer}.
\end{proof}

\begin{corollary}[Massless Ward-incidence sign on the star sector]
\label{cor:v18v5-star-wi}
Under (13.16), the exact massless Ward-incidence constant of the star
satisfies
\[
 K_{\mathrm{WI}}(d)=0.                             \tag{13.20}
\]
With the uniform mass-derivative hypothesis recorded in (2.7), this star
sector therefore has
\[
 K_{\mathrm{WI}}(a;d)=O(a^2)                       \tag{13.21}
\]
under the canonical regulator scaling.
\end{corollary}

\begin{proof}
A star is a tree, so its incidence map is surjective by Theorem
\ref{thm:v18v1-incidence-surjective}.  Equation (13.17) is therefore exactly
the massless incidence inequality.  Equations (2.7)--(2.8) give (13.21) when
their uniform derivative hypothesis is imposed.
\end{proof}

\section{V5 consequence and mandatory companion audit point}
\label{sec:v18v5-next}

Theorem \ref{thm:v18v5-star-sector} is the first result in the restarted
cycle that proves the complete massless conformal edge Hessian on a graph
with more than one edge.  Its sufficient condition scales with the star
degree through \(\sum_i\mathcal M_i\).  The canonical Standard Model
cofinal family has a fixed order-one quartic coefficient, so the theorem does
not yet give a degree-uniform lattice result.  The next direction must decide
whether the positive child stiffness terms discarded in (13.18) remove that
degree cost, or whether a bounded-degree and quantitative coupling condition
is necessary.

This is V5 of the eighteenth cycle.  Before V6 is chosen, the newest active
Yang--Mills and Navier--Stokes notes in the shared folder must be inspected as
required by the five-version audit rule.

\begin{assessment}[V5 acquisition and boundary]
\label{ass:v18v5}
V5 proves finiteness of the leaf mixed-transfer norm and uses it to close a
nonperturbative positive-Hessian star sector under the explicit strong-parent
condition (13.16).  This gives \(K_{\mathrm{WI}}=0\) on that sector and the
canonical \(O(a^2)\) conclusion when the previously isolated uniform mass
derivative is supplied.  The threshold has not been shown to hold for the
canonical order-one couplings on every regulator star; depth propagation,
cycles, the continuum state, and the coupled Einstein system remain open.
\end{assessment}

\begin{record}[V5 append record]
\label{rec:v18v5}
V5 is appended to the validated V4 whose installed SHA--256 was
\texttt{0570E7DA77A3CFFC55B39EDB1F7F4B17CA696680FDF18FDDE903D67873AE5B27}.
After installation only V5 is the active numeric SM note; the frozen
\texttt{V100\_f17} archive is unchanged.
\end{record}

\section{V6: companion audit and the stiffness-retaining star threshold}
\label{sec:v18v6}

V5 reached the compulsory five-version audit point.  The companion review is
recorded first.  Its main directional consequence is that the present
order-one massless curvature inequality cannot be deferred to a
weak-coupling factor that is only logarithmic in the lattice spacing.  The
rest of the section improves the V5 star criterion by retaining every
positive child stiffness term in the radial Schur denominator.

\begin{audit}[Post-V5 companion audit]
\label{aud:v18v6-companions}
The newest active companion files in the shared folder were inspected
read-only:
\[
\begin{array}{c|r|r}
\text{file}&\text{bytes}&\text{lines}\\
\hline
\texttt{Renewal\_Yang\_Mills\_Mass\_Gap\_Research\_Notes\_V1.tex}
&17560&466\\
\texttt{Renewal\_NS\_Stability\_Research\_Notes\_V26.tex}
&278283&5319
\end{array}                                                   \tag{14.1}
\]
Their SHA--256 digests are, respectively,
\[
\begin{split}
&\texttt{0EEF69B3521AAB3ED117685ABA372AB2117FCB9104F71DB978048633CA563261},\\
&\texttt{82C887F8C2C69E4F28FAD7C3DC8DE5B9099CB5A4BF431EBA1FFF3E91935978AD}.
                                                               \tag{14.2}
\end{split}
\]

The Yang--Mills programme has restarted after its tenth archive.  Its compact
V1 retains the exact degree-faithful forest bound and the finite reduced-fibre
torus estimates, then proves that a bound of order \(\beta^{-p}\) becomes
only \([\log(a_0/a)]^{-p}\) on its geometric cofinal clock.  It also proves
that such nonsummable forcing is harmless in a genuinely contracting
transverse recursion, while the complete scalar Ward mean and a contraction
factor below one remain open.  This supports a proof discipline needed here:
an inverse logarithm cannot be used as the \(a^2\) estimate demanded by
(2.8).

The Navier--Stokes V26 computes exact pointwise rank-one norms for the first
and second derivatives of its Oseen kernel.  The first derivative gains
nothing over the symmetric-traceless norm, while the second gains up to about
eleven percent pointwise in a finite radial window.  Those kernel constants
do not enter the Higgs radial message law.

No estimate or constant is imported from either companion.  The SM direction
is adjusted by keeping the massless star certificate at its natural
order-one scale and optimizing that certificate before considering cofinal
scaling.
\end{audit}

\subsection{The child stiffness omitted in V5}

For a pure quartic leaf \(i\), define
\[
\begin{split}
 D_i(q)&=\partial_{d_i}^2\ell_i(d_i,q),\\
 C_i(q)&=\partial_{d_i}\partial_q\ell_i(d_i,q),\\
 S_i(q)&=2\partial_q\ell_i(d_i,q)
          +4q\partial_q^2\ell_i(d_i,q).             \tag{14.3}
\end{split}
\]
Theorems \ref{thm:v18v2-global-leaf} and
\ref{thm:v18v3-closed-kernel} give
\[
 D_i(q)>0,
 \qquad S_i(q)>0.                                   \tag{14.4}
\]
For a star with root quartic coefficient \(\lambda_o\), equation (13.18)
is the exact identity
\[
 R_A(q)=12\lambda_oq+\sum_{i=1}^mS_i(q).            \tag{14.5}
\]

\begin{definition}[Stiffness-retaining thresholds]
\label{def:v18v6-threshold}
Define the collective threshold
\[
 \Lambda_\star:=
 \sup_{q>0}
 \frac1{12q}
 \left[
 \sum_{i=1}^m
 \left(\frac{4qC_i(q)^2}{D_i(q)}-S_i(q)\right)
 \right]_+,                                         \tag{14.6}
\]
and the separate leaf deficits
\[
 \mathcal N_i:=
 \sup_{q>0}
 \frac1{12q}
 \left[\frac{4qC_i(q)^2}{D_i(q)}-S_i(q)\right]_+.
                                                               \tag{14.7}
\]
Here \([x]_+=\max\{x,0\}\).
\end{definition}

\begin{proposition}[Finiteness and comparison]
\label{prop:v18v6-finite}
Every \(\mathcal N_i\) and \(\Lambda_\star\) is finite.  Moreover,
\[
 0\leq\Lambda_\star
 \leq\sum_{i=1}^m\mathcal N_i
 \leq\frac13\sum_{i=1}^m\mathcal M_i,              \tag{14.8}
\]
where \(\mathcal M_i\) is the V5 transfer norm (13.15).  These bounds are
uniform for fixed \(m\) on compact positive leaf-parameter sets.
\end{proposition}

\begin{proof}
At \(q=0\), continuity gives \(D_i(0)>0\), finite \(C_i(0)\), and
\[
 S_i(0)=2\partial_q\ell_i(d_i,0)>0.                 \tag{14.9}
\]
Hence the expression inside the positive part in (14.7) is negative for all
sufficiently small \(q\), so the quotient vanishes there.  As
\(q\to\infty\), Theorem \ref{thm:v18v5-transfer} gives
\[
 \frac{C_i(q)^2}{D_i(q)}=O(q^{-1}).                 \tag{14.10}
\]
Since \(S_i\geq0\), the quotient in (14.7) is bounded above by
\[
 \frac13\frac{C_i(q)^2}{D_i(q)}=O(q^{-1}).          \tag{14.11}
\]
It therefore tends to zero at infinity and is continuous on the remaining
compact interval.  This proves finiteness, including compact-parameter
uniformity by the uniform parts of Theorem
\ref{thm:v18v5-transfer} and Proposition
\ref{prop:v18v4-stiffness}.

The first inequality in (14.8) is immediate.  Subadditivity of the positive
part gives \(\Lambda_\star\leq\sum_i\mathcal N_i\).  Dropping the
nonnegative \(S_i\) in (14.7) gives
\[
 \mathcal N_i\leq
 \frac13\sup_{q>0}\frac{C_i(q)^2}{D_i(q)}
 =\frac13\mathcal M_i,                              \tag{14.12}
\]
which completes the proof.
\end{proof}

\begin{theorem}[Optimized pointwise-certificate star sector]
\label{thm:v18v6-optimized-star}
For the massless pure-quartic star of (13.13)--(13.14), if
\[
 \lambda_o\geq\Lambda_\star,                        \tag{14.13}
\]
then
\[
 \nabla_d^2F_\star(d)\succeq0.                     \tag{14.14}
\]
The value \(\Lambda_\star\) is the smallest root quartic coefficient for
which the pointwise Brascamp--Lieb Schur certificate (11.18) holds at every
\(q\).  It need not be the smallest coefficient for which the exact Hessian
in (14.14) is positive.
\end{theorem}

\begin{proof}
The pointwise star condition (11.18) is equivalent, using (14.5), to
\[
 \sum_{i=1}^m\frac{4qC_i(q)^2}{D_i(q)}
 \leq12\lambda_oq+\sum_{i=1}^mS_i(q).              \tag{14.15}
\]
For every \(q>0\), Definition \ref{def:v18v6-threshold} and
\(\lambda_o\geq\Lambda_\star\) give exactly (14.15).  At \(q=0\), the
rank-one term vanishes.  Theorem \ref{thm:v18v4-star} proves (14.14).

Conversely, if \(\lambda_o<\Lambda_\star\), the definition of the supremum
provides a \(q>0\) for which (14.15) fails, up to the usual limiting sequence
when the supremum is not attained.  Thus no smaller \(\lambda_o\) validates
the pointwise certificate for every \(q\).  Since Brascamp--Lieb is only an
upper bound on the actual score covariance, failure of this certificate is
not a counterexample to the exact Hessian sign.
\end{proof}

\begin{corollary}[Bounded-degree volume-uniform criterion]
\label{cor:v18v6-degree}
Suppose every star in a regulator family has degree at most
\(\Delta<\infty\), and all its leaf parameters lie in one compact positive
set.  Let \(\mathcal N_*\) be the common upper bound for (14.7).  Then
\[
 \lambda_o\geq\Delta\mathcal N_*                    \tag{14.16}
\]
is sufficient for (14.14) at every star, independently of the total graph
volume.
\end{corollary}

\begin{proof}
Proposition \ref{prop:v18v6-finite} gives
\(\Lambda_\star\leq\sum_i\mathcal N_i\leq
\Delta\mathcal N_*\).  Apply Theorem
\ref{thm:v18v6-optimized-star}.
\end{proof}

\section{V6 consequence and next decision}
\label{sec:v18v6-next}

The optimized threshold is an order-one, volume-uniform condition on a
bounded-degree regulator family.  It is therefore compatible in scale with
the canonical choice \(\lambda_H(a)=\bar\lambda_H\), but compatibility of
powers is not verification of the coefficient inequality.  The companion
audit rules out treating a logarithmic Yang--Mills weak-coupling gain as the
missing \(a^2\) estimate or as a substitute for (14.13).

The next task is to estimate \(\Lambda_\star\) analytically in terms of the
dimensionless ratios
\[
 \frac{\chi_i}{\sqrt{\lambda_i}},
 \qquad d_i,                                         \tag{15.1}
\]
and compare it with the canonical parent coefficient.  If a useful analytic
bound cannot be obtained, the proof should move upstream to the exact
one-edge or symmetric-star mixture rather than iterate a pointwise
certificate whose constants are unknown.

\begin{assessment}[V6 acquisition and boundary]
\label{ass:v18v6}
V6 completes the mandatory companion audit and improves the positive-star
sector from the coarse V5 threshold to the stiffness-retaining value
\(\Lambda_\star\).  The improved certificate is finite and volume-uniform
on bounded-degree compact regulator sectors.  It has not been verified for
the canonical Standard Model coefficients, it is not yet propagated through
tree depth, and no graph-with-cycles or continuum Einstein conclusion is
claimed.
\end{assessment}

\begin{record}[V6 append record]
\label{rec:v18v6}
V6 is appended to the validated V5 whose installed SHA--256 was
\texttt{9DD7ADC81F378155553608E695BFFAEE618001CBA277346A1D88320B1BE851F6}.
The audit used the exact companion fingerprints in (14.2).  After
installation only V6 is the active numeric SM note; the frozen
\texttt{V100\_f17} archive is unchanged.
\end{record}

\section{V7: exact dimensionless reduction of the star certificate}
\label{sec:v18v7}

V6 left the optimized threshold as a supremum over dimensional leaf
parameters.  A field dilation removes the quartic coefficient from the
integral and shows that every leaf threshold is a fixed amplitude times a
universal one-parameter function.  This identifies the coefficient
combination that must be tested in the canonical Standard Model regulator.

For a pure quartic leaf with parameters \(\lambda,\chi,d\), set
\[
 \rho=\frac{\chi e^{-d}}{\sqrt\lambda},
 \qquad
 \sigma=\frac{\chi^2}{\sqrt\lambda},
 \qquad
 Q=\sigma q=\frac{\chi^2}{\sqrt\lambda}|z|^2.       \tag{16.1}
\]
Define the universal radial free energy
\[
\begin{split}
 \mathscr L(\rho,Q)
 :={}&\frac Q\rho-\log\mathscr I(\rho,Q),\\
 \mathscr I(\rho,Q)
 :={}&\int_{\mathbb R^N}
 \exp\left[-|v|^4-\rho|v|^2
            +2\sqrt Q\,v_1\right]dv.                \tag{16.2}
\end{split}
\]
Rotation invariance makes the chosen first coordinate immaterial.

\begin{proposition}[Exact leaf scaling]
\label{prop:v18v7-leaf-scaling}
Up to an additive constant independent of \(d\) and \(q\),
\[
 \ell_{\lambda,\chi}(d,q)
 =\frac N4\log\lambda+\mathscr L(\rho,Q).           \tag{16.3}
\]
If
\[
\begin{split}
 \mathscr D(\rho,Q)
 &:=(\rho\partial_\rho)^2\mathscr L(\rho,Q),\\
 \mathscr C(\rho,Q)
 &:=-\rho\partial_\rho\partial_Q\mathscr L(\rho,Q),\\
 \mathscr S(\rho,Q)
 &:=2\partial_Q\mathscr L(\rho,Q)
   +4Q\partial_Q^2\mathscr L(\rho,Q),               \tag{16.4}
\end{split}
\]
then
\[
\begin{split}
 D(d,q)&=\partial_d^2\ell=\mathscr D(\rho,Q),\\
 C(d,q)&=\partial_d\partial_q\ell
          =\sigma\mathscr C(\rho,Q),\\
 S(d,q)&=2\partial_q\ell+4q\partial_q^2\ell
          =\sigma\mathscr S(\rho,Q).                \tag{16.5}
\end{split}
\]
\end{proposition}

\begin{proof}
In (5.4), substitute \(u=\lambda^{-1/4}v\) and rotate \(z\) to the first
coordinate.  Since
\[
 a|u|^2=\rho|v|^2,\qquad
 2\chi u\cdot z=2\sqrt Q\,v_1,\qquad
 bq=\frac Q\rho,                                    \tag{16.6}
\]
and \(du=\lambda^{-N/4}dv\), taking the negative logarithm gives
(16.3).  When \(d\) varies, \(Q\) is fixed and
\[
 \partial_d=-\rho\partial_\rho.                    \tag{16.7}
\]
When \(q\) varies, \(\partial_q=\sigma\partial_Q\).  Applying these two
identities to (16.3) proves (16.4)--(16.5).
\end{proof}

\begin{definition}[Universal transfer functions]
\label{def:v18v7-universal}
For \(\rho>0\), define
\[
 \widetilde{\mathcal M}_N(\rho)
 :=\sup_{Q\geq0}
 \frac{\mathscr C(\rho,Q)^2}{\mathscr D(\rho,Q)},   \tag{16.8}
\]
and
\[
 \widetilde{\mathcal N}_N(\rho)
 :=\sup_{Q>0}\frac1{12Q}
 \left[
 \frac{4Q\mathscr C(\rho,Q)^2}{\mathscr D(\rho,Q)}
 -\mathscr S(\rho,Q)
 \right]_+.                                         \tag{16.9}
\]
\end{definition}

\begin{theorem}[Exact factorization of the leaf thresholds]
\label{thm:v18v7-factorization}
The V5 and V6 leaf thresholds satisfy
\[
\begin{split}
 \mathcal M(\lambda,\chi,d)
 &=\frac{\chi^4}{\lambda}\,
   \widetilde{\mathcal M}_N
   \left(\frac{\chi e^{-d}}{\sqrt\lambda}\right),\\
 \mathcal N(\lambda,\chi,d)
 &=\frac{\chi^4}{\lambda}\,
   \widetilde{\mathcal N}_N
   \left(\frac{\chi e^{-d}}{\sqrt\lambda}\right).
                                                               \tag{16.10}
\end{split}
\]
Both universal functions are finite.  They are bounded when \(\rho\) ranges
over a compact subset of \((0,\infty)\), and
\[
 0\leq\widetilde{\mathcal N}_N(\rho)
 \leq\frac13\widetilde{\mathcal M}_N(\rho).         \tag{16.11}
\]
\end{theorem}

\begin{proof}
As \(q\) ranges over \([0,\infty)\), so does \(Q=\sigma q\).  Equations
(16.5) give
\[
 \frac{C(d,q)^2}{D(d,q)}
 =\sigma^2\frac{\mathscr C(\rho,Q)^2}
                    {\mathscr D(\rho,Q)}.           \tag{16.12}
\]
Since \(\sigma^2=\chi^4/\lambda\), taking the supremum proves the first
identity in (16.10).

For the V6 deficit, (16.5) and \(q=Q/\sigma\) give
\[
\begin{split}
 &\frac1{12q}
 \left[\frac{4qC(d,q)^2}{D(d,q)}-S(d,q)\right]_+\\
 &\qquad=
 \frac{\sigma^2}{12Q}
 \left[\frac{4Q\mathscr C(\rho,Q)^2}
 {\mathscr D(\rho,Q)}-\mathscr S(\rho,Q)\right]_+.
                                                               \tag{16.13}
\end{split}
\]
This proves the second identity.  Finiteness and compact-\(\rho\)
boundedness follow from Theorem \ref{thm:v18v5-transfer} and Proposition
\ref{prop:v18v6-finite} after fixing, for example, \(\lambda=1\) and
\(\chi e^{-d}=\rho\).  Inequality (16.11) is the scaled form of (14.8).
\end{proof}

\begin{corollary}[Dimensionless sufficient star condition]
\label{cor:v18v7-star-condition}
For a star with leaf parameters \((\lambda_i,\chi_i,d_i)\), the
stiffness-retaining sufficient condition
\[
 \lambda_o\geq
 \sum_{i=1}^m
 \frac{\chi_i^4}{\lambda_i}\,
 \widetilde{\mathcal N}_N
 \left(\frac{\chi_i e^{-d_i}}{\sqrt{\lambda_i}}\right)
                                                               \tag{16.14}
\]
implies
\[
 \nabla_d^2F_\star(d)\succeq0.                    \tag{16.15}
\]
The collective threshold \(\Lambda_\star\) from (14.6) can only improve
(16.14).
\end{corollary}

\begin{proof}
Theorem \ref{thm:v18v7-factorization} identifies each
\(\mathcal N_i\).  Proposition \ref{prop:v18v6-finite} gives
\(\Lambda_\star\leq\sum_i\mathcal N_i\), and Theorem
\ref{thm:v18v6-optimized-star} completes the proof.
\end{proof}

\begin{corollary}[Canonical bounded-sector form]
\label{cor:v18v7-canonical}
Suppose
\[
 |d_i|\leq L,\qquad
 0<\lambda_-\leq\lambda_i\leq\lambda_+,\qquad
 0<\chi_-\leq\chi_i\leq\chi_+,
 \qquad m\leq\Delta.                                \tag{16.16}
\]
Let
\[
 I_\rho=
 \left[
 \frac{\chi_-e^{-L}}{\sqrt{\lambda_+}},
 \frac{\chi_+e^{L}}{\sqrt{\lambda_-}}
 \right],
 \qquad
 N_*=\sup_{\rho\in I_\rho}
 \widetilde{\mathcal N}_N(\rho).                   \tag{16.17}
\]
Then the volume-uniform sufficient condition
\[
 \lambda_o\geq
 \Delta\,\frac{\chi_+^4}{\lambda_-}\,N_*            \tag{16.18}
\]
implies the star Hessian sign throughout the bounded conformal sector.
\end{corollary}

\begin{proof}
The interval \(I_\rho\) is a compact subset of \((0,\infty)\), so \(N_*<\infty\)
by Theorem \ref{thm:v18v7-factorization}.  Bound every summand in (16.14) by
\((\chi_+^4/\lambda_-)N_*\) and use \(m\leq\Delta\).
\end{proof}

\section{V7 consequence and next analytic target}
\label{sec:v18v7-next}

The relevant leaf threshold is not an arbitrary function of three
parameters.  It is
\[
 \frac{\chi^4}{\lambda}\,
 \widetilde{\mathcal N}_N(\rho),
 \qquad
 \rho=\frac{\chi e^{-d}}{\sqrt\lambda}.             \tag{17.1}
\]
Thus the canonical order-one comparison reduces to bounding one universal
function on the compact interval (16.17).  Equation (16.18) is rigorous but
still qualitative because no explicit upper bound for \(N_*\) has been
proved.

The next iteration should derive analytic bounds on
\(\widetilde{\mathcal N}_N(\rho)\) as \(\rho\) tends to zero and infinity,
then combine them with a compact-intermediate estimate.  This is the
upstream calculation needed before any claim that the physical Higgs
couplings satisfy the star certificate.

\begin{assessment}[V7 acquisition and boundary]
\label{ass:v18v7}
V7 reduces every pure-leaf transfer and deficit norm to a universal
one-parameter function and proves the exact amplitude
\(\chi^4/\lambda\).  It gives the explicit bounded-sector condition
(16.18).  The universal function has not yet been bounded by a closed
constant, the condition has not been checked for canonical coefficients,
and tree depth, cycles, cutoff removal, and the Einstein continuum remain
open.
\end{assessment}

\begin{record}[V7 append record]
\label{rec:v18v7}
V7 is appended to the validated V6 whose installed SHA--256 was
\texttt{33717BD31EC1111E7C900ADD9D6354C2EFD2F013CD7483F5FBDEC4DC45967421}.
After installation only V7 is the active numeric SM note; the frozen
\texttt{V100\_f17} archive is unchanged.
\end{record}

\section{V8: the large-boundary Schur separator}
\label{sec:v18v8}

V7 reduced the star threshold to a universal function but left open whether
the positive child stiffness in V6 might already pay the entire mixed Schur
cost.  This section answers that question negatively for the pointwise
Brascamp--Lieb certificate.  At large parent field, a pure quartic leaf
supplies asymptotically only one half of the stiffness required to absorb its
own mixed derivative.  The parent quartic term is therefore a structural
part of this certificate.

Retain one pure quartic leaf with
\[
 a=\chi e^{-d},\qquad b=\chi e^d,\qquad q=|z|^2,
                                                               \tag{18.1}
\]
and the functions
\[
 D(q)=\partial_d^2\ell(d,q),\qquad
 C(q)=\partial_d\partial_q\ell(d,q),\qquad
 S(q)=2\partial_q\ell(d,q)+4q\partial_q^2\ell(d,q).
                                                               \tag{18.2}
\]

\begin{lemma}[Differentiated large-field Laplace expansion]
\label{lem:v18v8-laplace}
As \(q\to\infty\),
\[
 \ell(d,q)=bq-Aq^{2/3}+Bq^{1/3}+O(\log q),          \tag{18.3}
\]
where
\[
 A=\frac{3\chi^{4/3}}{2^{4/3}\lambda^{1/3}},
 \qquad
 B=\frac{a\chi^{2/3}}{(2\lambda)^{2/3}}.            \tag{18.4}
\]
The expansion may be differentiated twice in \(q\) and twice in \(d\).
More precisely, on compact positive \((\lambda,\chi,d)\)-sets, the
corresponding remainder derivatives are
\[
\begin{array}{c|c}
\text{derivative}&\text{remainder bound}\\
\hline
1&O(\log q)\\
\partial_q&O(q^{-1})\\
\partial_q^2&O(q^{-2})\\
\partial_d&O(\log q)\\
\partial_d^2&O(\log q)\\
\partial_d\partial_q&O(q^{-1}).
\end{array}                                                   \tag{18.5}
\]
\end{lemma}

\begin{proof}
Write \(H=\chi\sqrt q\).  The nontrivial integral in (5.4) is
\[
 I(a,H)=\int_{\mathbb R^N}
 e^{-\lambda|u|^4-a|u|^2+2Hu_1}\,du.               \tag{18.6}
\]
After \(u=H^{1/3}v\),
\[
 I(a,H)=H^{N/3}\int_{\mathbb R^N}
 \exp\left[
 H^{4/3}\Phi(v)-aH^{2/3}|v|^2\right]dv,             \tag{18.7}
\]
where
\[
 \Phi(v)=2v_1-\lambda|v|^4.                         \tag{18.8}
\]
The function \(\Phi\) has the unique maximizer
\[
 v_0=(2\lambda)^{-1/3}e_1,                          \tag{18.9}
\]
and its Hessian at \(v_0\) is negative definite:
\[
 \nabla^2\Phi(v_0)
 =-4\lambda|v_0|^2I-8\lambda v_0v_0^{\mathsf T}.   \tag{18.10}
\]
The implicit-function theorem gives a unique maximizer of the exponent in
(18.7) in a fixed neighborhood of \(v_0\), equal to
\(v_0+O(H^{-2/3})\).  Taylor expansion at that point gives
\[
 \log I(a,H)
 =H^{4/3}\Phi(v_0)-aH^{2/3}|v_0|^2+O(\log H).       \tag{18.11}
\]
Outside a fixed neighborhood of \(v_0\), strict maximality and the quartic
tail give an exponentially smaller integral.  Inside it, the change of
variables by the square root of the negative Hessian reduces the integral
to a Gaussian integral times a factor with a complete expansion in
\(H^{-2/3}\).  Polynomial insertions obtained by differentiating in \(a\)
or \(H\) obey the same localization.  This proves (18.11) with the
differentiated remainder bounds in (18.5); all constants are uniform on the
stated compact parameter sets.

Finally,
\[
 \Phi(v_0)=\frac3{2^{4/3}\lambda^{1/3}},
 \qquad
 |v_0|^2=(2\lambda)^{-2/3}.                         \tag{18.12}
\]
Since \(\ell=bq-\log I(a,\chi\sqrt q)\), equations
(18.11)--(18.12) give (18.3)--(18.4).
\end{proof}

\begin{theorem}[Asymptotic factor-two deficit]
\label{thm:v18v8-factor-two}
For every fixed \(N\geq1\), \(\lambda,\chi>0\), and \(d\in\mathbb R\),
\[
\begin{split}
 D(q)&=bq+Bq^{1/3}+O(\log q),\\
 C(q)&=b-\frac13Bq^{-2/3}+O(q^{-1}),\\
 S(q)&=2b-\frac49Aq^{-1/3}
          -\frac29Bq^{-2/3}+O(q^{-1}).              \tag{18.13}
\end{split}
\]
Consequently
\[
 \frac{4qC(q)^2}{D(q)}=4b+O(q^{-2/3}),              \tag{18.14}
\]
and
\[
 \frac{4qC(q)^2}{D(q)}-S(q)
 =2b+\frac49Aq^{-1/3}+O(q^{-2/3})>0                \tag{18.15}
\]
for all sufficiently large \(q\).
\end{theorem}

\begin{proof}
The coefficient \(A\) is independent of \(d\), while
\[
 \partial_db=b,\qquad \partial_da=-a,\qquad
 \partial_dB=-B.                                    \tag{18.16}
\]
Apply the differentiated expansion in Lemma
\ref{lem:v18v8-laplace} to obtain the first two lines of (18.13).
Direct \(q\)-differentiation of (18.3) gives
\[
\begin{split}
 \partial_q\ell
 &=b-\frac23Aq^{-1/3}+\frac13Bq^{-2/3}+O(q^{-1}),\\
 \partial_q^2\ell
 &=\frac29Aq^{-4/3}-\frac29Bq^{-5/3}+O(q^{-2}).
                                                               \tag{18.17}
\end{split}
\]
Substitution in the definition of \(S\) yields its line in (18.13).
The first two lines of (18.13) give (18.14), and subtracting (18.13) from
(18.14) gives (18.15).
\end{proof}

\begin{corollary}[Strict positivity of the universal deficit]
\label{cor:v18v8-positive-deficit}
For every \(\rho>0\),
\[
 \widetilde{\mathcal N}_N(\rho)>0.                 \tag{18.18}
\]
More precisely, in the universal variables of V7,
\[
 \frac{4Q\mathscr C(\rho,Q)^2}{\mathscr D(\rho,Q)}
 -\mathscr S(\rho,Q)
 \longrightarrow\frac2\rho
 \qquad(Q\to\infty).                                \tag{18.19}
\]
\end{corollary}

\begin{proof}
Set \(\lambda=1\), \(\chi=1\), and \(d=-\log\rho\).  Then
\(a=\rho\), \(b=1/\rho\), \(q=Q\), and the functions in (18.2) are the
universal functions in (16.4).  Equation (18.15) gives (18.19).  The
expression inside the positive part in (16.9) is therefore positive for all
sufficiently large finite \(Q\), so its supremum is strictly positive.
\end{proof}

\begin{corollary}[Child stiffness alone cannot certify a star]
\label{cor:v18v8-no-child-only}
Consider the pointwise star Schur condition (14.15) with
\(\lambda_o=0\).  For any nonempty collection of pure quartic leaves, that
condition fails for all sufficiently large common parent radii \(q\).
\end{corollary}

\begin{proof}
For each leaf, (18.15) gives
\[
 \frac{4qC_i(q)^2}{D_i(q)}-S_i(q)\longrightarrow2b_i>0.
\]
The sum of these differences is therefore positive at large \(q\), whereas
(14.15) with \(\lambda_o=0\) requires the sum to be nonpositive.
\end{proof}

\section{V8 consequence and redirected next step}
\label{sec:v18v8-next}

The factor-two limit (18.14)--(18.15) rules out a hoped-for
degree-independent closure using child stiffness alone.  It does not disprove
the exact star Hessian, because the Brascamp--Lieb covariance bound may be
strict.  It proves that any continuation of the pointwise certificate must
retain the parent quartic term and quantify its intermediate-radius maximum.

There are now two honest routes.  One can derive a certified upper bound for
the one-parameter function \(\widetilde{\mathcal N}_N\), or one can move
upstream to the exact angular expansion of the complete star partition and
estimate its discrete score covariance without Brascamp--Lieb.  The latter
can improve the sufficient condition and can decide whether failure of the
pointwise certificate is merely an artifact.  V9 should formulate that exact
star mixture before choosing between the two routes.

\begin{assessment}[V8 acquisition and boundary]
\label{ass:v18v8}
V8 proves a differentiated large-boundary Laplace expansion and the
factor-two Schur deficit.  The universal stiffness deficit is strictly
positive for every \(\rho>0\), so the parent quartic cannot be deleted from
the pointwise certificate.  The exact star Hessian may still be positive
outside that certificate; no tree-depth, cycle, continuum, or Einstein
conclusion is asserted.
\end{assessment}

\begin{record}[V8 append record]
\label{rec:v18v8}
V8 is appended to the validated V7 whose installed SHA--256 was
\texttt{97CC2E168F018276983CAAFA129F21359543059E093CC9BD92A89C0A8AA9E55D}.
After installation only V8 is the active numeric SM note; the frozen
\texttt{V100\_f17} archive is unchanged.
\end{record}

\section{V9: the exact angular mixture of a complete star}
\label{sec:v18v9}

V8 showed that the pointwise Brascamp--Lieb certificate loses a factor two at
large boundary field.  To decide whether this loss belongs to the exact star
or only to that covariance estimate, this section integrates the entire star
without an intermediate inequality.  The result is a positive multi-index
mixture.  Every fixed-index component has a positive conformal Hessian; the
only remaining obstruction is a discrete score covariance.

Let \(o\) be the root of a star and \(x_1,\ldots,x_m\) its leaves.  Put
\[
 a_i=\chi_i e^{-d_i},\qquad
 b_i=\chi_i e^{d_i},\qquad
 A=\sum_{i=1}^m b_i,\qquad
 \alpha=\frac N2.                                   \tag{19.1}
\]
For \(s>0\), define the root and leaf radial integrals
\[
\begin{split}
 J_{0,s}(A)&=\int_0^\infty
 t^{s-1}e^{-\lambda_ot^2-At}\,dt,\\
 J_{i,s}(a_i)&=\int_0^\infty
 t^{s-1}e^{-\lambda_it^2-a_it}\,dt.                 \tag{19.2}
\end{split}
\]

\begin{theorem}[Exact positive star expansion]
\label{thm:v18v9-star-expansion}
There is a constant \(C_{N,m}>0\), independent of \(d\), such that the
massless star partition function is
\[
\begin{split}
 Z_\star(d)
 =C_{N,m}\sum_{\mathbf k\in\mathbb N^m}
 &J_{0,\alpha+|\mathbf k|}(A)\\
 &\times\prod_{i=1}^m
 \frac{\chi_i^{2k_i}J_{i,\alpha+k_i}(a_i)}
 {k_i!(\alpha)_{k_i}},
                                                               \tag{19.3}
\end{split}
\]
where \(|\mathbf k|=\sum_i k_i\).  The series and all its finite-order
conformal derivatives converge locally uniformly.
\end{theorem}

\begin{proof}
After tree gauge erasure, write the root field as \(z\) and leaf fields as
\(u_i\).  Expanding every hopping square gives
\[
\begin{split}
 S={}&\lambda_o|z|^4+A|z|^2\\
 &+\sum_{i=1}^m\left(
 \lambda_i|u_i|^4+a_i|u_i|^2-2\chi_i u_i\cdot z
 \right).                                           \tag{19.4}
\end{split}
\]
Condition on \(r=|z|\).  The spherical average in the \(i\)-th leaf is
\[
 \sum_{k_i=0}^\infty
 \frac{(\chi_i^2r^2|u_i|^2)^{k_i}}
 {k_i!(\alpha)_{k_i}},                              \tag{19.5}
\]
by (7.7).  All coefficients are nonnegative, so Tonelli's theorem permits
termwise integration of every leaf radius.  The product of (19.5) supplies
the root factor \(r^{2|\mathbf k|}\).  The substitution \(t=r^2\) in the
root integral then gives \(J_{0,\alpha+|\mathbf k|}(A)\), proving (19.3).

Local uniform convergence of conformal derivatives follows either by
repeating the quartic domination argument of Lemma
\ref{lem:v18v1-differentiability} before expanding, or by applying it to
partial sums and using monotone convergence followed by dominated
differentiation.  This also justifies the differentiations below.
\end{proof}

Normalize the positive summands in (19.3) to a probability law
\(p_d(\mathbf k)\).  For \(s_0=\alpha+|\mathbf k|\) and
\(s_i=\alpha+k_i\), set
\[
\begin{array}{lll}
 m_0=\dfrac{J_{0,s_0+1}}{J_{0,s_0}},&
 v_0=\dfrac{J_{0,s_0+2}}{J_{0,s_0}}-m_0^2,&
 c_0=A m_0-A^2v_0,\\
 m_i=\dfrac{J_{i,s_i+1}}{J_{i,s_i}},&
 v_i=\dfrac{J_{i,s_i+2}}{J_{i,s_i}}-m_i^2,&
 c_i=a_i m_i-a_i^2v_i.
                                                               \tag{19.6}
\end{array}
\]
Every quantity in (19.6) depends on \(\mathbf k\), and \(c_0,c_i>0\) by
the zero-boundary curvature argument (7.14).

\begin{definition}[Component free energy and score]
\label{def:v18v9-component}
Up to a term independent of \(d\), define
\[
 \phi_{\mathbf k}(d)=
 -\log J_{0,s_0}(A)-\sum_{i=1}^m\log J_{i,s_i}(a_i).
                                                               \tag{19.7}
\]
Its conformal score is the vector \(g(\mathbf k)\in\mathbb R^m\) with
\[
 g_i(\mathbf k)=b_i m_0-a_i m_i.                    \tag{19.8}
\]
\end{definition}

\begin{proposition}[Every angular component is strictly convex]
\label{prop:v18v9-component-convex}
For every \(\mathbf k\),
\[
 \nabla_d\phi_{\mathbf k}=g(\mathbf k),              \tag{19.9}
\]
and
\[
 \nabla_d^2\phi_{\mathbf k}
 =
 \operatorname{diag}(b_1m_0+c_1,\ldots,b_mm_0+c_m)
 -v_0bb^{\mathsf T}\succ0,                          \tag{19.10}
\]
where \(b=(b_1,\ldots,b_m)^{\mathsf T}\).
\end{proposition}

\begin{proof}
For \(f(a)=-\log J_s(a)\),
\[
 f'(a)=m_s,\qquad f''(a)=-v_s.                      \tag{19.11}
\]
Use \(\partial_{d_i}a_i=-a_i\),
\(\partial_{d_i}b_i=b_i\), and \(\partial_{d_i}A=b_i\).
This gives (19.8)--(19.9).  A second differentiation gives the root block
\[
 m_0\operatorname{diag}(b)-v_0bb^{\mathsf T}        \tag{19.12}
\]
and the leaf diagonal \(\operatorname{diag}(c_i)\), proving (19.10).

For \(x\in\mathbb R^m\), weighted Cauchy--Schwarz gives
\[
 (b^{\mathsf T}x)^2
 \leq A\sum_i b_ix_i^2.                             \tag{19.13}
\]
Therefore
\[
\begin{split}
 x^{\mathsf T}
 \left(m_0\operatorname{diag}(b)-v_0bb^{\mathsf T}\right)x
 &\geq(m_0-Av_0)\sum_i b_ix_i^2\\
 &=\frac{c_0}{A}\sum_i b_ix_i^2>0.                 \tag{19.14}
\end{split}
\]
Adding the positive leaf diagonal proves strict positivity in (19.10).
\end{proof}

\begin{theorem}[Necessary and sufficient discrete covariance gate]
\label{thm:v18v9-discrete-gate}
The exact star free-energy Hessian is
\[
 \nabla_d^2F_\star(d)
 =
 \mathbb E_{p_d}\nabla_d^2\phi_{\mathbf K}(d)
 -\operatorname{Cov}_{p_d}(g(\mathbf K)).           \tag{19.15}
\]
Consequently
\[
 \nabla_d^2F_\star(d)\succeq0                       \tag{19.16}
\]
if and only if
\[
 \operatorname{Cov}_{p_d}(g(\mathbf K))
 \preceq
 \mathbb E_{p_d}\left[
 \operatorname{diag}(b_im_0+c_i)-v_0bb^{\mathsf T}
 \right].                                           \tag{19.17}
\]
Thus every continuous fixed-component curvature is positive.  Any negative
star direction must be produced entirely by mixing the angular indices.
\end{theorem}

\begin{proof}
Equation (19.3) has the form
\[
 Z_\star(d)=\sum_{\mathbf k}c_{\mathbf k}
 e^{-\phi_{\mathbf k}(d)}
                                                               \tag{19.18}
\]
with \(c_{\mathbf k}>0\) independent of \(d\).  Differentiate
\(-\log Z_\star\) twice.  The standard log-sum identity gives the expectation
of the component Hessian minus the covariance of the component gradient.
Use Proposition \ref{prop:v18v9-component-convex} to obtain (19.15), and
rearrange it to obtain the equivalence (19.17).
\end{proof}

\begin{proposition}[Exact coordinate birth ratios]
\label{prop:v18v9-birth}
For \(\mathbf e_i\) the \(i\)-th coordinate vector, define
\[
 r_i(\mathbf k)=
 \frac{(k_i+1)p_d(\mathbf k+\mathbf e_i)}
 {p_d(\mathbf k)}.                                  \tag{19.19}
\]
Then
\[
 r_i(\mathbf k)
 =\chi_i^2\,m_0\,\frac{m_i}{s_i}.                   \tag{19.20}
\]
In particular,
\[
 0<r_i(\mathbf k)<b_i m_0,\qquad
 \sum_{i=1}^m r_i(\mathbf k)<A m_0<s_0.             \tag{19.21}
\]
\end{proposition}

\begin{proof}
Take the ratio of the two summands in (19.3).  Increasing \(k_i\) by one
raises both \(s_i\) and \(s_0\) by one.  The root radial ratio is \(m_0\),
the leaf radial ratio is \(m_i\), and
\[
 \frac{(\alpha)_{k_i+1}}{(\alpha)_{k_i}}=s_i.
\]
This proves (19.20).  The recurrence (7.12) gives
\[
 \frac{a_im_i}{s_i}<1,
\]
so \(\chi_i^2m_i/s_i=a_ib_i m_i/s_i<b_i\).  Summing gives the first strict
bound in (19.21).  Applying the root recurrence with quadratic coefficient
\(A\) gives \(A m_0<s_0\).
\end{proof}

\begin{remark}[Why the one-leaf ULC proof does not tensorize immediately]
\label{rem:v18v9-no-tensor}
For a single conditional leaf, the birth ratio decreases with its one
angular index.  In (19.20), increasing a different coordinate raises the
root shape \(s_0\) and hence raises \(m_0\).  The coordinates therefore have
a positive coupling through \(|\mathbf K|\).  A product
ultra-log-concavity argument would discard precisely the covariance that
must be controlled in (19.17).
\end{remark}

\section{V9 consequence and the discrete target}
\label{sec:v18v9-next}

The exact star problem is now finite-dimensional in a stronger sense than
the V6 pointwise certificate.  The continuous radial integrations enter only
through moment ratios \(m_i,v_i,m_0,v_0\); the sole unresolved sign is the
discrete covariance inequality (19.17).  The birth ratios (19.20) identify a
natural reversible birth--death structure on \(\mathbb N^m\).

V10 should construct its Dirichlet form and seek a matrix Poincare estimate
for the particular score \(g\), retaining the positive root block
(19.12).  A generic product Poincare inequality is not enough because the
root moment ratio couples all coordinates.  If the discrete estimate loses
the same factor two as V8, the programme should return to the exact
component-score increments before imposing absolute values.

\begin{assessment}[V9 acquisition and boundary]
\label{ass:v18v9}
V9 derives the exact positive angular expansion of a complete massless star,
proves strict convexity of every fixed angular component, and reduces the
full Hessian sign to the necessary-and-sufficient discrete covariance gate
(19.17).  The multi-index mixing law has explicit birth ratios but is not a
product law.  The discrete covariance gate, tree-depth propagation, cycles,
cutoff removal, and the coupled Einstein continuum remain open.
\end{assessment}

\begin{record}[V9 append record]
\label{rec:v18v9}
V9 is appended to the validated V8 whose installed SHA--256 was
\texttt{771DD707D075882423CC4D9A9B548A876B609E7D0C58AD25445765C35CBA615A}.
After installation only V9 is the active numeric SM note; the frozen
\texttt{V100\_f17} archive is unchanged.
\end{record}

\section{V10: subfactorial angular tails and a finite discrete certificate}
\label{sec:v18v10}

The exact gate (19.17) is an infinite discrete covariance inequality.  This
section proves that its angular mixing law has a subfactorial total-degree
tail.  The gate can therefore be reduced, with an explicit error, to a finite
matrix sum of radial moment ratios.  This does not yet verify the sign for
canonical parameters; it makes such a verification rigorous and finite.

Retain the star notation of V9.  Let
\[
 K=|\mathbf K|,\qquad
 P_n=\mathbb P_{p_d}(K=n),\qquad
 R(\mathbf k)=\sum_{i=1}^m r_i(\mathbf k),           \tag{20.1}
\]
where \(r_i\) is the exact coordinate birth ratio (19.20).

\begin{lemma}[Square-root bound for radial moment ratios]
\label{lem:v18v10-radial-bound}
For
\[
 J_s(a,\lambda)=\int_0^\infty
 t^{s-1}e^{-\lambda t^2-at}\,dt,
 \qquad
 m_s=\frac{J_{s+1}}{J_s},
\]
one has
\[
 0<m_s\leq\sqrt{\frac{s}{2\lambda}}.                \tag{20.2}
\]
\end{lemma}

\begin{proof}
The integration-by-parts recurrence is
\[
 s=a m_s+2\lambda m_sm_{s+1}.                       \tag{20.3}
\]
Strict moment-ratio monotonicity gives \(m_{s+1}\geq m_s\).  Dropping the
nonnegative first term in (20.3) yields
\(s\geq2\lambda m_s^2\), which is (20.2).
\end{proof}

Define
\[
 C_\star=
 \frac1{\sqrt\alpha}
 \sum_{i=1}^m
 \frac{\chi_i^2}{2\sqrt{\lambda_o\lambda_i}}.       \tag{20.4}
\]

\begin{proposition}[Sublinear total birth rate]
\label{prop:v18v10-birth-bound}
For every \(\mathbf k\in\mathbb N^m\),
\[
 R(\mathbf k)\leq C_\star\sqrt{\alpha+|\mathbf k|}.
                                                               \tag{20.5}
\]
\end{proposition}

\begin{proof}
Put \(s_0=\alpha+|\mathbf k|\) and \(s_i=\alpha+k_i\).  Equations
(19.20) and (20.2) give
\[
\begin{split}
 r_i(\mathbf k)
 &\leq
 \chi_i^2
 \frac{\sqrt{s_0/(2\lambda_o)}
       \sqrt{s_i/(2\lambda_i)}}{s_i}\\
 &=\frac{\chi_i^2}{2\sqrt{\lambda_o\lambda_i}}
   \sqrt{\frac{s_0}{s_i}}
 \leq
 \frac{\chi_i^2}{2\sqrt{\lambda_o\lambda_i\alpha}}
 \sqrt{s_0}.                                        \tag{20.6}
\end{split}
\]
Sum over \(i\) to obtain (20.5).
\end{proof}

\begin{theorem}[Subfactorial total-degree tail]
\label{thm:v18v10-tail}
For every \(n\geq0\),
\[
 P_n\leq P_0\,u_n,\qquad
 u_n:=\frac{C_\star^n\sqrt{(\alpha)_n}}{n!}.         \tag{20.7}
\]
In particular, \(\sum_n n^pu_n<\infty\) for every \(p\geq0\).
\end{theorem}

\begin{proof}
Detailed balance in coordinate \(i\) is
\[
 p_d(\mathbf k)r_i(\mathbf k)
 =(k_i+1)p_d(\mathbf k+\mathbf e_i).                \tag{20.8}
\]
Sum (20.8) over all \(i\) and all \(\mathbf k\) with
\(|\mathbf k|=n\).  Each state of total degree \(n+1\) is counted with
coefficient \(\sum_i k_i=n+1\), so
\[
 (n+1)P_{n+1}
 =\mathbb E\left[R(\mathbf K)\mathbf1_{\{K=n\}}\right]
 \leq C_\star\sqrt{\alpha+n}\,P_n.                 \tag{20.9}
\]
Iteration of (20.9) gives (20.7).  Finally,
\[
 \frac{u_{n+1}}{u_n}
 =\frac{C_\star\sqrt{\alpha+n}}{n+1}
 =O(n^{-1/2}),                                      \tag{20.10}
\]
so the ratio test remains valid after multiplication by any fixed power of
\(n\).
\end{proof}

\subsection{Polynomial envelopes for the exact score and component Hessian}

Set
\[
\begin{split}
 G_\star^2&=
 \sum_{i=1}^m
 \left(\frac{b_i}{\sqrt{2\lambda_o}}
       +\frac{a_i}{\sqrt{2\lambda_i}}\right)^2,\\
 H_\star&=
 \frac{\sum_i b_i}{\sqrt{2\lambda_o}}
 +\sum_{i=1}^m\frac{a_i}{\sqrt{2\lambda_i}}.        \tag{20.11}
\end{split}
\]
Let
\[
 H(\mathbf k)=\nabla_d^2\phi_{\mathbf k},
 \qquad g(\mathbf k)=\nabla_d\phi_{\mathbf k}.       \tag{20.12}
\]

\begin{lemma}[Score and Hessian envelopes]
\label{lem:v18v10-envelopes}
For every \(\mathbf k\),
\[
 |g(\mathbf k)|^2
 \leq G_\star^2(\alpha+|\mathbf k|),                \tag{20.13}
\]
and
\[
 0\prec H(\mathbf k),\qquad
 \|H(\mathbf k)\|_{\mathrm{op}}
 \leq H_\star\sqrt{\alpha+|\mathbf k|}.             \tag{20.14}
\]
\end{lemma}

\begin{proof}
Equation (19.8), the triangle inequality, and (20.2) give
\[
 |g_i(\mathbf k)|
 \leq
 \left(\frac{b_i}{\sqrt{2\lambda_o}}
       +\frac{a_i}{\sqrt{2\lambda_i}}\right)
 \sqrt{\alpha+|\mathbf k|}.
\]
Squaring and summing proves (20.13).

Strict positivity of \(H\) is Proposition
\ref{prop:v18v9-component-convex}.  A positive matrix has operator norm no
larger than its trace.  From (19.10) and \(c_i<a_im_i\),
\[
\begin{split}
 \operatorname{Tr}H
 &\leq m_0\sum_i b_i+\sum_i a_im_i\\
 &\leq H_\star\sqrt{\alpha+|\mathbf k|},            \tag{20.15}
\end{split}
\]
where (20.2) was used once more.  This proves (20.14).
\end{proof}

\subsection{A finite matrix certificate}

Let \(w_{\mathbf k}\) denote the unnormalized positive summand in (19.3),
with its common factor \(C_{N,m}\) removed, and put
\[
 w_{\mathbf0}=
 J_{0,\alpha}(A)\prod_{i=1}^mJ_{i,\alpha}(a_i).     \tag{20.16}
\]
For an integer \(L\geq0\), define
\[
\begin{split}
 \mathcal Q_L&=
 \sum_{|\mathbf k|\leq L}
 w_{\mathbf k}
 \left(H(\mathbf k)-g(\mathbf k)g(\mathbf k)^{\mathsf T}\right),\\
 \varepsilon_L&=
 w_{\mathbf0}G_\star^2
 \sum_{n>L}(\alpha+n)u_n.                           \tag{20.17}
\end{split}
\]
The scalar tail in (20.17) is explicit and convergent by Theorem
\ref{thm:v18v10-tail}.

\begin{theorem}[Finite sufficient certificate for the exact star Hessian]
\label{thm:v18v10-finite-certificate}
If, for some \(L\),
\[
 \mathcal Q_L\succeq\varepsilon_L I_m,              \tag{20.18}
\]
then
\[
 \nabla_d^2F_\star(d)\succeq0.                     \tag{20.19}
\]
All entries in (20.18) are finite sums of the one-dimensional radial
integrals (19.2).  Thus interval bounds for finitely many such integrals,
together with the explicit series tail in (20.17), give a rigorous
finite-dimensional certificate.
\end{theorem}

\begin{proof}
Let \(Z_0=\sum_{\mathbf k}w_{\mathbf k}\), so the common angular factor is
irrelevant and \(p_d(\mathbf k)=w_{\mathbf k}/Z_0\).  Multiplying (19.15) by
\(Z_0\) gives
\[
\begin{split}
 Z_0\nabla_d^2F_\star
 ={}&\sum_{\mathbf k}w_{\mathbf k}
 \left(H(\mathbf k)-g(\mathbf k)g(\mathbf k)^{\mathsf T}\right)\\
 &+\frac1{Z_0}
 \left(\sum_{\mathbf k}w_{\mathbf k}g(\mathbf k)\right)
 \left(\sum_{\mathbf k}w_{\mathbf k}g(\mathbf k)\right)^{\mathsf T}.
                                                               \tag{20.20}
\end{split}
\]
The second line is positive semidefinite.  The positive \(H\) part of the
tail may also be discarded in a lower bound.  By (20.13),
\[
\begin{split}
 \sum_{|\mathbf k|>L}
 w_{\mathbf k}g(\mathbf k)g(\mathbf k)^{\mathsf T}
 &\preceq
 G_\star^2
 \sum_{n>L}(\alpha+n)\sum_{|\mathbf k|=n}w_{\mathbf k}\,I_m.
                                                               \tag{20.21}
\end{split}
\]
The proof of Theorem \ref{thm:v18v10-tail} before normalization gives
\[
 \sum_{|\mathbf k|=n}w_{\mathbf k}
 \leq w_{\mathbf0}u_n.                              \tag{20.22}
\]
Equations (20.20)--(20.22) yield
\[
 Z_0\nabla_d^2F_\star\succeq\mathcal Q_L-\varepsilon_LI_m.
                                                               \tag{20.23}
\]
Condition (20.18) and \(Z_0>0\) prove (20.19).
\end{proof}

\begin{remark}[Strength and loss of the finite certificate]
\label{rem:v18v10-loss}
The certificate drops the positive full-score mean in (20.20) and the
positive component Hessian in the tail.  It is sufficient, not necessary.
Failure at a chosen \(L\) is therefore not a negative-curvature example.
Increasing \(L\) reduces the explicit tail and retains more of the exact
matrix.
\end{remark}

\section{V10 consequence and audit point}
\label{sec:v18v10-next}

The exact discrete gate can now be checked by a finite, auditable
calculation.  No unproved truncation or heuristic angular cutoff is needed:
the omitted mass has the subfactorial envelope (20.7), and the score has the
polynomial envelope (20.13).  The next analytic question is whether
(20.18) can be proved uniformly on a compact canonical parameter box without
an impractically large \(L\), or whether a sharper discrete Poincare
inequality should retain the positive mean term in (20.20).

This is V10 of the eighteenth cycle.  The newest active Yang--Mills and
Navier--Stokes notes must now be inspected before V11 is chosen.

\begin{assessment}[V10 acquisition and boundary]
\label{ass:v18v10}
V10 proves a subfactorial tail for the complete-star angular law and a finite
matrix certificate for the exact Hessian.  The certificate has not been
evaluated on a canonical parameter box and remains sufficient rather than
necessary.  Tree depth, holonomy cycles, cofinal state construction, and the
Einstein continuum remain open.
\end{assessment}

\begin{record}[V10 append record]
\label{rec:v18v10}
V10 is appended to the validated V9 whose installed SHA--256 was
\texttt{020BE2EDBA8DEF947D27EC4D6A5867CFD6628FFAA9C86131525E66530208FA55}.
After installation only V10 is the active numeric SM note; the frozen
\texttt{V100\_f17} archive is unchanged.
\end{record}

\section{V11: companion audit and a mean-retaining finite certificate}
\label{sec:v18v11}

The mandatory post-V10 companion audit is recorded first.  The mathematical
continuation then strengthens the finite angular certificate by retaining a
controlled part of the positive full-score mean in (20.20).  This removes a
deliberate loss in V10 without weakening the explicit subfactorial tail.

\begin{audit}[Post-V10 companion audit]
\label{aud:v18v11-companions}
The newest active companion files inspected read-only were
\[
\begin{array}{c|r|r}
\text{file}&\text{bytes}&\text{lines}\\
\hline
\texttt{Renewal\_Yang\_Mills\_Mass\_Gap\_Research\_Notes\_V2.tex}
&34462&910\\
\texttt{Renewal\_NS\_Stability\_Research\_Notes\_V26.tex}
&278283&5319.
\end{array}                                                   \tag{21.1}
\]
Their SHA--256 digests are
\[
\begin{split}
&\texttt{20687686090948DA5839DE226CCF845899B5CF0508058C020B8EEDA23016F34C},\\
&\texttt{82C887F8C2C69E4F28FAD7C3DC8DE5B9099CB5A4BF431EBA1FFF3E91935978AD}.
                                                               \tag{21.2}
\end{split}
\]

Yang--Mills V2 proves completeness of a commuting-toron Gaussian endpoint,
a floor-free multiplicative residual enclosure for complete trace words, and
an exact rational negative interval for its \(M=4\) quarter-turn grid.  It
also keeps the decisive boundary: the finite grid is not the continuous
torus Haar mean until a quantified alias tail is controlled.  The useful
lesson for the present programme is methodological.  One should retain the
complete signed score and its normalization before applying absolute values,
and then certify the omitted tail.  The Yang--Mills response values and
residual constants do not enter the Higgs star law.

Navier--Stokes V26 and its fingerprint are unchanged from the V6 audit.  Its
rank-one Oseen-kernel refinement supplies no estimate for the angular Higgs
mixture.  No companion theorem or constant is imported.
\end{audit}

\subsection{Tail data for normalization and score}

Keep the V9--V10 unnormalized weights \(w_{\mathbf k}\), score
\(g(\mathbf k)\), component Hessian \(H(\mathbf k)\), and cutoff \(L\).
In addition to \(\mathcal Q_L\) in (20.17), define
\[
\begin{split}
 W_L&=\sum_{|\mathbf k|\leq L}w_{\mathbf k},\\
 s_L&=\sum_{|\mathbf k|\leq L}w_{\mathbf k}g(\mathbf k),\\
 \tau_{0,L}&=\sum_{n>L}u_n,\\
 \tau_{1/2,L}&=\sum_{n>L}\sqrt{\alpha+n}\,u_n,\\
 \tau_{1,L}&=\sum_{n>L}(\alpha+n)u_n.               \tag{21.3}
\end{split}
\]
All three tails converge and tend to zero by (20.10).  Put
\[
 Z_L^+=W_L+w_{\mathbf0}\tau_{0,L},
 \qquad
 \eta_L=w_{\mathbf0}G_\star\tau_{1/2,L},
 \qquad
 \varepsilon_L=w_{\mathbf0}G_\star^2\tau_{1,L}.     \tag{21.4}
\]
The last quantity agrees with the V10 tail error.

\begin{lemma}[Normalization and vector-score tails]
\label{lem:v18v11-tail-data}
Let
\[
 Z_0=\sum_{\mathbf k}w_{\mathbf k},
 \qquad
 s=\sum_{\mathbf k}w_{\mathbf k}g(\mathbf k).
                                                               \tag{21.5}
\]
Then
\[
 W_L\leq Z_0\leq Z_L^+,
 \qquad
 |s-s_L|\leq\eta_L.                                 \tag{21.6}
\]
\end{lemma}

\begin{proof}
The unnormalized form of (20.7) is
\[
 \sum_{|\mathbf k|=n}w_{\mathbf k}\leq w_{\mathbf0}u_n.
                                                               \tag{21.7}
\]
Summing (21.7) for \(n>L\) gives the normalization bound.  Lemma
\ref{lem:v18v10-envelopes} gives
\[
\begin{split}
 |s-s_L|
 &\leq\sum_{n>L}\sum_{|\mathbf k|=n}
 w_{\mathbf k}|g(\mathbf k)|\\
 &\leq w_{\mathbf0}G_\star
 \sum_{n>L}\sqrt{\alpha+n}\,u_n=\eta_L,             \tag{21.8}
\end{split}
\]
which proves the score bound.
\end{proof}

\begin{lemma}[A robust rank-one lower bound]
\label{lem:v18v11-rank-one}
Let \(x,y\in\mathbb R^m\), \(0<\delta<1\), and let
\(0<Z_-\leq Z\leq Z_+\).  If \(|y|\leq\eta\), then
\[
 \frac1Z(x+y)(x+y)^{\mathsf T}
 \succeq
 \frac{1-\delta}{Z_+}xx^{\mathsf T}
 -\frac{\delta^{-1}-1}{Z_-}\eta^2I_m.              \tag{21.9}
\]
\end{lemma}

\begin{proof}
For every \(v\in\mathbb R^m\), Young's inequality gives
\[
\begin{split}
 [v\cdot(x+y)]^2
 &\geq(1-\delta)(v\cdot x)^2
 -(\delta^{-1}-1)(v\cdot y)^2\\
 &\geq(1-\delta)(v\cdot x)^2
 -(\delta^{-1}-1)\eta^2|v|^2.                     \tag{21.10}
\end{split}
\]
Divide the positive term by the upper bound \(Z_+\) and the negative term
by the lower bound \(Z_-\).  The resulting quadratic-form inequality is
(21.9).
\end{proof}

\begin{theorem}[Mean-retaining finite star certificate]
\label{thm:v18v11-mean-certificate}
If there exist an integer \(L\geq0\) and \(0<\delta<1\) such that
\[
\begin{split}
 \mathcal Q_L
 +\frac{1-\delta}{Z_L^+}s_Ls_L^{\mathsf T}
 \succeq
 \left[
 \varepsilon_L+
 \frac{\delta^{-1}-1}{W_L}\eta_L^2
 \right]I_m,                                        \tag{21.11}
\end{split}
\]
then
\[
 \nabla_d^2F_\star(d)\succeq0.                     \tag{21.12}
\]
Every term in (21.11) is a finite sum of radial integrals plus an explicit
positive scalar series tail.
\end{theorem}

\begin{proof}
Equation (20.20) is
\[
 Z_0\nabla_d^2F_\star
 =
 \sum_{\mathbf k}w_{\mathbf k}
 \left(H(\mathbf k)-g(\mathbf k)g(\mathbf k)^{\mathsf T}\right)
 +\frac1{Z_0}ss^{\mathsf T}.                        \tag{21.13}
\]
As in V10, positivity of every \(H(\mathbf k)\) and the score envelope give
\[
\begin{split}
 \sum_{\mathbf k}w_{\mathbf k}
 \left(H(\mathbf k)-g(\mathbf k)g(\mathbf k)^{\mathsf T}\right)
 \succeq \mathcal Q_L-\varepsilon_LI_m.             \tag{21.14}
\end{split}
\]
Apply Lemma \ref{lem:v18v11-rank-one} with
\[
 x=s_L,\quad y=s-s_L,\quad
 Z_-=W_L,\quad Z_+=Z_L^+,\quad\eta=\eta_L.          \tag{21.15}
\]
Lemma \ref{lem:v18v11-tail-data} verifies its hypotheses and yields
\[
 \frac1{Z_0}ss^{\mathsf T}
 \succeq
 \frac{1-\delta}{Z_L^+}s_Ls_L^{\mathsf T}
 -\frac{\delta^{-1}-1}{W_L}\eta_L^2I_m.             \tag{21.16}
\]
Combine (21.13)--(21.16).  Condition (21.11) makes the resulting lower bound
positive semidefinite.  Since \(Z_0>0\), (21.12) follows.
\end{proof}

\begin{corollary}[Strict improvement over the V10 certificate]
\label{cor:v18v11-improvement}
If the V10 condition
\[
 \mathcal Q_L\succeq\varepsilon_LI_m
                                                               \tag{21.17}
\]
holds with a positive spectral margin, then (21.11) holds for some
\(\delta\) sufficiently close to \(1\).  Conversely, (21.11) may hold even
when (21.17) fails, because it retains a positive rank-one part of the exact
mean.
\end{corollary}

\begin{proof}
As \(\delta\uparrow1\), both the retained rank-one coefficient
\(1-\delta\) and the new error coefficient \(\delta^{-1}-1\) tend to zero.
A positive margin in (21.17) therefore absorbs the latter for all
sufficiently large \(\delta<1\).  The second statement follows directly
from the additional positive term in (21.11); it is a possibility, not a
claim that a particular parameter card has already passed.
\end{proof}

\section{V11 consequence and next finite task}
\label{sec:v18v11-next}

The finite certificate now preserves a controlled fraction of every signed
term in the exact Hessian identity.  This is the appropriate form for a
rational or interval implementation: compute the finite radial integrals
with outward error bounds, evaluate (21.11) for a finite set of
\(\delta\)'s, and use the explicit \(u_n\) tail.  The companion audit
reinforces that the finite angular box is not by itself a continuum or
all-parameter result.

The next step should choose a declared compact parameter card from the
canonical Higgs regulator and construct an interval verifier for (21.11).
If the card fails, the returned smallest eigenvector and the separated tail
budgets will identify whether to enlarge \(L\), sharpen the total-birth
bound, or test the exact Hessian for a genuine negative direction.

\begin{assessment}[V11 acquisition and boundary]
\label{ass:v18v11}
V11 records the compulsory companion audit and strengthens the V10
finite-star certificate by retaining a rigorously controlled part of the
positive full-score mean.  No canonical parameter card has yet been supplied
or certified.  Tree-depth propagation, cycles, cofinal cutoff removal,
physical reconstruction, and the Einstein continuum remain open.
\end{assessment}

\begin{record}[V11 append record]
\label{rec:v18v11}
V11 is appended to the validated V10 whose installed SHA--256 was
\texttt{6F44DB391A9CED3D259B0DDDED93D38090C5EB90BE6177262F838C96D708AA57}.
The audit used the exact companion fingerprints in (21.2).  After
installation only V11 is the active numeric SM note; the frozen
\texttt{V100\_f17} archive is unchanged.
\end{record}

\section{V12: source parameter audit and a compact-box certificate}
\label{sec:v18v12}

V11 requested a canonical numerical Higgs card.  The supplied sources do not
define one.  This section records that scope before doing mathematics, then
replaces the unavailable point card by a rigorous compact-parameter
certificate.  Explicit radial moment envelopes control both interpolation
between certified parameter centres and the common angular tail.

\begin{audit}[The supplied action leaves the Higgs coefficients physical]
\label{aud:v18v12-source-card}
The Grand Tensor manuscript defines its regulated classical Standard Model
action for positive gauge couplings and positive Higgs quartic
\(\lambda_H\), with \(v_H\geq0\).  Its finite gauge--Higgs shell likewise
uses independent coefficients \(\kappa_H,m_H^2,\lambda_H\).  The accompanying
scope statement says that these numerical coefficients remain physical
parameters and are not fixed by the structural representation packet.

The exact filename of the originally supplied V30 SM--Einstein programme is
not present in the current Downloads state.  A newer matching V35 file is
present.  Its displayed tuple
\[
 (\beta_g,\kappa_H,\lambda_4,m^2,\ldots)
 =(3/5,2/5,7/20,1/2,\ldots)                         \tag{22.1}
\]
is explicitly the \(N=2\), 81-record representative of a finite obstruction
test.  It is not identified there as a cofinal physical Standard Model
parameter card.

Accordingly, no numerical \((\lambda_i,\chi_i,d_i)\) values are imported.
Doing so would turn an example into an unstated physical assumption.  The
correct output at this stage is a theorem uniform on any declared compact
positive box.
\end{audit}

\subsection{Uniform radial moment envelopes}

For \(s>0\), \(a,\lambda>0\), write
\[
 J_s(a,\lambda)=\int_0^\infty
 t^{s-1}e^{-\lambda t^2-at}\,dt,\qquad
 M_{s,r}=\frac{J_{s+r}}{J_s}
 =\mathbb E_s t^r.                                  \tag{22.2}
\]
Use the rising factorial
\[
 (s)_r=s(s+1)\cdots(s+r-1),\qquad (s)_0=1.          \tag{22.3}
\]

\begin{lemma}[All polynomial radial moments]
\label{lem:v18v12-moments}
For every integer \(r\geq0\),
\[
 0<M_{s,r}\leq
 \left(\frac{(s)_r}{(2\lambda)^r}\right)^{1/2}.      \tag{22.4}
\]
Moreover,
\[
\begin{split}
 \partial_a M_{s,r}
 &=-\operatorname{Cov}_s(t^r,t),\\
 \partial_\lambda M_{s,r}
 &=-\operatorname{Cov}_s(t^r,t^2),                 \tag{22.5}
\end{split}
\]
and hence
\[
\begin{split}
 |\partial_aM_{s,r}|
 &\leq M_{s,r+1}+M_{s,r}M_{s,1},\\
 |\partial_\lambda M_{s,r}|
 &\leq M_{s,r+2}+M_{s,r}M_{s,2}.                  \tag{22.6}
\end{split}
\]
\end{lemma}

\begin{proof}
Factor the moment ratio:
\[
 M_{s,r}=
 \prod_{j=0}^{r-1}\frac{J_{s+j+1}}{J_{s+j}}.
                                                               \tag{22.7}
\]
Lemma \ref{lem:v18v10-radial-bound} bounds the \(j\)-th factor by
\(\sqrt{(s+j)/(2\lambda)}\), which proves (22.4).  Differentiation under the
integral gives (22.5).  Finally,
\[
 |\operatorname{Cov}(X,Y)|
 \leq\mathbb E|XY|+\mathbb E|X|\,\mathbb E|Y|,
\]
and all variables here are nonnegative.  This yields (22.6).
\end{proof}

\begin{corollary}[Relative enclosure of every radial integral]
\label{cor:v18v12-J-enclosure}
Let
\[
 a\in[a_-,a_+],\qquad
 \lambda\in[\lambda_-,\lambda_+],\qquad
 a_-,\lambda_->0.                                   \tag{22.8}
\]
For two points \((a,\lambda)\) and \((a',\lambda')\) in this rectangle,
\[
\begin{split}
 \left|\log\frac{J_s(a',\lambda')}{J_s(a,\lambda)}\right|
 \leq{}&
 \sqrt{\frac{s}{2\lambda_-}}\,|a'-a|\\
 &+\frac{\sqrt{s(s+1)}}{2\lambda_-}\,
 |\lambda'-\lambda|.                                \tag{22.9}
\end{split}
\]
\end{corollary}

\begin{proof}
The logarithmic derivatives are
\[
 \partial_a\log J_s=-M_{s,1},\qquad
 \partial_\lambda\log J_s=-M_{s,2}.                 \tag{22.10}
\]
Use (22.4) with \(r=1,2\), then integrate along the line segment joining the
two parameter points.
\end{proof}

\subsection{A uniform angular tail on a parameter box}

Let \(\mathcal B\) be a rectangular star-parameter box on which
\[
\begin{split}
 &0<\lambda_{0,-}\leq\lambda_o,\qquad
 0<\lambda_{i,-}\leq\lambda_i,\\
 &0<\chi_{i,-}\leq\chi_i\leq\chi_{i,+},\qquad
 d_i\in[d_{i,-},d_{i,+}].                           \tag{22.11}
\end{split}
\]
Define the common birth constant
\[
 C_{\mathcal B}=
 \frac1{\sqrt\alpha}
 \sum_{i=1}^m
 \frac{\chi_{i,+}^2}
 {2\sqrt{\lambda_{0,-}\lambda_{i,-}}},              \tag{22.12}
\]
and
\[
 u_n^{\mathcal B}=
 \frac{C_{\mathcal B}^n\sqrt{(\alpha)_n}}{n!}.       \tag{22.13}
\]

\begin{proposition}[Uniform total-degree tail]
\label{prop:v18v12-uniform-tail}
For every parameter point in \(\mathcal B\),
\[
 \sum_{|\mathbf k|=n}w_{\mathbf k}
 \leq w_{\mathbf0}u_n^{\mathcal B}.                 \tag{22.14}
\]
For \(p\in\{0,\tfrac12,1\}\), put
\[
 v_{n,p}=(\alpha+n)^p u_n^{\mathcal B}.             \tag{22.15}
\]
If \(R\geq L+1\) and
\[
 r_{R,p}:=
 \sup_{n\geq R}
 C_{\mathcal B}\frac{\sqrt{\alpha+n}}{n+1}
 \left(\frac{\alpha+n+1}{\alpha+n}\right)^p<1,      \tag{22.16}
\]
then
\[
 \sum_{n\geq R}v_{n,p}
 \leq\frac{v_{R,p}}{1-r_{R,p}}.                    \tag{22.17}
\]
Thus every tail in (21.3) has an explicit finite-sum plus geometric
remainder bound uniform on \(\mathcal B\).
\end{proposition}

\begin{proof}
The proof of Theorem \ref{thm:v18v10-tail} uses the parameter-dependent
constant \(C_\star\).  Equation (22.11) gives
\(C_\star\leq C_{\mathcal B}\), so the same induction proves (22.14).
The consecutive ratio of (22.15) is the expression inside the supremum in
(22.16).  Once it is at most \(r_{R,p}<1\), comparison with a geometric
series gives (22.17).
\end{proof}

\subsection{Finite covering of a compact card}

For fixed angular cutoff \(L\) and \(0<\delta<1\), define the V11 certificate
matrix
\[
\begin{split}
 \mathscr C_{L,\delta}(\vartheta)
 :={}&\mathcal Q_L
 +\frac{1-\delta}{Z_L^+}s_Ls_L^{\mathsf T}\\
 &-\left[
 \varepsilon_L+
 \frac{\delta^{-1}-1}{W_L}\eta_L^2
 \right]I_m,                                        \tag{22.18}
\end{split}
\]
where \(\vartheta\) denotes all star parameters.  In (22.18), the tails may
be replaced by their uniform upper bounds from Proposition
\ref{prop:v18v12-uniform-tail}.

\begin{theorem}[Computable Lipschitz modulus on a finite angular box]
\label{thm:v18v12-lipschitz}
Fix a compact positive parameter box \(\mathcal B\), an angular cutoff
\(L\), and \(0<\delta<1\).  There is an explicitly computable finite constant
\(\mathscr L_{\mathcal B,L,\delta}\) such that
\[
 \|\mathscr C_{L,\delta}(\vartheta)
 -\mathscr C_{L,\delta}(\vartheta')\|_{\mathrm{op}}
 \leq
 \mathscr L_{\mathcal B,L,\delta}
 \|\vartheta-\vartheta'\|_\infty                   \tag{22.19}
\]
for all \(\vartheta,\vartheta'\in\mathcal B\).

The constant is obtained using only the endpoints in (22.11), the finite
shapes
\[
 \alpha\leq s_i\leq\alpha+L,\qquad
 \alpha\leq s_0\leq\alpha+L,                        \tag{22.20}
\]
the moment envelopes (22.4)--(22.6), and elementary differentiation of
\(a_i=\chi_ie^{-d_i}\), \(b_i=\chi_ie^{d_i}\).
\end{theorem}

\begin{proof}
Every entry of \(w_{\mathbf k}\), \(g(\mathbf k)\), and
\(H(\mathbf k)\) is a finite algebraic combination of
\[
 J_{j,s},\quad
 \frac{J_{j,s+1}}{J_{j,s}},\quad
 \frac{J_{j,s+2}}{J_{j,s}},\quad
 a_i,\quad b_i,\quad A=\sum_i b_i,                 \tag{22.21}
\]
with the shapes restricted by (22.20).  Corollary
\ref{cor:v18v12-J-enclosure} gives two-sided positive relative enclosures for
every \(J_{j,s}\).  Lemma \ref{lem:v18v12-moments}, with \(r\leq4\), bounds
the first derivatives of every ratio in (22.21).  On the compact box,
\[
 |\partial_{\chi_i}a_i|,\ |\partial_{\chi_i}b_i|,
 |\partial_{d_i}a_i|,\ |\partial_{d_i}b_i|
                                                               \tag{22.22}
\]
have explicit endpoint bounds.

Finite sums and products preserve bounded first derivatives.  Also
\(W_L\geq w_{\mathbf0}>0\), while (22.9) gives an explicit positive lower
bound for \(w_{\mathbf0}\) relative to its value at any chosen centre.
Therefore the divisions by \(W_L\) and \(Z_L^+\) have bounded derivatives.
The uniform tail series and their parameter derivatives are bounded by the
same ratio argument as (22.17), after one polynomial factor in \(n\) is
added.  Taking the sum of the resulting entrywise derivative bounds and
using \(\|M\|_{\mathrm{op}}\leq\|M\|_{\mathrm F}\) produces the declared
finite \(\mathscr L_{\mathcal B,L,\delta}\).  The mean-value theorem proves
(22.19).
\end{proof}

\begin{corollary}[Finite compact-card certification]
\label{cor:v18v12-cover}
Cover \(\mathcal B\) by finitely many closed \(\ell^\infty\)-balls
\(B(\vartheta_j,h_j)\).  Suppose outward interval evaluation at every centre
proves
\[
 \lambda_{\min}
 \bigl(\mathscr C_{L,\delta}(\vartheta_j)\bigr)
 \geq\mu_j
 >\mathscr L_{\mathcal B,L,\delta}h_j.              \tag{22.23}
\]
Then
\[
 \nabla_d^2F_\star(d;\vartheta)\succeq0
 \qquad\text{for every }\vartheta\in\mathcal B.     \tag{22.24}
\]
\end{corollary}

\begin{proof}
Weyl's eigenvalue inequality and (22.19) give, for
\(\vartheta\in B(\vartheta_j,h_j)\),
\[
 \lambda_{\min}\mathscr C_{L,\delta}(\vartheta)
 \geq\mu_j-\mathscr L_{\mathcal B,L,\delta}h_j>0.   \tag{22.25}
\]
The mean-retaining certificate, Theorem
\ref{thm:v18v11-mean-certificate}, then gives (22.24).
\end{proof}

\section{V12 consequence and next target}
\label{sec:v18v12-next}

The missing numerical card is now separated from the analytic issue.  A
declared compact box can be certified by finitely many outward interval
calculations, with no interpolation gap and no uncontrolled angular tail.
The source manuscripts do not select that box, so this note does not claim
that a particular physical coefficient set has passed.

The next useful analytic step is independent of that choice: improve the
finite certificate by retaining the positive component Hessian in the
angular tail.  V10--V11 discarded it completely.  A lower bound on the
component floor from (19.14) can offset part of the score tail and reduce the
angular cutoff required by (22.23).

\begin{assessment}[V12 acquisition and boundary]
\label{ass:v18v12}
V12 proves uniform radial moment and parameter-derivative envelopes, a common
subfactorial angular tail over compact positive boxes, and a finite covering
criterion for the V11 matrix certificate.  It also records why the available
finite \(N=2\) tuple is not a canonical continuum Higgs card.  No specific
physical box has been certified, and tree depth, cycles, continuum
reconstruction, and the Einstein coupling remain open.
\end{assessment}

\begin{record}[V12 append record]
\label{rec:v18v12}
V12 is appended to the validated V11 whose installed SHA--256 was
\texttt{7A943E6A5A646E76521210B220948A453909E44D7FD7D5FE77DB220088B80F57}.
After installation only V12 is the active numeric SM note; the frozen
\texttt{V100\_f17} archive is unchanged.
\end{record}

\section{V13: a translation-invariant pairwise angular certificate}
\label{sec:v18v13}

V12 proposed retaining the positive component Hessian in the omitted angular
tail.  An upper bound on tail mass alone cannot supply a lower bound on that
positive contribution.  The correct upstream move is to rewrite the score
covariance as an independent-pair difference.  The component score has
uniformly bounded increments in every angular coordinate, so the tail can be
estimated without paying for the absolute size of a common score.

Retain the exact star mixture of V9.  For
\[
 s_0=\alpha+|\mathbf k|,\qquad s_i=\alpha+k_i,
\]
write
\[
\begin{split}
 \Delta_0(\mathbf k)&=
 m_{0,s_0+1}-m_{0,s_0},\\
 \Delta_i(\mathbf k)&=
 m_{i,s_i+1}-m_{i,s_i}.                             \tag{23.1}
\end{split}
\]

\begin{lemma}[Uniform score increments]
\label{lem:v18v13-score-increments}
For every \(i,j\),
\[
 g_i(\mathbf k+\mathbf e_j)-g_i(\mathbf k)
 =b_i\Delta_0(\mathbf k)
 -\delta_{ij}a_i\Delta_i(\mathbf k).                \tag{23.2}
\]
Moreover,
\[
 0<\Delta_0<\frac1A,\qquad
 0<\Delta_i<\frac1{a_i},                            \tag{23.3}
\]
and therefore
\[
 |g_i(\mathbf k+\mathbf e_j)-g_i(\mathbf k)|
 \leq\frac{b_i}{A}+\delta_{ij}.                     \tag{23.4}
\]
If
\[
 L_g=\max_{1\leq j\leq m}
 \left[\sum_{i=1}^m
 \left(\frac{b_i}{A}+\delta_{ij}\right)^2\right]^{1/2},
                                                               \tag{23.5}
\]
then
\[
 |g(\mathbf k)-g(\mathbf l)|
 \leq L_g\,|\mathbf k-\mathbf l|_1
 \leq L_g\bigl(|\mathbf k|+|\mathbf l|\bigr).       \tag{23.6}
\]
\end{lemma}

\begin{proof}
Increasing \(k_j\) raises the root shape \(s_0\) by one and raises only the
\(j\)-th leaf shape.  Apply this observation to
\(g_i=b_im_0-a_im_i\) in (19.8) to get (23.2).  The sharp increment bound
(7.8), used with root quadratic coefficient \(A\) and leaf coefficient
\(a_i\), gives (23.3).  The triangle inequality gives (23.4).

Join \(\mathbf k\) to \(\mathbf l\) by a nearest-neighbor coordinate path of
length \(|\mathbf k-\mathbf l|_1\).  Summing the vector form of (23.4) along
that path proves the first inequality in (23.6); the second is elementary.
\end{proof}

\begin{lemma}[Pairwise covariance identity]
\label{lem:v18v13-pairwise}
If \(\mathbf K,\mathbf K'\) are independent with law \(p_d\), then
\[
 \operatorname{Cov}(g(\mathbf K))
 =\frac12\mathbb E
 \left[
 (g(\mathbf K)-g(\mathbf K'))
 (g(\mathbf K)-g(\mathbf K'))^{\mathsf T}
 \right].                                           \tag{23.7}
\]
\end{lemma}

\begin{proof}
Expand the right side.  The two diagonal expectations both equal
\(\mathbb E[gg^{\mathsf T}]\), while independence makes each cross term
\((\mathbb Eg)(\mathbb Eg)^{\mathsf T}\).
\end{proof}

\subsection{Explicit pair-tail data}

Use the unnormalized weights \(w_{\mathbf k}\) and the envelope \(u_n\) from
(20.7).  Define, for \(r=0,1,2\),
\[
 U_r=\sum_{n\geq0}n^ru_n,\qquad
 T_{r,L}=\sum_{n>L}n^ru_n,                          \tag{23.8}
\]
with \(0^0=1\) in \(U_0\).  These series converge by Theorem
\ref{thm:v18v10-tail} and admit geometric remainder bounds by Proposition
\ref{prop:v18v12-uniform-tail}.  Let
\[
\begin{split}
 \mathcal H_L&=
 \sum_{|\mathbf k|\leq L}w_{\mathbf k}H(\mathbf k),\\
 \mathcal V_L&=
 \frac12
 \sum_{\substack{|\mathbf k|\leq L\\|\mathbf l|\leq L}}
 w_{\mathbf k}w_{\mathbf l}
 (g(\mathbf k)-g(\mathbf l))
 (g(\mathbf k)-g(\mathbf l))^{\mathsf T},           \tag{23.9}\\
 \mathcal E_L&=
 \frac{L_g^2w_{\mathbf0}^2}{W_L^2}
 \left(T_{2,L}U_0+2T_{1,L}U_1+T_{0,L}U_2\right).
                                                               \tag{23.10}
\end{split}
\]

\begin{theorem}[Pairwise finite-star certificate]
\label{thm:v18v13-pair-certificate}
If
\[
 \frac{\mathcal H_L}{Z_L^+}
 -\frac{\mathcal V_L}{W_L^2}
 \succeq\mathcal E_L I_m,                           \tag{23.11}
\]
then
\[
 \nabla_d^2F_\star(d)\succeq0.                     \tag{23.12}
\]
The certificate is invariant under replacing every score
\(g(\mathbf k)\) by \(g(\mathbf k)+c\) for a fixed vector \(c\).
\end{theorem}

\begin{proof}
Let \(Z_0=\sum_{\mathbf k}w_{\mathbf k}\).  Theorems
\ref{thm:v18v9-discrete-gate} and Lemma
\ref{lem:v18v13-pairwise} give the exact identity
\[
\begin{split}
 \nabla_d^2F_\star
 ={}&\frac1{Z_0}\sum_{\mathbf k}w_{\mathbf k}H(\mathbf k)\\
 &-\frac1{2Z_0^2}
 \sum_{\mathbf k,\mathbf l}w_{\mathbf k}w_{\mathbf l}
 \Delta g_{\mathbf k\mathbf l}
 \Delta g_{\mathbf k\mathbf l}^{\mathsf T},
                                                               \tag{23.13}
\end{split}
\]
where \(\Delta g_{\mathbf k\mathbf l}=g(\mathbf k)-g(\mathbf l)\).
Drop the positive Hessian terms with \(|\mathbf k|>L\).  Since
\(Z_0\leq Z_L^+\), the retained positive part is bounded below by
\(\mathcal H_L/Z_L^+\).  For the negative pair sum with both indices in the
finite box, \(Z_0\geq W_L\) gives the upper bound
\(\mathcal V_L/W_L^2\).

It remains to bound pairs with at least one index outside.  Symmetry and the
factor \(1/2\) permit a union bound by the sum with
\(|\mathbf k|>L\) and arbitrary \(\mathbf l\).  Lemma
\ref{lem:v18v13-score-increments} gives
\[
 \Delta g_{\mathbf k\mathbf l}
 \Delta g_{\mathbf k\mathbf l}^{\mathsf T}
 \preceq
 L_g^2(|\mathbf k|+|\mathbf l|)^2I_m.               \tag{23.14}
\]
The unnormalized tail bound (20.22) then yields
\[
\begin{split}
 &\sum_{\substack{|\mathbf k|>L\\\mathbf l\in\mathbb N^m}}
 w_{\mathbf k}w_{\mathbf l}
 (|\mathbf k|+|\mathbf l|)^2\\
 &\qquad\leq w_{\mathbf0}^2
 \left(T_{2,L}U_0+2T_{1,L}U_1+T_{0,L}U_2\right).
                                                               \tag{23.15}
\end{split}
\]
Divide by \(Z_0^2\geq W_L^2\) and combine
(23.13)--(23.15).  The result is
\[
 \nabla_d^2F_\star
 \succeq
 \frac{\mathcal H_L}{Z_L^+}
 -\frac{\mathcal V_L}{W_L^2}
 -\mathcal E_LI_m.                                  \tag{23.16}
\]
Condition (23.11) proves (23.12).  Every occurrence of the score in
(23.9) and (23.13) is through a difference, proving translation invariance.
\end{proof}

\begin{proposition}[Why this is a different loss from V10]
\label{prop:v18v13-comparison}
The V10 score-tail error is controlled by absolute squares
\[
 |g(\mathbf k)|^2\leq G_\star^2(\alpha+|\mathbf k|).
                                                               \tag{23.17}
\]
The V13 tail is controlled by the increment constant \(L_g\) and pair
distances.  Therefore an arbitrarily large common vector added to all
component scores changes the V10 bound but leaves (23.11) unchanged.
Neither certificate uniformly dominates the other for all parameter values
and cutoffs.
\end{proposition}

\begin{proof}
The first statement is the definition of the two errors.  Translation
invariance of V13 was proved in Theorem
\ref{thm:v18v13-pair-certificate}.  The V10 error uses
\(|g+c|^2\), which can be made arbitrarily large with \(c\), while the V13
pair-distance envelope contains a higher degree moment.  These opposite
features prevent a parameter-free dominance statement.
\end{proof}

\section{V13 consequence and next direction}
\label{sec:v18v13-next}

The exact finite-star gate now has two complementary finite certificates:
V11 retains a controlled rank-one part of the score mean, while V13 retains
the cancellation of any common score exactly.  Both use the same positive
angular weights and explicit subfactorial tail.  A parameter-box verifier
should evaluate both and accept a cell when either certificate has a positive
outward margin.

The remaining analytic obstacle is no longer truncation.  It is whether a
declared order-one coupling box satisfies one of the exact-star
certificates.  In the absence of a physically selected numerical card, the
next theorem-level advance should address tree depth: derive how a certified
star block becomes an effective radial message while preserving a
quantitative matrix margin for descendant edge parameters.

\begin{assessment}[V13 acquisition and boundary]
\label{ass:v18v13}
V13 proves a uniform component-score increment bound and a
translation-invariant pairwise finite certificate for the exact star
Hessian.  It removes the absolute common-score loss of V10.  No parameter
box has yet been certified, and propagation through tree depth, cycle
holonomy, cofinal cutoff removal, reconstruction, and the Einstein coupling
remain open.
\end{assessment}

\begin{record}[V13 append record]
\label{rec:v18v13}
V13 is appended to the validated V12 whose installed SHA--256 was
\texttt{C0C4E824FF2ADF57CE05B42B1B15EE9147251B7422960516742E7DC3335866D3}.
After installation only V13 is the active numeric SM note; the frozen
\texttt{V100\_f17} archive is unchanged.
\end{record}

\section{V14: recursive rank-one Schur transport through tree depth}
\label{sec:v18v14}

The exact star certificates control one branching layer.  To advance through
tree depth, this section treats an internal vertex whose child messages
already depend on all conformal parameters below them.  Radiality forces
every field gradient of a descendant score into the same one-dimensional
radial direction.  The Brascamp--Lieb loss is therefore rank one in the full
descendant parameter space, independently of the number of edges below the
vertex.

Let \(x\) be a nonroot vertex with parent boundary field
\(y\in\mathbb R^N\).  Write \(z\in\mathbb R^N\) for the field at \(x\) and
\(q=|z|^2\), \(Q=|y|^2\).  Its outgoing edge has parameter \(d_0\) and
\[
 a_0=\chi_0e^{-d_0},\qquad b_0=\chi_0e^{d_0}.       \tag{24.1}
\]
Suppose the children \(j=1,\ldots,r\) send radial messages
\[
 \ell_j(\theta_j,q),                                \tag{24.2}
\]
where the vectors \(\theta_j\) collect disjoint descendant conformal
parameters.  Assume each \(\ell_j\) is twice differentiable, convex and
nondecreasing in \(q\).  Define
\[
\begin{split}
 \mathcal S(\eta,z;y)
 :={}&\lambda_xq^2+\sum_{j=1}^r\ell_j(\theta_j,q)\\
 &+a_0q+b_0Q-2\chi_0z\cdot y,                       \tag{24.3}\\
 \eta&=(d_0,\theta_1,\ldots,\theta_r),              \tag{24.4}\\
 \mathcal L_x(\eta,y)
 :={}&-\log\int_{\mathbb R^N}e^{-\mathcal S(\eta,z;y)}\,dz.
                                                               \tag{24.5}
\end{split}
\]

\begin{definition}[Direct curvature and radial score vector]
\label{def:v18v14-blocks}
For \(j\geq1\), put
\[
 H_j(\theta_j,q)=\nabla_{\theta_j}^2\ell_j(\theta_j,q),
 \qquad
 c_j(\theta_j,q)=\nabla_{\theta_j}\partial_q\ell_j(\theta_j,q).
                                                               \tag{24.6}
\]
Set
\[
\begin{split}
 T_0(q,Q)&=a_0q+b_0Q,\\
 D(\eta,q,Q)&=
 \operatorname{diag}\bigl(T_0,H_1,\ldots,H_r\bigr),\\
 c(\eta,q)&=
 \begin{pmatrix}-a_0&c_1^{\mathsf T}&\cdots&c_r^{\mathsf T}\end{pmatrix}^{\mathsf T},
                                                               \tag{24.7}\\
 R(\eta,q)&=
 2a_0+12\lambda_xq+
 \sum_{j=1}^r
 \left(2\partial_q\ell_j+4q\partial_q^2\ell_j\right).
                                                               \tag{24.8}
\end{split}
\]
\end{definition}

\begin{theorem}[Recursive Schur transport]
\label{thm:v18v14-transport}
The conditional law in (24.5) is strictly log-concave in \(z\), and
\[
 \nabla_\eta^2\mathcal L_x(\eta,y)
 \succeq
 \mathbb E_{\eta,y}\left[
 D(\eta,q,Q)-\frac{4q}{R(\eta,q)}c(\eta,q)c(\eta,q)^{\mathsf T}
 \right].                                           \tag{24.9}
\]
Consequently, the pointwise condition
\[
 D(\eta,q,Q)-\frac{4q}{R(\eta,q)}cc^{\mathsf T}
 \succeq0
 \qquad(q\geq0)                                     \tag{24.10}
\]
implies
\[
 \nabla_\eta^2\mathcal L_x(\eta,y)\succeq0.         \tag{24.11}
\]
\end{theorem}

\begin{proof}
For fixed parameters and \(y\), the Hessian of
\(z\mapsto\mathcal S(\eta,z;y)\) is positive definite.  Indeed, its
tangential eigenvalue is
\[
 2a_0+4\lambda_xq+
 2\sum_j\partial_q\ell_j>0,
                                                               \tag{24.12}
\]
and its radial eigenvalue is \(R(\eta,q)>0\).

Differentiate (24.5) twice in \(\eta\):
\[
 \nabla_\eta^2\mathcal L_x
 =
 \mathbb E_{\eta,y}\nabla_\eta^2\mathcal S
 -\operatorname{Cov}_{\eta,y}(\nabla_\eta\mathcal S).
                                                               \tag{24.13}
\]
There are no direct mixed second derivatives between disjoint child
parameter blocks or between a child block and \(d_0\), so
\[
 \nabla_\eta^2\mathcal S=D.                         \tag{24.14}
\]
For the outgoing edge,
\[
 \partial_{d_0}\mathcal S=-a_0q+b_0Q,\qquad
 \nabla_z\partial_{d_0}\mathcal S=-2a_0z.           \tag{24.15}
\]
For a descendant block,
\[
 \nabla_z\nabla_{\theta_j}\mathcal S
 =2z\,c_j^{\mathsf T}.                              \tag{24.16}
\]
Thus the complete field Jacobian of the parameter score is
\[
 \nabla_z\nabla_\eta\mathcal S=2z\,c^{\mathsf T}.   \tag{24.17}
\]

The matrix Brascamp--Lieb covariance inequality says
\[
 \operatorname{Cov}(\nabla_\eta\mathcal S)
 \preceq
 \mathbb E\left[
 (\nabla_z\nabla_\eta\mathcal S)^{\mathsf T}
 (\nabla_z^2\mathcal S)^{-1}
 (\nabla_z\nabla_\eta\mathcal S)\right].            \tag{24.18}
\]
The vector \(z\) lies in the radial eigenspace of
\(\nabla_z^2\mathcal S\), with eigenvalue \(R\).  Substitution of (24.17)
into (24.18) therefore gives
\[
 \operatorname{Cov}(\nabla_\eta\mathcal S)
 \preceq
 \mathbb E\frac{4q}{R}cc^{\mathsf T}.               \tag{24.19}
\]
Combining (24.13)--(24.14) with (24.19) proves (24.9), and (24.10) gives
(24.11).
\end{proof}

\begin{corollary}[Scalar inverse form of the recursive gate]
\label{cor:v18v14-inverse-gate}
Suppose \(T_0>0\) and every \(H_j\succ0\).  Then (24.10) is equivalent to
\[
 \frac{4q}{R(\eta,q)}
 \left[
 \frac{a_0^2}{T_0(q,Q)}
 +\sum_{j=1}^r
 c_j(\theta_j,q)^{\mathsf T}
 H_j(\theta_j,q)^{-1}c_j(\theta_j,q)
 \right]\leq1.                                      \tag{24.20}
\]
At \(q=0\), the rank-one correction vanishes and the condition holds by
continuity whenever the direct block is positive semidefinite.
\end{corollary}

\begin{proof}
For a positive definite block diagonal matrix \(D\),
\[
 D-\rho cc^{\mathsf T}\succeq0
 \quad\Longleftrightarrow\quad
 \rho\,c^{\mathsf T}D^{-1}c\leq1.                  \tag{24.21}
\]
Insert \(\rho=4q/R\) and the blocks from (24.7).
\end{proof}

\begin{corollary}[A depth-recursive matrix budget]
\label{cor:v18v14-budget}
Assume that for each child block there is a positive function
\(\Gamma_j(q)\) such that
\[
 H_j(\theta_j,q)\succ0,\qquad
 c_j^{\mathsf T}H_j^{-1}c_j\leq\Gamma_j(q).         \tag{24.22}
\]
Then the sufficient recursive condition is
\[
 4q\left[
 \frac{a_0^2}{a_0q+b_0Q}
 +\sum_{j=1}^r\Gamma_j(q)
 \right]
 \leq
 2a_0+12\lambda_xq+
 \sum_{j=1}^r
 \left(2\ell_{j,q}+4q\ell_{j,qq}\right).            \tag{24.23}
\]
All terms in (24.23) are local to the vertex, its parent boundary radius,
and the child message certificates.
\end{corollary}

\begin{proof}
Replace each exact quadratic form in (24.20) by its upper bound in (24.22)
and multiply by \(R\).
\end{proof}

\begin{remark}[The zero parent-boundary separator]
\label{rem:v18v14-zero-parent}
When \(Q=0\) and \(q>0\), the outgoing term in (24.20) is
\[
 \frac{a_0^2}{T_0}=\frac{a_0}{q}.                  \tag{24.24}
\]
Its contribution after multiplication by \(4q/R\) is \(4a_0/R\).  Since
\(R\) begins with \(2a_0\), the outgoing hopping alone again loses a factor
two.  Quartic and child radial stiffness must supply the difference.  This is
the depth-recursive analogue of the V8 separator.
\end{remark}

\section{V14 consequence and next local problem}
\label{sec:v18v14-next}

The tree-depth issue is now a local recursive matrix budget.  The dimension
of the negative correction does not grow with subtree size: it remains rank
one in parameter space because all conformal score gradients are radial.
The quantities that must be propagated are
\[
 H_j,\qquad c_j,\qquad
 2\ell_{j,q}+4q\ell_{j,qq}.                         \tag{24.25}
\]

What is not yet proved is closure of a uniform inequality of the form
(24.22) under the outgoing transform (24.5).  V15 should derive the mixed
boundary derivative of \(\mathcal L_x\) and determine whether the new
quadratic form \(c_x^{\mathsf T}H_x^{-1}c_x\) is bounded by a local function
with constants independent of subtree depth.

\begin{assessment}[V14 acquisition and boundary]
\label{ass:v18v14}
V14 proves the recursive rank-one Schur transport theorem for an internal
tree vertex and reduces its positivity to the local scalar budget (24.20).
This is the first certificate in the restarted cycle that carries all
descendant edge parameters through one further field integration.  Uniform
closure of the propagated budget, cycles, cofinal cutoff removal,
reconstruction, and the coupled Einstein continuum remain open.
\end{assessment}

\begin{record}[V14 append record]
\label{rec:v18v14}
V14 is appended to the validated V13 whose installed SHA--256 was
\texttt{A560CCBA726CA97FCD2358164EEEF432859294A68BA7F0661E01579C99766722}.
After installation only V14 is the active numeric SM note; the frozen
\texttt{V100\_f17} archive is unchanged.
\end{record}

\section{V15: exact mixed boundary response of a recursive message}
\label{sec:v18v15}

V14 identified the propagated quantity
\(c_x=\nabla_\eta\partial_Q\mathcal L_x\).  This section derives it exactly
as a conditional response covariance.  The formula has one form for
nonzero parent boundary and a separate isotropic limit at zero.  This
separation prevents a spurious \(Q^{-1}\) singularity from entering the
depth recursion.

Use the internal-vertex action (24.3).  Let
\[
 \overline z(\eta,y)=\mathbb E_{\eta,y}z.           \tag{25.1}
\]
Rotational covariance gives a scalar \(\mu(\eta,Q)\) such that
\[
 \overline z(\eta,y)=\mu(\eta,Q)y.                  \tag{25.2}
\]

\begin{proposition}[Boundary slope from the conditional mean]
\label{prop:v18v15-boundary-slope}
The outgoing message is radial,
\[
 \mathcal L_x(\eta,y)=\ell_x(\eta,Q),
\]
and
\[
 \partial_Q\ell_x(\eta,Q)=b_0-\chi_0\mu(\eta,Q).
                                                               \tag{25.3}
\]
\end{proposition}

\begin{proof}
Differentiate (24.5) in \(y\):
\[
 \nabla_y\mathcal L_x
 =\mathbb E_{\eta,y}(2b_0y-2\chi_0z)
 =2(b_0-\chi_0\mu)y.                                \tag{25.4}
\]
For a radial function the left side is \(2(\partial_Q\ell_x)y\), proving
(25.3) for \(y\ne0\).  Continuity gives the value at zero.
\end{proof}

\begin{theorem}[Exact nonzero-boundary mixed response]
\label{thm:v18v15-mixed-response}
Let \(\theta\) be any descendant conformal parameter and let \(d_0\) be the
outgoing-edge parameter.  For \(Q>0\),
\[
\begin{split}
 \partial_\theta\partial_Q\ell_x
 &=\frac{\chi_0}{Q}\,
 y\cdot\operatorname{Cov}_{\eta,y}
 \left(z,\partial_\theta\mathcal S\right),\\
 \partial_{d_0}\partial_Q\ell_x
 &=b_0+\frac{\chi_0}{Q}\,
 y\cdot\operatorname{Cov}_{\eta,y}
 \left(z,\partial_{d_0}\mathcal S\right).           \tag{25.5}
\end{split}
\]
\end{theorem}

\begin{proof}
For any parameter \(\xi\),
\[
 \partial_\xi\mathbb E_{\eta,y}z
 =-\operatorname{Cov}_{\eta,y}
 (z,\partial_\xi\mathcal S).                        \tag{25.6}
\]
Differentiate (25.2) while holding \(y\) fixed and take its scalar product
with \(y\):
\[
 Q\,\partial_\xi\mu
 =-y\cdot\operatorname{Cov}
 (z,\partial_\xi\mathcal S).                        \tag{25.7}
\]
For a descendant parameter, \(b_0\) is constant.  For \(d_0\),
\(\partial_{d_0}b_0=b_0\).  Differentiate (25.3) and use (25.7) to obtain
(25.5).
\end{proof}

\subsection{The finite isotropic limit}

At \(y=0\), let \(\nu_\eta\) be the radial law on \(z\) proportional to
\[
 \exp\left[-\lambda_x|z|^4-a_0|z|^2
 -\sum_{j=1}^r\ell_j(\theta_j,|z|^2)\right]dz.      \tag{25.8}
\]
Put \(t=|z|^2\) and \(m=\mathbb E_{\nu_\eta}t\).

\begin{theorem}[Zero-boundary transfer formula]
\label{thm:v18v15-zero-transfer}
At \(Q=0\),
\[
 \partial_Q\ell_x(\eta,0)
 =b_0-\frac{\chi_0^2}{\alpha}m.                    \tag{25.9}
\]
For a descendant parameter \(\theta\),
\[
 \partial_\theta\partial_Q\ell_x(\eta,0)
 =
 \frac{\chi_0^2}{\alpha}
 \operatorname{Cov}_{\nu_\eta}
 \left(t,\partial_\theta\mathcal S(\eta,z;0)\right),
                                                               \tag{25.10}
\]
while
\[
 \partial_{d_0}\partial_Q\ell_x(\eta,0)
 =
 b_0-\frac{\chi_0^2a_0}{\alpha}
 \operatorname{Var}_{\nu_\eta}(t).                 \tag{25.11}
\]
All three quantities are finite.
\end{theorem}

\begin{proof}
Expand the boundary tilt in (24.3):
\[
 e^{2\chi_0z\cdot y}
 =1+2\chi_0z\cdot y+
 2\chi_0^2(z\cdot y)^2+O(|y|^3).                   \tag{25.12}
\]
Odd terms integrate to zero under the radial law.  Isotropy gives
\[
 \mathbb E_{\nu_\eta}(z\cdot y)^2
 =\frac{Q}{N}\mathbb E_{\nu_\eta}|z|^2
 =\frac{Qm}{2\alpha}.                               \tag{25.13}
\]
Together with the factor \(e^{-b_0Q}\), equations
(25.12)--(25.13) give
\[
 M_x(\eta,y)
 =M_x(\eta,0)
 \left[1-\left(b_0-\frac{\chi_0^2m}{\alpha}\right)Q
 +O(Q^2)\right],                                    \tag{25.14}
\]
which proves (25.9).

For a descendant parameter,
\[
 \partial_\theta m
 =-\operatorname{Cov}_{\nu_\eta}
 \left(t,\partial_\theta\mathcal S\right).          \tag{25.15}
\]
Differentiate (25.9) to obtain (25.10).  For \(d_0\), only
\(a_0=\chi_0e^{-d_0}\) changes in the zero-boundary law, and
\[
 \partial_{d_0}m
 =a_0\operatorname{Var}_{\nu_\eta}(t).              \tag{25.16}
\]
Use \(\partial_{d_0}b_0=b_0\) in (25.9) to obtain (25.11).  Quartic
confinement gives every polynomial moment, proving finiteness.
\end{proof}

\subsection{A response bound away from zero}

Assume the pointwise V14 Schur condition (24.10).  Let
\[
 r_*=\inf_z\lambda_{\min}
 \nabla_z^2\mathcal S(\eta,z;y)>0.                  \tag{25.17}
\]

\begin{proposition}[Nonzero-boundary covariance bound]
\label{prop:v18v15-away-zero}
For every descendant parameter vector \(u\),
\[
\left|
 u^{\mathsf T}\nabla_\theta\partial_Q\ell_x
\right|^2
\leq
\frac{\chi_0^2}{Qr_*}\,
\mathbb E_{\eta,y}
\left[u^{\mathsf T}D_\theta(\eta,q)u\right],
\qquad Q>0,                                         \tag{25.18}
\]
where \(D_\theta=\operatorname{diag}(H_1,\ldots,H_r)\) is the descendant
block of (24.7).
\end{proposition}

\begin{proof}
Theorem \ref{thm:v18v15-mixed-response} and Cauchy--Schwarz give
\[
\begin{split}
 \left|u^{\mathsf T}\nabla_\theta\partial_Q\ell_x\right|^2
 &\leq\frac{\chi_0^2}{Q^2}
 \operatorname{Var}(y\cdot z)\,
 \operatorname{Var}(u^{\mathsf T}\nabla_\theta\mathcal S).
                                                               \tag{25.19}
\end{split}
\]
Brascamp--Lieb and (25.17) imply
\[
 \operatorname{Var}(y\cdot z)\leq\frac Q{r_*}.      \tag{25.20}
\]
The parameter-score covariance estimate in the proof of V14, together with
the pointwise Schur condition, gives
\[
 \operatorname{Var}(u^{\mathsf T}\nabla_\theta\mathcal S)
 \leq\mathbb E[u^{\mathsf T}D_\theta u].            \tag{25.21}
\]
Substitution into (25.19) proves (25.18).
\end{proof}

\begin{firewall}[The apparent singularity is not a zero-boundary estimate]
\label{fire:v18v15-singularity}
The right side of (25.18) behaves like \(Q^{-1}\).  It must not be used at
\(Q=0\).  The exact finite values there are (25.10)--(25.11).  A uniform
depth recursion therefore needs a small-\(Q\) analytic covariance estimate
joined to the away-from-zero response bound.
\end{firewall}

\section{V15 consequence and audit point}
\label{sec:v18v15-next}

The propagated mixed vector is now explicit in both boundary regimes.  The
remaining depth-uniform task is to bound the radial covariance in (25.10)
by the already propagated child Hessian margin, and to join that bound to
(25.18) at a positive radius without losing constants at each generation.

This is V15 of the eighteenth cycle.  The newest active Yang--Mills and
Navier--Stokes notes must be inspected before V16 is chosen.

\begin{assessment}[V15 acquisition and boundary]
\label{ass:v18v15}
V15 proves exact mixed-boundary response formulas for a recursive tree
message, including the finite isotropic zero-boundary value, and an
away-from-zero Brascamp--Lieb response bound.  It prevents a false
\(Q^{-1}\) singularity from entering the recursion.  A depth-uniform join,
cycles, cutoff removal, reconstruction, and the Einstein continuum remain
open.
\end{assessment}

\begin{record}[V15 append record]
\label{rec:v18v15}
V15 is appended to the validated V14 whose installed SHA--256 was
\texttt{A0AF5E53E0BA28FB135B9FFF892C81BF2108A30E6F422D33B82CFF4C25B90B5D}.
After installation only V15 is the active numeric SM note; the frozen
\texttt{V100\_f17} archive is unchanged.
\end{record}

\section{V16: companion audit and a two-region depth-transfer bound}
\label{sec:v18v16}

The post-V15 companion audit is recorded first.  The mathematical work then
joins the exact zero-boundary response of V15 to its nonzero-boundary
Brascamp--Lieb estimate.  The resulting theorem gives a finite bound for the
propagated quadratic transfer quotient whenever the V14 direct-curvature
ledger has uniform margins.

\begin{audit}[Post-V15 companion audit]
\label{aud:v18v16-companions}
The active companion files are unchanged from the V11 audit:
\[
\begin{array}{c|r|r|l}
\text{file}&\text{bytes}&\text{lines}&\text{SHA--256 prefix}\\
\hline
\texttt{Yang\_Mills\_Mass\_Gap\_Research\_Notes\_V2.tex}
&34462&910&\texttt{20687686090948DA}\\
\texttt{NS\_Stability\_Research\_Notes\_V26.tex}
&278283&5319&\texttt{82C887F8C2C69E4}.
\end{array}                                                   \tag{26.1}
\]
Their complete digests remain those in (21.2).  Yang--Mills V2 still leaves
its continuous alias tail open after a certified finite-grid sign, and
Navier--Stokes V26 still concerns rank-one Oseen constants.  Neither changes
the recursive Higgs message estimates, so no theorem or constant is
imported.
\end{audit}

\subsection{A zero-boundary covariance bound}

Use the zero-boundary law \(\nu_\eta\) from (25.8).  Let
\[
 D_\theta(q)=\operatorname{diag}(H_1,\ldots,H_r)     \tag{26.2}
\]
be the descendant direct-curvature block.  Assume the pointwise recursive
Schur condition (24.10).

\begin{proposition}[Finite anchor for the descendant mixed vector]
\label{prop:v18v16-zero-anchor}
At \(Q=0\),
\[
 c_\theta(0)c_\theta(0)^{\mathsf T}
 \preceq
 \frac{2\chi_0^4m}{\alpha^2a_0}\,
 \mathbb E_{\nu_\eta}D_\theta(q),                  \tag{26.3}
\]
where \(m=\mathbb E_{\nu_\eta}q\).  In addition,
\[
 m\leq\sqrt{\frac{\alpha}{2\lambda_x}},             \tag{26.4}
\]
and
\[
 |c_{d_0}(0)|
 \leq b_0+\frac{2\chi_0^2m}{\alpha}.                \tag{26.5}
\]
\end{proposition}

\begin{proof}
For a descendant vector \(u\), equation (25.10) and Cauchy--Schwarz give
\[
\begin{split}
 |u^{\mathsf T}c_\theta(0)|^2
 \leq\frac{\chi_0^4}{\alpha^2}
 \operatorname{Var}_{\nu_\eta}(q)\,
 \operatorname{Var}_{\nu_\eta}
 (u^{\mathsf T}\nabla_\theta\mathcal S).            \tag{26.6}
\end{split}
\]
The zero-boundary spatial potential has Hessian at least \(2a_0I\).
Poincare--Brascamp--Lieb applied to \(q=|z|^2\) gives
\[
 \operatorname{Var}_{\nu_\eta}(q)
 \leq\frac1{2a_0}\mathbb E|2z|^2
 =\frac{2m}{a_0}.                                   \tag{26.7}
\]
The proof of Proposition \ref{prop:v18v15-away-zero}, now at \(Q=0\), gives
\[
 \operatorname{Var}
 (u^{\mathsf T}\nabla_\theta\mathcal S)
 \leq
 \mathbb E[u^{\mathsf T}D_\theta u].                \tag{26.8}
\]
Equations (26.6)--(26.8) prove (26.3).

The radial potential in (25.8) is a quartic
\(\lambda_xq^2\), a quadratic \(a_0q\), and convex nondecreasing child
potentials.  Its recurrence (9.6), with \(s=\alpha\), implies
\(\alpha\geq2\lambda_xm^2\), proving (26.4).  Finally, (25.11) and
(26.7) give
\[
 |c_{d_0}(0)|
 \leq b_0+\frac{\chi_0^2a_0}{\alpha}\frac{2m}{a_0},
\]
which is (26.5).
\end{proof}

\subsection{Analytic control near zero}

Let \(c_x(Q)=\nabla_\eta\partial_Q\ell_x(\eta,Q)\) and
\[
 H_x(Q)=\nabla_\eta^2\ell_x(\eta,Q).                \tag{26.9}
\]

\begin{lemma}[Uniform analytic collar]
\label{lem:v18v16-collar}
Fix a compact positive set of local coefficients and descendant message
parameters for which the child score derivatives have a common polynomial
moment envelope under (25.8).  Then for every fixed \(Q_*>0\) there is a
finite computable constant \(L_c(Q_*)\) such that
\[
 \|c_x(Q)-c_x(0)\|
 \leq L_c(Q_*)Q,\qquad 0\leq Q\leq Q_*.             \tag{26.10}
\]
\end{lemma}

\begin{proof}
Angular integration of (24.5) gives the positive series
\[
 \sum_{k\geq0}
 \frac{(\chi_0^2Q)^k}{k!(\alpha)_k}
 J_{\alpha+k}^{\Psi_\theta}(a_0),                  \tag{26.11}
\]
where
\[
 \Psi_\theta(q)=\lambda_xq^2+
 \sum_j\ell_j(\theta_j,q).
\]
The quartic term and monotonicity of the child potentials give the moment
bound (22.4) with \(\lambda=\lambda_x\).  The assumed polynomial envelope
does the same after one or two descendant-score insertions.  Therefore
(26.11), its \(\eta\)-derivatives, and two \(Q\)-derivatives converge
uniformly on the declared compact set and on \(0\leq Q\leq Q_*\).
The derivative \(\partial_Qc_x\) consequently has a finite computable
supremum.  The mean-value theorem gives (26.10).
\end{proof}

\subsection{The joined transfer quotient}

Assume the V14 pointwise certificate has a uniform margin
\[
 H_x(Q)\succeq\mu_x I,\qquad Q\geq0,\qquad \mu_x>0. \tag{26.12}
\]
For \(Q>0\), define direct-score ceilings
\[
\begin{split}
 M_\theta(Q)&=
 \left\|\mathbb E_{\eta,y}D_\theta(q)\right\|_{\mathrm{op}},\\
 M_0(Q)&=\mathbb E_{\eta,y}T_0(q,Q).                \tag{26.13}
\end{split}
\]
The spatial Hessian in V14 obeys
\[
 r_*:=\inf_z\lambda_{\min}\nabla_z^2\mathcal S
 \geq2a_0.                                          \tag{26.14}
\]

\begin{theorem}[Two-region propagated transfer bound]
\label{thm:v18v16-two-region}
Let \(Q_*>0\), and set
\[
\begin{split}
 C_{\theta,0}^2&=
 \frac{2\chi_0^4}{\alpha^2a_0}
 \sqrt{\frac{\alpha}{2\lambda_x}}\,
 \left\|\mathbb E_{\nu_\eta}D_\theta\right\|_{\mathrm{op}},\\
 C_{0,0}&=
 b_0+\frac{2\chi_0^2}{\alpha}
 \sqrt{\frac{\alpha}{2\lambda_x}},\\
 C_{\mathrm{collar}}&=
 \sqrt{C_{\theta,0}^2+C_{0,0}^2}
 +L_c(Q_*)Q_*.                                      \tag{26.15}
\end{split}
\]
Then
\[
 c_x(Q)^{\mathsf T}H_x(Q)^{-1}c_x(Q)
 \leq\frac{C_{\mathrm{collar}}^2}{\mu_x},
 \qquad 0\leq Q\leq Q_*.                           \tag{26.16}
\]
For \(Q\geq Q_*\),
\[
\begin{split}
 c_x(Q)^{\mathsf T}H_x(Q)^{-1}c_x(Q)
 \leq\frac1{\mu_x}\Bigg[
 &\frac{\chi_0^2M_\theta(Q)}{Qr_*}\\
 &+\left(
 b_0+\chi_0\sqrt{\frac{M_0(Q)}{Qr_*}}
 \right)^2
 \Bigg].                                            \tag{26.17}
\end{split}
\]
In particular, if
\[
 \sup_{Q\geq Q_*}\frac{M_\theta(Q)}Q<\infty,
 \qquad
 \sup_{Q\geq Q_*}\frac{M_0(Q)}Q<\infty,             \tag{26.18}
\]
then the propagated transfer quotient has a finite global ceiling.
\end{theorem}

\begin{proof}
Proposition \ref{prop:v18v16-zero-anchor} bounds the descendant and outgoing
components of \(c_x(0)\) by the first two lines of (26.15).  Lemma
\ref{lem:v18v16-collar} therefore gives
\(|c_x(Q)|\leq C_{\mathrm{collar}}\) on the collar.  Equation (26.12) implies
\(H_x^{-1}\preceq\mu_x^{-1}I\), which proves (26.16).

For \(Q>0\), Proposition \ref{prop:v18v15-away-zero} gives the descendant
part
\[
 |c_\theta(Q)|^2\leq
 \frac{\chi_0^2M_\theta(Q)}{Qr_*}.                  \tag{26.19}
\]
For the outgoing component, use (25.5), Cauchy--Schwarz,
\(\operatorname{Var}(y\cdot z)\leq Q/r_*\), and
\(\operatorname{Var}(\partial_{d_0}\mathcal S)\leq M_0(Q)\):
\[
 |c_{d_0}(Q)|
 \leq b_0+\chi_0\sqrt{\frac{M_0(Q)}{Qr_*}}.         \tag{26.20}
\]
Add the squared component bounds and use
\(H_x^{-1}\preceq\mu_x^{-1}I\) to prove (26.17).  Conditions (26.18) make
its right side uniformly finite.
\end{proof}

\section{V16 consequence and next unpaid row}
\label{sec:v18v16-next}

The apparent zero-boundary singularity is now removed, and a global
propagated transfer ceiling follows from four named local ledgers:
\[
 \mu_x,\qquad L_c(Q_*),\qquad
 \sup_{Q\geq Q_*}\frac{M_\theta(Q)}Q,\qquad
 \sup_{Q\geq Q_*}\frac{M_0(Q)}Q.                   \tag{26.21}
\]
The quartic onsite term supplies the moment part of the first two-region
join.  What remains is to prove the two linear-growth ceilings in (26.18)
uniformly under tree recursion and to keep \(\mu_x\) from decaying with
depth.

\begin{assessment}[V16 acquisition and boundary]
\label{ass:v18v16}
V16 completes the mandatory companion audit and proves a finite
zero-boundary anchor, an analytic collar, and an away-from-zero transfer
bound for recursive messages.  Under the named uniform ledgers it gives a
finite global ceiling for \(c_x^{\mathsf T}H_x^{-1}c_x\).  Those ledgers are
not yet closed uniformly in depth; cycles, regulator removal,
reconstruction, and the coupled Einstein continuum remain open.
\end{assessment}

\begin{record}[V16 append record]
\label{rec:v18v16}
V16 is appended to the validated V15 whose installed SHA--256 was
\texttt{266FC8ACAE1563B786EEE59728D0834D56145E2733A1725A8FEBE7B5932DCEEE}.
The companion fingerprints are unchanged from (21.2).  After installation
only V16 is the active numeric SM note; the frozen \texttt{V100\_f17}
archive is unchanged.
\end{record}

\section{V17: quartic moment growth closes the two large-boundary ledgers}
\label{sec:v18v17}

V16 reduced the away-from-zero transfer ceiling to the two growth conditions
(26.18).  This section proves those conditions from quartic confinement,
provided the recursive Schur certificates keep the conformal Hessians
positive.  The key estimate is a \(Q^{1/3}\) bound for the mean squared
internal field, uniform under arbitrary convex nondecreasing child
potentials.

\begin{lemma}[Quartic response moment]
\label{lem:v18v17-moment}
Let \(z\in\mathbb R^N\) have density proportional to
\[
 \exp\left[
 -\lambda|z|^4-a|z|^2-\Phi(|z|^2)+2\chi z\cdot y
 \right],                                           \tag{27.1}
\]
where \(\lambda,a,\chi>0\) and \(\Phi\) is differentiable, convex, and
nondecreasing.  Put \(Q=|y|^2\) and \(m=\mathbb E|z|^2\).  Then
\[
 m\leq C_0+C_1Q^{1/3},                              \tag{27.2}
\]
where one admissible choice is
\[
\begin{split}
 C_0&=\sqrt{\frac{N}{3\lambda}},\\
 C_1&=
 \left[
 \frac{2^{-4/3}\,2^{2/3}\chi^{4/3}}
 {\lambda^{4/3}}
 \right]^{1/2}
 =
 \frac{\chi^{2/3}}{2^{1/3}\lambda^{2/3}}.           \tag{27.3}
\end{split}
\]
Any larger explicit constants are also valid.
\end{lemma}

\begin{proof}
Integration by parts with the dilation vector field \(z\) gives
\[
\begin{split}
 N={}&4\lambda\mathbb E|z|^4+2a\mathbb E|z|^2
 +2\mathbb E\left[|z|^2\Phi'(|z|^2)\right]
 -2\chi\mathbb E(z\cdot y).                         \tag{27.4}
\end{split}
\]
The middle two terms on the first line are nonnegative.  Jensen and
Cauchy--Schwarz therefore imply
\[
 4\lambda m^2
 \leq N+2\chi\sqrt Q\,\sqrt m.                      \tag{27.5}
\]
Write \(x=\sqrt m\).  Young's inequality in the form
\[
 uv\leq\frac{\varepsilon}{4}u^4
 +\frac34\varepsilon^{-1/3}v^{4/3}                 \tag{27.6}
\]
with \(u=x\), \(v=2\chi\sqrt Q\), and
\(\varepsilon=4\lambda\), gives
\[
 2\chi\sqrt Q\,x
 \leq\lambda x^4+
 \frac34(4\lambda)^{-1/3}(2\chi)^{4/3}Q^{2/3}.
                                                               \tag{27.7}
\]
Insert this in (27.5), divide by \(3\lambda\), and take a square root.
Using \(\sqrt{r+s}\leq\sqrt r+\sqrt s\) proves (27.2).  Simplifying the
second coefficient gives the displayed \(C_1\).
\end{proof}

\begin{remark}[Coefficient check]
\label{rem:v18v17-coefficient}
The particular \(C_1\) in (27.3) is a convenient enlargement of the
coefficient obtained directly from (27.7).  The exponent \(Q^{1/3}\), rather
than this numerical prefactor, is the part used below.
\end{remark}

\begin{definition}[Affine radial Hessian-growth class]
\label{def:v18v17-affine-class}
A radial message \(\ell(\theta,q)\) belongs to
\(\mathfrak A(A,B)\) if
\[
 0\preceq\nabla_\theta^2\ell(\theta,q)
 \preceq(A+Bq)I
 \qquad(q\geq0).                                    \tag{27.8}
\]
\end{definition}

\begin{proposition}[Pure quartic base case]
\label{prop:v18v17-base}
On every compact positive \((\lambda,\chi,d)\)-set, the pure quartic leaf
message belongs to \(\mathfrak A(A_{\rm leaf},B_{\rm leaf})\) for finite
uniform constants.
\end{proposition}

\begin{proof}
The parameter space of one leaf is one-dimensional, and its curvature is
\[
 D(q)=bq+a\mathbb E|u|^2-a^2\operatorname{Var}(|u|^2)
 \leq bq+a\mathbb E|u|^2.                           \tag{27.9}
\]
Apply Lemma \ref{lem:v18v17-moment} with \(\Phi=0\) and parent radius
\(q\).  On a compact positive coefficient set,
\[
 D(q)\leq C(1+q^{1/3})+bq\leq A_{\rm leaf}+B_{\rm leaf}q.
                                                               \tag{27.10}
\]
Positivity is Theorem \ref{thm:v18v2-global-leaf}.
\end{proof}

\begin{theorem}[Affine Hessian growth propagates through one vertex]
\label{thm:v18v17-propagation}
Use the recursive message (24.5), and suppose:
\begin{enumerate}[leftmargin=2.2em]
\item every child message belongs to
      \(\mathfrak A(A_j,B_j)\);
\item the V14 pointwise Schur condition holds, so the output conformal
      Hessian is positive semidefinite.
\end{enumerate}
Then the output message belongs to
\(\mathfrak A(A_x,B_x)\) for finite constants depending only on the local
coefficients, degree, and the child constants.  These constants are uniform
on compact positive local-parameter sets of bounded degree.
\end{theorem}

\begin{proof}
The exact marginal Hessian identity (24.13) gives
\[
 0\preceq H_x(Q)
 =\mathbb E D-\operatorname{Cov}(\nabla_\eta\mathcal S)
 \preceq\mathbb E D.                                \tag{27.11}
\]
The outgoing diagonal block of \(D\) is
\[
 T_0=a_0q+b_0Q.                                     \tag{27.12}
\]
The child blocks obey
\[
 H_j(q)\preceq(A_j+B_jq)I.                          \tag{27.13}
\]
Therefore
\[
 \|\,\mathbb E D\,\|_{\mathrm{op}}
 \leq C_A+C_B\,\mathbb E q+b_0Q                    \tag{27.14}
\]
for finite sums or maxima \(C_A,C_B\) determined by the child ledgers and
\(a_0\).

In the conditional \(z\)-law, the child potential
\(\Phi(q)=\sum_j\ell_j(\theta_j,q)\) is convex and nondecreasing by V3.
Lemma \ref{lem:v18v17-moment} thus gives
\[
 \mathbb E q\leq C_0+C_1Q^{1/3}.                   \tag{27.15}
\]
Insert (27.15) into (27.14) and use
\(Q^{1/3}\leq1+Q\).  This yields
\[
 H_x(Q)\preceq(A_x+B_xQ)I.                          \tag{27.16}
\]
All constants use only finite sums, compact endpoint bounds, and the bounded
degree, proving uniformity.
\end{proof}

\begin{corollary}[Closure of the V16 large-boundary conditions]
\label{cor:v18v17-growth-ledgers}
Under the hypotheses of Theorem \ref{thm:v18v17-propagation}, for every
\(Q_*>0\),
\[
 \sup_{Q\geq Q_*}\frac{M_\theta(Q)}Q<\infty,
 \qquad
 \sup_{Q\geq Q_*}\frac{M_0(Q)}Q<\infty.             \tag{27.17}
\]
Hence the away-from-zero right side of (26.17) has a finite global ceiling.
\end{corollary}

\begin{proof}
The child direct block satisfies
\[
 D_\theta(q)\preceq
 \left(C_A+C_Bq\right)I.
\]
Taking the conditional expectation and using (27.15) gives
\[
 M_\theta(Q)\leq C(1+Q^{1/3}).                      \tag{27.18}
\]
Also
\[
 M_0(Q)=a_0\mathbb E q+b_0Q
 \leq C(1+Q^{1/3})+b_0Q.                            \tag{27.19}
\]
Divide by \(Q\geq Q_*\) to prove (27.17), then invoke Theorem
\ref{thm:v18v16-two-region}.
\end{proof}

\begin{corollary}[Finite-depth closure conditional on local margins]
\label{cor:v18v17-finite-depth}
On any finite tree, suppose the V14 pointwise Schur condition holds with a
positive output margin at every internal vertex.  Then every recursive
message has:
\begin{enumerate}[leftmargin=2.2em]
\item a positive conformal Hessian;
\item an affine radial Hessian-growth bound;
\item a finite global transfer quotient
      \(c_x^{\mathsf T}H_x^{-1}c_x\).
\end{enumerate}
\end{corollary}

\begin{proof}
The pure quartic leaves satisfy the base case by Proposition
\ref{prop:v18v17-base}.  Apply Theorems
\ref{thm:v18v17-propagation} and
\ref{thm:v18v16-two-region} inductively toward the root.
\end{proof}

\section{V17 consequence and remaining depth issue}
\label{sec:v18v17-next}

The two growth ledgers in (26.18) are no longer open.  Conditional on local
positive Schur margins, the radial shape, conformal Hessian upper growth, and
mixed-transfer ceiling all propagate through every finite depth.

The remaining cofinal tree issue is the lower margin \(\mu_x\).  The present
induction allows it to shrink with depth.  A volume-uniform massless
incidence theorem requires either a contraction of the transfer budget or a
fixed-point cone whose lower margin is preserved at every branching step.
That is the next upstream target.

\begin{assessment}[V17 acquisition and boundary]
\label{ass:v18v17}
V17 proves the quartic \(Q^{1/3}\) response-moment bound, establishes the
pure-leaf affine Hessian-growth base case, and propagates that class through
recursive tree messages.  It closes both large-boundary growth assumptions
from V16 and yields finite-depth closure conditional on local positive
margins.  A depth-uniform lower margin, cycles, cofinal cutoff removal,
reconstruction, and the Einstein continuum remain open.
\end{assessment}

\begin{record}[V17 append record]
\label{rec:v18v17}
V17 is appended to the validated V16 whose installed SHA--256 was
\texttt{0754F6ADFCA258575C23B977A846A0F14D32902E263290DD8AAFA7883538E8D8}.
After installation only V17 is the active numeric SM note; the frozen
\texttt{V100\_f17} archive is unchanged.
\end{record}

\section{V18: the depth-margin product and its summability gate}
\label{sec:v18v18}

V17 closes the upper-growth ledgers, but a cofinal tree theorem also needs a
lower Hessian margin that does not vanish with depth.  This section proves a
uniform lower moment at every vertex and then tracks the precise
multiplicative loss created by the local rank-one Schur update.  The result
isolates a new necessary feature of this proof route: local relative losses
must have a nonzero infinite product.

\subsection{A depth-independent radial moment floor}

At an internal vertex, let
\[
 B_x=\sum_{j=1}^r b_j                              \tag{28.1}
\]
be the sum of the child boundary-slope ceilings.

\begin{lemma}[Uniform child-slope ceiling]
\label{lem:v18v18-slope}
Every child message generated by the recursive transform satisfies
\[
 0<\partial_q\ell_j(\theta_j,q)<b_j
 \qquad(q\geq0).                                    \tag{28.2}
\]
Consequently, after subtracting an irrelevant constant,
\[
 0\leq\sum_j\ell_j(\theta_j,q)
       -\sum_j\ell_j(\theta_j,0)
 \leq B_xq.                                         \tag{28.3}
\]
\end{lemma}

\begin{proof}
Every internal effective potential contains the nonconstant onsite quartic.
Proposition \ref{prop:v18v4-stiffness}, applied recursively, gives (28.2).
Integrate the derivative bounds from \(0\) to \(q\) and sum to obtain
(28.3).
\end{proof}

For \(c,\lambda>0\), define
\[
 m_\alpha(c,\lambda)=
 \frac{\displaystyle\int_0^\infty
 t^\alpha e^{-\lambda t^2-ct}\,dt}
 {\displaystyle\int_0^\infty
 t^{\alpha-1}e^{-\lambda t^2-ct}\,dt}.              \tag{28.4}
\]

\begin{proposition}[Internal second-moment floor]
\label{prop:v18v18-moment-floor}
For the conditional law in (24.5),
\[
 \mathbb E_{\eta,y}|z|^2
 \geq m_\alpha(a_0+B_x,\lambda_x)>0.                \tag{28.5}
\]
On compact positive coefficient sets of bounded degree, the right side has
a common positive lower bound.
\end{proposition}

\begin{proof}
First take \(y=0\).  In the radial variable \(t=|z|^2\), compare the actual
law with the reference law proportional to
\[
 t^{\alpha-1}e^{-\lambda_xt^2-(a_0+B_x)t}\,dt.      \tag{28.6}
\]
By (28.3), the Radon--Nikodym ratio of the actual law to (28.6), up to
normalization, is
\[
 \exp\left[
 B_xt-\sum_j(\ell_j(\theta_j,t)-\ell_j(\theta_j,0))
 \right],                                           \tag{28.7}
\]
which is nondecreasing in \(t\).  Chebyshev's covariance inequality for two
nondecreasing functions shows that the actual law stochastically increases
the expectation of \(t\).  This proves (28.5) at \(y=0\).

For \(y\ne0\), angular integration multiplies the zero-boundary radial
density by
\[
 {}_0F_1(; \alpha;\chi_0^2|y|^2t),                 \tag{28.8}
\]
which is positive and nondecreasing in \(t\).  The same covariance argument
shows that this tilt can only increase \(\mathbb Et\).  This proves (28.5)
for all \(y\).

The function in (28.4) is continuous and positive.  It decreases when
either \(c\) or \(\lambda\) increases, by the covariance derivative
identities.  Compact upper bounds on \(a_0+B_x\) and \(\lambda_x\) therefore
give a common positive lower bound.
\end{proof}

\subsection{Relative Schur slack}

At a vertex \(x\), the direct block \(D_x\) in (24.7) is positive definite
whenever \(q+Q>0\).  At the single point \((q,Q)=(0,0)\), its edge entry
\(T_0=a_0q+b_0Q\) vanishes, while the rank-one coefficient \(4q/R_x\)
also vanishes.  For \(q+Q>0\), define the pointwise relative load
\[
 \kappa_x(q,Q)=
 \frac{4q}{R_x(q)}
 c_x(q)^{\mathsf T}D_x(q,Q)^{-1}c_x(q).             \tag{28.9}
\]

\begin{lemma}[Relative rank-one margin]
\label{lem:v18v18-relative}
If
\[
 \kappa_x(q,Q)\leq1-\varepsilon_x
 \qquad\text{for all }q,Q\geq0\text{ with }q+Q>0  \tag{28.10}
\]
with \(0<\varepsilon_x\leq1\), then, for all \(q,Q\geq0\),
\[
 D_x-\frac{4q}{R_x}c_xc_x^{\mathsf T}
 \succeq\varepsilon_xD_x.                           \tag{28.11}
\]
\end{lemma}

\begin{proof}
For \(q+Q>0\), let
\[
 v=\left(\frac{4q}{R_x}\right)^{1/2}D_x^{-1/2}c_x.
\]
Then \(|v|^2=\kappa_x\leq1-\varepsilon_x\), so
\[
 I-vv^{\mathsf T}\succeq\varepsilon_xI.
\]
Conjugate this inequality by \(D_x^{1/2}\) to obtain (28.11).  At
\((q,Q)=(0,0)\), the rank-one term is zero, so (28.11) reduces to
\(D_x\succeq\varepsilon_xD_x\), which holds because \(D_x\) is positive
semidefinite.
\end{proof}

\begin{theorem}[One-step lower-margin recursion]
\label{thm:v18v18-one-step}
Assume (28.10), and suppose every child Hessian has a uniform lower bound
\[
 H_j(\theta_j,q)\succeq\mu_jI.                      \tag{28.12}
\]
Then the outgoing message Hessian satisfies
\[
 H_x(\eta,Q)\succeq
 \varepsilon_x
 \min\left\{
 a_0m_\alpha(a_0+B_x,\lambda_x),\
 \mu_1,\ldots,\mu_r
 \right\}I.                                         \tag{28.13}
\]
\end{theorem}

\begin{proof}
The recursive transport inequality (24.9) and Lemma
\ref{lem:v18v18-relative} give
\[
 H_x\succeq\varepsilon_x\mathbb E D_x.              \tag{28.14}
\]
The outgoing direct entry is
\[
 \mathbb ET_0=a_0\mathbb E q+b_0Q
 \geq a_0m_\alpha(a_0+B_x,\lambda_x)                \tag{28.15}
\]
by Proposition \ref{prop:v18v18-moment-floor}.  Every child block is at
least \(\mu_jI\).  Since \(D_x\) is block diagonal, (28.13) follows.
\end{proof}

\subsection{The cofinal product}

Consider a descendant edge whose conformal direction passes through
successive ancestors \(x_1,\ldots,x_h\).  Let \(\mu_{\rm birth}>0\) be its
margin when first created, including the local moment floor.

\begin{theorem}[Depth-product lower bound]
\label{thm:v18v18-product}
Under the relative slack conditions (28.10) along the path,
\[
 \mu_h\geq
 \mu_{\rm birth}\prod_{\ell=1}^h\varepsilon_{x_\ell}.
                                                               \tag{28.16}
\]
For an infinite cofinal depth sequence, this method gives a positive
depth-uniform inherited margin exactly when
\[
 \inf_h\prod_{\ell=1}^h\varepsilon_{x_\ell}>0,
                                                               \tag{28.17}
\]
equivalently,
\[
 \sum_{\ell=1}^\infty-\log\varepsilon_{x_\ell}<\infty.
                                                               \tag{28.18}
\]
In particular, if \(\varepsilon_{x_\ell}\leq\bar\varepsilon<1\) at every
level, the certified inherited margin tends to zero exponentially.
\end{theorem}

\begin{proof}
Apply Theorem \ref{thm:v18v18-one-step} at each ancestor.  In the block
corresponding to the chosen descendant direction, the previous margin is
multiplied by at least \(\varepsilon_{x_\ell}\).  Iteration gives (28.16).
The equivalence of (28.17) and (28.18) is the elementary infinite-product
criterion for numbers in \((0,1]\).  The final statement follows from
\(\prod_{\ell=1}^h\varepsilon_{x_\ell}\leq\bar\varepsilon^h\).
\end{proof}

\begin{corollary}[Summable relative loss is sufficient]
\label{cor:v18v18-summable}
If
\[
 \varepsilon_{x_\ell}=1-\delta_\ell,\qquad
 0\leq\delta_\ell\leq\delta_*<1,\qquad
 \sum_{\ell=1}^\infty\delta_\ell<\infty,            \tag{28.19}
\]
then the product in (28.17) is positive.
\end{corollary}

\begin{proof}
On \([0,\delta_*]\),
\[
 -\log(1-\delta)\leq\frac{\delta}{1-\delta_*}.
\]
Thus (28.19) implies (28.18).
\end{proof}

\section{V18 consequence and upstream decision}
\label{sec:v18v18-next}

The internal radial moment cannot collapse on a bounded-degree compact
coefficient sector; Proposition \ref{prop:v18v18-moment-floor} supplies a
fresh local floor at every vertex.  The inherited directions are different:
the present Schur induction multiplies their margins by the relative slacks.
A fixed local contraction away from one is therefore insufficient for a
cofinal tree.

The next upstream question is whether the canonical refinement makes the
relative losses \(\delta_\ell=1-\varepsilon_{x_\ell}\) summable, or whether
the exact discrete angular covariance of V9 must replace Brascamp--Lieb at
every level to avoid multiplicative loss.  Since the canonical Higgs hopping
and quartic coefficients are order one after rescaling, summability cannot be
inferred from powers of the lattice spacing without a new cancellation.

\begin{assessment}[V18 acquisition and boundary]
\label{ass:v18v18}
V18 proves a depth-independent internal second-moment floor, a relative
rank-one margin theorem, and the exact depth-product criterion for inherited
conformal directions.  It shows that fixed per-level Schur loss makes this
certificate degenerate at cofinal depth.  Summable loss or an exact
covariance cancellation remains unproved; cycles, cutoff removal,
reconstruction, and the Einstein continuum remain open.
\end{assessment}

\begin{record}[V18 append record]
\label{rec:v18v18}
V18 is appended to the validated V17 whose installed SHA--256 was
\texttt{E3E056601D3B7F932058E178DF3B8C5CD63D4CFA586C3356FD528E3ECD94D608}.
After installation only V18 is the active numeric SM note; the frozen
\texttt{V100\_f17} archive is unchanged.
\end{record}

% =====================================================================
% V19 APPENDIX: ADDITIVE TREE LEDGER AND GLOBAL BESSEL GATE
% =====================================================================

\section{V19: an additive tree ledger replaces the depth product}
\label{sec:v18v19}

V18 located the failure of the pointwise recursive certificate: applying a
relative Schur slack separately at every ancestor spends part of every
descendant margin repeatedly.  That bookkeeping is sufficient but too
expensive.  This note goes upstream and expands all recursive Schur losses
before comparing them with the positive terms.  The result is a single
finite-tree operator inequality.  Its continuum gate is a uniform Bessel
bound for a family of transported response vectors; no product of local
slacks occurs.

The result concerns the conformal edge parameters already used in V14--V18.
It does not add a new physical assumption.  It reorganizes the same
Brascamp--Lieb lower bound and therefore gives a sharper target for the
remaining proof.

\subsection{Conditional tree kernels}

Let \(\mathcal T_x\) be a finite rooted descendant tree with root \(x\).
Write \(e_v\) for the edge from a vertex \(v\) to its parent and \(d_v\)
for its conformal parameter.  The parameter space is
\[
 E_x=\mathbb R^{E(\mathcal T_x)}
     =\mathbb Re_x\oplus\bigoplus_{j=1}^{r_x}E_{x_j},             \tag{29.1}
\]
where \(x_1,\ldots,x_{r_x}\) are the children of \(x\).  Put
\[
 q_v=|z_v|^2,
 \qquad Q_v=|z_{p(v)}|^2,
 \qquad T_v=a_vq_v+b_vQ_v.                                     \tag{29.2}
\]
For the root, \(Q_x\) is the prescribed exterior boundary magnitude.

The normalized integral defining the message at \(v\) gives a conditional
probability kernel
\[
 K_v(dz_v\mid Q_v).                                             \tag{29.3}
\]
Composing these kernels from the root to the leaves gives a probability
measure \(\mathbb P_{\mathcal T_x,Q_x}\) on all internal fields.  This is
only the finite-tree Gibbs disintegration: conditional on \(z_v\), the
child subtrees are independent and their boundary magnitude is \(q_v\).
Consequently, repeated conditional expectations in the message recursion
are expectations under \(\mathbb P_{\mathcal T_x,Q_x}\) by the tower
property.

At a vertex \(v\), let \(c_v\in E_v\) be the mixed-response vector of
V14,
\[
 c_v=(-a_v,c_{v,1},\ldots,c_{v,r_v}),
 \qquad c_{v,j}=\partial_{\theta_jq_v}\ell_{v,j},                \tag{29.4}
\]
and extend it by zero from \(E_v\) to \(E_x\); denote the extension by
\(\widetilde c_v\).  Let
\[
 w_v=\frac{4q_v}{R_v(q_v)},                                     \tag{29.5}
\]
where \(R_v\) is the positive radial field Hessian in (24.6).  Thus the
local Schur loss is \(w_vc_vc_v^{\mathsf T}\).

\subsection{The unrolled lower bound}

Define two deterministic matrices on \(E_x\):
\[
 A_{mathcal T_x}(Q_x)
   =\operatorname{diag}_{e_v\in E(\mathcal T_x)}
      \mathbb E_{\mathcal T_x,Q_x}T_v,                           \tag{29.6}
\]
and
\[
 B_{mathcal T_x}(Q_x)
   =\mathbb E_{\mathcal T_x,Q_x}
      \sum_{v\in V(\mathcal T_x)}
          w_v\widetilde c_v\widetilde c_v^{\mathsf T}.          \tag{29.7}
\]
Both are positive semidefinite.  The first is the fresh edge ledger and the
second is the accumulated radial Schur ledger.

\begin{theorem}[Additive tree-Schur ledger]
\label{thm:v18v19-ledger}
Let \(H_{\mathcal T_x}(Q_x)\) be the Hessian of the outgoing message with
respect to all conformal parameters \((d_v)_{e_v\in E(\mathcal T_x)}\).
Under the hypotheses of the recursive transport theorem (24.9),
\[
 H_{\mathcal T_x}(Q_x)
 \succeq A_{\mathcal T_x}(Q_x)-B_{\mathcal T_x}(Q_x).            \tag{29.8}
\]
Equivalently, for every \(u\in E_x\),
\[
 u^{\mathsf T}H_{\mathcal T_x}(Q_x)u
 \geq
 \mathbb E_{\mathcal T_x,Q_x}
 \left[
   \sum_v T_vu_v^2
   -\sum_v w_v(\widetilde c_v^{\mathsf T}u)^2
 \right].                                                       \tag{29.9}
\]
\end{theorem}

\begin{proof}
We argue by induction on the height.  For a leaf, (24.9) is precisely
\[
 H_x(Q_x)\succeq
 \mathbb E_{x,Q_x}
 \left[T_x-w_xc_xc_x^{\mathsf T}\right],
\]
which is (29.8).

Suppose the assertion holds for every child tree.  The block form of
(24.9) gives
\[
 H_{\mathcal T_x}(Q_x)\succeq
 \mathbb E_{x,Q_x}
 \left[
  \operatorname{diag}
  \bigl(T_x,H_{\mathcal T_{x_1}}(q_x),\ldots,
            H_{\mathcal T_{x_{r_x}}}(q_x)\bigr)
  -w_xc_xc_x^{\mathsf T}
 \right].                                                       \tag{29.10}
\]
Insert the inductive lower bounds in the child blocks.  The positive terms
are \(T_x\) and all descendant \(T_v\); the negative terms are the root
rank-one loss and all descendant losses.  Conditional independence and the
tower property convert the iterated expectations to
\(\mathbb E_{\mathcal T_x,Q_x}\).  This gives (29.8).  Taking a quadratic
form gives (29.9).
\end{proof}

Theorem \ref{thm:v18v19-ledger} does not apply a scalar loss to an entire
child block.  Each radial elimination subtracts only its actual rank-one
form.  Sparse support and weak overlap among different response vectors can
therefore be used at the final comparison.

\subsection{Fresh positivity and the global Bessel gate}

For a class of trees, define
\[
 p_*=inf_v
 a_v m_{\alpha_v}(a_v+B_v,\lambda_v),
 \qquad B_v=\sum_{j=1}^{r_v}b_{v,j}.                             \tag{29.11}
\]
On a bounded-degree compact positive coefficient sector,
Proposition \ref{prop:v18v18-moment-floor} gives \(p_*>0\).

\begin{lemma}[Uniform fresh-ledger floor]
\label{lem:v18v19-fresh}
For every finite tree and every exterior boundary value,
\[
 A_{\mathcal T_x}(Q_x)\succeq p_*I.                             \tag{29.12}
\]
\end{lemma}

\begin{proof}
Condition on the parent field of \(v\).  Equations (28.5) and (29.2) give
\[
 \mathbb E[T_v\mid Q_v]
 =a_v\mathbb E[q_v\mid Q_v]+b_vQ_v
 \geq a_vm_{\alpha_v}(a_v+B_v,\lambda_v)
 \geq p_* .                                                     \tag{29.13}
\]
Another expectation preserves the inequality.  Since (29.6) is diagonal,
(29.12) follows.
\end{proof}

\begin{theorem}[Global Bessel criterion]
\label{thm:v18v19-bessel}
Suppose there is an \(\varepsilon>0\), independent of the finite tree and
its exterior boundary, such that
\[
 B_{\mathcal T_x}(Q_x)
 \preceq(1-\varepsilon)A_{\mathcal T_x}(Q_x).                   \tag{29.14}
\]
Then
\[
 H_{\mathcal T_x}(Q_x)
 \succeq\varepsilon A_{\mathcal T_x}(Q_x)
 \succeq\varepsilon p_*I                                      \tag{29.15}
\]
uniformly in depth.  In quadratic-form language, (29.14) is the Bessel
inequality
\[
 \mathbb E\sum_v w_v(\widetilde c_v^{\mathsf T}u)^2
 \leq(1-\varepsilon)
       \mathbb E\sum_vT_vu_v^2
 \qquad (u\in E_x).                                             \tag{29.16}
\]
\end{theorem}

\begin{proof}
Subtract (29.14) from \(A_{\mathcal T_x}\) in (29.8), then use
Lemma \ref{lem:v18v19-fresh}.
\end{proof}

Thus the cofinal problem has been reduced from a product of scalar margins
to one operator norm.  Define the normalized loss matrix
\[
 C_{\mathcal T_x}(Q_x)
 =A_{\mathcal T_x}(Q_x)^{-1/2}
  B_{\mathcal T_x}(Q_x)
  A_{\mathcal T_x}(Q_x)^{-1/2}.                                 \tag{29.17}
\]
The gate (29.14) is exactly
\[
 \lambda_{\max}(C_{\mathcal T_x}(Q_x))\leq1-\varepsilon.       \tag{29.18}
\]
Every entry in (29.17) is a finite-dimensional expectation and is
amenable to the radial tail, angular truncation, and finite-cover
certificates of V11--V13.

\begin{proposition}[Certified row-sum gate]
\label{prop:v18v19-row}
If, uniformly in the finite tree and the exterior boundary,
\[
 \max_e\sum_f
 \frac{|(B_{\mathcal T_x})_{ef}|}
      {\sqrt{(A_{\mathcal T_x})_{ee}(A_{\mathcal T_x})_{ff}}}
 \leq1-\varepsilon,                                             \tag{29.19}
\]
then (29.14) holds.
\end{proposition}

\begin{proof}
The left side of (29.19) is the maximum absolute row sum of the symmetric
matrix \(C_{\mathcal T_x}\).  Hence
\(\|C_{\mathcal T_x}\|_{2\to2}\leq
  \|C_{\mathcal T_x}\|_{\infty\to\infty}\leq1-\varepsilon\).
Since \(C_{\mathcal T_x}\) is positive semidefinite, this is (29.18).
\end{proof}

There is also a transparent decay version.  Let \(d_L(e,f)\) be distance
in the line graph of the tree, and suppose its maximum degree is \(D_L\).

\begin{corollary}[Ancestry-decay certificate]
\label{cor:v18v19-decay}
Assume \(D_L\geq2\) and, for constants \(\beta\geq0\) and
\(0\leq\rho<(D_L-1)^{-1}\),
\[
 |(C_{\mathcal T_x})_{ef}|
 \leq\beta\rho^{d_L(e,f)}.                                     \tag{29.20}
\]
If
\[
 \beta\left(1+\frac{D_L\rho}{1-(D_L-1)\rho}\right)
 \leq1-\varepsilon,                                             \tag{29.21}
\]
then (29.14) holds uniformly in the size and depth of the tree.
\end{corollary}

\begin{proof}
From a fixed edge there is one edge at distance zero and at most
\(D_L(D_L-1)^{k-1}\) edges at distance \(k\geq1\).  Sum (29.20) over the
distance shells.  The resulting geometric series is the left side of
(29.21), so Proposition \ref{prop:v18v19-row} applies.
\end{proof}

\section{V19 consequence and next target}
\label{sec:v18v19-next}

The depth product of V18 is a property of the local relative comparison,
not an unavoidable property of the tree Hessian.  The additive expansion
(29.8) retains every fresh edge contribution and every actual rank-one loss
until the final operator comparison.  A depth-uniform Ward-convexity margin
now follows from the single global frame bound (29.16).

The next analytic target is correspondingly precise: prove decay of the
transported mixed responses entering \(\widetilde c_v\), strong enough for
(29.20), or verify the row sums (29.19) with the rigorous finite-star
machinery.  The quartic response estimates of V15--V17 control the size at
large boundary; they must now be sharpened to control genealogical overlap.

\begin{assessment}[V19 acquisition and boundary]
\label{ass:v18v19}
V19 proves the additive finite-tree Schur ledger, a uniform positive fresh
ledger, the global Bessel criterion, and two verifiable spectral sufficient
conditions.  This removes the artificial multiplication of local slack
from the proof architecture.  The uniform Bessel bound itself remains to be
proved.  Cycles, cutoff removal, reconstruction, and the coupled Einstein
continuum remain open.
\end{assessment}

\begin{record}[V19 append record]
\label{rec:v18v19}
V19 is appended to the validated V18 whose installed SHA--256 was
\texttt{0316FEBAF9258CC84D5FA05DB96D0A50E90DC6EB8C267644F45374B78AD4609A}.
After installation only V19 is the active numeric SM note; the frozen
\texttt{V100\_f17} archive is unchanged.
\end{record}

% =====================================================================
% V20 APPENDIX: RADIAL MARKOV CONTRACTION AND ANCESTRY DECAY
% =====================================================================

\section{V20: an exact radial contraction coefficient}
\label{sec:v18v20}

V19 reduced the uniform tree Hessian to a Bessel bound for the transported
mixed-response vectors.  This note derives the transport mechanism itself.
After angular integration, every radial elimination is a one-dimensional
Markov operator.  Its norm on the Lipschitz seminorm has an exact cumulative
density formula.  If that norm is smaller than the inverse branching rate,
mixed responses decay down ancestry and the V19 Bessel gate follows with an
explicit constant.

\subsection{The radial Markov operator}

Fix a vertex type \(v\) and all parameters other than the exterior squared
magnitude \(Q\).  With \(t=|z_v|^2\), angular integration gives the radial
density
\[
 p_{v,Q}(t)=\frac{1}{Z_v(Q)}t^{\alpha_v-1}
 \exp\left\{-\lambda_vt^2-a_vt-
                 \sum_{j=1}^{r_v}\ell_{v,j}(\theta_j,t)\right\}
 {}_0F_1(;\alpha_v;\chi_v^2Qt),
 \qquad t\geq0.                                                  \tag{30.1}
\]
The corresponding Markov operator is
\[
 (P_vF)(Q)=\int_0^\infty F(t)p_{v,Q}(t)\,dt.                    \tag{30.2}
\]
Quartic confinement and the derivative bounds of V12 justify
differentiation under the integral for the functions used below.

Define the cumulative derivative
\[
 G_{v,Q}(s)=\int_0^s\partial_Qp_{v,Q}(t)\,dt                    \tag{30.3}
\]
and the radial response coefficient
\[
 \rho_v=\sup_{Q\geq0}\int_0^\infty|G_{v,Q}(s)|\,ds.            \tag{30.4}
\]
At \(Q=0\), \(\partial_Q\) means the right derivative.  The
hypergeometric series in (30.1) and the quartic tail make (30.3)--(30.4)
well defined whenever the displayed supremum is finite.

\begin{proposition}[Exact Lipschitz response norm]
\label{prop:v18v20-lip}
For every locally absolutely continuous \(F\) with
\(\|F'\|_\infty<\infty\),
\[
 \left|(P_vF)'(Q)\right|
 \leq \left(\int_0^\infty|G_{v,Q}(s)|\,ds\right)
       \|F'\|_\infty.                                           \tag{30.5}
\]
Moreover,
\[
 \sup_{\|F'\|_\infty\leq1}|(P_vF)'(Q)|
 =\int_0^\infty|G_{v,Q}(s)|\,ds.                               \tag{30.6}
\]
Thus \(\rho_v\) is exactly the operator norm of \(P_v\) from the
Lipschitz seminorm in \(t\) to the Lipschitz seminorm in \(Q\).
\end{proposition}

\begin{proof}
Normalization gives
\[
 \int_0^\infty\partial_Qp_{v,Q}(t)\,dt=0,
 \qquad G_{v,Q}(0)=G_{v,Q}(\infty)=0.                            \tag{30.7}
\]
Integration by parts therefore yields
\[
 (P_vF)'(Q)
 =\int_0^\infty F(t)\partial_Qp_{v,Q}(t)\,dt
 =-\int_0^\infty F'(s)G_{v,Q}(s)\,ds.                           \tag{30.8}
\]
The boundary term vanishes because \(F\) grows at most linearly and the
quartic density has faster decay.  Equation (30.5) follows.  For the reverse
inequality in (30.6), take an absolutely continuous function with
\[
 F'(s)=-\operatorname{sgn}G_{v,Q}(s),\qquad F(0)=0.              \tag{30.9}
\]
It is one-Lipschitz and (30.8) equals the right side of (30.6).
\end{proof}

Formula (30.4) is a one-dimensional quantity.  On a compact coefficient
box and a compact \(Q\)-interval it can be enclosed by the certified radial
quadrature and finite-cover method of V12.  The origin is controlled by the
analytic hypergeometric expansion, and the V17 quartic response estimates
provide the large-boundary tail split.  Therefore a strict numerical bound
on \(\rho_v\) can, in principle, be promoted to a finite proof certificate;
no uncountable scan is required.

\subsection{Propagation of a conformal response}

For an edge \(e\) in the subtree rooted at \(v\), define
\[
 r_{v,e}(Q)=\partial_{d_eQ}\ell_v(Q).                            \tag{30.10}
\]
If \(e\) lies in the child subtree rooted at \(v_j\), differentiation of
the log partition function gives
\[
 \partial_{d_e}\ell_v(Q)
 =P_v\left[\partial_{d_e}\ell_{v_j}\right](Q).                 \tag{30.11}
\]
The additive constant in the child score is immaterial to the Lipschitz
seminorm.

\begin{lemma}[One-step response contraction]
\label{lem:v18v20-step}
If \(e\in E(\mathcal T_{v_j})\), then
\[
 \|r_{v,e}\|_\infty
 \leq\rho_v\|r_{v_j,e}\|_\infty.                              \tag{30.12}
\]
\end{lemma}

\begin{proof}
Differentiate (30.11) in \(Q\) and apply (30.5) with
\(F(t)=\partial_{d_e}\ell_{v_j}(t)\).  Its derivative is
\(r_{v_j,e}(t)\).
\end{proof}

Let \(h(v,e)\) be zero when \(e=e_v\), and otherwise the number of
child-to-parent transports between the vertex carrying \(e\) and \(v\).
Recall from (29.4) that the component of \(c_v\) on its own edge is
\(-a_v\).

\begin{theorem}[Exponential ancestry response]
\label{thm:v18v20-response}
Suppose
\[
 a_v\leq a^*,\qquad \rho_v\leq\rho<1                           \tag{30.13}
\]
for every allowed vertex and boundary value.  Then every component of the
mixed-response vector satisfies
\[
 |(\widetilde c_v)_e|
 \leq a^*\rho^{h(v,e)}
 \quad\text{when }e\in E(\mathcal T_v),                         \tag{30.14}
\]
and is zero otherwise.
\end{theorem}

\begin{proof}
For \(e=e_v\), (30.14) is \(a_v\leq a^*\).  For a proper descendant,
the component \((\widetilde c_v)_e\) is the response (30.10) of the child
message appearing in (29.4).  Apply Lemma \ref{lem:v18v20-step} at every
edge on the ancestral path and finish with the base component \(-a_e\).
\end{proof}

\subsection{Closing the additive Bessel gate}

Let \(b\geq1\) bound the number of children of every vertex and assume
\(\lambda_v\geq\lambda_*>0\).  From (24.6),
\[
 0\leq w_v=\frac{4q_v}{R_v(q_v)}
 \leq\frac{1}{3\lambda_*},                                     \tag{30.15}
\]
where the value at \(q_v=0\) is zero.

\begin{theorem}[Quantitative ancestry gate]
\label{thm:v18v20-gate}
Assume the hypotheses of Theorem \ref{thm:v18v20-response}, the fresh
floor \(p_*>0\) of (29.11), and
\[
 b\rho<1.                                                       \tag{30.16}
\]
Set
\[
 \mathfrak G=
 \frac{(a^*)^2}
 {3\lambda_*p_*(1-b\rho)(1-\rho)}.                             \tag{30.17}
\]
If \(\mathfrak G<1\), then every finite descendant tree satisfies
\[
 B_{\mathcal T}(Q)\preceq\mathfrak G A_{\mathcal T}(Q),
 \qquad
 H_{\mathcal T}(Q)\succeq(1-\mathfrak G)p_*I,                 \tag{30.18}
\]
uniformly in its depth and exterior boundary.
\end{theorem}

\begin{proof}
Fix an edge \(e\).  By (29.7), only common ancestors of \(e\) and \(f\)
contribute to \((B_{\mathcal T})_{ef}\).  Equations (30.14)--(30.15) give
\[
 \sum_f |(B_{\mathcal T})_{ef}|
 \leq\frac{(a^*)^2}{3\lambda_*}
 \sum_{v\succeq e}\rho^{h(v,e)}
 \sum_{f\in E(\mathcal T_v)}\rho^{h(v,f)}.                    \tag{30.19}
\]
There are at most \(b^k\) descendant edges at height \(k\) from \(v\),
so the inner sum is at most \((1-b\rho)^{-1}\).  The ancestral sum is at
most \((1-\rho)^{-1}\).  Hence the right side of (30.19) is at most
\[
 \frac{(a^*)^2}
 {3\lambda_*(1-b\rho)(1-\rho)}.                                \tag{30.20}
\]
Lemma \ref{lem:v18v19-fresh} gives
\((A_{\mathcal T})_{ee}\geq p_*\) for every \(e\).  Divide (30.20) by
the normalized diagonal factors in (29.19).  The resulting row sum is at
most \(\mathfrak G\).  Proposition \ref{prop:v18v19-row} and Theorem
\ref{thm:v18v19-bessel} yield (30.18).
\end{proof}

\section{V20 consequence and certificate problem}
\label{sec:v18v20-next}

The unresolved infinite-depth estimate is now local and one-dimensional.
For finitely many vertex types, it is enough to certify a common bound
\(\rho_v\leq\rho\) from (30.3)--(30.4) and check the explicit inequalities
(30.16)--(30.17).  These conditions are sufficient and may be conservative;
if they fail, the exact matrix gate (29.18) can still hold.

This result also distinguishes the two large-boundary tasks.  The moment
growth bounds of V17 guarantee integrability and compactification of the
\(Q\)-axis.  The remaining numerical inequality concerns the derivative of
the entire normalized radial distribution, rather than a single moment.
That coefficient is the next object to enclose rigorously.

\begin{assessment}[V20 acquisition and boundary]
\label{ass:v18v20}
V20 gives an exact formula for the radial Markov operator's Lipschitz norm,
proves exponential decay of conformal mixed responses under a uniform local
contraction, and derives an explicit depth-independent Bessel criterion.
The strict coefficient bounds required by (30.13), (30.16), and (30.17)
have not yet been certified for the physical parameter card.  Cycles,
cutoff removal, reconstruction, and the Einstein continuum remain open.
\end{assessment}

\begin{record}[V20 append record]
\label{rec:v18v20}
V20 is appended to the validated V19 whose installed SHA--256 was
\texttt{88482CF7E73ACC94A7BBC1F92EBDA8FF3B240AF102A8917B24A6BF2476935D67}.
After installation only V20 is the active numeric SM note; the frozen
\texttt{V100\_f17} archive is unchanged.
\end{record}

% =====================================================================
% V21 APPENDIX: AUDIT AND MONOTONE RADIAL RESPONSE
% =====================================================================

\section{V21: companion audit and scalarization of the contraction gate}
\label{sec:v18v21}

V20 was the fifth version after the V15 audit, so the newest active
companion notes were inspected before this continuation.  The active YM
note was
\[
 \texttt{Renewal\_Yang\_Mills\_Mass\_Gap\_Research\_Notes\_V3.tex}
\]
with SHA--256
\[
 \texttt{1E0E12B0BBC002F5B652819D6ACB4C99AD6CB0B3312DBBEE2B6617F461E1858A}.
                                                               \tag{31.1}
\]
Its new relevant device is to multiply Laurent coefficients by a local
trial inverse before taking norms and to retain signed Taylor centers until
the final error estimate.  The active NS note was
\[
 \texttt{Renewal\_NS\_Stability\_Research\_Notes\_V26.tex}
\]
with SHA--256
\[
 \texttt{82C887F8C2C69E4F28FAD7C3DC8DE5B9099CB5A4BF431EBA1FFF3E91935978AD}.
                                                               \tag{31.2}
\]
It computes exact rank-one kernel norms and shows both that rank-one geometry
can improve a full operator norm and that the improvement need not occur in
every derivative channel.  Neither companion note proves an SM radial
contraction.  The methodological consequence here is to preserve the signed
cumulative radial density in the compact certificate and to test the exact
normalized loss matrix before falling back to the V20 row-sum bound.

This note now simplifies the V20 coefficient analytically.  Monotone
likelihood ratio structure forces its cumulative derivative to have one
sign.  The exact Lipschitz norm is therefore just the slope of the radial
second moment.

\subsection{Monotone likelihood ratio of the angular tilt}

Write
\[
 A_\alpha(s)={}_0F_1(;\alpha;s),
 \qquad
 h_Q(t)=\partial_Q\log A_\alpha(\chi^2Qt).                    \tag{31.3}
\]
The factors in (30.1) other than \(A_\alpha(\chi^2Qt)\) do not depend on
\(Q\).

\begin{lemma}[Monotonicity of the angular score]
\label{lem:v18v21-angular}
For every \(Q\geq0\), the function \(t\mapsto h_Q(t)\) is nondecreasing.
It is strictly increasing for \(\chi>0\).
\end{lemma}

\begin{proof}
For \(Q>0\), use the normalized spherical representation
\[
 A_\alpha(s)=\int e^{2\sqrt{s}\,X}\,d\sigma(X),                \tag{31.4}
\]
where \(X\) is the first coordinate on the sphere of dimension
\(2\alpha-1\).  Put
\[
 M(r)=\log\int e^{rX}\,d\sigma(X),
 \qquad r=2\chi\sqrt{Qt}.
\]
Symmetry and exponential tilting give
\[
 M'(r)=\mathbb E_rX\geq0,
 \qquad M''(r)=\operatorname{Var}_r(X)\geq0                    \tag{31.5}
\]
for \(r\geq0\).  Direct differentiation yields
\[
 h_Q(t)=\frac{rM'(r)}{2Q}.                                     \tag{31.6}
\]
Since \((rM'(r))'=M'(r)+rM''(r)\geq0\), (31.6) is nondecreasing in
\(t\).  Strictness follows for a nondegenerate sphere and \(\chi>0\).
At \(Q=0\), the hypergeometric series gives
\[
 h_0(t)=\frac{\chi^2}{\alpha}t,                               \tag{31.7}
\]
which has the same property.
\end{proof}

Let \(p_Q=p_{v,Q}\), suppressing the vertex label.  Differentiating its
normalization gives
\[
 \partial_Qp_Q(t)
 =p_Q(t)\left(h_Q(t)-\mathbb E_Qh_Q\right).                    \tag{31.8}
\]

\begin{proposition}[One-sign cumulative derivative]
\label{prop:v18v21-sign}
The cumulative derivative (30.3) satisfies
\[
 G_Q(s)=\int_0^s\partial_Qp_Q(t)\,dt\leq0
 \qquad(s\geq0).                                               \tag{31.9}
\]
\end{proposition}

\begin{proof}
By Lemma \ref{lem:v18v21-angular}, the conditional average of \(h_Q\) on
\([0,s]\) is no larger than its unconditional average.  Hence
\[
 G_Q(s)=\mathbb P_Q(t\leq s)
 \left(\mathbb E_Q[h_Q\mid t\leq s]-\mathbb E_Qh_Q\right)
 \leq0.                                                        \tag{31.10}
\]
The statement is trivial when the probability of the interval is zero.
\end{proof}

\subsection{The exact coefficient is a mean slope}

Put
\[
 m_v(Q)=\mathbb E_{v,Q}t.                                     \tag{31.11}
\]

\begin{theorem}[Scalarization of the V20 contraction coefficient]
\label{thm:v18v21-scalar}
For every allowed radial kernel,
\[
 \int_0^\infty|G_{v,Q}(s)|\,ds=m_v'(Q),                       \tag{31.12}
\]
and consequently
\[
 \boxed{\quad \rho_v=\sup_{Q\geq0}m_v'(Q).\quad}              \tag{31.13}
\]
In particular,
\[
 m_v'(Q)=\operatorname{Cov}_{v,Q}(t,h_Q(t))\geq0,              \tag{31.14}
\]
while at the origin
\[
 m_v'(0)=\frac{\chi_v^2}{\alpha_v}
          \operatorname{Var}_{v,0}(t).                         \tag{31.15}
\]
\end{theorem}

\begin{proof}
Proposition \ref{prop:v18v21-sign} gives
\[
 \int_0^\infty|G_Q(s)|\,ds=-\int_0^\infty G_Q(s)\,ds.         \tag{31.16}
\]
Integrate first on \([0,L]\).  Fubini's theorem gives
\[
 -\int_0^LG_Q(s)\,ds
 =-L\int_0^L\partial_Qp_Q(t)\,dt
   +\int_0^L t\partial_Qp_Q(t)\,dt.                            \tag{31.17}
\]
Quartic tails make the first term tend to zero and permit passage to the
limit in the second.  The limit is
\(\partial_Q\mathbb E_Qt=m_v'(Q)\), proving (31.12).  Equation (31.13)
follows from (30.4).  Formula (31.14) is (31.8) integrated against \(t\),
and (31.15) follows from (31.7).
\end{proof}

This identity is stronger than estimating the absolute cumulative density
directly.  The signed center in (31.8) is retained until monotonicity fixes
the sign, paralleling the useful certification lesson of the YM audit.

\subsection{A rigorous large-boundary cutoff}

The scalar slope also admits a direct tail estimate.  Let
\(y=\sqrt Q\,e\), where \(|e|=1\), and let \(\mathcal H_z\) denote the
field Hessian of the internal action.  Under the V14 hypotheses that the
child radial messages are convex and nondecreasing,
\[
 \mathcal H_z\succeq2a_vI,
 \qquad
 (\nabla t)^{\mathsf T}\mathcal H_z^{-1}\nabla t
 =\frac{4t}{R_v(t)}\leq\frac{1}{3\lambda_v}.                    \tag{31.18}
\]

\begin{proposition}[Large-boundary response tail]
\label{prop:v18v21-tail}
For \(Q>0\),
\[
 0\leq m_v'(Q)
 \leq\frac{\chi_v}{\sqrt{6a_v\lambda_vQ}}.                    \tag{31.19}
\]
At the origin,
\[
 0\leq m_v'(0)\leq\frac{\chi_v^2}{3\alpha_v\lambda_v}.        \tag{31.20}
\]
\end{proposition}

\begin{proof}
Differentiation in the radial boundary variable \(r=\sqrt Q\) gives
\[
 \frac{d}{dr}\mathbb E_rt
 =2\chi_v\operatorname{Cov}_r(t,z\cdot e),
\]
and hence
\[
 m_v'(Q)=\frac{\chi_v}{\sqrt Q}
          \operatorname{Cov}_{v,Q}(t,z\cdot e).                \tag{31.21}
\]
Brascamp--Lieb and (31.18) imply
\[
 \operatorname{Var}(t)\leq\frac{1}{3\lambda_v},
 \qquad
 \operatorname{Var}(z\cdot e)\leq\frac{1}{2a_v}.              \tag{31.22}
\]
Cauchy--Schwarz in (31.21) proves (31.19).  At \(Q=0\), combine
(31.15) with the first inequality in (31.22).
\end{proof}

\begin{corollary}[Compact certificate for a target contraction]
\label{cor:v18v21-compact}
On a coefficient sector with
\[
 \chi_v\leq\chi^*,\qquad a_v\geq a_*,\qquad
 \lambda_v\geq\lambda_*,                                      \tag{31.23}
\]
fix a target \(\rho_{\rm tar}>0\) and choose
\[
 Q_*\geq\frac{(\chi^*)^2}
 {6a_*\lambda_*\rho_{\rm tar}^2}.                             \tag{31.24}
\]
If certified one-dimensional quadrature proves
\[
 m_v'(Q)\leq\rho_{\rm tar}
 \qquad(0\leq Q\leq Q_*)                                      \tag{31.25}
\]
for every allowed vertex type and parameter box, then
\(\rho_v\leq\rho_{\rm tar}\) globally.
\end{corollary}

\begin{proof}
Equation (31.25) covers the compact interval, including the analytic anchor
(31.15).  Equations (31.19) and (31.24) cover \(Q\geq Q_*\).  Apply
Theorem \ref{thm:v18v21-scalar}.
\end{proof}

\section{V21 consequence and next computation}
\label{sec:v18v21-next}

The V20 coefficient no longer requires enclosing an absolute integral of a
cumulative density.  It is the maximum slope of the scalar function
\(Q\mapsto\mathbb E_{v,Q}|z_v|^2\).  Proposition
\ref{prop:v18v21-tail} turns the unbounded boundary domain into the compact
certificate (31.25).  On that compact set, (31.14) is a ratio of radial
integrals with the same quartic weight, so the V12 relative enclosures can
retain their common normalization and signed covariance center.

The immediate missing input is a physical compact coefficient card.  The
Grand Tensor source leaves the positive Higgs coefficients as physical
parameters, while the exact tuple in the related finite \(N=2\) obstruction
test is not a continuum card.  Until such a card is derived, the rigorous
next step is a symbolic interval certificate in the dimensionless
coefficient box, recording the subregion where (30.16)--(30.17) close.

\begin{assessment}[V21 acquisition and boundary]
\label{ass:v18v21}
V21 records the mandatory companion audit, proves monotone likelihood ratio
for the angular tilt, identifies the exact contraction coefficient with a
radial mean slope, and supplies explicit origin and large-boundary bounds.
No companion theorem was imported.  The compact strict inequality and a
physical continuum coefficient card remain unproved; cycles, cutoff
removal, reconstruction, and the Einstein continuum remain open.
\end{assessment}

\begin{record}[V21 append record]
\label{rec:v18v21}
V21 is appended to the validated V20 whose installed SHA--256 was
\texttt{693D975A2B6418254467E174DA7A1FA95EB9D138FAB3AD24A468223A11D25390}.
After installation only V21 is the active numeric SM note; the frozen
\texttt{V100\_f17} archive is unchanged.
\end{record}

% =====================================================================
% V22 APPENDIX: POSITIVE SERIES CERTIFICATE FOR THE RADIAL SLOPE
% =====================================================================

\section{V22: a finite positive certificate for the radial mean slope}
\label{sec:v18v22}

V21 reduced the local contraction gate to the maximum of
\(m'(Q)\).  Direct interval evaluation of the covariance (31.14) can lose
the strict margin by subtracting two close positive numbers.  This note
uses the angular expansion to rewrite the derivative as a sum of positive
pair determinants.  Truncation then gives a monotone lower stock and a
one-sided, subfactorial upper remainder.  This implements the
cancellation-retaining lesson of the V20 companion audit without importing
any YM statement.

\subsection{Angular mixture representation}

Suppress the vertex label and put
\[
 \Psi(t)=\sum_{j=1}^r\ell_j(\theta_j,t),
 \qquad
 J_s=\int_0^\infty
 t^{s-1}e^{-\lambda t^2-at-\Psi(t)}\,dt.                       \tag{32.1}
\]
For \(k\geq0\), define
\[
 c_k(Q)=\frac{(\chi^2Q)^k}{k!(\alpha)_k},
 \qquad
 W_k(Q)=c_k(Q)J_{\alpha+k},
 \qquad
 \mu_k=\frac{J_{\alpha+k+1}}{J_{\alpha+k}}.                  \tag{32.2}
\]
The hypergeometric expansion in (30.1) gives
\[
 Z(Q)=\sum_{k=0}^\infty W_k(Q),
 \qquad
 m(Q)=\frac{\sum_{k=0}^\infty c_k(Q)J_{\alpha+k+1}}{Z(Q)}.    \tag{32.3}
\]
Equivalently, if \(\pi_Q(k)=W_k(Q)/Z(Q)\), then
\[
 m(Q)=\mathbb E_{\pi_Q}\mu_K.                                \tag{32.4}
\]

\begin{lemma}[Strict increase of the shape mean]
\label{lem:v18v22-shape}
The sequence \(k\mapsto\mu_k\) is strictly increasing.
Consequently, for \(0\leq k<l\),
\[
 \Delta_{kl}:=
 J_{\alpha+k}J_{\alpha+l+1}
 -J_{\alpha+k+1}J_{\alpha+l}>0.                              \tag{32.5}
\]
\end{lemma}

\begin{proof}
For real \(s>0\), let \(\nu_s\) be the probability measure proportional to
the integrand in (32.1).  Differentiation in \(s\) gives
\[
 \frac{d}{ds}\frac{J_{s+1}}{J_s}
 =\operatorname{Cov}_{\nu_s}(t,\log t)>0,                    \tag{32.6}
\]
because both functions are strictly increasing and \(\nu_s\) is
nondegenerate.  Thus \(J_{s+1}/J_s\) is strictly increasing in \(s\).
Equation (32.5) is the same inequality after clearing positive
denominators.
\end{proof}

\begin{theorem}[Positive pair formula]
\label{thm:v18v22-pair}
For \(Q>0\),
\[
 m'(Q)=\frac{1}{Q}\operatorname{Cov}_{\pi_Q}(K,\mu_K)          \tag{32.7}
\]
and, more explicitly,
\[
 \boxed{
 m'(Q)=\frac{1}{QZ(Q)^2}
 \sum_{0\leq k<l}
 c_k(Q)c_l(Q)(l-k)\Delta_{kl}.}
                                                                    \tag{32.8}
\]
Every summand in (32.8) is positive.  After cancelling the displayed
factor \(Q\), the series extends continuously to \(Q=0\) and agrees with
(31.15).
\end{theorem}

\begin{proof}
Only \(c_k(Q)\) depends on \(Q\), and
\(c_k'(Q)=kc_k(Q)/Q\).  Differentiate (32.4) to obtain (32.7).  The pairwise
covariance identity gives
\[
 \operatorname{Cov}_{\pi_Q}(K,\mu_K)
 =\sum_{k<l}\pi_Q(k)\pi_Q(l)(l-k)(\mu_l-\mu_k).                \tag{32.9}
\]
Insert (32.2) and clear the moment-ratio denominators to obtain (32.8).
Positivity follows from Lemma \ref{lem:v18v22-shape}.  The least power in
the numerator is the \((k,l)=(0,1)\) term, which is proportional to \(Q\).
Its quotient at zero is
\[
 \frac{\chi^2}{\alpha}
 \left(\frac{J_{\alpha+2}}{J_\alpha}
 -\frac{J_{\alpha+1}^2}{J_\alpha^2}\right),                  \tag{32.10}
\]
the variance formula (31.15); all remaining terms have higher powers.
\end{proof}

\subsection{A one-sided truncation remainder}

For an integer \(L\geq1\), put
\[
 Z_L=\sum_{k=0}^LW_k,
 \qquad
 P_L=\sum_{0\leq k<l\leq L}
 c_kc_l(l-k)\Delta_{kl},                                     \tag{32.11}
\]
and
\[
 U_l=l c_lJ_{\alpha+l+1},
 \qquad
 R_L=\sum_{l>L}U_l.                                           \tag{32.12}
\]

\begin{lemma}[Positive omitted-pair bound]
\label{lem:v18v22-tailpair}
For \(Q>0\),
\[
 0\leq m'(Q)-\frac{P_L}{QZ(Q)^2}
 \leq\frac{R_L}{QZ(Q)}.                                      \tag{32.13}
\]
\end{lemma}

\begin{proof}
The omitted pairs in (32.9) have \(l>L\).  Since \(k<l\) and
\(0\leq\mu_k\leq\mu_l\),
\[
 (l-k)(\mu_l-\mu_k)\leq l\mu_l.                              \tag{32.14}
\]
Thus their total contribution to \(Qm'(Q)\) is at most
\[
 \sum_{l>L}\pi_Q(l)l\mu_l\sum_{k<l}\pi_Q(k)
 \leq\sum_{l>L}\pi_Q(l)l\mu_l
 =\frac{R_L}{Z(Q)}.                                           \tag{32.15}
\]
The lower inequality follows by discarding positive pairs.
\end{proof}

The moment-ratio envelope of V12 gives
\[
 \frac{J_{s+1}}{J_s}\leq\sqrt{\frac{s}{2\lambda}}.           \tag{32.16}
\]
Hence, on \(0\leq Q\leq Q_*\), \(\chi\leq\chi^*\), and
\(\lambda\geq\lambda_*\),
\[
 \frac{U_{l+1}}{U_l}
 \leq
 \frac{(\chi^*)^2Q_*}{l(\alpha+l)}
 \sqrt{\frac{\alpha+l+1}{2\lambda_*}}
 =:\vartheta_l.                                               \tag{32.17}
\]

\begin{proposition}[Geometric tail stock]
\label{prop:v18v22-geometric}
If
\[
 \vartheta_L^*:=\sup_{l\geq L+1}\vartheta_l<1,                \tag{32.18}
\]
then
\[
 R_L\leq\frac{U_{L+1}}{1-\vartheta_L^*}.                      \tag{32.19}
\]
Consequently,
\[
 m'(Q)\leq
 \frac{P_L}{QZ_L^2}
 +\frac{U_{L+1}}
 {QZ_L(1-\vartheta_L^*)}.                                     \tag{32.20}
\]
The right side has a continuous value at \(Q=0\) after cancelling the
known powers of \(Q\).
\end{proposition}

\begin{proof}
Equation (32.17) and (32.18) bound the successive terms of (32.12) by a
geometric sequence, proving (32.19).  In (32.13), use \(Z\geq Z_L\), then
apply (32.19).  In the first term of (32.20), each summand of \(P_L/Q\) is
proportional to \(Q^{k+l-1}\); in the second,
\(U_{L+1}/Q\) is proportional to \(Q^L\).  Thus neither quotient is
singular at zero.
\end{proof}

\subsection{Finite interval implementation}

\begin{theorem}[Finite certified contraction test]
\label{thm:v18v22-certificate}
Fix a compact positive coefficient box, a finite family of allowed vertex
types, a target \(\rho_{\rm tar}>0\), and \(Q_*\) satisfying (31.24).
Choose \(L\) so that (32.18) holds.  Suppose the positive relative
enclosures of V12, on a finite subdivision of the coefficient box and
\([0,Q_*]\), certify
\[
 \frac{P_L}{QZ_L^2}
 +\frac{U_{L+1}}
 {QZ_L(1-\vartheta_L^*)}
 \leq\rho_{\rm tar}.                                          \tag{32.21}
\]
Here the powers of \(Q\) are cancelled symbolically on cells meeting zero.
Then
\[
 \rho_v\leq\rho_{\rm tar}                                    \tag{32.22}
\]
for every allowed vertex.  If in addition (30.16)--(30.17) hold with
\(\rho=\rho_{\rm tar}\), the finite-tree conformal Hessians have the
uniform margin (30.18).
\end{theorem}

\begin{proof}
All quantities in (32.21) are finite sums or positive radial integrals.
The V12 derivative envelopes control their variation on each parameter
cell, while Proposition \ref{prop:v18v22-geometric} controls the omitted
angular tail uniformly.  Thus (32.21) and (32.20) give
\(m_v'(Q)\leq\rho_{\rm tar}\) on the compact boundary interval.  The
large-boundary estimate (31.19) and the choice (31.24) give the same bound
for \(Q\geq Q_*\).  Theorem \ref{thm:v18v21-scalar} proves (32.22), and
Theorem \ref{thm:v18v20-gate} gives the final assertion.
\end{proof}

The determinant form (32.5) should be enclosed as a unit, or with a common
relative center, rather than by independently subtracting four absolute
interval endpoints.  This is the radial analogue of applying the YM V3
residual weight before taking a norm.  If a coarse interval still crosses
zero, log-convexity supplies its nonnegative lower endpoint and subdivision
sharpens the upper endpoint.

\section{V22 consequence and remaining input}
\label{sec:v18v22-next}

The local contraction gate is now a finite proof object.  Given a compact
coefficient box, the certificate consists of finitely many enclosures of
\(J_s\), the positive pair stock \(P_L\), the positive normalizer stock
\(Z_L\), and the single geometric remainder (32.19).  The computation
never subtracts the two large terms in a covariance.

The current manuscripts still do not select a canonical numerical Higgs
coefficient card for the continuum.  V23 should therefore formulate the
dimensionless admissible region cut out by (31.24), (32.21), and
(30.16)--(30.17), and separate universal symbolic inequalities from the
eventual physical matching conditions.

\begin{assessment}[V22 acquisition and boundary]
\label{ass:v18v22}
V22 proves a positive angular-pair formula for the exact radial contraction
slope, a rigorous one-sided truncation remainder, and a finite interval
certificate which closes the tree Bessel gate whenever its strict
inequalities hold.  No numerical physical coefficient card is asserted,
and the strict certificate has not yet been executed on such a card.
Cycles, cutoff removal, reconstruction, and the Einstein continuum remain
open.
\end{assessment}

\begin{record}[V22 append record]
\label{rec:v18v22}
V22 is appended to the validated V21 whose installed SHA--256 was
\texttt{1361A6469364F5FB00F9EAA4768EE7742C05F86C892777927B725E2FCBBBE57C}.
After installation only V22 is the active numeric SM note; the frozen
\texttt{V100\_f17} archive is unchanged.
\end{record}

% =====================================================================
% V23 APPENDIX: DIMENSIONLESS ADMISSIBLE REGION
% =====================================================================

\section{V23: the dimensionless admissible region}
\label{sec:v18v23}

The certificate of V22 is written in coefficient variables whose numerical
values depend on the field unit.  This note removes that ambiguity and
defines the strict local region needed by the cofinal tree theorem.  It also
proves two structural facts: the region is open, so rational interval boxes
can lie strictly inside it, and its local radial gate is nonempty near a
weak-hopping quartic anchor.  Closure of a recursively generated message
family inside that region remains a separate condition.

\subsection{Common field dilation}

Fix a reference quartic scale \(\lambda_0>0\) and set
\[
 u=\sqrt{\lambda_0}\,t,
 \qquad V=\sqrt{\lambda_0}\,Q.                                \tag{33.1}
\]
Define dimensionless coefficients and potential
\[
 \widehat\lambda=\frac{\lambda}{\lambda_0},
 \qquad x=\frac{a}{\sqrt{\lambda_0}},
 \qquad g=\frac{\chi}{\sqrt{\lambda_0}},
 \qquad \Phi(u)=\Psi\left(\frac{u}{\sqrt{\lambda_0}}\right).  \tag{33.2}
\]
Apart from its normalizing constant, (30.1) becomes
\[
 \widehat p_V(u)
 \propto u^{\alpha-1}
 e^{-\widehat\lambda u^2-xu-\Phi(u)}
 {}_0F_1(;\alpha;g^2Vu).                                      \tag{33.3}
\]
Let \(\widehat m(V)=\mathbb E_Vu\).

\begin{proposition}[Scale invariance of the contraction]
\label{prop:v18v23-scale}
Under (33.1)--(33.3),
\[
 m(Q)=\lambda_0^{-1/2}\widehat m(V),
 \qquad m'(Q)=\widehat m'(V),
 \qquad \rho=\sup_{V\geq0}\widehat m'(V).                    \tag{33.4}
\]
The fresh edge quantity is also invariant:
\[
 am(Q)=x\widehat m(V).                                        \tag{33.5}
\]
Thus \(\rho\) and the product \(am\) depend only on the dimensionless
kernel data in (33.2)--(33.3).
\end{proposition}

\begin{proof}
The substitution \(t=u/\sqrt{\lambda_0}\) in the radial integral gives
the first identity in (33.4).  Since
\(dV/dQ=\sqrt{\lambda_0}\), the two scale factors cancel on
differentiation.  Taking the supremum proves the last identity, and
multiplication by \(a=x\sqrt{\lambda_0}\) gives (33.5).
\end{proof}

For example, the origin and tail estimates of V21 become
\[
 \widehat m'(0)\leq\frac{g^2}{3\alpha\widehat\lambda},
 \qquad
 \widehat m'(V)
 \leq\frac{g}{\sqrt{6x\widehat\lambda V}}
 \quad(V>0).                                                   \tag{33.6}
\]
These expressions contain no lattice or field-unit scale.

\subsection{Definition of the strict gate region}

Let \(\mathfrak K\) be a compact family of dimensionless kernels (33.3)
with finitely many values of \(\alpha\), uniformly positive \(x\) and
\(\widehat\lambda\), and a finite-dimensional compact family of convex,
nondecreasing \(\Phi\)'s satisfying a common quartic domination bound.
Put
\[
 \rho_{\mathfrak K}=\sup_{\kappa\in\mathfrak K}\sup_{V\geq0}
 \widehat m_\kappa'(V),                                       \tag{33.7}
\]
\[
 p_{\mathfrak K}=inf_{\kappa\in\mathfrak K}
 x_\kappa\widehat m_\kappa(V)>0,                             \tag{33.8}
\]
where the infimum in (33.8) also ranges over \(V\geq0\), and define the
coarse amplitude stock
\[
 \Xi_{\mathfrak K}
 =\frac{x_{max}^2}
 {3\widehat\lambda_{min}p_{\mathfrak K}}.                    \tag{33.9}
\]
The moment floor of V18 allows the infimum in (33.8) to be replaced by the
explicit lower stock
\[
 p_{\mathfrak K}^{\rm floor}
 =\inf_{\kappa\in\mathfrak K}
 x_\kappa\widehat m_{\alpha_\kappa}
 (x_\kappa+B_\kappa,\widehat\lambda_\kappa)>0.                \tag{33.10}
\]

\begin{definition}[Dimensionless radial admissible region]
\label{def:v18v23-region}
For a branching bound \(b\geq1\), the family \(\mathfrak K\) is in
\(\mathcal A_b\) if
\[
 b\rho_{\mathfrak K}<1,
 \qquad
 \frac{\Xi_{\mathfrak K}}
 {(1-b\rho_{\mathfrak K})(1-\rho_{\mathfrak K})}<1.           \tag{33.11}
\]
\end{definition}

The second inequality is precisely the V20 gate (30.17), expressed in the
common units (33.1).  Using (33.10) in place of (33.8) makes it more
conservative but directly certifiable.

\begin{theorem}[Uniform tree margin in the admissible region]
\label{thm:v18v23-region}
Suppose every conditional vertex kernel of every finite regulator belongs
to one fixed compact family \(\mathfrak K\in\mathcal A_b\), and the V14
recursive transport hypotheses hold.  Then there is a constant
\(\mu_{\mathfrak K}>0\), independent of the regulator depth and boundary,
such that
\[
 H_{\mathcal T}(Q)\succeq\mu_{\mathfrak K}I.                  \tag{33.12}
\]
One admissible value is
\[
 \mu_{\mathfrak K}
 =p_{\mathfrak K}
 \left[1-
 \frac{\Xi_{\mathfrak K}}
 {(1-b\rho_{\mathfrak K})(1-\rho_{\mathfrak K})}
 \right].                                                     \tag{33.13}
\]
\end{theorem}

\begin{proof}
Proposition \ref{prop:v18v23-scale} permits evaluation in the common
dimensionless variables.  The first inequality in (33.11) is (30.16), and
the fraction in its second inequality is the constant \(\mathfrak G\) in
(30.17).  Apply Theorem \ref{thm:v18v20-gate} and then return to the
original common field units.  The conformal parameters were not rescaled,
so the Hessian margin is (33.13).
\end{proof}

\subsection{Openness and a nonempty local anchor}

\begin{theorem}[Robustness of the strict gate]
\label{thm:v18v23-open}
Within the finite-dimensional kernel class just described, the maps
\[
 \kappa\longmapsto\rho_\kappa,
 \qquad
 \kappa\longmapsto\inf_{V\geq0}x_\kappa\widehat m_\kappa(V) \tag{33.14}
\]
are continuous on compact positive coefficient sets.  Consequently, the
strict inequalities (33.11) define an open set of kernel data.  Every
compact subset of that open set has a common positive margin in (33.13).
\end{theorem}

\begin{proof}
On a fixed compact \(V\)-interval, the quartic envelope permits dominated
differentiation of (33.3).  Hence \(\widehat m\) and
\(\widehat m'\) are jointly continuous in \((\kappa,V)\).  The tail bound
(33.6) tends to zero uniformly on a compact positive coefficient set as
\(V\to\infty\).  Therefore the supremum in (33.7) is the uniform limit of
suprema over compact intervals and is continuous.  The moment floor (28.5)
and the same dominated-convergence argument give continuity and strict
positivity of the second map in (33.14).  The quantities in (33.11) are
then continuous wherever their denominators are positive.  Strict
inequalities are stable under a sufficiently small perturbation, proving
the result.
\end{proof}

\begin{proposition}[Nonempty weak-hopping radial gate]
\label{prop:v18v23-nonempty}
Consider the homogeneous local kernel with
\[
 \widehat\lambda=1,
 \qquad \Phi=0,
 \qquad x>0,
 \qquad g\geq0.                                                \tag{33.15}
\]
For every fixed branching number \(b\), there are sufficiently small
positive \(x\) and \(g\) for which its one-type family satisfies (33.11).
Thus the dimensionless local radial gate is nonempty.
\end{proposition}

\begin{proof}
At \(g=0\), the density (33.3) is independent of \(V\), so
\(\rho=0\).  Let
\[
 \overline m_x=
 \frac{\int_0^\infty u^\alpha e^{-u^2-xu}\,du}
      {\int_0^\infty u^{\alpha-1}e^{-u^2-xu}\,du}.             \tag{33.16}
\]
Then \(p=x\overline m_x\) and
\[
 \Xi=\frac{x}{3\overline m_x}\longrightarrow0
 \qquad(x\downarrow0),                                       \tag{33.17}
\]
because \(\overline m_x\) converges to the positive pure-quartic moment
\[
 \frac{\Gamma((\alpha+1)/2)}{\Gamma(\alpha/2)}.               \tag{33.18}
\]
Choose \(x>0\) so that \(\Xi<1\).  By Theorem
\ref{thm:v18v23-open}, \(\rho\) and \(\Xi\) vary continuously with
\(g\); hence all sufficiently small positive \(g\) preserve both strict
inequalities in (33.11).
\end{proof}

The last proposition proves nonemptiness only of the local kernel gate.  It
does not assert that the physical SM coefficients lie there, nor that the
recursive message transform preserves a neighborhood of (33.15).  Those
are the matching and invariant-family problems that must be addressed next.

\section{V23 consequence and upstream split}
\label{sec:v18v23-next}

The cofinal tree analysis now has a unit-free target.  There are two
logically distinct remaining tasks:
\begin{enumerate}[leftmargin=2.2em]
\item derive the physical dimensionless coefficient box from the
      Standard Model--Einstein discretization and test it with V22;
\item prove that the recursive message transform maps a certified compact
      kernel family into itself, so the hypothesis of Theorem
      \ref{thm:v18v23-region} holds at arbitrary depth.
\end{enumerate}
The source audit shows that the first task cannot be replaced by the finite
\(N=2\) obstruction tuple.  The second is purely analytic and can proceed
symbolically: V24 should formulate an invariant cone for
\((\widehat m,\widehat m',\Phi',\Phi'')\) and test its one-step closure.

\begin{assessment}[V23 acquisition and boundary]
\label{ass:v18v23}
V23 proves scale invariance of the radial contraction, defines the strict
dimensionless admissible region, proves its robustness, and proves that its
local radial gate is nonempty near a weak-hopping quartic anchor.  Recursive
invariance and physical coefficient matching remain unproved.  Cycles,
cutoff removal, reconstruction, and the coupled Einstein continuum remain
open.
\end{assessment}

\begin{record}[V23 append record]
\label{rec:v18v23}
V23 is appended to the validated V22 whose installed SHA--256 was
\texttt{95E62F48751D94414EFFF63CEDE105B75367C3489A4F2EC239B31731D0FC0A67}.
After installation only V23 is the active numeric SM note; the frozen
\texttt{V100\_f17} archive is unchanged.
\end{record}

% =====================================================================
% V24 APPENDIX: ANALYTIC CONTRACTION FROM INHERITED STIFFNESS
% =====================================================================

\section{V24: analytic contraction from inherited radial stiffness}
\label{sec:v18v24}

V22 expressed the exact contraction as
\(Q^{-1}\operatorname{Cov}(K,\mu_K)\).  Two estimates already proved in
V3 can be combined to control this covariance without numerical
quadrature: the angular count is ultra-log-concave, and the radial shape
mean has a bounded increment.  The combination is new here and yields a
global response bound.  More importantly, any positive slope inherited
from child messages augments the linear stiffness and improves the bound at
every internal vertex.

\subsection{A covariance lemma for the angular mixture}

\begin{lemma}[Increasing discrete Lipschitz covariance]
\label{lem:v18v24-covariance}
Let \(K\) be square-integrable and let \(f:\mathbb Z_{\geq0}\to\mathbb R\)
be nondecreasing with
\[
 0\leq f(k+1)-f(k)\leq L.                                     \tag{34.1}
\]
Then
\[
 0\leq\operatorname{Cov}(K,f(K))
 \leq L\operatorname{Var}(K).                                \tag{34.2}
\]
\end{lemma}

\begin{proof}
For independent copies \(K,K'\),
\[
 \operatorname{Cov}(K,f(K))
 =\frac12\mathbb E[(K-K')(f(K)-f(K'))].                       \tag{34.3}
\]
Monotonicity makes the integrand nonnegative.  The increment bound gives
\(|f(K)-f(K')|\leq L|K-K'|\), so (34.3) is at most
\(L\mathbb E(K-K')^2/2=L\operatorname{Var}(K)\).
\end{proof}

Let the radial action be
\[
 \lambda t^2+at+\Psi(t),
 \qquad \Psi'(t)\geq\sigma\geq0.                              \tag{34.4}
\]
Write
\[
 A=a+\sigma,
 \qquad \widetilde\Psi(t)=\Psi(t)-\sigma t.                  \tag{34.5}
\]
Then \(\widetilde\Psi\) is convex and nondecreasing, and the V3 radial
recurrence applies with linear coefficient \(A\).

\begin{theorem}[Stiffness-improved radial contraction]
\label{thm:v18v24-stiffness}
For the radial kernel (30.1) satisfying (34.4),
\[
 0\leq m'(Q)\leq\frac{\chi^2}{(a+\sigma)^2}
 \qquad(Q\geq0).                                               \tag{34.6}
\]
Consequently,
\[
 \boxed{\quad\rho\leq
 \left(\frac{\chi}{a+\sigma}\right)^2.\quad}                 \tag{34.7}
\]
\end{theorem}

\begin{proof}
Use the angular law \(\pi_Q\) and means \(\mu_k\) of (32.2), now with
the decomposition (34.5).  Lemma \ref{lem:v18v3-recurrence}, with \(A\)
in place of \(a\), gives
\[
 0<\mu_{k+1}-\mu_k\leq\frac1A.                               \tag{34.8}
\]
The closed radial kernel theorem gives the angular count estimates
\[
 \operatorname{Var}_{\pi_Q}(K)leq\mathbb E_{\pi_Q}K
 \leq\frac{\chi^2Q}{A}.                                      \tag{34.9}
\]
For \(Q>0\), Theorem \ref{thm:v18v22-pair} and Lemma
\ref{lem:v18v24-covariance} imply
\[
 m'(Q)=\frac1Q\operatorname{Cov}_{\pi_Q}(K,\mu_K)
 \leq\frac{1}{AQ}\operatorname{Var}_{\pi_Q}(K)
 \leq\frac{\chi^2}{A^2}.                                    \tag{34.10}
\]
Analyticity at \(Q=0\) extends the inequality to the origin.  Taking the
supremum proves (34.7).
\end{proof}

The bound is exact for a Gaussian radial kernel with no quartic or child
curvature.  Positive quartic and child curvature may make the true
coefficient strictly smaller; V22 remains available when (34.7) is too
coarse.

\subsection{Stiffness generated by descendants}

At an internal vertex \(v\), the child part of the radial potential is
\[
 \Psi_v(t)=\sum_{j=1}^{r_v}\ell_{v,j}(t)+\lambda_vt^2.         \tag{34.11}
\]
The V4 boundary-stiffness estimate gives certified constants
\(\gamma_{v,j}>0\) such that
\[
 \ell_{v,j}'(t)\geq\gamma_{v,j}.                              \tag{34.12}
\]
Therefore
\[
 \Psi_v'(t)\geq\sigma_v,
 \qquad \sigma_v:=\sum_{j=1}^{r_v}\gamma_{v,j}.               \tag{34.13}
\]

\begin{corollary}[Internal-vertex contraction stock]
\label{cor:v18v24-internal}
Every internal message kernel obeys
\[
 \rho_v\leq
 \left(\frac{\chi_v}{a_v+\sum_j\gamma_{v,j}}\right)^2.       \tag{34.14}
\]
At a leaf the corresponding analytic stock is
\[
 \rho_v\leq\left(\frac{\chi_v}{a_v}\right)^2.                \tag{34.15}
\]
The radial convexity and slope hypotheses used here propagate through the
whole finite tree by Corollary \ref{cor:v18v3-tree-closure} and Proposition
\ref{prop:v18v4-stiffness}.
\end{corollary}

\begin{proof}
Apply Theorem \ref{thm:v18v24-stiffness} with (34.13).  At a leaf take
\(\sigma_v=0\).  The final statement is exactly the cited closure and
stiffness results.
\end{proof}

This is the first contraction stock in the current cycle which improves
automatically with the number and stiffness of child messages.  It also
removes a circularity: the V20 response decay assumed a radial contraction,
whereas (34.14) derives one from quantities already propagated by V3--V4.

\subsection{One exceptional terminal layer}

The leaf estimate (34.15) can be weaker than the internal estimate (34.14).
The global Bessel argument does not require spending the leaf coefficient
at every depth.

\begin{theorem}[Boundary-layer ancestry gate]
\label{thm:v18v24-boundary-layer}
Assume that every non-leaf vertex has at most \(b\) children and satisfies
\[
 \rho_v\leq\rho,\qquad 0<\rho<\frac1b,                       \tag{34.16}
\]
while every leaf satisfies \(\rho_v\leq\rho_\partial<\infty\).
Let
\[
 C_\partial=\max\left\{1,\frac{\rho_\partial}{\rho}\right\} \tag{34.17}
\]
and retain the coefficient stocks \(a^*,\lambda_*,p_*\) of V20.  If
\[
 \mathfrak G_\partial=
 \frac{C_\partial^2(a^*)^2}
 {3\lambda_*p_*(1-b\rho)(1-\rho)}<1,                           \tag{34.18}
\]
then every finite tree satisfies
\[
 H_{\mathcal T}(Q)
 \succeq(1-\mathfrak G_\partial)p_*I                          \tag{34.19}
\]
uniformly in its depth and exterior boundary.
\end{theorem}

\begin{proof}
Let \(e\) be carried by a vertex \(u\).  Along a path of \(h\geq1\)
transports, Lemma \ref{lem:v18v20-step} gives a factor at most \(\rho^h\)
when \(u\) is internal, and at most
\(\rho_\partial\rho^{h-1}\leq C_\partial\rho^h\) when \(u\) is a
leaf.  At \(h=0\), the direct component is bounded by
\(a^*\leq C_\partial a^*\).  Thus (30.14) is replaced by
\[
 |(\widetilde c_v)_e|
 \leq C_\partial a^*\rho^{h(v,e)}.                            \tag{34.20}
\]
Repeat the row-sum proof of Theorem \ref{thm:v18v20-gate}.  The two response
factors in each rank-one entry multiply the former constant by
\(C_\partial^2\), producing (34.18)--(34.19).
\end{proof}

More generally, any fixed number of exceptional boundary layers can be
absorbed into the amplitude constant by multiplying their response stocks.
Only the coefficient repeated through unboundedly many internal levels must
be below the inverse branching rate.

\subsection{Canonical edge relation and the remaining separator}

For the conformal edge parametrization used in V3--V4,
\[
 a_v=\chi_ve^{-d_v},
 \qquad b_v=\chi_ve^{d_v}.                                    \tag{34.21}
\]
The leaf stock (34.15) is therefore
\[
 \rho_v^{\rm leaf}\leq e^{2d_v}.                              \tag{34.22}
\]
Small \(\chi_v\) alone does not improve this analytic ratio because it
scales \(a_v\) at the same time.  At an internal vertex, however,
\[
 \rho_v^{\rm int}\leq
 \left(
  \frac{\chi_v}{\chi_ve^{-d_v}+\sum_j\gamma_{v,j}}
 \right)^2,                                                    \tag{34.23}
\]
so inherited stiffness is the mechanism capable of producing strict
ancestral contraction.  Equation (34.23), not a bare power of the lattice
spacing, is the next physical matching quantity.

\section{V24 consequence and next gate}
\label{sec:v18v24-next}

V24 converts the local contraction from a numerical maximum into the
analytic stock (34.14).  Together with the boundary-layer theorem, this
reduces recursive invariance to lower bounds on the child stiffnesses
\(\gamma_{v,j}\).  Those quantities already have the exact V4 formula
(11.2); the next step is to propagate a uniform stiffness floor through
the allowed coefficient sector and insert it into (34.18).

The bound may still be too conservative for the physical sector.  In that
case V22 supplies the exact positive-series coefficient and V19 supplies the
exact normalized loss matrix, so failure of (34.18) would reject only this
analytic sufficient gate.

\begin{assessment}[V24 acquisition and boundary]
\label{ass:v18v24}
V24 proves the global stiffness-improved contraction estimate, derives the
internal response stock from child-message slope floors, and proves that a
single weaker leaf layer can be absorbed without spoiling cofinal depth.
It also shows why canonical weak hopping alone does not create contraction.
A uniform physical lower stock for the inherited stiffness and the strict
gate (34.18) remain unproved.  Cycles, cutoff removal, reconstruction, and
the coupled Einstein continuum remain open.
\end{assessment}

\begin{record}[V24 append record]
\label{rec:v18v24}
V24 is appended to the validated V23 whose installed SHA--256 was
\texttt{0CCA78C3B9D9C8F30CC535CEDA248064C596B7F9B47FD5BA94EED4F00CD54541}.
After installation only V24 is the active numeric SM note; the frozen
\texttt{V100\_f17} archive is unchanged.
\end{record}

% =====================================================================
% V25 APPENDIX: CLOSED STIFFNESS FLOOR AND COEFFICIENT-ONLY TREE GATE
% =====================================================================

\section{V25: a closed stiffness floor from coefficient endpoints}
\label{sec:v18v25}

V24 reduced internal contraction to the child slope floors
\(\gamma_{v,j}\).  This note closes that recursion.  The onsite quartic
term supplies a positive contribution to the exact V4 stiffness formula at
every vertex, while the child slope ceilings prevent the radial measure from
moving all its mass to zero.  The resulting lower constant depends only on
the endpoints of a compact coefficient sector and not on depth or on the
particular descendant messages.

\subsection{Coefficient sector}

Let \(\mathfrak A\Subset(0,\infty)\) be the finite set of allowed
half-dimensions \(\alpha\).  Assume
\[
\begin{gathered}
 0<a_*\leq a_v\leq a^*,\qquad
 0<b_*\leq b_v\leq b^*,\qquad
 0<\lambda_*\leq\lambda_v\leq\lambda^*,\qquad
 0<\chi_v\leq\chi^*,                                           \tag{35.1}\
 r_v\leq\mathfrak b,
 \qquad
 \sum_{j=1}^{r_v}b_{v,j}\leq B^*.                             \tag{35.2}
\end{gathered}
\]
Here \(\mathfrak b\) is the branching bound and the symbols \(b_v\) and
\(b_{v,j}\) are edge-action coefficients.  Define the pure radial moment
\[
 \mathfrak m_s(c,\lambda)=
 \frac{\int_0^\infty t^s e^{-\lambda t^2-ct}\,dt}
      {\int_0^\infty t^{s-1}e^{-\lambda t^2-ct}\,dt}           \tag{35.3}
\]
and the endpoint stocks
\[
 M_*=min_{\alpha\in\mathfrak A}
 \mathfrak m_{\alpha+1}(a^*+B^*,\lambda^*),
 \qquad G_*=2\lambda_*M_*,                                    \tag{35.4}
\]
\[
 \gamma_*=\frac{b_*G_*}{a^*+G_*}>0.                           \tag{35.5}
\]

\begin{lemma}[Shape-\(\alpha+1\) moment floor]
\label{lem:v18v25-shape-floor}
Let
\[
 \Psi_v(t)=\lambda_vt^2+sum_{j=1}^{r_v}\ell_{v,j}(t).        \tag{35.6}
\]
Under the radial closure hypotheses of V3--V4,
\[
 \mathbb E_{\alpha_v+1}^{\Psi_v}t\geq M_* .                  \tag{35.7}
\]
\end{lemma}

\begin{proof}
The V4 slope ceiling gives
\[
 0\leq\ell_{v,j}(t)-\ell_{v,j}(0)\leq b_{v,j}t.               \tag{35.8}
\]
After dropping the irrelevant constants, the likelihood ratio of the
actual shape-\(\alpha_v+1\) radial law to the reference law with potential
\[
 \lambda_vt^2+(a_v+B_v)t,
 \qquad B_v=\sum_jb_{v,j},                                    \tag{35.9}
\]
is nondecreasing in \(t\).  Indeed, the derivative of its logarithm is
\(B_v-\sum_j\ell_{v,j}'(t)\geq0\).  Hence the actual law stochastically
dominates the reference law, and its mean is at least
\(\mathfrak m_{\alpha_v+1}(a_v+B_v,\lambda_v)\).  This moment decreases
when either coefficient increases.  Equations (35.1)--(35.4) prove
(35.7).
\end{proof}

\begin{theorem}[Uniform outgoing stiffness floor]
\label{thm:v18v25-stiffness}
Every outgoing child message in every finite tree in the coefficient sector
(35.1)--(35.2) satisfies
\[
 \ell_v'(Q)\geq\gamma_*
 \qquad(Q\geq0).                                               \tag{35.10}
\]
The constant is independent of the depth and of the descendant messages.
\end{theorem}

\begin{proof}
In the notation of Proposition \ref{prop:v18v4-stiffness},
\[
 G_{\alpha_v+1}^{\Psi_v}
 =\mathbb E_{\alpha_v+1}^{\Psi_v}\Psi_v'(t).                  \tag{35.11}
\]
Convexity and monotonicity of every child message give
\[
 \Psi_v'(t)=2\lambda_vt+sum_j\ell_{v,j}'(t)
 \geq2\lambda_vt.                                             \tag{35.12}
\]
Lemma \ref{lem:v18v25-shape-floor} therefore yields
\[
 G_{\alpha_v+1}^{\Psi_v}\geq2\lambda_*M_*=G_* .              \tag{35.13}
\]
The exact V4 formula (11.2) and the coefficient endpoints give
\[
 \ell_v'(Q)
 \geq b_v\frac{G_{\alpha_v+1}^{\Psi_v}}
                 {a_v+G_{\alpha_v+1}^{\Psi_v}}
 \geq b_*\frac{G_*}{a^*+G_*}=\gamma_* .                      \tag{35.14}
\]
This proves (35.10) at every vertex, including leaves.
\end{proof}

\subsection{Coefficient-only response stocks}

Every internal vertex has at least one child.  Equations (34.13) and
(35.10) give
\[
 \sigma_v\geq r_v\gamma_*\geq\gamma_* .                      \tag{35.15}
\]
Define
\[
 \overline\rho=
 \left(\frac{\chi^*}{a_*+\gamma_*}\right)^2,
 \qquad
 \rho_\partial=left(\frac{\chi^*}{a_*}\right)^2,            \tag{35.16}
\]
\[
 C_\partial=max\left\{1,
                    \frac{\rho_\partial}{\overline\rho}
              \right\}.                                      \tag{35.17}
\]
The V18 fresh moment floor has the coefficient-only lower stock
\[
 p_0=a_*min_{\alpha\in\mathfrak A}
 \mathfrak m_\alpha(a^*+B^*,\lambda^*)>0.                    \tag{35.18}
\]

\begin{corollary}[Closed response hierarchy]
\label{cor:v18v25-response}
Every internal vertex satisfies \(\rho_v\leq\overline\rho\), every leaf
satisfies \(\rho_v\leq\rho_\partial\), and every fresh edge ledger entry
is at least \(p_0\).
\end{corollary}

\begin{proof}
Use (35.15) in Corollary \ref{cor:v18v24-internal}; the leaf statement is
(34.15).  The final statement is Proposition
\ref{prop:v18v18-moment-floor} with the endpoint monotonicity used in the
proof of Lemma \ref{lem:v18v25-shape-floor}.
\end{proof}

\begin{theorem}[Coefficient-only cofinal tree gate]
\label{thm:v18v25-tree-gate}
Suppose
\[
 0<\overline\rho<\frac1{\mathfrak b}                          \tag{35.19}
\]
and
\[
 \mathfrak Q_0:=
 \frac{C_\partial^2(a^*)^2}
 {3\lambda_*p_0
  (1-\mathfrak b\overline\rho)(1-\overline\rho)}<1.          \tag{35.20}
\]
Then, for every finite rooted tree in the sector and every exterior boundary
value,
\[
 H_{\mathcal T}(Q)\succeq(1-\mathfrak Q_0)p_0I.               \tag{35.21}
\]
The lower margin is independent of the depth and the number of edges.
\end{theorem}

\begin{proof}
Corollary \ref{cor:v18v25-response} verifies the internal, terminal, and
fresh-floor hypotheses of Theorem \ref{thm:v18v24-boundary-layer} with
\(\rho=\overline\rho\).  Equation (35.20) is exactly its strict amplitude
gate.  Equation (34.19) then gives (35.21).
\end{proof}

Theorem \ref{thm:v18v25-tree-gate} is a finite list of explicit inequalities
in coefficient-box endpoints and one-dimensional pure radial integrals.
It does not require recursive numerical message evaluation.  If its coarse
stocks fail, the exact V22 contraction and V19 loss matrix remain valid
sharper alternatives.

\subsection{Canonical parameter box}

If the edge coefficients have the canonical form (34.21) and
\[
 \chi_-\leq\chi_v\leq\chi_+,
 \qquad d_-\leq d_v\leq d_+,                                  \tag{35.22}
\]
then admissible endpoint choices are
\[
\begin{split}
 a_*&=\chi_-e^{-d_+},& a^*&=\chi_+e^{-d_-},\\
 b_*&=\chi_-e^{d_-},& b^*&=\chi_+e^{d_+},\\
 \chi^*&=\chi_+ .                                             \tag{35.23}
\end{split}
\]
Thus (35.19)--(35.20) can be tested directly on a proposed physical box.
The present source manuscripts do not supply such a continuum box with
numerical endpoints, so no physical pass is asserted here.

\section{V25 consequence and mandatory audit point}
\label{sec:v18v25-next}

The tree-sector conformal Hessian now has a complete coefficient-only
sufficient theorem.  The theorem is conditional only on the explicit
strict inequalities (35.19)--(35.20) and the already proved radial closure
hypotheses.  This resolves the recursive stiffness circularity identified
in V24.

V25 is also the next mandatory companion-review point.  Before V26, the
newest active YM and NS notes must be inspected.  The subsequent direction
will depend on whether they contain a sharper frame estimate applicable to
the conservative amplitude stock \(\mathfrak Q_0\), or whether the programme
should proceed to physical coefficient matching and then to cycles.

\begin{assessment}[V25 acquisition and boundary]
\label{ass:v18v25}
V25 proves a depth-independent child-message stiffness floor and the fully
explicit coefficient-only cofinal tree gate (35.19)--(35.21).  The strict
inequalities have not been verified on a physical continuum coefficient
box, because the source manuscripts do not provide one.  The theorem still
concerns trees; cycle restoration, cutoff removal, reconstruction, and the
coupled Einstein continuum remain open.
\end{assessment}

\begin{record}[V25 append record]
\label{rec:v18v25}
V25 is appended to the validated V24 whose installed SHA--256 was
\texttt{728A6C18F40B045131C34EC0D1A73FCC47F9D79A9E6217CFC9EA087E4C3ACEAB}.
After installation only V25 is the active numeric SM note; the frozen
\texttt{V100\_f17} archive is unchanged.
\end{record}

% =====================================================================
% V26 APPENDIX: COMPANION AUDIT AND CYCLE-GLUING RESPONSE
% =====================================================================

\section{V26: companion audit and a cycle-gluing gate}
\label{sec:v18v26}

The mandatory V25 companion audit was completed before this continuation.
The newest active YM note was
\[
 \texttt{Renewal\_Yang\_Mills\_Mass\_Gap\_Research\_Notes\_V3.tex}
\]
with SHA--256
\[
 \texttt{ADC461AF3B3D8D1135F962F4DCAC52B187F5CE53FD3D5F45D87965D09D5CDADE}.
                                                               \tag{36.1}
\]
It now gives a residual-weighted cell-response theorem and an exact
sufficient average inequality for promoting a finite grid sign to a Haar
mean, while proving that a single global derivative majorant is too coarse
for its cells.  The active NS note remained
\[
 \texttt{Renewal\_NS\_Stability\_Research\_Notes\_V26.tex}
\]
with SHA--256
\[
 \texttt{82C887F8C2C69E4F28FAD7C3DC8DE5B9099CB5A4BF431EBA1FFF3E91935978AD}.
                                                               \tag{36.2}
\]
Its exact rank-one norms sometimes improve a full operator norm and
sometimes coincide with it.  No theorem is transferred between programmes.
The applicable design decision is to retain localized signed response
matrices during cycle restoration and postpone absolute operator bounds
until their spatial overlap is summed.

\subsection{Cutting the chord cross terms}

Let \(G=(V,E)\) be a finite connected regulator graph, choose a spanning
tree \(T\), and write \(\mathcal C=E\setminus T\) for the chords.  A chord
Higgs edge has the form
\[
 \chi_c\left|
 e^{-d_c/2}z_{x(c)}-e^{d_c/2}U_cz_{y(c)}
 \right|^2
 =a_c|z_{x(c)}|^2+b_c|z_{y(c)}|^2+W_c,                       \tag{36.3}
\]
where
\[
 a_c=\chi_ce^{-d_c},\qquad b_c=\chi_ce^{d_c},\qquad
 W_c=-2\chi_c z_{x(c)}\mathbin{\cdot}U_cz_{y(c)}.             \tag{36.4}
\]
The cross term \(W_c\) is independent of every conformal difference
parameter.  Define the cut action \(S_0\) by retaining the two endpoint
quadratic terms of every chord but omitting \(W_c\), and interpolate
\[
 S_s=S_0+s\sum_{c\in\mathcal C}W_c,
 \qquad 0\leq s\leq1.                                         \tag{36.5}
\]
Thus \(S_1\) is the full cyclic action.  The cut interaction graph is a
tree; the retained chord endpoint terms are radial onsite terms and can be
included in the coefficient sector of V25.

Let \(\vartheta=(d_e)_{e\in E}\) denote all conformal edge parameters and
put
\[
 F_s(\vartheta)=-\log\int e^{-S_s(\vartheta,X)}\,dX,
 \qquad H_s=\nabla_\vartheta^2F_s.                             \tag{36.6}
\]
For indices \(e,f\in E\), write
\[
 A_e=\partial_{d_e}S_0,
 \qquad A_{ef}=\partial_{d_e}\partial_{d_f}S_0.               \tag{36.7}
\]
These scores do not depend on \(s\).

\subsection{Exact Hessian response to gluing}

For random variables under the \(S_s\) law, define
\[
 \operatorname{Cum}_s(X,Y,Z)
 =\mathbb E_s[(X-\mathbb E_sX)(Y-\mathbb E_sY)
                         (Z-\mathbb E_sZ)].                    \tag{36.8}
\]

\begin{theorem}[Exact chord-insertion response]
\label{thm:v18v26-response}
For every \(e,f\in E\),
\[
 (H_s)_{ef}=\mathbb E_sA_{ef}
             -\operatorname{Cov}_s(A_e,A_f),                  \tag{36.9}
\]
and
\[
 \frac{d}{ds}(H_s)_{ef}
 =\sum_{c\in\mathcal C}(K_c(s))_{ef},                         \tag{36.10}
\]
where
\[
 (K_c(s))_{ef}
 =-\operatorname{Cov}_s(A_{ef},W_c)
  +\operatorname{Cum}_s(A_e,A_f,W_c).                         \tag{36.11}
\]
Consequently,
\[
 H_1=H_0+int_0^1\sum_{c\in\mathcal C}K_c(s)\,ds.            \tag{36.12}
\]
\end{theorem}

\begin{proof}
Twice differentiating the log partition function gives (36.9).  Since
\(W_c\) is independent of \(\vartheta\), the variables in (36.7) have no
explicit \(s\)-derivative.  For every integrable observable \(Y\),
\[
 \frac{d}{ds}\mathbb E_sY
 =-\sum_c\operatorname{Cov}_s(Y,W_c).                          \tag{36.13}
\]
Apply this first to \(A_{ef}\).  Direct expansion of the centered products
also gives
\[
 \frac{d}{ds}\operatorname{Cov}_s(A_e,A_f)
 =-\sum_c\operatorname{Cum}_s(A_e,A_f,W_c).                   \tag{36.14}
\]
Substitution into (36.9) proves (36.10)--(36.11), and integration proves
(36.12).
\end{proof}

This formula keeps the two contributions with their signs.  Bounding the
covariance and cumulant separately is permitted but can lose the same kind
of cancellation that the YM V3 audit found important.

\subsection{A volume-independent locality criterion}

Equip the edge set with a graph distance and let \(d(e,c)\) be the distance
from edge \(e\) to either endpoint edge of chord \(c\).  For
\(0<r<1\), define the geometric overlap stocks
\[
 M_{E\mathcal C}(r)=sup_{e\in E}
       \sum_{c\in\mathcal C}r^{d(e,c)},
 \qquad
 M_{\mathcal CE}(r)=sup_{c\in\mathcal C}
       \sum_{f\in E}r^{d(f,c)}.                               \tag{36.15}
\]
On a bounded-degree graph family with uniformly bounded chord density,
these stocks are finite uniformly in the volume whenever \(r\) is smaller
than the reciprocal shell-growth rate.

\begin{theorem}[Localized cycle-restoration gate]
\label{thm:v18v26-cycle-gate}
Assume the cut reference satisfies
\[
 H_0\succeq\mu_0I                                             \tag{36.16}
\]
with \(\mu_0>0\) independent of the regulator.  Suppose there are
\(\delta\geq0\) and \(0<r<1\), also regulator-independent, such that
\[
 |(K_c(s))_{ef}|
 \leq\delta r^{d(e,c)+d(f,c)}                                 \tag{36.17}
\]
for all \(c,e,f\) and \(0\leq s\leq1\).  Then
\[
 H_1\succeq
 \left[
  \mu_0-\delta M_{E\mathcal C}(r)M_{\mathcal CE}(r)
 \right]I.                                                     \tag{36.18}
\]
In particular, the full cyclic Hessian has a uniform positive margin if
\[
 \delta M_{E\mathcal C}(r)M_{\mathcal CE}(r)<\mu_0.           \tag{36.19}
\]
\end{theorem}

\begin{proof}
For a fixed row \(e\), (36.17) gives
\[
\begin{split}
 \sum_f\left|\sum_c(K_c(s))_{ef}\right|
 &\leq\delta\sum_c r^{d(e,c)}\sum_f r^{d(f,c)}\\
 &\leq\delta M_{E\mathcal C}(r)M_{\mathcal CE}(r).           \tag{36.20}
\end{split}
\]
The summed response matrix is symmetric, so its operator norm is bounded by
the maximum absolute row sum.  Equation (36.12) and (36.16) therefore give
(36.18).  Strict inequality (36.19) makes the bracket positive.
\end{proof}

\begin{corollary}[Finite-range chord response]
\label{cor:v18v26-finite-range}
If every \(K_c(s)\) is supported on edge parameters at distance at most
\(R\) from \(c\), has operator norm at most \(\delta_0\), and each edge
belongs to at most \(\Omega_R\) such supports, then
\[
 H_1\succeq(\mu_0-\Omega_R\delta_0)I.                          \tag{36.21}
\]
\end{corollary}

\begin{proof}
Coloring is unnecessary.  Extend each local response matrix by zero to the
full parameter space.  Its negative part is bounded below by
\(-\delta_0P_c\), where \(P_c\) is the projection onto its support.  The
overlap assumption gives \(\sum_cP_c\preceq\Omega_RI\).  Integrate (36.12).
\end{proof}

\subsection{What must be proved for the actual chord response}

The response matrices (36.11) contain one local insertion \(W_c\) and one
or two conformal scores.  The tree Markov contraction of V20 and the closed
coefficient stock of V25 make (36.17) plausible: influence must travel from
\(c\) to \(e\) and \(f\).  It is not enough, however, to cite two-point
decay.  The third cumulant in (36.11) requires a three-observable response
estimate uniform along the entire interpolation (36.5).

The next rigorous target is therefore an insertion version of the V20
Lipschitz contraction: a conditional covariance with a centered local
observable should lose one factor of the radial response coefficient each
time it crosses a tree edge.  Proving that estimate would yield (36.17)
with \(r\) determined by the internal stock \(\overline\rho\) of V25.

\begin{assessment}[V26 acquisition and boundary]
\label{ass:v18v26}
V26 records the mandatory companion audit, gives the exact conformal
Hessian response to gluing all chord cross terms, and proves a
volume-independent cycle-restoration criterion from localized response
matrices.  The required covariance/cumulant decay (36.17) has not yet been
proved, and the cut reference margin must satisfy the V25 coefficient gate.
Cutoff removal, reconstruction, and the coupled Einstein continuum remain
open.
\end{assessment}

\begin{record}[V26 append record]
\label{rec:v18v26}
V26 is appended to the validated V25 whose installed SHA--256 was
\texttt{ADFB67F7A55787F3ED5AD8B6AF1B0942EA843E32EC1481FA8A0875F62086C84F}.
After installation only V26 is the active numeric SM note; the frozen
\texttt{V100\_f17} archive is unchanged.
\end{record}

% =====================================================================
% V27 APPENDIX: FULL-FIELD CONTRACTION AND DOBRUSHIN PATH SUM
% =====================================================================

\section{V27: full-field contraction beyond the radial channel}
\label{sec:v18v27}

The V20 contraction acts on functions of the squared radius.  A chord
insertion \(W_c=-2\chi_c z_x\cdot U_cz_y\) also contains angular field
directions, so radial contraction alone cannot prove (36.17).  This note
derives the needed full-field estimate directly from strong convexity.  Its
one-step coefficient is the square root of the analytic radial coefficient
in V24.  The associated Dobrushin path sum is valid on graphs with cycles
and supplies a volume-independent two-observable influence bound.

\subsection{One-step Sobolev contraction}

Let
\[
 K_v(dz\mid y)=\frac1{Z_v(y)}
 \exp\{-U_v(z)+2\chi_v z\cdot y\}\,dz                          \tag{37.1}
\]
and suppose
\[
 \nabla_z^2U_v(z)\succeq2A_vI,
 \qquad A_v=a_v+\sigma_v>0.                                  \tag{37.2}
\]
Write \((P_vF)(y)=\int F(z)K_v(dz\mid y)\) and
\[
 \tau_v=\frac{\chi_v}{A_v}.                                  \tag{37.3}
\]

\begin{theorem}[Full-field Markov gradient contraction]
\label{thm:v18v27-gradient}
For every locally Lipschitz \(F\) with square-integrable gradient,
\[
 |\nabla_yP_vF(y)|^2
 \leq\tau_v^2(P_v|\nabla F|^2)(y).                            \tag{37.4}
\]
In particular,
\[
 \operatorname{Lip}(P_vF)\leq\tau_v\operatorname{Lip}(F).    \tag{37.5}
\]
\end{theorem}

\begin{proof}
For a vector \(h\), differentiation under the integral gives
\[
 D_hP_vF(y)=2\chi_v\operatorname{Cov}_{v,y}(F,z\cdot h).       \tag{37.6}
\]
Brascamp--Lieb, ordinary covariance Cauchy--Schwarz, and (37.2) give
\[
\begin{split}
 |\operatorname{Cov}(F,z\cdot h)|^2
 &\leq
 \mathbb E[\nabla F^{\mathsf T}(\nabla^2U_v)^{-1}\nabla F]
 \mathbb E[h^{\mathsf T}(\nabla^2U_v)^{-1}h]\\
 &\leq\frac{|h|^2}{4A_v^2}\mathbb E|\nabla F|^2.              \tag{37.7}
\end{split}
\]
Insert (37.7) in (37.6) and take the supremum over unit \(h\).  This proves
(37.4); bounded gradients give (37.5).
\end{proof}

For the recursive Higgs kernel, (34.13) and V25 give
\[
 A_v\geq a_v+\sum_j\gamma_{v,j}
 \geq a_*+\gamma_*                                             \tag{37.8}
\]
at every internal vertex.  Thus
\[
 \tau_v\leq\overline\tau:=
 \frac{\chi^*}{a_*+\gamma_*},
 \qquad \overline\rho=\overline\tau^2,                       \tag{37.9}
\]
which agrees with (35.16).

\begin{corollary}[Composed field response]
\label{cor:v18v27-composition}
For a chain of conditional kernels \(P_0\cdots P_{h-1}\),
\[
 |\nabla(P_0\cdots P_{h-1}F)|^2
 \leq\left(\prod_{j=0}^{h-1}\tau_j^2\right)
 P_0\cdots P_{h-1}|\nabla F|^2.                               \tag{37.10}
\]
\end{corollary}

\begin{proof}
Apply (37.4) successively and use positivity of every Markov operator.
\end{proof}

\subsection{Two-branch covariance decay on the cut tree}

Let \(u\) and \(v\) lie in distinct branches below their lowest common
ancestor \(o\).  Let \(F(z_u)\) and \(G(z_v)\) have square-integrable
gradients.  Conditional independence below \(o\) gives
\[
 \operatorname{Cov}(F(z_u),G(z_v))
 =\operatorname{Cov}(f_o(z_o),g_o(z_o)),                       \tag{37.11}
\]
where \(f_o\) and \(g_o\) are the conditional expectations transported to
\(o\).

\begin{proposition}[Two-branch Sobolev covariance]
\label{prop:v18v27-two-branch}
If every vertex on the two paths has strong-convexity constant at least
\(A_*>0\), then
\[
 |\operatorname{Cov}(F(z_u),G(z_v))|
 \leq\frac1{2A_*}
 \left(\prod_{e\in[o,u)}\tau_e\right)
 \left(\prod_{e\in[o,v)}\tau_e\right)
 \bigl(\mathbb E|\nabla F|^2\bigr)^{1/2}
 \bigl(\mathbb E|\nabla G|^2\bigr)^{1/2}.                    \tag{37.12}
\]
\end{proposition}

\begin{proof}
Apply Brascamp--Lieb at \(o\) to the right side of (37.11):
\[
 |\operatorname{Cov}(f_o,g_o)|
 \leq\frac1{2A_*}
 (\mathbb E|\nabla f_o|^2)^{1/2}
 (\mathbb E|\nabla g_o|^2)^{1/2}.                             \tag{37.13}
\]
Corollary \ref{cor:v18v27-composition} on each branch, followed by the
tower property, proves (37.12).
\end{proof}

This estimate applies to polynomial observables such as chord insertions:
their gradients are unbounded but have finite moments by quartic
confinement.  It therefore avoids replacing \(W_c\) by a global supremum.

\subsection{Dobrushin matrix on a cyclic graph}

For a general finite interaction graph, condition at one vertex while
holding all neighboring fields fixed.  Suppose changing the field at a
neighbor \(w\) changes the linear source at \(v\) through a coefficient
\(\chi_{vw}\) and an isometry.  Define
\[
 C_{vw}=\begin{cases}
 \chi_{vw}/A_v,& v\sim w,\\
 0,&v\not\sim w.
 \end{cases}                                                   \tag{37.14}
\]

\begin{lemma}[One-site Wasserstein influence]
\label{lem:v18v27-wasserstein}
For two neighbor configurations \(x,x'\),
\[
 W_1\bigl(K_v(\cdot\mid x),K_v(\cdot\mid x')\bigr)
 \leq\sum_w C_{vw}|x_w-x_w'|.                                \tag{37.15}
\]
\end{lemma}

\begin{proof}
Interpolate linearly between the two source fields.  By (37.5), every
one-Lipschitz test function changes by at most the right side of (37.15).
Kantorovich duality gives the Wasserstein bound.
\end{proof}

\begin{theorem}[Dobrushin path-sum influence]
\label{thm:v18v27-dobrushin}
Assume
\[
 \kappa:=\|C\|_{\infty\to\infty}
 =\max_v\sum_wC_{vw}<1.                                      \tag{37.16}
\]
Then the influence of a unit boundary displacement at \(w\) on the
expectation of a one-Lipschitz observable at \(v\) is at most
\[
 D_{vw}:=\sum_{n=0}^\infty(C^n)_{vw}.                          \tag{37.17}
\]
If interactions are nearest-neighbor and graph degree is at most
\(\Delta\), with \(C_{vw}\leq c_*\), then
\[
 D_{vw}\leq
 \frac{(\Delta c_*)^{d(v,w)}}{1-\Delta c_*}.                  \tag{37.18}
\]
The constants are independent of the volume and the presence of cycles.
\end{theorem}

\begin{proof}
Couple two single-site conditional laws optimally and update the sites in a
common order.  If \(d_v^{(n)}\) is the expected field discrepancy after
the \(n\)-th influence propagation, Lemma
\ref{lem:v18v27-wasserstein} gives the componentwise recursion
\[
 d^{(n+1)}\leq Cd^{(n)}+b,                                    \tag{37.19}
\]
where \(b\) is the boundary discrepancy vector.  Iteration and
\(\|C\|_\infty<1\) give
\[
 d\leq\sum_{n=0}^\infty C^nb.                                \tag{37.20}
\]
Kantorovich duality converts this coupling bound to observable influence,
proving (37.17).  A nonzero contribution to \((C^n)_{vw}\) requires a walk
of length \(n\geq d(v,w)\).  There are at most \(\Delta^n\) such walks and
each has weight at most \(c_*^n\).  Summing the geometric tail proves
(37.18).
\end{proof}

\subsection{Interface with the cycle response}

At the cut point \(s=0\), Proposition
\ref{prop:v18v27-two-branch} supplies the required distance factors for the
covariance term in (36.11) whenever the conformal scores are assigned to
their local edge supports.  Theorem \ref{thm:v18v27-dobrushin} supplies the
corresponding cycle-compatible influence mechanism if the one-site row sum
(37.16) holds uniformly along the interpolation.

Two gaps remain before (36.17) follows.  First, a three-observable cumulant
bound must preserve two distance factors rather than collapsing to a global
third moment.  Second, the local strong-convexity constants \(A_v\) used in
(37.14) must be certified uniformly for every \(s\in[0,1]\); chord cross
terms change the joint Hessian even though they do not change the one-site
onsite Hessian.  These are now explicit local estimates rather than a
volume-dependent graph problem.

\begin{assessment}[V27 acquisition and boundary]
\label{ass:v18v27}
V27 proves full-field Sobolev and Lipschitz contraction, its composition
law, a two-branch covariance estimate, and a Dobrushin path-sum influence
bound valid on cyclic graphs.  This controls the two-point part of the V26
chord response under a strict row-sum gate.  The three-point cumulant decay
and a physical verification of the row sum remain unproved.  Cutoff
removal, reconstruction, and the coupled Einstein continuum remain open.
\end{assessment}

\begin{record}[V27 append record]
\label{rec:v18v27}
V27 is appended to the validated V26 whose installed SHA--256 was
\texttt{144E8E5D0B73302F38AC2BC8F02F5120AF6D7A7F24B13FBC4B1B98D7A67C463F}.
After installation only V27 is the active numeric SM note; the frozen
\texttt{V100\_f17} archive is unchanged.
\end{record}

% =====================================================================
% V28 APPENDIX: THREE-BRANCH CUMULANT DECAY
% =====================================================================

\section{V28: three-branch cumulant decay}
\label{sec:v18v28}

The second term of the exact chord response (36.11) is a third cumulant.
Replacing it by a global third moment would destroy spatial decay.  On the
cut tree, conditional independence at the Steiner branch point reduces the
cumulant to three conditional expectations.  V27 contracts the gradient of
each conditional expectation along its own arm.  A strong-convexity
\(L^3\) Poincar\'e inequality at the branch point then preserves all three
distance factors.

\subsection{An \(L^p\) consequence of strong convexity}

\begin{lemma}[Strong-convexity \(L^p\) Poincar\'e bound]
\label{lem:v18v28-lp}
Let \(\nu(dz)\propto e^{-U(z)}dz\) with
\(\nabla^2U\succeq2AI\), \(A>0\).  For every locally Lipschitz \(f\) and
every \(p\geq2\),
\[
 \|f-\mathbb E_\nu f\|_{L^p(\nu)}
 \leq\frac{p}{\sqrt{2A}}\,
       \||\nabla f|\|_{L^p(\nu)}.                              \tag{38.1}
\]
In particular, for three such functions,
\[
 |\operatorname{Cum}_\nu(f_1,f_2,f_3)|
 \leq\frac{27}{(2A)^{3/2}}
       \prod_{i=1}^3\||\nabla f_i|\|_{L^3(\nu)}.            \tag{38.2}
\]
\end{lemma}

\begin{proof}
The Bakry--\'Emery criterion at curvature \(2A\) gives the log-Sobolev
inequality
\[
 \operatorname{Ent}_\nu(g^2)
 \leq\frac1A\int|\nabla g|^2\,d\nu.                            \tag{38.3}
\]
Its standard \(L^p\) Poincar\'e consequence, obtained by applying (38.3) to
smooth truncations of \(|f-\mathbb Ef|^{p/2}\) and integrating the resulting
moment differential inequality from \(2\) to \(p\), is
\[
 \|f-\mathbb Ef\|_p
 \leq\frac{p}{\sqrt{2A}}\||\nabla f|\|_p.
\]
Approximation removes the smoothness and boundedness truncations.  Apply
H\"older with exponents
\((3,3,3)\) to the centered-product definition (36.8), then use (38.1) with
\(p=3\) for each factor.  This gives (38.2).
\end{proof}

The constant in (38.1) is intentionally nonoptimal.  Its role is to retain
the three spatial factors; V22-style interval evaluation can later sharpen
the local moment constant.

\subsection{Gradient transport in \(L^3\)}

\begin{lemma}[Composed \(L^3\) gradient contraction]
\label{lem:v18v28-gradient3}
For a chain of kernels satisfying (37.2),
\[
 \left\|\left|
 \nabla(P_0\cdots P_{h-1}F)
 \right|\right\|_3
 \leq\left(\prod_{j=0}^{h-1}\tau_j\right)
       \||\nabla F|\|_3.                                      \tag{38.4}
\]
All norms are taken in the corresponding consistent marginal laws.
\end{lemma}

\begin{proof}
Equation (37.4) gives
\[
 |\nabla P_jF|^3
 \leq\tau_j^3(P_j|\nabla F|^2)^{3/2}
 \leq\tau_j^3P_j|\nabla F|^3,                                \tag{38.5}
\]
where the second inequality is Jensen's inequality.  Integrate and use the
tower property, then iterate along the chain.
\end{proof}

\subsection{The Steiner-point cumulant theorem}

Let \(o\) be a vertex of a tree.  Removing \(o\) produces branches.  Let
\(F_i\) be an observable supported in a branch \(\mathcal B_i\), with the
three branches distinct.  Write \(d_i\) for the length of the path from
\(o\) to the support of \(F_i\), and define
\[
 \tau[o,F_i]=\prod_{e\in[o,\operatorname{supp}F_i)}\tau_e.     \tag{38.6}
\]

\begin{theorem}[Three-branch cumulant contraction]
\label{thm:v18v28-three-branch}
Suppose the marginal potential at \(o\) has Hessian at least \(2A_*I\).
Then
\[
 |\operatorname{Cum}(F_1,F_2,F_3)|
 \leq\frac{27}{(2A_*)^{3/2}}
 \prod_{i=1}^3
 \left[
  \tau[o,F_i]\,
  \||\nabla F_i|\|_3
 \right].                                                       \tag{38.7}
\]
If every arm coefficient is at most \(\overline\tau<1\), the right side
contains the decay factor
\[
 \overline\tau^{d_1+d_2+d_3}.                                 \tag{38.8}
\]
\end{theorem}

\begin{proof}
Let
\[
 f_i(z_o)=\mathbb E[F_i\mid z_o].                              \tag{38.9}
\]
Conditional on \(z_o\), the three branches are independent.  Write
\(F_i=f_i(z_o)+\xi_i\) with
\(\mathbb E[\xi_i\mid z_o]=0\).  Expanding the centered triple product,
every term containing at least one \(\xi_i\) has zero conditional
expectation.  Hence
\[
 \operatorname{Cum}(F_1,F_2,F_3)
 =\operatorname{Cum}(f_1(z_o),f_2(z_o),f_3(z_o)).              \tag{38.10}
\]
Apply (38.2) at \(o\).  Lemma \ref{lem:v18v28-gradient3} on the three arms
gives
\[
 \||\nabla f_i|\|_3
 \leq\tau[o,F_i]\||\nabla F_i|\|_3.                        \tag{38.11}
\]
Equations (38.10)--(38.11) prove (38.7), and (38.8) is immediate.
\end{proof}

The same proof allows each \(F_i\) to occupy a finite terminal block.  One
first conditions that block onto the unique attachment field; the first
factor in (38.11) is then its block-to-boundary Sobolev response stock.

\subsection{Application at the cut-cycle anchor}

The conformal edge scores \(A_e\) and direct scores \(A_{ef}\) in (36.7)
are local polynomial observables.  The chord insertion \(W_c\) is a local
bilinear observable on its two endpoint fields.  Quartic confinement makes
their gradient \(L^3\) stocks finite.  Define a uniform local stock
\[
 \mathcal M_3=sup
 \left\{
  \||\nabla A_e|\|_3,
  \||\nabla A_{ef}|\|_3,
  \||\nabla W_c|\|_3
 \right\},                                                     \tag{38.12}
\]
where the supremum ranges over the permitted local types and boundary
boxes.

\begin{corollary}[Separated-support chord response at \(s=0\)]
\label{cor:v18v28-chord-anchor}
At the cut-tree measure \(s=0\), suppose the supports of
\(A_e,A_f,W_c\) lie in three distinct branches from their Steiner vertex
\(o\).  Then
\[
 |\operatorname{Cum}_0(A_e,A_f,W_c)|
 \leq\frac{27\mathcal M_3^3}{(2A_*)^{3/2}}
 \overline\tau^{d(o,e)+d(o,f)+d(o,c)}.                         \tag{38.13}
\]
If \(A_{ef}\) and \(W_c\) lie in distinct branches, the two-branch estimate
of V27 gives
\[
 |\operatorname{Cov}_0(A_{ef},W_c)|
 \leq\frac{\mathcal M_2^2}{2A_*}
 \overline\tau^{d(o,ef)+d(o,c)},                               \tag{38.14}
\]
with the analogous local \(L^2\) gradient stock \(\mathcal M_2\).
Thus the exact response matrix (36.11) has the required product decay at
\(s=0\) whenever its three local supports separate at a tree branch point.
\end{corollary}

\begin{proof}
Apply Theorem \ref{thm:v18v28-three-branch} and Proposition
\ref{prop:v18v27-two-branch}, then use the definitions of the uniform local
moment stocks.
\end{proof}

\subsection{Collision patterns and interpolation}

The separated-support result does not cover configurations in which two
observables lie in the same branch or overlap.  Those collision patterns
are local: after following their common path to the first separation point,
one combines the colliding observables into a finite-block response stock.
On a bounded-degree finite-range graph there are only finitely many such
local patterns up to type, so they can be certified by the V12 finite-cover
method.  The long arms still acquire the factors in (38.11).

The remaining global issue is uniformity for \(s\in(0,1]\), where chord
insertions destroy tree conditional independence.  The Dobrushin path sum
of V27 is the appropriate replacement, but a third-order Dobrushin response
theorem must be proved rather than assumed.  Its expected structure is a
sum over three-arm connected trees, with one resolvent factor
\((I-C)^{-1}\) on each arm and a local collision stock at the branching
core.

\begin{assessment}[V28 acquisition and boundary]
\label{ass:v18v28}
V28 proves an \(L^3\) strong-convexity cumulant bound, its contraction under
composed field kernels, and the three-branch Steiner decay needed for the
separated-support part of the chord response at the cut-tree anchor.  Local
collision patterns and the third-order Dobrushin estimate uniform in the
gluing parameter remain unproved.  Physical coefficient matching, cutoff
removal, reconstruction, and the coupled Einstein continuum remain open.
\end{assessment}

\begin{record}[V28 append record]
\label{rec:v18v28}
V28 is appended to the validated V27 whose installed SHA--256 was
\texttt{887B640D5F937DCE5F064FF168ECA28047102FFD109E4BD2383C2230AA21E720}.
After installation only V28 is the active numeric SM note; the frozen
\texttt{V100\_f17} archive is unchanged.
\end{record}

% =====================================================================
% V29 APPENDIX: STEINER-TREE RESPONSE GEOMETRY
% =====================================================================

\section{V29: the correct geometry of a three-point chord response}
\label{sec:v18v29}

The product-distance hypothesis (36.17) is a valid sufficient condition,
but V28 shows that it is stronger than the decay naturally produced by a
third cumulant.  If two score observables form a nearby cluster far from a
chord insertion, their connected response crosses the long arm once, not
twice.  The invariant geometric quantity is the length of the minimal
Steiner subtree joining the three supports.  This note proves the resulting
cut-tree bound and replaces the next working target by a summable Steiner
stock.

\subsection{Steiner length for local supports}

For finite vertex or edge sets \(S_1,S_2,S_3\) in a tree \(T\), let
\[
 \mathsf L_T(S_1,S_2,S_3)                                    \tag{39.1}
\]
be the number of tree edges in the smallest connected subtree meeting all
three sets.  For two sets, use \(\mathsf L_T(S_1,S_2)=d_T(S_1,S_2)\).
When the sets are single vertices \(u_1,u_2,u_3\), their unique median
\(o\) lies on all three pairwise geodesics and
\[
 \mathsf L_T(u_1,u_2,u_3)
 =d(o,u_1)+d(o,u_2)+d(o,u_3)
 =\frac12\sum_{i<j}d(u_i,u_j).                                \tag{39.2}
\]

Suppose every local observable used below has support diameter at most
\(R_0\).  Select one representative vertex in each support.  Replacing the
support Steiner length by the representative Steiner length changes it by
at most \(3R_0\).

\begin{theorem}[Steiner cumulant decay on the cut tree]
\label{thm:v18v29-steiner}
Let \(F_1,F_2,F_3\) be local observables with support diameter at most
\(R_0\).  Assume the tree has maximum degree \(\Delta\), only finitely
many local interaction types occur, and all one-step field contractions on the connecting subtree
are at most \(0<\overline\tau<1\), all branch-point strong-convexity
constants are at least \(A_*>0\), and
\[
 \||\nabla F_i|\|_3\leq M_{3,i}.                              \tag{39.3}
\]
Then there is a finite local collision constant
\(C_{\rm loc}(R_0,A_*,\Delta)\), independent of the volume and mutual distances,
such that
\[
 |\operatorname{Cum}(F_1,F_2,F_3)|
 \leq C_{\rm loc}(R_0,A_*,\Delta)
       \left(\prod_{i=1}^3M_{3,i}\right)
       \overline\tau^{\mathsf L_T(S_1,S_2,S_3)}.              \tag{39.4}
\]
For three disjoint branches with point supports, one may take
\[
 C_{\rm loc}(0,A_*,\Delta)=\frac{27}{(2A_*)^{3/2}}.                  \tag{39.5}
\]
\end{theorem}

\begin{proof}
First take point supports.  Let \(o\) be their median.  If none of the
observables is supported at \(o\), removing \(o\) places them in three
distinct branches.  Theorem \ref{thm:v18v28-three-branch} and (39.2) give
(39.4) with (39.5).

If one observable is supported at \(o\), condition the other two branches
onto \(z_o\).  The centered triple product is then bounded by H\"older at
\(o\), the \(L^3\) Poincar\'e estimate for the two transported factors, and
the local \(L^3\) stock of the observable at \(o\).  The two transported
gradients carry exactly
\(\overline\tau^{d(o,u_j)+d(o,u_k)}\), which is the Steiner length because
the third arm has length zero.  This only changes the local prefactor.

For supports of diameter at most \(R_0\), expose the first attachment point
of each support to the connecting subtree and condition the finite support
block onto that attachment field.  Collision of two or three blocks can
occur only within a subtree of diameter at most \(3R_0\).  Apply the local
\(L^3\) bound to that finite block and V28 to every remaining long arm.  At
worst, replacing attachment distances by support distances removes
\(3R_0\) contraction factors.  Absorb the finite multiplier
\(\overline\tau^{-3R_0}\) and the finitely many local block types into
\(C_{\rm loc}(R_0,A_*,\Delta)\).  The remaining exponent is (39.1).
\end{proof}

The same median argument with two observables and Proposition
\ref{prop:v18v27-two-branch} gives
\[
 |\operatorname{Cov}(F_1,F_2)|
 \leq C_{2,\rm loc}M_{2,1}M_{2,2}
       \overline\tau^{d_T(S_1,S_2)}.                           \tag{39.6}
\]

\subsection{The cut-anchor chord response}

Let \(S_e,S_f,S_c\) be the supports of \(A_e,A_f,W_c\).  The direct score
\(A_{ef}\) is zero unless the two conformal parameters occur in a common
local action term; when nonzero its support has uniformly bounded diameter.
Using the local moment stocks of (38.12), define a constant
\(\delta_{\rm St}\) large enough to dominate the local prefactors in
(39.4) and (39.6).

\begin{corollary}[Steiner response envelope at the cut anchor]
\label{cor:v18v29-anchor}
At \(s=0\), the exact chord-response matrix satisfies
\[
 |(K_c(0))_{ef}|
 \leq\delta_{\rm St}
 \left[
  \overline\tau^{d_T(\operatorname{supp}A_{ef},S_c)}
  \mathbf1_{\{A_{ef}\ne0\}}
  +\overline\tau^{\mathsf L_T(S_e,S_f,S_c)}
 \right].                                                       \tag{39.7}
\]
The constant is independent of the regulator volume.
\end{corollary}

\begin{proof}
Apply (39.6) to the covariance in (36.11) and Theorem
\ref{thm:v18v29-steiner} to its third cumulant.  The polynomial local
observables have uniform finite \(L^2\) and \(L^3\) gradient stocks under
the assumed quartic moment bounds, so all nongeometric factors are absorbed
in \(\delta_{\rm St}\).
\end{proof}

\subsection{A summable Steiner gate}

Define
\[
 \mathcal S_3(r)=sup_e
 \sum_{c\in\mathcal C}\sum_f
 r^{\mathsf L_T(S_e,S_f,S_c)},                                 \tag{39.8}
\]
and the two-point direct-score stock
\[
 \mathcal S_2(r)=sup_e
 \sum_{c\in\mathcal C}\sum_{f:A_{ef}\ne0}
 r^{d_T(\operatorname{supp}A_{ef},S_c)}.                       \tag{39.9}
\]

\begin{theorem}[Steiner-summed infinitesimal cycle gate]
\label{thm:v18v29-sum}
If \(\mathcal S_2(\overline\tau)\) and
\(\mathcal S_3(\overline\tau)\) are finite uniformly in the regulator,
then
\[
 \left\|\sum_{c\in\mathcal C}K_c(0)\right\|_{2\to2}
 \leq\delta_{\rm St}
 \left[
  \mathcal S_2(\overline\tau)+
  \mathcal S_3(\overline\tau)
 \right].                                                       \tag{39.10}
\]
Therefore the cut reference retains a positive Hessian margin for an
infinitesimal gluing interval \(0\leq s\leq s_0\) whenever a uniform
neighborhood of the cut anchor obeys the same bound and
\[
 s_0\delta_{\rm St}
 \left[
  \mathcal S_2(\overline\tau)+
  \mathcal S_3(\overline\tau)
 \right]<\mu_0.                                                \tag{39.11}
\]
\end{theorem}

\begin{proof}
Sum (39.7) first over \(c,f\) for a fixed row \(e\).  Equations
(39.8)--(39.9) bound the maximum absolute row sum.  The summed matrix is
symmetric, so its operator norm is at most that row sum, proving (39.10).
Integrate (36.12) on \([0,s_0]\) and use the uniform neighborhood bound to
obtain (39.11).
\end{proof}

\begin{proposition}[Simple graph-growth criterion]
\label{prop:v18v29-growth}
Suppose every radius-\(n\) tree ball contains at most \(C_g g^n\) local
edge and chord supports, with \(g\geq1\).  Then
\[
 \mathcal S_2(r)<\infty\quad\text{if }gr<1,
 \qquad
 \mathcal S_3(r)<\infty\quad\text{if }g^2r<1.                 \tag{39.12}
\]
For polynomial-growth regulator families, both stocks are finite for every
\(0<r<1\).
\end{proposition}

\begin{proof}
The number of two-support configurations in a distance shell of size
\(n\) is bounded by a constant times \(g^n\).  With one row fixed, the
number of pairs \((f,c)\) whose connecting Steiner subtree has length at
most \(n\) is bounded by a constant times \(g^{2n}\).  Summing the geometric
shell bounds proves (39.12).  Replacing the exponential ball bound by a
polynomial makes a polynomial times \(r^n\), which is summable for every
\(r<1\).
\end{proof}

\subsection{Correction of the working target}

The product-distance envelope (36.17) remains a sufficient hypothesis of
Theorem \ref{thm:v18v26-cycle-gate}; no earlier theorem is withdrawn.  It is
not the estimate naturally generated by the third cumulant when two score
supports share a long arm.  The correct next target is instead a version of
(39.7), uniform for \(0\leq s\leq1\), together with the Steiner summability
(39.8).  A third-order Dobrushin tree-graph inequality should propagate the
cut-anchor estimate through the gluing interpolation.

\begin{assessment}[V29 acquisition and boundary]
\label{ass:v18v29}
V29 proves Steiner-subtree decay for all three-observable configurations on
the cut tree, derives the corresponding exact chord-response envelope at
\(s=0\), and proves a volume-independent infinitesimal cycle gate.  It also
corrects the working geometry from a double-counted pairwise exponent to
Steiner length.  Uniform propagation to \(s=1\), physical coefficient
matching, cutoff removal, reconstruction, and the coupled Einstein
continuum remain open.
\end{assessment}

\begin{record}[V29 append record]
\label{rec:v18v29}
V29 is appended to the validated V28 whose installed SHA--256 was
\texttt{684CA4593D864FA40DBAAC7097CCE79F1A2EA23A5F30A3DA345BCC7E5EE6DD1C}.
After installation only V29 is the active numeric SM note; the frozen
\texttt{V100\_f17} archive is unchanged.
\end{record}

% =====================================================================
% V30 APPENDIX: LOCAL SECOND RESPONSE AND THREE-ARM RESOLVENT
% =====================================================================

\section{V30: local second response and the three-arm resolvent}
\label{sec:v18v30}

V29 identified the Steiner geometry required of a third-order Dobrushin
bound.  This note derives its local nonlinear coefficient and proves the
abstract resolvent theorem which propagates that coefficient through a
contractive cyclic specification.  The global majorant is a sum of three
resolvent arms meeting at one local branching site.  On a tree metric it has
exactly the Steiner decay of V29.

\subsection{Second source derivative of one conditional kernel}

Retain the conditional kernel (37.1) and the strong-convexity bound
\(\nabla^2U_v\succeq2A_vI\).  For source directions \(h,k\), write
\(D_h,D_k\) for derivatives in the boundary field \(y\).

\begin{proposition}[Exact local second response]
\label{prop:v18v30-local}
For every locally Lipschitz \(F\) with
\(\||\nabla F|\|_3<\infty\),
\[
 D_hD_kP_vF(y)
 =4\chi_v^2
 \operatorname{Cum}_{v,y}(F,z\cdot h,z\cdot k),               \tag{40.1}
\]
and
\[
 |D_hD_kP_vF(y)|
 \leq\beta_v\||\nabla F|\|_3|h||k|,                         \tag{40.2}
\]
where
\[
 \beta_v=\frac{108\chi_v^2}{(2A_v)^{3/2}}.                   \tag{40.3}
\]
\end{proposition}

\begin{proof}
The first derivative is (37.6).  Differentiating its covariance in the
direction \(k\) inserts the centered variable \(2\chi_vz\cdot k\), which
gives (40.1).  Apply Lemma \ref{lem:v18v28-lp} to the three variables
\(F,z\cdot h,z\cdot k\).  Their gradient \(L^3\) norms are respectively
\(\||\nabla F|\|_3,|h|,|k|\).  Thus the cumulant is at most
\[
 \frac{27}{(2A_v)^{3/2}}
 \||\nabla F|\|_3|h||k|.                                     \tag{40.4}
\]
Multiplication by \(4\chi_v^2\) proves (40.2)--(40.3).
\end{proof}

On the V25 coefficient sector one may use the uniform stock
\[
 \beta_*=
 \frac{108(\chi^*)^2}{[2(a_*+\gamma_*)]^{3/2}}               \tag{40.5}
\]
at every internal vertex.  A separate finite boundary-layer stock handles
leaves, as in V24.

\subsection{A nonlinear response fixed-point lemma}

Let \(C\) be a finite nonnegative matrix with
\(\|C\|_{\infty\to\infty}<1\), and put
\[
 D=(I-C)^{-1}=\sum_{n=0}^\infty C^n.                          \tag{40.6}
\]
Let \(B=\operatorname{diag}(\beta_v)\).  Products of vectors below are
componentwise.

\begin{lemma}[First and second response closure]
\label{lem:v18v30-fixed-point}
Suppose nonnegative first-response vectors \(x^{(p)}\) and a nonnegative
second connected-response vector \(y^{(p,q)}\) satisfy
\[
 x^{(p)}\leq e_p+Cx^{(p)},
 \qquad
 x^{(q)}\leq e_q+Cx^{(q)},                                    \tag{40.7}
\]
\[
 y^{(p,q)}leq Cy^{(p,q)}
     +B\bigl(x^{(p)}x^{(q)}\bigr).                            \tag{40.8}
\]
Then
\[
 x^{(p)}\leq De_p,qquad x^{(q)}\leq De_q,                    \tag{40.9}
\]
and
\[
 y^{(p,q)}_r
 \leq\sum_vD_{rv}\beta_vD_{vp}D_{vq}.                        \tag{40.10}
\]
\end{lemma}

\begin{proof}
Iterating (40.7) and using \(C^n\to0\) gives (40.9).  Move the first term
on the right of (40.8) through the positive resolvent:
\[
 y^{(p,q)}\leq DB(x^{(p)}x^{(q)})
 \leq DB[(De_p)(De_q)].                                       \tag{40.11}
\]
Its \(r\)-component is (40.10).
\end{proof}

\begin{theorem}[Three-arm response of a contractive specification]
\label{thm:v18v30-three-arm}
Consider a differentiable finite-volume heat-bath specification whose
first variation seminorms are dominated by the Dobrushin matrix \(C\) of
(37.14), and whose local second variation is dominated by (40.2).  Suppose
the heat-bath iteration converges in the same seminorm and differentiation
commutes with that limit.  Then the connected second response to source
perturbations at \(p,q\), observed at \(r\), is bounded by
\[
 \mathcal T_{r;pq}
 :=\sum_vD_{rv}\beta_vD_{vp}D_{vq}.                           \tag{40.12}
\]
\end{theorem}

\begin{proof}
At every finite heat-bath iteration, the chain rule and the local first
variation bound give (40.7).  Differentiating once more gives a transported
second variation plus the local bilinear second response; their seminorms
obey (40.8), with Proposition \ref{prop:v18v30-local} supplying \(B\).
Lemma \ref{lem:v18v30-fixed-point} bounds every iterate uniformly.  Pass to
the differentiable limit to obtain (40.12).
\end{proof}

The differentiable-limit hypothesis is not an extra physical assertion.  It
is a precise analytic gate.  It follows, for example, from uniform
integrability of the first two source derivatives together with contraction
of the heat-bath iteration.  Quartic moment bounds supply the former once a
uniform source box is fixed.

\subsection{Steiner decay of the resolvent tree}

Assume the underlying metric graph has a unique tree geodesic for the
present estimate; this includes the cut reference tree.  Suppose
\[
 D_{uv}\leq d_0\zeta^{d(u,v)},
 \qquad 0<\zeta<1,                                                 \tag{40.13}
\]
and \(\beta_v\leq\beta_*\).  Let \(o=o(r,p,q)\) be the median of the three
sites and \(\mathsf L_T(r,p,q)\) their Steiner length.

\begin{theorem}[Resolvent Steiner bound]
\label{thm:v18v30-steiner}
If
\[
 \mathcal V(\zeta):=\sup_o\sum_v\zeta^{d(v,o)}<\infty,                \tag{40.14}
\]
then
\[
 \mathcal T_{r;pq}
 \leq\beta_*d_0^3\mathcal V(\zeta)
 \zeta^{\mathsf L_T(r,p,q)}.                                      \tag{40.15}
\]
For a graph with at most \(C_gg^n\) vertices in a radius-\(n\) ball,
\(\mathcal V(\zeta)<\infty\) whenever \(g\zeta<1\); on polynomial-growth graphs it
is finite for every \(\zeta<1\).
\end{theorem}

\begin{proof}
For every vertex \(v\) in a tree,
\[
 d(v,r)+d(v,p)+d(v,q)
 \geq\mathsf L_T(r,p,q)+d(v,o).                               \tag{40.16}
\]
Indeed, moving away from the median increases the sum of the three
distances by at least the distance moved; an excursion off the Steiner
subtree increases it three times.  Insert (40.13) into (40.12), apply
(40.16), and sum (40.14).  The ball-growth statements follow by summing
the corresponding distance shells.
\end{proof}

Compared with the crude pair count in (39.12), the three-arm resolvent keeps
the common propagation before branching.  Its sufficient exponential
growth condition is \(g\zeta<1\), rather than \(g^2\zeta<1\).

\subsection{Interface with chord gluing}

The third cumulant in (36.11) is a connected second response of one score
to two local insertions after polarization and finite-block conditioning.
Theorems \ref{thm:v18v30-three-arm} and
\ref{thm:v18v30-steiner} therefore provide the correct majorant once the
score and chord blocks are assigned their local first- and second-variation
stocks.  Collision configurations contribute only to \(\beta_v\) in the
finite branching core; long arms are carried by \(D\).

What remains is to write the conformal-score/chord-insertion polarization
explicitly and verify the differentiable heat-bath limit uniformly for
\(0\leq s\leq1\).  This is narrower than proving a new cluster expansion:
the contraction, local second derivative, and summable three-arm kernel are
now available separately.

\begin{assessment}[V30 acquisition and boundary]
\label{ass:v18v30}
V30 proves an explicit local second source-response bound, the nonlinear
first/second response resolvent lemma, and a three-arm Steiner bound with a
volume-independent growth criterion.  These are the analytic components of
the third-order Dobrushin estimate sought in V29.  Uniform differentiability
of the cyclic heat-bath limit and the observable polarization for (36.11)
remain to be certified.  Physical matching, cutoff removal, reconstruction,
and the coupled Einstein continuum remain open.
\end{assessment}

\begin{record}[V30 append record]
\label{rec:v18v30}
V30 is appended to the validated V29 whose installed SHA--256 was
\texttt{9576CA1B8DAC464FB67E3A2CBBCE2491B67DAB1C3F422358348811663E93DF6C}.
After installation only V30 is the active numeric SM note; the frozen
\texttt{V100\_f17} archive is unchanged.
\end{record}

% =====================================================================
% V31 APPENDIX: AUDIT, INSERTION CALCULUS, AND C2 FIXED POINT
% =====================================================================

\section{V31: companion audit and compiled insertion response}
\label{sec:v18v31}

The mandatory V30 companion audit was completed before this continuation.
The active YM note was
\[
 \texttt{Renewal\_Yang\_Mills\_Mass\_Gap\_Research\_Notes\_V4.tex}
\]
with SHA--256
\[
 \texttt{A0EECCBF3384C6B55D28E9115B3DCB1803D9E78407A39CCFBF13C42A1A94EA3B}.
                                                               \tag{41.1}
\]
Its new result enforces the moving coarse-link constraint before inversion,
reduces the action-dependent Schur correction to rank three, and compiles
complete signed response words before applying interval norms.  The active
NS note remained
\[
 \texttt{Renewal\_NS\_Stability\_Research\_Notes\_V26.tex}
\]
with SHA--256
\[
 \texttt{82C887F8C2C69E4F28FAD7C3DC8DE5B9099CB5A4BF431EBA1FFF3E91935978AD}.
                                                               \tag{41.2}
\]
Its rank-one comparison is unchanged.  Neither companion result is imported
as an SM theorem.  The useful decision is to keep the two signed terms in
(36.11) together inside their finite collision core and transport the
compiled core only after its exact local response has been formed.

\subsection{General local insertion derivatives}

Let
\[
 d\nu_{\lambda,\mu}(z)=\frac1{Z_{\lambda,\mu}}
 e^{-U(z)-\lambda G(z)-\mu H(z)}\,dz                          \tag{41.3}
\]
and write \(P_{\lambda,\mu}F=\mathbb E_{\lambda,\mu}F\).  The
observables below are independent of \(\lambda,\mu\).

\begin{proposition}[Exact insertion calculus]
\label{prop:v18v31-insertion}
At \((\lambda,\mu)=(0,0)\),
\[
 \partial_\lambda PF=-\operatorname{Cov}(F,G),                \tag{41.4}
\]
\[
 \partial_\lambda\partial_\mu PF
 =\operatorname{Cum}(F,G,H).                                  \tag{41.5}
\]
If \(\nabla^2U\succeq2AI\), then
\[
 |\partial_\lambda PF|
 \leq\frac1{2A}
 \||\nabla F|\|_2\||\nabla G|\|_2,                              \tag{41.6}
\]
\[
 |\partial_\lambda\partial_\mu PF|
 \leq\frac{27}{(2A)^{3/2}}
 \||\nabla F|\|_3
 \||\nabla G|\|_3
 \||\nabla H|\|_3.                                           \tag{41.7}
\]
\end{proposition}

\begin{proof}
Differentiate the normalized integral in (41.3).  The first derivative is
the centered insertion (41.4); differentiating it in \(\mu\) gives the
centered triple product (41.5).  Brascamp--Lieb and covariance
Cauchy--Schwarz give (41.6).  Lemma \ref{lem:v18v28-lp} with \(p=3\) gives
(41.7).
\end{proof}

Unlike Proposition \ref{prop:v18v30-local}, this result permits arbitrary
local polynomial insertions.  It therefore applies directly to the
conformal scores \(A_e,A_f,A_{ef}\) and the bilinear chord insertion
\(W_c\), without representing them as linear source directions.

\subsection{The signed collision-core compiler}

For a finite collision block \(B\), condition all exterior fields and let
\(\nu_{B,s}\) be its conditional law along the gluing interpolation.  Define
the complete local response word
\[
 \mathscr R_{efc}^{B}(s)
 :=-\operatorname{Cov}_{B,s}(A_{ef},W_c)
   +\operatorname{Cum}_{B,s}(A_e,A_f,W_c).                    \tag{41.8}
\]
This is the local version of the full matrix entry (36.11).

\begin{proposition}[Rigorous coarse collision stock]
\label{prop:v18v31-core}
Suppose every allowed collision block has
\(\nabla^2U_{B,s}\succeq2A_BI\) for \(0\leq s\leq1\).  Then
\[
 |\mathscr R_{efc}^{B}(s)|
 \leq \mathfrak r_{efc}^{B},                                  \tag{41.9}
\]
where
\[
\begin{split}
 \mathfrak r_{efc}^{B}:={}&
 \frac1{2A_B}
 \||\nabla A_{ef}|\|_2\||\nabla W_c|\|_2\\
 &+\frac{27}{(2A_B)^{3/2}}
 \||\nabla A_e|\|_3
 \||\nabla A_f|\|_3
 \||\nabla W_c|\|_3.                                      \tag{41.10}
\end{split}
\]
All norms are conditional block norms and may be replaced by uniform
quartic moment envelopes.
\end{proposition}

\begin{proof}
Apply (41.6) to the first term of (41.8) and (41.7) to the second, then use
the triangle inequality once, after each complete centered response has
been formed.
\end{proof}

The stock (41.10) is intentionally safe.  A finite interval certificate
should evaluate the signed quantity (41.8) as a unit; any cancellation
between its covariance and cumulant is then retained.  This is the precise
SM analogue of the compilation order suggested by the YM V4 audit.  It is
also finite: bounded interaction range, bounded degree, and a fixed
collision radius leave only finitely many typed blocks.

\subsection{Differentiability of a contractive fixed point}

The differentiable-limit hypothesis in V30 can be replaced by a standard
uniform contraction criterion.

\begin{theorem}[Uniform \(C^2\) contraction fixed point]
\label{thm:v18v31-c2-fixed}
Let \(X\) and \(P\) be Banach spaces, let \(I\subset P\) be open, and let
\(T:I\times X\to X\) be twice continuously Fr\'echet differentiable.
Assume a closed set \(\mathcal B\subset X\) is invariant and
\[
 \sup_{(s,x)\in I\times\mathcal B}\|D_xT(s,x)\|\leq\kappa<1. \tag{41.11}
\]
Then the unique fixed point \(x(s)=T(s,x(s))\) is \(C^2\) and
\[
 D_sx=(I-D_xT)^{-1}D_sT,                                     \tag{41.12}
\]
\[
 D_s^2x=(I-D_xT)^{-1}
 \left[D_s^2T+2D_{sx}^2T[D_sx]
       +D_x^2T[D_sx,D_sx]\right].                            \tag{41.13}
\]
Moreover,
\[
 \|(I-D_xT)^{-1}\|\leq\frac1{1-\kappa}.                     \tag{41.14}
\]
\end{theorem}

\begin{proof}
The contraction theorem gives a unique fixed point in \(\mathcal B\).
The map
\[
 \mathcal F(s,x)=x-T(s,x)                                    \tag{41.15}
\]
has derivative \(D_x\mathcal F=I-D_xT\).  The Neumann series converges by
(41.11), proving invertibility and (41.14).  The Banach implicit-function
theorem makes \(x(s)\) twice continuously differentiable.  Differentiate
\(x=T(s,x)\) once and twice and solve with the same inverse to obtain
(41.12)--(41.13).
\end{proof}

\begin{corollary}[Heat-bath differentiability gate]
\label{cor:v18v31-heatbath}
Suppose the cyclic heat-bath update acts on a weighted observable-response
Banach space, preserves a uniform quartic-moment ball, is \(C^2\) there, and
its field derivative is dominated by a Dobrushin matrix with row sum
\(\kappa<1\) uniformly for \(0\leq s\leq1\).  Then its stationary response
is \(C^2\) in the chord source parameters, and the differentiable-limit
hypothesis of Theorem \ref{thm:v18v30-three-arm} holds.
\end{corollary}

\begin{proof}
Use Theorem \ref{thm:v18v27-dobrushin} to obtain (41.11) in the weighted
response norm.  Quartic domination and Proposition
\ref{prop:v18v31-insertion} give the required first and second derivative
continuity on the invariant moment ball.  Apply Theorem
\ref{thm:v18v31-c2-fixed}.
\end{proof}

\subsection{Compiled three-arm cycle gate}

Let \(\mathfrak r_*\) be a certified supremum of the complete local core
response (41.8), using either the coarse stock (41.10) or a sharper signed
interval evaluation.  Let \(D=(I-C)^{-1}\) be the Dobrushin resolvent.

\begin{theorem}[Conditional full-interpolation chord response]
\label{thm:v18v31-chord}
Assume the heat-bath hypotheses of Corollary
\ref{cor:v18v31-heatbath}, a finite family of collision cores, and a
uniform compiled core stock \(\mathfrak r_*\).  Then every connected
three-leg chord response is bounded by
\[
 \mathfrak r_*
 \sum_vD_{rv}D_{vp}D_{vq},                                   \tag{41.16}
\]
with the legs attached to the three relevant local supports.  If the
resolvent satisfies (40.13)--(40.14), this is at most
\[
 \mathfrak r_*d_0^3\mathcal V(\zeta)
 \zeta^{\mathsf L_T(r,p,q)}                                  \tag{41.17}
\]
uniformly for \(0\leq s\leq1\).
\end{theorem}

\begin{proof}
The fixed-point derivatives satisfy the response inequalities
(40.7)--(40.8).  The local inhomogeneous term in (40.8) is the compiled
block derivative (41.8), bounded by \(\mathfrak r_*\).  Lemma
\ref{lem:v18v30-fixed-point} gives the three-resolvent sum (41.16), and
Theorem \ref{thm:v18v30-steiner} gives (41.17).
\end{proof}

The theorem is conditional on a concrete Banach-space verification, not on
an unspecified differentiable limit.  The remaining analytic work is to
choose the weighted polynomial-response norm, prove that the heat-bath map
preserves its quartic moment ball throughout the gluing path, and compute
the finite core stock.  Those requirements are local and certificate-ready.

\begin{assessment}[V31 acquisition and boundary]
\label{ass:v18v31}
V31 records the mandatory companion audit, proves general polynomial
insertion derivatives, defines a signed finite collision-core compiler, and
replaces V30's differentiable-limit assumption by a uniform \(C^2\)
contraction fixed-point theorem.  It yields the desired full-interpolation
three-arm response under an explicit weighted heat-bath gate.  Preservation
of that weighted moment ball and a physical strict coefficient card remain
unproved.  Cutoff removal, reconstruction, and the coupled Einstein
continuum remain open.
\end{assessment}

\begin{record}[V31 append record]
\label{rec:v18v31}
V31 is appended to the validated V30 whose installed SHA--256 was
\texttt{DBD7A5C78AF216B2730894555E3B5B243840F50C8425904E5562D1C3FA28A4CC}.
After installation only V31 is the active numeric SM note; the frozen
\texttt{V100\_f17} archive is unchanged.
\end{record}

% =====================================================================
% V32 APPENDIX: UNIFORM QUARTIC MOMENT BALL ALONG CHORD GLUING
% =====================================================================

\section{V32: a uniform polynomial-response moment ball}
\label{sec:v18v32}

V31 reduced differentiability of the cyclic heat-bath fixed point to
preservation of a weighted quartic-moment ball.  This note constructs the
required ball for the polynomial degrees occurring in the conformal scores
and chord insertions.  The estimate is uniform in the volume and in the
gluing parameter \(0\leq s\leq1\); it does not require a small chord
coupling.

\subsection{Conditional quartic moment estimate}

Consider a one-site conditional law
\[
 d\nu_h(z)=\frac1{Z_h}
 \exp\{-\lambda|z|^4-A|z|^2-\Phi(z)+2h\cdot z\}\,dz,           \tag{42.1}
\]
where \(z\in\mathbb R^N\), \(\lambda\geq\lambda_*>0\),
\(A\geq0\), and
\[
 z\cdot\nabla\Phi(z)\geq0.                                   \tag{42.2}
\]

\begin{lemma}[Quartic fourth-moment stock]
\label{lem:v18v32-fourth}
There is a numerical constant \(C_Y=3\,2^{-5/3}\) such that
\[
 \mathbb E_h|z|^4
 \leq\frac{N}{2\lambda_*}
 +\frac{C_Y}{2\lambda_*^{4/3}}|h|^{4/3}.                     \tag{42.3}
\]
Consequently, for every \(0<p\leq4\), there are explicit finite constants
\(C_{p,0},C_{p,1}\), depending only on \(N,p,\lambda_*\), with
\[
 \mathbb E_h|z|^p
 \leq C_{p,0}+C_{p,1}|h|^{p/3}.                               \tag{42.4}
\]
\end{lemma}

\begin{proof}
Integration of
\(\nabla_z\cdot(ze^{-\lambda|z|^4-A|z|^2-\Phi(z)+2h\cdot z})\)
gives
\[
 N=4\lambda\mathbb E|z|^4+2A\mathbb E|z|^2
   +\mathbb E[z\cdot\nabla\Phi(z)]-2h\cdot\mathbb Ez.        \tag{42.5}
\]
By (42.2),
\[
 4\lambda\mathbb E|z|^4
 \leq N+2|h|\mathbb E|z|.                                    \tag{42.6}
\]
Young's inequality in the form
\[
 2Hr\leq2\lambda r^4+C_Y\lambda^{-1/3}H^{4/3}               \tag{42.7}
\]
then gives (42.3).  Finally,
\(\mathbb E|z|^p\leq(\mathbb E|z|^4)^{p/4}\); subadditivity of
\(x\mapsto x^{p/4}\) on \([0,\infty)\) gives (42.4).
\end{proof}

For \(p=3\), record the simpler form
\[
 \mathbb E_h|z|^3\leq C_{3,0}+C_{3,1}|h|.                    \tag{42.8}
\]

\subsection{Closing the graph-wide third moment}

Along the interpolation (36.5), the conditional source at a vertex has the
form
\[
 h_v(s)=h_v^{\rm ext}
 +\sum_{w\sim v}s_{vw}\chi_{vw}U_{vw}z_w,
 \qquad 0\leq s_{vw}\leq1,                                   \tag{42.9}
\]
where every \(U_{vw}\) is an isometry.  Assume degree at most \(\Delta\),
\(\chi_{vw}\leq\chi^*\), and
\(|h_v^{\rm ext}|\leq h_*\).

\begin{theorem}[Uniform third-moment ball]
\label{thm:v18v32-moment-ball}
Let
\[
 K_0=C_{3,0}+C_{3,1}h_*,
 \qquad K_1=C_{3,1}\Delta\chi^*,                              \tag{42.10}
\]
and
\[
 M_{3,*}=\max\left\{2K_0,(2K_1)^{3/2}\right\}.               \tag{42.11}
\]
Then every finite-volume Gibbs law associated with (42.9) satisfies
\[
 \sup_v\mathbb E_s|z_v|^3\leq M_{3,*}                         \tag{42.12}
\]
for all \(0\leq s\leq1\), uniformly in the volume.
\end{theorem}

\begin{proof}
Let \(M=\max_v\mathbb E_s|z_v|^3\), which is finite in every finite volume
by quartic confinement.  Condition at \(v\), apply (42.8), and then use
Minkowski and H\"older:
\[
\begin{split}
 \mathbb E_s|z_v|^3
 &\leq C_{3,0}+C_{3,1}\mathbb E_s|h_v(s)|\\
 &\leq K_0+C_{3,1}\sum_{w\sim v}\chi_{vw}\mathbb E_s|z_w|\\
 &\leq K_0+K_1M^{1/3}.                                        \tag{42.13}
\end{split}
\]
Thus \(M\leq K_0+K_1M^{1/3}\).  If \(M\leq2K_0\), (42.11) applies.  If
\(M>2K_0\), then
\(M/2<K_1M^{1/3}\), so \(M<(2K_1)^{3/2}\).  This proves (42.12).
\end{proof}

The same argument with (42.4) supplies volume-uniform moments of every
order \(p\leq4\).  Higher polynomial orders can be obtained by multiplying
(42.5) by powers of \(|z|\) and iterating; the present chord-response
calculus needs only the third moment.

\subsection{Explicit score-gradient stocks}

For an edge \(e=(x,y)\), the conformal score and direct score are
\[
 A_e=-a_e|z_x|^2+b_e|z_y|^2,
 \qquad
 A_{ee}=a_e|z_x|^2+b_e|z_y|^2,                                \tag{42.14}
\]
and \(A_{ef}=0\) for distinct independently parametrized edges.  The chord
insertion is
\[
 W_e=-2\chi_e z_x\cdot U_ez_y.                                \tag{42.15}
\]
Let \(a_e\leq a^*\), \(b_e\leq b^*\), and \(\chi_e\leq\chi^*\).

\begin{proposition}[Uniform local response gradients]
\label{prop:v18v32-gradient-stocks}
Under Theorem \ref{thm:v18v32-moment-ball},
\[
 \||\nabla A_e|\|_3,
 \ \||\nabla A_{ee}|\|_3
 \leq2(a^*+b^*)M_{3,*}^{1/3},                                \tag{42.16}
\]
\[
 \||\nabla W_e|\|_3
 \leq4\chi^*M_{3,*}^{1/3}.                                   \tag{42.17}
\]
The corresponding \(L^2\) bounds follow by monotonicity of probability
norms.
\end{proposition}

\begin{proof}
Differentiate (42.14)--(42.15) in their endpoint fields.  Pointwise,
\[
 |\nabla A_e|,|\nabla A_{ee}|
 \leq2a^*|z_x|+2b^*|z_y|,                                    \tag{42.18}
\]
\[
 |\nabla W_e|
 \leq2\chi^*(|z_x|+|z_y|).                                   \tag{42.19}
\]
Minkowski's inequality and (42.12) give (42.16)--(42.17).
\end{proof}

Inserting these stocks into (41.10) gives a completely explicit coarse
collision constant once the finite-block strong-convexity floor \(A_B\) is
known.  A signed interval evaluation of (41.8) can only improve that stock.

\subsection{Consequences for the heat-bath gate}

\begin{corollary}[Uniform integrability of first and second insertions]
\label{cor:v18v32-uniform-integrability}
On the compact coefficient sector above, all first and second derivatives
in Proposition \ref{prop:v18v31-insertion} involving
\(A_e,A_{ee},W_c\) are uniformly integrable for
\(0\leq s\leq1\), with constants independent of the volume.
\end{corollary}

\begin{proof}
Equations (42.16)--(42.17) give the \(L^2\) and \(L^3\) gradient stocks in
(41.6)--(41.7).  The observables themselves have degree at most two, so the
fourth-moment part of Lemma \ref{lem:v18v32-fourth} supplies their required
uniform integrability.  The estimates are uniform in \(s\) by Theorem
\ref{thm:v18v32-moment-ball}.
\end{proof}

This proves the moment-ball and dominated-differentiation part of Corollary
\ref{cor:v18v31-heatbath}.  The remaining part is functional: choose a
weighted response norm in which the complete heat-bath update has derivative
strictly below one.  The one-site coefficients are already the Dobrushin
stocks (37.14); V33 should construct the weight and state its exact finite
matrix certificate.

\begin{assessment}[V32 acquisition and boundary]
\label{ass:v18v32}
V32 proves a volume- and interpolation-uniform third-moment ball from the
quartic integration-by-parts identity, derives explicit gradient stocks for
all conformal scores and chord insertions, and verifies uniform
integrability of the first two insertion derivatives.  The weighted
Dobrushin contraction and finite-block strong-convexity certificate remain
unproved.  Physical matching, cutoff removal, reconstruction, and the
coupled Einstein continuum remain open.
\end{assessment}

\begin{record}[V32 append record]
\label{rec:v18v32}
V32 is appended to the validated V31 whose installed SHA--256 was
\texttt{31211C5713D96723B4130630F542F50B6DD89E74DE0007CB228EA8E5E3DAECCE}.
After installation only V32 is the active numeric SM note; the frozen
\texttt{V100\_f17} archive is unchanged.
\end{record}

% =====================================================================
% V33 APPENDIX: WEIGHTED DOBRUSHIN AND NORMALIZED HOPPING SPECTRUM
% =====================================================================

\section{V33: one spectral certificate for contraction and block convexity}
\label{sec:v18v33}

V31--V32 left two related hypotheses: a weighted norm in which the cyclic
heat-bath derivative is strictly contractive, and a strong-convexity floor
for every finite collision block.  Both follow from the same normalized
hopping matrix.  This note proves the equivalence and gives a rational
trial-weight certificate which avoids diagonalizing a regulator-sized
matrix.

\subsection{The normalized hopping matrix}

Let \(A_v>0\) be a certified lower half-Hessian of the one-site potential at
vertex \(v\):
\[
 \nabla_{z_v}^2U_v\succeq2A_vI.                               \tag{43.1}
\]
For an undirected edge \(v\sim w\), let \(\chi_{vw}=\chi_{wv}\geq0\), and
define
\[
 C_{vw}=\frac{\chi_{vw}}{A_v},                                \tag{43.2}
\]
\[
 S_{vw}=\frac{\chi_{vw}}{\sqrt{A_vA_w}},
 \qquad S_{vv}=0.                                             \tag{43.3}
\]
Then \(S\) is symmetric and
\[
 S=A^{1/2}CA^{-1/2},
 \qquad A=\operatorname{diag}(A_v).                           \tag{43.4}
\]

\begin{theorem}[Weighted Dobrushin--spectral equivalence]
\label{thm:v18v33-equivalence}
For a finite connected graph, the following are equivalent:
\begin{enumerate}[leftmargin=2.2em]
\item there are weights \(w_v>0\) and \(0\leq\kappa<1\) such that
      \[
      \sum_wC_{vw}w_w\leq\kappa w_v
      \quad\text{for every }v;                               \tag{43.5}
      \]
\item the Perron root satisfies \(\rho(C)<1\);
\item the symmetric normalized hopping matrix satisfies
      \[
      \lambda_{\max}(S)<1.                                  \tag{43.6}
      \]
\end{enumerate}
For weights satisfying (43.5), the heat-bath influence is a contraction in
the weighted supremum norm
\[
 \|x\|_w=\max_v\frac{|x_v|}{w_v}                              \tag{43.7}
\]
with contraction factor at most \(\kappa\).
\end{theorem}

\begin{proof}
Equation (43.5) is \(Cw\leq\kappa w\).  The Collatz--Wielandt formula for a
nonnegative irreducible matrix gives
\[
 \rho(C)=\inf_{w>0}\max_v\frac{(Cw)_v}{w_v}.                  \tag{43.8}
\]
This proves the equivalence of the first two statements; reducible graphs
follow componentwise.  Similarity (43.4) proves
\(\rho(C)=\rho(S)=\lambda_{\max}(S)\), since a symmetric nonnegative matrix
has Perron root equal to its operator spectral radius.  Finally,
\[
 \|Cx\|_w
 \leq\max_v\frac{\sum_wC_{vw}w_w}{w_v}\|x\|_w
 \leq\kappa\|x\|_w,                                         \tag{43.9}
\]
which proves the contraction claim.
\end{proof}

\subsection{The same certificate controls the joint field Hessian}

Let \(U_{vw}\) be orthogonal parallel transports with
\(U_{wv}=U_{vw}^{\mathsf T}\).  Along the gluing interpolation, the field
action has cross terms
\[
 -2s_{vw}\chi_{vw}z_v\cdot U_{vw}z_w,
 \qquad 0\leq s_{vw}\leq1.                                  \tag{43.10}
\]

\begin{theorem}[Normalized-hopping strong convexity]
\label{thm:v18v33-convexity}
If
\[
 \lambda_{\max}(S)\leq\kappa<1,                              \tag{43.11}
\]
then the complete field Hessian satisfies
\[
 \nabla_z^2S_s\succeq
 2(1-\kappa)\operatorname{diag}(A_vI)                         \tag{43.12}
\]
for every choice \(0\leq s_{vw}\leq1\).  In particular,
\[
 \nabla_z^2S_s\succeq2A_{\rm joint}I,
 \qquad
 A_{\rm joint}:=(1-\kappa)A_{\min}.                          \tag{43.13}
\]
Every conditional finite collision block has the same lower floor.
\end{theorem}

\begin{proof}
For a field variation \(\xi=(\xi_v)\), put
\(r_v=\sqrt{A_v}|\xi_v|\).  Equations (43.1) and (43.10), together with
orthogonality of \(U_{vw}\), give
\[
\begin{split}
 \frac12\xi^{\mathsf T}(\nabla_z^2S_s)\xi
 &\geq\sum_vA_v|\xi_v|^2
       -2\sum_{v<w}s_{vw}\chi_{vw}|\xi_v||\xi_w|\\
 &\geq r^{\mathsf T}(I-S)r
 \geq(1-\kappa)|r|^2.                                        \tag{43.14}
\end{split}
\]
This is (43.12)--(43.13).  A conditional block Hessian is a principal
submatrix after exterior fields have become linear sources, so the same
quadratic-form lower bound applies.
\end{proof}

Thus the block constant \(A_B\) in (41.9)--(41.10) can be chosen as
\(A_{\rm joint}\), uniformly in the collision type, volume, and gluing
parameter.

\subsection{Finite certificates without spectral diagonalization}

\begin{proposition}[Positive trial-weight certificate]
\label{prop:v18v33-trial}
Let \(w_v>0\) be any proposed weights.  If rigorous intervals prove
\[
 \max_v\frac1{A_vw_v}
 \sum_{w\sim v}\chi_{vw}w_w\leq\kappa<1,                     \tag{43.15}
\]
then both the weighted heat-bath contraction and the joint Hessian floor
(43.12) hold.  Rational coefficient bounds and rational trial weights make
(43.15) a finite exact certificate.
\end{proposition}

\begin{proof}
Equation (43.15) is (43.5).  Apply Theorems
\ref{thm:v18v33-equivalence} and \ref{thm:v18v33-convexity}.
\end{proof}

\begin{corollary}[Unweighted row-sum gate]
\label{cor:v18v33-row}
The simpler condition
\[
 \max_v\frac1{A_v}\sum_{w\sim v}\chi_{vw}\leq\kappa<1      \tag{43.16}
\]
is sufficient.
\end{corollary}

\begin{proof}
Use \(w_v=1\) in Proposition \ref{prop:v18v33-trial}.
\end{proof}

The weighted test can be strictly sharper on inhomogeneous graphs.  It also
retains the complete neighbor sum before taking its maximum, consistent with
the compilation lesson of the V30 companion audit.

\subsection{Closing the analytic hypotheses of the compiled response}

\begin{theorem}[Spectral closure of the V31 heat-bath gate]
\label{thm:v18v33-closure}
Assume the compact coefficient sector and quartic moment hypotheses of V32,
and suppose (43.15) holds uniformly over the allowed regulator family with
one \(\kappa<1\).  Then:
\begin{enumerate}[leftmargin=2.2em]
\item the heat-bath update is contractive in the weighted norm (43.7);
\item the uniform third-moment ball (42.12) supplies the first two insertion
      derivatives;
\item every collision block has strong-convexity constant
      \(A_B\geq(1-\kappa)A_{\min}\);
\item the stationary response is \(C^2\) and the compiled three-arm bound
      (41.16)--(41.17) holds uniformly for \(0\leq s\leq1\).
\end{enumerate}
\end{theorem}

\begin{proof}
Items one and three are Proposition \ref{prop:v18v33-trial} and Theorem
\ref{thm:v18v33-convexity}.  Item two is Corollary
\ref{cor:v18v32-uniform-integrability}.  These verify every hypothesis of
Corollary \ref{cor:v18v31-heatbath}; Theorem
\ref{thm:v18v31-chord} then proves item four.
\end{proof}

This theorem closes the analytic cycle-response mechanism under one finite
strict matrix inequality.  It still does not prove that the physical
Standard Model coefficient sector satisfies (43.15), nor that the integrated
response loss stays below the cut-tree margin.  Those are now two explicit
numerical-symbolic comparisons:
\[
 \kappa<1,
 \qquad
 \int_0^1\left\|\sum_cK_c(s)\right\|ds<\mu_0.                 \tag{43.17}
\]

\begin{assessment}[V33 acquisition and boundary]
\label{ass:v18v33}
V33 proves equivalence of weighted Dobrushin contraction and a normalized
hopping spectral bound, proves that the same bound gives uniform joint and
collision-block strong convexity, and supplies a finite rational
trial-weight certificate.  Together with V31--V32 it closes the
full-interpolation response machinery under (43.15).  The physical strict
card and the final integrated cycle-margin comparison in (43.17) remain
unproved.  Cutoff removal, reconstruction, and the coupled Einstein
continuum remain open.
\end{assessment}

\begin{record}[V33 append record]
\label{rec:v18v33}
V33 is appended to the validated V32 whose installed SHA--256 was
\texttt{0724889C70F0E181431F6D91CDEEC558F3EA2CD5DF44760A036806F4369CBA00}.
After installation only V33 is the active numeric SM note; the frozen
\texttt{V100\_f17} archive is unchanged.
\end{record}

% =====================================================================
% V34 APPENDIX: INTEGRATED CYCLE MARGIN
% =====================================================================

\section{V34: summing the full chord response}
\label{sec:v18v34}

V33 closes the local and differentiability hypotheses of the chord-response
calculus under one normalized hopping certificate.  This note performs the
remaining global sum.  The two-point part of (36.11) propagates through two
Dobrushin resolvent arms; the connected three-point part propagates through
three.  Positivity and Fubini factor their total row sums into elementary
resolvent stocks.  The result is an explicit sufficient margin for the full
cyclic conformal Hessian.

\subsection{Attachment sets and resolvent stocks}

Choose one attachment site for each bounded local support of an edge score
and one for each chord collision block.  Enlarging a support by a uniformly
bounded radius changes only the local constants below.  Let
\[
 D=(I-C)^{-1}                                                   \tag{44.1}
\]
be the nonnegative Dobrushin resolvent.  Write \(\mathcal E\) for score
attachments and \(\mathcal C\) for chord attachments.  Define
\[
 M_{\rm in}=\sup_{e\in\mathcal E}\sum_vD_{ev},                \tag{44.2}
\]
\[
 M_E=\sup_v\sum_{f\in\mathcal E}D_{vf},
 \qquad
 M_C=\sup_v\sum_{c\in\mathcal C}D_{vc}.                      \tag{44.3}
\]
For a finite graph these quantities are finite.  The required condition is
that they remain bounded uniformly along the regulator family.

Define the two-arm and three-arm kernels
\[
 \mathcal U_{e;c}=\sum_vD_{ev}D_{vc},                         \tag{44.4}
\]
\[
 \mathcal T_{e;fc}=\sum_vD_{ev}D_{vf}D_{vc}.                  \tag{44.5}
\]

\begin{lemma}[Resolvent arm summation]
\label{lem:v18v34-arms}
The positive kernels satisfy
\[
 \sup_e\sum_c\mathcal U_{e;c}
 \leq M_{\rm in}M_C,                                         \tag{44.6}
\]
\[
 \sup_e\sum_f\sum_c\mathcal T_{e;fc}
 \leq M_{\rm in}M_EM_C.                                      \tag{44.7}
\]
\end{lemma}

\begin{proof}
All summands are nonnegative, so Tonelli's theorem gives
\[
 \sum_c\mathcal U_{e;c}
 =\sum_vD_{ev}\sum_cD_{vc}
 \leq M_C\sum_vD_{ev}\leq M_{\rm in}M_C.                    \tag{44.8}
\]
Similarly,
\[
 \sum_{f,c}\mathcal T_{e;fc}
 =\sum_vD_{ev}
       \left(\sum_fD_{vf}\right)
       \left(\sum_cD_{vc}\right)
 \leq M_EM_C\sum_vD_{ev}.                                    \tag{44.9}
\]
This proves (44.7).
\end{proof}

If the weighted contraction of V33 uses weights with
\[
 0<w_*\leq w_v\leq w^*<\infty                                \tag{44.10}
\]
uniformly, then (43.5) implies the crude but useful bounds
\[
 M_{\rm in},M_E\leq\frac{w^*}{w_*(1-\kappa)}.                 \tag{44.11}
\]
A uniformly bounded chord density gives the analogous bound for \(M_C\).
Sharper spatial decay from (40.13) can be inserted directly.

\subsection{Compiled local coefficients}

Let \(\alpha_*\) be a uniform bound for the complete local two-observable
response associated with
\(-\operatorname{Cov}(A_{ef},W_c)\), after all colliding supports have been
compiled.  Since \(A_{ef}=0\) for distinct independent edge parameters,
this part is diagonal in \((e,f)\) before propagation.  Let
\(\mathfrak r_*\) be the uniform complete three-observable core stock of
(41.8).  V31--V33 give the safe choices
\[
 \alpha_*\leq\frac{M_{2,A''}M_{2,W}}{2A_{\rm joint}},         \tag{44.12}
\]
\[
 \mathfrak r_*\leq
 \frac{27M_{3,A}^2M_{3,W}}{(2A_{\rm joint})^{3/2}},           \tag{44.13}
\]
where the moment-gradient stocks are explicitly bounded in V32 and
\[
 A_{\rm joint}=(1-\kappa)A_{\min}.                            \tag{44.14}
\]
A signed finite-block evaluation of (41.8) may replace the separate safe
stocks (44.12)--(44.13).

\begin{proposition}[Uniform chord-response row bound]
\label{prop:v18v34-row}
Under the V33 spectral closure hypotheses, suppose the compiled response
decomposition obeys
\[
 |(K_c(s))_{ef}|
 \leq\alpha_*\mathbf1_{\{e=f\}}\mathcal U_{e;c}
       +\mathfrak r_*\mathcal T_{e;fc}                        \tag{44.15}
\]
for \(0\leq s\leq1\), with bounded support enlargements absorbed into
\(\alpha_*,\mathfrak r_*\).  Then
\[
 \left\|\sum_cK_c(s)\right\|_{2\to2}
 \leq\mathcal L_{\rm cyc}                                    \tag{44.16}
\]
uniformly in \(s\) and the regulator, where
\[
 \mathcal L_{\rm cyc}
 :=\alpha_*M_{\rm in}M_C
   +\mathfrak r_*M_{\rm in}M_EM_C.                            \tag{44.17}
\]
\end{proposition}

\begin{proof}
For a fixed row \(e\), sum (44.15) over \(f,c\), then apply Lemma
\ref{lem:v18v34-arms}.  The resulting maximum absolute row sum is
\(\mathcal L_{\rm cyc}\).  The matrix \(\sum_cK_c(s)\) is symmetric, so
its Euclidean operator norm is at most the same quantity.
\end{proof}

\subsection{The full cyclic conformal margin}

Let \(\mu_0>0\) be the cut-reference conformal Hessian margin.  Under the
coefficient-only tree theorem of V25 one may take
\[
 \mu_0=(1-\mathfrak Q_0)p_0                                  \tag{44.18}
\]
when (35.19)--(35.20) hold, after duplicating each cut chord parameter on
its two terminal stubs and restricting the positive doubled Hessian to the
diagonal shared-parameter subspace.

\begin{theorem}[Full cycle-restoration margin]
\label{thm:v18v34-cycle-margin}
Assume:
\begin{enumerate}[leftmargin=2.2em]
\item the cut reference has \(H_0\succeq\mu_0I\);
\item the normalized hopping certificate (43.15) holds with
      \(\kappa<1\);
\item the uniform quartic moment ball of V32 and the compiled local response
      bounds (44.12)--(44.15) hold;
\item the resolvent stocks (44.2)--(44.3) are uniformly finite;
\item
      \[
      \mathcal L_{\rm cyc}<\mu_0.                            \tag{44.19}
      \]
\end{enumerate}
Then the full cyclic conformal Hessian satisfies
\[
 \boxed{\quad
 H_1\succeq(\mu_0-\mathcal L_{\rm cyc})I>0
 \quad}                                                       \tag{44.20}
\]
uniformly in the regulator volume.
\end{theorem}

\begin{proof}
Theorems \ref{thm:v18v33-closure} and
\ref{thm:v18v31-chord} justify the response decomposition throughout
\(0\leq s\leq1\).  Proposition \ref{prop:v18v34-row} gives
\[
 \left\|\frac{dH_s}{ds}\right\|_{2\to2}
 =\left\|\sum_cK_c(s)\right\|_{2\to2}
 \leq\mathcal L_{\rm cyc}.                                   \tag{44.21}
\]
Integrate (36.12):
\[
 H_1\succeq H_0-int_0^1
 \left\|\frac{dH_s}{ds}\right\|_{2\to2}I\,ds
 \succeq(\mu_0-\mathcal L_{\rm cyc})I.                      \tag{44.22}
\]
Equation (44.19) proves strict positivity.
\end{proof}

\begin{corollary}[Explicit coarse cycle card]
\label{cor:v18v34-card}
If weights satisfy (44.10), chord attachments have uniform weighted density
\(M_C\), and the V32 stocks are inserted in (44.12)--(44.13), then
(44.19) is a finite inequality involving only
\[
 p_0,\ \mathfrak Q_0,\ \kappa,\ A_{\min},\
 M_{3,*},\ a^*,\ b^*,\ \chi^*,\ M_C,\ w_*/w^*.              \tag{44.23}
\]
Every entry is either a coefficient-box endpoint, a one-dimensional radial
moment, or a finite graph-overlap stock.
\end{corollary}

\begin{proof}
Use (44.11)--(44.14), the gradient bounds (42.16)--(42.17), and the margin
(44.18) in (44.17)--(44.19).
\end{proof}

\subsection{Position of the cycle sector}

Theorem \ref{thm:v18v34-cycle-margin} completes the analytic passage from a
tree conformal Hessian to a cyclic one under explicit strict gates.  It does
not assert that the physical Standard Model coefficient card passes them.
The source manuscripts still leave the relevant positive Higgs coefficients
as physical data rather than a certified compact numerical box.  The next
upstream task is therefore a matching theorem: express every stock in
(44.23) in the canonical lattice parameters and determine which are fixed by
the continuum renormalization conditions rather than chosen freely.

\begin{assessment}[V34 acquisition and boundary]
\label{ass:v18v34}
V34 proves volume-independent summability of the two- and three-arm
Dobrushin responses and the full cycle-restoration margin (44.20).  This
closes the cyclic conformal-Hessian module conditional on explicit finite
coefficient, spectral, overlap, and response inequalities.  Those strict
inequalities have not been verified on a physical continuum card.  Gauge
and fermion coupling, cutoff removal, reconstruction, and the coupled
Einstein continuum remain open.
\end{assessment}

\begin{record}[V34 append record]
\label{rec:v18v34}
V34 is appended to the validated V33 whose installed SHA--256 was
\texttt{5FD523F838BA0DD09502D1CBAA110846A06EFA2178A49D4671F636F3BDB4D436}.
After installation only V34 is the active numeric SM note; the frozen
\texttt{V100\_f17} archive is unchanged.
\end{record}

% =====================================================================
% V35 APPENDIX: SOURCE AUDIT AND CYCLIC GAUSSIANIZATION SEPARATOR
% =====================================================================

\section{V35: physical-source audit and the cyclic Gaussianization separator}
\label{sec:v18v35}

V34 closes the finite-regulator cyclic conformal-Hessian module under
explicit strict inequalities.  Before treating that module as a continuum
interaction mechanism, the attached physical-source manuscripts were read
again at their newest available versions.  The audit changes the next
direction.

The audited sources were:
\[
\begin{array}{ll}
\text{Grand Tensor V067:}&
\texttt{FACD058EC8BB210AD8B296C9D7F4D78A229D08D7A2AE378808700655CB8BC0FE},\\
\text{physical source metric V35:}&
\texttt{AE1FBC0A8125C37465E38E7649E2992725F4579A86F0BB17483B9AE829CFB17C},\\
\text{uniform interacting continuum V53:}&
\texttt{C5932EA71AE54F440F721309505B1EA386B6CC2D1627CDE5D13AC1F6F18CBFD7}.
\end{array}                                                     \tag{45.1}
\]
The first two sources still do not supply a canonical numerical compact box
for the positive Higgs coefficients entering (44.23).  The third source does
supply an exact theorem relevant to interpretation: on its selected
summably coupled elementary-Higgs shell branch, every scalar normalization
is either zero, Gaussian, or non-tight.  It explicitly requires a separately
populated renormalized composite, a scale-critical replacement law, a
thermodynamic branch, or another protected operator sector for non-Gaussian
interaction survival.

This note proves the corresponding scaling mechanism for any four-dimensional
cyclic regulator satisfying uniform connected-cumulant summability.  It is
not a repetition of the source theorem: the hypotheses are stated directly
in terms of the volume-uniform cycle-response bounds developed in V26--V34.

\subsection{Four-dimensional cumulant counting}

Let \(\Lambda_n\) be a four-dimensional regulator of mesh \(h_n\downarrow0\)
in a fixed compact volume, with
\[
 |\Lambda_n|\leq C_\Lambda h_n^{-4}.                           \tag{45.2}
\]
Let \(X_{n,z}\) be centered scalar components of a local Higgs coordinate.
For \(r\geq2\), assume the connected cumulants obey
\[
 \sup_{n}\sup_{z_1\in\Lambda_n}
 \sum_{z_2,\ldots,z_r\in\Lambda_n}
 \left|operatorname{Cum}
 (X_{n,z_1},\ldots,X_{n,z_r})\right|
 \leq K_r<\infty.                                             \tag{45.3}
\]
For a bounded test array \(f=(f_z)\), define the critical pairing
\[
 Z_n(f)=h_n^2\sum_{z\in\Lambda_n}f_zX_{n,z}.                  \tag{45.4}
\]

\begin{theorem}[Connected-cumulant scaling]
\label{thm:v18v35-cumulant-scaling}
Under (45.2)--(45.3),
\[
 \left|operatorname{Cum}
 (Z_n(f_1),\ldots,Z_n(f_r))\right|
 \leq C_\Lambda K_r
 h_n^{2r-4}\prod_{j=1}^r\|f_j\|_\infty.                     \tag{45.5}
\]
In particular, every fixed connected critical cumulant of order
\(r\geq3\) tends to zero.
\end{theorem}

\begin{proof}
Multilinearity of cumulants gives
\[
 \operatorname{Cum}(Z_n(f_1),\ldots,Z_n(f_r))
 =h_n^{2r}\sum_{z_1,\ldots,z_r}
 \left(\prod_{j=1}^rf_{j,z_j}\right)
 \operatorname{Cum}(X_{n,z_1},\ldots,X_{n,z_r}).              \tag{45.6}
\]
Take absolute values, fix \(z_1\), apply (45.3) to the remaining sums, and
then use (45.2).  This gives (45.5).  Since \(2r-4>0\) for \(r\geq3\), the
last assertion follows.
\end{proof}

For vector-valued Higgs fields, apply the theorem to arbitrary fixed real
linear combinations of components.  Cram\'er--Wold then gives the joint
statement.

\begin{corollary}[Cubic interaction vanishes under the V34 cycle gate]
\label{cor:v18v35-cubic}
Suppose the local third-gradient stocks of V32, the uniform Dobrushin
certificate of V33, and the summed three-arm response stock of V34 hold on a
four-dimensional regulator family.  Then (45.3) holds for \(r=3\), and
\[
 \operatorname{Cum}(Z_n(f_1),Z_n(f_2),Z_n(f_3))=O(h_n^2).      \tag{45.7}
\]
Thus the elementary critical field cannot retain a nonzero cubic connected
continuum vertex inside this uniformly summable branch.
\end{corollary}

\begin{proof}
The compiled local third-cumulant stock is finite by V31--V32.  The uniform
three-arm resolvent sum in V34 makes the sum over the two relative insertion
positions finite independently of \(n\), which is (45.3) for \(r=3\).
Apply (45.5).
\end{proof}

\subsection{Gaussian limit under all-order summability}

\begin{theorem}[Cyclic elementary-field Gaussianization]
\label{thm:v18v35-gaussianization}
Assume (45.3) for every fixed \(r\geq2\), uniform local moments of every
fixed order, and convergence of the critical covariance forms
\[
 \operatorname{Cov}(Z_n(f),Z_n(g))\longrightarrow\mathcal Q(f,g). \tag{45.8}
\]
Then every finite collection \((Z_n(f_1),\ldots,Z_n(f_m))\) converges in
law to the centered Gaussian vector with covariance \(\mathcal Q\).
If, in addition, the usual covariance-weighted Sobolev trace is uniformly
finite, the fields are tight in \(H^{-s}\), \(s>2\), and converge there to
the corresponding Gaussian distribution.
\end{theorem}

\begin{proof}
Theorem \ref{thm:v18v35-cumulant-scaling} makes every joint cumulant of order
at least three tend to zero, while (45.8) supplies the limiting second
cumulants and centering supplies the first.  The moment--cumulant partition
formula therefore makes every fixed joint moment converge to its Wick value.
The assumed local moment bounds and (45.3) give uniform bounds for each fixed
moment of the critical pairings, hence tightness of every finite-dimensional
family.  Every subsequential limit has the Gaussian moment sequence; that
sequence is determinate, so the limit is the stated Gaussian vector.

For the distributional statement, expand in a Laplace eigenbasis.  The
uniform covariance-weighted trace bounds the expected \(H^{-s_0}\)-norm for
some \(s_0>2\).  Compact embedding into \(H^{-s}\), \(s>s_0\), gives
tightness, and the finite-dimensional limit identifies every subsequential
limit.
\end{proof}

\subsection{The scalar-normalization trichotomy}

Define the unscaled elementary distribution by
\[
 \Phi_n(f)=h_n^4\sum_zf_zX_{n,z},
 \qquad Z_n(f)=h_n^{-2}\Phi_n(f).                             \tag{45.9}
\]
Let \(a_n>0\) be any scalar amplitude and put
\[
 b_n=a_nh_n^2,
 \qquad a_n\Phi_n(f)=b_nZ_n(f).                               \tag{45.10}
\]

\begin{corollary}[Cyclic scalar-scaling trichotomy]
\label{cor:v18v35-trichotomy}
Assume Theorem \ref{thm:v18v35-gaussianization} and choose one test \(f\)
with \(\mathcal Q(f,f)>0\).  Along every subsequence on which \(b_n\) has a
limit in \([0,\infty]\), exactly one of the following holds:
\begin{enumerate}[leftmargin=2.2em]
\item if \(b_n\to0\), then \(a_n\Phi_n\to0\) in probability;
\item if \(b_n\to b\in(0,\infty)\), then
      \(a_n\Phi_n\) converges to the Gaussian field \(b\mathcal W_{\mathcal Q}\);
\item if \(b_n\to\infty\), then the pairing
      \(a_n\Phi_n(f)\) is not tight.
\end{enumerate}
No scalar normalization yields a nonzero non-Gaussian elementary continuum
inside this branch.
\end{corollary}

\begin{proof}
Equation (45.10) and Slutsky's theorem prove the first two cases.  In the
third, \(Z_n(f)\) converges to a nondegenerate Gaussian.  For every \(M>0\),
\[
 \mathbb P(|b_nZ_n(f)|\leq M)
 =\mathbb P(|Z_n(f)|\leq M/b_n)\longrightarrow0,              \tag{45.11}
\]
by convergence in distribution and continuity of the Gaussian law at zero.
Thus the pairings escape every compact interval.  The cases exhaust all
subsequential limits of \(b_n\).
\end{proof}

\subsection{Upstream consequence for the SM--Einstein programme}

The finite-regulator convexity theorem of V34 and the continuum interaction
problem are logically different.  If the V33--V34 response bounds extend to
all cumulant orders with uniform summability, they place the scalar-rescaled
elementary Higgs field in the Gaussian trichotomy above.  Stronger uniform
mixing would therefore stabilize the regulator while simultaneously ruling
out elementary non-Gaussian interaction survival.

The programme must now branch explicitly:
\begin{enumerate}[leftmargin=2.2em]
\item retain V34 as the finite-regulator stability and Ward-convexity module;
\item populate a renormalized enhanced coordinate or composite generated by
      the common action, with its subtraction, normalization, and source
      ancestry fixed before testing non-Gaussianity; or
\item replace the summably coupled shell law by a scale-critical interacting
      trajectory and prove tightness by a method which does not assume
      all-order cumulant summability.
\end{enumerate}
The attached V53 manuscript names the first alternative as its active bulk
request.  V36 should therefore start from the literal quartic action score
and determine which enhanced combinations escape the additive
Gaussianization theorem without being post hoc observables.

\begin{assessment}[V35 acquisition and boundary]
\label{ass:v18v35}
V35 records the physical-source audit, proves four-dimensional cumulant
scaling for cyclic regulators, proves vanishing of the critical cubic
cumulant under the current cycle gate, and proves the complete scalar
Gaussianization trichotomy under all-order summability.  It shows that V34
alone cannot supply an interacting elementary continuum.  An occurring
renormalized enhanced coordinate or a scale-critical replacement law remains
unpopulated.  Gauge/fermion coupling, cutoff removal, reconstruction, and the
coupled Einstein continuum remain open.
\end{assessment}

\begin{record}[V35 append record]
\label{rec:v18v35}
V35 is appended to the validated V34 whose installed SHA--256 was
\texttt{FC9C27ADF2806BA65AC69948BD9414CA3131B6EEFB57FB8C74F13C7F6F6EF808}.
After installation only V35 is the active numeric SM note; the frozen
\texttt{V100\_f17} archive is unchanged.
\end{record}

\section{V36: compiled local-score banks and the interaction-survival threshold}
\label{sec:v18v36}

\subsection{Mandatory companion audit and choice of direction}

After V35, the newest active companion notes are
\[
\begin{array}{c|c|c}
\hbox{programme} & \hbox{active file} & \hbox{SHA--256}\\ \hline
\hbox{Yang--Mills}
&
\texttt{Renewal\_Yang\_Mills\_Mass\_Gap\_Research\_Notes\_V5.tex}
&
\texttt{53D900C8A36175A6702FA4D1AC2D65DD748D7C7988BC389967CD1E16E008D89A}
\\
\hbox{Navier--Stokes}
&
\texttt{Renewal\_NS\_Stability\_Research\_Notes\_V26.tex}
&
\texttt{82C887F8C2C69E4F28FAD7C3DC8DE5B9099CB5A4BF431EBA1FFF3E91935978AD}.
\end{array}                                                     \tag{46.1}
\]
The Yang--Mills V5 note contains a useful methodological correction.  Its
seven response words form one exact second derivative of a compiled scalar,
and large direct and rank--three contributions may cancel.  Therefore one
must contract the complete response word before applying absolute-value
bounds.  This is not imported as an SM theorem; below it is implemented as
the definition of the SM cumulant stock.  The Navier--Stokes V26 rank--one
refinement gives a modest gain in some regimes and no gain in others, so it
does not alter the present scaling obstruction.

The attached interacting-continuum manuscript also makes the literal
quartic action score an unavoidable first candidate.  That candidate is
generated by the common action and hence is not post hoc.  V36 determines
whether a finite bank containing it can escape the V35 Gaussianization
mechanism.

\subsection{The compiled common-action bank}

Let \(m<\infty\) be fixed.  At every regulator cell \(z\in\Lambda_n\), let
\[
 {\bf Y}_{n,z}
 =
 \bigl(Y_{n,z}^{1},\ldots,Y_{n,z}^{m}\bigr)                   \tag{46.2}
\]
be a centered bank of local observables generated by derivatives of the
same regulated action.  It may contain elementary field components,
centered mass and quartic scores, plaquette scores, and finitely many
collision-core counterterms.  For a deterministic direction
\(u_n\in{\mathbb R}^{m}\), compile first:
\[
 X_{n,z}[u_n]=u_n\mathbin{\cdot}{\bf Y}_{n,z}.                \tag{46.3}
\]
The relevant connected stock is the scalar quantity
\[
 {\cal K}_{r,n}[u_n]
 =
 \sup_{z_1\in\Lambda_n}
 \sum_{z_2,\ldots,z_r\in\Lambda_n}
 \left|
 \operatorname{Cum}\bigl(
 X_{n,z_1}[u_n],\ldots,X_{n,z_r}[u_n]\bigr)
 \right|.                                                     \tag{46.4}
\]
Definition (46.4), rather than a termwise norm of the uncontracted tensor,
preserves every cancellation which is fixed by the common action.

\begin{lemma}[Componentwise stock is only a sufficient majorant]
\label{lem:v18v36-component-majorant}
Suppose
\[
 {\cal B}_{r,n}
 =
 \max_{1\leq i_1,\ldots,i_r\leq m}
 \sup_{z_1}
 \sum_{z_2,\ldots,z_r}
 \left|
 \operatorname{Cum}\bigl(
 Y_{n,z_1}^{i_1},\ldots,Y_{n,z_r}^{i_r}\bigr)
 \right|<\infty.                                             \tag{46.5}
\]
Then
\[
 {\cal K}_{r,n}[u_n]\leq
 \|u_n\|_1^r{\cal B}_{r,n}.                                  \tag{46.6}
\]
The converse need not hold.
\end{lemma}

\begin{proof}
Expand the cumulant in (46.4) by multilinearity and apply the triangle
inequality after the complete contraction.  Summing the absolute
coefficients gives \(\|u_n\|_1^r\), proving (46.6).  The converse fails
whenever distinct component words cancel after contraction; no lower bound
on (46.5) follows from the contracted scalar.
\end{proof}

\subsection{Finite-bank Gaussianization}

Assume as in V35 that
\[
 |\Lambda_n|\leq C_\Lambda h_n^{-4},\qquad h_n\downarrow0,    \tag{46.7}
\]
and define the critical compiled pairing
\[
 {\cal Z}_n^{u}(f)
 =
 h_n^2\sum_{z\in\Lambda_n}f_zX_{n,z}[u_n].                   \tag{46.8}
\]

\begin{theorem}[Compiled finite-bank cumulant bound]
\label{thm:v18v36-compiled-bound}
For every \(r\geq2\) and bounded test arrays \(f_1,\ldots,f_r\),
\[
\begin{split}
 &\left|
 \operatorname{Cum}\bigl(
 {\cal Z}_n^{u}(f_1),\ldots,{\cal Z}_n^{u}(f_r)\bigr)
 \right|                                                        \\
 &\hspace{2em}\leq
 C_\Lambda h_n^{2r-4}{\cal K}_{r,n}[u_n]
 \prod_{j=1}^{r}\|f_j\|_\infty .                              \tag{46.9}
\end{split}
\]
Consequently, if
\(\sup_n{\cal K}_{r,n}[u_n]<\infty\) for every fixed \(r\geq2\),
all connected cumulants of order \(r\geq3\) vanish.
\end{theorem}

\begin{proof}
By multilinearity, the left side before taking absolute values equals
\[
 h_n^{2r}
 \sum_{z_1,\ldots,z_r}
 \left(\prod_{j=1}^{r}f_{j,z_j}\right)
 \operatorname{Cum}\bigl(
 X_{n,z_1}[u_n],\ldots,X_{n,z_r}[u_n]\bigr).                 \tag{46.10}
\]
Bound the tests by their sup norms, fix \(z_1\), and use (46.4) for all
remaining sums.  The final \(z_1\)-sum is bounded by (46.7).  This proves
(46.9).  For \(r\geq3\), the remaining power \(h_n^{2r-4}\) tends to zero.
\end{proof}

\begin{corollary}[No-go theorem for a fixed local score bank]
\label{cor:v18v36-finite-bank-nogo}
Assume the hypotheses of Theorem
\ref{thm:v18v36-compiled-bound}, uniform moments of every fixed order, and
convergence of the covariance forms of (46.8).  Then every
finite-dimensional limit of \({\cal Z}_n^{u}\) is Gaussian.  The same
conclusion holds uniformly for directions satisfying
\(\sup_n\|u_n\|_1<\infty\) if
\(\sup_n{\cal B}_{r,n}<\infty\) for each \(r\).

Thus adjoining any fixed finite list of summably correlated local
common-action scores does not by itself produce an interacting scalar
continuum.
\end{corollary}

\begin{proof}
Theorem \ref{thm:v18v36-compiled-bound} kills every joint connected
cumulant of order at least three.  Covariance convergence fixes the second
cumulants, centering fixes the first, and the moment--cumulant formula gives
the Wick moments.  Uniform moments give finite-dimensional tightness, and
the Gaussian moment sequence is determinate.  The uniform-direction
statement follows from Lemma \ref{lem:v18v36-component-majorant}.
\end{proof}

This corollary extends the V35 obstruction from an elementary field to an
entire finite local bank.  It also explains why exact compilation is
essential: a componentwise estimate may fail to prove summability even when
the scalar score selected by the action has a small contracted stock.

\subsection{The literal quartic score does not escape}

For the four-real-component pure quartic one-cell reference law, write
\[
 R_z=\|H_z\|^4,\qquad Y_z=R_z-1.                              \tag{46.11}
\]
The radial change of variables recorded in the attached
interacting-continuum manuscript gives \(R_z\sim\operatorname{Exp}(1)\).
Hence
\[
 \operatorname{Cum}_1(Y_z)=0,\qquad
 \operatorname{Cum}_r(Y_z)=(r-1)!\quad(r\geq2).               \tag{46.12}
\]
The score is an exact derivative of the local quartic action, so it is an
occurring common-action observable.

\begin{proposition}[One-cell non-Gaussianity is erased by four-dimensional
critical summation]
\label{prop:v18v36-quartic-score}
For independent cells with the law (46.11), set
\[
 Q_n(f)=h_n^2\sum_{z\in\Lambda_n}f_zY_z.                     \tag{46.13}
\]
Then
\[
 \operatorname{Cum}\bigl(Q_n(f_1),\ldots,Q_n(f_r)\bigr)
 =
 (r-1)!h_n^{2r}\sum_{z\in\Lambda_n}
 \prod_{j=1}^{r}f_{j,z},                                    \tag{46.14}
\]
and therefore
\[
 \left|\operatorname{Cum}_3(Q_n)\right|=O(h_n^2),
 \qquad
 \left|\operatorname{Cum}_4(Q_n)\right|=O(h_n^4).             \tag{46.15}
\]
The same conclusion holds for an interacting regulator whenever the
contracted quartic-score stocks (46.4) are uniformly summable.
\end{proposition}

\begin{proof}
Independence makes a joint cumulant vanish unless all cell labels agree.
The surviving diagonal cumulant is (46.12), which proves (46.14).
Equation (46.7) then yields the two rates in (46.15).  In the interacting
case, use Theorem \ref{thm:v18v36-compiled-bound}.
\end{proof}

Thus a non-Gaussian one-cell law and common-action ancestry are both
insufficient.  A fixed local polynomial score remains inside the same
central-limit basin as the elementary field when its connected correlations
are summable.

\subsection{A quantitative interaction-survival threshold}

The obstruction above can be reversed into a necessary diagnostic.  Define
\[
 \Xi_{r,n}[u_n]
 =
 h_n^{2r-4}{\cal K}_{r,n}[u_n].                              \tag{46.16}
\]

\begin{theorem}[Necessary loss of cumulant summability]
\label{thm:v18v36-survival-threshold}
Fix \(r\geq3\), tests with \(\|f_j\|_\infty\leq1\), and amplitudes
\(0<b_-\leq b_n\leq b_+<\infty\).  If
\[
 \limsup_{n\to\infty}
 \left|
 \operatorname{Cum}\bigl(
 b_n{\cal Z}_n^{u}(f_1),\ldots,
 b_n{\cal Z}_n^{u}(f_r)\bigr)
 \right|\geq c>0,                                           \tag{46.17}
\]
then along a subsequence
\[
 {\cal K}_{r,n}[u_n]
 \geq
 \frac{c}{2C_\Lambda b_+^r}\,h_n^{-(2r-4)},                 \tag{46.18}
\]
or equivalently
\[
 \limsup_{n\to\infty}\Xi_{r,n}[u_n]>0.                       \tag{46.19}
\]
In particular, survival of a cubic connected vertex requires a contracted
three-point stock of order at least \(h_n^{-2}\), while survival of a
quartic connected vertex requires order at least \(h_n^{-4}\).
\end{theorem}

\begin{proof}
Apply (46.9) to the variables multiplied by \(b_n\).  The upper bound is
\[
 C_\Lambda b_n^r h_n^{2r-4}{\cal K}_{r,n}[u_n].
\]
Condition (46.17) supplies a subsequence on which the left side is at least
\(c/2\).  Since \(b_n\leq b_+\), rearrangement gives (46.18), and (46.19)
is its dimensionless form.
\end{proof}

\begin{corollary}[Finite local grammar cannot cross the threshold]
\label{cor:v18v36-finite-grammar}
Suppose an enhanced coordinate is an additive sum over cells of a fixed
finite collection of uniformly local common-action words, with bounded
compiled coefficients.  If their connected cumulants have a uniform
anchor-sum bound, then \(\Xi_{r,n}\to0\) for every \(r\geq3\).  Such a
coordinate cannot have a nonzero connected continuum vertex.
\end{corollary}

\begin{proof}
Uniform locality and the assumed anchor-sum bound give
\({\cal K}_{r,n}=O(1)\).  Equation (46.16) then tends to zero, contradicting
the necessary condition (46.19).
\end{proof}

\subsection{Revised continuum branch}

The next construction must therefore be specified before another stability
estimate is attempted.  A viable common-action enhancement must cross
(46.19) while retaining a finite nonzero covariance and distributional
tightness.  There are only three structural ways for the present proof to
fail:
\begin{enumerate}[leftmargin=2.2em]
\item the physical correlation range grows so that the contracted
      anchor-sum stock reaches the rate (46.18);
\item the coordinate uses a scale-dependent or nonlocal renormalized
      grammar whose number of effective words grows with the cutoff; or
\item the regulator action itself follows a scale-critical trajectory, so
      the V33--V34 uniform summability gate is replaced by a critical
      tightness argument.
\end{enumerate}
This is a sharper target than asking merely for an enhanced observable.
The literal quartic score, any bounded finite linear combination of local
action scores, and any other fixed finite-range additive polynomial are now
excluded under the existing summability hypotheses.  The next note should
construct the smallest scale-dependent common-action coordinate, calculate
its covariance normalization, and test its contracted \(\Xi_{3,n}\) and
\(\Xi_{4,n}\) before importing gauge or fermion sectors.

\begin{assessment}[V36 acquisition and boundary]
\label{ass:v18v36}
V36 records the mandatory YM/NS audit, imports its compilation lesson without
importing a companion theorem, and proves a finite-bank Gaussianization
no-go theorem.  It verifies explicitly that the literal quartic action score
does not evade that theorem.  It then proves the quantitative survival
threshold \({\cal K}_{r,n}\gtrsim h_n^{-(2r-4)}\) for a nonzero connected
\(r\)-vertex.  No coordinate meeting that threshold with finite covariance
and tightness has yet been constructed.  Gauge/fermion coupling, cutoff
removal, reconstruction, and the coupled Einstein continuum remain open.
\end{assessment}

\begin{record}[V36 append record]
\label{rec:v18v36}
V36 is appended to the validated V35 whose installed SHA--256 was
\texttt{DAB63E06CF12C5384F54ECDD1B769DE878BFCDEFC69F9DC8E65EF46258A1D53C}.
After installation only V36 is the active numeric SM note; the frozen
\texttt{V100\_f17} archive is unchanged.
\end{record}


\section{V37: mesoscopic reblocking, bounded filters, and the false-enhancement obstruction}
\label{sec:v18v37}

\subsection{Source check and the next unclosed case}

The attached uniform interacting-continuum V53 manuscript already proves
that the literal quartic score and its first quadratic-composite partner have
a joint Gaussian white-noise landing, and that their first direct-minus-Wick
excess is a diagonal contact.  Those results are not repeated here.  V36
strengthened the obstruction to a finite compiled bank.  The remaining
elementary proposal is to let the observable itself depend on the scale by
averaging over mesoscopic cells.  V37 tests that proposal exactly.

\subsection{Bounded deterministic filters cannot create interaction}

Let \(X_{n,z}\) be a centered scalar common-action score and define its
unfiltered critical pairing by
\[
 Z_n(f)=h_n^2\sum_{z\in\Lambda_n}f_zX_{n,z}.                 \tag{47.1}
\]
For a deterministic kernel \(k_n(x,z)\), put
\[
 ({\cal K}_n^*f)_z=\sum_{x\in\Lambda_n}k_n(x,z)f_x,
 \qquad
 L_n=\sup_z\sum_x|k_n(x,z)|,                                 \tag{47.2}
\]
and define the filtered pairing
\[
\begin{split}
 Z_n^{{\cal K}}(f)
 &=h_n^2\sum_x f_x\sum_zk_n(x,z)X_{n,z}\\
 &=Z_n({\cal K}_n^*f).                                      \tag{47.3}
\end{split}
\]

\begin{theorem}[Quasilocal-filter cumulant bound]
\label{thm:v18v37-filter}
Let
\[
 K_{r,n}
 =
 \sup_{z_1}\sum_{z_2,\ldots,z_r}
 \left|
 \operatorname{Cum}(X_{n,z_1},\ldots,X_{n,z_r})
 \right|.                                                    \tag{47.4}
\]
Under \(|\Lambda_n|\leq C_\Lambda h_n^{-4}\),
\[
\begin{split}
 &\left|
 \operatorname{Cum}\bigl(
 Z_n^{\cal K}(f_1),\ldots,Z_n^{\cal K}(f_r)\bigr)
 \right|\\
 &\hspace{2em}\leq
 C_\Lambda h_n^{2r-4}K_{r,n}L_n^r
 \prod_{j=1}^r\|f_j\|_\infty .                              \tag{47.5}
\end{split}
\]
In particular, a uniformly bounded filter,
\(\sup_nL_n<\infty\), preserves the V36 Gaussianization conclusion whenever
\(\sup_nK_{r,n}<\infty\).
\end{theorem}

\begin{proof}
Equation (47.3) rewrites the filtered field as the old field evaluated on a
new test array.  By (47.2),
\[
 \|{\cal K}_n^*f\|_\infty\leq L_n\|f\|_\infty.               \tag{47.6}
\]
Apply the connected-cumulant estimate to the \(r\) transformed tests and use
(47.6).
\end{proof}

\begin{corollary}[Required filter amplification]
\label{cor:v18v37-filter-amplification}
If, for some \(r\geq3\) and tests of sup norm at most one,
\[
 \limsup_n\left|
 \operatorname{Cum}\bigl(
 Z_n^{\cal K}(f_1),\ldots,Z_n^{\cal K}(f_r)\bigr)
 \right|\geq c>0,                                           \tag{47.7}
\]
then along a subsequence
\[
 L_n^rK_{r,n}
 \geq\frac{c}{2C_\Lambda}h_n^{-(2r-4)}.                     \tag{47.8}
\]
Thus a summably correlated parent score requires
\[
 L_n\gtrsim h_n^{-(2-4/r)}.                                 \tag{47.9}
\]
\end{corollary}

\begin{proof}
Rearrange (47.5) on a subsequence on which the left side is at least
\(c/2\).  If \(K_{r,n}=O(1)\), take the \(r\)-th root.
\end{proof}

This separates quasilocal smearing from a genuine renormalized enhancement.
A positive normalized average has \(L_n\leq1\); it cannot satisfy (47.9).
A signed filter can have large \(L_n\), but then its covariance and topology
must be renormalized and checked independently.

\subsection{Exact invariance under mesoscopic reblocking}

Partition \(\Lambda_n\) into disjoint blocks \(B\), each containing
\(N_n\) cells.  Let
\[
 \delta_n=h_nN_n^{1/4}                                      \tag{47.10}
\]
be the four-dimensional block mesh, and define the variance-scale block
variable
\[
 U_{n,B}=N_n^{-1/2}\sum_{z\in B}X_{n,z}.                    \tag{47.11}
\]
For a test array \(g=(g_B)\) on blocks, form
\[
 {\cal U}_n(g)=\delta_n^2\sum_Bg_BU_{n,B}.                   \tag{47.12}
\]

\begin{theorem}[Reblocking identity]
\label{thm:v18v37-reblocking}
For every block test \(g\),
\[
 \boxed{
 {\cal U}_n(g)
 =
 h_n^2\sum_{z\in\Lambda_n}g_{B(z)}X_{n,z}.}                 \tag{47.13}
\]
Consequently, if \(\delta_n\to0\) and the blockwise constant tests
\(g_{B(z)}\) approximate a continuum test \(f\) uniformly, mesoscopic
reblocking has exactly the same finite-dimensional limit as \(Z_n(f)\).
\end{theorem}

\begin{proof}
By (47.10),
\[
 \delta_n^2N_n^{-1/2}
 =h_n^2N_n^{1/2}N_n^{-1/2}=h_n^2.                            \tag{47.14}
\]
Insert (47.11) into (47.12), interchange the two finite sums, and use
(47.14).  This gives (47.13).  Uniform approximation of the test arrays then
identifies the limits whenever the critical pairings are tight.
\end{proof}

\begin{corollary}[Averaging scale is not a new normalization]
\label{cor:v18v37-no-new-normalization}
Choosing \(N_n\to\infty\) with \(\delta_n\to0\) does not create a new
continuum branch.  The factor \(N_n^{-1/2}\) in the block variable and the
factor \(N_n^{1/2}\) in the block critical normalization cancel identically.
\end{corollary}

\subsection{Nonlinear block scores still central-limit}

One might instead apply a fixed nonlinear function after forming a block
variable.  Let \(W_{n,B}\) be any centered common-action block score, possibly
a polynomial of \(U_{n,B}\), and set
\[
 {\cal W}_n(g)=\delta_n^2\sum_{B\in{\cal P}_n}g_BW_{n,B}.     \tag{47.15}
\]
Assume
\[
 |{\cal P}_n|\leq C_{\cal P}\delta_n^{-4}                    \tag{47.16}
\]
and define the block anchored stock
\[
 J_{r,n}
 =
 \sup_{B_1}
 \sum_{B_2,\ldots,B_r}
 \left|
 \operatorname{Cum}(W_{n,B_1},\ldots,W_{n,B_r})
 \right|.                                                    \tag{47.17}
\]

\begin{theorem}[Mesoscopic nonlinear-score dichotomy]
\label{thm:v18v37-block-score}
For every \(r\geq2\),
\[
\begin{split}
 &\left|
 \operatorname{Cum}\bigl(
 {\cal W}_n(g_1),\ldots,{\cal W}_n(g_r)\bigr)
 \right|\\
 &\hspace{2em}\leq
 C_{\cal P}\delta_n^{2r-4}J_{r,n}
 \prod_{j=1}^r\|g_j\|_\infty.                               \tag{47.18}
\end{split}
\]
If \(\delta_n\to0\) and \(J_{r,n}=O(1)\) for every fixed \(r\), every
subsequential limit with convergent covariance and uniform moments is
Gaussian.  If a connected \(r\)-vertex survives, then necessarily
\[
 J_{r,n}\gtrsim\delta_n^{-(2r-4)}.                           \tag{47.19}
\]
\end{theorem}

\begin{proof}
Expand the cumulant of (47.15), fix the first block, use (47.17), and then
sum over at most \(C_{\cal P}\delta_n^{-4}\) anchors.  This proves
(47.18).  The Gaussian statement follows from the moment--cumulant formula,
as in V36.  Rearranging (47.18) proves (47.19).
\end{proof}

\begin{corollary}[Mesoscopic false-enhancement trichotomy]
\label{cor:v18v37-mesoscopic-trichotomy}
A fixed nonlinear score of normalized blocks has only the following
possibilities:
\begin{enumerate}[leftmargin=2.2em]
\item if \(\delta_n\to0\) and the block stocks are uniformly summable, its
      nonzero tight limit is Gaussian;
\item if \(\delta_n\not\to0\), the construction retains a positive physical
      block size and does not define a local continuum field on arbitrarily
      fine tests;
\item if a non-Gaussian local limit survives, at least one block stock loses
      summability at the quantitative rate (47.19).
\end{enumerate}
\end{corollary}

\begin{proof}
The first and third alternatives are Theorem
\ref{thm:v18v37-block-score}.  In the second, there is an
\(\varepsilon>0\) and a subsequence on which the partition mesh is at least
\(\varepsilon\).  Tests varying below that scale cannot be approximated by
blockwise constants, so local continuum resolution fails.
\end{proof}

\subsection{Consequence for the enhanced SM--Einstein packet}

Neither a normalized moving average, ordinary block-spin regrouping, nor a
fixed polynomial applied independently to shrinking blocks crosses the V36
interaction threshold.  The smallest remaining source-native candidate must
therefore contain a scale-dependent \emph{spectral} mode or a genuinely
long-range connected kernel.  It is no longer enough to enlarge the local
polynomial bank.

The next exact calculation will use a quadratic common-action coordinate.
For a Gaussian reference vector \(G_n\) and a symmetric action-derived
kernel \(A_n\), the centered quadratic score
\[
 G_n^{\mathsf T}A_nG_n-\operatorname{tr}(A_n\operatorname{Cov}G_n)
\]
has cumulants determined by the eigenvalues of the whitened kernel.  This
will decide whether a scale-dependent quadratic packet Gaussianizes by
growing effective rank or retains a non-Gaussian mode because one
source-visible eigenvalue remains macroscopic.

\begin{assessment}[V37 acquisition and boundary]
\label{ass:v18v37}
V37 proves that every uniformly bounded deterministic filter preserves the
V36 Gaussianization obstruction, gives the exact four-dimensional
mesoscopic reblocking identity, and proves the same obstruction for fixed
nonlinear scores on shrinking blocks.  It reduces a viable enhancement to a
diverging connected stock or a macroscopic spectral mode.  No such physical
mode, its action ancestry, or its coupled Einstein transport has yet been
populated.
\end{assessment}

\begin{record}[V37 append record]
\label{rec:v18v37}
V37 is appended to the validated V36 whose installed SHA--256 was
\texttt{BEB5A6E5D55D309F176AE63185096FE4840D6B5E4B76EE7791D44B389DACA8C9}.
After installation only V37 is the active numeric SM note; the frozen
\texttt{V100\_f17} archive is unchanged.
\end{record}


\section{V38: quadratic action scores, effective rank, and the first spectral survivor criterion}
\label{sec:v18v38}

\subsection{The scale-dependent quadratic candidate}

V37 excludes bounded linear filters and fixed nonlinear scores on shrinking
independent or summably coupled blocks.  A scale-dependent quadratic form is
different because its kernel can carry a coherent eigenmode across many
cells.  This section gives an exact finite-dimensional criterion for that
mode.

Let \(G_n\) be a centered real Gaussian vector in a finite-dimensional space
\({\cal H}_n\), with covariance \(C_n\succeq0\).  Let \(A_n=A_n^*\) be a
deterministic symmetric kernel occurring as a second-order coefficient in
the regulated action.  On the support of \(C_n\), define the whitened kernel
\[
 B_n=C_n^{1/2}A_nC_n^{1/2}.                                  \tag{48.1}
\]
The centered quadratic action score is
\[
 Q_n(A_n)
 =
 \langle G_n,A_nG_n\rangle-\operatorname{tr}(A_nC_n).         \tag{48.2}
\]
After writing \(G_n=C_n^{1/2}g_n\) with \(g_n\) standard Gaussian, (48.2)
has the same law as
\[
 \langle g_n,B_ng_n\rangle-\operatorname{tr}B_n.              \tag{48.3}
\]

\subsection{Exact cumulant compiler}

\begin{theorem}[Quadratic Gaussian cumulants]
\label{thm:v18v38-quadratic-cumulants}
For every integer \(r\geq2\),
\[
 \boxed{
 \operatorname{Cum}_r\bigl(Q_n(A_n)\bigr)
 =
 2^{r-1}(r-1)!\operatorname{tr}(B_n^r).}                     \tag{48.4}
\]
In particular,
\[
 \operatorname{Var}Q_n(A_n)=2\operatorname{tr}(B_n^2).       \tag{48.5}
\]
\end{theorem}

\begin{proof}
Diagonalize \(B_n\), with eigenvalues \(\lambda_{n,j}\).  For
\(|t|<(2\|B_n\|_{\rm op})^{-1}\),
\[
\begin{split}
 \log\mathbb E e^{tQ_n}
 &=-t\operatorname{tr}B_n
   -\frac12\log\det(I-2tB_n)\\
 &=\sum_{r=2}^{\infty}
   \frac{2^{r-1}}{r}t^r\operatorname{tr}(B_n^r).              \tag{48.6}
\end{split}
\]
The coefficient of \(t^r\) is
\(\operatorname{Cum}_r(Q_n)/r!\), which proves (48.4).
\end{proof}

Assume \(\operatorname{tr}(B_n^2)>0\) and normalize
\[
 W_n
 =
 \frac{Q_n(A_n)}
 {\sqrt{2\operatorname{tr}(B_n^2)}}.                         \tag{48.7}
\]
Then
\[
 \operatorname{Var}W_n=1,
 \qquad
 \operatorname{Cum}_4(W_n)
 =
 12\frac{\operatorname{tr}(B_n^4)}
 {(\operatorname{tr}(B_n^2))^2}.                            \tag{48.8}
\]
Define
\[
 \rho_n
 =
 \frac{\|B_n\|_{\rm op}}
 {\sqrt{\operatorname{tr}(B_n^2)}},
 \qquad
 r_{\rm eff}(B_n)
 =
 \frac{(\operatorname{tr}(B_n^2))^2}
 {\operatorname{tr}(B_n^4)}.                                \tag{48.9}
\]
The second quantity is the quadratic effective rank.

\begin{theorem}[Spectral Gaussianization equivalence]
\label{thm:v18v38-spectral-equivalence}
The following are equivalent:
\begin{enumerate}[leftmargin=2.2em]
\item \(W_n\) converges in law to \(N(0,1)\);
\item \(\operatorname{Cum}_4(W_n)\to0\);
\item \(r_{\rm eff}(B_n)\to\infty\);
\item \(\rho_n\to0\).
\end{enumerate}
If these conditions hold, every cumulant of \(W_n\) of order at least three
converges to zero.
\end{theorem}

\begin{proof}
Let \(\lambda_{n,j}\) be the eigenvalues of \(B_n\), and put
\[
 a_{n,j}
 =
 \frac{\lambda_{n,j}}
 {\sqrt{2\sum_k\lambda_{n,k}^2}}.
\]
Then
\[
 W_n=\sum_ja_{n,j}(g_{n,j}^2-1),\qquad
 2\sum_ja_{n,j}^2=1.                                       \tag{48.10}
\]
The condition \(\rho_n\to0\) is precisely the infinitesimal-coefficient
condition for this independent triangular array.  Since
\(g^2-1\) has moments of every order, the Lindeberg theorem gives
\(W_n\Rightarrow N(0,1)\).

Equations (48.8)--(48.9) make the equivalence of the second and third
conditions immediate.  Moreover,
\[
 \rho_n^4
 \leq
 \frac{\operatorname{tr}(B_n^4)}
 {(\operatorname{tr}(B_n^2))^2}
 \leq
 \rho_n^2.                                                   \tag{48.11}
\]
Thus the third and fourth conditions are equivalent.

If \(W_n\Rightarrow N(0,1)\), Gaussian-chaos hypercontractivity gives a
uniform moment bound of some order strictly above four.  Hence fourth powers
are uniformly integrable, the fourth moments converge to \(3\), and (48.8)
gives the second condition.  Finally, from (48.4),
\[
 \left|\operatorname{Cum}_r(W_n)\right|
 \leq c_r\rho_n^{r-2}\qquad(r\geq3),                         \tag{48.12}
\]
because
\(
 |\operatorname{tr}(B_n^r)|
 \leq\|B_n\|_{\rm op}^{r-2}\operatorname{tr}(B_n^2)
\).
\end{proof}

\subsection{The macroscopic-eigenmode criterion}

\begin{corollary}[Necessary and sufficient spectral survivor]
\label{cor:v18v38-macroscopic-mode}
A variance-normalized quadratic Gaussian action score fails to Gaussianize
exactly when, along some subsequence, at least one eigenvalue obeys
\[
 |\lambda_{n,j_n}|
 \geq c\left(\sum_j\lambda_{n,j}^2\right)^{1/2}              \tag{48.13}
\]
for some \(c>0\).  Equivalently, its quadratic effective rank does not
diverge.
\end{corollary}

\begin{proof}
The left side of (48.13), maximized over \(j\), divided by the right-hand
Hilbert--Schmidt scale is \(\rho_n\).  Apply Theorem
\ref{thm:v18v38-spectral-equivalence}.
\end{proof}

\begin{example}[Rank-one survivor]
\label{ex:v18v38-rank-one}
If \(B_n=\lambda_nv_nv_n^{\mathsf T}\) with \(\|v_n\|=1\) and
\(\lambda_n\ne0\), then \(r_{\rm eff}(B_n)=1\) and
\[
 W_n
 =
 \frac{\operatorname{sgn}(\lambda_n)}{\sqrt2}
 \bigl(\langle g_n,v_n\rangle^2-1\bigr),                    \tag{48.14}
\]
so \(\operatorname{Cum}_4(W_n)=12\).  The mode is non-Gaussian at every
cutoff.  This is a spectral possibility, not yet a physical SM coordinate:
the vector \(v_n\), kernel \(A_n\), source occurrence, and cutoff transport
must all be supplied by the common action.
\end{example}

\begin{corollary}[Diffuse local quadratic packets Gaussianize]
\label{cor:v18v38-diffuse-packets}
Suppose \(B_n\) has rank tending to infinity, uniformly bounded eigenvalues,
and
\[
 \operatorname{tr}(B_n^2)\longrightarrow\infty.              \tag{48.15}
\]
Then the variance-normalized score is Gaussian in the limit.  In particular,
a diagonal or finite-range block kernel with order-one weights on
\(\asymp h_n^{-4}\) independent modes has
\[
 r_{\rm eff}(B_n)\asymp h_n^{-4},
 \qquad
 \operatorname{Cum}_4(W_n)=O(h_n^4).                         \tag{48.16}
\]
\end{corollary}

\begin{proof}
The hypotheses give
\(\rho_n\leq\sup_n\|B_n\|_{\rm op}/\sqrt{\operatorname{tr}(B_n^2)}
\to0\).  The first claim follows from Theorem
\ref{thm:v18v38-spectral-equivalence}.  For comparable order-one local
weights on \(\asymp h_n^{-4}\) modes, both
\(\operatorname{tr}(B_n^2)\) and \(\operatorname{tr}(B_n^4)\) are of that
order.  Equations (48.8)--(48.9) yield (48.16).
\end{proof}

\subsection{Joint quadratic packets}

For a finite test bank \(f_1,\ldots,f_\ell\), suppose the action supplies
symmetric kernels \(A_n(f_j)\), linear in the test.  Write
\[
 B_n(c)=C_n^{1/2}
 \left(\sum_{j=1}^{\ell}c_jA_n(f_j)\right)C_n^{1/2}.          \tag{48.17}
\]

\begin{theorem}[Cram\'er--Wold spectral criterion]
\label{thm:v18v38-joint}
Assume the covariance matrices of the centered quadratic scores converge.
If
\[
 \frac{\|B_n(c)\|_{\rm op}}
 {\sqrt{\operatorname{tr}(B_n(c)^2)}}\longrightarrow0        \tag{48.18}
\]
for every \(c\in{\mathbb R}^{\ell}\) for which the limiting variance is
nonzero, then the score vector converges jointly to a centered Gaussian
vector.  Conversely, a non-Gaussian joint limit with uniformly nondegenerate
second-chaos normalization forces (48.18) to fail for at least one fixed
linear combination.
\end{theorem}

\begin{proof}
Every linear combination of the score vector is the quadratic score with
kernel \(\sum_jc_jA_n(f_j)\).  Apply Theorem
\ref{thm:v18v38-spectral-equivalence} and the Cram\'er--Wold theorem.  If
(48.18) held for every fixed combination, all combinations would have
Gaussian limits, which characterizes a Gaussian vector.
\end{proof}

\subsection{Relation to the V36 cumulant threshold}

For a quadratic score, V36's abstract survival stock has now acquired a
spectral mechanism.  Diffuse eigenvalues create many weak independent chaos
coordinates and force the normalized fourth cumulant to zero.  A
macroscopic eigenvalue prevents that dilution.  In particular,
\[
 \boxed{
 \text{quadratic interaction survival}
 \ \Longrightarrow\
 \limsup_n
 \frac{\|C_n^{1/2}A_nC_n^{1/2}\|_{\rm op}^2}
 {\operatorname{tr}((C_n^{1/2}A_nC_n^{1/2})^2)}>0.}          \tag{48.19}
\]
This condition is invariant under orthogonal changes of regulator
coordinates and therefore is not a choice of local polynomial basis.

The condition also exposes the next ancestry question.  A macroscopic
quadratic eigenmode may already be a deterministic follower of the
primitive source and physical reflected clock.  Only its component outside
the complete source-cyclic follower range is eligible to define a new
enhanced coordinate.  V39 will perform that dynamic short directly on the
quadratic kernel and determine which spectral mass remains.

\begin{assessment}[V38 acquisition and boundary]
\label{ass:v18v38}
V38 supplies the exact cumulants of a Gaussian quadratic action score and
proves an if-and-only-if Gaussianization criterion in terms of effective
rank.  It identifies a macroscopic eigenvalue of the whitened action kernel
as the first precise spectral mechanism able to evade the V36--V37 central
limit obstruction.  No claim is made that the regulated SM action contains
such a mode or that it remains after source-cyclic shortening.
\end{assessment}

\begin{record}[V38 append record]
\label{rec:v18v38}
V38 is appended to the validated V37 whose installed SHA--256 was
\texttt{C4440040CEB657EC349EB85C8B9DA60CB0B6E091B93ACB69D1A3009E4FEE54D4}.
After installation only V38 is the active numeric SM note; the frozen
\texttt{V100\_f17} archive is unchanged.
\end{record}


\section{V39: the dynamic quadratic short and the source-new spectral residue}
\label{sec:v18v39}

\subsection{Why the linear short must be lifted}

V38 identifies a macroscopic eigenmode of a whitened quadratic kernel as a
possible non-Gaussian survivor.  That criterion alone does not establish a
new physical coordinate.  The mode may be generated entirely by the
primitive source under the completed reflected clock.  The attached
continuum programme already defines the corresponding first-chaos
source-cyclic carrier.  V39 lifts that carrier to the quadratic Gaussian
chaos and performs the minimum follower short there.

Let \({\cal H}_n\) be the real whitened one-particle space, \(H_n\geq0\) the
completed reflected Hamiltonian, and
\[
 S_n:E_n\longrightarrow{\cal H}_n                               \tag{49.1}
\]
the primitive source synthesis.  Let
\[
 {\cal C}_n
 =
 \overline{\operatorname{span}}
 \{e^{-tH_n}S_nu:t\geq0,\ u\in E_n\},\qquad
 P_n=P_{{\cal C}_n},\qquad R_n=I-P_n.                         \tag{49.2}
\]
Only the definition (49.2) is imported; the linear source ancestry theorems
are not reproved.

\subsection{Second-chaos geometry}

For a symmetric Hilbert--Schmidt operator \(B\) on \({\cal H}_n\), write
\[
 I_2(B)=\langle g,Bg\rangle-\operatorname{tr}B,               \tag{49.3}
\]
where \(g\) is standard Gaussian.  The quadratic-chaos covariance is
\[
 \mathbb E[I_2(A)I_2(B)]
 =
 2\operatorname{tr}(AB).                                     \tag{49.4}
\]
Thus symmetric Hilbert--Schmidt kernels, with inner product
\[
 \langle A,B\rangle_{(2)}
 =
 2\operatorname{tr}(AB),                                     \tag{49.5}
\]
are an isometric model of the second Gaussian chaos.

\begin{definition}[Minimum primitive quadratic follower space]
\label{def:v18v39-quadratic-follower}
The primitive quadratic follower space is
\[
 {\cal F}_n^{(2)}
 =
 \{P_nAP_n:A=A^*,\ A\ {\rm Hilbert--Schmidt}\}.               \tag{49.6}
\]
It is the closed span of centered products of two first-chaos coordinates
from \({\cal C}_n\).  No mixed primitive--residual product is included unless
the physical source grammar separately declares it.
\end{definition}

\begin{theorem}[Exact quadratic follower projection]
\label{thm:v18v39-quadratic-projection}
The orthogonal projection in (49.5) onto \({\cal F}_n^{(2)}\) is
\[
 \Pi_n^{(2)}B=P_nBP_n.                                       \tag{49.7}
\]
Consequently the dynamically shorted quadratic kernel is
\[
 \boxed{
 B_n^\perp=B_n-P_nB_nP_n,}                                   \tag{49.8}
\]
and its residual variance is
\[
 \boxed{
 {\cal R}_n^{(2)}
 =
 \mathbb E[I_2(B_n^\perp)^2]
 =
 2\|B_n^\perp\|_{\rm HS}^2.}                                 \tag{49.9}
\]
\end{theorem}

\begin{proof}
For any symmetric Hilbert--Schmidt \(A\),
\[
 \operatorname{tr}\bigl((B-PBP)PAP\bigr)
 =
 \operatorname{tr}(PBPAP)-\operatorname{tr}(PBPAP)=0.        \tag{49.10}
\]
Thus \(B-PBP\) is orthogonal to (49.6), while \(PBP\) lies in (49.6).
This proves (49.7)--(49.8).  Equation (49.9) follows from (49.4).
\end{proof}

Decomposing relative to
\({\cal H}_n={\cal C}_n\oplus{\cal C}_n^\perp\) gives
\[
\begin{split}
 B_n
 &=
 P_nB_nP_n
 +(P_nB_nR_n+R_nB_nP_n)
 +R_nB_nR_n,                                                 \tag{49.11}\\
 B_n^\perp
 &=
 (P_nB_nR_n+R_nB_nP_n)+R_nB_nR_n.                           \tag{49.12}
\end{split}
\]

\begin{proposition}[Orthogonal variance ledger]
\label{prop:v18v39-variance-ledger}
The three kernels in (49.11) are pairwise orthogonal in the second-chaos
inner product, and
\[
\begin{split}
 \|B_n\|_{\rm HS}^2
 &=
 \|P_nB_nP_n\|_{\rm HS}^2
 +2\|P_nB_nR_n\|_{\rm HS}^2
 +\|R_nB_nR_n\|_{\rm HS}^2,                                 \tag{49.13}\\
 {\cal R}_n^{(2)}
 &=
 4\|P_nB_nR_n\|_{\rm HS}^2
 +2\|R_nB_nR_n\|_{\rm HS}^2.                                \tag{49.14}
\end{split}
\]
\end{proposition}

\begin{proof}
Every mixed trace contains a factor \(P_nR_n=0\).  Also
\[
 \|PBR+RBP\|_{\rm HS}^2=2\|PBR\|_{\rm HS}^2,                 \tag{49.15}
\]
because the two off-diagonal blocks are adjoints and orthogonal.
Equations (49.13)--(49.14) follow from (49.9).
\end{proof}

\begin{firewall}[Follower grammar]
\label{fire:v18v39-follower-grammar}
The pure follower block \(P_nB_nP_n\) is removed because it is generated by
products of two declared primitive clock followers.  The mixed block
\(P_nB_nR_n+R_nB_nP_n\) is retained: it contains a residual leg and is not a
function of the primitive first-chaos carrier alone.  If a larger physical
grammar declares mixed products as followers, that larger synthesis must be
written explicitly and projected before assigning a new source price.
\end{firewall}

\subsection{Spectral survival after the dynamic short}

Assume \(\|B_n^\perp\|_{\rm HS}>0\) and define
\[
 W_n^\perp
 =
 \frac{I_2(B_n^\perp)}
 {\sqrt2\|B_n^\perp\|_{\rm HS}},
 \qquad
 \rho_n^\perp
 =
 \frac{\|B_n^\perp\|_{\rm op}}
 {\|B_n^\perp\|_{\rm HS}}.                                  \tag{49.16}
\]

\begin{theorem}[Source-new quadratic survivor criterion]
\label{thm:v18v39-source-new-survivor}
The dynamically shorted, variance-normalized quadratic score
\(W_n^\perp\) converges to \(N(0,1)\) if and only if
\[
 \rho_n^\perp\longrightarrow0.                               \tag{49.17}
\]
Its normalized fourth cumulant is
\[
 \boxed{
 \operatorname{Cum}_4(W_n^\perp)
 =
 12
 \frac{\operatorname{tr}((B_n^\perp)^4)}
 {\|B_n^\perp\|_{\rm HS}^4}.}                                \tag{49.18}
\]
Hence a non-Gaussian quadratic continuum coordinate is eligible as
source-new only if both
\[
 \liminf_n{\cal R}_n^{(2)}>0
 \quad\hbox{and}\quad
 \limsup_n\rho_n^\perp>0                                    \tag{49.19}
\]
hold after transport to common physical units.
\end{theorem}

\begin{proof}
Apply the V38 spectral equivalence theorem to the symmetric kernel
\(B_n^\perp\).  Formula (49.18) is its exact fourth-cumulant formula.
The first condition in (49.19) prevents the source-new quadratic variance
from collapsing; the second prevents Gaussianization of its normalized
shape.
\end{proof}

\begin{corollary}[Three distinct quadratic outcomes]
\label{cor:v18v39-three-outcomes}
After the minimum dynamic short, exactly one of the following occurs along
each convergent subsequence:
\begin{enumerate}[leftmargin=2.2em]
\item \({\cal R}_n^{(2)}\to0\): the quadratic score is a collapsing follower
      modulo the primitive quadratic grammar;
\item \(\liminf{\cal R}_n^{(2)}>0\) and
      \(\rho_n^\perp\to0\): a source-new quadratic variance survives, but its
      normalized law is Gaussian;
\item \(\liminf{\cal R}_n^{(2)}>0\) and
      \(\limsup\rho_n^\perp>0\): a non-Gaussian source-new spectral branch is
      possible and requires a retained macroscopic residual eigenmode.
\end{enumerate}
\end{corollary}

\subsection{Exact rank-one ancestry calculation}

Let
\[
 B_n=\lambda_nv_nv_n^{\mathsf T},
 \qquad \|v_n\|=1,\qquad
 p_n=P_nv_n,\qquad q_n=R_nv_n,                               \tag{49.20}
\]
and write \(\beta_n=\|q_n\|^2\).

\begin{proposition}[Rank-one follower leakage]
\label{prop:v18v39-rank-one-leakage}
For (49.20),
\[
 \boxed{
 \|B_n^\perp\|_{\rm HS}^2
 =
 \lambda_n^2(2\beta_n-\beta_n^2),}                           \tag{49.21}
\]
and therefore
\[
 \boxed{
 {\cal R}_n^{(2)}
 =
 2\lambda_n^2(2\beta_n-\beta_n^2).}                          \tag{49.22}
\]
If \(\beta_n=0\), the entire macroscopic quadratic mode is a primitive
follower.  If \(|\lambda_n|\geq\lambda_*>0\) and
\(\beta_n\geq\beta_*>0\), a positive residual variance survives and
\(B_n^\perp\) has rank at most two; hence its effective rank is uniformly
bounded and the normalized residual cannot Gaussianize.
\end{proposition}

\begin{proof}
Since \(P_nv_n=p_n\),
\[
 B_n^\perp
 =
 \lambda_n(v_nv_n^{\mathsf T}-p_np_n^{\mathsf T}).           \tag{49.23}
\]
Now
\[
\begin{split}
 \lambda_n^{-2}\|B_n^\perp\|_{\rm HS}^2
 &=
 \|v_nv_n^{\mathsf T}\|_{\rm HS}^2
 +\|p_np_n^{\mathsf T}\|_{\rm HS}^2
 -2\operatorname{tr}
   (v_nv_n^{\mathsf T}p_np_n^{\mathsf T})\\
 &=1+\|p_n\|^4-2\|p_n\|^4
 =1-(1-\beta_n)^2,
\end{split}                                                   \tag{49.24}
\]
which is (49.21).  Equation (49.22) follows from (49.9).
The rank assertion follows from (49.23), and a uniformly bounded effective
rank contradicts the V38 Gaussianization criterion.
\end{proof}

\subsection{Cutoff stability of the quadratic short}

To compare two cutoffs on one transported Hilbert space, let \(B,B'\) be
symmetric Hilbert--Schmidt kernels and \(P,P'\) orthogonal projections.  Put
\[
 {\cal S}(B,P)=B-PBP.                                        \tag{49.25}
\]

\begin{lemma}[Perturbation bound for the lifted short]
\label{lem:v18v39-short-perturbation}
One has
\[
\boxed{
 \|{\cal S}(B',P')-{\cal S}(B,P)\|_{\rm HS}
 \leq
 2\|B'-B\|_{\rm HS}
 +2\|P'-P\|_{\rm op}\|B\|_{\rm HS}.}                         \tag{49.26}
\]
\end{lemma}

\begin{proof}
Insert and subtract \(P'BP'\):
\[
\begin{split}
 {\cal S}(B',P')-{\cal S}(B,P)
 ={}&(B'-B)-P'(B'-B)P'\\
 &-(P'-P)BP'-PB(P'-P).                                      \tag{49.27}
\end{split}
\]
Compression does not increase the Hilbert--Schmidt norm.  The two remaining
terms are bounded by
\(\|P'-P\|_{\rm op}\|B\|_{\rm HS}\), proving (49.26).
\end{proof}

\begin{corollary}[Cofinal stability criterion]
\label{cor:v18v39-cofinal-stability}
Suppose transported kernels and source-cyclic projections satisfy
\[
 \|B_{n+1}-B_n\|_{\rm HS}
 +\|P_{n+1}-P_n\|_{\rm op}\|B_n\|_{\rm HS}
 \longrightarrow0.                                         \tag{49.28}
\]
Then \(B_n^\perp\) is Cauchy whenever the corresponding defect series is
summable.  If additionally
\(\inf_n\|B_n^\perp\|_{\rm HS}>0\), the normalized residual kernels and their
fourth-cumulant ratios are stable under the same transport.
\end{corollary}

\begin{proof}
Apply (49.26) on adjacent cutoffs and sum.  A positive lower
Hilbert--Schmidt norm makes normalization locally Lipschitz.  Operator norm
and traces of second and fourth powers are continuous under
Hilbert--Schmidt convergence in finite dimension on each transported
support, so the ratio in (49.18) is stable.
\end{proof}

\subsection{The first admissible enhanced packet}

The smallest remaining quadratic packet is now explicit.  It must contain:
\begin{enumerate}[leftmargin=2.2em]
\item an action-derived symmetric kernel \(A_n\) and physical covariance
      \(C_n\), giving the whitened \(B_n=C_n^{1/2}A_nC_n^{1/2}\);
\item the completed reflected clock \(H_n\) and primitive source synthesis
      \(S_n\), fixing \(P_n\) before the quadratic terminal is inspected;
\item the residual kernel \(B_n^\perp=B_n-P_nB_nP_n\);
\item a positive transported variance floor for (49.9);
\item a positive macroscopic-mode floor for (49.18), or the honest Gaussian
      outcome; and
\item the common-action metric variation that identifies the residual
      quadratic score as a stress or current component rather than an
      arbitrarily selected spectral statistic.
\end{enumerate}

The last item is the immediate bottleneck.  V40 will derive the exact metric
source-response identities, including their seagull/contact terms, so the
quadratic residual is tested on separated supports before it is admitted as
an Einstein source.

\begin{assessment}[V39 acquisition and boundary]
\label{ass:v18v39}
V39 constructs the minimum second-chaos lift of the primitive source-cyclic
follower space, proves the exact follower projection and orthogonal variance
ledger, and gives the source-new spectral survival criterion.  It also
computes rank-one leakage exactly and proves a cutoff perturbation bound.
No literal SM metric-response kernel has yet been inserted into this
compiler, so no physical non-Gaussian stress coordinate is claimed.
\end{assessment}

\begin{record}[V39 append record]
\label{rec:v18v39}
V39 is appended to the validated V38 whose installed SHA--256 was
\texttt{8F2DBC10411E37E65B4591A7B9E359DBCD9698BB949C9BBA4EE75AE69ABA2905}.
After installation only V39 is the active numeric SM note; the frozen
\texttt{V100\_f17} archive is unchanged.
\end{record}


\section{V40: common-action metric responses, contact separation, and the first spectral stress witness}
\label{sec:v18v40}

\subsection{Metric-source setup}

Let \(q_1,\ldots,q_r\) be compactly supported real symmetric metric
variations on the regulated four-dimensional geometry, and let
\[
 g({\bf s})=g+\sum_{i=1}^rs_iq_i.                            \tag{50.1}
\]
Write \(S_n({\bf s},\phi)\) for the complete regulated
Einstein--Standard-Model action, including every determinant, gauge-fixing,
ghost, and counterterm contribution admitted by the chosen finite
regulator.  Assume differentiation under the finite-dimensional integral is
justified up to the order used below.  Define
\[
 Z_n({\bf s})
 =
 \int e^{-S_n({\bf s},\phi)}\,d\nu_n(\phi),
 \qquad
 W_n({\bf s})=-\log Z_n({\bf s}).                            \tag{50.2}
\]
For each nonempty \(B\subseteq\{1,\ldots,r\}\), put
\[
 S_{n,B}
 =
 \left.
 \prod_{i\in B}\partial_{s_i}S_n({\bf s},\phi)
 \right|_{{\bf s}=0}.                                       \tag{50.3}
\]
The singleton \(S_{n,\{i\}}\) is the literal integrated stress score in
direction \(q_i\).  Terms \(S_{n,B}\) with \(|B|\geq2\) are metric seagulls
or contact insertions fixed by the same action.

\subsection{The complete response partition formula}

\begin{theorem}[Common-action response compiler]
\label{thm:v18v40-response-compiler}
Let \({\cal P}([r])\) be the set of set partitions of
\([r]=\{1,\ldots,r\}\).  Then
\[
\boxed{
 \left.
 \partial_{s_1}\cdots\partial_{s_r}W_n({\bf s})
 \right|_{{\bf s}=0}
 =
 \sum_{\pi\in{\cal P}([r])}
 (-1)^{|\pi|+1}
 \operatorname{Cum}_n\bigl(S_{n,B}:B\in\pi\bigr).}           \tag{50.4}
\]
Here a one-entry cumulant is expectation.  Thus the direct connected stress
response and all contact terms are compiled before any norm or quotient is
taken.
\end{theorem}

\begin{proof}
For observables \(A_1,\ldots,A_k\) depending differentiably on a parameter
\(s\), differentiation of the tilted cumulant generating function gives
\[
 \partial_s\operatorname{Cum}_s(A_1,\ldots,A_k)
 =
 \sum_{j=1}^k
 \operatorname{Cum}_s(A_1,\ldots,\partial_sA_j,\ldots,A_k)
 -
 \operatorname{Cum}_s(A_1,\ldots,A_k,\partial_sS_n).         \tag{50.5}
\]
For \(r=1\), (50.4) is
\(\partial_{s_1}W_n=\mathbb E S_{n,\{1\}}\).
Assume (50.4) at order \(r\) and differentiate in \(s_{r+1}\).
The first sum in (50.5) inserts \(r+1\) into exactly one existing block of
each partition; the last term adds the singleton block \(\{r+1\}\) and
reverses the sign.  Every partition of \([r+1]\) is produced exactly once,
with coefficient \((-1)^{|\pi|+1}\).  Induction proves (50.4).
\end{proof}

At second and third order, (50.4) reads
\[
 \partial_{12}W_n
 =
 \mathbb ES_{n,12}
 -
 \operatorname{Cov}(S_{n,1},S_{n,2}),                        \tag{50.6}
\]
and
\[
\begin{split}
 \partial_{123}W_n
 ={}&
 \mathbb ES_{n,123}
 -\operatorname{Cov}(S_{n,12},S_{n,3})
 -\operatorname{Cov}(S_{n,13},S_{n,2})\\
 &-\operatorname{Cov}(S_{n,23},S_{n,1})
 +\operatorname{Cum}(S_{n,1},S_{n,2},S_{n,3}).              \tag{50.7}
\end{split}
\]

\subsection{Separated supports remove the metric contacts}

\begin{definition}[Metric locality at the regulator]
\label{def:v18v40-metric-locality}
The action is metric-local to order \(r\) if, for every
\(B\subseteq[r]\) with \(|B|\geq2\),
\[
 \operatorname{supp}S_{n,B}
 \subseteq
 \bigcap_{i\in B}\operatorname{supp}q_i                     \tag{50.8}
\]
up to already declared cell-boundary contact collars which shrink to the
same intersection in physical units.
\end{definition}

\begin{theorem}[Separated stress-response identity]
\label{thm:v18v40-separated-response}
Assume metric locality to order \(r\), and suppose the supports of
\(q_1,\ldots,q_r\) are pairwise disjoint with positive physical separation.
At every sufficiently fine cutoff,
\[
\boxed{
 \partial_{s_1}\cdots\partial_{s_r}W_n(0)
 =
 (-1)^{r+1}
 \operatorname{Cum}_n(S_{n,1},\ldots,S_{n,r}).}              \tag{50.9}
\]
In particular,
\[
 \partial_{123}W_n(0)
 =
 \operatorname{Cum}_n(S_{n,1},S_{n,2},S_{n,3}).              \tag{50.10}
\]
\end{theorem}

\begin{proof}
Positive separation makes every intersection in (50.8) empty after the
contact collars become sufficiently small.  Hence \(S_{n,B}=0\) whenever
\(|B|\geq2\).  In (50.4), only the partition into \(r\) singleton blocks
remains, and its sign is \((-1)^{r+1}\).
\end{proof}

\begin{corollary}[Seagull/contact-scheme invariance]
\label{cor:v18v40-contact-invariance}
Changing the representation of the mixed metric derivatives
\(S_{n,B}\), \(|B|\geq2\), by terms supported on collision diagonals leaves
(50.9) unchanged on pairwise positively separated tests.  Therefore a
nonzero limit of (50.10) cannot be manufactured or removed by a seagull or
contact-basis reparametrization.  The singleton stress score
\(S_{n,\{i\}}\) remains fixed by the chosen common action.
\end{corollary}

\begin{proof}
Every mixed variation supported on a collision diagonal vanishes when its
test supports have empty intersection.  Apply Theorem
\ref{thm:v18v40-separated-response}.
\end{proof}

\subsection{Gaussian quadratic prediction on separated supports}

Suppose the centered linear metric scores have quadratic Gaussian
representatives
\[
 S_{n,i}^{\rm G}
 =
 \langle g_n,B_{n,i}g_n\rangle-\operatorname{tr}B_{n,i},      \tag{50.11}
\]
where \(g_n\) is standard Gaussian and \(B_{n,i}=B_{n,i}^*\) is the
whitened common-action kernel.

\begin{lemma}[Joint quadratic cumulant]
\label{lem:v18v40-joint-quadratic}
For symmetric matrices \(B_1,\ldots,B_r\),
\[
\boxed{
 \operatorname{Cum}\bigl(I_2(B_1),\ldots,I_2(B_r)\bigr)
 =
 \frac{2^{r-1}}{r}
 \sum_{\sigma\in{\mathfrak S}_r}
 \operatorname{tr}
 (B_{\sigma(1)}\cdots B_{\sigma(r)}).}                       \tag{50.12}
\]
For \(r=3\),
\[
\boxed{
 \operatorname{Cum}\bigl(I_2(B_1),I_2(B_2),I_2(B_3)\bigr)
 =
 4\operatorname{tr}(B_1B_2B_3+B_1B_3B_2).}                  \tag{50.13}
\]
\end{lemma}

\begin{proof}
Apply the scalar identity
\[
 \operatorname{Cum}_r I_2\left(\sum_it_iB_i\right)
 =
 2^{r-1}(r-1)!
 \operatorname{tr}\left(\sum_it_iB_i\right)^r               \tag{50.14}
\]
and extract the coefficient of \(t_1\cdots t_r\).
Multilinearity supplies \(r!\) copies of the symmetric joint cumulant on the
left; ordered products on the right are indexed by
\({\mathfrak S}_r\).  Division by \(r!\) gives (50.12).
For symmetric matrices, cyclicity and transposition divide the six
third-order words into the two terms in (50.13), each with multiplicity
three.
\end{proof}

Define the unshorted Gaussian prediction
\[
 {\cal G}_{3,n}^{\rm raw}
 =
 4\operatorname{tr}\left(
 B_{n,1}B_{n,2}B_{n,3}
 +
 B_{n,1}B_{n,3}B_{n,2}
 \right).                                                    \tag{50.15}
\]

\begin{definition}[Separated raw direct-minus-Gaussian stress excess]
\label{def:v18v40-raw-stress-excess}
On pairwise separated metric tests, define
\[
\boxed{
 {\cal E}_{3,n}^{\rm raw}(q_1,q_2,q_3)
 =
 \partial_{123}W_n(0)-{\cal G}_{3,n}^{\rm raw}.}             \tag{50.16}
\]
\end{definition}

\begin{theorem}[Source-native raw cubic stress witness]
\label{thm:v18v40-raw-stress-witness}
Assume the metric variations are pairwise positively separated, the direct
action and Gaussian comparator use the same physical covariance, source
frame, and metric units, and both terms in (50.16) have cofinal limits.
Then \(\lim_n{\cal E}_{3,n}^{\rm raw}\) is independent of the
seagull/contact representation.  If
\[
 \left|\lim_n{\cal E}_{3,n}^{\rm raw}(q_1,q_2,q_3)\right|>0, \tag{50.17}
\]
the unshorted stress packet is not quasi-free on these tests.
\end{theorem}

\begin{proof}
Theorem \ref{thm:v18v40-separated-response} identifies the first term in
(50.16) with the direct connected three-point stress tensor.  Lemma
\ref{lem:v18v40-joint-quadratic} identifies the second with the complete
matched Gaussian prediction.  Corollary
\ref{cor:v18v40-contact-invariance} gives the stated invariance.  A
quasi-free quadratic packet obeys (50.12), so a nonzero excess contradicts
quasi-freeness.
\end{proof}

\subsection{The same follower short on both sides}

The raw witness must still pass the V39 ancestry short.  Let
\({\mathfrak P}_n^{(2)}\) be the orthogonal projection in
\(L_0^2(\mu_n)\) onto the declared primitive quadratic follower space, and
put
\[
 T_{n,i}=S_{n,i}-\mathbb ES_{n,i},
 \qquad
 T_{n,i}^{\perp}
 =(I-{\mathfrak P}_n^{(2)})T_{n,i}.                           \tag{50.18}
\]
For the Gaussian comparator, let \(P_n\) be the primitive dynamic
source-cyclic projection of V39 and put
\[
 B_{n,i}^{\perp}=B_{n,i}-P_nB_{n,i}P_n.                      \tag{50.19}
\]
Its shorted third cumulant is
\[
 {\cal G}_{3,n}^{\perp}
 =
 4\operatorname{tr}\left(
 B_{n,1}^{\perp}B_{n,2}^{\perp}B_{n,3}^{\perp}
 +
 B_{n,1}^{\perp}B_{n,3}^{\perp}B_{n,2}^{\perp}
 \right).                                                    \tag{50.20}
\]

\begin{definition}[Dynamically eligible stress excess]
\label{def:v18v40-stress-excess}
\[
\boxed{
 {\cal E}_{3,n}^{\perp}(q_1,q_2,q_3)
 =
 \operatorname{Cum}_n
 (T_{n,1}^{\perp},T_{n,2}^{\perp},T_{n,3}^{\perp})
 -
 {\cal G}_{3,n}^{\perp}.}                                    \tag{50.21}
\]
\end{definition}

\begin{theorem}[Source-new cubic stress witness]
\label{thm:v18v40-stress-witness}
Assume the direct and Gaussian follower spaces in (50.18)--(50.19) are the
two realizations of the same predeclared primitive quadratic grammar, and
all terms in (50.21) have cofinal limits.  If
\[
 \left|\lim_n{\cal E}_{3,n}^{\perp}(q_1,q_2,q_3)\right|>0,   \tag{50.22}
\]
the dynamically residual stress packet is not quasi-free on these tests.
\end{theorem}

\begin{proof}
The first term is the connected three-point tensor after orthogonal removal
of the declared direct follower grammar.  By V39 and Lemma
\ref{lem:v18v40-joint-quadratic}, the second term is precisely the matched
Gaussian tensor after removal of the same grammar.  Equality is compulsory
for a quasi-free realization, so (50.22) excludes one.
\end{proof}

\subsection{Interaction-threshold consequence}

Suppose the separated residual score insertions can be written as critical
pairings of a local stress density.  Then V36 applies directly.

\begin{corollary}[Stress survival requires critical long-range stock]
\label{cor:v18v40-stress-threshold}
If the anchored third connected stock of the dynamically residual local
stress density is \(K_{3,n}^{T}\), then
\[
 \left|
 \operatorname{Cum}_n
 (T_{n,1}^{\perp},T_{n,2}^{\perp},T_{n,3}^{\perp})
 \right|
 \leq
 C h_n^2K_{3,n}^{T}
 \prod_{i=1}^3\|q_i\|_\infty.                               \tag{50.23}
\]
Therefore a nonzero separated cubic residual response requires
\[
 K_{3,n}^{T}\gtrsim h_n^{-2}.                               \tag{50.24}
\]
Uniform Dobrushin summability forces the separated residual response to
vanish.
\end{corollary}

\begin{proof}
Apply the V36 compiled cumulant estimate at \(r=3\) to the residual stress
density.  Rearrangement gives (50.24).
\end{proof}

This locates the precise conflict between the V34 convex-mixing branch and
an interacting stress source.  The finite-regulator stability theorem
remains useful, but an all-scale uniform summability bound on the same
dynamically residual metric score would kill its separated cubic Einstein
source.

\subsection{Next Einstein step}

The physical packet is now testable without fitting a terminal stress
tensor.  One varies the common action in three disjoint metric directions,
compiles the complete response by (50.4), evaluates the raw witness
(50.16), and then projects both direct and Gaussian packets through the same
declared quadratic follower grammar before evaluating (50.21).  A positive
eligible excess is an interaction witness; zero excess sends the programme
to the scale-critical covariance and higher-response rows.

Before continuing, the V40 checkpoint requires the scheduled inspection of
the active Yang--Mills and Navier--Stokes notes.  V41 will record that audit
and then derive the diffeomorphism Ward identity needed to turn the
metric-response packet into a conserved distributional Einstein source.

\begin{assessment}[V40 acquisition and boundary]
\label{ass:v18v40}
V40 proves the complete set-partition formula for metric derivatives of the
regulated effective action, separates all seagull/contact terms on disjoint
supports, polarizes the quadratic spectral cumulant compiler, and defines
both raw and dynamically shorted direct-minus-Gaussian cubic stress
witnesses.  It proves that a surviving separated cubic residual response
requires third-cumulant stock of order at least \(h_n^{-2}\).  The literal
regulated SM metric kernels and their Ward-conserving continuum limit remain
unpopulated.
\end{assessment}

\begin{record}[V40 append record]
\label{rec:v18v40}
V40 is appended to the validated V39 whose installed SHA--256 was
\texttt{53071DAA0D9B98DF7564A7FFBD67B5BA59BA7BE0C03E793D7FC91E92C3136A20}.
After installation only V40 is the active numeric SM note; the frozen
\texttt{V100\_f17} archive is unchanged.
\end{record}


\section{V41: Ward-defect descent and the conserved Einstein-source handoff}
\label{sec:v18v41}

\subsection{Mandatory V40 companion audit}

The newest active companion notes at this checkpoint are
\[
\begin{array}{c|c|c|c}
\hbox{programme}&\hbox{version}&\hbox{bytes}&\hbox{SHA--256}\\ \hline
\hbox{Yang--Mills}&V7&124028&
\texttt{1CF27C0B8913EABDBE1D1BA99290C487176D65E8ECAFD051DD8CE06B33CD761B}\\
\hbox{Navier--Stokes}&V26&278283&
\texttt{82C887F8C2C69E4F28FAD7C3DC8DE5B9099CB5A4BF431EBA1FFF3E91935978AD}.
\end{array}                                                     \tag{51.1}
\]
Yang--Mills V7 constructs a rational pivot atlas and spends a rigorously
bounded part of its error on a small toron collar before treating the
nonsingular complement.  The transferable principle is to budget a
regulator Ward defect as collision collar plus regular complement, with the
complete response compiled before absolute values.  Its Haar-sign and mass
gap conclusions are not imported.  Navier--Stokes V26 remains unchanged;
its rank-one refinement is empty at first derivative and modest at second,
so it does not change the SM direction.

\subsection{Finite-regulator Ward identity}

Let \(\xi\) be a smooth compactly supported vector field.  Denote the metric
variation by
\[
 q_\xi={\cal L}_\xi g                                        \tag{51.2}
\]
and the corresponding infinitesimal transformation of all regulated matter,
gauge, ghost, and auxiliary variables by \({\cal R}_{n,\xi}\phi\).  Let
\({\cal J}_{n,\xi}\) be the logarithmic Jacobian derivative of the regulated
reference measure under this field transformation.  Define the complete
diffeomorphism Ward defect
\[
\boxed{
 {\cal A}_{n}(\xi)
 =
 D_gS_n[q_\xi]
 +
 D_\phi S_n[{\cal R}_{n,\xi}\phi]
 -
 {\cal J}_{n,\xi}.}                                         \tag{51.3}
\]
This definition keeps action, gauge-fixing, determinant/Pfaffian, ghost,
measure, and regulator terms on one ledger.

\begin{theorem}[Inserted finite-regulator Ward identity]
\label{thm:v18v41-inserted-ward}
Let \(O_n(\phi)\) be differentiable and integrable together with the terms
below.  Then
\[
\boxed{
 \mathbb E_n\!\left[O_nD_gS_n[q_\xi]\right]
 =
 \mathbb E_n[O_n{\cal A}_n(\xi)]
 -
 \mathbb E_n[D_\phi O_n[{\cal R}_{n,\xi}\phi]].}             \tag{51.4}
\]
For \(O_n=1\),
\[
\boxed{
 \mathbb E_nD_gS_n[q_\xi]
 =
 \mathbb E_n{\cal A}_n(\xi).}                                \tag{51.5}
\]
\end{theorem}

\begin{proof}
Apply the infinitesimal field change
\(\phi\mapsto\phi+\varepsilon{\cal R}_{n,\xi}\phi\) to the finite regulated
integral with insertion \(O_n\).  Differentiation at zero gives
\[
 0=
 \mathbb E_n\!\left[
 D_\phi O_n[{\cal R}_{n,\xi}\phi]
 -O_nD_\phi S_n[{\cal R}_{n,\xi}\phi]
 +O_n{\cal J}_{n,\xi}
 \right].                                                    \tag{51.6}
\]
Solve (51.3) for
\(D_\phi S_n[{\cal R}_{n,\xi}\phi]-{\cal J}_{n,\xi}\), insert it into
(51.6), and rearrange.  This proves (51.4); taking \(O_n=1\) proves
(51.5).
\end{proof}

\begin{firewall}[Meaning of zero defect]
\label{fire:v18v41-zero-defect}
The condition \({\cal A}_n(\xi)=0\) is not imposed by naming the regulator
diffeomorphism invariant.  It must follow from the combined variation of
every term in (51.3).  A nonzero Jacobian, gauge-fixing defect, boundary
collar, or chiral anomaly remains in \({\cal A}_n\) until it is cancelled or
shown to vanish in the continuum topology.
\end{firewall}

\subsection{Distributional stress conservation}

Normalize the stress tensor by
\[
 D_gS_n[q]
 =
 \frac12\langle T_n,q\rangle.                                \tag{51.7}
\]
For a symmetric tensor and a boundary-free compactly supported test,
\[
 \frac12\langle T_n,{\cal L}_\xi g\rangle
 =
 -\langle\operatorname{div}_gT_n,\xi\rangle.                 \tag{51.8}
\]

\begin{corollary}[Exact regulated divergence identity]
\label{cor:v18v41-divergence}
The mean stress satisfies
\[
\boxed{
 \left\langle
 \operatorname{div}_g\mathbb E_nT_n,\xi
 \right\rangle
 =
 -\mathbb E_n{\cal A}_n(\xi).}                               \tag{51.9}
\]
If the complete regulator has zero Ward defect, its mean stress is exactly
distributionally conserved.
\end{corollary}

\begin{proof}
Combine (51.5), (51.7), and (51.8).
\end{proof}

\begin{theorem}[Ward-defect continuum descent]
\label{thm:v18v41-ward-descent}
Let \(s>0\).  Suppose
\[
 \mathbb E_nT_n\longrightarrow T
 \quad\hbox{in }H^{-s}_{\rm loc},                            \tag{51.10}
\]
and, for every compact set \(K\),
\[
 \sup_{\substack{\xi\in C_c^\infty(K;TM)\\
                  \|\xi\|_{H^{s+1}}\leq1}}
 \left|\mathbb E_n{\cal A}_n(\xi)\right|
 \longrightarrow0.                                         \tag{51.11}
\]
Then
\[
\boxed{\operatorname{div}_gT=0}
\quad\hbox{in }{\cal D}'(M;T^*M).                            \tag{51.12}
\]
\end{theorem}

\begin{proof}
For every test vector field \(\xi\), convergence in (51.10) permits passage
to the limit in
\(\langle\mathbb E_nT_n,\nabla^{(g)}\xi\rangle\).
Equation (51.9) and (51.11) show that the limit is zero.  This is precisely
the distributional definition of (51.12).
\end{proof}

\subsection{Collision-collar plus complement budget}

Let \({\cal C}_n\) be the set of regulator cells meeting a collision,
chart, gauge-fixing, or boundary collar, and let
\({\cal C}_n^c\) be its complement.  Suppose the mean Ward defect has a
cell decomposition
\[
 \mathbb E_n{\cal A}_n(\xi)
 =
 \sum_{z\in{\cal C}_n}a_{n,z}(\xi)
 +
 \sum_{z\in{\cal C}_n^c}a_{n,z}(\xi).                        \tag{51.13}
\]

\begin{theorem}[Ward collar budget]
\label{thm:v18v41-collar-budget}
Assume, for all tests supported in a fixed compact set,
\[
\begin{split}
 \sum_{z\in{\cal C}_n}|a_{n,z}(\xi)|
 &\leq M_n\,v_n\,\|\xi\|_{H^{s+1}},\\
 \left|
 \sum_{z\in{\cal C}_n^c}a_{n,z}(\xi)
 \right|
 &\leq\varepsilon_n\|\xi\|_{H^{s+1}},                        \tag{51.14}
\end{split}
\]
where \(v_n\) is the physical collar-volume budget.  Then
\[
 \|\mathbb E_n{\cal A}_n\|_{H^{-(s+1)}}
 \leq M_nv_n+\varepsilon_n.                                 \tag{51.15}
\]
In particular, (51.11) follows if
\[
 M_nv_n+\varepsilon_n\longrightarrow0.                       \tag{51.16}
\]
\end{theorem}

\begin{proof}
Apply the triangle inequality only after the two regions in (51.13) have
been separately compiled.  Equations (51.14) give
\[
 |\mathbb E_n{\cal A}_n(\xi)|
 \leq(M_nv_n+\varepsilon_n)\|\xi\|_{H^{s+1}}.
\]
Taking the dual supremum proves (51.15), and (51.16) implies (51.11).
\end{proof}

\begin{assessment}[Use of the YM V7 transfer]
\label{ass:v18v41-ym-transfer}
The collar may have a large pointwise defect \(M_n\); small volume alone is
not enough.  The required quantity is the paid product \(M_nv_n\).
Likewise, the complement must be certified independently by
\(\varepsilon_n\).  This is the precise Ward analogue of the YM V7
collar/complement split.
\end{assessment}

\subsection{Inserted conservation away from operator supports}

Let \(O_n\) be centered and suppose the support of \(\xi\) is disjoint from
the physical support of \(O_n\).  Local covariance of the insertion gives
\[
 D_\phi O_n[{\cal R}_{n,\xi}\phi]=0.                          \tag{51.17}
\]

\begin{corollary}[Separated inserted Ward identity]
\label{cor:v18v41-inserted-separated}
Under (51.17),
\[
\boxed{
 \mathbb E_n\!\left[O_nD_gS_n[{\cal L}_\xi g]\right]
 =
 \mathbb E_n[O_n{\cal A}_n(\xi)].}                           \tag{51.18}
\]
Thus connected stress correlators are transverse away from insertion
contacts whenever the inserted Ward defect tends to zero uniformly.
\end{corollary}

\begin{proof}
Equation (51.18) is (51.4) with the last term zero.  Since \(O_n\) is
centered, the left side is the connected stress insertion.
\end{proof}

This statement applies to the separated cubic witness of V40 by taking
\(O_n\) to be a product of the other two centered, follower-shorted stress
insertions, provided their supports and the vector-field support are
separated.  Conservation and non-Gaussianity are therefore independent
checks: a positive cubic excess must also satisfy (51.18).

\subsection{Einstein residual handoff}

Let \(g_n\) be the regulated metric and set
\[
 {\cal E}_n
 =
 G(g_n)+\Lambda g_n-\kappa\,\mathbb E_nT_n,                  \tag{51.19}
\]
where all terms are transported to the same tensor-distribution space and
physical units.  Let \({\cal B}_n\) denote the regulated Bianchi defect:
\[
 \left\langle
 \operatorname{div}_{g_n}(G(g_n)+\Lambda g_n),\xi
 \right\rangle
 =
 {\cal B}_n(\xi).                                            \tag{51.20}
\]

\begin{theorem}[Conserved continuum Einstein residual]
\label{thm:v18v41-einstein-handoff}
Assume, on every compact set:
\begin{enumerate}[leftmargin=2.2em]
\item \(g_n\to g\) strongly enough that the transported Einstein tensors
      converge in \(H^{-s}\) and divergence against smooth tests passes to
      the limit;
\item \(\mathbb E_nT_n\to T\) in \(H^{-s}\);
\item \({\cal B}_n\to0\) and
      \(\mathbb E_n{\cal A}_n\to0\) in \(H^{-(s+1)}\).
\end{enumerate}
Then
\[
 {\cal E}_n\longrightarrow
 {\cal E}=G(g)+\Lambda g-\kappa T
 \quad\hbox{in }H^{-s},                                     \tag{51.21}
\]
and
\[
\boxed{\operatorname{div}_g{\cal E}=0.}                     \tag{51.22}
\]
If additionally \(\|{\cal E}_n\|_{H^{-s}}\to0\), then
\[
\boxed{G(g)+\Lambda g=\kappa T}
\quad\hbox{in }H^{-s}.                                      \tag{51.23}
\]
\end{theorem}

\begin{proof}
The first two hypotheses give (51.21).  By (51.9)--(51.20),
\[
 \langle\operatorname{div}{\cal E}_n,\xi\rangle
 =
 {\cal B}_n(\xi)
 +
 \kappa\,\mathbb E_n{\cal A}_n(\xi),                         \tag{51.24}
\]
up to the stated transport identification.  The third hypothesis and
passage to the limit give (51.22).  If the residual norms tend to zero,
(51.21) has zero limit, which is (51.23).
\end{proof}

\begin{firewall}[Conservation is not the Einstein equation]
\label{fire:v18v41-conservation-not-equation}
Equation (51.22) is a compatibility condition.  It does not imply
\({\cal E}=0\).  The residual norm in the last hypothesis of Theorem
\ref{thm:v18v41-einstein-handoff} is a separate analytic gate and remains
unpopulated for the literal full SM regulator.
\end{firewall}

\subsection{Revised acquisition order}

The next physical work has three independent rows:
\begin{enumerate}[leftmargin=2.2em]
\item compute the literal scalar, gauge, fermion, ghost, and measure
      contributions to (51.3) and prove a collar/complement bound of the form
      (51.16);
\item populate the dynamically shorted metric kernels and cubic excess of
      V40, then verify the inserted Ward identity (51.18);
\item establish tightness of the same stress tensor and a vanishing Einstein
      residual norm, rather than inferring the field equation from
      conservation alone.
\end{enumerate}
The scalar metric score is the smallest row whose action derivative,
Gaussian kernel, Ward transformation, and continuum scaling can all be
written explicitly; it is the next target.

\begin{assessment}[V41 acquisition and boundary]
\label{ass:v18v41}
V41 records the mandatory companion audit, proves the inserted
finite-regulator Ward identity with the complete Jacobian and regulator
defect, gives a quantitative collar/complement descent criterion, and proves
a conditional conserved Einstein-residual handoff.  It also isolates the
additional norm gate required for the Einstein equation itself.  No
full-theory Ward-defect estimate, stress tightness theorem, or vanishing
Einstein residual has yet been supplied.
\end{assessment}

\begin{record}[V41 append record]
\label{rec:v18v41}
V41 is appended to the validated V40 whose installed SHA--256 was
\texttt{526A9AD57AAC405CA38604ABCDA925D7F4EFAB70AD532D0890E6218B1E993544}.
After installation only V41 is the active numeric SM note; the frozen
\texttt{V100\_f17} archive is unchanged.
\end{record}


\section{V42: the literal Higgs metric score and its complete Gaussian chaos decomposition}
\label{sec:v18v42}

\subsection{Finite common-action scalar sector}

Fix one regulated geometry and gauge background.  Write the four real Higgs
components as
\[
 h=(h_z)_{z\in\Lambda_n}\in{\mathbb R}^{4|\Lambda_n|}.        \tag{52.1}
\]
The most economical finite scalar action containing hopping, mass,
curvature coupling, and quartic potential is
\[
\boxed{
 S_{H,n}(g,U,h)
 =
 \frac12 h^{\mathsf T}K_n(g,U)h
 +
 \lambda_n\sum_{z\in\Lambda_n}w_{n,z}(g)\,|h_z|^4.}          \tag{52.2}
\]
Here \(K_n\) contains the gauge-covariant hopping matrix, the quadratic
potential, and any admitted \(\xi R|H|^2\) term; \(w_{n,z}\) is the physical
cell volume.  No continuum stress formula is fitted independently of
(52.2).

For a metric direction \(q\), let
\[
 \dot K_n[q]=D_gK_n[q],
 \qquad
 \dot w_{n,z}[q]=D_gw_{n,z}[q].                              \tag{52.3}
\]

\begin{theorem}[Literal scalar metric score]
\label{thm:v18v42-literal-score}
The scalar contribution to the V40 singleton metric insertion is
\[
\boxed{
 S_{H,n}[q]
 =
 \frac12h^{\mathsf T}\dot K_n[q]h
 +
 \lambda_n\sum_z\dot w_{n,z}[q]|h_z|^4.}                    \tag{52.4}
\]
Its centered stress score is
\[
 T_{H,n}[q]
 =
 S_{H,n}[q]-\mathbb E_nS_{H,n}[q].                           \tag{52.5}
\]
All scalar seagulls are obtained by further metric differentiation of
\(K_n\) and \(w_{n,z}\); they are therefore already included in the V40
partition compiler.
\end{theorem}

\begin{proof}
Differentiate (52.2) in the declared metric direction.  The field and gauge
coordinates are held fixed in a metric source derivative.  This gives
(52.4), and centering gives (52.5).  Repeated differentiation gives the
mixed insertions \(S_{n,B}\) of V40 with no additional stochastic input.
\end{proof}

\begin{firewall}[Sector scope]
\label{fire:v18v42-sector-scope}
Equation (52.4) is the scalar-sector metric score.  The full stress source
also contains gauge, fermion, ghost, gauge-fixing, measure, and
gravitational counterterm contributions.  Those terms must be added before
the V41 Ward defect is evaluated.
\end{firewall}

\subsection{Matched Gaussian scalar packet}

Assume first that the finite gauge-invariant scalar law is centered, \(\mathbb E_nh=0\).
Let \(G_n\) be a centered Gaussian vector with the same scalar covariance
\[
 C_n=\mathbb E_n[h h^{\mathsf T}]                            \tag{52.6}
\]
on its supported subspace, and write \(G_n=C_n^{1/2}g_n\) with \(g_n\)
standard Gaussian.  The kinetic/quadratic part of the centered Gaussian
score is
\[
 I_2(B_{n}^{\rm kin}[q]),
 \qquad
 B_n^{\rm kin}[q]
 =
 \frac12C_n^{1/2}\dot K_n[q]C_n^{1/2}.                       \tag{52.7}
\]
The quartic volume score cannot in general be discarded or represented by
a quadratic kernel.

For an exact closed formula, assume at first that every one-cell covariance
is gauge-isotropic:
\[
 \mathbb E[G_{n,z}G_{n,z}^{\mathsf T}]
 =
 \sigma_{n,z}^2I_4.                                         \tag{52.8}
\]
Define the normalized cell radius
\[
 R_{n,z}=\frac{|G_{n,z}|^2}{\sigma_{n,z}^2}.                 \tag{52.9}
\]

\begin{lemma}[Four-component radial Wick decomposition]
\label{lem:v18v42-radial-wick}
For a standard four-dimensional Gaussian vector,
\[
 \mathbb ER=4,\qquad \mathbb ER^2=24,                        \tag{52.10}
\]
and the fourth-chaos radial polynomial is
\[
 {\cal H}_{4}(R)=R^2-12R+24.                                \tag{52.11}
\]
Consequently,
\[
\boxed{
 R^2-\mathbb ER^2
 =
 {\cal H}_{4}(R)+12(R-4).}                                  \tag{52.12}
\]
The two terms on the right belong respectively to the fourth and second
Gaussian chaoses and are orthogonal in \(L^2\).
\end{lemma}

\begin{proof}
The squared radius is chi-square with four degrees of freedom, giving
(52.10).  The polynomial (52.11) is orthogonal to constants because its
expectation is \(24-48+24=0\).  It is orthogonal to \(R-4\) because
\[
 \mathbb E[R^3]-16\mathbb E[R^2]+72\mathbb E[R]-96
 =
 192-384+288-96=0.                                         \tag{52.13}
\]
It is a degree-four polynomial in the Gaussian coordinates with its degree
zero and degree two contractions removed, hence lies in the fourth chaos.
Rearrangement proves (52.12).
\end{proof}

Let \(P_{n,z}\) be the orthogonal projection in the whitened one-particle
space onto the four normalized coordinates of cell \(z\).  Then
\[
 R_{n,z}-4=I_2(P_{n,z}).                                    \tag{52.14}
\]
Define
\[
\begin{split}
 B_n^{H}[q]
 &=
 B_n^{\rm kin}[q]
 +
 12\lambda_n\sum_z
 \dot w_{n,z}[q]\sigma_{n,z}^4P_{n,z},                       \tag{52.15}\\
 I_4(D_n^{H}[q])
 &=
 \lambda_n\sum_z
 \dot w_{n,z}[q]\sigma_{n,z}^4
 {\cal H}_4(R_{n,z}).                                       \tag{52.16}
\end{split}
\]

\begin{theorem}[Complete Gaussian scalar stress decomposition]
\label{thm:v18v42-complete-gaussian}
Under (52.8), the centered matched-Gaussian version of (52.4) is exactly
\[
\boxed{
 T_{H,n}^{\rm G}[q]
 =
 I_2(B_n^{H}[q])+I_4(D_n^{H}[q]).}                           \tag{52.17}
\]
Its covariance is
\[
\boxed{
 \operatorname{Cov}
 (T_{H,n}^{\rm G}[q],T_{H,n}^{\rm G}[r])
 =
 2\operatorname{tr}(B_n^{H}[q]B_n^{H}[r])
 +
 \mathbb E[I_4(D_n^{H}[q])I_4(D_n^{H}[r])].}                \tag{52.18}
\]
\end{theorem}

\begin{proof}
The quadratic part is (52.7).  At each cell,
\[
 |G_{n,z}|^4-\mathbb E|G_{n,z}|^4
 =
 \sigma_{n,z}^4
 \bigl({\cal H}_4(R_{n,z})+12(R_{n,z}-4)\bigr)               \tag{52.19}
\]
by Lemma \ref{lem:v18v42-radial-wick}.  Summation with the coefficients in
(52.4) yields (52.15)--(52.17).  Distinct homogeneous Gaussian chaoses are
orthogonal, so the cross covariance between \(I_2\) and \(I_4\) vanishes,
which proves (52.18).
\end{proof}

\begin{corollary}[Quadratic-only comparison is incomplete]
\label{cor:v18v42-quadratic-incomplete}
If
\[
 \lambda_n\dot w_{n,z}[q]\ne0                               \tag{52.20}
\]
for some cell of nonzero scalar variance, the matched Gaussian stress
contains a nonzero fourth-chaos insertion.  Comparing the direct scalar
stress only with the V40 quadratic trace formula can therefore create a
false excess.  The Gaussian comparator must contain both terms in (52.17).
\end{corollary}

\begin{proof}
The radial fourth Hermite polynomial has positive variance, and its
coefficient in (52.16) is nonzero under (52.20).  It cannot be represented
inside the orthogonal second chaos.
\end{proof}

\begin{firewall}[Nonzero scalar mean]
\label{fire:v18v42-nonzero-mean}
If a symmetry-breaking gauge-fixed chart has nonzero scalar mean, its quartic
metric score also has first- and third-chaos components.  They must be added
to (52.17) before forming a comparator; the centered gauge-invariant branch
is the only branch treated by the displayed decomposition.
\end{firewall}

\subsection{Full cubic Gaussian prediction}

For three metric directions, set
\[
 X_{n,i}^{\rm G}
 =
 I_2(B_n^{H}[q_i])+I_4(D_n^{H}[q_i]).                        \tag{52.21}
\]

\begin{definition}[Complete scalar Gaussian cubic]
\label{def:v18v42-complete-cubic}
The matched Gaussian prediction is
\[
\boxed{
 {\cal G}_{3,n}^{H}(q_1,q_2,q_3)
 =
 \operatorname{Cum}_{\rm G}
 (X_{n,1}^{\rm G},X_{n,2}^{\rm G},X_{n,3}^{\rm G}).}         \tag{52.22}
\]
Equivalently, it is the sum of the eight second/fourth-chaos choices
\[
 \sum_{\alpha_1,\alpha_2,\alpha_3\in\{2,4\}}
 \operatorname{Cum}_{\rm G}
 \bigl(I_{\alpha_1}(D_{n,1}^{(\alpha_1)}),
       I_{\alpha_2}(D_{n,2}^{(\alpha_2)}),
       I_{\alpha_3}(D_{n,3}^{(\alpha_3)})\bigr),              \tag{52.23}
\]
with \(D_{n,i}^{(2)}=B_n^H[q_i]\) and
\(D_{n,i}^{(4)}=D_n^H[q_i]\).  Every term is fixed by Wick contraction of
the same covariance \(C_n\).
\end{definition}

\begin{theorem}[Correct scalar direct-minus-Gaussian witness]
\label{thm:v18v42-correct-witness}
On pairwise positively separated metric supports, define
\[
\boxed{
 {\cal E}_{3,n}^{H}
 =
 \operatorname{Cum}_n
 (T_{H,n}[q_1],T_{H,n}[q_2],T_{H,n}[q_3])
 -
 {\cal G}_{3,n}^{H}(q_1,q_2,q_3).}                           \tag{52.24}
\]
If this quantity has a nonzero cofinal limit after applying the same
declared follower projection to the direct and Gaussian packets, then the
scalar stress sector is not quasi-free modulo that follower grammar.
\end{theorem}

\begin{proof}
The first term is the direct connected tensor of the literal score (52.5).
Theorem \ref{thm:v18v42-complete-gaussian} and Wick's theorem make
(52.22)--(52.23) the complete Gaussian prediction, including the quartic
volume term.  Equal follower projections remove the same predeclared
ancestry on both sides.  A quasi-free realization has exactly the Gaussian
Wick contractions, so a nonzero residual difference excludes it.
\end{proof}

\subsection{Follower projection must also include fourth chaos}

Let \({\cal C}_n\) be the primitive first-chaos source-cyclic carrier of V39.
The minimum scalar Gaussian follower space through degree four is
\[
 {\cal F}_{n,H}^{\leq4}
 =
 \operatorname{Sym}^2({\cal C}_n)
 \oplus
 \operatorname{Sym}^4({\cal C}_n).                           \tag{52.25}
\]
Let \(P_n^{\otimes_{\rm s}k}\) denote the symmetric tensor power of the
first-chaos projection.

\begin{proposition}[Degree-complete scalar follower short]
\label{prop:v18v42-degree-complete-short}
The orthogonal Gaussian follower projection of (52.17) is
\[
\boxed{
 I_2(P_nB_n^H[q]P_n)
 +
 I_4\bigl((P_n^{\otimes_{\rm s}4})D_n^H[q]\bigr).}           \tag{52.26}
\]
Hence the minimum residual packet is
\[
\boxed{
\begin{split}
 (T_{H,n}^{\rm G}[q])^\perp
 ={}&
 I_2(B_n^H[q]-P_nB_n^H[q]P_n)\\
 &+
 I_4\bigl(
 D_n^H[q]-(P_n^{\otimes_{\rm s}4})D_n^H[q]
 \bigr).
\end{split}}                                                 \tag{52.27}
\]
\end{proposition}

\begin{proof}
Homogeneous chaoses of different degrees are orthogonal.  Within degree
\(k\), the Gaussian-chaos isometry identifies the closed span generated by
\({\cal C}_n\) with \(\operatorname{Sym}^k({\cal C}_n)\), whose orthogonal
projection is \(P_n^{\otimes_{\rm s}k}\).  Apply this separately at degrees
two and four.
\end{proof}

\begin{firewall}[Mixed residual legs]
\label{fire:v18v42-mixed-legs}
As in V39, (52.27) removes only tensors whose every leg lies in the primitive
carrier.  A fourth-chaos word with at least one residual leg remains unless
a larger physical follower grammar explicitly synthesizes it.
\end{firewall}

\subsection{What summable locality still rules out}

\begin{corollary}[Local quartic volume scores remain in the V36 basin]
\label{cor:v18v42-local-quartic-basin}
Suppose the cellwise quartic-score field in (52.16), after the declared
follower short, has uniformly summable connected cumulants of every fixed
order and bounded physical coefficients.  Then its critically summed
continuum limit is Gaussian.  A non-Gaussian scalar stress limit must
therefore arise from a scale-critical covariance of the kinetic and quartic
packets, a macroscopic residual chaos tensor, or their non-summable mixed
connected contractions.
\end{corollary}

\begin{proof}
The fourth-chaos insertion is still a centered additive local score over
cells.  The V36 finite-bank cumulant theorem applies to its compiled scalar
direction.  It kills all cumulants of order at least three under the stated
summability hypothesis.
\end{proof}

\subsection{Next sector}

The scalar metric occurrence and its complete degree-two/four Gaussian
comparator are now fixed.  The next note should add the gauge kinetic metric
score and verify that the scalar--gauge sum, rather than either sector
separately, satisfies the diffeomorphism Ward identity: their divergences
exchange the gauge Lorentz-force term.  Only after that cancellation may the
combined bosonic stress enter the Einstein handoff of V41.

\begin{assessment}[V42 acquisition and boundary]
\label{ass:v18v42}
V42 derives the literal regulated Higgs metric score from the common action,
proves its exact second-plus-fourth Gaussian chaos decomposition, corrects
the quadratic-only comparator, and lifts the primitive follower short
through degree four.  The full scalar direct-minus-Gaussian cubic is now
well-defined.  Its cofinal limit, a macroscopic residual chaos mode, and the
gauge/fermion Ward cancellations remain unproved.
\end{assessment}

\begin{record}[V42 append record]
\label{rec:v18v42}
V42 is appended to the validated V41 whose installed SHA--256 was
\texttt{BA4F66AB69D96E4A94824FB1A99C06D03507A61B43020CF354EE1FBC1548C95E}.
After installation only V42 is the active numeric SM note; the frozen
\texttt{V100\_f17} archive is unchanged.
\end{record}


\section{V43: the literal gauge metric score and bosonic Ward exchange}
\label{sec:v18v43}

\subsection{Common-action gauge occurrence}

Let \(G\) be one compact Standard-Model gauge factor with a fixed invariant
inner product on its Lie algebra.  On a Riemannian regulator the continuum
notation for its kinetic action is
\[
 S_{G,n}(g,A)
 =
 \frac{1}{4g_{G,n}^2}
 \int_M\sqrt g\,
 \langle F_{\mu\nu},F^{\mu\nu}\rangle\,dx,                   \tag{53.1}
\]
with the integral understood as its declared finite plaquette
discretization.  Equivalently, on the actual lattice,
\[
\boxed{
 S_{G,n}(g,U)
 =
 \sum_{p\in{\cal P}_n}
 \beta_{n,p}(g)\,
 \ell_p(U),
 \qquad
 \ell_p(U)=
 1-\frac{1}{d_R}\operatorname{Re}\operatorname{Tr}_R U_p.}  \tag{53.2}
\]
All anisotropy, dual-area, and physical-unit factors are contained in
\(\beta_{n,p}(g)\).

\begin{theorem}[Literal gauge metric score]
\label{thm:v18v43-literal-gauge-score}
For a metric direction \(q\), the gauge contribution to the V40 singleton
score is
\[
\boxed{
 S_{G,n}[q]
 =
 \sum_{p\in{\cal P}_n}
 \dot\beta_{n,p}[q]\,\ell_p(U).}                             \tag{53.3}
\]
Its centered stress score is
\[
 T_{G,n}[q]
 =
 S_{G,n}[q]-\mathbb E_nS_{G,n}[q].                           \tag{53.4}
\]
Every gauge seagull is a further metric derivative of the same
\(\beta_{n,p}\) and therefore enters the V40 partition compiler without a
new source.
\end{theorem}

\begin{proof}
The plaquette holonomy is held fixed in a metric source derivative; all
explicit metric dependence in (53.2) is in \(\beta_{n,p}\).
Differentiation gives (53.3), centering gives (53.4), and repeated
differentiation gives the seagulls.
\end{proof}

Summing (53.2)--(53.4) over
\[
 G=SU(3),\quad SU(2),\quad U(1)                              \tag{53.5}
\]
gives the gauge part of the bosonic Standard-Model metric score.

\subsection{Continuum stress tensor fixed by the same variation}

Retain the V41 convention
\[
 D_gS_G[q]
 =
 \frac12\int_M\sqrt g\,
 \Theta_G^{\mu\nu}q_{\mu\nu}\,dx.                            \tag{53.6}
\]

\begin{proposition}[Gauge metric tensor]
\label{prop:v18v43-gauge-tensor}
For (53.1),
\[
\boxed{
 \Theta_G^{\mu\nu}
 =
 \frac{1}{g_{G,n}^2}
 \left(
 \frac14g^{\mu\nu}
 \langle F_{\alpha\beta},F^{\alpha\beta}\rangle
 -
 \langle F^{\mu}{}_{\rho},F^{\nu\rho}\rangle
 \right).}                                                  \tag{53.7}
\]
This is the metric-derivative convention used throughout V41--V43.
\end{proposition}

\begin{proof}
For a covariant metric variation \(q_{\mu\nu}\),
\[
 \delta\sqrt g=\frac12\sqrt g\,g^{\mu\nu}q_{\mu\nu},
 \qquad
 \delta g^{\mu\nu}=-q^{\mu\nu}.                              \tag{53.8}
\]
The two inverse metrics contracting \(F_{\mu\nu}F_{\alpha\beta}\) give
\[
 \delta\langle F_{\alpha\beta},F^{\alpha\beta}\rangle
 =
 -2q_{\mu\nu}
 \langle F^{\mu}{}_{\rho},F^{\nu\rho}\rangle.                \tag{53.9}
\]
Insert (53.8)--(53.9) into (53.1) and compare with (53.6).
\end{proof}

\subsection{Gauge invariance and locality of the score}

\begin{lemma}[Gauge invariance of the metric insertion]
\label{lem:v18v43-gauge-invariance}
For every gauge transformation \(V=(V_z)\),
\[
 S_{G,n}[q](U^V)=S_{G,n}[q](U).                              \tag{53.10}
\]
Thus the metric score descends to the finite gauge-orbit quotient before
any Gaussian comparison or continuum limit.
\end{lemma}

\begin{proof}
Plaquette holonomy transforms by conjugation at its base vertex.  The real
trace in \(\ell_p\) is conjugation invariant, and
\(\dot\beta_{n,p}[q]\) is deterministic.
\end{proof}

\begin{corollary}[Uniformly summable plaquette scores Gaussianize]
\label{cor:v18v43-plaquette-gaussianization}
Suppose the compiled connected cumulants of the centered plaquette metric
density in (53.4) have uniform anchored sums.  Then its
four-dimensional critical continuum field is Gaussian and every connected
vertex of order at least three vanishes.  A non-Gaussian gauge stress
continuum therefore requires a scale-critical Wilson-loop stock or a
macroscopic residual spectral mode.
\end{corollary}

\begin{proof}
Equation (53.4) is an additive common-action score over a finite local
plaquette grammar.  Apply the V36 finite-bank theorem.
\end{proof}

The Yang--Mills V7 companion computation supplies a finite-regulator
quadrature and collar method for one Wilson response, but it does not
supply the scale-critical stock required by this corollary.

\subsection{Quadratic gauge comparator}

Choose a gauge-fixed local chart around a background connection and write
the quadratic part of the complete gauge-plus-gauge-fixing action as
\[
 \frac12a^{\mathsf T}K_{G,n}(g)a.                            \tag{53.11}
\]
Let \(C_{G,n}\) be the matched covariance on the physical supported
subspace.  The quadratic part of the centered Gaussian metric score is
\[
 I_2(B_{G,n}[q]),
 \qquad
 B_{G,n}[q]
 =
 \frac12C_{G,n}^{1/2}\dot K_{G,n}[q]C_{G,n}^{1/2}.            \tag{53.12}
\]

\begin{proposition}[Gauge spectral test]
\label{prop:v18v43-gauge-spectral}
After the V39 source-cyclic quadratic short,
\[
 B_{G,n}^{\perp}[q]
 =
 B_{G,n}[q]-P_nB_{G,n}[q]P_n.                               \tag{53.13}
\]
The variance-normalized quadratic gauge stress is non-Gaussian along a
subsequence only if
\[
 \limsup_n
 \frac{\|B_{G,n}^{\perp}[q]\|_{\rm op}}
 {\|B_{G,n}^{\perp}[q]\|_{\rm HS}}>0.                        \tag{53.14}
\]
This criterion concerns only the quadratic chart.  The full compact Wilson
score (53.3) remains the direct comparator.
\end{proposition}

\begin{proof}
Equation (53.13) is the V39 follower projection.  Apply the V38
effective-rank equivalence to the residual kernel.
\end{proof}

\begin{firewall}[Gauge fixing and ghosts]
\label{fire:v18v43-gauge-fixing}
The covariance in (53.12) is not defined on gauge zero modes before a
physical quotient or a gauge-fixing/ghost completion.  The metric
variations of gauge fixing, Faddeev--Popov determinant, and any regulator
measure must enter the complete V41 Ward defect.  They cannot be omitted
because (53.7) is gauge invariant.
\end{firewall}

\subsection{Exact scalar--gauge divergence exchange}

The scalar stress from V42 and the gauge stress (53.7) are not separately
conserved when the Higgs field is gauge coupled.  This can be proved without
choosing coordinate formulas for the force term.

Let
\[
 {\cal E}_{H,n}=D_hS_{H,n},
 \qquad
 {\cal J}_{H,n}=D_AS_{H,n},
 \qquad
 {\cal E}_{G,n}=D_AS_{G,n}.                                 \tag{53.15}
\]
Use the gauge-covariant infinitesimal diffeomorphism
\[
 \delta_\xi^Ah=\iota_\xi D_Ah,
 \qquad
 \delta_\xi^AA=\iota_\xi F_A.                               \tag{53.16}
\]
It differs from the ordinary Lie derivative by an infinitesimal gauge
transformation, which leaves both actions invariant.

\begin{theorem}[Bosonic off-shell Noether identities]
\label{thm:v18v43-bosonic-noether}
For compactly supported \(\xi\),
\[
\boxed{
 \langle\operatorname{div}\Theta_{H,n},\xi\rangle
 =
 {\cal E}_{H,n}[\delta_\xi^Ah]
 +
 {\cal J}_{H,n}[\delta_\xi^AA],}                             \tag{53.17}
\]
\[
\boxed{
 \langle\operatorname{div}\Theta_{G,n},\xi\rangle
 =
 {\cal E}_{G,n}[\delta_\xi^AA].}                             \tag{53.18}
\]
Consequently,
\[
\boxed{
 \langle\operatorname{div}
 (\Theta_{H,n}+\Theta_{G,n}),\xi\rangle
 =
 {\cal E}_{H,n}[\delta_\xi^Ah]
 +
 ({\cal J}_{H,n}+{\cal E}_{G,n})[\delta_\xi^AA].}            \tag{53.19}
\]
\end{theorem}

\begin{proof}
Diffeomorphism covariance and gauge invariance of the scalar action give
\[
 D_gS_{H,n}[{\cal L}_\xi g]
 +
 {\cal E}_{H,n}[\delta_\xi^Ah]
 +
 {\cal J}_{H,n}[\delta_\xi^AA]
 =0.                                                        \tag{53.20}
\]
By the metric convention and integration by parts,
\[
 D_gS_{H,n}[{\cal L}_\xi g]
 =
 -\langle\operatorname{div}\Theta_{H,n},\xi\rangle.          \tag{53.21}
\]
This proves (53.17).  The same calculation for the pure gauge action gives
\[
 D_gS_{G,n}[{\cal L}_\xi g]
 +
 {\cal E}_{G,n}[\delta_\xi^AA]=0,                            \tag{53.22}
\]
which proves (53.18).  Add the two identities.
\end{proof}

\begin{corollary}[On-shell Lorentz-force cancellation]
\label{cor:v18v43-force-cancellation}
If
\[
 {\cal E}_{H,n}=0,
 \qquad
 {\cal E}_{G,n}+{\cal J}_{H,n}=0,                            \tag{53.23}
\]
then
\[
\boxed{
 \operatorname{div}
 (\Theta_{H,n}+\Theta_{G,n})=0.}                             \tag{53.24}
\]
When \({\cal E}_{H,n}=0\), the two separate divergences in
(53.17)--(53.18) are equal and opposite after the gauge equation is used.
\end{corollary}

\begin{proof}
Insert (53.23) into (53.19).  The second assertion follows from
(53.17)--(53.18).
\end{proof}

\begin{countertheorem}[Sectorwise conservation is the wrong gate]
\label{cth:v18v43-sectorwise}
If the gauge current is nonzero, the conditions
\(\operatorname{div}\Theta_{H,n}=0\) and
\(\operatorname{div}\Theta_{G,n}=0\) generally fail even on a solution of
the coupled equations.  Requiring them separately can therefore reject a
correct interacting scalar--gauge configuration.
\end{countertheorem}

\begin{proof}
On the scalar equation, (53.17) reduces to the current force
\({\cal J}_{H,n}[\iota_\xi F]\).  On the coupled gauge equation,
(53.18) is its negative.  Unless that force vanishes, neither sector is
conserved separately while their sum is conserved by (53.24).
\end{proof}

\subsection{Quantum Ward handoff}

At finite cutoff, field integration replaces the equations of motion in
(53.19) by the change-of-variables terms in V41.  Gauge-fixing, ghost,
Jacobian, chart, and boundary contributions combine into a bosonic Ward
defect \({\cal A}_{B,n}(\xi)\).  Therefore
\[
\boxed{
 \left\langle
 \operatorname{div}\mathbb E_n(T_{H,n}+T_{G,n}),\xi
 \right\rangle
 =
 -\mathbb E_n{\cal A}_{B,n}(\xi).}                           \tag{53.25}
\]
The collar/complement theorem of V41 applies to this \emph{sum}.  A bound
for only one sector cannot establish the Einstein-source conservation row.

\subsection{Next sector}

The combined bosonic metric source is now fixed at the action level, and its
classical force exchange is exact.  The next missing term is the fermion
determinant/Pfaffian metric derivative.  It is nonlocal after fermions are
integrated out, so its score must be expressed as a resolvent trace on the
same determinant line and then compiled with gauge, ghost, and measure
terms before the V41 Ward defect is bounded.

\begin{assessment}[V43 acquisition and boundary]
\label{ass:v18v43}
V43 derives the literal plaquette metric score, its continuum metric tensor,
its quadratic spectral kernel, and the exact scalar--gauge divergence
exchange.  It proves that conservation is a property of the combined
bosonic source and excludes sectorwise conservation as a gate.  The
scale-critical Wilson stock, fermion/ghost contributions, anomaly
cancellation, and cofinal stress limit remain unproved.
\end{assessment}

\begin{record}[V43 append record]
\label{rec:v18v43}
V43 is appended to the validated V42 whose installed SHA--256 was
\texttt{111E8013B1B886043F95331EFE1D6E4ADAA7D0BF9583E2A66D447EAED328D3F5}.
After installation only V43 is the active numeric SM note; the frozen
\texttt{V100\_f17} archive is unchanged.
\end{record}


\section{V44: fermion determinant metric responses and the anomaly-arithmetic gate}
\label{sec:v18v44}

\subsection{Determinant-line local chart}

Let
\[
 D_n(g,U,H):{\cal V}_{n,+}\longrightarrow{\cal V}_{n,-}      \tag{54.1}
\]
be the finite chiral Dirac matrix on one declared determinant-line chart,
after the regulator representation, boundary condition, and zero-mode
handling have been fixed.  Assume \(D_n\) is invertible on this chart and
write
\[
 G_n=D_n^{-1}.                                                \tag{54.2}
\]
For a complex Dirac determinant the local fermion effective action is
\[
 S_{F,n}=-\log\det D_n.                                      \tag{54.3}
\]
For a Majorana/Pfaffian block with invertible antisymmetric matrix
\({\cal D}_n\), replace (54.3) by
\[
 S_{F,n}^{\rm Pf}=-\log\operatorname{Pf}{\cal D}_n
 =-\frac12\log\det{\cal D}_n.                                \tag{54.4}
\]
All formulas below acquire a factor \(1/2\) on a Pfaffian block.

\begin{firewall}[Local logarithm]
\label{fire:v18v44-local-log}
Equations (54.3)--(54.4) are chartwise identities.  They do not trivialize
the determinant line, choose a global phase, or remove zero modes.  Chart
transition phases and Berry holonomy remain separate data.
\end{firewall}

For metric directions \(q_i\), abbreviate
\[
 D_i=D_gD_n[q_i],
 \quad
 D_{ij}=D_g^2D_n[q_i,q_j],
 \quad
 D_{ijk}=D_g^3D_n[q_i,q_j,q_k].                              \tag{54.5}
\]

\subsection{Exact resolvent response compiler}

\begin{theorem}[First three determinant metric responses]
\label{thm:v18v44-det-responses}
On the invertible chart,
\[
\boxed{
 \partial_iS_{F,n}
 =
 -\operatorname{Tr}(G_nD_i),}                               \tag{54.6}
\]
\[
\boxed{
 \partial_{ij}S_{F,n}
 =
 \operatorname{Tr}(G_nD_iG_nD_j)
 -
 \operatorname{Tr}(G_nD_{ij}),}                             \tag{54.7}
\]
and
\[
\boxed{
\begin{split}
 \partial_{ijk}S_{F,n}
 ={}&
 -\operatorname{Tr}(G_nD_iG_nD_jG_nD_k)
 -\operatorname{Tr}(G_nD_iG_nD_kG_nD_j)\\
 &+\operatorname{Tr}(G_nD_{ij}G_nD_k)
  +\operatorname{Tr}(G_nD_{ik}G_nD_j)
  +\operatorname{Tr}(G_nD_{jk}G_nD_i)\\
 &-\operatorname{Tr}(G_nD_{ijk}).
\end{split}}                                                 \tag{54.8}
\]
The two cubic loop orientations and every seagull word are included before
any absolute value is taken.
\end{theorem}

\begin{proof}
Jacobi's formula gives
\[
 \partial_i\log\det D_n=\operatorname{Tr}(G_nD_i),
 \qquad
 \partial_iG_n=-G_nD_iG_n.                                  \tag{54.9}
\]
Equation (54.6) follows.  Differentiate it once and use (54.9) to obtain
(54.7).  Differentiate (54.7) in the third direction.  The two inverse
factors in its first trace produce the two negatively signed cubic loop
words; differentiating \(D_i,D_j,D_{ij}\) produces the three double-vertex
words and the final triple seagull.  Cyclicity of trace gives (54.8).
\end{proof}

\begin{corollary}[Real score and Berry current]
\label{cor:v18v44-real-phase}
On a chart where the logarithm is continuous,
\[
\boxed{
 T_{F,n}^{\rm amp}[q]
 =
 -\operatorname{Re}\operatorname{Tr}(G_nD_q),\qquad
 J_{F,n}^{\rm ph}[q]
 =
 -\operatorname{Im}\operatorname{Tr}(G_nD_q).}               \tag{54.10}
\]
The first is the real determinant-amplitude metric score and the second is
the determinant-line phase connection in direction \(q\).  They must be
transported separately.
\end{corollary}

\begin{proof}
Take real and imaginary parts of (54.6).
\end{proof}

\subsection{Gap-dependent response bounds}

\begin{theorem}[Three-response resolvent stock]
\label{thm:v18v44-resolvent-stock}
Suppose
\[
 \|G_n\|_{\rm op}\leq\gamma_n^{-1}.                          \tag{54.11}
\]
Using the Frobenius norm for \(D_i,D_{ij}\) and the trace norm for
\(D_{ijk}\), one has
\[
 |\partial_iS_{F,n}|
 \leq\gamma_n^{-1}\|D_i\|_1,                                \tag{54.12}
\]
\[
 |\partial_{ij}S_{F,n}|
 \leq
 \gamma_n^{-2}\|D_i\|_{\rm F}\|D_j\|_{\rm F}
 +
 \gamma_n^{-1}\|D_{ij}\|_1,                                 \tag{54.13}
\]
and
\[
\boxed{
\begin{split}
 |\partial_{ijk}S_{F,n}|
 \leq{}&
 2\gamma_n^{-3}
 \|D_i\|_{\rm F}\|D_j\|_{\rm F}\|D_k\|_{\rm F}\\
 &+\gamma_n^{-2}\bigl(
 \|D_{ij}\|_{\rm F}\|D_k\|_{\rm F}
 +\|D_{ik}\|_{\rm F}\|D_j\|_{\rm F}
 +\|D_{jk}\|_{\rm F}\|D_i\|_{\rm F}
 \bigr)\\
 &+\gamma_n^{-1}\|D_{ijk}\|_1.
\end{split}}                                                 \tag{54.14}
\]
\end{theorem}

\begin{proof}
The trace duality bound gives
\(|\operatorname{Tr}(GA)|\leq\|G\|_{\rm op}\|A\|_1\).
For two insertions use
\[
 |\operatorname{Tr}(GD_iGD_j)|
 \leq
 \|GD_i\|_{\rm F}\|GD_j\|_{\rm F}.                           \tag{54.15}
\]
For three insertions, cyclically group the product and apply
\[
 |\operatorname{Tr}(XYZ)|
 \leq\|X\|_{\rm F}\|Y\|_{\rm F}\|Z\|_{\rm op}
 \leq\|X\|_{\rm F}\|Y\|_{\rm F}\|Z\|_{\rm F}.                \tag{54.16}
\]
Each resolvent contributes at most \(\gamma_n^{-1}\).  Apply these bounds
termwise to the already compiled identities (54.6)--(54.8).
\end{proof}

\begin{assessment}[Gap implication]
\label{ass:v18v44-gap}
A uniform fermion gap controls every fixed response word, but a collapsing
\(\gamma_n\) can amplify the cubic loop like \(\gamma_n^{-3}\).  This is a
possible route to the V36 critical stock, accompanied by an independent
tightness problem.  A response floor cannot be inferred from the upper
bounds (54.12)--(54.14).
\end{assessment}

\subsection{Separated supports and fermion loops}

If \(D_i\) is local in the support of \(q_i\), the triple product in (54.8)
is the first term able to join three pairwise separated metric supports
without a contact insertion.  The \(D_{ij}\) and \(D_{ijk}\) terms vanish
for disjoint supports under metric locality.

\begin{corollary}[Separated fermion metric loop]
\label{cor:v18v44-separated-loop}
For pairwise positively separated metric directions and a metric-local
Dirac regulator,
\[
\boxed{
 \partial_{ijk}S_{F,n}
 =
 -\operatorname{Tr}(G_nD_iG_nD_jG_nD_k)
 -\operatorname{Tr}(G_nD_iG_nD_kG_nD_j).}                   \tag{54.17}
\]
This two-orientation sum is invariant under cyclic reordering and is the
complete chartwise separated fermion response.
\end{corollary}

\begin{proof}
Metric locality makes \(D_{ij}=D_{ik}=D_{jk}=D_{ijk}=0\) when the relevant
support intersections are empty.  Apply (54.8).
\end{proof}

\begin{firewall}[No termwise lower bound]
\label{fire:v18v44-loop-cancellation}
The two loop orientations in (54.17) may cancel after representation,
chirality, and phase are included.  Neither orientation is an interaction
certificate by itself.  The complete trace must be evaluated before its
absolute value, as in the compilation lesson recorded at V36.
\end{firewall}

\subsection{One-generation Standard-Model anomaly arithmetic}

Write every fermion as a left-handed Weyl field.  The hypercharges and
multiplicities in one generation are
\[
\begin{array}{c|c|c}
\hbox{field}&Y&\hbox{multiplicity}\\ \hline
Q_L&1/6&6\\
u_R^c&-2/3&3\\
d_R^c&1/3&3\\
L_L&-1/2&2\\
e_R^c&1&1.
\end{array}                                                   \tag{54.18}
\]
A right-handed neutrino, when present, has \(Y=0\) and changes none of the
sums.

\begin{theorem}[Perturbative anomaly cancellation per generation]
\label{thm:v18v44-anomaly-arithmetic}
The Standard-Model representation (54.18) satisfies
\[
\boxed{
 \sum_{\rm Weyl}Y=0,\qquad
 \sum_{\rm Weyl}Y^3=0,}                                     \tag{54.19}
\]
\[
\boxed{
 {\cal A}_{SU(2)^2U(1)}=0,\qquad
 {\cal A}_{SU(3)^2U(1)}=0,\qquad
 {\cal A}_{SU(3)^3}=0.}                                     \tag{54.20}
\]
It also contains four \(SU(2)\) doublets per generation, so the parity
count for the \(SU(2)\) global doublet anomaly is even.
\end{theorem}

\begin{proof}
The mixed gravitational--hypercharge sum is
\[
 6\left(\frac16\right)
 +3\left(-\frac23\right)
 +3\left(\frac13\right)
 +2\left(-\frac12\right)+1=0.                               \tag{54.21}
\]
The cubic hypercharge sum is
\[
 6\left(\frac16\right)^3
 +3\left(-\frac23\right)^3
 +3\left(\frac13\right)^3
 +2\left(-\frac12\right)^3+1
 =
 \frac{1-32+4-9+36}{36}=0.                                  \tag{54.22}
\]
After removing common positive Dynkin-index factors,
\[
 {\cal A}_{SU(2)^2U(1)}
 =
 3\left(\frac16\right)-\frac12=0,                            \tag{54.23}
\]
\[
 {\cal A}_{SU(3)^2U(1)}
 =
 2\left(\frac16\right)-\frac23+\frac13=0.                    \tag{54.24}
\]
For \(SU(3)^3\), the two fundamental colour copies in \(Q_L\) are cancelled
by the two antifundamentals \(u_R^c,d_R^c\).  Finally, the three coloured
weak doublets plus \(L_L\) give four weak doublets.
\end{proof}

\begin{corollary}[Local anomaly polynomial is not the remaining arithmetic
obstruction]
\label{cor:v18v44-local-anomaly}
For three generations, every coefficient in (54.19)--(54.20) is still zero.
Thus a nonzero V41 Ward defect cannot be attributed to a failure of the
usual one-generation perturbative Standard-Model charge sums.
\end{corollary}

\begin{proof}
Three identical zero sums remain zero.
\end{proof}

\begin{firewall}[What the arithmetic does not prove]
\label{fire:v18v44-anomaly-scope}
Theorem \ref{thm:v18v44-anomaly-arithmetic} does not construct a
nonperturbative chiral regulator, orient the global determinant/Pfaffian
line, prove reflection positivity, control zero-mode crossings, or show that
the finite regulator Jacobian in V41 vanishes.  It removes only the listed
representation-theoretic anomaly coefficients.
\end{firewall}

\subsection{Fermion contribution to the Ward ledger}

Let \(q_\xi={\cal L}_\xi g\).  On a determinant chart, the fermion amplitude
contribution to the V41 metric score is
\[
 T_{F,n}^{\rm amp}[q_\xi]
 =
 -\operatorname{Re}\operatorname{Tr}
 (D_n^{-1}D_gD_n[q_\xi]).                                   \tag{54.25}
\]
Let \({\cal J}_{F,n}(\xi)\) contain the regulated chiral measure Jacobian
and determinant-line transition connection.  Define
\[
 {\cal A}_{F,n}(\xi)
 =
 T_{F,n}^{\rm amp}[q_\xi]
 +D_{\psi}S_{F,n}[{\cal R}_{n,\xi}\psi]
 -{\cal J}_{F,n}(\xi),                                      \tag{54.26}
\]
before integrating out the fermions; (54.25) is its Schur-complement
representation after fermion integration.

\begin{theorem}[Complete finite-regulator Ward sum]
\label{thm:v18v44-complete-ward-sum}
With scalar, gauge, fermion, ghost/gauge-fixing, measure, and gravitational
regulator pieces retained, the V41 defect decomposes as
\[
\boxed{
 {\cal A}_n(\xi)
 =
 {\cal A}_{H,n}(\xi)
 +{\cal A}_{G,n}(\xi)
 +{\cal A}_{F,n}(\xi)
 +{\cal A}_{{\rm gh+gf},n}(\xi)
 +{\cal A}_{{\rm meas},n}(\xi)
 +{\cal A}_{{\rm grav\,reg},n}(\xi).}                        \tag{54.27}
\]
Only this compiled sum enters the conservation criterion (51.11).
\end{theorem}

\begin{proof}
The complete regulated action and logarithmic Jacobian are sums of the
listed contributions.  Directional differentiation is linear, so (51.3)
splits as (54.27).  The Ward identity itself applies after the sum is
reassembled.
\end{proof}

\begin{corollary}[No sectorwise anomaly norm]
\label{cor:v18v44-no-sectorwise-norm}
Large terms in (54.27) may cancel by the charge identities
(54.19)--(54.24), by ghost/gauge-fixing pairing, or by determinant-line
transition.  Bounding their absolute values separately before assembly can
destroy the cancellation and cannot prove a sharp Ward collar budget.
\end{corollary}

\subsection{Next acquisition}

The fermion metric score and its first three response words are now exact,
and the perturbative charge arithmetic closes.  The next remaining Ward row
is analytic: choose the literal chiral regulator, prove a common inverse
bound or controlled zero-mode split for \(D_n\), assemble (54.27) on each
chart, and bound the chart/collision collar plus complement as in V41.
The determinant phase must be transported independently of the real
amplitude.

\begin{assessment}[V44 acquisition and boundary]
\label{ass:v18v44}
V44 derives the exact determinant/Pfaffian metric response through third
order, gives a gap-dependent response stock, isolates the complete
two-orientation separated fermion loop, and proves the Standard-Model
anomaly arithmetic per generation.  It assembles the correct full Ward
ledger.  A global determinant line, chiral regulator, inverse/zero-mode
control, and vanishing full Ward defect remain unproved.
\end{assessment}

\begin{record}[V44 append record]
\label{rec:v18v44}
V44 is appended to the validated V43 whose installed SHA--256 was
\texttt{FAE34BC2467BF82F42787E01C31309BFACC0BE339C7285FD8CAA8081171B1CA2}.
After installation only V44 is the active numeric SM note; the frozen
\texttt{V100\_f17} archive is unchanged.
\end{record}


\section{V45: the Dirac hard--soft atlas and a compiled Ward-collar criterion}
\label{sec:v18v45}

\subsection{Why a global inverse bound is too strong}

The V44 response stock uses
\(\|D_n^{-1}\|_{\rm op}\leq\gamma_n^{-1}\).  A chiral theory with
topological sectors or physical light modes cannot be expected to have one
uniform \(\gamma_n>0\) on its full configuration space.  The correct
replacement is a singular-value atlas which isolates a finite soft sector
and controls the complement uniformly.

Let
\[
 Q_n=D_n^*D_n\succeq0.                                      \tag{55.1}
\]
For a threshold \(\eta>0\) which is not a singular value, let
\[
 P_{n,\eta}=1_{[0,\eta^2)}(Q_n),
 \qquad
 R_{n,\eta}=I-P_{n,\eta}.                                   \tag{55.2}
\]
The hard inverse is
\[
 G_{n,\eta}^{\rm hard}
 =
 R_{n,\eta}D_n^{-1}R_{n,\eta},                              \tag{55.3}
\]
with the domain/codomain projections transported through the polar
decomposition when \(D_n\) is not normal.

\begin{lemma}[Hard inverse floor]
\label{lem:v18v45-hard-floor}
On the hard singular subspace,
\[
\boxed{
 \|G_{n,\eta}^{\rm hard}\|_{\rm op}\leq\eta^{-1}.}            \tag{55.4}
\]
\end{lemma}

\begin{proof}
Every retained singular value of \(D_n\) is at least \(\eta\).  The
singular values of its inverse on that subspace are their reciprocals.
\end{proof}

\subsection{Exact amplitude splitting}

Assume first that \(D_n\) is invertible.  The determinant amplitude is
\[
 S_{F,n}^{\rm amp}
 =
 -\log|\det D_n|
 =
 -\frac12\log\det Q_n.                                      \tag{55.5}
\]

\begin{theorem}[Soft/hard determinant decomposition]
\label{thm:v18v45-det-split}
For every admissible threshold,
\[
\boxed{
 S_{F,n}^{\rm amp}
 =
 -\frac12\log\det(Q_n|_{\operatorname{Ran}P_{n,\eta}})
 -\frac12\log\det(Q_n|_{\operatorname{Ran}R_{n,\eta}}).}     \tag{55.6}
\]
For a parameter direction in which the rank of \(P_{n,\eta}\) is locally
constant,
\[
\boxed{
\begin{split}
 D S_{F,n}^{\rm amp}
 ={}&
 -\frac12\operatorname{Tr}
 (P_{n,\eta}Q_n^{-1}P_{n,\eta}\,DQ_n)\\
 &-\frac12\operatorname{Tr}
 (R_{n,\eta}Q_n^{-1}R_{n,\eta}\,DQ_n).
\end{split}}                                                 \tag{55.7}
\]
No derivative of the spectral projection remains in the trace.
\end{theorem}

\begin{proof}
The spectral projection \(P_{n,\eta}\) commutes with \(Q_n\), so the
determinant is the product of its soft and hard spectral determinants.  This
proves (55.6).  Jacobi's formula applied to each moving invariant subspace
gives a trace of the compressed logarithmic derivative plus a frame
commutator.  The commutator has zero trace.  Equivalently, differentiate
\[
 -\frac12\operatorname{Tr}\log Q_n
 =
 -\frac12\sum_j\log\lambda_{n,j};                            \tag{55.8}
\]
the eigenvalue derivatives split into the two fixed-rank clusters and yield
(55.7).
\end{proof}

\begin{firewall}[True zero modes]
\label{fire:v18v45-zero-modes}
If \(Q_n\) has a zero eigenvalue, the unsaturated determinant amplitude in
(55.5) is infinite and (55.7) is not valid on that mode.  The corresponding
sector must be removed from the unsaturated vacuum packet or populated by
the declared fermion-source saturation form.  A Moore--Penrose inverse does
not by itself replace the missing determinant weight.
\end{firewall}

\subsection{Stability of the soft chart}

\begin{theorem}[Spectral-projector transport bound]
\label{thm:v18v45-projector-bound}
Let \(Q,Q'\) be positive matrices on one transported Hilbert space.  Suppose
the circle or contour separating the soft cluster from the hard spectrum
has distance at least \(\delta>0\) from both spectra.  Then their soft
spectral projections satisfy
\[
\boxed{
 \|P'-P\|_{\rm op}
 \leq
 \frac{L_\Gamma}{2\pi\delta^2}\,
 \|Q'-Q\|_{\rm op},}                                        \tag{55.9}
\]
where \(L_\Gamma\) is the length of the chosen separating contour.
\end{theorem}

\begin{proof}
The Riesz formula gives
\[
 P'-P
 =
 \frac{1}{2\pi i}
 \int_\Gamma
 (z-Q')^{-1}(Q'-Q)(z-Q)^{-1}\,dz.                           \tag{55.10}
\]
Both resolvents have norm at most \(\delta^{-1}\) on \(\Gamma\).
Taking the operator norm inside the contour integral proves (55.9).
\end{proof}

\begin{corollary}[Fixed-rank cofinal soft bundle]
\label{cor:v18v45-soft-bundle}
If the adjacent defects in the right side of (55.9) are summable and one
common separating contour persists, then the transported soft projections
are Cauchy in operator norm and have constant rank.  They define a cofinal
finite-dimensional soft bundle.
\end{corollary}

\begin{proof}
Sum (55.9) over adjacent cutoffs.  Operator-norm distance below one between
orthogonal projections forces equality of their finite ranks.
\end{proof}

\subsection{Hard response stock}

Compress every metric derivative \(D_i,D_{ij},D_{ijk}\) from V44 to the hard
singular subspace and denote the corresponding norms by
\[
 a_i=\|D_i\|_{\rm F},\quad
 a_{ij}=\|D_{ij}\|_{\rm F},\quad
 a_{ijk}=\|D_{ijk}\|_1.                                     \tag{55.11}
\]

\begin{corollary}[Uniform hard-chart third response]
\label{cor:v18v45-hard-response}
The complete hard contribution to the third determinant response obeys
\[
\boxed{
\begin{split}
 |\partial_{ijk}S_{F,n}^{\rm hard}|
 \leq{}&
 2\eta^{-3}a_ia_ja_k\\
 &+\eta^{-2}(a_{ij}a_k+a_{ik}a_j+a_{jk}a_i)
 +\eta^{-1}a_{ijk}.
\end{split}}                                                 \tag{55.12}
\]
\end{corollary}

\begin{proof}
Use Lemma \ref{lem:v18v45-hard-floor} in the V44 three-response stock.
\end{proof}

The soft response is finite-dimensional but may be large.  It must be
computed from the explicit soft matrix or source-saturation form; it is not
bounded by replacing its small singular values with \(\eta\).

\subsection{Configuration-space collar budget}

Let \(\Gamma_n(\phi)=\sigma_{\min}(D_n(\phi))\), after removing exactly
saturated zero modes.  For a chosen \(\eta_n\), split the regulated
configuration space into
\[
 {\cal C}_n=\{\Gamma_n<\eta_n\},
 \qquad
 {\cal H}_n=\{\Gamma_n\geq\eta_n\}.                           \tag{55.13}
\]
Let \({\cal A}_n(\xi)\) be the \emph{fully compiled} Ward defect of V44.

\begin{theorem}[Singular-value Ward-collar criterion]
\label{thm:v18v45-ward-collar}
Suppose for tests with \(\|\xi\|_{H^{s+1}}\leq1\),
\[
\begin{split}
 \left|
 \mathbb E_n[{\cal A}_n(\xi)1_{{\cal H}_n}]
 \right|
 &\leq\varepsilon_n,\\
 \mathbb E_n\!\left[
 |{\cal A}_n(\xi)|1_{{\cal C}_n}
 \right]
 &\leq u_n.                                                  \tag{55.14}
\end{split}
\]
If
\[
 \varepsilon_n+u_n\longrightarrow0,                          \tag{55.15}
\]
then the Ward-defect descent hypothesis (51.11) holds.
\end{theorem}

\begin{proof}
Split the expectation over (55.13) after assembling all sectors:
\[
 |\mathbb E_n{\cal A}_n(\xi)|
 \leq
 \left|\mathbb E_n[{\cal A}_n(\xi)1_{{\cal H}_n}]\right|
 +
 \mathbb E_n[|{\cal A}_n(\xi)|1_{{\cal C}_n}].
\]
Apply (55.14)--(55.15) and take the dual supremum.
\end{proof}

\begin{corollary}[A usable soft uniform-integrability bound]
\label{cor:v18v45-soft-ui}
If for some \(p>1\),
\[
 \sup_{\|\xi\|_{H^{s+1}}\leq1}
 \mathbb E_n|{\cal A}_n(\xi)|^p\leq M_p^p                  \tag{55.16}
\]
and
\[
 \mathbb P_n({\cal C}_n)\leq p_n,                            \tag{55.17}
\]
then
\[
 u_n\leq M_pp_n^{1-1/p}.                                    \tag{55.18}
\]
Thus \(p_n\to0\), together with the uniform \(L^p\) stock, closes the soft
collar.
\end{corollary}

\begin{proof}
H\"older's inequality gives
\[
 \mathbb E[|{\cal A}|1_{\cal C}]
 \leq
 \|{\cal A}\|_{L^p}
 \mathbb P({\cal C})^{1-1/p}.
\]
Use (55.16)--(55.17).
\end{proof}

\begin{countertheorem}[Small collar probability is insufficient]
\label{cth:v18v45-small-probability}
The condition \(\mathbb P_n({\cal C}_n)\to0\) alone does not imply
\(u_n\to0\).
\end{countertheorem}

\begin{proof}
On a probability space take
\({\cal A}_n=p_n^{-1}1_{{\cal C}_n}\) with
\(\mathbb P({\cal C}_n)=p_n\downarrow0\).  Then
\[
 \mathbb E[|{\cal A}_n|1_{{\cal C}_n}]=1.                   \tag{55.19}
\]
\end{proof}

\subsection{Anomaly cancellation must precede the norm}

Let \(f\) range over all left-handed Weyl multiplets, and suppose the local
hard-chart anomaly has the decomposition
\[
 {\cal A}_{f,n}^{\rm loc}(\xi)
 =
 c_f\,{\cal I}_n(\xi)+r_{f,n}(\xi),                          \tag{55.20}
\]
where the common density \({\cal I}_n\) carries one of the representation
structures in (54.19)--(54.24), and \(c_f\) is the corresponding charge
coefficient.

\begin{theorem}[Compiled anomaly remainder]
\label{thm:v18v45-compiled-anomaly}
For the complete Standard-Model generation,
\[
\boxed{
 \sum_f{\cal A}_{f,n}^{\rm loc}(\xi)
 =
 \sum_fr_{f,n}(\xi).}                                       \tag{55.21}
\]
Therefore
\[
 \left|\sum_f{\cal A}_{f,n}^{\rm loc}(\xi)\right|
 \leq\sum_f|r_{f,n}(\xi)|,                                  \tag{55.22}
\]
with no debit proportional to \(|{\cal I}_n|\).
\end{theorem}

\begin{proof}
Sum (55.20).  The appropriate coefficient sum \(\sum_fc_f\) is zero by the
V44 anomaly arithmetic.  Only after this cancellation is performed is the
triangle inequality applied, giving (55.22).
\end{proof}

\begin{assessment}[Why termwise norms fail]
\label{ass:v18v45-termwise}
The bound
\(\sum_f|c_f{\cal I}_n+r_{f,n}|\) generally retains a large multiple of
\(|{\cal I}_n|\) and loses the anomaly cancellation.  The complement error
\(\varepsilon_n\) in (55.14) must be estimated from the compiled generation,
not from separately absolutized species.
\end{assessment}

\subsection{Cofinal determinant criterion}

\begin{theorem}[Conditional determinant Ward landing]
\label{thm:v18v45-det-landing}
Suppose:
\begin{enumerate}[leftmargin=2.2em]
\item exactly saturated zero modes form a transported finite-rank bundle;
\item a threshold sequence \(\eta_n\) gives the hard inverse floor (55.4);
\item the compiled hard-chart Ward remainder tends to zero;
\item the soft collar satisfies either (55.14) directly or
      (55.16)--(55.18);
\item determinant-line chart transitions and Berry currents cancel in the
      complete anomaly-free representation or converge to an explicitly
      retained current.
\end{enumerate}
Then the fermion contribution satisfies the V41 Ward-defect descent
criterion.  Combined with the scalar, gauge, ghost, measure, and
gravitational rows of (54.27), it yields a conserved continuum stress tensor
whenever their compiled residuals also tend to zero.
\end{theorem}

\begin{proof}
The hard and soft estimates give (55.15), hence Theorem
\ref{thm:v18v45-ward-collar}.  The zero-mode and chart assumptions make the
decomposition globally consistent.  Adding the other compiled rows and
applying the V41 descent theorem gives conservation.
\end{proof}

\begin{firewall}[Conditional status]
\label{fire:v18v45-conditional}
None of the five hypotheses of Theorem
\ref{thm:v18v45-det-landing} is asserted for a literal full SM cutoff family.
The theorem replaces an impossible global-gap request by explicit hard,
soft, saturation, and phase obligations.
\end{firewall}

\subsection{Checkpoint and next step}

V45 reaches the next five-version boundary.  The active YM and NS notes must
now be inspected before V46.  The next mathematical target after that audit
is a joint tightness criterion for the compiled scalar--gauge--fermion stress
packet in a negative Sobolev space, with its Ward defect controlled in the
matching dual topology.

\begin{assessment}[V45 acquisition and boundary]
\label{ass:v18v45}
V45 constructs a Dirac singular-value atlas, proves exact hard/soft
determinant splitting and projector transport, gives the hard response
stock, and proves a soft-collar uniform-integrability criterion.  It also
implements anomaly cancellation before norms.  The required literal
singular-value probability, hard remainder, source saturation, phase, and
full-sector Ward estimates remain unpopulated.
\end{assessment}

\begin{record}[V45 append record]
\label{rec:v18v45}
V45 is appended to the validated V44 whose installed SHA--256 was
\texttt{5F6C411F3B1CEA2D656BB2716495E84AED26A8B38125898E0A3D31890B626BFC}.
After installation only V45 is the active numeric SM note; the frozen
\texttt{V100\_f17} archive is unchanged.
\end{record}


\section{V46: compiled stress tightness and the regular-metric continuum handoff}
\label{sec:v18v46}

\subsection{Mandatory V45 companion audit}

At this checkpoint the active companion files are
\[
\begin{array}{c|c|c|c}
\hbox{programme}&\hbox{version}&\hbox{bytes}&\hbox{SHA--256}\\ \hline
\hbox{Yang--Mills}&V7&124277&
\texttt{46E8AA5EFAAB03E4141A6D7F5AB16D6C5FEE0D0E1F7530CB54CBAED4EEB9DE43}\\
\hbox{Navier--Stokes}&V26&278283&
\texttt{82C887F8C2C69E4F28FAD7C3DC8DE5B9099CB5A4BF431EBA1FFF3E91935978AD}.
\end{array}                                                     \tag{56.1}
\]
The corrected Yang--Mills V7 note defines its complement remainder only
after subtracting the finite-grid reference, so collar and complement form
one exact error identity.  The transferable point is that the \emph{total}
stress covariance and Ward defect must be assembled before their norm is
taken.  Its remaining quadrature target is not imported.  Navier--Stokes
V26 remains unchanged and supplies no stronger common estimate.

\subsection{The compiled random stress packet}

Let \(M\) be a compact four-dimensional reference manifold.  Transport every
regulated sector stress to one fixed real tensor bundle and physical unit,
and write
\[
\boxed{
 {\mathsf T}_n
 =
 {\mathsf T}_{H,n}
 +{\mathsf T}_{G,n}
 +{\mathsf T}_{F,n}
 +{\mathsf T}_{{\rm gh+gf},n}
 +{\mathsf T}_{{\rm meas},n}
 +{\mathsf T}_{{\rm grav\,reg},n}.}                          \tag{56.2}
\]
Let
\[
 \widetilde{\mathsf T}_n
 =
 {\mathsf T}_n-\mathbb E{\mathsf T}_n.                       \tag{56.3}
\]
The cross covariances in (56.2) are retained.

Let \((e_k)_{k\geq0}\) be an orthonormal eigenbasis of a fixed positive
tensor Laplacian,
\[
 (I+\Delta)e_k=(1+\lambda_k)e_k.                             \tag{56.4}
\]
Define the compiled modal variance
\[
 V_{n,k}
 =
 \mathbb E\left|
 \langle\widetilde{\mathsf T}_n,e_k\rangle
 \right|^2.                                                  \tag{56.5}
\]

\begin{theorem}[Negative-Sobolev stress tightness]
\label{thm:v18v46-stress-tightness}
Suppose for some \(\alpha\geq0\),
\[
 V_{n,k}\leq C(1+\lambda_k)^\alpha                          \tag{56.6}
\]
uniformly in \(n,k\), and
\[
 \sup_n
 \|\mathbb E{\mathsf T}_n\|_{H^{-s_0}}<\infty               \tag{56.7}
\]
for one \(s_0>\alpha+2\).  Then
\[
 \sup_n\mathbb E
 \|{\mathsf T}_n\|_{H^{-s_0}}^2<\infty,                     \tag{56.8}
\]
and the laws of \({\mathsf T}_n\) are tight in \(H^{-s}\) for every
\(s>s_0\).
\end{theorem}

\begin{proof}
By (56.5)--(56.6),
\[
\begin{split}
 \mathbb E
 \|\widetilde{\mathsf T}_n\|_{H^{-s_0}}^2
 &=
 \sum_k(1+\lambda_k)^{-s_0}V_{n,k}\\
 &\leq
 C\sum_k(1+\lambda_k)^{-(s_0-\alpha)}.                       \tag{56.9}
\end{split}
\]
Weyl summability in dimension four makes the last series finite precisely
when \(s_0-\alpha>2\).  Add the deterministic mean bound to obtain (56.8).
For \(R>0\), Markov's inequality puts probability at least \(1-C/R^2\) in
the radius-\(R\) ball of \(H^{-s_0}\).  The embedding
\(H^{-s_0}\hookrightarrow H^{-s}\) is compact for \(s>s_0\), so these balls
have compact closure in \(H^{-s}\).  This proves tightness.
\end{proof}

\begin{corollary}[White-noise stress threshold]
\label{cor:v18v46-white-threshold}
If the compiled covariance is uniformly white at high frequency, so that
\(\alpha=0\), the stress packet is tight in every \(H^{-s}\) with
\(s>2\), after choosing an intermediate \(s_0\in(2,s)\).
\end{corollary}

\subsection{Why sectorwise bounds are not the criterion}

Writing sector labels as \(a,b\), the exact modal variance is
\[
 V_{n,k}
 =
 \sum_{a,b}
 \operatorname{Cov}\left(
 \langle{\mathsf T}_{a,n},e_k\rangle,
 \langle{\mathsf T}_{b,n},e_k\rangle
 \right).                                                    \tag{56.10}
\]

\begin{proposition}[Compiled covariance advantage]
\label{prop:v18v46-compiled-covariance}
Condition (56.6) may hold even when the sum of absolute sector covariance
entries grows faster than \((1+\lambda_k)^\alpha\).  It is sufficient to
certify the positive quadratic form (56.10) directly.  A sectorwise estimate
\[
 V_{n,k}\leq m\sum_a
 \mathbb E|\langle\widetilde{\mathsf T}_{a,n},e_k\rangle|^2  \tag{56.11}
\]
is valid but can lose exact gauge--ghost, anomaly, and force-exchange
cancellations.
\end{proposition}

\begin{proof}
Equation (56.10) is the variance of the sum and includes signed cross
terms.  Cauchy--Schwarz proves (56.11), but replacing every cross term by
its absolute value can strictly increase the bound.  The V43 force exchange
and V44--V45 anomaly compilation exhibit mechanisms producing such signed
cross terms.
\end{proof}

\subsection{Passage of the Ward identity}

Let
\[
 {\cal D}:H^{-s}(S^2T^*M)\longrightarrow
 H^{-(s+1)}(T^*M),
 \qquad {\cal D}T=\operatorname{div}_gT                      \tag{56.12}
\]
be the distributional divergence for a fixed smooth metric \(g\).

\begin{theorem}[Closed Ward graph]
\label{thm:v18v46-closed-ward}
Suppose
\[
 {\mathsf T}_n\Rightarrow{\mathsf T}
 \quad\hbox{in }H^{-s},                                     \tag{56.13}
\]
and there are random Ward residuals \({\mathsf A}_n\in H^{-(s+1)}\) such
that
\[
 {\cal D}{\mathsf T}_n={\mathsf A}_n,
 \qquad
 {\mathsf A}_n\longrightarrow0
 \quad\hbox{in probability in }H^{-(s+1)}.                  \tag{56.14}
\]
Then
\[
\boxed{{\cal D}{\mathsf T}=0}
\quad\hbox{almost surely.}                                  \tag{56.15}
\]
The same statement for deterministic means follows with ordinary weak or
strong convergence.
\end{theorem}

\begin{proof}
The map \({\cal D}\) is continuous between the displayed Sobolev spaces.
The continuous mapping theorem sends (56.13) to
\({\cal D}{\mathsf T}_n\Rightarrow{\cal D}{\mathsf T}\).
By (56.14), the same variables converge in probability, hence in law, to
zero.  Uniqueness of distributional limits gives
\({\cal D}{\mathsf T}=0\) almost surely.
\end{proof}

\begin{assessment}[Operator versus mean Ward identity]
\label{ass:v18v46-operator-mean}
V41 proves the mean identity directly.  The stronger random form (56.14)
requires inserted Ward identities dense enough to identify the
distribution-valued divergence.  Mean conservation alone does not imply
almost-sure conservation of a random stress realization.
\end{assessment}

\subsection{Passage of a non-Gaussian stress witness}

Let \(q_1,q_2,q_3\) be the separated metric tests of V40, and set
\[
 X_{n,i}=\langle{\mathsf T}_n,q_i\rangle.                    \tag{56.16}
\]

\begin{theorem}[Cubic-witness stability]
\label{thm:v18v46-cubic-stability}
Suppose
\[
 (X_{n,1},X_{n,2},X_{n,3})
 \Rightarrow
 (X_1,X_2,X_3),                                              \tag{56.17}
\]
and for some \(\varepsilon>0\),
\[
 \sup_n\mathbb E|X_{n,i}|^{3+\varepsilon}<\infty
 \qquad(i=1,2,3).                                            \tag{56.18}
\]
Then
\[
 \operatorname{Cum}(X_{n,1},X_{n,2},X_{n,3})
 \longrightarrow
 \operatorname{Cum}(X_1,X_2,X_3).                            \tag{56.19}
\]
In particular, a nonzero limiting direct-minus-Gaussian cubic excess
survives in the continuum and excludes a quasi-free limit, provided the
Gaussian comparator converges on the same follower-shorted packet.
\end{theorem}

\begin{proof}
H\"older's inequality and (56.18) give uniform integrability of every
monomial of total degree at most three in the three variables.  Convergence
in law therefore implies convergence of the first, second, and third mixed
moments.  The third cumulant is a fixed polynomial in those moments, proving
(56.19).  A quasi-free limit has the matched Wick value, so a nonzero excess
excludes it.
\end{proof}

\begin{firewall}[Tightness does not create interaction]
\label{fire:v18v46-tightness-not-interaction}
The second-moment bound (56.6) establishes compactness only.  It does not
give the positive cubic floor required by V40.  Conversely, a finite-cutoff
positive cubic value does not survive without the uniform integrability
(56.18) and convergence of its comparator.
\end{firewall}

\subsection{Regular random-metric branch}

The V41 Einstein handoff used deterministic metrics.  The following gives a
rigorous random-metric branch at regularity high enough for the nonlinear
Einstein tensor to be continuous.

\begin{theorem}[Joint regular-metric compactness]
\label{thm:v18v46-joint-compactness}
Let \(0<\beta<\alpha<1\).  Suppose the random metrics satisfy, with
probability tending uniformly to one on large bounded sets,
\[
 g_n\in C^{2,\alpha},\qquad
 c\,g_0\leq g_n\leq C\,g_0,                                 \tag{56.20}
\]
and their laws are tight in \(C^{2,\beta}\).  Suppose also that
\({\mathsf T}_n\) is tight in \(H^{-s}\).  Then the joint laws of
\((g_n,{\mathsf T}_n)\) are tight in
\[
 C^{2,\beta}\times H^{-s}.                                  \tag{56.21}
\]
Along every convergent subsequence,
\[
 G(g_n)+\Lambda g_n-\kappa{\mathsf T}_n
 \Rightarrow
 G(g)+\Lambda g-\kappa{\mathsf T}
 \quad\hbox{in }H^{-s}.                                     \tag{56.22}
\]
\end{theorem}

\begin{proof}
Tightness of both marginals gives tightness of the joint laws on the product
of two Polish spaces.  The uniform ellipticity in (56.20) persists on
compact subsets of the metric space.  In local coordinates, the Einstein
tensor is a smooth algebraic expression in
\(g^{-1},\partial g,\partial^2g\).  Therefore
\[
 g\longmapsto G(g)+\Lambda g
\]
is continuous from uniformly elliptic \(C^{2,\beta}\) metrics to
\(C^{0,\beta}\), hence to \(H^{-s}\).  Combine this with the continuous
linear stress term and apply the continuous mapping theorem.
\end{proof}

\begin{corollary}[Conditional random Einstein equation]
\label{cor:v18v46-random-einstein}
Under Theorem \ref{thm:v18v46-joint-compactness}, if the residuals
\[
 {\mathsf E}_n
 =
 G(g_n)+\Lambda g_n-\kappa{\mathsf T}_n                     \tag{56.23}
\]
converge to zero in probability in \(H^{-s}\), then every subsequential
limit satisfies
\[
\boxed{
 G(g)+\Lambda g=\kappa{\mathsf T}}
\quad\hbox{almost surely in }H^{-s}.                         \tag{56.24}
\]
If the matching Ward/Bianchi residuals vanish, the two sides are
distributionally conserved.
\end{corollary}

\begin{proof}
Equation (56.22) identifies the law of the residual limit, while the
hypothesis identifies it with zero.  The conservation statement follows
from Theorem \ref{thm:v18v46-closed-ward} and the Bianchi analogue.
\end{proof}

\begin{firewall}[Regularity branch]
\label{fire:v18v46-regularity-branch}
The \(C^{2,\beta}\) metric hypothesis is a strong regular branch.  V46 does
not prove it for a quantum gravitational measure.  If the metric limit is
rougher, a weak curvature compactness theorem and a definition of the
nonlinear Einstein tensor are required before (56.22) is meaningful.
\end{firewall}

\subsection{Concrete remaining stocks}

The continuum handoff now reduces to measurable quantities:
\begin{enumerate}[leftmargin=2.2em]
\item the compiled modal variance \(V_{n,k}\) and mean stock
      (56.6)--(56.7);
\item the dense inserted Ward residual in the matching
      \(H^{-(s+1)}\) topology;
\item the \(L^{3+\varepsilon}\) stock for a separated cubic witness;
\item either regular metric tightness (56.20) or a weaker curvature
      compactness replacement; and
\item the actual Einstein residual norm (56.23).
\end{enumerate}
The next note should derive a practical sufficient covariance bound for
(56.6) from a source kernel or heat-semigroup estimate, rather than assume
modal control independently at every eigenfunction.

\begin{assessment}[V46 acquisition and boundary]
\label{ass:v18v46}
V46 records the mandatory audit, proves a compiled negative-Sobolev
tightness criterion for the full stress packet, passes Ward conservation and
a cubic interaction witness to the continuum, and proves a conditional
joint regular-metric Einstein handoff.  None of the required physical modal
variance, dense Ward, cubic UI, metric tightness, or Einstein residual stocks
has yet been populated for the literal full cutoff family.
\end{assessment}

\begin{record}[V46 append record]
\label{rec:v18v46}
V46 is appended to the validated V45 whose installed SHA--256 was
\texttt{91385E1121EB07FA7FEEB284CAD78D791E5F2C170054E49C1C481C36F36BC9A9}.
After installation only V46 is the active numeric SM note; the frozen
\texttt{V100\_f17} archive is unchanged.
\end{record}


\section{V47: heat-smoothed covariance stocks for the compiled stress packet}
\label{sec:v18v47}

\subsection{From modal assumptions to one operator estimate}

Let
\[
 L=I+\Delta
\]
be the positive tensor Laplacian used in V46, and let \(C_n^T\) be the
covariance operator of the compiled centered stress:
\[
 \langle f,C_n^Tg\rangle
 =
 \operatorname{Cov}
 (\langle\widetilde{\mathsf T}_n,f\rangle,
  \langle\widetilde{\mathsf T}_n,g\rangle).                  \tag{57.1}
\]
It is a positive quadratic form on smooth tensor tests.  For \(0<t\leq1\),
define the heat-smoothed covariance
\[
 C_{n,t}^T
 =
 e^{-tL/2}C_n^Te^{-tL/2}.                                   \tag{57.2}
\]

\begin{theorem}[Heat stock implies the V46 modal bound]
\label{thm:v18v47-heat-to-modal}
Suppose, for some \(\alpha\geq0\),
\[
\boxed{
 \|C_{n,t}^T\|_{\rm op}
 \leq A\,t^{-\alpha}
 \qquad(0<t\leq1)}                                          \tag{57.3}
\]
uniformly in \(n\).  Then every Laplace eigentensor \(e_k\) satisfies
\[
\boxed{
 \mathbb E|\langle\widetilde{\mathsf T}_n,e_k\rangle|^2
 \leq eA(1+\lambda_k)^\alpha.}                              \tag{57.4}
\]
Consequently the stress laws are tight in \(H^{-s}\) for every
\(s>\alpha+2\), provided the mean has the matching bound.
\end{theorem}

\begin{proof}
Since \(Le_k=(1+\lambda_k)e_k\),
\[
 \langle e_k,C_{n,t}^Te_k\rangle
 =
 e^{-t(1+\lambda_k)}
 \langle e_k,C_n^Te_k\rangle.                               \tag{57.5}
\]
Use (57.3) and choose \(t=(1+\lambda_k)^{-1}\leq1\).  Then
\[
 \langle e_k,C_n^Te_k\rangle
 \leq
 A t^{-\alpha}e^{t(1+\lambda_k)}
 =
 eA(1+\lambda_k)^\alpha.
\]
The final statement is Theorem \ref{thm:v18v46-stress-tightness}.
\end{proof}

\begin{corollary}[Uniform covariance operator]
\label{cor:v18v47-white}
If \(\sup_n\|C_n^T\|_{\rm op}<\infty\), then (57.3) holds with
\(\alpha=0\), and the compiled stress is tight in every \(H^{-s}\),
\(s>2\).
\end{corollary}

\begin{proof}
The heat semigroup is an \(L^2\) contraction, so
\(\|e^{-tL/2}C_n^Te^{-tL/2}\|_{\rm op}\leq\|C_n^T\|_{\rm op}\).
\end{proof}

\subsection{A physical-space Schur criterion}

Assume \(C_{n,t}^T\) has a matrix-valued integral kernel
\(K_{n,t}(x,y)\) relative to the fixed tensor bundle.  Let
\(\|\cdot\|_{\rm op,fib}\) denote the fibre operator norm.

\begin{theorem}[Heat-kernel Schur stock]
\label{thm:v18v47-schur}
If
\[
\sup_x\int_M
\|K_{n,t}(x,y)\|_{\rm op,fib}\,dy
\leq A t^{-\alpha},                                        \tag{57.6}
\]
and
\[
\sup_y\int_M
\|K_{n,t}(x,y)\|_{\rm op,fib}\,dx
\leq A t^{-\alpha},                                        \tag{57.7}
\]
then (57.3) holds.
\end{theorem}

\begin{proof}
For scalar nonnegative majorant
\(k(x,y)=\|K_{n,t}(x,y)\|_{\rm op,fib}\), the vector-valued Schur test gives
\[
 \|C_{n,t}^T\|_{\rm op}
 \leq
 \left(
 \sup_x\int k(x,y)\,dy
 \right)^{1/2}
 \left(
 \sup_y\int k(x,y)\,dx
 \right)^{1/2}.                                             \tag{57.8}
\]
Apply (57.6)--(57.7).
\end{proof}

This test allows contact singularities: heat smoothing is performed before
the absolute kernel integral.  Taking the absolute value of an unsmoothed
distributional covariance kernel would generally be meaningless.

\subsection{Source-factor criterion}

Suppose the compiled stress is generated from a centered innovation
\(\eta_n\) in a real Hilbert space \({\cal U}_n\) by a deterministic source
operator
\[
 B_n:{\cal U}_n\longrightarrow{\cal D}'(S^2T^*M),
 \qquad
 \widetilde{\mathsf T}_n=B_n\eta_n,                          \tag{57.9}
\]
and the innovation covariance obeys
\[
 0\preceq\operatorname{Cov}(\eta_n)\preceq I.                \tag{57.10}
\]

\begin{theorem}[Smoothed source-synthesis stock]
\label{thm:v18v47-source-factor}
If
\[
\boxed{
 \|e^{-tL/2}B_n\|_{\rm op}^2
 \leq A t^{-\alpha}
 \qquad(0<t\leq1),}                                         \tag{57.11}
\]
then the heat covariance bound (57.3) holds.
\end{theorem}

\begin{proof}
The covariance factors as
\[
 C_n^T
 =
 B_n\operatorname{Cov}(\eta_n)B_n^*.                        \tag{57.12}
\]
Hence, using (57.10),
\[
\begin{split}
 \|e^{-tL/2}C_n^Te^{-tL/2}\|_{\rm op}
 &\leq
 \|e^{-tL/2}B_n\|_{\rm op}^2\\
 &\leq At^{-\alpha}.
\end{split}
\]
\end{proof}

\begin{corollary}[Finite follower banks]
\label{cor:v18v47-follower-bank}
If
\[
 B_n=\sum_{a=1}^mB_{a,n}                                    \tag{57.13}
\]
is a fixed finite common source bank, it is sufficient to estimate the
compiled operator in (57.11).  The sectorwise majorant
\[
 \|e^{-tL/2}B_n\|_{\rm op}^2
 \leq
 m\sum_{a=1}^m
 \|e^{-tL/2}B_{a,n}\|_{\rm op}^2                            \tag{57.14}
\]
is valid but may lose source-follower cancellations.
\end{corollary}

\begin{proof}
The first statement is Theorem \ref{thm:v18v47-source-factor}.
Cauchy--Schwarz for the finite operator sum gives (57.14).
\end{proof}

\subsection{Direct covariance-form criterion}

The linear source factorization (57.9) may be unavailable for a nonlinear
stress packet.  One can instead work with its covariance form.  Define the
heat-smoothed test
\[
 f_t=e^{-tL/2}f.                                             \tag{57.15}
\]

\begin{theorem}[Smoothed pairing criterion]
\label{thm:v18v47-pairing}
The operator bound (57.3) is equivalent to
\[
\boxed{
 \operatorname{Var}
 \langle\widetilde{\mathsf T}_n,f_t\rangle
 \leq
 At^{-\alpha}\|f\|_{L^2}^2}                                 \tag{57.16}
\]
for every smooth tensor \(f\), uniformly in \(n\).
\end{theorem}

\begin{proof}
By (57.1)--(57.2), the left side is
\[
 \langle f,C_{n,t}^Tf\rangle.
\]
Since \(C_{n,t}^T\) is positive self-adjoint, the supremum of this quadratic
form over unit vectors is its operator norm.
\end{proof}

Thus a Brascamp--Lieb, spectral-gap, martingale, or source-response estimate
may be applied directly to the single smoothed stress pairing.

\subsection{A conditional Poincar\'e route}

Let the regulated field law satisfy the Poincar\'e inequality
\[
 \operatorname{Var}_{\mu_n}(F)
 \leq
 \rho_n^{-1}
 \mathbb E_{\mu_n}\|\nabla_\phi F\|_{{\cal H}_{\phi,n}}^2.   \tag{57.17}
\]

\begin{theorem}[Poincar\'e-to-heat stock]
\label{thm:v18v47-poincare}
If the compiled stress satisfies
\[
 \mathbb E_{\mu_n}
 \left\|
 \nabla_\phi
 \langle\widetilde{\mathsf T}_n,e^{-tL/2}f\rangle
 \right\|_{{\cal H}_{\phi,n}}^2
 \leq
 A\rho_n t^{-\alpha}\|f\|_{L^2}^2,                          \tag{57.18}
\]
then (57.16), hence V46 stress tightness, holds.
\end{theorem}

\begin{proof}
Apply (57.17) to the smoothed stress pairing and then use (57.18).
\end{proof}

\begin{assessment}[Convexity branch and interaction branch]
\label{ass:v18v47-convexity-tension}
The V34 convexity certificate can supply a Poincar\'e constant on a
finite-regulator branch.  To use it here, one still needs the compiled
gradient stock (57.18).  If the same certificate also gives all-order
anchored cumulant summability for the residual stress, V36 forces its
non-Gaussian vertices to vanish.  Tightness and interaction survival
therefore require separate estimates and must not be identified.
\end{assessment}

\subsection{Transport stability}

Let \(U_{n+1,n}\) transport tensor tests between adjacent cutoffs, and define
the covariance transport defect
\[
 {\cal R}_{n}^{C}
 =
 C_{n+1}^T
 -
 U_{n+1,n}C_n^TU_{n+1,n}^*.                                \tag{57.19}
\]

\begin{proposition}[Adjacent heat-stock transport]
\label{prop:v18v47-transport}
Assume \(U_{n+1,n}\) commutes with \(L\) on the common test core and is an
\(L^2\) contraction.  Then
\[
\begin{split}
 \|e^{-tL/2}C_{n+1}^Te^{-tL/2}\|_{\rm op}
 \leq{}&
 \|e^{-tL/2}C_n^Te^{-tL/2}\|_{\rm op}\\
 &+
 \|e^{-tL/2}{\cal R}_n^Ce^{-tL/2}\|_{\rm op}.                \tag{57.20}
\end{split}
\]
If the second terms have a summable uniform bound
\[
 \sum_n
 \sup_{0<t\leq1}
 t^\alpha
 \|e^{-tL/2}{\cal R}_n^Ce^{-tL/2}\|_{\rm op}<\infty,         \tag{57.21}
\]
one finite-level heat stock propagates uniformly along the cofinal tail.
\end{proposition}

\begin{proof}
Insert (57.19), commute the heat operators through \(U_{n+1,n}\), and use
the contraction property for the transported old covariance.  This gives
(57.20).  Multiply by \(t^\alpha\), take the supremum, and telescope using
(57.21).
\end{proof}

\subsection{Concrete V48 target}

The modal covariance row of V46 can now be populated by any one of:
\[
\boxed{
 \text{heat operator bound (57.3), Schur stock (57.6)--(57.7),
 source stock (57.11), or gradient stock (57.18).}}          \tag{57.22}
\]
The next unresolved compactness row is the
\(L^{3+\varepsilon}\) uniform integrability needed to preserve the V40
cubic interaction witness.  V48 should derive it for finite Gaussian chaos
packets by hypercontractivity and state the additional weighted-gradient
condition needed for the direct interacting measure.

\begin{assessment}[V47 acquisition and boundary]
\label{ass:v18v47}
V47 replaces the mode-by-mode V46 covariance hypothesis by equivalent or
sufficient heat-smoothed operator, Schur-kernel, source-factor, and
Poincar\'e gradient stocks, and proves their cutoff transport.  No one of
these stocks is yet evaluated on the complete literal SM stress source.
\end{assessment}

\begin{record}[V47 append record]
\label{rec:v18v47}
V47 is appended to the validated V46 whose installed SHA--256 was
\texttt{5A006013F4ACD6B03CDB337F3602BA5B2F739BD55BF3B4DF09ADE4F8F26B93A9}.
After installation only V47 is the active numeric SM note; the frozen
\texttt{V100\_f17} archive is unchanged.
\end{record}


\section{V48: cubic uniform integrability for Gaussian and direct stress packets}
\label{sec:v18v48}

\subsection{Finite-chaos comparator}

Let \((\Omega,\gamma)\) carry an isonormal Gaussian process, and let
\[
 F=\sum_{k=0}^{d}I_k(f_k)                                   \tag{58.1}
\]
belong to the sum of homogeneous Wiener chaoses through degree \(d\).
The Higgs comparator of V42 has \(d=4\); the quadratic gauge comparator of
V43 has \(d=2\).

\begin{theorem}[Finite-chaos hypercontractive stock]
\label{thm:v18v48-hypercontractive}
For every \(p\geq2\),
\[
\boxed{
 \|F-\mathbb EF\|_{L^p(\gamma)}
 \leq
 (p-1)^{d/2}
 \|F-\mathbb EF\|_{L^2(\gamma)}.}                           \tag{58.2}
\]
Consequently, a uniform covariance bound for a finite test bank gives
uniform moments of every fixed order for its Gaussian comparator.
\end{theorem}

\begin{proof}
Let \(P_t\) be the Ornstein--Uhlenbeck semigroup.  It acts on the \(k\)-th
chaos by multiplication with \(e^{-kt}\).  Nelson
hypercontractivity gives
\[
 \|P_tG\|_p\leq\|G\|_2
 \quad\hbox{when}\quad
 e^{2t}\geq p-1.                                             \tag{58.3}
\]
For \(F-\mathbb EF=\sum_{k=1}^dI_k(f_k)\), apply (58.3) to
\[
 G=\sum_{k=1}^de^{kt}I_k(f_k).
\]
Then \(P_tG=F-\mathbb EF\), while orthogonality of chaoses gives
\[
 \|G\|_2^2
 =
 \sum_{k=1}^de^{2kt}\|I_k(f_k)\|_2^2
 \leq
 e^{2dt}\|F-\mathbb EF\|_2^2.                               \tag{58.4}
\]
Choose \(e^{2t}=p-1\) to obtain (58.2).
\end{proof}

\begin{corollary}[V42 comparator uniform integrability]
\label{cor:v18v48-higgs-ui}
For the complete degree-two/four Gaussian Higgs stress score
\[
 T_{H,n}^{\rm G}[q]
 =
 I_2(B_n^H[q])+I_4(D_n^H[q]),
\]
one has
\[
 \|T_{H,n}^{\rm G}[q]\|_{L^p}
 \leq
 (p-1)^2
 \|T_{H,n}^{\rm G}[q]\|_{L^2}.                              \tag{58.5}
\]
Thus a uniform covariance bound closes the
\(L^{3+\varepsilon}\) requirement of V46 for every fixed Gaussian test
bank.
\end{corollary}

\begin{proof}
Apply Theorem \ref{thm:v18v48-hypercontractive} with \(d=4\).
\end{proof}

\begin{corollary}[Follower-shorted Gaussian packet]
\label{cor:v18v48-gaussian-short}
Orthogonal projection onto or away from a fixed sum of homogeneous Gaussian
chaoses does not increase the \(L^2\) norm and does not increase the maximum
chaos degree.  Hence (58.2) holds for the V39--V42 follower-shorted Gaussian
packet with the same degree \(d\).
\end{corollary}

\begin{proof}
The follower projection acts orthogonally within each homogeneous chaos.
Apply (58.2) to the projected kernel.
\end{proof}

\subsection{An \(L^p\) gradient criterion for the direct measure}

Let \(\mu_n\) be the full finite interacting regulator measure on its
field-coordinate manifold, and assume the logarithmic Sobolev inequality
\[
 \operatorname{Ent}_{\mu_n}(u^2)
 \leq
 \frac{2}{\rho_n}
 \mathbb E_{\mu_n}\|\nabla_\phi u\|^2.                       \tag{58.6}
\]

\begin{lemma}[\(L^p\) consequence of logarithmic Sobolev]
\label{lem:v18v48-lp-gradient}
For every locally Lipschitz \(F\) and \(p\geq2\),
\[
\boxed{
 \|F-\mathbb EF\|_{L^p(\mu_n)}
 \leq
 \sqrt{\frac{2p}{\rho_n}}\,
 \bigl\|\|\nabla_\phi F\|\bigr\|_{L^p(\mu_n)}.}              \tag{58.7}
\]
\end{lemma}

\begin{proof}
First apply (58.6) to the smooth truncations
\[
 u_\varepsilon=(|F-\mathbb EF|^2+\varepsilon)^{p/4}.
\]
The chain rule bounds
\[
 \|\nabla u_\varepsilon\|
 \leq
 \frac p2
 (|F-\mathbb EF|^2+\varepsilon)^{(p-2)/4}
 \|\nabla F\|.                                               \tag{58.8}
\]
The entropy inequality for successive exponents, integrated from \(2\) to
\(p\), yields
\[
 \|F-\mathbb EF\|_p^2
 \leq
 \frac{2p}{\rho_n}
 \|\|\nabla F\|\|_p^2.                                      \tag{58.9}
\]
This is the standard \(L^p\) logarithmic-Sobolev interpolation; (58.8)
justifies it for the truncations.  Letting the truncation level tend to
infinity and then \(\varepsilon\downarrow0\) proves (58.7).
\end{proof}

\begin{theorem}[Direct stress uniform-integrability stock]
\label{thm:v18v48-direct-ui}
Let \(p=3+\varepsilon>3\).  If, for every test \(q\) in the frozen cubic
bank,
\[
\boxed{
 \bigl\|
 \|\nabla_\phi\langle{\mathsf T}_n,q\rangle\|
 \bigr\|_{L^p(\mu_n)}
 \leq
 C_q\sqrt{\rho_n},}                                         \tag{58.10}
\]
then
\[
 \sup_n
 \|\langle\widetilde{\mathsf T}_n,q\rangle\|_{L^p}
 \leq
 \sqrt{2p}\,C_q.                                             \tag{58.11}
\]
Therefore the direct cubic stress cumulants are uniformly integrable and
the V46 cubic-witness stability theorem applies.
\end{theorem}

\begin{proof}
Use Lemma \ref{lem:v18v48-lp-gradient} with
\(F=\langle{\mathsf T}_n,q\rangle\), then insert (58.10).
\end{proof}

\begin{assessment}[Degenerating log-Sobolev constant]
\label{ass:v18v48-lsi-degeneration}
The factor \(\sqrt{\rho_n}\) in (58.10) is essential.  If the
logarithmic-Sobolev constant tends to zero on a scale-critical branch, a
uniform unweighted gradient moment does not imply (58.11).  The gradient
must shrink at the matching rate or a different uniform-integrability
argument is needed.
\end{assessment}

\subsection{Literal scalar gradient}

For the Higgs metric score (52.4), differentiation in the scalar field gives
\[
\boxed{
 \nabla_hS_{H,n}[q]
 =
 \dot K_n[q]h
 +
 4\lambda_n
 \bigl(
 \dot w_{n,z}[q]|h_z|^2h_z
 \bigr)_{z\in\Lambda_n}.}                                   \tag{58.12}
\]

\begin{proposition}[Scalar \(L^p\) gradient bound]
\label{prop:v18v48-scalar-gradient}
Let
\[
 k_{n,q}=\|\dot K_n[q]\|_{\rm op},
 \qquad
 w_{n,q}=\sup_z|\dot w_{n,z}[q]|.
\]
Then
\[
\boxed{
 \|\nabla_hS_{H,n}[q]\|
 \leq
 k_{n,q}\|h\|
 +
 4|\lambda_n|w_{n,q}
 \left(\sum_z|h_z|^6\right)^{1/2}.}                          \tag{58.13}
\]
Thus (58.10) for the scalar row follows from the corresponding
\(L^p\) moments of \(\|h\|\) and
\((\sum_z|h_z|^6)^{1/2}\), with the physical coefficients included.
\end{proposition}

\begin{proof}
Apply the operator norm to the first term of (58.12).  The squared norm of
the second term is bounded by
\[
 16\lambda_n^2w_{n,q}^2\sum_z|h_z|^6.
\]
Use the triangle inequality.
\end{proof}

\subsection{Compact gauge-link gradient}

For the plaquette score (53.3), let \(\nabla_{U_e}\) be the invariant
gradient on the compact link group.  Each link occurs in a uniformly finite
number of plaquettes on a bounded-valence regulator.

\begin{proposition}[Gauge metric-score gradient]
\label{prop:v18v48-gauge-gradient}
If
\[
 \sup_{p,e,U}\|\nabla_{U_e}\ell_p(U)\|\leq c_R              \tag{58.14}
\]
and each link lies in at most \(N_p\) plaquettes, then
\[
\boxed{
 \|\nabla_US_{G,n}[q]\|^2
 \leq
 c_R^2N_p
 \sum_p|\dot\beta_{n,p}[q]|^2.}                             \tag{58.15}
\]
The gauge row of (58.10) is therefore deterministic once the physical
coefficient stock is bounded.
\end{proposition}

\begin{proof}
For each link,
\[
 \nabla_{U_e}S_{G,n}[q]
 =
 \sum_{p\ni e}\dot\beta_{n,p}[q]\nabla_{U_e}\ell_p.
\]
Cauchy--Schwarz over at most \(N_p\) plaquettes and summation over links give
(58.15), after absorbing the uniformly bounded number of links per
plaquette into \(N_p\).
\end{proof}

\subsection{Fermion resolvent gradient}

Let
\[
 F_{F,n}[q]
 =
 -\operatorname{Re}\operatorname{Tr}(G_nD_q)                \tag{58.16}
\]
be the determinant-amplitude score from V44.  For a field-coordinate
direction \(v\), set
\[
 D_v=D_\phi D_n[v],
 \qquad
 D_{qv}=D_\phi D_q[v].
\]

\begin{proposition}[Fermion score derivative]
\label{prop:v18v48-fermion-gradient}
One has
\[
\boxed{
 D_\phi F_{F,n}[q][v]
 =
 \operatorname{Re}\operatorname{Tr}(G_nD_vG_nD_q)
 -
 \operatorname{Re}\operatorname{Tr}(G_nD_{qv}).}             \tag{58.17}
\]
On a hard singular-value chart with floor \(\eta_n\),
\[
 |D_\phi F_{F,n}[q][v]|
 \leq
 \eta_n^{-2}\|D_v\|_{\rm F}\|D_q\|_{\rm F}
 +
 \eta_n^{-1}\|D_{qv}\|_1.                                  \tag{58.18}
\]
\end{proposition}

\begin{proof}
Differentiate (58.16) and use
\(D_\phi G_n[v]=-G_nD_vG_n\).  This gives (58.17).
The Frobenius and trace duality bounds used in V44, together with the hard
inverse floor, give (58.18).
\end{proof}

On the soft collar, (58.18) is unavailable.  The V45
uniform-integrability debit must control the full compiled gradient, not
only the determinant score itself.

\subsection{Follower projection on the direct packet}

Let \({\mathfrak P}_n\) be the direct follower projection used in V40.  An
orthogonal \(L^2\) projection need not be a contraction on \(L^p\),
\(p>2\).

\begin{theorem}[Projected direct UI]
\label{thm:v18v48-projected-ui}
Suppose
\[
 \sup_n
 \|I-{\mathfrak P}_n\|_{L^p(\mu_n)\to L^p(\mu_n)}
 \leq M_p<\infty.                                           \tag{58.19}
\]
Then (58.11) implies
\[
 \sup_n
 \|(I-{\mathfrak P}_n)
 \langle\widetilde{\mathsf T}_n,q\rangle\|_{L^p}
 \leq
 M_p\sqrt{2p}\,C_q.                                         \tag{58.20}
\]
If \({\mathfrak P}_n\) is conditional expectation onto a follower
\(\sigma\)-algebra, one may take
\[
 M_p\leq2.                                                   \tag{58.21}
\]
\end{theorem}

\begin{proof}
Apply the operator bound (58.19) to (58.11).  Conditional expectation is an
\(L^p\) contraction, so
\(\|I-\mathbb E[\cdot\mid{\cal F}]\|_{p\to p}\leq2\).
\end{proof}

\subsection{Cubic landing theorem}

\begin{theorem}[Conditional interacting stress landing]
\label{thm:v18v48-cubic-landing}
Assume:
\begin{enumerate}[leftmargin=2.2em]
\item the heat covariance stock of V47 gives stress tightness;
\item the compiled gradient stock (58.10), including the V45 soft collar,
      holds for a separated three-test bank;
\item the follower projections obey (58.19);
\item the direct and complete Gaussian comparator finite-dimensional laws
      converge.
\end{enumerate}
Then the follower-shorted cubic direct-minus-Gaussian excess converges.  A
nonzero limit is a continuum non-quasi-free stress witness; a zero limit is
an honest Gaussian/Wick outcome for this test bank.
\end{theorem}

\begin{proof}
Tightness supplies convergent subsequences.  Theorem
\ref{thm:v18v48-direct-ui} and Theorem
\ref{thm:v18v48-projected-ui} give the direct
\(L^{3+\varepsilon}\) stock.  Corollaries
\ref{cor:v18v48-higgs-ui}--\ref{cor:v18v48-gaussian-short} give it for the
Gaussian comparator.  The V46 cubic-witness stability theorem passes every
moment entering the excess to the limit.
\end{proof}

\begin{assessment}[V48 acquisition and boundary]
\label{ass:v18v48}
V48 proves uniform integrability for the complete finite-chaos Gaussian
stress comparator and gives a logarithmic-Sobolev gradient criterion for the
direct packet.  It calculates the literal scalar, gauge, and fermion
gradient stocks and isolates the \(L^p\) follower-projection requirement.
The scale-critical log-Sobolev/gradient balance and the fermion soft-collar
gradient moment remain unpopulated.
\end{assessment}

\begin{record}[V48 append record]
\label{rec:v18v48}
V48 is appended to the validated V47 whose installed SHA--256 was
\texttt{B3DA5E921702700C400E1AEC5277B0F55964629F6846ADFFAF07984A13979394}.
After installation only V48 is the active numeric SM note; the frozen
\texttt{V100\_f17} archive is unchanged.
\end{record}


\section{V49: effective-action stationarity and quantitative Einstein-residual closure}
\label{sec:v18v49}

\subsection{The gauge-fixed metric chart}

The Einstein tensor cannot be obtained by varying over raw metric
coordinates because diffeomorphism directions are redundant and the
Euclidean conformal direction is not globally convex.  Fix a local physical
metric slice \({\cal X}\) through a reference metric, with Hilbert norm
\(\|\cdot\|_{\cal X}\), and let \({\cal X}^*\) be its dual.  Every statement
below is restricted to a convex ball
\[
 {\cal B}_R=\{x\in{\cal X}:\|x\|_{\cal X}\leq R\}             \tag{59.1}
\]
on which the metric remains nondegenerate and the slice is valid.

Let \(g(x)\) be the represented metric.  Define the complete regulated
effective action
\[
\boxed{
 \Gamma_n(x)
 =
 S_{{\rm grav},n}(g(x))
 +
 W_{{\rm SM},n}(g(x))
 +
 S_{{\rm ct},n}(g(x))
 +
 S_{{\rm slice},n}(x),}                                     \tag{59.2}
\]
where \(W_{{\rm SM},n}=-\log Z_{{\rm SM},n}\) uses the literal common
Standard-Model action.  The determinant phase, if nontrivial, remains a
separate line/current; \(\Gamma_n\) denotes the real effective action.

\begin{definition}[Gauge-fixed Einstein residual]
\label{def:v18v49-residual}
The residual \({\cal E}_n(x)\in{\cal X}^*\) is defined by
\[
\boxed{
 D\Gamma_n(x)[q]
 =
 \frac{1}{2\kappa_n}
 \langle{\cal E}_n(x),q\rangle
 \qquad(q\in{\cal X}).}                                     \tag{59.3}
\]
On physical slice directions, after counterterm and unit transport,
\[
 {\cal E}_n
 =
 G(g_n)+\Lambda_ng_n-\kappa_n\mathbb E_n{\mathsf T}_n
 +{\cal E}_{{\rm reg},n}.                                   \tag{59.4}
\]
\end{definition}

\begin{proposition}[The matter term is the literal stress source]
\label{prop:v18v49-matter-gradient}
For every slice test \(q\),
\[
 DW_{{\rm SM},n}(g)[q]
 =
 \mathbb E_nD_gS_{{\rm SM},n}(g,\phi)[q]
 =
 \frac12
 \langle\mathbb E_n{\mathsf T}_n,q\rangle.                  \tag{59.5}
\]
Thus the stress appearing in (59.4) is the compiled action derivative of
V40--V44 and is not an independently fitted tensor.
\end{proposition}

\begin{proof}
Differentiate the finite partition function under the integral:
\[
 D(-\log Z_n)[q]
 =
 \frac{\int D_gS_n[q]e^{-S_n}\,d\nu_n}
 {\int e^{-S_n}\,d\nu_n}.
\]
Use the V41 stress normalization.
\end{proof}

\subsection{Strong monotonicity on the physical slice}

\begin{definition}[Uniform slice margin]
\label{def:v18v49-slice-margin}
The regulated gradients have margin \(\lambda>0\) on \({\cal B}_R\) if
\[
\boxed{
 \langle
 D\Gamma_n(x)-D\Gamma_n(y),x-y
 \rangle
 \geq
 \lambda\|x-y\|_{\cal X}^2}                                 \tag{59.6}
\]
for every \(x,y\in{\cal B}_R\), uniformly in \(n\).
\end{definition}

\begin{lemma}[Residual error controls metric error]
\label{lem:v18v49-error-bound}
If \(x_n^*\in{\cal B}_R\) satisfies
\[
 D\Gamma_n(x_n^*)=0,                                        \tag{59.7}
\]
then every \(x\in{\cal B}_R\) obeys
\[
\boxed{
 \|x-x_n^*\|_{\cal X}
 \leq
 \lambda^{-1}\|D\Gamma_n(x)\|_{{\cal X}^*}.}                 \tag{59.8}
\]
\end{lemma}

\begin{proof}
Apply (59.6) to \(x,x_n^*\) and use (59.7):
\[
 \lambda\|x-x_n^*\|_{\cal X}^2
 \leq
 \langle D\Gamma_n(x),x-x_n^*\rangle
 \leq
 \|D\Gamma_n(x)\|_{{\cal X}^*}\|x-x_n^*\|_{\cal X}.
\]
Cancel the distance unless it is zero.
\end{proof}

\begin{theorem}[Existence and uniqueness on a coercive slice]
\label{thm:v18v49-existence}
Suppose \(\Gamma_n\) is lower semicontinuous, satisfies (59.6), and has a
sublevel set compactly contained in the interior of \({\cal B}_R\).  Then
\(\Gamma_n\) has a unique minimizer \(x_n^*\) in that sublevel set, and it
satisfies (59.7).
\end{theorem}

\begin{proof}
Compactness and lower semicontinuity give a minimizer.  Since it lies in the
interior, its first variation vanishes.  If \(x,y\) were two critical
points, (59.6) would give
\(\lambda\|x-y\|_{\cal X}^2\leq0\), hence \(x=y\).
\end{proof}

\begin{firewall}[Scope of the V34 margin]
\label{fire:v18v49-v34-scope}
V34 supplies only a conditional finite-regulator cyclic Hessian margin.
It can populate (59.6) only after its coordinates are identified with the
physical metric slice in (59.2), every omitted sector is added, and the
margin is uniform in the physical \({\cal X}\)-norm.  No global convexity of
the Euclidean Einstein--Hilbert functional is asserted.
\end{firewall}

\subsection{Cofinal convergence of critical metrics}

Assume there is a differentiable limiting functional
\(\Gamma:{\cal B}_R\to{\mathbb R}\) and define the gradient defect
\[
 \delta_n
 =
 \sup_{x\in{\cal B}_R}
 \|D\Gamma_n(x)-D\Gamma(x)\|_{{\cal X}^*}.                  \tag{59.9}
\]

\begin{theorem}[Quantitative critical-point convergence]
\label{thm:v18v49-critical-convergence}
Suppose every \(\Gamma_n\) satisfies (59.6), \(\delta_n\to0\), and
\(x_n^*\in{\cal B}_R\) solves (59.7).  If \(x^*\in{\cal B}_R\) satisfies
\[
 D\Gamma(x^*)=0,                                             \tag{59.10}
\]
then it is unique and
\[
\boxed{
 \|x_n^*-x^*\|_{\cal X}
 \leq\lambda^{-1}\delta_n.}                                 \tag{59.11}
\]
Conversely, any convergent subsequence of \(x_n^*\) has limit \(x^*\)
satisfying (59.10).
\end{theorem}

\begin{proof}
Apply Lemma \ref{lem:v18v49-error-bound} with \(x=x^*\):
\[
 \|x_n^*-x^*\|_{\cal X}
 \leq
 \lambda^{-1}\|D\Gamma_n(x^*)-D\Gamma(x^*)\|_{{\cal X}^*}
 \leq\lambda^{-1}\delta_n.
\]
Strong monotonicity passes to the uniform gradient limit, so the limiting
zero is unique.  For the converse,
\[
 \|D\Gamma(x_n^*)\|_{{\cal X}^*}
 \leq\delta_n,
\]
and continuity of \(D\Gamma\) identifies every convergent limit as a zero.
\end{proof}

\begin{corollary}[Approximate stationarity]
\label{cor:v18v49-approximate}
If instead
\[
 \|D\Gamma_n(\widehat x_n)\|_{{\cal X}^*}\leq\varepsilon_n, \tag{59.12}
\]
then
\[
\boxed{
 \|\widehat x_n-x^*\|_{\cal X}
 \leq
 \lambda^{-1}(\varepsilon_n+\delta_n).}                     \tag{59.13}
\]
\end{corollary}

\begin{proof}
Use (59.6) between \(\widehat x_n\) and \(x^*\), and bound
\[
 \|D\Gamma_n(x^*)\|_{{\cal X}^*}\leq\delta_n.
\]
The same duality calculation as in Lemma
\ref{lem:v18v49-error-bound} gives (59.13).
\end{proof}

\subsection{Einstein equation on the limiting slice}

\begin{theorem}[Effective-action Einstein closure]
\label{thm:v18v49-einstein-closure}
Assume Theorem \ref{thm:v18v49-critical-convergence}, and suppose:
\begin{enumerate}[leftmargin=2.2em]
\item the gravitational first variations converge to
      \(G(g)+\Lambda g\) in \({\cal X}^*\);
\item the compiled mean stresses converge to \(T\) in \({\cal X}^*\);
\item \(\kappa_n\to\kappa>0\), and the regulator first-variation residual
      tends to zero.
\end{enumerate}
Then the limiting critical metric satisfies
\[
\boxed{
 G(g(x^*))+\Lambda g(x^*)=\kappa T}
 \quad\hbox{on every physical slice test}.                  \tag{59.14}
\]
If the V41 Ward and Bianchi defects vanish, this equality is conserved and
extends from the slice to the full gauge-equivalence class.
\end{theorem}

\begin{proof}
At \(x_n^*\), equation (59.3) gives
\({\cal E}_n(x_n^*)=0\) on slice tests.  Pass to the limit using the three
hypotheses and (59.11), obtaining (59.14).  Vanishing Ward and Bianchi
defects make both sides transverse.  Pure diffeomorphism directions are
zero in the quotient, so equality on the physical slice represents equality
of the gauge classes.
\end{proof}

\begin{assessment}[What gradient convergence contains]
\label{ass:v18v49-gradient-content}
The defect \(\delta_n\) is not one scalar estimate.  By
(59.2)--(59.5), it contains gravitational counterterm convergence, the full
compiled stress convergence of V46--V48, determinant-line transport, and
the coordinate/slice transport.  Strong convexity cannot replace any of
these rows.
\end{assessment}

\subsection{Compatibility with an interacting stress limit}

\begin{proposition}[Metric convexity does not force Gaussian matter]
\label{prop:v18v49-convexity-interaction}
The slice margin (59.6) constrains the second metric derivative of the
\emph{total} effective action.  It does not imply that the follower-shorted
third stress cumulant or the V40 direct-minus-Gaussian excess vanishes.
Therefore (59.6) is logically compatible with a nonzero continuum cubic
stress witness.
\end{proposition}

\begin{proof}
Strong monotonicity is a lower bound on
\[
 D^2\Gamma_n(x)[q,q].
\]
The V40 witness is a third connected response, after contact and follower
shorts, compared with its Gaussian value.  No bound or identity for that
trilinear form follows from a positive lower bound on the quadratic form.
\end{proof}

\begin{firewall}[Uniform matter mixing remains restrictive]
\label{fire:v18v49-mixing}
If the proof of (59.6) additionally imposes uniform all-order anchored
cumulant summability on the same residual stress coordinate, V36 still
forces its non-Gaussian continuum vertices to vanish.  The desired branch
needs metric-slice stability together with scale-critical matter
correlations, not an all-order uniform mixing estimate for every sector.
\end{firewall}

\subsection{A posteriori residual certificate}

The effective-action formulation makes a numerical or interval certificate
mathematically meaningful.  If a candidate metric \(\widehat x_n\) is
returned with:
\[
 \|D\Gamma_n(\widehat x_n)\|_{{\cal X}^*}\leq\varepsilon_n,
 \qquad
 D^2\Gamma_n\succeq\lambda I
 \quad\hbox{on }{\cal B}_R,                                  \tag{59.15}
\]
then (59.13) gives a rigorous distance to the unique physical-slice
solution.  The residual must be computed from the complete action derivative
(59.5); a fitted stress tensor would not certify (59.15).

\subsection{Next cutoff theorem}

V49 converts a uniform physical slice margin and gradient convergence into
the Einstein equation.  The next step is to replace the strong uniform
gradient convergence (59.9) by a variational convergence condition stable
under rougher approximations.  V50 should prove a Mosco/graph convergence
theorem for the strongly convex effective actions and identify precisely
which stress-source convergence is sufficient.

\begin{assessment}[V49 acquisition and boundary]
\label{ass:v18v49}
V49 defines the gauge-fixed effective action from the literal SM partition
function, identifies its gradient with the Einstein residual, and proves
existence, uniqueness, quantitative approximate-stationarity, and cofinal
Einstein closure under a uniform slice margin and gradient convergence.  The
physical slice margin and full gradient defect remain conditional; global or
rough quantum-metric closure is not claimed.
\end{assessment}

\begin{record}[V49 append record]
\label{rec:v18v49}
V49 is appended to the validated V48 whose installed SHA--256 was
\texttt{196A380EC896B6CCD57DA8B9DD4C86A169FD38D766DB47ADE32A98FAD07A428A}.
After installation only V49 is the active numeric SM note; the frozen
\texttt{V100\_f17} archive is unchanged.
\end{record}


\section{V50: Mosco convergence of the effective action and graph-stable Einstein sources}
\label{sec:v18v50}

\subsection{Variational replacement for uniform gradient convergence}

Let \({\cal X}\) be the physical metric-slice Hilbert space of V49.  A
sequence of proper lower semicontinuous functionals
\(\Gamma_n:{\cal X}\to(-\infty,\infty]\) converges to \(\Gamma\) in the
Mosco sense if:
\begin{enumerate}[leftmargin=2.2em]
\item whenever \(x_n\rightharpoonup x\) weakly,
\[
 \Gamma(x)\leq\liminf_n\Gamma_n(x_n);                        \tag{60.1}
\]
\item for every \(x\in{\cal X}\), there are \(y_n\to x\) strongly with
\[
 \Gamma_n(y_n)\longrightarrow\Gamma(x).                     \tag{60.2}
\]
\end{enumerate}
This permits nonsmooth limits and does not require the uniform gradient
defect (59.9).

Assume each \(\Gamma_n\) is \(\lambda\)-strongly convex:
\[
 \Gamma_n((1-t)x+ty)
 \leq
 (1-t)\Gamma_n(x)+t\Gamma_n(y)
 -\frac{\lambda}{2}t(1-t)\|x-y\|_{\cal X}^2                 \tag{60.3}
\]
for \(0\leq t\leq1\), with one \(\lambda>0\).

\begin{lemma}[Strong convexity passes to the Mosco limit]
\label{lem:v18v50-limit-convexity}
Under (60.1)--(60.3), \(\Gamma\) is \(\lambda\)-strongly convex.
\end{lemma}

\begin{proof}
Choose strong recovery sequences \(x_n\to x\) and \(y_n\to y\).  Apply
(60.3) to them.  The convex combinations converge strongly to
\((1-t)x+ty\); use the Mosco liminf on the left and the recovery limits on
the right.
\end{proof}

\subsection{Strong convergence of effective metrics}

\begin{theorem}[Mosco convergence of minimizers]
\label{thm:v18v50-minimizers}
Suppose \(\Gamma_n\) Mosco-converges to \(\Gamma\), (60.3) holds, and each
\(\Gamma_n\) has a minimizer \(x_n\).  Suppose also that \((x_n)\) is
bounded and that \(\Gamma\) has a finite minimizer \(x_*\).  Then \(x_*\) is
unique and
\[
\boxed{
 x_n\longrightarrow x_*
 \quad\hbox{strongly in }{\cal X},\qquad
 \Gamma_n(x_n)\longrightarrow\Gamma(x_*).}                  \tag{60.4}
\]
\end{theorem}

\begin{proof}
By boundedness, every subsequence has a weakly convergent subsubsequence
\(x_{n_j}\rightharpoonup\bar x\).  Let \(y_n\to x_*\) be a recovery
sequence.  Minimality and (60.1) give
\[
 \Gamma(\bar x)
 \leq\liminf_j\Gamma_{n_j}(x_{n_j})
 \leq\limsup_j\Gamma_{n_j}(x_{n_j})
 \leq\lim_j\Gamma_{n_j}(y_{n_j})
 =\Gamma(x_*).                                               \tag{60.5}
\]
Thus \(\bar x=x_*\) by Lemma
\ref{lem:v18v50-limit-convexity}, and the middle energies converge.

At a minimizer of a \(\lambda\)-strongly convex functional,
\[
 \Gamma_n(y_n)
 \geq
 \Gamma_n(x_n)+\frac{\lambda}{2}\|y_n-x_n\|_{\cal X}^2.      \tag{60.6}
\]
The energy difference tends to zero by (60.5), so
\(\|y_n-x_n\|_{\cal X}\to0\).  Since \(y_n\to x_*\), strong convergence
follows.  Every subsequence has the same conclusion.
\end{proof}

\begin{corollary}[Approximate minimizers]
\label{cor:v18v50-approximate-minimizers}
If
\[
 \Gamma_n(\widehat x_n)
 \leq
 \inf\Gamma_n+\varepsilon_n,\qquad\varepsilon_n\to0,         \tag{60.7}
\]
and \((\widehat x_n)\) is bounded, then
\[
 \widehat x_n\longrightarrow x_*
 \quad\hbox{strongly}.                                      \tag{60.8}
\]
\end{corollary}

\begin{proof}
Let \(x_n\) be the exact minimizer.  Strong convexity gives
\[
 \frac{\lambda}{2}\|\widehat x_n-x_n\|_{\cal X}^2
 \leq\varepsilon_n.                                         \tag{60.9}
\]
Use Theorem \ref{thm:v18v50-minimizers}.
\end{proof}

\subsection{Resolvent and graph convergence}

For \(\tau>0\) and \(y\in{\cal X}\), define the metric resolvent
\[
 J_{n,\tau}y
 =
 \operatorname*{arg\,min}_{x\in{\cal X}}
 \left\{
 \Gamma_n(x)+\frac{1}{2\tau}\|x-y\|_{\cal X}^2
 \right\}.                                                   \tag{60.10}
\]

\begin{theorem}[Strong resolvent landing]
\label{thm:v18v50-resolvent}
Under the hypotheses of Theorem \ref{thm:v18v50-minimizers}, for every
\(\tau>0\) and \(y_n\to y\) strongly,
\[
\boxed{
 J_{n,\tau}y_n\longrightarrow J_\tau y
 \quad\hbox{strongly in }{\cal X}.}                          \tag{60.11}
\]
\end{theorem}

\begin{proof}
The perturbed functionals
\[
 x\longmapsto
 \Gamma_n(x)+\frac{1}{2\tau}\|x-y_n\|^2                     \tag{60.12}
\]
Mosco-converge to the analogous functional with \(\Gamma,y\): the norm term
is weakly lower semicontinuous for (60.1), and converges along every strong
recovery sequence for (60.2).  They are
\((\lambda+\tau^{-1})\)-strongly convex.  Their minimizers are bounded by
comparison with one recovery point.  Apply Theorem
\ref{thm:v18v50-minimizers}.
\end{proof}

The subgradient relation
\[
 \frac{y-J_{n,\tau}y}{\tau}
 \in\partial\Gamma_n(J_{n,\tau}y)                            \tag{60.13}
\]
shows why (60.11) is a graph-stable replacement for coefficientwise
gradient convergence.

\begin{theorem}[Subgradient graph closure]
\label{thm:v18v50-graph-closure}
Suppose \(\Gamma_n\) Mosco-converges to \(\Gamma\), and
\[
 x_n\to x,\qquad
 \xi_n\rightharpoonup\xi,\qquad
 \xi_n\in\partial\Gamma_n(x_n),                              \tag{60.14}
\]
with
\[
 \Gamma_n(x_n)\longrightarrow\Gamma(x).                     \tag{60.15}
\]
Then
\[
\boxed{\xi\in\partial\Gamma(x).}                             \tag{60.16}
\]
\end{theorem}

\begin{proof}
For arbitrary \(y\in{\cal X}\), choose a strong recovery sequence
\(y_n\to y\).  The subgradient inequality gives
\[
 \Gamma_n(y_n)
 \geq
 \Gamma_n(x_n)+\langle\xi_n,y_n-x_n\rangle.                 \tag{60.17}
\]
Pass to the limit using (60.2), (60.14), and (60.15):
\[
 \Gamma(y)\geq\Gamma(x)+\langle\xi,y-x\rangle.
\]
This is (60.16).
\end{proof}

\subsection{Identification of the matter stress subgradient}

Split the effective action as
\[
 \Gamma_n=A_n+W_n,                                          \tag{60.18}
\]
where \(A_n\) contains the gravitational, slice, and admitted deterministic
counterterm terms, and \(W_n=W_{{\rm SM},n}\) is the real Standard-Model
effective action.

\begin{theorem}[Graph-stable stress source]
\label{thm:v18v50-stress-graph}
Suppose:
\begin{enumerate}[leftmargin=2.2em]
\item \(x_n\to x\) strongly and
      \(0\in DA_n(x_n)+\partial W_n(x_n)\);
\item \(W_n\) Mosco-converges to \(W\);
\item \(W_n(x_n)\to W(x)\);
\item \(DA_n(x_n)\rightharpoonup a\) in \({\cal X}^*\).
\end{enumerate}
Then
\[
\boxed{-a\in\partial W(x).}                                 \tag{60.19}
\]
If \(W\) is differentiable at \(x\), then
\[
 DW(x)=-a,                                                   \tag{60.20}
\]
and this derivative is the limiting compiled mean stress source.
\end{theorem}

\begin{proof}
Choose \(\xi_n\in\partial W_n(x_n)\) with
\(\xi_n=-DA_n(x_n)\).  Apply Theorem
\ref{thm:v18v50-graph-closure} to \(W_n\).  Differentiability makes the
subdifferential a singleton.
\end{proof}

\begin{corollary}[Variational Einstein equation]
\label{cor:v18v50-variational-einstein}
If \(a\) is the limiting gravitational first variation and
\[
 DW(x)[q]=\frac12\langle T,q\rangle,                         \tag{60.21}
\]
then (60.20) is the weak gauge-slice Einstein equation with source \(T\),
using the normalization of V49.  Vanishing Ward and Bianchi defects extend
it to the metric gauge class.
\end{corollary}

\begin{proof}
Insert the gravitational and matter first-variation formulas into
\(a+DW(x)=0\).  The V41 conservation theorem supplies the gauge extension.
\end{proof}

\begin{firewall}[Mosco convergence of the sum is insufficient for source
identification]
\label{fire:v18v50-sum-insufficient}
Mosco convergence of \(\Gamma_n=A_n+W_n\) determines the limiting total
variational problem and its minimizer.  It does not by itself identify the
separate limit of \(DW_n\).  The energy and subgradient graph hypotheses in
Theorem \ref{thm:v18v50-stress-graph} are required to prevent an unobserved
transfer between gravitational counterterms and the matter stress.
\end{firewall}

\subsection{Ward-constrained Mosco limit}

Let \({\cal X}_n^{\rm phys}\subset{\cal X}\) be the transported finite
physical slices with orthogonal projections \(\Pi_n\), and assume
\[
 \|\Pi_n-\Pi\|_{\rm op}\longrightarrow0.                     \tag{60.22}
\]
Extend \(\Gamma_n\) by \(+\infty\) outside \({\cal X}_n^{\rm phys}\).

\begin{proposition}[Stable physical-slice constraint]
\label{prop:v18v50-slice-constraint}
If the unrestricted local representatives Mosco-converge and (60.22) holds,
then the constrained functionals Mosco-converge on
\({\cal X}^{\rm phys}=\operatorname{Ran}\Pi\), provided their recovery
sequences may be projected without changing the limiting energy.
\end{proposition}

\begin{proof}
If \(x_n\in\operatorname{Ran}\Pi_n\) and \(x_n\rightharpoonup x\), then
\[
 (I-\Pi)x
 =
 \lim_n(I-\Pi_n)x_n=0
\]
weakly, so \(x\in\operatorname{Ran}\Pi\); the unrestricted liminf applies.
For \(x\in\operatorname{Ran}\Pi\), take an unrestricted recovery sequence
\(y_n\to x\) and replace it by \(\Pi_ny_n\).  Equation (60.22) gives
\(\Pi_ny_n\to x\), and the stated energy stability gives recovery.
\end{proof}

\subsection{Compatibility with the stress interaction witness}

\begin{proposition}[First-variation landing does not determine higher
connected response]
\label{prop:v18v50-higher-response}
The graph convergence in Theorem \ref{thm:v18v50-stress-graph} identifies
the limiting mean stress, but does not determine the V40 cubic
direct-minus-Gaussian excess.  That excess still requires the V47 covariance
and V48 uniform-integrability stocks.
\end{proposition}

\begin{proof}
Subgradient convergence concerns one first derivative of \(W_n\) at
\(x_n\).  Two sequences of convex functionals can have the same value and
first derivative at a point but different third derivatives there.  Hence
the cubic response is independent data.
\end{proof}

\subsection{Checkpoint}

V50 reaches the next five-version boundary.  Before V51, the active
Yang--Mills and Navier--Stokes notes must be inspected.  After that audit,
the next target is a single closure theorem combining:
\[
\boxed{
 \text{Mosco metric landing, compiled stress tightness, Ward descent,
 and a surviving separated interaction witness}.}           \tag{60.23}
\]
Such a theorem will state exactly which still-unpopulated physical inputs
separate the present programme from a full SM--Einstein continuum.

\begin{assessment}[V50 acquisition and boundary]
\label{ass:v18v50}
V50 replaces uniform effective-action gradient convergence by Mosco
convergence, proves strong convergence of minimizers and resolvents, and
gives a graph theorem that preserves the identity of the matter stress
source.  It also treats moving physical slices and keeps higher connected
responses separate.  The required physical Mosco, energy, subgradient, and
slice-projection convergence remain unproved.
\end{assessment}

\begin{record}[V50 append record]
\label{rec:v18v50}
V50 is appended to the validated V49 whose installed SHA--256 was
\texttt{3388F2A6298A3A849F6D24863591A781DA52281D696D06469987038925B95C1E}.
After installation only V50 is the active numeric SM note; the frozen
\texttt{V100\_f17} archive is unchanged.
\end{record}


\section{V51: conditional regular-branch SM--Einstein continuum closure and the exact remaining capsule}
\label{sec:v18v51}

\subsection{Mandatory V50 companion audit}

The active companion notes are
\[
\begin{array}{c|c|c|c}
\hbox{programme}&\hbox{version}&\hbox{bytes}&\hbox{SHA--256}\\ \hline
\hbox{Yang--Mills}&V8&145033&
\texttt{DAC23C20F3B2984995D0AB9B9B0735F0921D60E441157861F825253D93D294D2}\\
\hbox{Navier--Stokes}&V26&278283&
\texttt{82C887F8C2C69E4F28FAD7C3DC8DE5B9099CB5A4BF431EBA1FFF3E91935978AD}.
\end{array}                                                     \tag{61.1}
\]
Yang--Mills V8 replaces a failed pointwise enclosure by a live-phase
frozen-inverse scalar.  It integrates the signed trial polynomial exactly
and norms only the inverse defect; it proves that finite rational
certificates exist at arbitrary accuracy, while explicitly leaving
feasibility and the sign computation open.  The SM transfer is to compile
each finite stress/Ward/interaction panel with its phases and reference
subtractions intact before adding an error radius.  No Yang--Mills sign,
cofinal state, or mass-gap conclusion is imported.  Navier--Stokes V26 is
unchanged and does not alter the direction.

\subsection{Operational definition of the target}

For this programme, a \emph{regular-branch SM--Einstein continuum} means a
cutoff-independent limiting Schwinger functional with:
\begin{enumerate}[leftmargin=2.2em]
\item the Standard-Model gauge, Higgs, and three-generation chiral fermion
      field content on a nondegenerate limiting metric;
\item gauge invariance or an equivalent physical BRST quotient, a globally
      consistent determinant/Pfaffian line, and the anomaly-free
      representation of V44;
\item reflection positivity and the remaining reconstruction hypotheses on
      a declared physical time reflection;
\item a conserved compiled stress tensor and the distributional Einstein
      equation;
\item one follower-shorted, contact-quotiented, separated-support
      direct-minus-Gaussian connected response with a nonzero limit; and
\item independence of the chosen cofinal cutoff route.
\end{enumerate}
The fifth item prevents a Gaussian or merely contact-non-Gaussian theory
from being called interacting.

\subsection{The closure capsule}

Let \({\cal A}_0\) be a countable nuclear test algebra containing elementary
bosonic fields, gauge-invariant Wilson tests, fermionic source tests, the
enhanced scalar stress packet of V42, and the full stress tests of
V43--V44.  Let \({\cal S}_n\) be the normalized finite-cutoff Schwinger
functional on \({\cal A}_0\).

\begin{definition}[Regular closure hypotheses]
\label{def:v18v51-hypotheses}
The following rows form the regular closure capsule.
\begin{enumerate}[leftmargin=2.5em]
\item[\(\mathbf{C1}\)] \textbf{Common action and line.}
The scalar, gauge, fermion, ghost, measure, and metric terms arise from one
regulated action; the chiral determinant/Pfaffian charts glue to a fixed
global line with controlled Berry current and exactly saturated zero modes.

\item[\(\mathbf{C2}\)] \textbf{Positivity and gauge quotient.}
The finite Schwinger functionals are normalized and reflection positive on
the physical gauge/BRST quotient, and their polynomial growth bounds are
uniform on every finite test panel.

\item[\(\mathbf{C3}\)] \textbf{Joint compactness.}
The bosonic fields and enhanced stress packet are tight in declared
distribution spaces; in particular the V47 heat stock and the metric
compactness hypothesis of V46 hold.  Fermionic source coefficients have
uniform bounds sufficient for diagonal convergence of all Grassmann
Schwinger forms.

\item[\(\mathbf{C4}\)] \textbf{Moment uniform integrability.}
The V48 \(L^{3+\varepsilon}\) stock holds for the selected interaction panel
and the higher fixed moments needed by \({\cal A}_0\).

\item[\(\mathbf{C5}\)] \textbf{Ward and anomaly descent.}
The completely assembled defect (54.27), including hard/soft Dirac collars,
converges to zero in the dual topology of V41, and the inserted identities
are dense enough to give the closed Ward graph of V46.

\item[\(\mathbf{C6}\)] \textbf{Metric variational landing.}
The gauge-fixed effective actions Mosco-converge with one physical-slice
strong-convexity margin, their minimizers are bounded, and the matter
subgradient satisfies the energy/graph conditions of V50.

\item[\(\mathbf{C7}\)] \textbf{Einstein residual.}
The regulated Bianchi defect vanishes and the full effective-action
critical residual tends to zero in the same distribution topology.

\item[\(\mathbf{C8}\)] \textbf{Interaction survival.}
For three positively separated tests, the raw response is compiled by the
V40 partition formula, the full Gaussian scalar comparator includes the
V42 fourth chaos, both sides undergo the same dynamic follower short, and
\[
 \lim_n{\cal E}_{3,n}^{\perp}=e_*\ne0.                       \tag{61.2}
\]

\item[\(\mathbf{C9}\)] \textbf{Route closure.}
Adjacent and independently retained direct transports of fields, source
lines, covariances, Ward defects, follower projections, stress, and metric
forms have a common summable defect tail.

\item[\(\mathbf{C10}\)] \textbf{Reconstruction/cluster row.}
The limiting Schwinger functional satisfies symmetry, regularity,
reflection, cluster, and domain hypotheses sufficient for the declared
Osterwalder--Schrader reconstruction, with the massless photon sector
separated from any claimed massive channel.
\end{enumerate}
\end{definition}

\subsection{Conditional closure theorem}

\begin{theorem}[Regular-branch SM--Einstein continuum]
\label{thm:v18v51-conditional-closure}
If one literal cofinal cutoff family satisfies
\(\mathbf{C1}\)--\(\mathbf{C10}\), then it has a
cutoff-independent regular-branch SM--Einstein continuum in the sense
defined above.  More precisely:
\begin{enumerate}[leftmargin=2.2em]
\item the limiting Schwinger functional exists on \({\cal A}_0\) and
      extends by the stated growth bounds;
\item reflection positivity and the physical gauge quotient pass to the
      limit, so the reconstruction produces a positive physical Hilbert
      space and the declared field operators;
\item the compiled continuum stress is conserved and satisfies
\[
 \boxed{G(g)+\Lambda g=\kappa T}                             \tag{61.3}
\]
in the declared distribution space;
\item the limit is not quasi-free modulo contacts and declared followers,
      because its separated cubic excess equals \(e_*\ne0\); and
\item every admissible cofinal route gives the same functional and metric
      class.
\end{enumerate}
\end{theorem}

\begin{proof}
By \(\mathbf{C3}\), Prokhorov compactness for the bosonic and enhanced
coordinates, together with diagonal extraction on the countable fermionic
test algebra, gives a subsequential limiting Schwinger functional.
Condition \(\mathbf{C4}\) passes all fixed moments used in the selected
panels to the limit.

Reflection positivity is a family of closed nonnegative quadratic
inequalities on finite positive-time test panels.  Conditions
\(\mathbf{C2}\)--\(\mathbf{C4}\) pass each inequality to the limit.
The gauge/BRST null space is quotiented before completion, and
\(\mathbf{C1}\) gives a globally consistent chiral line and saturated
sector decomposition.  The reconstruction hypotheses in \(\mathbf{C10}\)
then give the physical Hilbert-space realization.

Conditions \(\mathbf{C5}\) and \(\mathbf{C7}\), through the V41 and V46
closed-Ward theorems, give distributional stress conservation and matching
Bianchi conservation.  Condition \(\mathbf{C6}\), through V50 minimizer and
subgradient graph convergence, identifies the limiting matter stress as the
same action derivative appearing in the metric variational equation.
The residual clause in \(\mathbf{C7}\) gives (61.3).

By \(\mathbf{C4}\), the cubic-witness stability theorem applies.
Condition \(\mathbf{C8}\) therefore yields the nonzero separated continuum
excess \(e_*\).  Contacts vanish on these supports and declared followers
have already been projected out, so the limit cannot be quasi-free modulo
those ranges.

Finally, telescope the summable adjacent defects in \(\mathbf{C9}\).
Agreement with each retained direct route makes every cofinal subsequence
have the same finite-panel limits.  Density of \({\cal A}_0\) and the
uniform growth bounds identify the full limiting functional and metric
class, proving cutoff independence.
\end{proof}

\begin{firewall}[Logical status]
\label{fire:v18v51-logical-status}
Theorem \ref{thm:v18v51-conditional-closure} is an implication, not a proof
that a literal cutoff family satisfies \(\mathbf{C1}\)--\(\mathbf{C10}\).
Several rows remain unpopulated.  The full SM--Einstein continuum goal has
therefore not yet been reached.
\end{firewall}

\subsection{Unconditional obstruction inherited from V36}

\begin{theorem}[Scale-critical necessity for the closure capsule]
\label{thm:v18v51-critical-necessity}
Suppose \(\mathbf{C3}\), \(\mathbf{C4}\), and the nonzero cubic conclusion
(61.2) hold with bounded critical normalization.  Let \({\cal K}_{3,n}^{\rm dir}\) and
\({\cal K}_{3,n}^{\rm G}\) be the anchored connected stocks of the direct
and matched-Gaussian follower-shorted stress densities.  Then necessarily
\[
\boxed{
 \max\{
 {\cal K}_{3,n}^{\rm dir},{\cal K}_{3,n}^{\rm G}
 \}\gtrsim h_n^{-2}}                                        \tag{61.4}
\]
along a subsequence.  In particular, no branch on which both compiled
stocks are uniformly summable can satisfy \(\mathbf{C8}\).
\end{theorem}

\begin{proof}
Apply the V36 cumulant estimate separately to the direct and
matched-Gaussian local densities.  For unit-bounded tests,
\[
 |{\cal E}_{3,n}^{\perp}|
 \leq
 C h_n^2
 \bigl(
 {\cal K}_{3,n}^{\rm dir}+{\cal K}_{3,n}^{\rm G}
 \bigr).
\]
On a subsequence where the left side is bounded below, at least one of the
two stocks has the rate (61.4).  If both were uniformly summable, both
third cumulants and their difference would instead tend to zero.
\end{proof}

\begin{assessment}[Interpretation of the necessity theorem]
\label{ass:v18v51-necessity}
The programme cannot close by extending the V34 Dobrushin estimate
uniformly to every order and every residual stress coordinate.  It must
retain stability of the metric slice while allowing a specific physical
matter channel to develop the scale-critical connected stock (61.4).
\end{assessment}

\subsection{Finite-panel certification architecture}

For a frozen finite test panel, write the target scalar as
\[
 {\cal Q}_n
 =
 {\cal Q}_n^{\rm trial}
 +
 {\cal R}_n^{\rm inv}
 +
 {\cal R}_n^{\rm collar}
 +
 {\cal R}_n^{\rm route}.                                    \tag{61.5}
\]
Here \({\cal Q}_n^{\rm trial}\) retains all action words, determinant phases,
direct-minus-Gaussian subtractions, and follower projections before
evaluation.  The three residuals are respectively inverse, hard/soft collar,
and transport errors.

\begin{proposition}[Sound compiled panel certificate]
\label{prop:v18v51-panel-certificate}
If a rational or interval computation proves
\[
 {\cal Q}_n^{\rm trial}\in I_n,
 \qquad
 |{\cal R}_n^{\rm inv}|\leq r_n^{\rm inv},
 \quad
 |{\cal R}_n^{\rm collar}|\leq r_n^{\rm collar},
 \quad
 |{\cal R}_n^{\rm route}|\leq r_n^{\rm route},               \tag{61.6}
\]
then
\[
\boxed{
 {\cal Q}_n
 \in
 I_n+
 [-r_n,r_n],
 \qquad
 r_n=r_n^{\rm inv}+r_n^{\rm collar}+r_n^{\rm route}.}        \tag{61.7}
\]
A positive interaction floor is certified if the right side of (61.7)
stays a positive distance from zero uniformly on a cofinal tail.
\end{proposition}

\begin{proof}
Equation (61.5) is an exact compiled identity.  Apply the triangle
inequality only to its three error terms and use interval addition.
\end{proof}

\begin{firewall}[Finite certificate versus continuum row]
\label{fire:v18v51-finite-versus-continuum}
Existence of arbitrarily accurate certificates at each fixed cutoff does not
give a uniform cofinal certificate, a feasible complexity bound, or the
positive distance from zero required in (61.7).  These are separate
quantitative rows.
\end{firewall}

\subsection{Exact unresolved capsule after V51}

The results through V51 reduce the first useful physical submission to the
following data from one literal regulator family:
\begin{enumerate}[leftmargin=2.2em]
\item a global anomaly-free chiral determinant/Pfaffian line with saturated
      zero modes and reflection-positive physical quotient;
\item one compiled heat covariance stock (57.3) and direct
      \(L^{3+\varepsilon}\) gradient/collar stock (58.10);
\item one dynamically follower-shorted, complete-chaos, separated cubic
      stress excess with a uniform nonzero interval as in (61.7);
\item the full Ward collar/complement estimate for the assembled defect
      (54.27);
\item Mosco convergence of the physical-slice effective action together
      with the matter subgradient graph row; and
\item a vanishing Einstein residual and summable direct/staged route defect.
\end{enumerate}
These are not formal placeholders: every item has a precise norm,
projection, or interval target in V39--V50.  The most upstream unresolved
item is still the literal global chiral regulator and determinant line; the
most direct analytic interaction item is the scale-critical stress stock
(61.4).

\subsection{Next attack}

The next note should go upstream to the chiral line rather than assume
\(\mathbf{C1}\).  On a finite family of overlap/Ginsparg--Wilson Dirac
matrices, it should derive the spectral-projector connection, its curvature,
and the precise condition under which the anomaly-free representation
trivializes the determinant line on the admissible gauge-field component.
If global trivialization cannot be proved, the first nonzero holonomy is the
correct obstruction witness.

\begin{assessment}[V51 acquisition and boundary]
\label{ass:v18v51}
V51 records the mandatory companion audit, states and proves a conditional
regular-branch SM--Einstein continuum theorem, proves the scale-critical
necessity of any surviving cubic stress interaction, and supplies a sound
compiled finite-panel certificate architecture.  It lists the exact
unpopulated capsule and selects the global chiral determinant line as the
next upstream target.  The full continuum has not yet been constructed.
\end{assessment}

\begin{record}[V51 append record]
\label{rec:v18v51}
V51 is appended to the validated V50 whose installed SHA--256 was
\texttt{9746E38099290F46E42868AB01530B3AB6D8867C164E6E82FFF807CC675EA4E4}.
After installation only V51 is the active numeric SM note; the frozen
\texttt{V100\_f17} archive is unchanged.
\end{record}


\section{V52: overlap chiral projectors, determinant-line curvature, and global measure obstructions}
\label{sec:v18v52}

\subsection{Finite overlap chirality}

Let \(D_n(U)\) be a finite gauge-covariant lattice Dirac matrix satisfying
\[
 D_n^*=\gamma_5D_n\gamma_5,
 \qquad
 \gamma_5D_n+D_n\gamma_5
 =
 a_nD_n\gamma_5D_n.                                        \tag{62.1}
\]
Define
\[
 \widehat\gamma_{5,n}
 =
 \gamma_5(I-a_nD_n),
 \qquad
 \widehat P_{-,n}
 =
 \frac12(I-\widehat\gamma_{5,n}).                            \tag{62.2}
\]

\begin{lemma}[Exact Ginsparg--Wilson projector]
\label{lem:v18v52-gw-projector}
Under (62.1),
\[
 \widehat\gamma_{5,n}^*
 =
 \widehat\gamma_{5,n},
 \qquad
 \widehat\gamma_{5,n}^2=I,                                  \tag{62.3}
\]
and consequently
\[
\boxed{
 \widehat P_{-,n}^*
 =
 \widehat P_{-,n},
 \qquad
 \widehat P_{-,n}^2=\widehat P_{-,n}.}                       \tag{62.4}
\]
\end{lemma}

\begin{proof}
The first identity in (62.1) gives
\[
 \widehat\gamma_{5,n}^*
 =
 (I-a_nD_n^*)\gamma_5
 =
 \gamma_5(I-a_nD_n).
\]
Multiplying the Ginsparg--Wilson relation by \(\gamma_5\) on the left gives
\[
 D_n+D_n^*
 =
 a_nD_n^*D_n.                                                \tag{62.5}
\]
Therefore
\[
\begin{split}
 \widehat\gamma_{5,n}^2
 &=
 \gamma_5(I-a_nD_n)\gamma_5(I-a_nD_n)\\
 &=
 (I-a_nD_n^*)(I-a_nD_n)\\
 &=
 I-a_n(D_n+D_n^*)+a_n^2D_n^*D_n
 =I.
\end{split}
\]
Equation (62.4) follows from (62.2)--(62.3).
\end{proof}

Let \({\cal U}_{n}^{\rm adm}\) be one connected component of admissible
gauge fields on which the rank of \(\widehat P_{-,n}(U)\) is constant.
The left chiral spaces form the finite-rank Hermitian vector bundle
\[
 E_{-,n}
 =
 \operatorname{Ran}\widehat P_{-,n}
 \longrightarrow{\cal U}_{n}^{\rm adm}.                     \tag{62.6}
\]
Its determinant line is
\[
 {\cal L}_{n}
 =
 \det E_{-,n}.                                               \tag{62.7}
\]
For a full chiral determinant with a varying antifermion bundle, one tensors
(62.7) with the corresponding dual determinant line.  The formulas below
then apply to the signed sum of the two curvatures.

\subsection{Local frame connection}

On a chart \({\cal O}\subset{\cal U}_{n}^{\rm adm}\), choose an orthonormal
frame
\[
 v=(v_1,\ldots,v_N),
 \qquad
 v^*v=I,
 \qquad
 vv^*=\widehat P_{-,n}.                                     \tag{62.8}
\]
Define the real Berry connection one-form
\[
 {\cal A}_v
 =
 -i\,\operatorname{Tr}(v^*dv).                              \tag{62.9}
\]

\begin{lemma}[Frame-change law]
\label{lem:v18v52-frame-change}
If \(v'=vg\) for a smooth \(g:{\cal O}\to U(N)\), then
\[
\boxed{
 {\cal A}_{v'}
 =
 {\cal A}_v-i\,\operatorname{Tr}(g^{-1}dg).}                 \tag{62.10}
\]
Equivalently, the connection changes by the logarithmic derivative of the
determinant transition function.
\end{lemma}

\begin{proof}
Since \(v'^*dv'=g^{-1}v^*(dv)g+g^{-1}dg\), trace invariance gives
\[
 \operatorname{Tr}(v'^*dv')
 =
 \operatorname{Tr}(v^*dv)+\operatorname{Tr}(g^{-1}dg).
\]
Multiply by \(-i\).
\end{proof}

\begin{theorem}[Projector curvature formula]
\label{thm:v18v52-curvature}
The chartwise two-form \(d{\cal A}_v\) is frame independent and equals
\[
\boxed{
 {\cal F}_n
 =
 -i\,\operatorname{Tr}
 \left(
 \widehat P_{-,n}\,
 d\widehat P_{-,n}\wedge d\widehat P_{-,n}
 \right).}                                                  \tag{62.11}
\]
For tangent directions \(X,Y\),
\[
 {\cal F}_n(X,Y)
 =
 -i\,\operatorname{Tr}
 \left(
 \widehat P_{-,n}
 [D_X\widehat P_{-,n},D_Y\widehat P_{-,n}]
 \right).                                                    \tag{62.12}
\]
\end{theorem}

\begin{proof}
Differentiate \(v^*v=I\) and \(P=vv^*\).  The curvature of the induced
connection \(v^*dv\) on \(E_-\) is
\[
 d(v^*dv)+(v^*dv)\wedge(v^*dv)
 =
 dv^*(I-vv^*)\wedge dv.                                     \tag{62.13}
\]
Taking the trace kills
\(\operatorname{Tr}((v^*dv)\wedge(v^*dv))\), because it is the trace of a
matrix commutator in each pair of tangent directions.  A direct expansion
of \(P\,dP\wedge dP\) using \(P=vv^*\) gives the trace of the right side of
(62.13).  Multiplication by \(-i\) proves (62.11).  Evaluation on \(X,Y\)
gives (62.12).  Frame independence also follows from (62.10), since the
exterior derivative of the logarithmic transition term is zero.
\end{proof}

\subsection{Spectral formula and gap bound}

For an overlap construction, let \(H_n(U)=H_n(U)^*\) be the Hermitian
kernel whose negative spectral subspace defines the chiral projector:
\[
 \widehat P_{-,n}
 =
 1_{(-\infty,0)}(H_n).                                      \tag{62.14}
\]
Assume a contour \(\Gamma\) separates the negative and positive spectra and
has distance at least \(\delta_n>0\) from \(\operatorname{spec}H_n\).

\begin{theorem}[Riesz derivative of the chiral projector]
\label{thm:v18v52-riesz-derivative}
For every tangent direction \(X\),
\[
\boxed{
 D_X\widehat P_{-,n}
 =
 \frac{1}{2\pi i}
 \int_\Gamma
 (z-H_n)^{-1}(D_XH_n)(z-H_n)^{-1}\,dz,}                     \tag{62.15}
\]
and
\[
\boxed{
 \|D_X\widehat P_{-,n}\|_{\rm op}
 \leq
 \frac{L_\Gamma}{2\pi\delta_n^2}
 \|D_XH_n\|_{\rm op}.}                                      \tag{62.16}
\]
\end{theorem}

\begin{proof}
Use the Riesz formula
\[
 \widehat P_{-,n}
 =
 \frac{1}{2\pi i}
 \int_\Gamma(z-H_n)^{-1}\,dz.
\]
The derivative of the resolvent is
\[
 D_X(z-H_n)^{-1}
 =
 (z-H_n)^{-1}(D_XH_n)(z-H_n)^{-1}.
\]
This proves (62.15).  Both resolvents have norm at most
\(\delta_n^{-1}\) on \(\Gamma\); integration gives (62.16).
\end{proof}

\begin{corollary}[Curvature envelope]
\label{cor:v18v52-curvature-envelope}
If \(r_n=\operatorname{rank}\widehat P_{-,n}\), then
\[
\boxed{
 |{\cal F}_n(X,Y)|
 \leq
 2r_n
 \left(\frac{L_\Gamma}{2\pi\delta_n^2}\right)^2
 \|D_XH_n\|_{\rm op}\|D_YH_n\|_{\rm op}.}                   \tag{62.17}
\]
\end{corollary}

\begin{proof}
From (62.12),
\[
 |\operatorname{Tr}(P[A,B])|
 \leq
 r_n\|[A,B]\|_{\rm op}
 \leq2r_n\|A\|_{\rm op}\|B\|_{\rm op}.
\]
Apply (62.16).
\end{proof}

The bound deteriorates like \(\delta_n^{-4}\).  A shrinking overlap-kernel
gap must therefore enter the V45 hard/soft collar; it cannot be hidden by
the perturbative charge sums.

\subsection{Chern-period obstruction}

\begin{theorem}[Necessary topological triviality criterion]
\label{thm:v18v52-chern-obstruction}
For every smooth closed oriented two-cycle
\(\Sigma\subset{\cal U}_{n}^{\rm adm}\),
\[
\boxed{
 c_{1,n}[\Sigma]
 =
 \frac{1}{2\pi}\int_\Sigma{\cal F}_n
 \in{\mathbb Z}.}                                           \tag{62.18}
\]
If \(c_{1,n}[\Sigma]\ne0\) for one cycle, the determinant line
\({\cal L}_n\) has no global nowhere-zero section on that component.
\end{theorem}

\begin{proof}
On an oriented chart cover of \(\Sigma\), Stokes' theorem converts the
integral of the local curvatures into integrals of the logarithmic
transition derivatives (62.10) around chart boundaries.  Their sum is the
winding number of the \(U(1)\) determinant transition functions, hence an
integer.  A nowhere-zero global section would trivialize all transition
functions and force every such winding number to vanish.
\end{proof}

\begin{corollary}[Sufficient topological criterion]
\label{cor:v18v52-topological-sufficiency}
If the admissible component is paracompact and has the homotopy type of a
CW complex, then
\[
 c_1({\cal L}_n)=0\in H^2({\cal U}_{n}^{\rm adm};{\mathbb Z}) \tag{62.19}
\]
is equivalent to topological triviality of the complex determinant line.
\end{corollary}

\begin{proof}
Complex line bundles over such a base are classified by their first Chern
class.
\end{proof}

\begin{firewall}[Perturbative cancellation is not (62.19)]
\label{fire:v18v52-local-not-global}
The V44 hypercharge and mixed-anomaly sums cancel the leading local
continuum anomaly polynomial.  They do not by themselves prove that every
finite-regulator Chern period (62.18) vanishes on every admissible gauge
component.  The total projector curvature must first be assembled over the
complete chiral representation and then integrated on the actual cycles.
\end{firewall}

\subsection{Flat holonomy is a separate issue}

For a closed loop \(\ell\) covered by determinant charts, define the
connection holonomy
\[
 {\rm Hol}_n(\ell)
 =
 \exp\left(i\oint_\ell{\cal A}\right),                       \tag{62.20}
\]
including the determinant transition functions at chart changes.

\begin{proposition}[Contractible-loop formula]
\label{prop:v18v52-stokes-holonomy}
If \(\ell=\partial\Sigma\) and the determinant line is trivialized over
\(\Sigma\), then
\[
\boxed{
 {\rm Hol}_n(\ell)
 =
 \exp\left(i\int_\Sigma{\cal F}_n\right).}                   \tag{62.21}
\]
\end{proposition}

\begin{proof}
Apply Stokes' theorem chartwise.  Boundary transition terms cancel those
inserted in the definition of (62.20).
\end{proof}

\begin{countertheorem}[Zero Chern class need not give a parallel measure]
\label{cth:v18v52-flat-holonomy}
There are topologically trivial determinant lines with zero curvature but
nontrivial connection holonomy.
\end{countertheorem}

\begin{proof}
On the circle take the trivial line with connection
\[
 {\cal A}=\frac{\theta}{2\pi}\,d\varphi,
 \qquad \theta\notin2\pi{\mathbb Z}.                         \tag{62.22}
\]
Its curvature and first Chern class vanish, but
\[
 {\rm Hol}(S^1)=e^{i\theta}\ne1.                             \tag{62.23}
\]
Thus a global nonzero section exists, but no global parallel section exists
for this connection.
\end{proof}

\begin{assessment}[Three distinct fermion-line gates]
\label{ass:v18v52-three-gates}
The physical measure problem contains three logically distinct rows:
\begin{enumerate}[leftmargin=2.2em]
\item local anomaly cancellation, controlled in the continuum by the V44
      representation arithmetic;
\item topological triviality of the finite determinant line, tested by all
      Chern periods (62.18);
\item compatibility of the chosen measure connection with gauge loops,
      including flat holonomy on noncontractible cycles.
\end{enumerate}
Only the first row is presently populated.
\end{assessment}

\subsection{Gauge-equivariant measure}

Let \({\cal G}_n\) act on \({\cal U}_n^{\rm adm}\), and let
\(\ell_\omega(U)\) be the gauge-orbit loop or path generated by a gauge
transformation \(\omega\).  A physical chiral measure is a section of
\({\cal L}_n\) whose transformation phase matches the regulated fermion
Jacobian.

\begin{theorem}[Necessary gauge-loop integrability]
\label{thm:v18v52-gauge-loop}
If a single-valued gauge-equivariant chiral measure exists on an admissible
component, then for every closed gauge-orbit loop \(\ell\),
\[
\boxed{
 {\rm Hol}_n(\ell)
 =
 \exp\bigl(i\,{\cal J}_{F,n}(\ell)\bigr),}                   \tag{62.24}
\]
where \({\cal J}_{F,n}(\ell)\) is the integrated regulated fermion Jacobian
fixed by the representation.  In an anomaly-free convention with zero
target Jacobian, this reduces to
\[
 {\rm Hol}_n(\ell)=1.                                       \tag{62.25}
\]
\end{theorem}

\begin{proof}
Transport the proposed section around \(\ell\).  Parallel transport gives
the left side of (62.24), while gauge equivariance prescribes the right
side.  Single-valuedness requires equality.
\end{proof}

\begin{firewall}[Necessity versus construction]
\label{fire:v18v52-integrability-scope}
Equations (62.18) and (62.24) are necessary.  Even when they hold, the
measure section must still be constructed with the required locality,
smoothness, reflection, sector transport, and continuum bounds.
\end{firewall}

\subsection{Pfaffian orientation}

For a real or pseudoreal chiral block, the Pfaffian takes values in a real
orientation line.  Its obstruction is
\[
 w_1({\cal L}^{\rm Pf}_n)\in
 H^1({\cal U}_{n}^{\rm adm};{\mathbb Z}_2).                 \tag{62.26}
\]

\begin{proposition}[Even-doublet representation test]
\label{prop:v18v52-pfaffian-parity}
Tensoring two identical Pfaffian orientation lines adds their
Stiefel--Whitney classes modulo two.  Therefore an even number of identical
\(SU(2)\) doublet blocks cancels the representation-level
\({\mathbb Z}_2\) orientation obstruction.
\end{proposition}

\begin{proof}
For real line bundles,
\[
 w_1(L_1\otimes L_2)=w_1(L_1)+w_1(L_2)
 \quad\hbox{in }H^1(\,\cdot\,;{\mathbb Z}_2).                \tag{62.27}
\]
An even number of equal classes sums to zero.  V44 counts four weak
doublets per generation.
\end{proof}

\begin{firewall}[Finite-family realization]
\label{fire:v18v52-pfaffian-scope}
The even representation count is necessary arithmetic.  One must still
show that the regulated Pfaffian blocks are paired on the same admissible
component with compatible boundary conditions and sector transports.
\end{firewall}

\subsection{Next acquisition}

The global chiral-line problem is now reduced to computable finite-family
objects:
\[
\boxed{
 \widehat P_{-,n},\quad
 {\cal F}_n=-i\operatorname{Tr}(P_ndP_n\wedge dP_n),\quad
 c_{1,n}[\Sigma],\quad
 {\rm Hol}_n(\ell).}                                        \tag{62.28}
\]
V53 should descend these quantities from the admissible gauge-field space
to the gauge-orbit quotient.  It must distinguish vertical gauge curvature,
horizontal physical curvature, and noncontractible gauge-loop holonomy,
then state a sufficient equivariant trivialization theorem.

\begin{assessment}[V52 acquisition and boundary]
\label{ass:v18v52}
V52 constructs the exact Ginsparg--Wilson chiral projector, determinant-line
connection and curvature, Riesz derivative and gap-dependent envelope,
Chern-period obstruction, flat-holonomy obstruction, gauge-loop
integrability condition, and Pfaffian parity test.  It proves that the V44
local anomaly arithmetic is not yet a global fermion measure.  No actual
finite-regulator Chern periods or holonomies have been evaluated.
\end{assessment}

\begin{record}[V52 append record]
\label{rec:v18v52}
V52 is appended to the validated V51 whose installed SHA--256 was
\texttt{6675BE7BDCF33CBA0BAB7DE7D299A6E2A8FDCF0A614EBF656B9BEDC62E605FDC}.
After installation only V52 is the active numeric SM note; the frozen
\texttt{V100\_f17} archives are unchanged.
\end{record}


\section{V53: equivariant determinant-line descent to the gauge-orbit quotient}
\label{sec:v18v53}

\subsection{Free gauge stratum and the combined line}

Let
\[
 X_n={\cal U}_{n}^{\rm adm}
\]
be one admissible gauge-field component, and let the finite lattice gauge
group \({\cal G}_n\) act smoothly and properly.  Begin on a stratum
\(X_n^{\rm fr}\) where the effective action is free, so
\[
 \pi_n:X_n^{\rm fr}\longrightarrow
 B_n^{\rm fr}=X_n^{\rm fr}/{\cal G}_n                       \tag{63.1}
\]
is a principal bundle.

Let \({\cal L}_{f,n}\) be the determinant or Pfaffian line for each chiral
multiplet.  Include conjugate, ghost, and measure lines with their physical
tensor powers and define the combined fermionic line
\[
\boxed{
 {\cal L}_{{\rm tot},n}
 =
 \bigotimes_f{\cal L}_{f,n}^{\otimes\epsilon_f}.}            \tag{63.2}
\]
Its connection and curvature are the signed compiled sums
\[
 {\cal A}_{{\rm tot},n}
 =
 \sum_f\epsilon_f{\cal A}_{f,n},
 \qquad
 {\cal F}_{{\rm tot},n}
 =
 \sum_f\epsilon_f{\cal F}_{f,n}.                             \tag{63.3}
\]
The sums are formed before norms or period bounds.

\subsection{Basic-connection criterion}

For \(\omega\) in the gauge Lie algebra, write
\(\omega^\#\) for the fundamental vertical vector field on \(X_n^{\rm fr}\).
A form \(\alpha\) is \emph{basic} if it is gauge invariant and horizontal:
\[
 R_g^*\alpha=\alpha,
 \qquad
 \iota_{\omega^\#}\alpha=0.                                 \tag{63.4}
\]

\begin{theorem}[Connection descent]
\label{thm:v18v53-connection-descent}
Assume \({\cal L}_{{\rm tot},n}\) has a
\({\cal G}_n\)-equivariant structure.  Its connection descends to a
connection on a line
\[
 \overline{\cal L}_{{\rm tot},n}\longrightarrow B_n^{\rm fr} \tag{63.5}
\]
if and only if its connection form is basic in every equivariant local
frame.  In that case there is a unique quotient curvature
\(\overline{\cal F}_{{\rm tot},n}\) satisfying
\[
\boxed{
 \pi_n^*\overline{\cal F}_{{\rm tot},n}
 =
 {\cal F}_{{\rm tot},n}.}                                   \tag{63.6}
\]
\end{theorem}

\begin{proof}
If the connection is pulled back from the quotient, its form is invariant
and vanishes on vertical vectors, so it is basic.  Conversely, a basic
connection assigns the same horizontal covariant derivative to every two
points in one gauge orbit.  It therefore defines a unique connection on the
associated quotient line.  Exterior differentiation preserves basic forms,
and pullback is injective on forms under a surjective submersion, giving the
unique curvature in (63.6).
\end{proof}

The vertical connection component is the infinitesimal fermion Jacobian.
With the convention of V52, write
\[
 \iota_{\omega^\#}{\cal A}_{{\rm tot},n}
 =
 {\cal J}_{{\rm tot},n}(\omega).                             \tag{63.7}
\]

\begin{corollary}[Infinitesimal anomaly as failure of horizontality]
\label{cor:v18v53-vertical-anomaly}
In a zero-target-Jacobian convention, connection descent requires
\[
\boxed{
 {\cal J}_{{\rm tot},n}(\omega)=0
 \quad\hbox{for every }\omega.}                              \tag{63.8}
\]
If a prescribed local counterterm line has connection
\({\cal A}_{{\rm ct},n}\), the correct condition is that
\({\cal A}_{{\rm tot},n}+{\cal A}_{{\rm ct},n}\) be basic, not that either
summand be basic separately.
\end{corollary}

\begin{proof}
Horizontality in (63.4) is exactly (63.8) by (63.7).  Connections on tensor
products add, so a counterterm may cancel the vertical component only in the
compiled sum.
\end{proof}

\begin{proposition}[Vertical contraction of curvature]
\label{prop:v18v53-curvature-contraction}
If \({\cal A}_{{\rm tot},n}\) is gauge invariant, then
\[
\boxed{
 \iota_{\omega^\#}{\cal F}_{{\rm tot},n}
 =
 -d{\cal J}_{{\rm tot},n}(\omega).}                          \tag{63.9}
\]
Thus vanishing infinitesimal Jacobian makes the curvature horizontal.
\end{proposition}

\begin{proof}
Cartan's formula gives
\[
 0={\cal L}_{\omega^\#}{\cal A}
 =d(\iota_{\omega^\#}{\cal A})
 +\iota_{\omega^\#}d{\cal A}.
\]
Use \(d{\cal A}={\cal F}\) and (63.7).
\end{proof}

\subsection{Equivariant trivialization}

\begin{theorem}[Gauge-equivariant measure criterion on a free stratum]
\label{thm:v18v53-equivariant-trivialization}
Assume the combined connection descends as in Theorem
\ref{thm:v18v53-connection-descent}.  Then:
\begin{enumerate}[leftmargin=2.2em]
\item a smooth nowhere-zero gauge-equivariant measure section exists if and
      only if the quotient line
      \(\overline{\cal L}_{{\rm tot},n}\) is topologically trivial;
\item for a paracompact CW-type quotient this is equivalent to
\[
 c_1(\overline{\cal L}_{{\rm tot},n})=0;                    \tag{63.10}
\]
\item a gauge-equivariant section parallel for the descended connection
      exists if and only if
\[
 \overline{\cal F}_{{\rm tot},n}=0
 \quad\hbox{and}\quad
 {\rm Hol}_{\overline{\cal A}}(\ell)=1
 \ \hbox{for every closed loop }\ell\subset B_n^{\rm fr}.    \tag{63.11}
\]
\end{enumerate}
\end{theorem}

\begin{proof}
Equivariant sections of a descended line are in one-to-one correspondence
with sections of the quotient line.  A complex line has a nowhere-zero
section exactly when it is trivial, and CW-type complex lines are classified
by \(c_1\), proving the first two statements.

A parallel section forces zero curvature and trivial holonomy.  Conversely,
if (63.11) holds, choose a nonzero vector in one fibre and parallel
transport it to every point.  Zero curvature gives local path independence,
and trivial holonomy gives global path independence.  The resulting
nowhere-zero parallel section pulls back equivariantly.
\end{proof}

\begin{firewall}[A physical measure need not be parallel]
\label{fire:v18v53-parallel}
A gauge-equivariant fermion measure requires an equivariant nonzero section
with the prescribed anomaly connection.  It need not make that connection
zero.  The stronger parallel criterion (63.11) is appropriate only after
all prescribed Jacobian and counterterm connections have been absorbed into
the combined line.
\end{firewall}

\subsection{Stabilizers and singular orbit strata}

If the action of \({\cal G}_n\) is not free, let
\[
 {\cal G}_{n,U}
 =
 \{g\in{\cal G}_n:g\cdot U=U\}                              \tag{63.12}
\]
be the stabilizer.

\begin{theorem}[Stabilizer-character obstruction]
\label{thm:v18v53-stabilizer}
If the combined determinant line descends to an honest line over a
neighbourhood of the orbit \([U]\), every element of
\({\cal G}_{n,U}\) must act trivially on the fibre
\(({\cal L}_{{\rm tot},n})_U\).  A nontrivial stabilizer character is an
explicit obstruction to an honest quotient measure.
\end{theorem}

\begin{proof}
Every stabilizer element represents the identity morphism of the base point
\([U]\).  In a pullback of an honest quotient line, it must therefore act as
the identity on the fibre.  A nontrivial character contradicts this.
\end{proof}

If the character is nontrivial, one may obtain an orbifold line rather than
an honest line.  That object does not satisfy \(\mathbf{C1}\) of V51 unless
the physical quotient and its reconstruction are reformulated accordingly.

\subsection{Spanning-tree reduction of finite lattice gauge space}

Let \(\mathsf G=(V,E)\) be a finite connected oriented lattice graph and
let \(K\) be a compact gauge group.  The link space is
\[
 X=K^E.                                                      \tag{63.13}
\]
Choose a root \(v_0\) and let the based gauge group be
\[
 {\cal G}_0
 =
 \{(g_v)_{v\in V}\in K^V:g_{v_0}=1\}.                       \tag{63.14}
\]
It acts by
\[
 U_{(v,w)}\longmapsto g_vU_{(v,w)}g_w^{-1}.                 \tag{63.15}
\]

\begin{theorem}[Exact tree gauge slice]
\label{thm:v18v53-tree-slice}
Fix a spanning tree \(T\subset E\).  Every based gauge orbit contains a
unique representative satisfying
\[
 U_e=1\qquad(e\in T).                                       \tag{63.16}
\]
Consequently,
\[
\boxed{
 K^E/{\cal G}_0
 \cong
 K^{E\setminus T}
 \cong K^{\,|E|-|V|+1}.}                                    \tag{63.17}
\]
\end{theorem}

\begin{proof}
Starting at the root, define \(g_v\) recursively along the unique tree path
so that every transformed tree edge is the identity.  This gives existence.
If two based gauge transformations preserve (63.16), (63.15) forces their
values to agree across every tree edge; connectedness and the root value
make both transformations identical.  The non-tree links are unconstrained
coordinates, proving (63.17).
\end{proof}

\begin{corollary}[Residual global conjugation]
\label{cor:v18v53-global-conjugation}
Quotienting by the full gauge group adds the residual action
\[
 (U_e)_{e\notin T}
 \longmapsto
 (gU_eg^{-1})_{e\notin T}                                   \tag{63.18}
\]
of one global \(K\).  Its reducible configurations create precisely the
stabilizer strata of Theorem \ref{thm:v18v53-stabilizer}.
\end{corollary}

\begin{proof}
The quotient \(K^V/{\cal G}_0\) is the root copy of \(K\), and it acts on
the tree-gauge coordinates by simultaneous conjugation.
\end{proof}

\subsection{Topological consequences and the admissibility caveat}

For \(K=SU(N)\), one has
\[
 H^1(SU(N);{\mathbb Z})=H^2(SU(N);{\mathbb Z})=0.            \tag{63.19}
\]
K\"unneth therefore gives the same vanishing in degree two for a finite
product of \(SU(N)\) factors.

\begin{corollary}[Unrestricted based semisimple link space]
\label{cor:v18v53-semimple-triviality}
On the unrestricted based quotient
\[
 SU(3)^{b}\times SU(2)^{b},\qquad b=|E|-|V|+1,               \tag{63.20}
\]
every complex determinant line is topologically trivial.  For the
hypercharge factor \(U(1)^b\), degree-two classes may occur when \(b\geq2\).
\end{corollary}

\begin{proof}
Equation (63.19) and the degree-two K\"unneth formula prove the semisimple
claim.  For the torus,
\[
 H^2(U(1)^b;{\mathbb Z})
 \cong\bigwedge\nolimits^2{\mathbb Z}^b,                    \tag{63.21}
\]
which is nonzero for \(b\geq2\).
\end{proof}

\begin{firewall}[Admissibility changes topology]
\label{fire:v18v53-admissibility}
The overlap admissibility condition removes gauge configurations from the
unrestricted product (63.17).  The remaining component can acquire new
cycles and distinct topological sectors.  Corollary
\ref{cor:v18v53-semimple-triviality} therefore does not prove triviality of
the actual overlap determinant line.
\end{firewall}

\subsection{Cofinal equivariant transport}

Let \(T_{n+1,n}:X_n^{\rm adm}\to X_{n+1}^{\rm adm}\) be the declared gauge
field transport, compatible with gauge homomorphisms.  A determinant-line
transport is a fibrewise unitary map
\[
 \tau_{n+1,n}:{\cal L}_{{\rm tot},n}
 \longrightarrow
 T_{n+1,n}^*{\cal L}_{{\rm tot},n+1}.                       \tag{63.22}
\]
Define its connection defect
\[
 {\cal R}_{n}^{\cal A}
 =
 \tau_{n+1,n}^{-1}
 T_{n+1,n}^*{\cal A}_{{\rm tot},n+1}
 \tau_{n+1,n}
 -
 {\cal A}_{{\rm tot},n}.                                    \tag{63.23}
\]

\begin{proposition}[Holonomy transport bound]
\label{prop:v18v53-holonomy-transport}
For a transported loop \(\ell_n\),
\[
\boxed{
 \left|
 {\rm arg\,Hol}_{n+1}(T_{n+1,n}\ell_n)
 -
 {\rm arg\,Hol}_{n}(\ell_n)
 \right|
 \leq
 \int_{\ell_n}|{\cal R}_n^{\cal A}|
 \quad\operatorname{mod}2\pi.}                              \tag{63.24}
\]
If the right sides are summable uniformly on a generating loop bank,
the corresponding holonomies have cofinal limits.
\end{proposition}

\begin{proof}
The ratio of the two parallel transports is the exponential of the
integral of the connection difference (63.23).  Taking its argument gives
(63.24), and summability makes the sequence Cauchy on every generator.
\end{proof}

\subsection{Next acquisition}

The determinant-line problem has now been moved to the finite based quotient
and separated into:
\[
\boxed{
 \text{vertical Jacobian, quotient Chern periods, stabilizer characters,
 quotient holonomies, and cofinal connection defects}.}     \tag{63.25}
\]
V54 should construct a finite generating cycle/loop panel on one admissible
component and prove that vanishing on that panel is sufficient under an
explicit homology-generation hypothesis.  This will make a negative
holonomy or Chern period an executable obstruction rather than an abstract
warning.

\begin{assessment}[V53 acquisition and boundary]
\label{ass:v18v53}
V53 proves the basic-connection descent criterion, the quotient
trivialization and holonomy theorem, the stabilizer obstruction, the exact
spanning-tree gauge slice, and a cofinal holonomy transport bound.  It
identifies the simpler topology of the unrestricted based semisimple link
space while retaining the overlap-admissibility caveat.  No actual
admissible-component generator panel or determinant-line period has yet
been computed.
\end{assessment}

\begin{record}[V53 append record]
\label{rec:v18v53}
V53 is appended to the validated V52 whose installed SHA--256 was
\texttt{0912C0DAC74091CE9265B5AAB49F272CCE47CB402505916F9095D6C49283F59E}.
After installation only V53 is the active numeric SM note; the frozen
\texttt{V100\_f17} archives are unchanged.
\end{record}


\section{V54: finite Cech and holonomy certificates for the chiral measure line}
\label{sec:v18v54}

\subsection{Finite-topology hypothesis}

Let \(B_n\) be one free admissible gauge-orbit component from V53 and let
\[
 \overline{\cal L}_n\longrightarrow B_n                    \tag{64.1}
\]
be the descended compiled determinant line.  Assume \(B_n\) has the
homotopy type of a finite CW complex and choose a finite good cover
\[
 {\mathfrak U}_n=\{U_i\}_{i=1}^{N_n},                        \tag{64.2}
\]
so every nonempty finite intersection is contractible.  This is an explicit
hypothesis; no finite triangulation of the actual overlap-admissible quotient
has yet been supplied.

Choose unit local frames \(s_i\) with transition functions
\[
 s_j=s_i g_{ij},
 \qquad
 g_{ij}:U_i\cap U_j\longrightarrow U(1).                    \tag{64.3}
\]
On each nonempty double intersection choose a real lift
\[
 g_{ij}=e^{i\theta_{ij}}.                                   \tag{64.4}
\]

\subsection{The integer Cech cocycle}

\begin{lemma}[Triple-overlap integer]
\label{lem:v18v54-triple-integer}
On every nonempty triple intersection,
\[
\boxed{
 n_{ijk}
 =
 \frac{\theta_{ij}+\theta_{jk}+\theta_{ki}}{2\pi}
 \in{\mathbb Z}.}                                           \tag{64.5}
\]
The integer is constant on each connected triple intersection.
\end{lemma}

\begin{proof}
The transition cocycle identity gives
\[
 g_{ij}g_{jk}g_{ki}=1.
\]
Using (64.4), the numerator in (64.5) is an integer multiple of \(2\pi\).
It is continuous and integer valued on a connected set, hence constant.
\end{proof}

\begin{theorem}[Finite Cech representative of the determinant line]
\label{thm:v18v54-cech-class}
The integers \(n=(n_{ijk})\) form a Cech two-cocycle.  Its class
\[
\boxed{
 [n]\in\check H^2(B_n;{\mathbb Z})
 \cong H^2(B_n;{\mathbb Z})}                                \tag{64.6}
\]
is the first Chern class \(c_1(\overline{\cal L}_n)\).  In particular,
\[
\boxed{
 \overline{\cal L}_n\ \hbox{is topologically trivial}
 \quad\Longleftrightarrow\quad
 [n]=0.}                                                    \tag{64.7}
\]
\end{theorem}

\begin{proof}
Changing a logarithm lift by
\(\theta_{ij}\mapsto\theta_{ij}+2\pi m_{ij}\) changes \(n\) by the integer
Cech coboundary \(\delta m\).  Changing local frames multiplies the
transition functions by a Cech coboundary and produces the same change of
representative.  Thus \([n]\) is intrinsic.  The exponential exact sequence
\[
 0\longrightarrow{\mathbb Z}
 \longrightarrow{\mathbb R}
 \stackrel{e^{2\pi i(\cdot)}}{\longrightarrow}
 U(1)\longrightarrow0                                      \tag{64.8}
\]
sends the transition one-cocycle to \([n]\); this connecting class is
precisely \(c_1\).  Complex line bundles over a paracompact CW-type base are
classified by \(c_1\), proving (64.7).
\end{proof}

\begin{corollary}[Finite integer decision]
\label{cor:v18v54-smith}
After replacing the good cover by its finite nerve, whether \([n]=0\) is
decidable by integer linear algebra: \(n\) is trivial exactly when the
integer system
\[
 \delta m=n                                                  \tag{64.9}
\]
has a solution.  Smith normal form of the finite coboundary matrix gives a
certificate and detects both free and torsion components of \(c_1\).
\end{corollary}

\begin{proof}
The nerve theorem identifies the Cech complex of the good cover with the
cochain complex of a finite simplicial complex.  Equation (64.9) is a finite
integer linear system.  Smith normal form decides membership in the image
of the coboundary map, including its torsion quotient.
\end{proof}

\subsection{Why curvature periods alone miss torsion}

The universal-coefficient sequence contains
\[
 0\longrightarrow
 \operatorname{Ext}(H_1(B_n;{\mathbb Z}),{\mathbb Z})
 \longrightarrow
 H^2(B_n;{\mathbb Z})
 \longrightarrow
 \operatorname{Hom}(H_2(B_n;{\mathbb Z}),{\mathbb Z})
 \longrightarrow0.                                         \tag{64.10}
\]

\begin{proposition}[Period limitation]
\label{prop:v18v54-period-limitation}
Integrals of the curvature over closed two-cycles determine the image of
\(c_1\) in
\(\operatorname{Hom}(H_2,{\mathbb Z})\), but do not in general detect the
torsion subgroup
\(\operatorname{Ext}(H_1,{\mathbb Z})\).  The integer cocycle test
(64.9) detects both.
\end{proposition}

\begin{proof}
The de Rham class of curvature is the real image of \(c_1\).  Tensoring with
\({\mathbb R}\) kills torsion, so curvature periods see only the free
component.  The integral Cech class retains torsion.
\end{proof}

\begin{countertheorem}[All curvature periods may vanish on a nontrivial
line]
\label{cth:v18v54-torsion-line}
If \(H^2(B_n;{\mathbb Z})\) contains nonzero torsion, there is a nontrivial
complex line with a flat connection.  Its curvature and every curvature
period vanish, while its Cech class is nonzero.
\end{countertheorem}

\begin{proof}
Choose a nonzero torsion class in \(H^2(B_n;{\mathbb Z})\) and the
corresponding line bundle.  Torsion classes have zero real image and admit a
flat unitary connection.  Classification by \(c_1\) keeps the line
topologically nontrivial.
\end{proof}

\subsection{Connection compatibility with the Cech data}

Let \({\cal A}_i\) be the real connection one-form in frame \(s_i\).  From
V52,
\[
 {\cal A}_j-{\cal A}_i=d\theta_{ij}.                         \tag{64.11}
\]
Hence
\[
 d{\cal A}_i=d{\cal A}_j={\cal F}_n                         \tag{64.12}
\]
on overlaps.

\begin{theorem}[Differential-character loop formula]
\label{thm:v18v54-loop-formula}
Let \(\ell\) be a piecewise smooth closed loop subdivided into arcs
\(\ell_a\subset U_{i_a}\), with transition point \(x_a\) from
\(U_{i_a}\) to \(U_{i_{a+1}}\).  Then
\[
\boxed{
 {\rm Hol}_n(\ell)
 =
 \exp\left(
 i\sum_a\int_{\ell_a}{\cal A}_{i_a}
 +
 i\sum_a\theta_{i_ai_{a+1}}(x_a)
 \right).}                                                  \tag{64.13}
\]
The result is independent of subdivision, charts, logarithm lifts, and
local frames.
\end{theorem}

\begin{proof}
Parallel transport on each arc contributes
\(\exp(i\int{\cal A}_{i_a})\); changing frames at \(x_a\) contributes the
transition phase.  Equation (64.11) cancels changes caused by moving a
transition point.  The cocycle identity (64.5) makes refinement across a
triple overlap invariant, and integer changes of the logarithm do not change
the exponential.  Frame changes telescope around the closed loop.
\end{proof}

\subsection{Finite flat-holonomy panel}

Assume the compiled descended connection is flat:
\[
 {\cal F}_n=0.                                               \tag{64.14}
\]
Then holonomy is a character
\[
 \chi_n:H_1(B_n;{\mathbb Z})\longrightarrow U(1).           \tag{64.15}
\]

\begin{theorem}[Finite generator sufficiency for a parallel measure]
\label{thm:v18v54-finite-holonomy}
Let
\[
 [\ell_1],\ldots,[\ell_m]
\]
generate the finitely generated abelian group \(H_1(B_n;{\mathbb Z})\).
Under (64.14), the following are equivalent:
\begin{enumerate}[leftmargin=2.2em]
\item the connection has a global nonzero parallel section;
\item
\[
 {\rm Hol}_n(\ell_j)=1
 \qquad(1\leq j\leq m).                                     \tag{64.16}
\]
\end{enumerate}
\end{theorem}

\begin{proof}
A parallel section has trivial holonomy on every loop, so it implies
(64.16).  Conversely, flat holonomy factors through (64.15).  A character
which is one on a generating set is the trivial character.  The V53
parallel-section theorem then applies.
\end{proof}

\begin{corollary}[Prescribed gauge Jacobian]
\label{cor:v18v54-prescribed-character}
If the desired gauge transformation law gives a character
\(\chi_n^{\rm phys}\), tensor the determinant line with the inverse
one-dimensional Jacobian line.  The combined flat line admits a parallel
equivariant section exactly when
\[
 {\rm Hol}_n(\ell_j)=\chi_n^{\rm phys}(\ell_j)
 \quad(1\leq j\leq m).                                      \tag{64.17}
\]
\end{corollary}

\begin{proof}
The holonomy of a tensor product is the product of the two holonomies.
Apply Theorem \ref{thm:v18v54-finite-holonomy} to the combined line.
\end{proof}

\subsection{Free Chern-period panel}

Let \([\Sigma_1],\ldots,[\Sigma_b]\) map to a basis of the free part of
\(H_2(B_n;{\mathbb Z})\).  Define
\[
 c_{n,j}
 =
 \frac{1}{2\pi}\int_{\Sigma_j}{\cal F}_n\in{\mathbb Z}.      \tag{64.18}
\]

\begin{proposition}[Certified vanishing of free Chern data]
\label{prop:v18v54-integer-interval}
If a rigorous interval computation gives
\[
 c_{n,j}\in I_{n,j}\subset\left(-\frac12,\frac12\right)      \tag{64.19}
\]
for every \(j\), then
\[
 c_{n,j}=0
\qquad(1\leq j\leq b).                                      \tag{64.20}
\]
\end{proposition}

\begin{proof}
Each \(c_{n,j}\) is an integer by the V52 Chern-period theorem.  Zero is the
only integer in the interval (64.19).
\end{proof}

\begin{firewall}[Free-period certificate is incomplete with torsion]
\label{fire:v18v54-free-period}
Equations (64.19)--(64.20) prove only the free part of \(c_1=0\).  The Smith
certificate (64.9) or an equivalent torsion computation remains compulsory.
\end{firewall}

\subsection{Obstruction certificates are one-sided under intervals}

\begin{proposition}[Finite obstruction witnesses]
\label{prop:v18v54-obstruction}
Each of the following proves failure of the corresponding global measure
condition:
\begin{enumerate}[leftmargin=2.2em]
\item a Smith-normal-form certificate that \(n\notin\operatorname{im}\delta\);
\item an interval for one Chern number contained in
      \((k-\tfrac12,k+\tfrac12)\) with \(k\ne0\);
\item a certified loop phase interval disjoint from
      \(2\pi{\mathbb Z}\) after subtracting the prescribed physical
      Jacobian phase;
\item a nontrivial stabilizer character from V53.
\end{enumerate}
\end{proposition}

\begin{proof}
The first contradicts (64.7), the second gives a nonzero Chern period, the
third contradicts (64.17), and the fourth invokes the V53 stabilizer
obstruction.
\end{proof}

\begin{firewall}[An interval containing zero is not an exact holonomy proof]
\label{fire:v18v54-zero-interval}
A numerical phase interval which contains \(0\) does not prove that the
holonomy is exactly one.  Exact trivial holonomy requires a symbolic
identity, a topological argument, or a convergent cofinal estimate whose
limiting statement is all that the continuum theorem needs.
\end{firewall}

\subsection{Cofinal generator transport}

Suppose generator panels are transported between cutoffs:
\[
 T_{n+1,n}(\Sigma_{n,j})\sim\Sigma_{n+1,j},
 \qquad
 T_{n+1,n}(\ell_{n,k})\sim\ell_{n+1,k}.                      \tag{64.21}
\]

\begin{theorem}[Stable finite-panel landing]
\label{thm:v18v54-panel-landing}
Assume:
\begin{enumerate}[leftmargin=2.2em]
\item the transported cycles and loops generate the stable free
      \(H_2\) and full \(H_1\) groups;
\item the integral Cech torsion classes transport isomorphically;
\item the V53 connection defects are summable on the loop bank;
\item curvature transport defects have summable integrals on the cycle bank.
\end{enumerate}
Then the Chern integers are eventually constant and the generator
holonomies have unique cofinal limits.  A zero Cech class and trivial
limiting physical holonomy yield a cofinally consistent measure line on
this stable component.
\end{theorem}

\begin{proof}
An integer sequence whose adjacent differences tend to zero is eventually
constant; the summable curvature transport makes the Chern-period
differences tend to zero.  The Cech assumption controls torsion.  The V53
holonomy transport bound makes each loop phase Cauchy.  Generation then
extends the limiting identities from the panel to all stable cycles and
loops.
\end{proof}

\subsection{Next finite-family input}

The global line row is now decidable once the following are supplied for one
literal admissible component:
\[
\boxed{
 \text{finite good cover, transition phases, nerve coboundary matrices,
 free two-cycles, and first-homology loop generators}.}      \tag{64.22}
\]
V55 should investigate whether the overlap admissibility region for the
actual finite SM lattice admits a uniform finite-cover complexity and a
stable generator transport.  If that geometric input is unavailable, it
should return the first topology-change or unbounded-complexity witness
rather than assume a stable quotient.

\begin{assessment}[V54 acquisition and boundary]
\label{ass:v18v54}
V54 constructs a finite integral Cech certificate for determinant-line
triviality, proves why curvature periods miss torsion, gives an exact
chartwise holonomy formula and finite generator theorem, and supplies
one-sided obstruction certificates plus cofinal transport conditions.  No
good cover, homology generator bank, Smith certificate, or holonomy value
has yet been produced for the literal SM overlap family.
\end{assessment}

\begin{record}[V54 append record]
\label{rec:v18v54}
V54 is appended to the validated V53 whose installed SHA--256 was
\texttt{B1E892DC727282799621FD2E2E4D1938F85B7BEDB5D16238EE0611C9F31A7293}.
After installation only V54 is the active numeric SM note; the frozen
\texttt{V100\_f17} archives are unchanged.
\end{record}


\section{V55: fixed-cutoff semialgebraic decidability of the overlap measure line}
\label{sec:v18v55}

\subsection{Algebraic regulator hypothesis}

Fix one finite lattice graph.  Embed each compact gauge factor in a real
matrix space by real and imaginary parts.  Assume:
\begin{enumerate}[leftmargin=2.2em]
\item the link group \(K\) is a compact real algebraic matrix group;
\item all regulator coefficients and admissibility thresholds are real
      algebraic numbers;
\item the Hermitian overlap kernel \(H_n(U)\) has polynomial or rational
      matrix entries in the link coordinates, with denominators nonzero on
      the admissible set;
\item admissibility is specified by finitely many polynomial strict
      inequalities and includes a spectral gap
      \(0\notin\operatorname{spec}H_n(U)\).
\end{enumerate}
These hypotheses cover a finite Wilson-kernel construction with a
polynomial plaquette admissibility condition.  They are not asserted
uniformly in the cutoff.

\subsection{Semialgebraic gauge space}

\begin{theorem}[Fixed-cutoff admissible space is semialgebraic]
\label{thm:v18v55-admissible-semialgebraic}
Under the algebraic regulator hypothesis, the link space, admissible subset,
and each of its connected components are semialgebraic sets.
\end{theorem}

\begin{proof}
For \(U(N)\), the equations
\[
 U^*U=I                                                        \tag{65.1}
\]
are polynomial in the real and imaginary matrix entries.  For \(SU(N)\),
add the two polynomial equations
\[
 \operatorname{Re}\det U=1,\qquad
 \operatorname{Im}\det U=0.                                 \tag{65.2}
\]
The circle \(U(1)\) is polynomially defined by \(x^2+y^2=1\).
A finite product over links remains algebraic.  Plaquette products are
polynomial matrix expressions, and the declared admissibility and
denominator conditions are finite Boolean combinations of polynomial
equalities and inequalities.  The gap condition can be written
\[
 \det H_n(U)\ne0.                                            \tag{65.3}
\]
Hence the admissible set is semialgebraic.  A semialgebraic set has finitely
many connected components, each semialgebraic.
\end{proof}

\begin{corollary}[Finite CW type]
\label{cor:v18v55-finite-cw}
Every fixed-cutoff admissible component has the homotopy type of a finite CW
complex.  Its integral homology groups are finitely generated.
\end{corollary}

\begin{proof}
The semialgebraic triangulation theorem gives a finite triangulation by
locally closed simplices.  The resulting finite simplicial complex gives
the stated homotopy type and finitely generated cellular homology.
\end{proof}

\subsection{Semialgebraic based quotient}

Choose the V53 spanning tree and impose tree gauge.  The based quotient is a
semialgebraic subset of
\[
 K^{\,|E|-|V|+1}.                                           \tag{65.4}
\]
The remaining global \(K\)-action is simultaneous conjugation.

\begin{theorem}[Semialgebraic orbit quotient]
\label{thm:v18v55-orbit-quotient}
For a compact real algebraic group acting algebraically on a semialgebraic
set \(X\), there are finitely many invariant polynomials
\[
 p_1,\ldots,p_N                                             \tag{65.5}
\]
whose joint map separates compact orbits.  The orbit space is homeomorphic
to the semialgebraic image
\[
 p(X)\subset{\mathbb R}^N.                                  \tag{65.6}
\]
Consequently every fixed-cutoff gauge-orbit stratum has finite CW type and
algorithmically computable integral homology.
\end{theorem}

\begin{proof}
Finite generation of the invariant polynomial algebra gives (65.5).
For a compact group all orbits are closed, and invariant polynomials
separate distinct closed orbits.  Thus the induced map from \(X/K\) to
\(p(X)\) is continuous and injective.  Properness of the compact action
gives the quotient topology and makes this map a homeomorphism onto its
image.  Tarski--Seidenberg elimination makes \(p(X)\) semialgebraic.
Apply semialgebraic triangulation and cellular integer homology.
\end{proof}

\begin{firewall}[Strata]
\label{fire:v18v55-strata}
The full quotient may have stabilizer singularities.  The determinant line
descends as an honest line only on strata passing the V53 stabilizer test.
Semialgebraic decidability does not remove an obstructing stabilizer
character.
\end{firewall}

\subsection{Semialgebraic spectral projector}

The overlap projector is an analytic function of \(H_n\), but its graph has
an algebraic characterization.

\begin{theorem}[Semialgebraic projector graph]
\label{thm:v18v55-projector-graph}
On a constant-index gapped component, the graph of the negative spectral
projector
\[
 P_n(U)=1_{(-\infty,0)}(H_n(U))                              \tag{65.7}
\]
is semialgebraic.  It is the unique matrix \(P\) satisfying
\[
\boxed{
 P^*=P,\quad P^2=P,\quad PH_n=H_nP,}                         \tag{65.8}
\]
together with
\[
 -PH_nP\succ0\ \hbox{on }\operatorname{Ran}P,\qquad
 (I-P)H_n(I-P)\succ0\ \hbox{on }\operatorname{Ran}(I-P).     \tag{65.9}
\]
\end{theorem}

\begin{proof}
The spectral projector satisfies (65.8)--(65.9).  Conversely, commutation
splits the space into the ranges of \(P\) and \(I-P\); the two positivity
conditions force these to be precisely the negative and positive spectral
subspaces.  Hermiticity, idempotence, and commutation are polynomial
equations.  Positive definiteness on a fixed-rank range can be expressed by
the positivity of the relevant principal minors after a finite pivot
decomposition, hence semialgebraically.  Uniqueness and
Tarski--Seidenberg elimination show that the graph of \(U\mapsto P_n(U)\)
is semialgebraic.
\end{proof}

\subsection{A finite pivot atlas for the chiral bundle}

Let \(r=\operatorname{rank}P_n\), and for every \(r\)-element coordinate
set \(I\), let \(E_I:{\mathbb C}^r\to{\mathbb C}^N\) be the coordinate
inclusion.  Define
\[
 G_I(U)=E_I^*P_n(U)E_I.                                     \tag{65.10}
\]
On
\[
 {\cal O}_I=\{U:\det G_I(U)>0\},                             \tag{65.11}
\]
put
\[
 v_I(U)
 =
 P_n(U)E_I\,G_I(U)^{-1/2}.                                  \tag{65.12}
\]

\begin{theorem}[Finite chiral-frame atlas]
\label{thm:v18v55-pivot-atlas}
The nonempty sets \({\cal O}_I\) form a finite semialgebraic cover of the
constant-rank component.  Each \(v_I\) is an orthonormal frame for
\(\operatorname{Ran}P_n\), and on overlaps
\[
\boxed{
 g_{IJ}=v_I^*v_J\in U(r)}                                   \tag{65.13}
\]
is the transition matrix.  The determinant transition
\[
 \det g_{IJ}:{\cal O}_I\cap{\cal O}_J\longrightarrow U(1)   \tag{65.14}
\]
defines the determinant line of V52.
\end{theorem}

\begin{proof}
Since \(P_n\) has rank \(r\), at least one \(r\times r\) coordinate minor of
an orthonormal frame is nonsingular; equivalently some Gram matrix
\(G_I\) is positive definite.  Hence the sets cover.  There are only
\(\binom Nr\) choices.  Equation (65.12) gives
\[
 v_I^*v_I
 =
 G_I^{-1/2}E_I^*P_n^2E_IG_I^{-1/2}=I,                       \tag{65.15}
\]
and its range is the range of \(P_n\).  Two orthonormal frames of the same
space differ by the unitary matrix (65.13).  Positive square root and
inverse are semialgebraic on the positive-definite cone through their
unique defining equations, so the atlas data are semialgebraic.
\end{proof}

\begin{corollary}[Finite refinement to the V54 data]
\label{cor:v18v55-good-refinement}
Semialgebraic triangulation and barycentric subdivision refine the pivot
atlas to a finite cover on which all required transition logarithms,
connected triple-overlap components, cellular two-cocycles, and homology
generators of V54 can be represented by finite data.
\end{corollary}

\begin{proof}
Triangulate compatibly with the finite family of semialgebraic pivot sets
and their intersections.  A sufficiently fine barycentric star refinement
has contractible simplicial intersections on the finite CW model.  The
transition maps restrict to these cells, and cellular chains provide finite
integer boundary matrices and homology generators.
\end{proof}

\subsection{Fixed-cutoff decision theorem}

\begin{theorem}[Decidability of the finite chiral-line gate]
\label{thm:v18v55-fixed-decision}
Under the algebraic regulator hypothesis, there is a finite exact procedure
which, on each fixed admissible free-orbit stratum:
\begin{enumerate}[leftmargin=2.2em]
\item computes a finite CW model and integer boundary matrices;
\item constructs the pivot transition functions (65.13)--(65.14);
\item computes the integral Cech class of the compiled determinant line;
\item decides whether that class is zero by Smith normal form;
\item computes a finite generating bank for \(H_1\) and \(H_2\);
\item reduces prescribed gauge-loop integrability and flat holonomy to the
      finite V54 panels.
\end{enumerate}
It therefore returns either a finite topological obstruction or a
topologically trivial finite-cutoff determinant line.
\end{theorem}

\begin{proof}
Quantifier elimination and semialgebraic triangulation provide items one
and five.  Theorem \ref{thm:v18v55-pivot-atlas} provides finite transition
data.  The V54 Cech construction converts them into a finite integer
two-cocycle, and Smith normal form decides triviality, proving items
three--four.  The V54 generator theorems give item six.  Every step is
finite at fixed matrix size and graph size.
\end{proof}

\begin{firewall}[Topological decision is not a physical measure]
\label{fire:v18v55-topology-not-measure}
A zero Cech class supplies a global nonzero line section.  The section must
still be chosen gauge-equivariantly with the correct Jacobian, locality,
smoothness, reflection, and cutoff-transport bounds.  V55 decides the
topological gate only.
\end{firewall}

\subsection{No uniform complexity follows}

Let \(N_n^{\rm cell}\) be the number of cells in a compatible triangulation,
\(N_n^{\rm piv}\) the number of active pivot charts, and
\[
 {\cal C}_n^{\rm top}
 =
 N_n^{\rm cell}+N_n^{\rm piv}
 +\|\partial_{1,n}\|_0+\|\partial_{2,n}\|_0                 \tag{65.16}
\]
be a crude certificate-size stock, where \(\|\cdot\|_0\) counts nonzero
integer boundary entries.

\begin{countertheorem}[Per-cutoff decidability gives no cofinal bound]
\label{cth:v18v55-no-uniform-complexity}
Finiteness of \({\cal C}_n^{\rm top}\) for every \(n\) does not imply
\[
 \sup_n{\cal C}_n^{\rm top}<\infty                           \tag{65.17}
\]
or any computable growth rate sufficient for cofinal transport.
\end{countertheorem}

\begin{proof}
The assertion that each term of an arbitrary integer sequence is finite
places no bound on its growth.  Semialgebraic triangulation algorithms also
have complexity depending on the number of variables, degrees, and
inequalities, all of which grow with the lattice.  Therefore fixed-level
termination alone implies neither (65.17) nor a uniform rate.
\end{proof}

\begin{assessment}[Transfer from YM V8]
\label{ass:v18v55-ym-transfer}
V55 is the topological analogue of the YM V8 existence/feasibility split.
It proves that a finite exact certificate exists at every fixed cutoff under
the algebraic hypotheses.  It does not produce a feasible certificate or a
uniform cofinal ledger.
\end{assessment}

\subsection{A cofinal sufficiency criterion}

\begin{theorem}[Uniform topology-and-line capsule]
\label{thm:v18v55-uniform-capsule}
Suppose, on a cofinal sequence:
\begin{enumerate}[leftmargin=2.2em]
\item the admissible components and free-orbit strata admit compatible
      finite CW models with generator transports;
\item their cellular boundary matrices stabilize up to elementary integer
      equivalence on the relevant degrees;
\item the compiled determinant Cech classes vanish;
\item global sections have summable connection, locality, and reflection
      transport defects;
\item stabilizer characters and physical gauge-loop residuals vanish.
\end{enumerate}
Then the chiral measure lines form a cofinally consistent topologically
trivial equivariant family on those components.
\end{theorem}

\begin{proof}
The first two hypotheses identify the stable homology and Cech groups.
The third gives a global section at every level by V54.  The fourth makes
the transported sections Cauchy after their phases are fixed on one base
fibre.  The fifth permits descent across gauge orbits and stabilizer strata.
The limiting compatible family is therefore an equivariant trivialization
with the stated locality and reflection control.
\end{proof}

\begin{firewall}[Unpopulated uniform rows]
\label{fire:v18v55-uniform-unpopulated}
No compatible CW models, stabilized boundary matrices, zero compiled Cech
classes, or summable section defects have been computed for the literal
cofinal Standard-Model overlap family.  Theorem
\ref{thm:v18v55-uniform-capsule} is conditional.
\end{firewall}

\subsection{Checkpoint and next step}

V55 reaches the next five-version boundary.  The active Yang--Mills and
Navier--Stokes notes must be inspected before V56.  After that audit, the
next upstream target is the analytic construction of a global section from
a zero Cech class with quantitative derivative bounds.  Topological
triviality without a controlled connection norm is insufficient for the V45
Ward collar and the V51 continuum capsule.

\begin{assessment}[V55 acquisition and boundary]
\label{ass:v18v55}
V55 proves fixed-cutoff semialgebraicity, finite CW type, a finite
chiral-projector pivot atlas, and decidability of the integral
determinant-line class, including torsion.  It proves that this fixed-level
result supplies neither feasible complexity nor cofinal control, and states
the exact uniform capsule that would.  No literal certificate has yet been
executed.
\end{assessment}

\begin{record}[V55 append record]
\label{rec:v18v55}
V55 is appended to the validated V54 whose installed SHA--256 was
\texttt{601EC5620996874DB38D71AEECE46CB8C2C561A05C6F5B4539A4423EBC0D44D8}.
After installation only V55 is the active numeric SM note; the frozen
\texttt{V100\_f17} archives are unchanged.
\end{record}

\section{V56: quantitative trivialization of the compiled chiral measure line}
\label{sec:v18v56}

V55 made the finite-cutoff topological question decidable under an
algebraic regulator hypothesis.  This note performs the next analytic
step.  On a quotient stratum where the compiled determinant line is
topologically trivial, it constructs a global unit section by an explicit
Cech homotopy and estimates the resulting Berry connection.  It then
separates two inputs which must not be conflated: vanishing of the integral
line class, and uniform analytic control of a connection representing that
class.

\subsection{Mandatory V55 companion audit}

The checkpoint audit used the active companion files
\[
 \texttt{Renewal\_Yang\_Mills\_Mass\_Gap\_Research\_Notes\_V8.tex},
 \qquad
 \texttt{Renewal\_NS\_Stability\_Research\_Notes\_V26.tex}.             \tag{66.1}
\]
Their byte counts and SHA--256 hashes were respectively
\[
\begin{split}
 &(145306,
 \texttt{23EC89B8E57800E424713FE450882682DF6009D431C38B25A6D6E0F281A9B0C0}),\\
 &(278283,
 \texttt{82C887F8C2C69E4F28FAD7C3DC8DE5B9099CB5A4BF431EBA1FFF3E91935978AD}).
                                                                    \tag{66.2}
\end{split}
\]
YM V8 retains the live complex phase, pairs the reflection-related boxes,
and applies triangle estimates only to the inverse-defect radii.  The
transferable rule is to compile cancelling chiral phases before estimating
their size.  Its finite rational certificate remains conditional on
feasibility and sign, so no YM conclusion is imported.  NS V26 supplies no
stronger estimate for the present line-bundle problem: its rank-one gains
remain too small to alter the direction.

\subsection{Good-cover data and a real Cech homotopy}

Let \(B\) be a compact smooth free-orbit stratum of the admissible gauge
quotient.  Let \(L\to B\) be the compiled Hermitian determinant line, and
choose a finite good cover \(\{U_i\}_{i=1}^M\), local unit sections \(s_i\),
and a smooth partition of unity \(\{\chi_i\}\) with
\(\operatorname{supp}\chi_i\Subset U_i\).  On a nonempty overlap choose
antisymmetric real lifts
\[
 s_j=s_i e^{i\theta_{ij}},\qquad
 \theta_{ji}=-\theta_{ij}.                                  \tag{66.3}
\]
On every connected triple overlap,
\[
 n_{ijk}=\frac{\theta_{ij}+\theta_{jk}+\theta_{ki}}{2\pi}
 \in\mathbb Z                                                     \tag{66.4}
\]
is the V54 integral two-cocycle.

\begin{theorem}[Constructive zero-class trivialization]
\label{thm:v18v56-cech-homotopy}
Suppose \([n]=0\) in \(\check H^2(B;\mathbb Z)\).  Choose an
integer one-cochain \(m_{ij}=-m_{ji}\) such that
\[
 n_{ijk}=m_{ij}+m_{jk}+m_{ki},                               \tag{66.5}
\]
and put
\[
 \widetilde\theta_{ij}=\theta_{ij}-2\pi m_{ij}.              \tag{66.6}
\]
Then
\[
 \widetilde\theta_{ij}+\widetilde\theta_{jk}
 +\widetilde\theta_{ki}=0.                                  \tag{66.7}
\]
The functions
\[
 \boxed{
 \phi_i=\sum_{k=1}^M\chi_k\widetilde\theta_{ki}
 \quad\hbox{on }U_i}                                        \tag{66.8}
\]
are smooth and satisfy
\[
 \phi_j-\phi_i=\widetilde\theta_{ij}
 \quad\hbox{on }U_i\cap U_j.                               \tag{66.9}
\]
Consequently
\[
 t|_{U_i}=s_i e^{-i\phi_i}                                  \tag{66.10}
\]
defines a global unit section of \(L\).
\end{theorem}

\begin{proof}
Equations (66.4)--(66.6) give (66.7).  In (66.8), the summand with
index \(k\) is needed only where \(\chi_k\ne0\), a set compactly contained
in \(U_k\).  Extending that product by zero therefore gives a smooth
function on \(U_i\).  On \(U_i\cap U_j\), (66.7) with indices
\((k,i,j)\) gives
\[
 \widetilde\theta_{kj}-\widetilde\theta_{ki}
 =\widetilde\theta_{ij}.                                    \tag{66.11}
\]
Since \(\sum_k\chi_k=1\), summing (66.11) proves (66.9).  Finally,
using (66.3), (66.6), and the fact that \(e^{-2\pi i m_{ij}}=1\),
\[
 s_j e^{-i\phi_j}
 =s_i e^{i\widetilde\theta_{ij}}e^{-i\phi_j}
 =s_i e^{-i\phi_i}.                                        \tag{66.12}
\]
Thus the local formulas glue and have unit norm.
\end{proof}

\begin{remark}[What Smith normal form must return]
\label{rem:v18v56-smith-witness}
The V55 decision gate should retain the integer cochain \(m\), not merely
the Boolean answer \([n]=0\).  Row and column operations in Smith reduction
recover a solution of (66.5); this solution is the discrete witness used in
(66.6)--(66.8).
\end{remark}

\subsection{A first derivative ledger}

Let
\[
\begin{split}
 \Theta&=\max_{i,k}\|\widetilde\theta_{ki}\|_{L^\infty(U_i\cap U_k)},\\
 L_\chi&=\sup_{x\in B}\sum_k |d\chi_k(x)|,\\
 L_\theta&=\max_i\sup_{x\in U_i}
 \sum_{k:\,x\in U_k}\chi_k(x)|d\widetilde\theta_{ki}(x)|.
                                                                  \tag{66.13}
\end{split}
\]
Terms on empty overlaps are omitted.  If the local Berry connection is
\[
 A_i=-i\langle s_i,ds_i\rangle,                              \tag{66.14}
\]
then \(A_j-A_i=d\widetilde\theta_{ij}\).

\begin{theorem}[Quantitative partition-of-unity bound]
\label{thm:v18v56-quantitative-section}
The section (66.10) has a globally defined real connection form
\[
 A^t|_{U_i}=A_i-d\phi_i,                                    \tag{66.15}
\]
and
\[
 \boxed{
 \|d\phi_i\|_{L^\infty(U_i)}
 \le \Theta L_\chi+L_\theta .}                             \tag{66.16}
\]
In particular, if
\(A_0=\max_i\|A_i\|_{L^\infty(U_i)}\), then
\[
 \|A^t\|_{L^\infty(B)}
 \le A_0+\Theta L_\chi+L_\theta.                           \tag{66.17}
\]
\end{theorem}

\begin{proof}
Differentiate (66.8):
\[
 d\phi_i
 =\sum_k(d\chi_k)\widetilde\theta_{ki}
  +\sum_k\chi_k d\widetilde\theta_{ki}.                   \tag{66.18}
\]
The definitions in (66.13) give (66.16).  Under a unit-frame change
\(s\mapsto se^{ig}\), the convention (66.14) gives \(A\mapsto A+dg\).
Hence (66.10) gives (66.15).  Equations (66.9) and
\(A_j-A_i=d\widetilde\theta_{ij}\) show that its two local expressions
agree on every overlap.  The triangle inequality then gives (66.17).
\end{proof}

\begin{corollary}[Equivariant lift from the quotient]
\label{cor:v18v56-equivariant-lift}
Suppose the stabilizer character vanishes and the V53 determinant line
descends to \(L\to B=X/\mathcal G\).  The pullback of (66.10) is a
\(\mathcal G\)-equivariant global unit section on \(X\), with the pullback
of the estimate (66.17).
\end{corollary}

\begin{proof}
A section of an associated line over the orbit quotient is equivalent to
an equivariant section of its pullback.  Pullback commutes with the local
frame changes, the partition-of-unity identity, and the connection formula.
The stabilizer hypothesis makes the fibre action unambiguous.
\end{proof}

\begin{firewall}[The cover stock is analytic data]
\label{fire:v18v56-cover-stock}
The integers \(m_{ij}\) can force large lifted phases in (66.6), and a
refining cover can force \(L_\chi\) to diverge.  Thus a zero Cech class does
not bound \(\Theta L_\chi+L_\theta\).  The quantities in (66.13) must be
entered in the cofinal ledger rather than hidden inside the word
trivialization.
\end{firewall}

\subsection{Compile the anomaly-cancelling phase before estimating}

Let the Standard-Model chiral sectors give Hermitian lines \(L_a\), signs
\(\sigma_a\in\{1,-1\}\), local unit frames \(s_{a,i}\), transition lifts
\(\theta^a_{ij}\), and local connection forms \(A^a_i\).  The compiled
line is
\[
 L_{\rm tot}=\bigotimes_a L_a^{\otimes\sigma_a},             \tag{66.19}
\]
where exponent \(-1\) means the dual line.

\begin{proposition}[Compile-then-estimate principle]
\label{prop:v18v56-compile-first}
In the tensor-product frame, the transition lift and connection are
\[
 \theta^{\rm tot}_{ij}=\sum_a\sigma_a\theta^a_{ij},
 \qquad
 A^{\rm tot}_i=\sum_a\sigma_a A^a_i.                        \tag{66.20}
\]
If the total integral class vanishes, Theorem
\ref{thm:v18v56-cech-homotopy} applies directly to these total quantities.
Its estimate is therefore
\[
 \|A^{t}_{\rm tot}\|_\infty
 \le A^{\rm tot}_0+\Theta_{\rm tot}L_\chi
      +L_{\theta,\rm tot},                                  \tag{66.21}
\]
with stocks formed after the signed sums in (66.20).  Neither sectorwise
triviality nor the larger bound
\(\sum_a(\|A^a\|+\Theta_aL_\chi+L_{\theta,a})\) is required.
\end{proposition}

\begin{proof}
Tensor products multiply transition functions, so their real phases add;
dualization changes the sign.  The Berry connection obeys the same
addition law.  Apply Theorems \ref{thm:v18v56-cech-homotopy} and
\ref{thm:v18v56-quantitative-section} only after forming (66.20).
Cancellation then occurs before any absolute value is taken.  This proves
(66.21).
\end{proof}

\begin{assessment}[Use of the YM V8 lesson]
\label{ass:v18v56-ym-lesson}
Proposition \ref{prop:v18v56-compile-first} is the determinant-line version
of the YM V8 phase-preserving pairing rule.  It prevents a false loss of
the Standard-Model anomaly cancellation through premature sectorwise
triangle estimates.
\end{assessment}

\subsection{Intrinsic Coulomb representative}

The cover estimate is explicit but may be inefficient.  A second
construction replaces cover geometry by curvature, holonomy, and the
Hodge constants of the quotient.  Assume in this subsection that \(B\) is
a connected closed Riemannian manifold and that \(L\) is topologically
trivial.  In a global unit frame write the connection as a real one-form
\(A\), with curvature \(F=dA\).  Let
\[
 \Lambda_B=2\pi\,\operatorname{im}
 \bigl(H^1(B;\mathbb Z)\longrightarrow\mathcal H^1(B)\bigr) \tag{66.22}
\]
be the lattice of harmonic forms with integral periods, and let \(h(A)\)
be the harmonic term in the Hodge decomposition of \(A\).  Define the flat
residual
\[
 \eta_B(A)=\inf_{\lambda\in\Lambda_B}
 \|h(A)-\lambda\|_{H^1(B)}.                                 \tag{66.23}
\]

\begin{theorem}[Curvature--holonomy gauge bound]
\label{thm:v18v56-hodge-bound}
There is a constant \(C_B<\infty\), depending only on the Riemannian Hodge
problem on \(B\), and a unitary gauge \(u:B\to U(1)\) such that
\[
 \delta A^u=0,
 \qquad
 \boxed{
 \|A^u\|_{H^1(B)}
 \le C_B\|F\|_{L^2(B)}+\eta_B(A).}                          \tag{66.24}
\]
If \(F=0\) and every flat holonomy is trivial, one may choose \(A^u=0\).
\end{theorem}

\begin{proof}
Hodge decomposition gives
\[
 A=df+\delta\beta+h,                                        \tag{66.25}
\]
with \(h\in\mathcal H^1(B)\).  The gauge \(e^{-if}\) removes \(df\).
Every map \(B\to U(1)\) contributes a closed one-form whose periods lie in
\(2\pi\mathbb Z\), and every class in \(2\pi H^1(B;\mathbb Z)\) arises in
this way.  A large gauge can therefore replace \(h\) by
\(h-\lambda\), where \(\lambda\in\Lambda_B\) is chosen in a nearest
fundamental cell.  The remaining coexact term
\(a=\delta\beta\) satisfies
\[
 \delta a=0,\qquad a\perp\mathcal H^1(B),\qquad da=F.        \tag{66.26}
\]
The elliptic estimate for the Hodge Laplacian on coexact one-forms gives
\[
 \|a\|_{H^1(B)}\le C_B\|da\|_{L^2(B)}
 =C_B\|F\|_{L^2(B)}.                                       \tag{66.27}
\]
The triangle bound proves (66.24).  If \(F=0\), the coexact term vanishes.
Trivial flat holonomy says that the harmonic class lies in \(\Lambda_B\),
so a large gauge also removes it.
\end{proof}

\begin{remark}[Finite flat panel]
\label{rem:v18v56-flat-panel}
The number \(\eta_B(A)\) is determined by the holonomies on a finite
integral basis for the free part of \(H_1(B;\mathbb Z)\).  The torsion line
class has already been removed by the V54--V55 Cech gate; it cannot be
detected by the curvature term in (66.24).
\end{remark}

\begin{countertheorem}[Topological zero does not imply analytic control]
\label{cth:v18v56-topology-not-bound}
There is a sequence of connections on a fixed topologically trivial line
whose connection forms have no uniformly bounded \(W^{1,p}\) gauge,
for any fixed \(1\le p\le\infty\).
\end{countertheorem}

\begin{proof}
On the unit disk \(D\subset\mathbb R^2\), take the trivial line and
\[
 A_N=N x\,dy,
 \qquad F_N=dA_N=N\,dx\wedge dy.                             \tag{66.28}
\]
Every unitary gauge changes \(A_N\) by a closed form and hence preserves
\(F_N\).  The exterior derivative is a bounded map
\(W^{1,p}(D;T^*D)\to L^p(D;\Lambda^2T^*D)\), so
\[
 N\|dx\wedge dy\|_{L^p(D)}
 =\|dA_N^u\|_{L^p(D)}
 \le C_p\|A_N^u\|_{W^{1,p}(D)}.                             \tag{66.29}
\]
The right side cannot be bounded uniformly in \(N\).  Yet
\(H^2(D;\mathbb Z)=0\), so every one of these lines has zero Cech class.
\end{proof}

\begin{firewall}[Sobolev control is not locality]
\label{fire:v18v56-sobolev-not-locality}
The estimate (66.24) controls a global Sobolev norm on a smooth compact
quotient.  It does not by itself prove exponential lattice locality,
reflection positivity, a bound uniform across singular orbit strata, or
the pointwise regularity needed by the V45 Ward collar.  Those require
separate kernel and quotient-geometry estimates.
\end{firewall}

\subsection{Cofinal quantitative capsule}

For cutoff \(n\), write \(B_n\) for a regular quotient stratum,
\(L_n\to B_n\) for the compiled line, \(F_n\) for its curvature, and
\(\eta_n\) for (66.23).  Let \(C_n\) be the Hodge constant in (66.24).
Suppose the cofinal maps carry unitary line identifications \(J_n\), and
let \(\tau_n\) measure the chosen section-transport defect in a complete
comparison norm.

\begin{theorem}[Quantitative chiral-section capsule]
\label{thm:v18v56-cofinal-capsule}
Assume along a cofinal sequence that:
\begin{enumerate}[leftmargin=2.2em]
\item the stabilizer characters vanish and the compiled integral Cech
      classes are zero;
\item the regular strata admit the Hodge problem above, with
\[
 \varepsilon_n
 :=C_n\|F_n\|_{L^2(B_n)}+\eta_n\longrightarrow0;            \tag{66.30}
\]
\item gauges can be chosen so that their transported unit sections obey
\[
 \|J_n t_n-t_{n+1}\|\le\tau_n,
 \qquad \sum_n\tau_n<\infty;                                \tag{66.31}
\]
\item the same gauges have summable locality, reflection, and Ward-collar
      defects in the norms required by V45--V51.
\end{enumerate}
Then the \(t_n\) form a cofinally consistent equivariant unit section, and
their compiled Berry connections converge to zero in the transported
Hodge norm.  Consequently the measure-line contribution to the limiting
Ward residual vanishes in every comparison norm dominated by the ledger in
item four.
\end{theorem}

\begin{proof}
Item one and Theorem \ref{thm:v18v56-cech-homotopy} give an equivariant
unit section at each level.  Theorem \ref{thm:v18v56-hodge-bound} and
(66.30) choose gauges with \(\|A_n^{t_n}\|_{H^1}\le\varepsilon_n\), hence
the connection term tends to zero.  By telescoping (66.31), for \(m>n\),
\[
 \|J_{m-1}\cdots J_n t_n-t_m\|
 \le\sum_{k=n}^{m-1}\tau_k,                                 \tag{66.32}
\]
so completeness gives a compatible cofinal limit.  Item four transports
the equivariance, locality, and reflection properties and makes the
measure-line term in the V45 Ward collar tend to zero.  The V46--V51
handoff then accepts this term as a vanishing component of its explicitly
listed residual; it does not remove any other residual.
\end{proof}

\begin{firewall}[The literal Standard-Model rows remain empty]
\label{fire:v18v56-unpopulated}
No calculation in V56 establishes zero integral Cech class for the literal
cofinal overlap Standard Model, a uniform Hodge constant \(C_n\), decay of
\(F_n\) and \(\eta_n\), or the summability in (66.31).  Orbit-space
singularities also prevent automatic application of the smooth Hodge
theorem.  Therefore Theorem \ref{thm:v18v56-cofinal-capsule} is a proved
implication with unpopulated physical hypotheses, not a construction of
the full SM--Einstein continuum.
\end{firewall}

\subsection{Acquisition and next bottleneck}

V56 converts the V55 topological decision into an explicit global section
and exposes the exact analytic stocks which the decision does not control.
The next upstream bottleneck is the compiled overlap curvature itself:
one must express \(F_n=-i\operatorname{Tr}(P_n\,dP_n\wedge dP_n)\) after
the full Standard-Model signed sum and prove a regulator-uniform estimate
for its coexact part and flat holonomy residual.  The rational direction is
therefore to derive a resolvent-level curvature identity in which the
anomaly-cancelling representation sum is performed before operator norms.

\begin{assessment}[V56 acquisition and boundary]
\label{ass:v18v56}
V56 proves an explicit Cech-homotopy trivialization, a quantitative
first-derivative bound, an equivariant quotient lift, the
compile-then-estimate rule, and an intrinsic curvature--holonomy gauge
bound.  It also proves by a fixed-base counterexample that topological
triviality alone cannot give uniform analytic control.  The remaining
physical gap is now a concrete compiled-curvature and cofinal-geometry
estimate.
\end{assessment}

\begin{record}[V56 append record]
\label{rec:v18v56}
V56 is appended to the validated V55 whose installed SHA--256 was
\texttt{33E1B2256D948CE6725C2D0255D44F973E87F09EE1C6A8FDBA80C0B953FAF3F0}.
After installation only V56 is the active numeric SM note; the frozen
\texttt{V100\_f2} through \texttt{V100\_f17} archives are unchanged.
\end{record}

\section{V57: vertical Ward curvature and the compiled resolvent target}
\label{sec:v18v57}

V56 gave a strong sufficient condition which makes the entire compiled
Berry connection small on a smooth quotient.  That condition is more than
the gauge Ward identity requires.  A nonzero horizontal Berry curvature
can survive on the physical orbit space without producing a gauge anomaly.
This note isolates the vertical curvature, proves its exact finite-matrix
moment-map identity, and identifies the resolvent expression on which the
Standard-Model representation sum must be compiled.

\subsection{Covariant finite-cutoff family}

Let \(X_n\) be a smooth admissible configuration stratum, let
\(\mathcal G_n\) act on it, and let \(H_{a,n}(U)\) be the Hermitian gapped
kernel for chiral sector \(a\).  Write
\[
 P_{a,n}=1_{(-\infty,0)}(H_{a,n}),\qquad
 F_{a,n}=-i\operatorname{Tr}
 \bigl(P_{a,n}\,dP_{a,n}\wedge dP_{a,n}\bigr).               \tag{67.1}
\]
The signed compiled quantities are
\[
 F_n^{\rm tot}=\sum_a\sigma_aF_{a,n},\qquad
 \sigma_a\in\{1,-1\}.                                      \tag{67.2}
\]
For a fixed infinitesimal gauge parameter \(\omega\), suppose covariance
is represented on the one-particle space by an anti-Hermitian matrix
\(\Omega_{a,n}(\omega)\), independent of \(U\), so that
\[
 d_{\omega^\#}H_{a,n}=[\Omega_{a,n}(\omega),H_{a,n}],
 \qquad
 d_{\omega^\#}P_{a,n}=[\Omega_{a,n}(\omega),P_{a,n}].        \tag{67.3}
\]

\begin{theorem}[Exact vertical moment-map identity]
\label{thm:v18v57-moment-map}
Define
\[
 J_{a,n}(\omega)
 =-i\operatorname{Tr}\bigl(P_{a,n}\Omega_{a,n}(\omega)\bigr),
 \qquad
 J_n^{\rm tot}(\omega)=\sum_a\sigma_aJ_{a,n}(\omega).        \tag{67.4}
\]
Then
\[
 \boxed{
 \iota_{\omega^\#}F_{a,n}=-dJ_{a,n}(\omega),\qquad
 \iota_{\omega^\#}F_n^{\rm tot}=-dJ_n^{\rm tot}(\omega).}    \tag{67.5}
\]
\end{theorem}

\begin{proof}
Fix a tangent vector \(X\) and abbreviate \(P=P_{a,n}\),
\(\Omega=\Omega_{a,n}(\omega)\), and \(P_X=d_XP\).  From (67.1) and
(67.3),
\[
 F_{a,n}(\omega^\#,X)
 =-i\operatorname{Tr}
 \bigl(P[[\Omega,P],P_X]\bigr).                              \tag{67.6}
\]
Relative to \(\operatorname{Ran}P\oplus\ker P\), both
\([\Omega,P]\) and \(P_X\) are off diagonal because
\(P P_XP=(I-P)P_X(I-P)=0\).  Direct block multiplication, followed by
cyclicity of trace, gives
\[
 \operatorname{Tr}\bigl(P[[\Omega,P],P_X]\bigr)
 =-\operatorname{Tr}(P_X\Omega).                            \tag{67.7}
\]
Since \(\Omega\) is independent of \(U\),
\[
 d_XJ_{a,n}(\omega)=-i\operatorname{Tr}(P_X\Omega).          \tag{67.8}
\]
Equations (67.6)--(67.8) give the first identity in (67.5); signed
summation gives the second.
\end{proof}

\begin{remark}[Configuration-dependent generators]
\label{rem:v18v57-dependent-generator}
If a chosen gauge slice makes \(\Omega(\omega)\) depend on \(U\), then
\(dJ\) contains the additional term
\(-i\operatorname{Tr}(P\,d\Omega)\).  Equation (67.5) is valid in the
fixed one-particle representation used in (67.3).  A moving-frame
description must retain that extra term.
\end{remark}

\subsection{What gauge-anomaly cancellation does and does not say}

\begin{proposition}[Basicness criterion for the compiled curvature]
\label{prop:v18v57-basicness}
If \(J_n^{\rm tot}(\omega)\) is constant on each connected component for
every \(\omega\), then
\[
 \iota_{\omega^\#}F_n^{\rm tot}=0.                           \tag{67.9}
\]
If the connection is gauge invariant, (67.9) makes
\(F_n^{\rm tot}\) basic and hence the pullback of a two-form on the
free-orbit quotient.  If the constants in \(J_n^{\rm tot}\) are normalized
to zero and all gauge-loop holonomies vanish, the connection itself
descends.
\end{proposition}

\begin{proof}
The first assertion is (67.5).  Gauge invariance gives
\(\mathcal L_{\omega^\#}F_n^{\rm tot}=0\); horizontality and invariance are
the definition of a basic form.  For the connection, zero infinitesimal
vertical component removes connected gauge-orbit phases, while the
gauge-loop condition removes the remaining global monodromy.  These are
exactly the local and global descent conditions of V53--V54.
\end{proof}

\begin{countertheorem}[Anomaly cancellation does not flatten the quotient]
\label{cth:v18v57-horizontal-survival}
Vanishing of every vertical contraction
\(\iota_{\omega^\#}F\) does not imply \(F=0\), even when the underlying
line is topologically trivial.
\end{countertheorem}

\begin{proof}
Let \(B\) be the unit disk with the trivial line and connection
\(A=x\,dy\), so \(F=dx\wedge dy\ne0\).  Put
\(X=B\times\mathcal G\) and let \(\mathcal G\) act on the second factor.
Pull back \(A\) from \(B\).  Its curvature is horizontal, so every vertical
contraction vanishes, but its restriction to the quotient \(B\) remains
\(dx\wedge dy\).  The line is topologically trivial because \(B\) is
contractible.
\end{proof}

\begin{assessment}[Correction to the V56 strong capsule]
\label{ass:v18v57-correction}
The decay condition \(C_n\|F_n\|+\eta_n\to0\) in V56 is a valid sufficient
condition for a small full connection.  It is not necessary for the gauge
Ward identity.  The sharper Ward target is control of
\(J_n^{\rm tot}\), \(dJ_n^{\rm tot}\), and gauge-loop holonomy; horizontal
quotient curvature may remain.
\end{assessment}

\subsection{Exact resolvent formulas}

Fix a positively oriented contour \(\Gamma\) enclosing the negative
spectrum of every sector kernel and no positive spectrum.  Put
\[
 R_{a,n}(z)=(z-H_{a,n})^{-1}.                                \tag{67.10}
\]
Then
\[
 P_{a,n}=\frac1{2\pi i}\oint_\Gamma R_{a,n}(z)\,dz,\qquad
 d_XP_{a,n}=\frac1{2\pi i}\oint_\Gamma
 R_{a,n}(z)(d_XH_{a,n})R_{a,n}(z)\,dz.                      \tag{67.11}
\]

\begin{theorem}[Compiled vertical resolvent identity]
\label{thm:v18v57-compiled-resolvent}
For every \(X\) and \(\omega\),
\[
\boxed{
 d_XJ_n^{\rm tot}(\omega)
 =
 -\frac1{2\pi}
 \oint_\Gamma
 \sum_a\sigma_a
 \operatorname{Tr}\!\left(
 R_{a,n}(z)(d_XH_{a,n})R_{a,n}(z)
 \Omega_{a,n}(\omega)\right)dz .}                           \tag{67.12}
\]
The signed sector sum in (67.12) is inside the contour integral.  It is
therefore legitimate to exploit exact algebraic cancellation in that sum
before taking an absolute value or an operator norm.
\end{theorem}

\begin{proof}
Differentiate (67.4), insert the second formula in (67.11), and use
\(-i/(2\pi i)=-1/(2\pi)\).  Finite dimensionality permits signed summation
and contour integration in either order.  No estimate has yet been made,
so cancellations in the summed trace are exact.
\end{proof}

\begin{corollary}[Honest noncancelling fallback bound]
\label{cor:v18v57-fallback}
If \(\ell(\Gamma)\le L\),
\(\operatorname{dist}(\Gamma,\operatorname{spec}H_{a,n})\ge r>0\), and
\[
 Q_X=\max_a\|d_XH_{a,n}\|,\qquad
 W_\omega=\max_a\|\Omega_{a,n}(\omega)\|,\qquad
 N_\Sigma=\sum_a\dim\mathcal H_{a,n},                        \tag{67.13}
\]
then
\[
 |d_XJ_n^{\rm tot}(\omega)|
 \le \frac{L}{2\pi}\,N_\Sigma r^{-2}Q_XW_\omega.             \tag{67.14}
\]
\end{corollary}

\begin{proof}
On \(\Gamma\), \(\|R_{a,n}\|\le r^{-1}\).  Apply
\(|\operatorname{Tr}M|\le(\dim\mathcal H_{a,n})\|M\|\) to each sector only
after the exact identity (67.12).  This deliberately crude estimate is
the fallback when no compiled cancellation certificate is available.
\end{proof}

\begin{firewall}[Sectorwise bounds erase the Standard-Model arithmetic]
\label{fire:v18v57-sectorwise-loss}
Applying (67.14) before forming the signed sum replaces every representation
coefficient by its absolute size.  The cubic, mixed, and gravitational
hypercharge cancellations verified in V44 then become invisible.  Such a
bound can establish finite-cutoff smoothness, but cannot certify anomaly
cancellation.
\end{firewall}

\subsection{A finite weak-field cancellation certificate}

The resolvent identity also yields an exact remainder theorem.  Suppress
\(n\) and suppose on the contour that
\[
 H_a=H_0+V_a,\qquad R_0(z)=(z-H_0)^{-1},\qquad
 q:=\max_{a,z\in\Gamma}\|R_0(z)V_a\|<1.                     \tag{67.15}
\]
For an integer \(M\ge1\), define
\[
 S_{a,M}(z)=\sum_{k=0}^{M-1}(R_0(z)V_a)^kR_0(z).             \tag{67.16}
\]

\begin{lemma}[Neumann remainder]
\label{lem:v18v57-neumann}
If \(\|R_0(z)\|\le r^{-1}\), then
\[
 R_a=S_{a,M}+E_{a,M},\qquad
 \|E_{a,M}\|\le r^{-1}\frac{q^M}{1-q},\qquad
 \|R_a\|,\|S_{a,M}\|\le r^{-1}(1-q)^{-1}.                   \tag{67.17}
\]
\end{lemma}

\begin{proof}
The identity
\[
 (I-R_0V_a)^{-1}
 =\sum_{k=0}^{M-1}(R_0V_a)^k
 +(R_0V_a)^M(I-R_0V_a)^{-1}                                \tag{67.18}
\]
multiplied on the right by \(R_0\) gives the decomposition.  The geometric
series bound gives all three estimates.
\end{proof}

\begin{theorem}[Truncated compiled certificate]
\label{thm:v18v57-truncated-certificate}
Assume (67.13), (67.15), and that the finite algebraic identity
\[
 \sum_a\sigma_a\operatorname{Tr}\!\left(
 S_{a,M}(z)(d_XH_a)S_{a,M}(z)\Omega_a(\omega)\right)=0       \tag{67.19}
\]
holds for every \(z\in\Gamma\) and for the prescribed \(X,\omega\).
Then
\[
\boxed{
 |d_XJ^{\rm tot}(\omega)|
 \le
 \frac{L}{\pi}\,
 N_\Sigma Q_XW_\omega r^{-2}
 \frac{q^M}{(1-q)^2}.}                                      \tag{67.20}
\]
\end{theorem}

\begin{proof}
Write \(R_a=S_a+E_a\).  The exact telescoping identity
\[
 R_a(d_XH_a)R_a-S_a(d_XH_a)S_a
 =
 E_a(d_XH_a)R_a+S_a(d_XH_a)E_a                            \tag{67.21}
\]
and (67.17) give
\[
 \|R_a(d_XH_a)R_a-S_a(d_XH_a)S_a\|
 \le 2Q_Xr^{-2}\frac{q^M}{(1-q)^2}.                         \tag{67.22}
\]
Insert (67.19) into (67.12), use (67.22), and sum the trace-dimension
bounds.  The contour length contributes \(L/(2\pi)\), proving (67.20).
\end{proof}

\begin{remark}[Finite exact nature of the certificate]
\label{rem:v18v57-finite-certificate}
For fixed \(M\), expansion of (67.19) contains finitely many traces of
words in \(R_0,V_a,d_XH_a,\Omega_a\).  Whenever the background part is
representation blind and the remaining factors split into lattice kernels
and representation matrices, the signed representation traces can be
checked exactly.  V44 supplies the familiar low-degree Standard-Model
trace cancellations.  Equation (67.19), however, is stronger and must be
verified for every word actually produced at the chosen order.
\end{remark}

\begin{firewall}[Perturbative anomaly arithmetic is not (67.19)]
\label{fire:v18v57-not-automatic}
The continuum anomaly polynomial cancels a specified finite family of
local representation traces.  It does not imply that the full
finite-cutoff resolvent integrand, or every coefficient in an arbitrary
Neumann truncation, vanishes.  The smallness condition \(q<1\) is also only
a weak-field chart.  Thus (67.20) is a rigorous conditional estimate, not
a completed Standard-Model certificate.
\end{firewall}

\subsection{The sharpened cofinal Ward target}

\begin{theorem}[Vertical measure-line capsule]
\label{thm:v18v57-vertical-capsule}
Suppose along a cofinal cutoff sequence:
\begin{enumerate}[leftmargin=2.2em]
\item the compiled line passes the V53--V55 stabilizer, integral Cech, and
      gauge-loop tests;
\item for a norming bank \(\mathfrak B_n\) of gauge parameters and tangent
      vectors,
\[
 \alpha_n
 =\sup_{\omega,X\in\mathfrak B_n}
 |d_XJ_n^{\rm tot}(\omega)|\longrightarrow0;                \tag{67.23}
\]
\item the constants \(J_n^{\rm tot}(\omega)\) are normalized compatibly
      and their transport defects are summable;
\item the locality, reflection, and Ward comparison defects appearing in
      V45--V51 are summable.
\end{enumerate}
Then the vertical curvature contribution of the chiral measure line
vanishes in the continuum Ward residual.  No decay of the horizontal
quotient curvature is required.
\end{theorem}

\begin{proof}
By (67.5), (67.23) is precisely uniform decay of the curvature with one
gauge-vertical leg on the norming bank.  Density and the assumed comparison
control extend it to the completed Ward test space.  The first and third
items eliminate stabilizer phases and global gauge-orbit monodromy, so the
infinitesimal estimate is not defeated by a flat global obstruction.
The fourth item passes this vanishing term through the V45 Ward collar and
the V46--V51 continuum handoff.  Horizontal components have no
gauge-vertical leg and therefore do not enter this conclusion.
\end{proof}

\begin{firewall}[Unpopulated resolvent row]
\label{fire:v18v57-unpopulated}
No uniform contour, weak-field radius \(q<1\), finite identity (67.19), or
nonperturbative substitute has been established for the literal cofinal
Standard-Model overlap kernels.  Nor has a norming bank satisfying
(67.23) been constructed.  These are the current measure-sector rows.
\end{firewall}

\subsection{Acquisition and next bottleneck}

V57 replaces the unnecessarily strong goal of flattening the entire
quotient connection by the exact vertical target \(dJ_n^{\rm tot}\).
It also supplies a finite weak-field certificate with an explicit
geometric-series remainder.  The next upstream task is to construct a
chart-independent cancellation scheme: cover the admissible configuration
space by resolvent neighborhoods, transport the compiled moment map across
overlaps, and determine whether gap and locality stocks permit a cofinal
chain of such neighborhoods.

\begin{assessment}[V57 acquisition and boundary]
\label{ass:v18v57}
V57 proves the vertical moment-map identity, distinguishes gauge basicness
from horizontal flatness, places the signed Standard-Model sum inside an
exact resolvent integral, and proves the quantitative truncated-certificate
bound (67.20).  The unresolved step is a literal uniform cancellation
certificate beyond one weak-field chart.
\end{assessment}

\begin{record}[V57 append record]
\label{rec:v18v57}
V57 is appended to the validated V56 whose installed SHA--256 was
\texttt{C4C284AB2C89D9241A3239040D76FEBC6EAC3755EC4E10E819BAC880A08DECFF}.
After installation only V57 is the active numeric SM note; the frozen
\texttt{V100\_f2} through \texttt{V100\_f17} archives are unchanged.
\end{record}

\section{V58: compact resolvent atlases and finite-center Ward certification}
\label{sec:v18v58}

V57 reduced the chiral gauge problem to the compiled scalar one-form
\(dJ_n^{\rm tot}\).  A weak-field Neumann chart is insufficient on a large
admissible component, but the spectral gap produces a finite resolvent
atlas at every fixed cutoff.  This note proves that construction and shows
how a finite bank of center certificates controls the whole component.
The analytic error is a maximum over charts; it does not acquire a factor
equal to the number of charts.

\subsection{Gap-preserving neighborhoods}

Fix one cutoff and suppress \(n\).  Let \(X\) be a compact connected
Riemannian configuration stratum.  For every chiral sector \(a\), let
\(H_a:X\to\operatorname{Herm}(\mathcal H_a)\) be \(C^2\), and assume
\[
 g=\inf_{U\in X}\min_a
 \operatorname{dist}\bigl(0,\operatorname{spec}H_a(U)\bigr)>0.          \tag{68.1}
\]
Let
\[
 L_H=\max_a\sup_{U\in X}\|dH_a(U)\|_{\rm op}.                \tag{68.2}
\]

\begin{lemma}[Uniform gap neighborhood]
\label{lem:v18v58-gap-neighborhood}
If \(d(U,V)<\rho\) and \(L_H\rho\le g/4\), then
\[
 \|H_a(U)-H_a(V)\|\le g/4,\qquad
 \operatorname{dist}\bigl(0,\operatorname{spec}H_a(U)\bigr)
 \ge 3g/4                                                     \tag{68.3}
\]
whenever the center \(V\) has gap at least \(g\).
\end{lemma}

\begin{proof}
Integrate \(dH_a\) along a minimizing path from \(V\) to \(U\) to obtain
the first bound.  Weyl's eigenvalue perturbation inequality moves every
eigenvalue by at most \(\|H_a(U)-H_a(V)\|\), which gives the second.
\end{proof}

\begin{proposition}[Finite resolvent atlas]
\label{prop:v18v58-finite-atlas}
Put \(\rho=g/(4L_H)\) when \(L_H>0\), and take any positive \(\rho\) when
\(L_H=0\).  There are finitely many centers
\(\{U_\alpha\}_{\alpha=1}^{N}\) such that
\[
 X=\bigcup_{\alpha=1}^{N}B(U_\alpha,\rho).                  \tag{68.4}
\]
On every ball, the negative spectral rank is constant and the Riesz
projector admits a center contour, common to all points of that ball, and
a convergent resolvent perturbation expansion.
\end{proposition}

\begin{proof}
Compactness gives a finite \(\rho\)-subcover.  Lemma
\ref{lem:v18v58-gap-neighborhood} prevents an eigenvalue from crossing
zero, hence fixes the negative rank.  Around each finite-dimensional
center spectrum choose a contour at distance at least \(g/2\) from the
spectrum which separates its negative and positive parts; its outer sides
may be placed beyond the finite spectral radius.  On the ball,
\[
 \|R_{a,\alpha}(z)(H_a(U)-H_a(U_\alpha))\|
 \le (2/g)(g/4)=1/2.                                      \tag{68.4a}
\]
The V57 Neumann identity therefore converges uniformly.
\end{proof}

\begin{firewall}[Finite at one cutoff is not uniform cofinally]
\label{fire:v18v58-atlas-complexity}
The covering number \(N\), radius \(\rho\), contour lengths, and resolvent
distances depend on the cutoff.  Compactness proves finite termination for
each \(n\), but supplies no cofinal bound on any of these stocks.
\end{firewall}

\subsection{A resolvent Lipschitz bound for the compiled Ward one-form}

Work in one atlas ball and use one contour \(\Gamma\) for two points
\(U,V\).  Suppose
\[
\begin{split}
 &\ell(\Gamma)\le L_\Gamma,\qquad
 \|R_a(U,z)\|,\|R_a(V,z)\|\le r^{-1},\\
 &\|H_a(U)-H_a(V)\|\le L_Hd(U,V),\\
 &\|d_XH_a(U)\|\le Q,\qquad
 \|d_XH_a(U)-d_XH_a(V)\|\le S\,d(U,V),\\
 &\|\Omega_a(\omega)\|\le W.                                \tag{68.5}
\end{split}
\]
Here \(X\) is a unit tangent field transported inside the coordinate ball,
and
\[
 N_\Sigma=\sum_a\dim\mathcal H_a.                            \tag{68.6}
\]

\begin{lemma}[Resolvent displacement]
\label{lem:v18v58-resolvent-displacement}
Under (68.5),
\[
 \|R_a(U,z)-R_a(V,z)\|
 \le r^{-2}L_Hd(U,V).                                       \tag{68.7}
\]
\end{lemma}

\begin{proof}
The second resolvent identity gives
\[
 R_a(U,z)-R_a(V,z)
 =R_a(U,z)\bigl(H_a(U)-H_a(V)\bigr)R_a(V,z).                \tag{68.8}
\]
Take norms and use (68.5).
\end{proof}

\begin{theorem}[Local Lipschitz stock for \(dJ^{\rm tot}\)]
\label{thm:v18v58-local-lipschitz}
Under (68.5),
\[
\begin{split}
 &\left|
 d_XJ^{\rm tot}(U;\omega)-d_XJ^{\rm tot}(V;\omega)
 \right|\\
 &\qquad\le
 K_J\,d(U,V),\\
 &K_J=
 \frac{L_\Gamma}{2\pi}N_\Sigma W
 \left(2r^{-3}L_HQ+r^{-2}S\right).                          \tag{68.9}
\end{split}
\]
\end{theorem}

\begin{proof}
Write \(R_U=R_a(U,z)\), \(R_V=R_a(V,z)\),
\(Q_U=d_XH_a(U)\), and \(Q_V=d_XH_a(V)\).  The exact three-term
telescoping identity is
\[
\begin{split}
 R_UQ_UR_U-R_VQ_VR_V
 ={}&(R_U-R_V)Q_UR_U\\
 &+R_V(Q_U-Q_V)R_U
 +R_VQ_V(R_U-R_V).                                         \tag{68.10}
\end{split}
\]
Lemmas \ref{lem:v18v58-resolvent-displacement} and (68.5) bound its norm by
\[
 \left(2r^{-3}L_HQ+r^{-2}S\right)d(U,V).                    \tag{68.11}
\]
Insert (68.10) into the compiled V57 resolvent identity.  Only after the
signed difference has been formed, use
\(|\operatorname{Tr}M|\le(\dim\mathcal H_a)\|M\|\), sum over sectors, and
integrate the contour length.  This gives (68.9).
\end{proof}

\begin{remark}[Where cancellation is retained]
\label{rem:v18v58-center-versus-remainder}
The value at the center \(dJ^{\rm tot}(U_\alpha)\) is compiled before it is
estimated and can use Standard-Model trace cancellation.  The displacement
from the center uses the noncancelling fallback stock \(K_J\).  Thus only a
small geometric remainder pays the sectorwise triangle bound.
\end{remark}

\subsection{Finite-center global theorem}

Let \(\{B(U_\alpha,\delta_\alpha)\}\) cover \(X\), with the local constants
\(K_{J,\alpha}\) from (68.9).  Suppose an exact or interval certificate at
each center gives
\[
 \sup_{\omega,X\in\mathfrak B_\alpha}
 |d_XJ^{\rm tot}(U_\alpha;\omega)|\le\kappa_\alpha,          \tag{68.12}
\]
for a norming bank transported throughout its ball.

\begin{theorem}[Global finite-center certificate]
\label{thm:v18v58-global-certificate}
Under the stated cover and transport assumptions,
\[
\boxed{
 \sup_{U\in X}\sup_{\omega,X\in\mathfrak B_U}
 |d_XJ^{\rm tot}(U;\omega)|
 \le
 \max_\alpha\bigl(\kappa_\alpha+
 K_{J,\alpha}\delta_\alpha\bigr).}                          \tag{68.13}
\]
In particular, the right side has no multiplicative covering-number
factor.
\end{theorem}

\begin{proof}
For a given \(U\), choose one center \(\alpha\) with
\(d(U,U_\alpha)<\delta_\alpha\).  The triangle inequality is used once:
\[
 |dJ^{\rm tot}(U)|
 \le |dJ^{\rm tot}(U_\alpha)|
 +|dJ^{\rm tot}(U)-dJ^{\rm tot}(U_\alpha)|
 \le\kappa_\alpha+K_{J,\alpha}\delta_\alpha.                \tag{68.14}
\]
Taking the supremum over \(U\) gives the maximum in (68.13).  There is no
chain of charts and hence no sum over \(\alpha\).
\end{proof}

\begin{corollary}[Overlap consistency panel]
\label{cor:v18v58-overlap-panel}
Let \(\widetilde K_\alpha\) be a certified approximation to
\(dJ^{\rm tot}\) on chart \(\alpha\), with uniform error
\(\varepsilon_\alpha\).  On an overlap,
\[
 \|\widetilde K_\alpha-\widetilde K_\beta\|
 \le\varepsilon_\alpha+\varepsilon_\beta.                   \tag{68.15}
\]
Thus every edge of the finite atlas nerve supplies a finite consistency
test.
\end{corollary}

\begin{proof}
Both chart approximants are compared with the same intrinsic one-form
\(dJ^{\rm tot}\).  Apply the triangle inequality.
\end{proof}

\begin{assessment}[No partition loss for the vertical observable]
\label{ass:v18v58-no-partition-loss}
Unlike the determinant-line section in V56, the moment-map derivative is
already a global intrinsic one-form.  Local certified bounds are combined
by selection of one covering chart, not by a partition of unity.  Cover
derivatives and the number of charts affect certificate construction, but
not the analytic supremum in (68.13).
\end{assessment}

\subsection{Cofinal order schedule for weak-field remainders}

Let the V57 truncated certificate be used at every center of cutoff \(n\).
Absorb its contour length, total matrix dimension, insertion norms, and
generator norms into
\[
 K_n=\max_\alpha
 \frac{L_{\Gamma,n,\alpha}}{\pi}
 N_{\Sigma,n}Q_{n,\alpha}W_{n,\alpha}.                      \tag{68.16}
\]
Let \(r_n\) be the smallest center-contour distance, let
\(q_n<1\) be the largest Neumann ratio, and let \(M_n\) be the smallest
truncation order.  The center remainder obeys
\[
 \kappa_n
 \le K_nr_n^{-2}\frac{q_n^{M_n}}{(1-q_n)^2}.                \tag{68.17}
\]

\begin{proposition}[Necessary and sufficient growth test for this bound]
\label{prop:v18v58-order-schedule}
Assume \(K_n,r_n,q_n>0\).  The right side of (68.17) tends to zero if and
only if
\[
 M_n|\log q_n|
 -2\log r_n^{-1}
 -\log K_n
 -2\log(1-q_n)^{-1}
 \longrightarrow+\infty.                                   \tag{68.18}
\]
Given any \(s_n\to+\infty\), it is sufficient to choose
\[
 M_n\ge
 \frac{
 2\log r_n^{-1}+\log K_n+2\log(1-q_n)^{-1}+s_n
 }{|\log q_n|}.                                             \tag{68.19}
\]
\end{proposition}

\begin{proof}
Take the logarithm of the positive right side of (68.17).  Its convergence
to \(-\infty\) is exactly (68.18).  Substitution of (68.19) bounds that
logarithm by \(-s_n\).
\end{proof}

\begin{firewall}[Order growth does not create algebraic cancellation]
\label{fire:v18v58-order-not-cancellation}
Increasing \(M_n\) reduces the Neumann remainder only after the entire
truncated compiled identity has been certified through that order.  The
number of representation words generally grows with \(M_n\), and the
center resolvent is representation dependent away from a symmetric
background.  Formula (68.19) is an analytic schedule, not a proof that the
required trace identities exist.
\end{firewall}

\subsection{Cofinal atlas capsule}

\begin{theorem}[Finite-center vertical Ward capsule]
\label{thm:v18v58-cofinal-atlas}
For each cutoff \(n\), suppose:
\begin{enumerate}[leftmargin=2.2em]
\item the relevant admissible component is compact and has a finite
      gap-preserving resolvent atlas;
\item every center has a compiled certificate (68.12);
\item the tangent and gauge-parameter banks transport across each chart;
\item
\[
 \epsilon_n=
 \max_\alpha
 \bigl(\kappa_{n,\alpha}
 +K_{J,n,\alpha}\delta_{n,\alpha}\bigr)\longrightarrow0;    \tag{68.20}
\]
\item the V57 topological, gauge-loop, locality, reflection, and cutoff
      transport rows hold.
\end{enumerate}
Then the chiral measure-line vertical curvature contribution vanishes in
the continuum Ward residual.
\end{theorem}

\begin{proof}
Theorem \ref{thm:v18v58-global-certificate} gives
\(\|dJ_n^{\rm tot}\|\le\epsilon_n\) on the entire component.  Equation
(67.5) converts this to uniform decay of curvature with one gauge-vertical
leg.  Item five supplies the global phase and passage through the V45--V51
continuum comparison maps.  Apply the V57 vertical measure-line capsule.
\end{proof}

\begin{firewall}[Literal status]
\label{fire:v18v58-literal-status}
For the literal cofinal Standard-Model overlap system, V58 has not supplied
a regulator-uniform gap \(g_n\), the center certificates
\(\kappa_{n,\alpha}\), the second-derivative stock \(S_n\), or the decay
(68.20).  The theorem proves that a growing number of charts is not itself
an analytic obstruction, but checking all centers can still be an
unbounded computational task.
\end{firewall}

\subsection{Acquisition and next bottleneck}

V58 globalizes local vertical-anomaly estimates without a chart-count
penalty and gives the exact truncation-order schedule needed when gaps
shrink.  The next upstream bottleneck lies at the chart centers: ordinary
continuum anomaly arithmetic does not automatically cancel a resolvent
built from a nontrivial background.  The next note must replace
center-by-center brute force by a covariance or transgression identity
which transports the compiled moment map from a reference configuration
with a controlled local remainder.

\begin{assessment}[V58 acquisition and boundary]
\label{ass:v18v58}
V58 proves existence of finite fixed-cutoff resolvent atlases, an explicit
Lipschitz stock for \(dJ^{\rm tot}\), a finite-center global certificate
with no covering-number loss, and the cofinal order criterion (68.18).
The remaining gap is the production of the compiled center certificates
on nontrivial backgrounds.
\end{assessment}

\begin{record}[V58 append record]
\label{rec:v18v58}
V58 is appended to the validated V57 whose installed SHA--256 was
\texttt{A3FFE24F32B2175F2ED937324DC2AC99A87AFE5AD10B9A5A0D8F9B42EB215DBF}.
After installation only V58 is the active numeric SM note; the frozen
\texttt{V100\_f2} through \texttt{V100\_f17} archives are unchanged.
\end{record}

\section{V59: equivariant moment-map reduction and Kato transgression}
\label{sec:v18v59}

V58 left a finite bank of center certificates on the configuration space.
Gauge covariance makes many of those centers redundant: the compiled
moment map transforms by the coadjoint action, and its norm is constant on
gauge orbits.  This note proves that reduction, derives an exact Kato
transport formula between physical backgrounds, and identifies the
quotient-diameter stock needed to turn small \(dJ^{\rm tot}\) into a
uniformly normalized small \(J^{\rm tot}\).

\subsection{Coadjoint covariance}

Let a compact gauge group \(\mathcal G\) act on \(X\) from the left.  For
sector \(a\), assume a unitary representation \(G_a(g)\) satisfies
\[
\begin{split}
 H_a(g\cdot U)&=G_a(g)H_a(U)G_a(g)^{-1},\\
 P_a(g\cdot U)&=G_a(g)P_a(U)G_a(g)^{-1},\\
 \Omega_a(\operatorname{Ad}_g\omega)
 &=G_a(g)\Omega_a(\omega)G_a(g)^{-1}.                       \tag{69.1}
\end{split}
\]
Regard
\[
 J^{\rm tot}(U)\in\mathfrak g^*,\qquad
 \langle J^{\rm tot}(U),\omega\rangle
 =-i\sum_a\sigma_a\operatorname{Tr}
 \bigl(P_a(U)\Omega_a(\omega)\bigr).                         \tag{69.2}
\]

\begin{theorem}[Exact coadjoint equivariance]
\label{thm:v18v59-coadjoint}
The compiled moment map obeys
\[
\boxed{
 \langle J^{\rm tot}(g\cdot U),\omega\rangle
 =
 \langle J^{\rm tot}(U),\operatorname{Ad}_{g^{-1}}\omega\rangle.}       \tag{69.3}
\]
For \(X_U\in T_UX\),
\[
\begin{split}
 &d\langle J^{\rm tot},\omega\rangle_{g\cdot U}
 \bigl(dg(X_U)\bigr)\\
 &\qquad =
 d\langle J^{\rm tot},\operatorname{Ad}_{g^{-1}}\omega\rangle_U
 (X_U).                                                       \tag{69.4}
\end{split}
\]
\end{theorem}

\begin{proof}
Insert (69.1) into (69.2), move \(G_a(g)\) through the trace, and use the
third covariance identity with \(g^{-1}\).  This gives (69.3).
Differentiate (69.3) in the \(U\) variable to obtain (69.4).
\end{proof}

\begin{corollary}[Orbitwise reduction of the certificate bank]
\label{cor:v18v59-orbit-reduction}
Equip \(\mathfrak g\) with an \(\operatorname{Ad}\)-invariant norm and
\(\mathfrak g^*\) with its dual norm, and use a gauge-invariant metric on
\(X\).  Then
\[
 \|J^{\rm tot}(g\cdot U)\|_{\mathfrak g^*}
 =\|J^{\rm tot}(U)\|_{\mathfrak g^*},\qquad
 \|dJ^{\rm tot}(g\cdot U)\|=\|dJ^{\rm tot}(U)\|.             \tag{69.5}
\]
It is therefore enough to certify one center per gauge orbit, or one
finite atlas on a regular quotient stratum.
\end{corollary}

\begin{proof}
The coadjoint action is an isometry for the dual invariant norm, proving
the first equality.  Equation (69.4), together with invariance of the
tangent norm, proves the second.
\end{proof}

\begin{proposition}[Two-vertical-leg identity]
\label{prop:v18v59-two-vertical}
For \(\omega,\eta\in\mathfrak g\),
\[
 d\langle J^{\rm tot},\omega\rangle(\eta^\#)
 =\langle J^{\rm tot},[\omega,\eta]\rangle,                  \tag{69.6}
\]
and hence
\[
 F^{\rm tot}(\omega^\#,\eta^\#)
 =-\langle J^{\rm tot},[\omega,\eta]\rangle.                 \tag{69.7}
\]
\end{proposition}

\begin{proof}
Put \(g(t)=\exp(t\eta)\) in (69.3).  Since
\[
 \left.\frac d{dt}\right|_{t=0}
 \operatorname{Ad}_{g(t)^{-1}}\omega
 =-[\eta,\omega]=[\omega,\eta],                              \tag{69.8}
\]
differentiation gives (69.6).  Apply the V57 identity
\(\iota_{\omega^\#}F^{\rm tot}
=-d\langle J^{\rm tot},\omega\rangle\) to obtain (69.7).
\end{proof}

\begin{remark}[Abelian and nonabelian vertical planes]
\label{rem:v18v59-vertical-planes}
For an abelian factor, (69.7) vanishes identically.  For a nonabelian
factor, vanishing of \(J^{\rm tot}\) implies vanishing on every
two-vertical plane.  Neither statement controls a horizontal--horizontal
curvature component.
\end{remark}

\subsection{Kato transport along a physical path}

Let \(U(t)\), \(0\le t\le1\), be a \(C^1\) path inside one constant-rank
gapped component, and abbreviate \(P_a(t)=P_a(U(t))\).  Define the
anti-Hermitian Kato generator
\[
 K_a(t)=[\dot P_a(t),P_a(t)]                                \tag{69.9}
\]
and the unitary solution
\[
 \dot W_a(t)=K_a(t)W_a(t),\qquad W_a(0)=I.                  \tag{69.10}
\]

\begin{lemma}[Kato intertwining]
\label{lem:v18v59-kato}
For every \(t\),
\[
 P_a(t)=W_a(t)P_a(0)W_a(t)^*.                               \tag{69.11}
\]
\end{lemma}

\begin{proof}
Differentiating \(P_a^2=P_a\) gives
\(P_a\dot P_aP_a=(I-P_a)\dot P_a(I-P_a)=0\).  Direct block multiplication
relative to \(P_a\oplus(I-P_a)\) yields
\[
 [[\dot P_a,P_a],P_a]=\dot P_a.                             \tag{69.12}
\]
Both sides of (69.11) therefore solve the same linear differential
equation \(\dot Q=[K_a,Q]\) with the same initial value.
\end{proof}

\begin{theorem}[Exact compiled Kato transgression]
\label{thm:v18v59-kato-transgression}
For every fixed gauge parameter \(\omega\),
\[
\boxed{
\begin{split}
 &\langle J^{\rm tot}(U(1))-J^{\rm tot}(U(0)),\omega\rangle\\
 &\quad =
 -i\int_0^1\sum_a\sigma_a
 \operatorname{Tr}\!\left(
 P_a(0)W_a(t)^*
 [\Omega_a(\omega),K_a(t)]
 W_a(t)\right)dt .
\end{split}}                                                \tag{69.13}
\]
The signed representation sum occurs before any norm.
\end{theorem}

\begin{proof}
By Lemma \ref{lem:v18v59-kato} and cyclicity,
\[
 \langle J^{\rm tot}(U(t)),\omega\rangle
 =-i\sum_a\sigma_a\operatorname{Tr}
 \bigl(P_a(0)W_a(t)^*\Omega_a(\omega)W_a(t)\bigr).           \tag{69.14}
\]
Since \(K_a^*=-K_a\), (69.10) gives
\[
 \frac d{dt}\bigl(W_a^*\Omega_aW_a\bigr)
 =W_a^*[\Omega_a,K_a]W_a.                                  \tag{69.15}
\]
Differentiate (69.14) and integrate from zero to one.
\end{proof}

\begin{corollary}[Fallback path bound]
\label{cor:v18v59-kato-fallback}
If \(r_a=\operatorname{rank}P_a\), then
\[
\begin{split}
 &|\langle J^{\rm tot}(U(1))-J^{\rm tot}(U(0)),\omega\rangle|\\
 &\qquad\le
 2\int_0^1\sum_a r_a
 \|\Omega_a(\omega)\|\,\|K_a(t)\|\,dt.                      \tag{69.16}
\end{split}
\]
\end{corollary}

\begin{proof}
Apply \(\|[\Omega,K]\|\le2\|\Omega\|\|K\|\) and
\(|\operatorname{Tr}(PM)|\le r\|M\|\) only after writing the compiled
identity (69.13).
\end{proof}

\begin{firewall}[Kato transport preserves phase but does not cancel it]
\label{fire:v18v59-kato-not-cancellation}
Equation (69.13) is an exact phase-preserving transgression formula.  It
does not make its signed integrand small.  A useful Standard-Model
certificate must exploit representation identities inside that integrand
or bound a compiled remainder after those identities are applied.
\end{firewall}

\subsection{Quotient diameter and normalization of \(J^{\rm tot}\)}

Let \(B=X/\mathcal G\) be a connected regular quotient with its path
metric.  Assume horizontal lifts exist and preserve path length.  Let
\[
 D_B=\sup_{b_0,b_1\in B}d_B(b_0,b_1)                        \tag{69.17}
\]
be its diameter.  By (69.5), the norm of \(J^{\rm tot}\) and the operator
norm of its horizontal derivative are well defined on \(B\).

\begin{theorem}[Reference-to-global moment-map bound]
\label{thm:v18v59-diameter-bound}
Fix \(b_*\in B\).  If
\[
 \|J^{\rm tot}(b_*)\|_{\mathfrak g^*}\le\beta,\qquad
 \sup_{b\in B}\|d_{\rm hor}J^{\rm tot}(b)\|\le\alpha,        \tag{69.18}
\]
then
\[
\boxed{
 \sup_{b\in B}\|J^{\rm tot}(b)\|_{\mathfrak g^*}
 \le\beta+D_B\alpha.}                                      \tag{69.19}
\]
\end{theorem}

\begin{proof}
Join \(b_*\) to \(b\) by paths whose lengths decrease to
\(d_B(b_*,b)\), and horizontally lift them.  The fundamental theorem of
calculus and (69.18) give
\[
 \|J^{\rm tot}(b)-J^{\rm tot}(b_*)\|
 \le\alpha d_B(b_*,b).                                     \tag{69.20}
\]
Take the supremum and use (69.17).
\end{proof}

\begin{countertheorem}[Small derivative without diameter control]
\label{cth:v18v59-diameter-no-go}
The condition \(\alpha_n\to0\) alone does not imply
\(\sup\|J_n^{\rm tot}\|\to0\), even if \(J_n^{\rm tot}\) is normalized to
zero at one reference point.
\end{countertheorem}

\begin{proof}
Take \(B_n=[0,n]\), \(J_n(x)=x/n\), and reference point zero.  Then
\(\|dJ_n\|_\infty=1/n\to0\), while
\(\sup_{B_n}|J_n|=1\).  Here \(D_{B_n}\alpha_n=1\), exactly the term retained
in (69.19).
\end{proof}

\begin{remark}[Components and stabilizers]
\label{rem:v18v59-components}
One reference normalization is required on every connected quotient
component.  Singular strata must separately pass the V53 stabilizer
character test; coadjoint covariance on the regular stratum cannot remove
a nontrivial isotropy action on the determinant fibre.
\end{remark}

\subsection{A quotient-level finite-center certificate}

Let \(\{b_\alpha\}\) be centers of a finite atlas on \(B\), and lift one
representative \(U_\alpha\) of each center.  Suppose V58 gives
\[
 \|dJ^{\rm tot}\|_{B(U_\alpha,\delta_\alpha)}
 \le\kappa_\alpha+K_{J,\alpha}\delta_\alpha.                \tag{69.21}
\]

\begin{proposition}[Gauge-saturated lift]
\label{prop:v18v59-saturated-lift}
The bound (69.21) holds with the same right side on the gauge saturation
of the lifted chart.  Hence no gauge-orbit volume, orbit diameter, or
number of gauge copies enters the V58 maximum.
\end{proposition}

\begin{proof}
Equation (69.5) makes the norm of \(dJ^{\rm tot}\) constant along every
orbit.  Every point in the saturation is a gauge transform of a point in
the lifted slice, so (69.21) transfers without loss.
\end{proof}

\begin{theorem}[Normalized cofinal moment-map capsule]
\label{thm:v18v59-cofinal-normalization}
At cutoff \(n\), let
\[
 \alpha_n=
 \max_\alpha(\kappa_{n,\alpha}
 +K_{J,n,\alpha}\delta_{n,\alpha}),                         \tag{69.22}
\]
let \(D_n\) be the regular quotient diameter, and let \(\beta_n\) bound the
compiled moment map at one reference point of each component.  If
\[
 \beta_n+D_n\alpha_n\longrightarrow0,                       \tag{69.23}
\]
and the V57--V58 topological, gauge-loop, locality, reflection, and
transport hypotheses hold, then both the infinitesimal vertical curvature
and the normalized gauge moment map vanish in the continuum Ward
residual.
\end{theorem}

\begin{proof}
V58 gives \(\|dJ_n^{\rm tot}\|\le\alpha_n\) on the quotient atlas, and
Proposition \ref{prop:v18v59-saturated-lift} extends it to the full regular
stratum.  Theorem \ref{thm:v18v59-diameter-bound} and (69.23) give
\(\|J_n^{\rm tot}\|\to0\).  The V57 identity converts the first bound to
vertical-curvature decay.  The remaining stated hypotheses carry both
terms through the Ward comparison and exclude global flat phases.
\end{proof}

\begin{firewall}[Unpopulated geometry row]
\label{fire:v18v59-unpopulated}
No uniform quotient diameter \(D_n\), reference bound \(\beta_n\), or
compiled Kato-integrand cancellation has been proved for the literal
cofinal Standard-Model overlap system.  The gauge reduction removes
redundant orbit copies from the certificate bank, but it does not control
the size or geometry of the physical quotient.
\end{firewall}

\subsection{Acquisition and next bottleneck}

V59 reduces V58 to a quotient atlas and adds the missing normalization
ledger \(\beta_n+D_n\alpha_n\).  The next upstream target is local rather
than topological: express the Kato generator through the resolvent,
separate its ultralocal principal symbol from its gap-controlled tail, and
test the V44 representation identities on the principal compiled
transgression before bounding the tail.

\begin{assessment}[V59 acquisition and boundary]
\label{ass:v18v59}
V59 proves exact coadjoint covariance, the two-vertical-leg identity, Kato
intertwining and compiled transgression, quotient-level certificate
reduction, and the diameter-sensitive normalization theorem (69.19).  The
remaining gap is a regulator-uniform compiled bound on the Kato or
resolvent transgression over nontrivial physical backgrounds.
\end{assessment}

\begin{record}[V59 append record]
\label{rec:v18v59}
V59 is appended to the validated V58 whose installed SHA--256 was
\texttt{EAF0C653D060EAD29FD37EF097F0D80C733DDC62A9AF3AEB50B4F07B549F7339}.
After installation only V59 is the active numeric SM note; the frozen
\texttt{V100\_f2} through \texttt{V100\_f17} archives are unchanged.
\end{record}

\section{V60: finite-range sign surrogates and the anomaly--locality tradeoff}
\label{sec:v18v60}

V59 identified the compiled Kato transgression but did not isolate a
finite local principal part.  Directly truncating a resolvent kernel leaves
coefficients which still depend on the global inverse.  V60 instead
approximates the gapped sign function by a polynomial.  The resulting
projector derivative is exactly finite range, its signed Standard-Model
trace is a finite algebraic word bank, and the discarded tail has an
explicit divided-difference bound.

\subsection{A \(C^1\) polynomial sign approximation}

Fix \(0<g<M\) and the disconnected spectral set
\[
 K_{g,M}=[-M,-g]\cup[g,M].                                  \tag{70.1}
\]

\begin{lemma}[Geometric \(C^1\) sign approximation]
\label{lem:v18v60-sign-polynomial}
There are constants \(C_\kappa<\infty\) and \(\rho_\kappa>1\), depending
only on \(\kappa=M/g\), and real odd polynomials \(p_m\) of degree at most
\(2m+1\) such that
\[
\boxed{
 \sup_{x\in K_{g,M}}|p_m(x)-\operatorname{sign}x|
 +g\sup_{x\in K_{g,M}}|p_m'(x)|
 \le C_\kappa\rho_\kappa^{-m}.}                             \tag{70.2}
\]
\end{lemma}

\begin{proof}
On \([g^2,M^2]\), the function \(f(y)=y^{-1/2}\) extends holomorphically
to a complex neighborhood whose relative size depends only on
\(\kappa\).  Bernstein's polynomial approximation theorem, applied on
the affinely rescaled interval, gives polynomials \(q_m\) for which both
\(q_m-f\) and \(q_m'-f'\) decay geometrically.  The derivative estimate
follows either by differentiating the Chebyshev series on a slightly
smaller Bernstein ellipse or by Cauchy's estimate there.  A polynomial
factor in \(m\) is absorbed by replacing the ellipse parameter by a
slightly smaller number still greater than one.

Set
\[
 p_m(x)=xq_m(x^2).                                          \tag{70.3}
\]
This polynomial is odd and has degree at most \(2m+1\).  Since
\(xf(x^2)=\operatorname{sign}x\) on \(K_{g,M}\), the zeroth-order error is
bounded by \(M\|q_m-f\|_\infty\).  Moreover,
\[
 p_m'(x)=q_m(x^2)+2x^2q_m'(x^2),                            \tag{70.4}
\]
while \(f(y)+2yf'(y)=0\).  The \(C^1\) approximation of \(q_m\) therefore
gives the second term in (70.2), after rescaling constants by powers of
\(g\) and \(M\).  Those ratios depend only on \(\kappa\).
\end{proof}

\subsection{Fréchet derivative control}

For a scalar function \(f\) on \(K_{g,M}\), write
\[
 f^{[1]}(x,y)=
 \begin{cases}
 \dfrac{f(x)-f(y)}{x-y},&x\ne y,\\
 f'(x),&x=y.
 \end{cases}                                                \tag{70.5}
\]

\begin{lemma}[Divided-difference tail]
\label{lem:v18v60-divided-tail}
Let \(e_m=p_m-\operatorname{sign}\) on \(K_{g,M}\).  Then
\[
 \sup_{x,y\in K_{g,M}}|e_m^{[1]}(x,y)|
 \le \frac{C_\kappa}{g}\rho_\kappa^{-m}.                    \tag{70.6}
\]
\end{lemma}

\begin{proof}
If \(x,y\) lie in the same connected interval of (70.1), the mean-value
theorem and (70.2) give
\[
 |e_m^{[1]}(x,y)|\le\sup_{K_{g,M}}|p_m'|
 \le C_\kappa g^{-1}\rho_\kappa^{-m}.                       \tag{70.7}
\]
If they lie in opposite intervals, then \(|x-y|\ge2g\), while
\(|e_m(x)-e_m(y)|\le2C_\kappa\rho_\kappa^{-m}\).  This gives the same
bound.
\end{proof}

\begin{theorem}[Projector derivative tail]
\label{thm:v18v60-projector-tail}
Let \(H\) be Hermitian with spectrum in \(K_{g,M}\), put
\[
 P=\frac{I-\operatorname{sign}H}{2},\qquad
 P_m=\frac{I-p_m(H)}2,                                     \tag{70.8}
\]
and let \(Q\) be a Hermitian perturbation.  In Hilbert--Schmidt norm,
\[
\boxed{
 \|DP_H[Q]-D(P_m)_H[Q]\|_{\rm HS}
 \le
 \frac{C_\kappa}{2g}\rho_\kappa^{-m}\|Q\|_{\rm HS}.}         \tag{70.9}
\]
\end{theorem}

\begin{proof}
Diagonalize \(H\), with eigenvalues \(\lambda_i\).  The finite-dimensional
Fréchet functional-calculus formula gives
\[
 \bigl(Df_H[Q]\bigr)_{ij}
 =f^{[1]}(\lambda_i,\lambda_j)Q_{ij}.                        \tag{70.10}
\]
Apply this to \(f=p_m-\operatorname{sign}\).  Squaring and summing the
matrix entries, then using Lemma \ref{lem:v18v60-divided-tail}, gives
\[
 \|D(p_m-\operatorname{sign})_H[Q]\|_{\rm HS}
 \le C_\kappa g^{-1}\rho_\kappa^{-m}\|Q\|_{\rm HS}.          \tag{70.11}
\]
The factor \(1/2\) in (70.8) proves (70.9).
\end{proof}

\subsection{Compiled local principal trace}

For sector \(a\), take \(H_a\) with the common spectral enclosure
\(K_{g,M}\), put \(Q_a=d_XH_a\), and define
\[
\begin{split}
 {\cal W}^{(m)}(X,\omega)
 &=-i\sum_a\sigma_a\operatorname{Tr}
 \bigl(D(P_{m,a})_{H_a}[Q_a]\,\Omega_a(\omega)\bigr),\\
 {\cal R}^{(m)}(X,\omega)
 &=d_XJ^{\rm tot}(\omega)-{\cal W}^{(m)}(X,\omega).          \tag{70.12}
\end{split}
\]

\begin{theorem}[Principal-plus-tail decomposition]
\label{thm:v18v60-principal-tail}
The decomposition (70.12) is exact and
\[
\boxed{
 |{\cal R}^{(m)}(X,\omega)|
 \le
 \frac{C_\kappa}{2g}\rho_\kappa^{-m}
 \sum_a
 \|Q_a\|_{\rm HS}\|\Omega_a(\omega)\|_{\rm HS}.}             \tag{70.13}
\]
If the finite compiled identity
\[
 {\cal W}^{(m)}(X,\omega)=0                                 \tag{70.14}
\]
holds, the right side of (70.13) bounds the full vertical curvature
\(|d_XJ^{\rm tot}(\omega)|\).
\end{theorem}

\begin{proof}
The first line of (70.12) is obtained from the exact formula
\(d_XJ^{\rm tot}=-i\sum_a\sigma_a
\operatorname{Tr}(DP_{a,H_a}[Q_a]\Omega_a)\) by replacing \(P_a\) with
\(P_{m,a}\).  Their difference is the second line.  Apply
Cauchy--Schwarz for the Hilbert--Schmidt inner product sector by sector,
then use (70.9).  Equation (70.14) leaves only the tail.
\end{proof}

\begin{assessment}[Where the absolute value is taken]
\label{ass:v18v60-compile-first}
The principal term in (70.12) is summed with the chiral signs before the
condition (70.14) is tested.  Only the approximation tail is bounded
sectorwise.  Thus V60 follows the phase-preserving rule extracted from YM
V8 and avoids erasing the V44 representation arithmetic.
\end{assessment}

\subsection{Exact finite range of the principal term}

Let the one-particle Hilbert space be carried by a finite graph.  Suppose
\(H_a\) has kernel range at most \(R_H\), \(Q_a\) has range at most \(R_Q\),
and \(\Omega_a(\omega)\) is diagonal in the site variable.

\begin{lemma}[Derivative of a matrix polynomial]
\label{lem:v18v60-polynomial-derivative}
If \(p(x)=\sum_{k=0}^dc_kx^k\), then
\[
 Dp_H[Q]
 =\sum_{k=1}^dc_k\sum_{j=0}^{k-1}
 H^jQH^{k-1-j}.                                             \tag{70.15}
\]
Consequently \(Dp_H[Q]\) has range at most
\[
 (d-1)R_H+R_Q.                                              \tag{70.16}
\]
\end{lemma}

\begin{proof}
Extract the coefficient linear in \(t\) from
\((H+tQ)^k\):
\[
 \left.\frac d{dt}\right|_{t=0}(H+tQ)^k
 =\sum_{j=0}^{k-1}H^jQH^{k-1-j}.                            \tag{70.17}
\]
This proves (70.15).  Kernel ranges add under matrix multiplication, which
gives (70.16).
\end{proof}

\begin{corollary}[Finite local word bank]
\label{cor:v18v60-local-word-bank}
For the polynomial \(p_m\) of Lemma
\ref{lem:v18v60-sign-polynomial}, the compiled principal trace
\({\cal W}^{(m)}\) is a finite sum of closed lattice walks contained within
range
\[
 R_{\rm prin}(m)=2mR_H+R_Q.                                 \tag{70.18}
\]
Each summand is a product of a lattice/spin kernel and a finite
representation trace.  Hence (70.14) is a finite exact algebraic
certificate at fixed cutoff and fixed \(m\).
\end{corollary}

\begin{proof}
Use (70.15) with degree at most \(2m+1\), then multiply by the site-diagonal
\(\Omega_a\) and take the trace.  A nonzero diagonal trace closes the
corresponding kernel walk.  Expanding gauge-link matrices in their sector
representations separates each finite walk coefficient into its
lattice/spin factor and its representation trace.
\end{proof}

\begin{firewall}[The V44 identities do not cancel every word]
\label{fire:v18v60-word-growth}
V44 verifies the Standard-Model linear, cubic, and mixed anomaly traces.
The bank in Corollary \ref{cor:v18v60-local-word-bank} generally contains
longer background-dependent representation words as \(m\) grows.
Anomaly-polynomial cancellation does not automatically imply (70.14).
The finite word bank is a decidable target, not a claimed zero.
\end{firewall}

\subsection{Cofinal degree versus physical locality}

At cutoff \(n\), absorb the tail prefactor into
\[
 K_n=
 \frac{C_{\kappa_n}}{2g_n}
 \sup_{X,\omega}
 \sum_a\|d_XH_{a,n}\|_{\rm HS}
 \|\Omega_{a,n}(\omega)\|_{\rm HS},                         \tag{70.19}
\]
and suppose \(\rho_n=\rho_{\kappa_n}>1\).  If the principal identity is
certified at degree \(m_n\), then
\[
 \|dJ_n^{\rm tot}\|\le K_n\rho_n^{-m_n}.                    \tag{70.20}
\]
Let \(a_n\) be the lattice spacing.  The physical range of the principal
word bank is bounded by
\[
 \ell_n^{\rm prin}
 =a_n(2m_nR_{H,n}+R_{Q,n}).                                 \tag{70.21}
\]

\begin{proposition}[Simultaneous tail decay and shrinking range]
\label{prop:v18v60-locality-tradeoff}
Assume \(K_n\ge1\) and put
\[
 A_n=\frac{a_nR_{H,n}}{\log\rho_n}.                          \tag{70.22}
\]
There exist integers \(m_n\ge1\) such that
\[
 K_n\rho_n^{-m_n}\to0,\qquad
 a_nm_nR_{H,n}\to0                                          \tag{70.23}
\]
if
\[
 a_nR_{H,n}\to0,\qquad
 A_n\to0,\qquad
 A_n\log K_n\to0.                                           \tag{70.24}
\]
Conversely, (70.23) forces
\[
 a_nR_{H,n}\to0,\qquad A_n\log K_n\to0.                     \tag{70.25}
\]
\end{proposition}

\begin{proof}
Under (70.24), choose \(s_n\to\infty\) so slowly that
\(A_ns_n\to0\), and set
\[
 m_n=
 \left\lceil\frac{\log K_n+s_n}{\log\rho_n}\right\rceil.     \tag{70.26}
\]
Then \(K_n\rho_n^{-m_n}\le e^{-s_n}\).  Also
\[
 a_nR_{H,n}m_n
 \le A_n(\log K_n+s_n)+a_nR_{H,n}\to0.                      \tag{70.27}
\]
This proves sufficiency.  Conversely, the first relation in (70.23)
implies
\(m_n\log\rho_n-\log K_n\to+\infty\), hence
\[
 A_n\log K_n
 \le a_nR_{H,n}m_n\to0.                                    \tag{70.28}
\]
Since \(m_n\ge1\) eventually, the second relation in (70.23) also forces
\(a_nR_{H,n}\to0\).
\end{proof}

\begin{remark}[Condition-number bottleneck]
\label{rem:v18v60-condition-number}
When \(g_n/M_n\) collapses, \(\rho_n\downarrow1\), so \(A_n\) can fail to
vanish even though \(a_n\to0\).  A shrinking lattice spacing therefore
does not by itself reconcile a small sign-approximation tail with a
shrinking physical word range.
\end{remark}

\subsection{A polynomial vertical-Ward capsule}

\begin{theorem}[Finite-range compiled certificate capsule]
\label{thm:v18v60-polynomial-capsule}
Suppose on a cofinal sequence:
\begin{enumerate}[leftmargin=2.2em]
\item every sector spectrum lies in \(K_{g_n,M_n}\);
\item the finite local identity (70.14) is certified on the V58 quotient
      center bank at degrees \(m_n\);
\item the center displacement, reference normalization, quotient diameter,
      and transport rows of V58--V59 tend to zero;
\item the locality tradeoff (70.23) holds;
\item the V53--V57 topology, gauge-loop, reflection, and Ward-collar rows
      hold.
\end{enumerate}
Then the compiled measure-line contribution to the gauge Ward residual
vanishes, and the finite principal certificate has shrinking physical
range.
\end{theorem}

\begin{proof}
Theorem \ref{thm:v18v60-principal-tail} and item two bound every center
residual by (70.20).  Items three and the V58 finite-center theorem extend
the bound over each quotient component and normalize \(J_n^{\rm tot}\).
The first relation in (70.23) makes the vertical curvature vanish; the
second gives shrinking physical support for the certified principal bank.
Item five passes the result through the global determinant-line and
continuum Ward handoffs.
\end{proof}

\begin{firewall}[Unpopulated finite word identities]
\label{fire:v18v60-unpopulated}
No identities (70.14) have been checked for the literal background-dependent
Standard-Model word banks, and no cofinal values of
\(g_n,M_n,\rho_n,K_n\), or \(m_n\) have been supplied.  The main unresolved
question is whether gauge covariance and the anomaly identities collapse
the long word bank to the V44 trace basis, or whether additional local
counterterm data are necessary.
\end{firewall}

\subsection{Checkpoint and next step}

V60 reaches the next five-version boundary.  The active Yang--Mills and
Navier--Stokes notes must be inspected before V61.  After that audit, the
next step is to classify the representation words in (70.15) under gauge
covariance and cyclic trace, separating the universal anomaly basis from
gauge-invariant horizontal responses which do not enter the vertical Ward
obstruction.

\begin{assessment}[V60 acquisition and boundary]
\label{ass:v18v60}
V60 proves a geometrically convergent \(C^1\) polynomial approximation to
the gapped sign projector, a Hilbert--Schmidt derivative-tail estimate, an
exact finite-range compiled principal word bank, and the sharp
tail-versus-physical-range criterion (70.24)--(70.25).  The literal finite
word cancellation remains unproved.
\end{assessment}

\begin{record}[V60 append record]
\label{rec:v18v60}
V60 is appended to the validated V59 whose installed SHA--256 was
\texttt{85BDF1A3B90472B364C2AF92AA9FF59D84AC94BB98DEF2D2DB551DE54A6269F3}.
After installation only V60 is the active numeric SM note; the frozen
\texttt{V100\_f2} through \texttt{V100\_f17} archives are unchanged.
\end{record}

\section{V61: live link-word expansion and a finite smooth-chart anomaly pilot}
\label{sec:v18v61}

V60 reached a checkpoint with a finite-range polynomial word bank, but
the Standard-Model anomaly identities do not automatically cancel its
long background-dependent representation words.  The mandatory companion
audit changes the next move: as in YM V9, retain a finite live expansion
to sufficiently high order and bound only its exact tail.  V61 develops
that device for lattice link words and states the finite pilot certificate
which must be executed before a global atlas is attempted.

\subsection{Mandatory V60 companion audit}

The active companion notes at this checkpoint were
\[
\begin{array}{c|c|c}
\hbox{file}&\hbox{bytes}&\hbox{SHA--256}\\ \hline
\texttt{Renewal\_Yang\_Mills\_Mass\_Gap\_Research\_Notes\_V9.tex}
&160441&
\texttt{C97496D124E981A1BFC9A23C13B346F82A368D7E46FDB1B5560D0711104F5D86}\\
\texttt{Renewal\_NS\_Stability\_Research\_Notes\_V26.tex}
&278283&
\texttt{82C887F8C2C69E4F28FAD7C3DC8DE5B9099CB5A4BF431EBA1FFF3E91935978AD}.
\end{array}                                                   \tag{71.1}
\]
YM V9 proves on one exact rational pilot box that retaining higher
Neumann powers changes the first passing radius from \(2^{-22}\) at first
order to \(2^{-8}\) at order \(47\).  It also proves that this improvement
does not yet make a uniform four-dimensional cover feasible: coefficient
growth and chart conditioning remain live.  Its prescribed next move is
to test preconditioned high-alias charts before building a cover.  The
transfer to the present programme is exact: test a compiled SM link-word
pilot with a live higher-order tail before expanding the V58 center bank.
NS V26 is unchanged and supplies no better estimate for this task.

\subsection{A product-exponential remainder}

For a matrix \(X\), put
\[
 T_D(X)=\sum_{j=0}^{D}\frac{X^j}{j!}.                        \tag{71.2}
\]

\begin{lemma}[Single exponential tail]
\label{lem:v18v61-single-exp}
If \(\|X\|\le\eta\), then
\[
 \|e^X-T_D(X)\|
 \le e^\eta\frac{\eta^{D+1}}{(D+1)!},\qquad
 \|e^X\|,\|T_D(X)\|\le e^\eta.                              \tag{71.3}
\]
\end{lemma}

\begin{proof}
The norm of the exponential tail is at most
\(\sum_{j>D}\eta^j/j!\).  Since
\[
 \sum_{j>D}\frac{\eta^j}{j!}
 =\frac{\eta^{D+1}}{(D+1)!}
 \sum_{k\ge0}
 \frac{\eta^k}{(D+2)\cdots(D+1+k)}
 \le e^\eta\frac{\eta^{D+1}}{(D+1)!},                       \tag{71.4}
\]
the first estimate follows.  The other two follow from the full and
truncated exponential series.
\end{proof}

\begin{lemma}[Live product tail]
\label{lem:v18v61-product-tail}
Let \(B_0,\ldots,B_k\) and \(X_1,\ldots,X_k\) be matrices, with
\(\|X_j\|\le\eta\).  Define
\[
\begin{split}
 W&=B_0e^{X_1}B_1\cdots e^{X_k}B_k,\\
 W_D&=B_0T_D(X_1)B_1\cdots T_D(X_k)B_k,\\
 B_*&=\prod_{j=0}^{k}\|B_j\|.
                                                                  \tag{71.5}
\end{split}
\]
Then
\[
\boxed{
 \|W-W_D\|
 \le
 B_*\,k\,e^{k\eta}
 \frac{\eta^{D+1}}{(D+1)!}.}                                \tag{71.6}
\]
\end{lemma}

\begin{proof}
Replace the \(k\) exponential factors one at a time.  The exact telescoping
sum has \(k\) terms.  In each term, Lemma
\ref{lem:v18v61-single-exp} bounds the replaced factor by
\(e^\eta\eta^{D+1}/(D+1)!\), and bounds each of the other \(k-1\) live or
truncated factors by \(e^\eta\).  Multiplication by all \(B_j\) and summing
the \(k\) terms gives (71.6).
\end{proof}

\begin{corollary}[Trace tail]
\label{cor:v18v61-trace-tail}
On a \(d\)-dimensional representation space,
\[
 |\operatorname{Tr}(W-W_D)|
 \le dB_*\,k\,e^{k\eta}
 \frac{\eta^{D+1}}{(D+1)!}.                                 \tag{71.7}
\]
\end{corollary}

\begin{proof}
Use \(|\operatorname{Tr}A|\le d\|A\|\) in (71.6).
\end{proof}

\subsection{Link logarithm charts}

Let a local word use oriented links \(\ell_1,\ldots,\ell_k\).  In a chart
centered at \(U_0\), write
\[
 U_{\ell_j}=U_{0,\ell_j}e^{X_{\ell_j}},\qquad
 \|d\rho_a(X_{\ell_j})\|\le\eta,                            \tag{71.8}
\]
where \(\rho_a\) is the representation of sector \(a\).  Orientation
reversal replaces the corresponding factor by
\[
 \rho_a(U_{\ell_j}^{-1})
 =e^{-d\rho_a(X_{\ell_j})}\rho_a(U_{0,\ell_j}^{-1}),         \tag{71.9}
\]
so Lemma \ref{lem:v18v61-product-tail} applies to both orientations.

\begin{proposition}[Finite live polynomial for a link word]
\label{prop:v18v61-link-polynomial}
For fixed \(D\), replacing every exponential in a representation word by
\(T_D\) produces a finite polynomial in the matrices
\(d\rho_a(X_\ell)\).  Its coefficients retain the ordered center factors
\(\rho_a(U_{0,\ell}^{\pm1})\).  The trace error obeys (71.7), with \(B_*\)
equal to the product of the fixed insertion and center-factor norms.
\end{proposition}

\begin{proof}
Equations (71.8)--(71.9) put every variable link factor into the form
required by Lemma \ref{lem:v18v61-product-tail}.  Each \(T_D\) is a finite
matrix polynomial, so their ordered product is finite.  The lemma and
Corollary \ref{cor:v18v61-trace-tail} give the error.
\end{proof}

\begin{remark}[Identity and nontrivial centers]
\label{rem:v18v61-centers}
At the identity center, every fixed link factor is the identity, and the
polynomial coefficients reduce to ordinary ordered Lie-algebra traces.
At a nontrivial center they contain holonomy insertions
\(\rho_a(U_{0,\ell})\).  The V44 anomaly arithmetic applies directly only
after those inserted words have been reduced to its universal trace basis.
\end{remark}

\subsection{Application to the V60 principal bank}

Write the finite V60 principal term at one quotient center as
\[
 {\cal W}^{(m)}(X,\omega)
 =\sum_{w\in{\cal P}_m}\sum_a
 \sigma_a c_{a,w}\operatorname{Tr}_{R_a}W_{a,w},             \tag{71.10}
\]
where \({\cal P}_m\) is the finite closed-walk bank and the coefficients
\(c_{a,w}\) include the lattice and spin contractions.  Let
\({\cal W}^{(m,D)}\) be obtained by the link replacement of Proposition
\ref{prop:v18v61-link-polynomial}.  For a word \(w\), denote its number of
variable exponential factors by \(k_w\), its representation dimension by
\(d_a\), and its fixed-factor product by \(B_{a,w,*}\).

\begin{theorem}[Compiled polynomial plus live link tail]
\label{thm:v18v61-compiled-link-tail}
The identity
\[
 {\cal W}^{(m)}
 ={\cal W}^{(m,D)}+{\cal T}^{(m,D)}                          \tag{71.11}
\]
is exact, and
\[
\boxed{
 |{\cal T}^{(m,D)}|
 \le
 \sum_{w\in{\cal P}_m}\sum_a
 |c_{a,w}|d_aB_{a,w,*}\,
 k_we^{k_w\eta}
 \frac{\eta^{D+1}}{(D+1)!}.}                               \tag{71.12}
\]
The polynomial \({\cal W}^{(m,D)}\) is compiled with the chiral signs
\(\sigma_a\) before its coefficients are tested.  Absolute values occur
only in the tail (71.12).
\end{theorem}

\begin{proof}
Apply Proposition \ref{prop:v18v61-link-polynomial} to every word in the
finite bank (71.10), preserving its signed sector sum.  Subtraction defines
the exact tail in (71.11).  Corollary \ref{cor:v18v61-trace-tail} and one
final triangle inequality give (71.12).
\end{proof}

\begin{assessment}[Two independent approximation levels]
\label{ass:v18v61-two-level}
The sign-polynomial degree \(m\) controls the V60 functional-calculus tail
and spatial range.  The live link order \(D\) controls the expansion of the
resulting finite local words inside one chart.  Increasing \(D\) does not
increase their spatial support, but it does increase the coefficient bank.
\end{assessment}

\subsection{Finite exact pilot certificate}

Expand the compiled polynomial into a fixed monomial basis,
\[
 {\cal W}^{(m,D)}
 =\sum_{\nu\in{\cal I}_{m,D}}C_\nu(U_0)\,M_\nu(X,\omega).    \tag{71.13}
\]
The coefficients \(C_\nu(U_0)\) are signed sums of finite representation
traces, including center holonomies when \(U_0\ne I\).

\begin{theorem}[One-chart pilot enclosure]
\label{thm:v18v61-pilot}
Suppose exact symbolic arithmetic or outward-rounded interval arithmetic
certifies
\[
 \sup_{\|X_\ell\|\le\eta,\,\omega,X\in\mathfrak B}
 |{\cal W}^{(m,D)}(X,\omega)|\le\zeta_{m,D}.                 \tag{71.14}
\]
Let \(T_{m,D}\) denote the right side of (71.12), and let
\[
 S_m=
 \frac{C_\kappa}{2g}\rho_\kappa^{-m}
 \sup_{\omega,X}\sum_a
 \|d_XH_a\|_{\rm HS}\|\Omega_a(\omega)\|_{\rm HS}            \tag{71.15}
\]
be the V60 sign tail.  Then throughout that chart,
\[
\boxed{
 |d_XJ^{\rm tot}(\omega)|
 \le\zeta_{m,D}+T_{m,D}+S_m.}                               \tag{71.16}
\]
\end{theorem}

\begin{proof}
The V60 principal-plus-tail theorem bounds
\(dJ^{\rm tot}-{\cal W}^{(m)}\) by \(S_m\).  Theorem
\ref{thm:v18v61-compiled-link-tail} bounds
\({\cal W}^{(m)}-{\cal W}^{(m,D)}\) by \(T_{m,D}\).  Apply the certified
bound (71.14) and add the three terms.
\end{proof}

\begin{remark}[Exact zero and interval alternatives]
\label{rem:v18v61-exact-or-interval}
At a symmetric center, trace identities and lattice reflection may prove
\(\zeta_{m,D}=0\) coefficient by coefficient.  At a general algebraic
center, rational interval enclosures can instead prove a small nonzero
\(\zeta_{m,D}\).  Both are finite, reviewable certificates; neither is
asserted here for the literal Standard Model.
\end{remark}

\begin{firewall}[Counterterms do not erase the curvature target]
\label{fire:v18v61-counterterm}
A change of determinant-line section adds an exact one-form to its
connection but does not change the curvature or
\(dJ^{\rm tot}=-\iota_{\omega^\#}F^{\rm tot}\).  Therefore a long word in
(71.13) cannot simply be declared removable by a local phase counterterm.
It must cancel in the compiled vertical curvature or enter a quantitatively
vanishing lattice remainder.
\end{firewall}

\subsection{Cofinal factorial ledger}

Define
\[
 K_{m,n}^{\rm word}
 =
 \sum_{w\in{\cal P}_{m_n}}\sum_a
 |c_{a,w}|d_aB_{a,w,*},\qquad
 k_n=\max_{w\in{\cal P}_{m_n}}k_w.                           \tag{71.17}
\]
Then
\[
 T_{m_n,D_n}
 \le K_{m,n}^{\rm word}k_ne^{k_n\eta_n}
 \frac{\eta_n^{D_n+1}}{(D_n+1)!}.                           \tag{71.18}
\]

\begin{proposition}[Exact tail-decay test]
\label{prop:v18v61-tail-test}
For positive stocks, the right side of (71.18) tends to zero if and only if
\[
\begin{split}
 &\log K_{m,n}^{\rm word}+\log k_n+k_n\eta_n\\
 &\quad +(D_n+1)\log\eta_n-\log((D_n+1)!)
 \longrightarrow-\infty.                                  \tag{71.19}
\end{split}
\]
\end{proposition}

\begin{proof}
Take logarithms of the positive expression in (71.18).
\end{proof}

\begin{corollary}[Uniform small-link chart]
\label{cor:v18v61-small-link}
If \(0<\eta_n\le\eta_*<1\), \(k_n\eta_n=O(1)\), and
\(\log(K_{m,n}^{\rm word}k_n)=o(D_n)\), then any \(D_n\to\infty\) satisfying
the last relation makes the tail vanish.
\end{corollary}

\begin{proof}
Already
\((D_n+1)\log\eta_*\to-\infty\); the factorial supplies additional decay.
The coefficient hypothesis and bounded \(k_n\eta_n\) make the positive
terms lower order.
\end{proof}

\begin{firewall}[Rough admissible links]
\label{fire:v18v61-rough-links}
Plaquette admissibility does not imply that every link has a uniformly
small logarithm in one global gauge.  Thus \(\eta_n\to0\) is natural on
smooth continuum charts but is not automatic on the full admissible
component.  Gauge fixing, high-alias coordinates, or a different center
family is required before (71.19) can be used globally.
\end{firewall}

\subsection{Pilot-first acquisition order}

The YM V9 audit rules out forming a large atlas before one chart is shown
feasible.  The next concrete SM acquisition order is:
\begin{enumerate}[leftmargin=2.2em]
\item choose the identity or a flat commuting background and a finite V60
      degree \(m\);
\item generate the closed-walk word bank \({\cal P}_m\) without taking
      sectorwise absolute values;
\item reduce the coefficients of (71.13) by cyclic trace, orientation
      reversal, lattice reflection, and the V44 representation identities;
\item certify \(\zeta_{m,D}\), \(T_{m,D}\), and \(S_m\) for increasing
      \(D\), recording the first order at which (71.16) meets a prescribed
      Ward budget;
\item only after a feasible radius is found, test preconditioned
      nontrivial centers and the V58 covering ledger.
\end{enumerate}

\begin{firewall}[Pilot not executed]
\label{fire:v18v61-pilot-open}
V61 supplies the exact finite algorithm and its error formulas but does not
generate the literal overlap word bank or execute the interval certificate.
The full Standard-Model chiral measure, its cofinal Ward limit, and the
SM--Einstein continuum therefore remain open.
\end{firewall}

\begin{assessment}[V61 acquisition and boundary]
\label{ass:v18v61}
V61 incorporates the YM V9 higher-order lesson, proves the matrix
product-exponential tail, combines it with the V60 sign tail, and gives the
finite pilot enclosure (71.16) and exact factorial ledger (71.19).  The
next bottleneck is executable generation and reduction of the first
compiled overlap word bank.
\end{assessment}

\begin{record}[V61 append record]
\label{rec:v18v61}
V61 is appended to the validated V60 whose installed SHA--256 was
\texttt{B905855A73FD853773D0AAF7E3E6AE017DF8F47E1D0B7DC0E266F0764211450D}.
After installation only V61 is the active numeric SM note; the frozen
\texttt{V100\_f2} through \texttt{V100\_f17} archives are unchanged.
\end{record}

\section{V62: exact hypercharge word pilot and the coefficientwise-zero obstruction}
\label{sec:v18v62}

V61 called for the first executable compiled word pilot.  The attached
sources do not specify a literal finite Wilson kernel with enough numerical
data to enumerate its lattice walks, but they do specify the generated
one-generation chiral charge packet exactly.  This determines the
representation factor of every commuting central word.  V62 computes that
factor with exact integers and proves that the strong V60 condition of
coefficientwise vanishing cannot follow from Standard-Model charge
arithmetic alone.

\subsection{The generated left-chiral packet}

Use the integer normalization \(y=6Y\).  After replacing right-handed
fields by their left-handed conjugates, the five fermion rows are
\[
\begin{array}{c|ccccc}
 a&Q&U^c&D^c&L&E^c\\ \hline
 d_a&6&3&3&2&1\\
 y_a&1&-4&2&-3&6 .
\end{array}                                                   \tag{72.1}
\]
Here \(d_a\) includes colour and weak multiplicity.  Define the exact
central word moments
\[
 S_r=\sum_a d_ay_a^r
 =6+3(-4)^r+3(2)^r+2(-3)^r+6^r.                            \tag{72.2}
\]

\begin{theorem}[Exact moment panel]
\label{thm:v18v62-moments}
For one generation,
\[
\begin{array}{c|rrrrrrrrrr}
 r&0&1&2&3&4&5&6&7&8&9\\ \hline
 S_r&15&0&120&0&2280&4320&60600&226800&1890120&9253440 .
\end{array}                                                   \tag{72.3}
\]
In particular,
\[
 S_1=S_3=0,\qquad S_2=120,\quad S_4=2280,\quad S_5=4320.     \tag{72.4}
\]
For three identical generations every entry is multiplied by three.
\end{theorem}

\begin{proof}
Substitute the five integer pairs \((d_a,y_a)\) from (72.1) into (72.2).
All operations are in \(\mathbb Z\).  For example,
\[
\begin{split}
 S_1&=6-12+6-6+6=0,\\
 S_3&=6-192+24-54+216=0,\\
 S_5&=6-3072+96-486+7776=4320.                              \tag{72.5}
\end{split}
\]
The remaining entries are obtained by the same exact evaluation and are
independently asserted by the executable recorded below.
\end{proof}

\subsection{Local and global anomaly checks}

In the normalization (72.1), the mixed anomaly coefficients reduce to
\[
\begin{split}
 A_{SU(3)^2U(1)}&=2y_Q+y_{U^c}+y_{D^c}=2-4+2=0,\\
 A_{SU(2)^2U(1)}&=3y_Q+y_L=3-3=0.                           \tag{72.6}
\end{split}
\]
The pure colour coefficient and the weak global parity are
\[
 2-1-1=0,\qquad (3+1)\bmod2=0.                              \tag{72.7}
\]

\begin{proposition}[Generated anomaly panel vanishes]
\label{prop:v18v62-anomaly-panel}
The one-generation packet has zero mixed gravitational \(U(1)\), cubic
\(U(1)\), mixed colour, mixed weak, pure colour, and mod-two weak anomaly.
Every displayed field row, including the Higgs row, descends through the
connected \(\mathbb Z_6\) quotient.
\end{proposition}

\begin{proof}
The gravitational and cubic coefficients are \(S_1\) and \(S_3\), which
vanish by Theorem \ref{thm:v18v62-moments}.  Equations (72.6)--(72.7) give
the other local and global coefficients.  For triality \(t\), weak parity
\(w\), and charge \(y\), descent requires
\[
 6\mid 2t+3w+y.                                              \tag{72.8}
\]
For \(Q,U^c,D^c,L,E^c,H\), the six exponents are respectively
\[
 6,\ -6,\ 0,\ 0,\ 6,\ 6,                                   \tag{72.9}
\]
so each is divisible by six.  These are the same finite identities
formalized in
\(\texttt{GeneratedMatterAnomalyCancellationExact.lean}\) and
\(\texttt{AnomalyFreeFieldShellGaugeQuotientExact.lean}\).
\end{proof}

\subsection{A no-go for raw coefficientwise cancellation}

\begin{countertheorem}[Anomaly freedom does not kill generic word moments]
\label{cth:v18v62-coefficientwise-no-go}
The Standard-Model identities in Proposition
\ref{prop:v18v62-anomaly-panel} do not imply
\[
 S_r=0\quad\hbox{for every positive word degree }r.          \tag{72.10}
\]
This failure occurs already at \(r=2\), and the first odd failure above the
cubic anomaly degree occurs at \(r=5\).
\end{countertheorem}

\begin{proof}
Equation (72.4) gives \(S_2=120\ne0\) and \(S_5=4320\ne0\), while all the
anomaly coefficients in Proposition \ref{prop:v18v62-anomaly-panel}
vanish.
\end{proof}

\begin{corollary}[Correction to the V60 principal identity]
\label{cor:v18v62-v60-correction}
The V60 condition
\[
 {\cal W}^{(m)}(X,\omega)=0                                 \tag{72.11}
\]
remains a valid sufficient certificate for vertical-curvature
cancellation.  It cannot be proved merely by replacing every
representation word with a generic statement that all signed
Standard-Model traces vanish.  A term whose representation factor is
\(S_2,S_4\), or \(S_5\) needs an independent spin, lattice, orientation, or
scaling cancellation.
\end{corollary}

\begin{proof}
The V60 theorem is an implication and is unaffected.  The proposed generic
charge argument would imply (72.10), contradicted by Theorem
\ref{thm:v18v62-moments}.
\end{proof}

\begin{assessment}[Physical even words are retained]
\label{ass:v18v62-even-words}
The nonzero even moments are not defects by themselves.  They occur in
orientation-even gauge response, including vacuum-polarization and kinetic
channels.  Erasing them would also erase physical Standard-Model response.
The vertical anomaly pilot must first project to the orientation-odd Ward
channel.
\end{assessment}

\subsection{Exact orientation pairing}

For a common commuting central phase \(\theta\), define the signed
orientation-odd packet
\[
 \Phi(\theta)
 =\frac1{2i}\sum_a d_a
 \left(e^{iy_a\theta}-e^{-iy_a\theta}\right)
 =\sum_a d_a\sin(y_a\theta).                                \tag{72.12}
\]
Also define the absolute moments
\[
 A_r=\sum_a d_a|y_a|^r.                                     \tag{72.13}
\]
Exact evaluation gives
\[
 A_7=333852.                                                 \tag{72.14}
\]

\begin{theorem}[Fifth-order onset of the paired central word]
\label{thm:v18v62-fifth-onset}
For every real \(\theta\),
\[
\boxed{
 \Phi(\theta)=36\theta^5+R_7(\theta),\qquad
 |R_7(\theta)|
 \le\frac{27821}{420}|\theta|^7.}                           \tag{72.15}
\]
Consequently, if
\[
 0<\theta^2<\frac{15120}{27821},                             \tag{72.16}
\]
then \(\Phi(\theta)\) has the sign of \(\theta\) and is nonzero.
\end{theorem}

\begin{proof}
Taylor's theorem through degree six gives, for real \(x\),
\[
 \sin x=x-\frac{x^3}{3!}+\frac{x^5}{5!}+r_7(x),
 \qquad |r_7(x)|\le\frac{|x|^7}{7!}.                        \tag{72.17}
\]
Sum (72.17) with weights \(d_a\) and arguments \(y_a\theta\).
The \(S_1\) and \(S_3\) terms vanish, while
\(S_5/5!=4320/120=36\).  The absolute remainder is bounded by
\[
 \frac{A_7}{7!}|\theta|^7
 =\frac{333852}{5040}|\theta|^7
 =\frac{27821}{420}|\theta|^7,                              \tag{72.18}
\]
which proves (72.15).  Under (72.16), the right side of (72.18) is
strictly smaller than \(36|\theta|^5\), proving the sign statement.
\end{proof}

\begin{corollary}[Three generations]
\label{cor:v18v62-three-generations}
For three identical generations,
\[
 \Phi_3(\theta)=108\theta^5+R_{7,3}(\theta),\qquad
 |R_{7,3}(\theta)|
 \le\frac{27821}{140}|\theta|^7,                            \tag{72.19}
\]
with the same sign radius (72.16).
\end{corollary}

\begin{proof}
Multiply (72.15) by three.
\end{proof}

\begin{firewall}[The fifth-order term is not yet an overlap anomaly]
\label{fire:v18v62-not-overlap-anomaly}
The function \(\Phi\) is the exact representation factor of one common
commuting orientation-paired phase.  V62 has not proved that the literal
overlap vertical curvature contains this word with a nonzero spin/lattice
coefficient.  The theorem rules out a proposed blanket cancellation; it
does not prove a gauge anomaly in the Standard Model.  In a correct overlap
calculation the fifth-order word may cancel across paths or be suppressed
by powers of the lattice spacing.
\end{firewall}

\subsection{A revised principal decomposition}

Let \(\Pi_{\rm odd}\) pair every path with its reversed orientation before
any norm, and let \(\Pi_{\rm an}\) extract the degree-one, degree-three,
mixed, and pure anomaly trace basis verified in Proposition
\ref{prop:v18v62-anomaly-panel}.  On a finite V61 polynomial bank write
\[
 {\cal W}^{(m,D)}
 =
 {\cal W}_{\rm even}
 +{\cal W}_{\rm an}
 +{\cal W}_{\ge5},                                          \tag{72.20}
\]
where
\[
\begin{split}
 {\cal W}_{\rm even}&=(I-\Pi_{\rm odd}){\cal W}^{(m,D)},\\
 {\cal W}_{\rm an}&=\Pi_{\rm an}\Pi_{\rm odd}{\cal W}^{(m,D)},\\
 {\cal W}_{\ge5}&=(I-\Pi_{\rm an})\Pi_{\rm odd}{\cal W}^{(m,D)}.
                                                                  \tag{72.21}
\end{split}
\]

\begin{proposition}[Correct finite pilot logic]
\label{prop:v18v62-revised-pilot}
Suppose the path bank is closed under orientation reversal.  Then:
\begin{enumerate}[leftmargin=2.2em]
\item the even and odd parts in (72.20) are exact complementary
      projections;
\item the V44 and V62 anomaly arithmetic kills the representation
      coefficients of \({\cal W}_{\rm an}\);
\item the remaining vertical Ward budget must bound
      \({\cal W}_{\ge5}\) together with the V60 sign tail and V61 live link
      tail;
\item \({\cal W}_{\rm even}\) is not charged to the vertical
      orientation-odd obstruction.
\end{enumerate}
\end{proposition}

\begin{proof}
Orientation reversal is an involution, so its symmetric and antisymmetric
projectors are complementary.  The anomaly projection selects exactly the
finite trace coefficients in Proposition
\ref{prop:v18v62-anomaly-panel}, proving item two.  The decomposition
(72.20) then leaves only the stated odd complement in the vertical bank.
The even projection is disjoint from that bank by construction.
\end{proof}

\begin{firewall}[Projection compatibility must be checked]
\label{fire:v18v62-projection-compatibility}
For the literal overlap projector, orientation reversal also acts on spin
matrices, link order, basepoint transport, and the gauge insertion.
One must prove that the actual vertical Ward trace intertwines this
involution before discarding \({\cal W}_{\rm even}\).  V62 proves the
finite representation calculation, not that intertwining.
\end{firewall}

\subsection{Executable certificate}

\begin{record}[Exact V62 verifier]
\label{rec:v18v62-verifier}
The script
\[
 \texttt{scripts/verify\_sm\_v62\_hypercharge\_word\_pilot.py}          \tag{72.22}
\]
uses integer arithmetic, checks (72.3)--(72.9), identifies the first
nonzero higher odd moment, and records the absolute seventh-moment tail.
Its SHA--256 is
\[
 \texttt{43D34599851A04D9816051E8E23D8DBFED95368C0DB3BBD939E361AEB949B5D5}.
                                                                  \tag{72.23}
\]
The canonical JSON payload
\[
 \texttt{artifacts/sm\_v62\_hypercharge\_word\_pilot.json}             \tag{72.24}
\]
has SHA--256
\[
 \texttt{4986E115A97C5E1D51EF7E4333A00F8E57C9E8B2C7A4027260024F09BBF43BB0}.
                                                                  \tag{72.25}
\]
\end{record}

\subsection{Acquisition and next bottleneck}

V62 executes the first exact representation-word pilot and falsifies the
overstrong coefficientwise-zero route.  The next upstream step is to pass
from the single common phase (72.12) to a multivariate commuting link
word.  Its fifth derivative is a symmetric charge tensor.  Contracting
that tensor with the actual orientation-paired lattice/spin coefficient
will decide whether the first nonanomalous term is killed by geometry or
must be carried as an \(a^5\) irrelevant remainder.

\begin{assessment}[V62 acquisition and boundary]
\label{ass:v18v62}
V62 provides an executed exact certificate for the generated SM charge
packet, proves the higher-word no-go, proves the fifth-order paired-phase
formula with a rational sign radius, and replaces blanket word
cancellation by the anomaly/even/higher-odd decomposition (72.20).
The overlap-specific orientation intertwining and lattice coefficient are
still uncomputed.
\end{assessment}

\begin{record}[V62 append record]
\label{rec:v18v62}
V62 is appended to the validated V61 whose installed SHA--256 was
\texttt{4E6D4A50B9F9417DB1A9750F1F6080C0034BF51A55E5F43F4D11FCFFBB746553}.
After installation only V62 is the active numeric SM note; the frozen
\texttt{V100\_f2} through \texttt{V100\_f17} archives are unchanged.
\end{record}

\section{V63: multivariate incidence tensors for the first surviving central word}
\label{sec:v18v63}

V62 proved that the orientation-odd central packet begins at fifth order,
but treated one common phase.  A lattice path sees a signed sum of link
phases.  V63 converts the whole finite path bank into symmetric incidence
tensors, proves exact fifth- and seventh-order criteria, and supplies an
executable checker.  This isolates the missing overlap datum: the
coefficient with which a path and its reverse enter after the spin trace.

\subsection{Path incidence}

Fix \(E\) oriented Abelian links with phase vector
\(\theta\in\mathbb R^E\).  A closed path word \(w\) has integer incidence
vector
\[
 q_w=(q_{w,e})_{e\in E}\in\mathbb Z^E,\qquad
 \Theta_w(\theta)=q_w\cdot\theta.                            \tag{73.1}
\]
Repeated traversal is recorded by the integer \(q_{w,e}\).  Let
\({\cal B}=\{(c_w,q_w)\}\) be a finite real coefficient bank and define
\[
 {\cal V}_{\cal B}(\theta)
 =\sum_{w\in{\cal B}}c_w\Phi(q_w\cdot\theta),                \tag{73.2}
\]
where \(\Phi\) is the exact one-generation packet (72.12).

For \(r\ge1\), put
\[
 P_r^{\cal B}(\theta)
 =\sum_{w\in{\cal B}}c_w(q_w\cdot\theta)^r,\qquad
 T_r^{\cal B}=\sum_{w\in{\cal B}}c_wq_w^{\otimes r}.         \tag{73.3}
\]

\begin{lemma}[Polynomial--tensor equivalence]
\label{lem:v18v63-tensor-equivalence}
For every \(r\),
\[
 P_r^{\cal B}(\theta)
 =T_r^{\cal B}[\theta,\ldots,\theta].                        \tag{73.4}
\]
Moreover,
\[
 P_r^{\cal B}\equiv0
 \quad\Longleftrightarrow\quad
 T_r^{\cal B}=0.                                             \tag{73.5}
\]
\end{lemma}

\begin{proof}
Multilinearity gives (73.4).  The tensor in (73.3) is symmetric.  If its
diagonal polynomial vanishes, the real polarization identity reconstructs
every mixed value from diagonal values:
\[
 T[v_1,\ldots,v_r]
 =\frac1{2^rr!}
 \sum_{\varepsilon_j=\pm1}
 \left(\prod_{j=1}^r\varepsilon_j\right)
 P\!\left(\sum_{j=1}^r\varepsilon_jv_j\right).              \tag{73.6}
\]
Thus \(P\equiv0\) implies \(T=0\); the reverse implication follows from
(73.4).
\end{proof}

\subsection{Fifth-order tensor and rational tail}

Define the absolute incidence stock
\[
 B_r({\cal B})
 =\sum_{w\in{\cal B}}|c_w|\,\|q_w\|_1^r.                    \tag{73.7}
\]

\begin{theorem}[Multivariate fifth-order decomposition]
\label{thm:v18v63-fifth-tensor}
For every \(\theta\in\mathbb R^E\),
\[
\boxed{
 {\cal V}_{\cal B}(\theta)
 =36P_5^{\cal B}(\theta)+{\cal E}_7^{\cal B}(\theta),}       \tag{73.8}
\]
where
\[
 |{\cal E}_7^{\cal B}(\theta)|
 \le\frac{27821}{420}B_7({\cal B})\|\theta\|_\infty^7.       \tag{73.9}
\]
The fifth-order term vanishes for every phase vector if and only if
\[
\boxed{T_5^{\cal B}=0.}                                     \tag{73.10}
\]
\end{theorem}

\begin{proof}
Apply the V62 formula (72.15) with scalar argument
\(q_w\cdot\theta\), multiply by \(c_w\), and sum.  Since
\[
 |q_w\cdot\theta|
 \le\|q_w\|_1\|\theta\|_\infty,                              \tag{73.11}
\]
the sum of the absolute remainders is (73.9).  Lemma
\ref{lem:v18v63-tensor-equivalence} gives the equivalence (73.10).
\end{proof}

\begin{corollary}[A bankwise sign test]
\label{cor:v18v63-bank-sign}
If for some \(\theta\),
\[
 36|P_5^{\cal B}(\theta)|
 >
 \frac{27821}{420}B_7({\cal B})\|\theta\|_\infty^7,          \tag{73.12}
\]
then \({\cal V}_{\cal B}(\theta)\ne0\) and has the sign of
\(P_5^{\cal B}(\theta)\).
\end{corollary}

\begin{proof}
The tail in (73.8) is strictly smaller than its principal term.
\end{proof}

\subsection{Seventh-order refinement}

The next signed hypercharge moment is
\[
 \frac{S_7}{7!}=\frac{226800}{5040}=45,                     \tag{73.13}
\]
and the sine sign makes its contribution negative.  The ninth absolute
moment is
\[
 A_9=10905036,\qquad
 \frac{A_9}{9!}=\frac{908753}{30240}.                       \tag{73.14}
\]

\begin{theorem}[Fifth--seventh expansion]
\label{thm:v18v63-seventh-tensor}
For every finite bank,
\[
\boxed{
 {\cal V}_{\cal B}(\theta)
 =
 36P_5^{\cal B}(\theta)
 -45P_7^{\cal B}(\theta)
 +{\cal E}_9^{\cal B}(\theta),}                             \tag{73.15}
\]
with
\[
 |{\cal E}_9^{\cal B}(\theta)|
 \le
 \frac{908753}{30240}
 B_9({\cal B})\|\theta\|_\infty^9.                          \tag{73.16}
\]
If \(T_5^{\cal B}=0\), the first displayed candidate is the seventh-order
tensor \(-45T_7^{\cal B}\).
\end{theorem}

\begin{proof}
Taylor's theorem through degree eight gives
\[
 \sin x=x-\frac{x^3}{3!}+\frac{x^5}{5!}
 -\frac{x^7}{7!}+r_9(x),\qquad
 |r_9(x)|\le\frac{|x|^9}{9!}.                               \tag{73.17}
\]
Sum over the charge packet.  The \(S_1,S_3\) terms vanish, (72.4) gives
the coefficient \(36\), and (73.13) gives \(-45\).  Sum the absolute
ninth-order remainder and use (73.11), proving (73.16).
\end{proof}

\subsection{Orientation reversal}

Let \(\bar w\) be the reversed path, so
\[
 q_{\bar w}=-q_w,\qquad
 \Phi(q_{\bar w}\cdot\theta)=-\Phi(q_w\cdot\theta).          \tag{73.18}
\]

\begin{proposition}[The reversal coefficient decides cancellation]
\label{prop:v18v63-reversal}
For a pair \(w,\bar w\):
\begin{enumerate}[leftmargin=2.2em]
\item if \(c_{\bar w}=c_w\), its contribution to
      \({\cal V}_{\cal B}\) vanishes identically, and every odd incidence
      tensor cancels;
\item if \(c_{\bar w}=-c_w\), its contribution is
      \(2c_w\Phi(q_w\cdot\theta)\), and every odd incidence tensor doubles.
\end{enumerate}
\end{proposition}

\begin{proof}
Substitute (73.18).  Since
\((-q_w)^{\otimes r}=-q_w^{\otimes r}\) for odd \(r\), the tensor statements
follow as well.
\end{proof}

\begin{firewall}[The spin trace supplies the coefficient sign]
\label{fire:v18v63-spin-sign}
Path reversal changes more than Abelian incidence: it reverses the ordered
Wilson hopping matrices, acts on gamma matrices, and moves the gauge
insertion.  V63 does not determine whether the literal overlap spin/lattice
coefficient satisfies \(c_{\bar w}=c_w\) or
\(c_{\bar w}=-c_w\).  Proposition \ref{prop:v18v63-reversal} proves that
this sign cannot be omitted.
\end{firewall}

\subsection{Executed tensor panels}

The exact verifier uses the plaquette incidence
\[
 q=(1,1,-1,-1).                                             \tag{73.19}
\]
It evaluates all monomial coefficients of \(P_5\) and \(P_7\) for three
banks:
\[
\begin{array}{c|c|cc|cc}
\hbox{bank}&(c,q)\hbox{ data}&T_5=0&T_7=0&B_7&B_9\\ \hline
\hbox{single}&(1,q)&\hbox{no}&\hbox{no}&16384&262144\\
\hbox{balanced}&(1,q),(1,-q)&\hbox{yes}&\hbox{yes}&32768&524288\\
\hbox{antisymmetric}&(1,q),(-1,-q)&\hbox{no}&\hbox{no}&32768&524288 .
\end{array}                                                   \tag{73.20}
\]

\begin{record}[Exact V63 verifier]
\label{rec:v18v63-verifier}
The script
\[
 \texttt{scripts/verify\_sm\_v63\_multivariate\_phase\_tensor.py}       \tag{73.21}
\]
has SHA--256
\[
 \texttt{58D1D126D7BA58AEBAB9039761BA11BEFAAC1EAB64BAFB540D58CA5C83F8010F}.
                                                                  \tag{73.22}
\]
Its canonical payload
\[
 \texttt{artifacts/sm\_v63\_multivariate\_phase\_tensor.json}          \tag{73.23}
\]
has SHA--256
\[
 \texttt{E76428AAC3C0DA46475463B20B1B70B88F9E2BFBB95ACC71D062DA7C284193A6}.
                                                                  \tag{73.24}
\]
The panels are structural tests of the tensor compiler; none is asserted
to be the literal overlap path bank.
\end{record}

\subsection{Cofinal irrelevant-word ledger}

For cutoff \(n\), let \({\cal B}_n\) be the actual orientation-odd bank and
let \(\eta_n\) bound the link phases.  Theorem
\ref{thm:v18v63-fifth-tensor} gives
\[
 \|{\cal V}_{{\cal B}_n}\|
 \le36\|T_{5,n}\|_{\rm diag}\eta_n^5
 +\frac{27821}{420}B_7({\cal B}_n)\eta_n^7,                 \tag{73.25}
\]
where
\[
 \|T_{5,n}\|_{\rm diag}
 =\sup_{\|\theta\|_\infty\le1}
 |T_{5,n}[\theta,\ldots,\theta]|.                           \tag{73.26}
\]

\begin{theorem}[Fifth-order cofinal sufficiency]
\label{thm:v18v63-cofinal-fifth}
If
\[
 \|T_{5,n}\|_{\rm diag}\eta_n^5\longrightarrow0,\qquad
 B_7({\cal B}_n)\eta_n^7\longrightarrow0,                   \tag{73.27}
\]
then the commuting central higher-word contribution vanishes.  If
\(T_{5,n}=0\) exactly, it is enough to replace the first condition by
\[
 \|T_{7,n}\|_{\rm diag}\eta_n^7\to0,\qquad
 B_9({\cal B}_n)\eta_n^9\to0.                               \tag{73.28}
\]
\end{theorem}

\begin{proof}
Equation (73.25) proves the first assertion.  Under exact fifth-tensor
cancellation, apply (73.15)--(73.16) and take the supremum on the
\(\eta_n\)-box.
\end{proof}

\begin{remark}[Smooth-field scaling]
\label{rem:v18v63-smooth-scaling}
If \(\eta_n\lesssim a_n\|A\|_\infty\) on a smooth chart, (73.27) is an
explicit fifth-order irrelevance test.  If a word traverses
\(\ell_{w,n}\) links, then
\(|q_w\cdot\theta|\le\ell_{w,n}\eta_n\); the actual small parameter is the
physical path length \(a_n\ell_{w,n}\|A\|_\infty\), consistent with V60.
\end{remark}

\begin{firewall}[Literal incidence bank remains absent]
\label{fire:v18v63-literal-bank}
The source material gives the overlap sign projector and abstract
finite-range surrogates, but not an emitted list of the actual
orientation-paired coefficients \(c_w\).  Therefore V63 has not evaluated
\(T_{5,n}\) or the stocks in (73.27) for the literal regulator.
\end{firewall}

\subsection{Acquisition and next bottleneck}

V63 reduces the first surviving central word to a finite symmetric tensor.
The next step is to determine the reversal coefficient from the Wilson
spin kernel.  An antiunitary or reflection intertwiner for \(H_W\), its
sign polynomial, and the gauge insertion would fix the relation between
\(c_w\) and \(c_{\bar w}\) before the tensor bank is generated.

\begin{assessment}[V63 acquisition and boundary]
\label{ass:v18v63}
V63 proves the multivariate fifth- and seventh-order incidence-tensor
expansions, exact reversal criteria, and cofinal irrelevance ledger.  Its
executed compiler verifies the tensor arithmetic on three finite panels.
The literal overlap spin/reflection intertwiner is the remaining input.
\end{assessment}

\begin{record}[V63 append record]
\label{rec:v18v63}
V63 is appended to the validated V62 whose installed SHA--256 was
\texttt{F920A525B81A487B73D92F5640BA28748B92CE09055B81CB4AC10F61261CB518}.
After installation only V63 is the active numeric SM note; the frozen
\texttt{V100\_f2} through \texttt{V100\_f17} archives are unchanged.
\end{record}

\section{V64: antiunitary sign transport and the exact reversal-parity ledger}
\label{sec:v18v64}

V63 isolated the coefficient relating a path to its reverse.  The attached
V14 source proves a polar criterion for positivity of a chronological
boundary crossing, but explicitly leaves the selected Wilson matrix and
physical reflection unpopulated.  That criterion does not determine the
charge/path-reversal parity of the vertical Ward trace.  V64 supplies the
missing functional-calculus theorem and a quantitative defect version.

\subsection{Antiunitary functional calculus}

Let \({\cal H}\) be finite dimensional and let \(C:{\cal H}\to{\cal H}\)
be antiunitary.  For a linear operator \(A\), write
\[
 {\cal C}(A)=CAC^{-1}.                                      \tag{74.1}
\]
This map is conjugate linear on scalars, multiplicative on operators, and
preserves adjoints and operator norms.

\begin{lemma}[Real functional calculus intertwines]
\label{lem:v18v64-functional-calculus}
Let \(H=H^*\) be gapped and let \(f\) be real valued on its spectrum.  Then
\[
 {\cal C}(f(H))=f({\cal C}(H)).                              \tag{74.2}
\]
If \(f\) is \(C^1\) on a neighborhood of the spectrum and \(Q=Q^*\), then
\[
 {\cal C}(Df_H[Q])
 =Df_{{\cal C}(H)}[{\cal C}(Q)].                            \tag{74.3}
\]
\end{lemma}

\begin{proof}
Diagonalize \(H=\sum_\lambda\lambda E_\lambda\).  Antiunitarity sends this
to
\({\cal C}(H)=\sum_\lambda\lambda{\cal C}(E_\lambda)\), because the
eigenvalues are real.  This proves (74.2).  In eigenbases, the derivative
is
\[
 Df_H[Q]
 =\sum_{\lambda,\mu}
 f^{[1]}(\lambda,\mu)E_\lambda QE_\mu.                      \tag{74.4}
\]
The divided differences are real, so applying \({\cal C}\) term by term
gives (74.3).
\end{proof}

\begin{corollary}[Sign projector transport]
\label{cor:v18v64-sign-transport}
For \(P_-(H)=1_{(-\infty,0)}(H)\),
\[
 {\cal C}(P_-(H))=P_-({\cal C}(H)),\qquad
 {\cal C}(DP_{-,H}[Q])
 =DP_{-,{\cal C}(H)}[{\cal C}(Q)].                          \tag{74.5}
\]
Every real sign polynomial \(p_m\) and its V60 derivative obey the same
identities.
\end{corollary}

\begin{proof}
The spectral gap makes the indicator \(C^1\) on a disconnected
neighborhood of the spectrum.  Apply Lemma
\ref{lem:v18v64-functional-calculus}.  A real polynomial is the elementary
special case.
\end{proof}

\subsection{Parity of the vertical Ward trace}

Let \(\vartheta\mapsto-\vartheta\) reverse the chosen link-phase chart.
Assume
\[
\begin{split}
 H(-\vartheta)&={\cal C}(H(\vartheta)),\\
 Q(-\vartheta)&=\epsilon_Q{\cal C}(Q(\vartheta)),\\
 \Omega&=\epsilon_\Omega{\cal C}(\Omega),
 \qquad \epsilon_Q,\epsilon_\Omega\in\{1,-1\}.              \tag{74.6}
\end{split}
\]
Here \(Q\) is the tangent insertion and \(\Omega\) is anti-Hermitian.  Put
\[
 {\cal W}(\vartheta)
 =-i\operatorname{Tr}
 \bigl(DP_{-,H(\vartheta)}[Q(\vartheta)]\Omega\bigr).         \tag{74.7}
\]

\begin{theorem}[Exact antiunitary parity]
\label{thm:v18v64-parity}
The number \({\cal W}(\vartheta)\) is real and
\[
\boxed{
 {\cal W}(-\vartheta)
 =-\epsilon_Q\epsilon_\Omega{\cal W}(\vartheta).}            \tag{74.8}
\]
In particular, if both the differentiated link insertion and the
anti-Hermitian central generator are odd under \(C\), then
\[
 \epsilon_Q=\epsilon_\Omega=-1,\qquad
 {\cal W}(-\vartheta)=-{\cal W}(\vartheta).                  \tag{74.9}
\]
\end{theorem}

\begin{proof}
The operator \(DP_{-,H}[Q]\) is Hermitian.  For Hermitian \(A\) and
anti-Hermitian \(\Omega\),
\[
 \overline{\operatorname{Tr}(A\Omega)}
 =\operatorname{Tr}(\Omega^*A^*)
 =-\operatorname{Tr}(\Omega A)
 =-\operatorname{Tr}(A\Omega),                              \tag{74.10}
\]
so (74.7) is real.  Antiunitary conjugation satisfies
\[
 \operatorname{Tr}({\cal C}(B))
 =\overline{\operatorname{Tr}(B)}.                          \tag{74.11}
\]
By Corollary \ref{cor:v18v64-sign-transport} and (74.6),
\[
\begin{split}
 &\operatorname{Tr}
 \bigl(DP_{-,H(-\vartheta)}[Q(-\vartheta)]\Omega\bigr)\\
 &\quad=
 \epsilon_Q\epsilon_\Omega
 \overline{
 \operatorname{Tr}
 \bigl(DP_{-,H(\vartheta)}[Q(\vartheta)]\Omega\bigr)}.
                                                                  \tag{74.12}
\end{split}
\]
The trace on the right is purely imaginary by (74.10).  Multiplication by
\(-i\) therefore gives (74.8).  The last statement is substitution.
\end{proof}

\begin{lemma}[Derivative insertion parity]
\label{lem:v18v64-derivative-parity}
Suppose the first identity in (74.6) holds and
\[
 Q_e(\vartheta)=\partial_{\vartheta_e}H(\vartheta).
                                                                  \tag{74.13}
\]
Then
\[
 Q_e(-\vartheta)=-{\cal C}(Q_e(\vartheta)),                  \tag{74.14}
\]
so \(\epsilon_Q=-1\).
\end{lemma}

\begin{proof}
Differentiate
\(H(-\vartheta)={\cal C}(H(\vartheta))\) with respect to
\(\vartheta_e\).  The chain rule contributes the minus sign on the left.
\end{proof}

\begin{remark}[Central generator parity]
\label{rem:v18v64-central-parity}
For a \(U(1)\) representation with anti-Hermitian generator
\(\Omega_y=iyI\), ordinary complex conjugation gives
\({\cal C}(\Omega_y)=-\Omega_y\).  Thus the natural central link derivative
and generator have the parity pair in (74.9), provided the full Wilson spin
kernel has the first intertwining identity in (74.6).
\end{remark}

\subsection{Laurent coefficient reversal}

Suppose a finite principal trace has a Laurent expansion
\[
 {\cal W}(\vartheta)=
 \sum_{q\in{\cal Q}}a_qe^{iq\cdot\vartheta},\qquad
 {\cal Q}=-{\cal Q}.                                       \tag{74.15}
\]

\begin{proposition}[Coefficient relations]
\label{prop:v18v64-laurent}
If \({\cal W}\) is real, then
\[
 a_{-q}=\overline{a_q}.                                     \tag{74.16}
\]
If it also has parity
\({\cal W}(-\vartheta)=\epsilon_{\cal W}{\cal W}(\vartheta)\),
then
\[
 a_{-q}=\epsilon_{\cal W}a_q.                               \tag{74.17}
\]
On the odd branch \(\epsilon_{\cal W}=-1\), the nonzero coefficients are
purely imaginary and every orbit \(\{q,-q\}\) contributes one real sine
word.  The incidence compiler should therefore retain one representative
per reversal orbit.
\end{proposition}

\begin{proof}
Complex conjugation of (74.15) and uniqueness of Laurent coefficients give
(74.16).  Substitution of \(-\vartheta\) and reindexing \(q\mapsto-q\)
give (74.17).  If the parity is odd, the two identities imply
\(\overline{a_q}=-a_q\).  Pairing the two exponential terms converts them
to a real multiple of \(\sin(q\cdot\vartheta)\).
\end{proof}

\begin{assessment}[Resolution of the V63 double-counting ambiguity]
\label{ass:v18v64-canonical-pair}
Once (74.9) is verified, the correct V63 bank is the quotient by
\(q\sim-q\), with the sine coefficient read from the imaginary Laurent
coefficient.  Adding both representatives again can falsely cancel or
double the word.  The antiunitary theorem fixes the canonical bookkeeping;
it does not determine the numerical sine coefficient.
\end{assessment}

\subsection{Approximate polynomial intertwining}

For a real polynomial
\[
 p(x)=\sum_{k=0}^db_kx^k
                                                                  \tag{74.18}
\]
and a spectral-radius bound \(M\), define
\[
\begin{split}
 L_1(p;M)&=\sum_{k=1}^dk|b_k|M^{k-1},\\
 L_2(p;M)&=\sum_{k=2}^dk(k-1)|b_k|M^{k-2}.                  \tag{74.19}
\end{split}
\]
Let \(P_p(H)=(I-p(H))/2\).

\begin{lemma}[Polynomial derivative stability]
\label{lem:v18v64-polynomial-stability}
If \(\|H\|,\|\widetilde H\|\le M\) and
\(\|Q\|,\|\widetilde Q\|\le Q_*\), then
\[
\boxed{
 \|DP_{p,H}[Q]-DP_{p,\widetilde H}[\widetilde Q]\|
 \le\frac12\left(
 L_1(p;M)\|Q-\widetilde Q\|
 +L_2(p;M)Q_*\|H-\widetilde H\|\right).}                    \tag{74.20}
\]
\end{lemma}

\begin{proof}
The polynomial derivative formula (70.15) gives
\[
 \|Dp_H[Q]\|\le L_1(p;M)\|Q\|.                              \tag{74.21}
\]
First replace \(Q\) by \(\widetilde Q\), producing the first term.
For each degree \(k\), its derivative contains \(k\) words and each word
contains \(k-1\) copies of \(H\).  Replacing those copies one at a time gives
at most
\(k(k-1)M^{k-2}Q_*\|H-\widetilde H\|\).  Sum with \(|b_k|\) and include
the factor \(1/2\).
\end{proof}

Let reflected data \(H^R,Q^R,\Omega^R\) obey
\[
\begin{split}
 \delta_H&=\|H^R-{\cal C}(H)\|,\\
 \delta_Q&=\|Q^R-\epsilon_Q{\cal C}(Q)\|,\\
 \delta_\Omega&=\|\Omega^R-\epsilon_\Omega{\cal C}(\Omega)\|,
                                                                  \tag{74.22}
\end{split}
\]
with norms bounded by \(M,Q_*,\Omega_*\) as above.

\begin{theorem}[Quantitative parity residual]
\label{thm:v18v64-parity-residual}
For
\[
 {\cal W}_p(H,Q,\Omega)
 =-i\operatorname{Tr}(DP_{p,H}[Q]\Omega),                   \tag{74.23}
\]
one has
\[
\boxed{
\begin{split}
 &|{\cal W}_p(H^R,Q^R,\Omega^R)
 +\epsilon_Q\epsilon_\Omega
 {\cal W}_p(H,Q,\Omega)|\\
 &\quad\le\frac{N}{2}\left[
 \Omega_*\bigl(L_1\delta_Q+L_2Q_*\delta_H\bigr)
 +L_1Q_*\delta_\Omega\right],
\end{split}}                                                \tag{74.24}
\]
where \(N=\dim{\cal H}\) and \(L_j=L_j(p;M)\).
\end{theorem}

\begin{proof}
Compare \(DP_{p,H^R}[Q^R]\) with
\(\epsilon_Q{\cal C}(DP_{p,H}[Q])\) using Lemmas
\ref{lem:v18v64-functional-calculus} and
\ref{lem:v18v64-polynomial-stability}.  This costs
\(\frac12(L_1\delta_Q+L_2Q_*\delta_H)\).  Replacing
\(\Omega^R\) by \(\epsilon_\Omega{\cal C}(\Omega)\) costs at most
\(\frac12L_1Q_*\delta_\Omega\) by (74.21).  Multiply by the remaining
operator norm and use
\(|\operatorname{Tr}A|\le N\|A\|\).  The exact antiunitary trace relation
supplies the sign on the left.
\end{proof}

\begin{corollary}[Adding the V60 sign tail]
\label{cor:v18v64-sign-tail}
If the exact sign derivative is approximated at both reflected points by
the same V60 polynomial with Ward tail at most \(S_m\), then its parity
residual is bounded by the right side of (74.24) plus \(2S_m\).
\end{corollary}

\begin{proof}
Insert the polynomial trace at each endpoint and apply the triangle
inequality once to the two approximation errors.
\end{proof}

\begin{firewall}[Dimension loss]
\label{fire:v18v64-dimension-loss}
The factor \(N\) in (74.24) is a crude full-matrix trace bound.  A local
gauge insertion should replace it by a support rank or trace-norm stock.
Without that localization, (74.24) is a finite-cutoff test but may be
useless cofinally.
\end{firewall}

\subsection{Combined reversal and incidence capsule}

\begin{theorem}[Conditional odd-bank reduction]
\label{thm:v18v64-odd-bank-capsule}
Suppose along a cofinal sequence:
\begin{enumerate}[leftmargin=2.2em]
\item the Wilson kernels, tangent insertions, and gauge generators satisfy
      (74.6) exactly, or the right side of (74.24) plus \(2S_m\) tends to
      zero;
\item the Laurent path bank is finite at each cutoff and closed under the
      induced reversal;
\item the canonical sine bank satisfies the V63 fifth-order conditions
      (73.27), or the stronger cancellation and seventh-order conditions
      (73.28);
\item all V57--V61 topology, gauge-loop, locality, atlas, and transport
      rows hold.
\end{enumerate}
Then the orientation-even leakage and the higher central
orientation-odd contribution both vanish in the continuum vertical Ward
residual.
\end{theorem}

\begin{proof}
Theorem \ref{thm:v18v64-parity} eliminates the even part exactly; in the
approximate case, Theorem \ref{thm:v18v64-parity-residual} bounds it by a
quantity tending to zero.  Proposition \ref{prop:v18v64-laurent}
canonicalizes the remaining sine bank.  V63 makes its fifth/seventh
remainder vanish under item three.  Item four transports the resulting
vertical-curvature bound through the determinant-line and Ward collars.
\end{proof}

\begin{firewall}[Literal reflection remains unpopulated]
\label{fire:v18v64-literal-reflection}
The attached source explicitly states that the selected RPESM Wilson
matrix, split, physical reflection, one-letter bank, and line orientation
remain unevaluated.  V64 therefore cannot assert (74.6) for the literal
history.  The V14 polar boundary-reflection residual and the V64
antiunitary path-reversal residual are separate tests and neither implies
the other.
\end{firewall}

\subsection{Acquisition and next bottleneck}

V64 supplies the exact parity theorem and its approximate residual.
The next concrete acquisition is one declared Wilson kernel card:
its matrix entries, link-phase chart, antiunitary \(C\), tangent insertion,
and central generator.  Those data make every term in (74.22) executable
and determine whether the V63 sine bank is physically present.

\begin{assessment}[V64 acquisition and boundary]
\label{ass:v18v64}
V64 proves antiunitary transport for the sign projector and its derivative,
the exact Ward parity sign, Laurent reversal relations, a canonical
one-representative sine bank, and the quantitative polynomial residual
(74.24).  It also proves that the source's polar crossing reflection does
not populate the missing path-reversal card.
\end{assessment}

\begin{record}[V64 append record]
\label{rec:v18v64}
V64 is appended to the validated V63 whose installed SHA--256 was
\texttt{7CFAE645BDD92D93A7840FA0A16A4ECD47A3951D2079748EAEC543DA4BE27304}.
After installation only V64 is the active numeric SM note; the frozen
\texttt{V100\_f2} through \texttt{V100\_f17} archives are unchanged.
\end{record}

\section{V65: exact balanced Wilson reversal cards and a non-identifiability theorem}
\label{sec:v18v65}

V64 proved that the Ward parity depends on both the reflected tangent
insertion and the anti-Hermitian gauge generator.  V65 now gives an exact
balanced Wilson collar in which all kernel, gap, projector, and crossing
data are held fixed while the generator is changed.  The two resulting
Ward traces have opposite reversal parity.  This proves that the missing
physical generator cannot be inferred from the spectral or polar data in
the attached V14 card.

\subsection{A one-link balanced collar}

Let \(m,t>0\), \(E=(m^2+t^2)^{1/2}\), and
\[
 H(\theta)=
 \begin{pmatrix}
 -m&-te^{-i\theta}\\
 -te^{i\theta}&m
 \end{pmatrix}.                                             \tag{75.1}
\]
This is the scalar \(B=te^{-i\theta}\) instance of the attached balanced
Wilson boundary collar.  Let \(C\) be ordinary complex conjugation.

\begin{lemma}[Exact sign and reversal data]
\label{lem:v18v65-balanced-card}
The family (75.1) satisfies
\[
\begin{split}
 H(\theta)^2&=E^2I,\qquad
 H(-\theta)=CH(\theta)C^{-1},\\
 P_-(\theta)&=\frac12\left(I-\frac{H(\theta)}E\right),\\
 Q(\theta):=\partial_\theta H(\theta),\qquad
 Q(-\theta)&=-CQ(\theta)C^{-1}.                             \tag{75.2}
\end{split}
\]
Its spectrum is \(\{-E,E\}\), its gap is \(E\), and the singular value of
the off-diagonal sign crossing is \(t/E\).
\end{lemma}

\begin{proof}
Direct multiplication cancels the off-diagonal terms and gives
\(H^2=(m^2+t^2)I\).  Functional calculus therefore gives
\(\operatorname{sign}H=H/E\) and the displayed projector.  Complex
conjugation replaces \(\theta\) by \(-\theta\), and differentiation gives
the tangent identity.  The off-diagonal entries of
\(\operatorname{sign}H\) have modulus \(t/E\).
\end{proof}

\begin{corollary}[Exact projector derivative]
\label{cor:v18v65-projector-derivative}
Since \(E\) is independent of \(\theta\),
\[
 DP_{-,H(\theta)}[Q(\theta)]
 =\partial_\theta P_-(\theta)
 =-\frac{Q(\theta)}{2E}.                                    \tag{75.3}
\]
\end{corollary}

\begin{proof}
Differentiate the projector formula in (75.2).
\end{proof}

\subsection{Two generator completions}

Let
\[
 \Omega_{\rm o}=i\sigma_1
 =\begin{pmatrix}0&i\\i&0\end{pmatrix},\qquad
 \Omega_{\rm e}=i\sigma_2
 =\begin{pmatrix}0&1\\-1&0\end{pmatrix}.                    \tag{75.4}
\]
Both are anti-Hermitian and have operator norm one, but
\[
 C\Omega_{\rm o}C^{-1}=-\Omega_{\rm o},\qquad
 C\Omega_{\rm e}C^{-1}=+\Omega_{\rm e}.                     \tag{75.5}
\]
Define
\[
 {\cal W}_{j}(\theta)
 =-i\operatorname{Tr}
 \left(DP_{-,H(\theta)}[Q(\theta)]\Omega_j\right),
 \qquad j\in\{{\rm o},{\rm e}\}.                            \tag{75.6}
\]

\begin{theorem}[Odd and even Ward completions of one kernel]
\label{thm:v18v65-two-completions}
For the common balanced kernel (75.1),
\[
\boxed{
 {\cal W}_{\rm o}(\theta)
 =-\frac{t}{E}\sin\theta,\qquad
 {\cal W}_{\rm e}(\theta)
 = \frac{t}{E}\cos\theta.}                                  \tag{75.7}
\]
Consequently
\[
 {\cal W}_{\rm o}(-\theta)=-{\cal W}_{\rm o}(\theta),\qquad
 {\cal W}_{\rm e}(-\theta)= {\cal W}_{\rm e}(\theta).        \tag{75.8}
\]
\end{theorem}

\begin{proof}
Writing \(Q\) in Pauli form gives
\[
 Q(\theta)=t\sin\theta\,\sigma_1-t\cos\theta\,\sigma_2.      \tag{75.9}
\]
Equations (75.3) and
\(\operatorname{Tr}(\sigma_j\sigma_k)=2\delta_{jk}\) yield
\[
\begin{split}
 {\cal W}_{\rm o}
 &=-\frac1{2E}\operatorname{Tr}(Q\sigma_1)
 =-\frac{t}{E}\sin\theta,\\
 {\cal W}_{\rm e}
 &=-\frac1{2E}\operatorname{Tr}(Q\sigma_2)
 = \frac{t}{E}\cos\theta.                                  \tag{75.10}
\end{split}
\]
The parity relations follow.  They also agree with the V64 rule because
\(\epsilon_Q=-1\), while
\(\epsilon_{\Omega_{\rm o}}=-1\) and
\(\epsilon_{\Omega_{\rm e}}=+1\).
\end{proof}

\begin{countertheorem}[Spectral and polar data do not determine Ward parity]
\label{cth:v18v65-nonidentifiability}
The following data do not determine whether the vertical Ward trace is odd
or even under link reversal:
\[
 H(\theta),\quad\operatorname{spec}H(\theta),\quad
 P_-(\theta),\quad \|T_+\operatorname{sign}H(\theta)T_-\|,
 \quad C.                                                    \tag{75.11}
\]
The physical anti-Hermitian generator, including its \(C\)-parity, is an
independent required datum.
\end{countertheorem}

\begin{proof}
The two cards in Theorem \ref{thm:v18v65-two-completions} have identical
data (75.11).  Their generators have equal norm, but one Ward trace is odd
and the other is even and nonzero at \(\theta=0\).  Hence no function of
the common data can determine the parity.
\end{proof}

\begin{assessment}[Relation to the attached polar criterion]
\label{ass:v18v65-polar-relation}
The V14 polar criterion determines whether the represented chronological
crossing is positive after a physical time reflection is supplied.  The
countertheorem above concerns the charge/path-reversal parity of a
projector derivative paired with a gauge generator.  Equal gap and crossing
singular values leave this parity undetermined, so neither test can replace
the other.
\end{assessment}

\subsection{Exact rational instance}

Choose
\[
 m=3,\qquad t=4,\qquad E=5,\qquad
 \cos\theta=\frac35,\qquad\sin\theta=\frac45.               \tag{75.12}
\]
Then
\[
 H=
 \begin{pmatrix}
 -3&-\frac{12}{5}+\frac{16}{5}i\\
 -\frac{12}{5}-\frac{16}{5}i&3
 \end{pmatrix},\qquad H^2=25I,                              \tag{75.13}
\]
and
\[
 P_-=
 \begin{pmatrix}
 \frac45&\frac6{25}-\frac8{25}i\\
 \frac6{25}+\frac8{25}i&\frac15
 \end{pmatrix},\qquad P_-^2=P_-.                            \tag{75.14}
\]
The two exact Ward readings are
\[
\begin{array}{c|cc}
 &\theta&-\theta\\ \hline
 \Omega_{\rm o}&-16/25&16/25\\
 \Omega_{\rm e}&12/25&12/25 .
\end{array}                                                   \tag{75.15}
\]

\begin{record}[Exact V65 verifier]
\label{rec:v18v65-verifier}
The Gaussian-rational script
\[
 \texttt{scripts/verify\_sm\_v65\_balanced\_reversal\_cards.py}         \tag{75.16}
\]
checks (75.13)--(75.15), the projector identity, both generator parities,
and the common crossing value \(4/5\).  Its SHA--256 is
\[
 \texttt{1CE4E606777255352A80CC47F9FBCD8213397C6D2192568ACE89F46B39E59E82}.
                                                                  \tag{75.17}
\]
The canonical payload
\[
 \texttt{artifacts/sm\_v65\_balanced\_reversal\_cards.json}             \tag{75.18}
\]
has SHA--256
\[
 \texttt{DB55EAC024F6F244CE759600FEF1D4F88B52D991878CFC6CF93C0DC4662DE3B4}.
                                                                  \tag{75.19}
\]
\end{record}

\subsection{Minimal literal reversal card}

\begin{proposition}[Sufficient finite data panel]
\label{prop:v18v65-minimal-card}
At one cutoff and one link-phase chart, the V63--V64 reversal decision is
computable from the finite tuple
\[
\boxed{
 {\mathfrak R}
 =(H(\vartheta),\,C,\,Q_e(\vartheta),\,\Omega(\omega),\,
   \hbox{trace convention},\,\hbox{path coefficient bank}).}          \tag{75.20}
\]
The tuple must specify the same carrier and basis identifications for every
entry.  Given these data, one can:
\begin{enumerate}[leftmargin=2.2em]
\item verify the three intertwining relations (74.6);
\item compute the exact or polynomial Ward trace (74.7);
\item extract the Laurent reversal coefficients (74.16)--(74.17);
\item populate the incidence tensors and tails of V63.
\end{enumerate}
\end{proposition}

\begin{proof}
All four operations are finite matrix arithmetic and finite polynomial or
Laurent coefficient extraction once the entries of (75.20) are supplied.
The shared carrier and basis convention make the products and
antiunitary conjugations well typed.
\end{proof}

\begin{firewall}[The source does not populate (75.20)]
\label{fire:v18v65-source-missing}
The attached SM programme states that the selected RPESM Wilson matrix,
split, physical reflection, one-letter bank, and line orientation remain
unevaluated.  It also does not give a literal matrix for the generator in
the same boundary basis.  The exact benchmark proves what must be read; it
does not substitute its \(2\times2\) matrices for the physical card.
\end{firewall}

\subsection{Checkpoint and next step}

V65 reaches the next five-version boundary.  The active Yang--Mills and
Navier--Stokes notes must be inspected before V66.  After that audit, the
next upstream task is to determine whether an existing same-history source
or formal module actually contains all entries of (75.20).  If none does,
the programme should develop bounds conditional on a measured reversal
residual rather than infer a parity from incomplete spectral data.

\begin{assessment}[V65 acquisition and boundary]
\label{ass:v18v65}
V65 proves by exact balanced Wilson cards that common kernel, gap,
projector, antiunitary, and crossing data do not determine Ward reversal
parity without the physical generator.  It supplies an executable rational
witness and the minimal finite data tuple (75.20).  The literal tuple
remains unpopulated.
\end{assessment}

\begin{record}[V65 append record]
\label{rec:v18v65}
V65 is appended to the validated V64 whose installed SHA--256 was
\texttt{8C6D8ECB3350AD2483803B112285F2558A5DE96067BF406C2AB315C92B8B2AA7}.
After installation only V65 is the active numeric SM note; the frozen
\texttt{V100\_f2} through \texttt{V100\_f17} archives are unchanged.
\end{record}

\section{V66: compiled used-image control of the reversal residual}
\label{sec:v18v66}

V65 reached a checkpoint after proving that the physical generator is an
independent datum.  The companion audit supplies a sharper quantitative
rule for incomplete cards: propagate only the perturbation images actually
seen by the final trace, and compile their signed sum before taking
absolute values.  V66 implements that rule through an exact path derivative
of the polynomial Ward functional.

\subsection{Mandatory V65 companion audit}

The active companion files were
\[
\begin{array}{c|c|c}
\hbox{file}&\hbox{bytes}&\hbox{SHA--256}\\ \hline
\texttt{Renewal\_Yang\_Mills\_Mass\_Gap\_Research\_Notes\_V10.tex}
&174700&
\texttt{341694A2E6AE928C7AF8AA30B335B9F2FFB85EBD6A614ADE0F6EFAF2D01D5B9D}\\
\texttt{Renewal\_NS\_Stability\_Research\_Notes\_V26.tex}
&278283&
\texttt{82C887F8C2C69E4F28FAD7C3DC8DE5B9099CB5A4BF431EBA1FFF3E91935978AD}.
\end{array}                                                   \tag{76.1}
\]
YM V10 proves with exact fractions that a frozen high-alias
preconditioner improves its first-order pilot by seven dyadic powers and
reduces the Neumann order at radius \(2^{-8}\) from \(47\) to \(21\).
Its next step replaces coarse product bounds by the exact images used by
the final factor.  The transferable point is not the numerical
preconditioner itself; it is the used-image principle.  NS V26 is
unchanged.  Its restriction to a used rank-one input gives the same
qualitative lesson, but no new SM estimate.

\subsection{The polynomial Ward functional}

Let \(p\) be a real polynomial, put
\[
 P_p(H)=\frac{I-p(H)}2,                                     \tag{76.2}
\]
and define
\[
 {\cal W}_p(H,Q,\Omega)
 =-i\operatorname{Tr}\bigl(DP_{p,H}[Q]\Omega\bigr),          \tag{76.3}
\]
where \(H,Q\) are Hermitian and \(\Omega\) is anti-Hermitian.

\begin{lemma}[Second polynomial derivative]
\label{lem:v18v66-second-derivative}
If \(p(x)=\sum_{k=0}^db_kx^k\), then
\[
 D^2p_H[A,Q]
 =
 \sum_{k=2}^db_k
 \sum_{\substack{0\le r,s\le k-1\\r\ne s}}
 {\cal M}_{k;r,s}(H;A,Q),                                   \tag{76.4}
\]
where \({\cal M}_{k;r,s}\) is the word of length \(k\) obtained by placing
\(A\) and \(Q\) in the ordered slots \(r,s\) and \(H\) in every other
slot.  In particular the degree-\(k\) term contains exactly \(k(k-1)\)
ordered words, and
\[
 D^2p_H[A,Q]=D^2p_H[Q,A].                                   \tag{76.5}
\]
\end{lemma}

\begin{proof}
Differentiate the monomial derivative (70.15) in the direction \(A\).
Each of its \(k\) words contains \(k-1\) copies of \(H\); differentiating
one of them places \(A\) in one of the remaining ordered slots.  This
produces every ordered pair of distinct slots once.  Interchanging the
labels of the two distinguished slots proves symmetry.
\end{proof}

\begin{lemma}[Exact path derivative]
\label{lem:v18v66-path-derivative}
For \(C^1\) affine paths \(H_s,Q_s,\Omega_s\),
\[
\boxed{
\begin{split}
 \frac d{ds}{\cal W}_p(H_s,Q_s,\Omega_s)
 =-i\operatorname{Tr}\bigl(&
 D^2P_{p,H_s}[\dot H_s,Q_s]\Omega_s\\
 &+DP_{p,H_s}[\dot Q_s]\Omega_s
 +DP_{p,H_s}[Q_s]\dot\Omega_s\bigr).
\end{split}}                                                \tag{76.6}
\]
\end{lemma}

\begin{proof}
Differentiate (76.3).  The derivative of
\(DP_{p,H_s}[Q_s]\) is
\[
 D^2P_{p,H_s}[\dot H_s,Q_s]
 +DP_{p,H_s}[\dot Q_s],                                     \tag{76.7}
\]
and product differentiation supplies the last term in (76.6).
\end{proof}

\subsection{Ideal and measured reflected cards}

For sector \(a\), let the ideal reflected tuple be
\[
\begin{split}
 H_{a,0}^R&={\cal C}_a(H_a),\\
 Q_{a,0}^R&=\epsilon_Q{\cal C}_a(Q_a),\\
 \Omega_{a,0}^R&=\epsilon_\Omega{\cal C}_a(\Omega_a),        \tag{76.8}
\end{split}
\]
and let the measured tuple be
\[
 H_{a,1}^R=H_{a,0}^R+\Delta H_a,\quad
 Q_{a,1}^R=Q_{a,0}^R+\Delta Q_a,\quad
 \Omega_{a,1}^R=\Omega_{a,0}^R+\Delta\Omega_a.              \tag{76.9}
\]
Interpolate affinely in \(s\in[0,1]\).  Assume the two parities are common
to the compiled sectors; sectors with different parities are first split
into separate banks.

Define the compiled used-image integrand
\[
\boxed{
\begin{split}
 {\cal I}_p(s)=-i\sum_a\sigma_a\operatorname{Tr}\bigl(&
 D^2P_{p,H_{a,s}^R}[\Delta H_a,Q_{a,s}^R]\Omega_{a,s}^R\\
 &+DP_{p,H_{a,s}^R}[\Delta Q_a]\Omega_{a,s}^R\\
 &+DP_{p,H_{a,s}^R}[Q_{a,s}^R]\Delta\Omega_a\bigr).
\end{split}}                                                \tag{76.10}
\]

\begin{theorem}[Exact compiled reversal residual]
\label{thm:v18v66-exact-used-residual}
Let
\[
 {\cal W}^{\rm tot}_p
 =\sum_a\sigma_a{\cal W}_{p,a}.
                                                                  \tag{76.11}
\]
Then
\[
\boxed{
\begin{split}
 &{\cal W}^{\rm tot}_p(H_1^R,Q_1^R,\Omega_1^R)
 +\epsilon_Q\epsilon_\Omega
 {\cal W}^{\rm tot}_p(H,Q,\Omega)\\
 &\qquad=\int_0^1{\cal I}_p(s)\,ds.
\end{split}}                                                \tag{76.12}
\]
Consequently the exact used-image certificate is
\[
 \Delta_p^{\rm used}
 :=\left|\int_0^1{\cal I}_p(s)\,ds\right|.                  \tag{76.13}
\]
It contains no ambient dimension factor and preserves cancellation among
sectors and among the three defect channels in (76.10).
\end{theorem}

\begin{proof}
The V64 antiunitary theorem gives
\[
 {\cal W}^{\rm tot}_p(H_0^R,Q_0^R,\Omega_0^R)
 =-\epsilon_Q\epsilon_\Omega
 {\cal W}^{\rm tot}_p(H,Q,\Omega).                          \tag{76.14}
\]
Apply Lemma \ref{lem:v18v66-path-derivative} sectorwise to the affine
interpolation (76.8)--(76.9), form the signed sum before integrating, and
use the fundamental theorem of calculus.  This is (76.12).  Taking the
absolute value after the compiled integral gives (76.13).
\end{proof}

\begin{corollary}[Finite exact computability]
\label{cor:v18v66-finite-computability}
If \(p\) has finite degree and all entries of the two card tuples lie in an
exact computable field, then \({\cal I}_p(s)\) is a scalar polynomial in
\(s\) over that field.  Its signed integral in (76.13) is exactly
computable.  With interval card entries, outward polynomial interval
arithmetic gives a certified enclosure.
\end{corollary}

\begin{proof}
Lemma \ref{lem:v18v66-second-derivative} and (70.15) express every term as a
finite matrix word.  The affine entries are polynomials in \(s\);
multiplication, trace, finite summation, and integration preserve exact
computability.
\end{proof}

\begin{assessment}[Why this is stronger than the V64 fallback]
\label{ass:v18v66-stronger}
The V64 estimate separately bounds
\(\|\Delta H\|,\|\Delta Q\|,\|\Delta\Omega\|\), replaces each used trace by
a full operator-norm product, and pays a dimension factor.  Formula
(76.13) instead evaluates the images of those defects under the actual
second derivative, first derivative, generator, and signed sector sum.
The V64 formula remains a fallback when the finite card entries are not
available.
\end{assessment}

\subsection{A trace-blind nonzero defect}

Use the exact V65 rational card, its odd generator
\(\Omega_{\rm o}=i\sigma_1\), and add at the reflected endpoint
\[
 \Delta\Omega=\frac73i\sigma_3,\qquad
 \|\Delta\Omega\|_{\rm op}=\frac73.                         \tag{76.15}
\]
Keep \(H^R\) and \(Q^R\) at their ideal values.

\begin{theorem}[Nonzero norm defect with zero used residual]
\label{thm:v18v66-blind-defect}
For the exact sign polynomial \(p(x)=x/5\) on the V65 card,
\[
\begin{split}
 {\cal W}_p(\theta)&=-\frac{16}{25},\\
 {\cal W}_p^R(\Omega_{\rm o})&=\frac{16}{25},\\
 {\cal W}_p^R(\Omega_{\rm o}+\Delta\Omega)&=\frac{16}{25},\\
 \Delta_p^{\rm used}&=0.                                   \tag{76.16}
\end{split}
\]
The generator-defect term in the V64 operator-norm fallback is instead
\[
 \frac N2L_1Q_*\|\Delta\Omega\|
 =\frac22\frac15\,4\,\frac73=\frac{28}{15}>0.               \tag{76.17}
\]
\end{theorem}

\begin{proof}
For this card,
\(DP_{-,H^R}[Q^R]\) is off diagonal, whereas
\(\Delta\Omega\) is diagonal.  Therefore
\[
 \operatorname{Tr}
 \bigl(DP_{-,H^R}[Q^R]\Delta\Omega\bigr)=0.                 \tag{76.18}
\]
The measured and ideal reflected Ward traces coincide, so their exact odd
parity residual is zero.  The remaining values are the V65 readings.
For \(p(x)=x/5\), one has \(L_1=1/5\), \(N=2\), and \(Q_*=4\), which gives
(76.17).
\end{proof}

\begin{record}[Exact V66 verifier]
\label{rec:v18v66-verifier}
The Gaussian-rational script
\[
 \texttt{scripts/verify\_sm\_v66\_used\_image\_residual.py}             \tag{76.19}
\]
has SHA--256
\[
 \texttt{ECF7B6B9D422454E11DE7E85D9BC505E3077FEB265981E077A95B00F95171D7F}.
                                                                  \tag{76.20}
\]
Its canonical payload
\[
 \texttt{artifacts/sm\_v66\_used\_image\_residual.json}                \tag{76.21}
\]
has SHA--256
\[
 \texttt{23DD8B43558BC01ADE07460CF3A001A69E6BD23A70E03FE6F99AFF57A0EB2FE8}.
                                                                  \tag{76.22}
\]
\end{record}

\subsection{Adding the sign-functional tail}

Let the exact sign Ward trace and its degree-\(m\) polynomial surrogate
differ by at most \(S_{m,a}\) at each endpoint in sector \(a\), as in V60.
Put
\[
 S_m^{\rm tot}=\sum_a S_{m,a}.                              \tag{76.23}
\]

\begin{corollary}[Exact-sign reversal certificate]
\label{cor:v18v66-exact-sign-certificate}
The exact compiled sign-projector reversal residual is bounded by
\[
\boxed{
 \Delta_{\rm sign}^{\rm rev}
 \le \Delta_{p_m}^{\rm used}+2S_m^{\rm tot}.}                \tag{76.24}
\]
\end{corollary}

\begin{proof}
Insert the polynomial Ward trace at the original and reflected endpoints.
Their compiled residual is (76.13); the two functional-calculus
replacement errors contribute at most \(2S_m^{\rm tot}\).
\end{proof}

\begin{firewall}[Sectorwise sign tails remain coarse]
\label{fire:v18v66-sign-tail-coarse}
The first term in (76.24) is fully compiled, but
\(S_m^{\rm tot}\) is still a sectorwise sum.  A common sign-polynomial
remainder representation could permit further Standard-Model
cancellation.  No such compiled remainder has yet been constructed.
\end{firewall}

\subsection{Cofinal used-image capsule}

\begin{theorem}[Used-image parity sufficiency]
\label{thm:v18v66-cofinal-used}
Suppose the literal finite cards exist along a cofinal sequence and
\[
 \Delta_{p_{m_n},n}^{\rm used}
 +2S_{m_n,n}^{\rm tot}\longrightarrow0.                    \tag{76.25}
\]
If the canonical odd bank also satisfies the V63 incidence conditions and
the V57--V61 topology, locality, atlas, and transport rows hold, then the
reversal-even leakage and higher central odd word both vanish in the
continuum vertical Ward residual.
\end{theorem}

\begin{proof}
Corollary \ref{cor:v18v66-exact-sign-certificate} makes the failure of the
required odd parity vanish.  V64 then reduces the trace to the canonical
sine bank, and V63 makes its higher-word contribution vanish.  The
remaining hypotheses perform the determinant-line and continuum Ward
handoff.
\end{proof}

\begin{firewall}[Literal used images remain unavailable]
\label{fire:v18v66-unpopulated}
The attached files still do not provide the tuple (75.20) for a selected
physical Wilson history, so (76.10) cannot yet be evaluated literally.
V66 proves that incomplete-card norm defects need not be charged in full,
but it does not infer their used images.
\end{firewall}

\subsection{Acquisition and next bottleneck}

V66 incorporates the YM V10 used-image strategy and removes a provably
spurious norm debit on the exact benchmark.  The next upstream task is to
localize (76.10) before a literal full-volume card exists: when
\(\Omega(\omega)\) has finite support, cyclicity of trace should restrict
every polynomial word to a finite neighborhood.  A support-local theorem
would replace the missing global matrix by a finite collar plus an explicit
boundary-crossing tail.

\begin{assessment}[V66 acquisition and boundary]
\label{ass:v18v66}
V66 proves the exact compiled trace-path formula (76.12), its finite
computability, and a rational example where a generator defect of norm
\(7/3\) has zero used residual while the V64 fallback charges \(28/15\).
The literal local collar entries remain the next missing data.
\end{assessment}

\begin{record}[V66 append record]
\label{rec:v18v66}
V66 is appended to the validated V65 whose installed SHA--256 was
\texttt{9BC97364E80090F34AB58938656133857A6F0C88EB71218F19BEC9A4C05FCB3E}.
After installation only V66 is the active numeric SM note; the frozen
\texttt{V100\_f2} through \texttt{V100\_f17} archives are unchanged.
\end{record}

\section{V67: exact support-localization of polynomial Ward and reversal words}
\label{sec:v18v67}

V66 left the literal full-volume reflected matrices unpopulated.  That is
not an obstruction to evaluating a finite-degree polynomial Ward trace.
If the gauge generator is supported on a finite set, cyclicity anchors
every contributing matrix walk there.  Finite propagation then confines
the complete principal trace and the V66 used-image residual to an explicit
collar.  Only the sign-functional approximation tail still uses global
spectral data.

\subsection{Finite-propagation notation}

Let \(G=(V,E)\) be a finite graph with distance \(d_G\), and write
\[
 {\cal H}=\bigoplus_{x\in V}{\cal H}_x.                     \tag{77.1}
\]
An operator \(A\) has range at most \(R_A\) if
\[
 A_{xy}=0\quad\hbox{whenever }d_G(x,y)>R_A.                 \tag{77.2}
\]
For \(S\subseteq V\) and \(R\ge0\), put
\[
 S^{[R]}=\{x\in V:d_G(x,S)\le R\},                          \tag{77.3}
\]
and let \(\Pi_\Lambda\) be the projection onto
\(\bigoplus_{x\in\Lambda}{\cal H}_x\).  The compression of \(A\) is
\[
 A_\Lambda=\Pi_\Lambda A\Pi_\Lambda.                        \tag{77.4}
\]
Assume the gauge generator is supported on \(S\):
\[
 \Omega=\Pi_S\Omega\Pi_S.                                   \tag{77.5}
\]

\begin{lemma}[Anchored closed-walk localization]
\label{lem:v18v67-anchored-walk}
Let \(A_1,\ldots,A_k\) have ranges \(R_1,\ldots,R_k\), and let
\[
 R=R_1+\cdots+R_k.                                          \tag{77.6}
\]
Then
\[
\boxed{
 \operatorname{Tr}(A_1\cdots A_k\Omega)
 =
 \operatorname{Tr}\bigl(
 (A_1)_{S^{[R]}}\cdots(A_k)_{S^{[R]}}\Omega\bigr).}          \tag{77.7}
\]
\end{lemma}

\begin{proof}
Expand the trace in site blocks:
\[
 \sum_{x_0,\ldots,x_k}
 \operatorname{tr}\bigl(
 (A_1)_{x_0x_1}\cdots(A_k)_{x_{k-1}x_k}
 \Omega_{x_kx_0}\bigr).                                    \tag{77.8}
\]
Condition (77.5) forces \(x_0,x_k\in S\).  A nonzero summand satisfies
\(d_G(x_{j-1},x_j)\le R_j\).  Hence every intermediate \(x_j\) lies at
distance at most \(R_1+\cdots+R_j\le R\) from \(S\).  All indices therefore
belong to \(S^{[R]}\), so inserting the collar projection between every
factor changes no summand.
\end{proof}

\begin{remark}[The bound is deliberately pathwise]
\label{rem:v18v67-pathwise}
The radius (77.6) is a worst-case sum.  A generated word bank can use the
actual ordered partial sums and often requires a smaller anisotropic
collar.  The theorem retains the simple isotropic radius because it is
valid before the word list is known.
\end{remark}

\subsection{Polynomial Ward collar}

Let \(H\) have range \(R_H\), let \(Q\) have range \(R_Q\), and let
\(p(x)=\sum_{k=0}^db_kx^k\).  Define
\[
 R_{\rm W}(d)=(d-1)R_H+R_Q,\qquad
 \Lambda_{\rm W}=S^{[R_{\rm W}(d)]}.                        \tag{77.9}
\]

\begin{theorem}[Exact polynomial Ward compression]
\label{thm:v18v67-ward-compression}
Under (77.5),
\[
\boxed{
 {\cal W}_p(H,Q,\Omega)
 =
 {\cal W}_p(H_{\Lambda_{\rm W}},
            Q_{\Lambda_{\rm W}},\Omega).}                   \tag{77.10}
\]
Thus the degree-\(d\) principal Ward trace depends only on matrix entries
inside the finite collar \(\Lambda_{\rm W}\).
\end{theorem}

\begin{proof}
By (70.15), every degree-\(k\) summand of \(Dp_H[Q]\) has the form
\[
 H^jQH^{k-1-j},\qquad 0\le j\le k-1.                        \tag{77.11}
\]
Its total propagation range is at most
\((k-1)R_H+R_Q\le R_{\rm W}(d)\).  Apply Lemma
\ref{lem:v18v67-anchored-walk} to its trace against \(\Omega\), then sum all
finite polynomial terms.  The compressed products are exactly the
polynomial derivative formed from the compressed matrices.
\end{proof}

\begin{corollary}[No volume multiplicity]
\label{cor:v18v67-no-volume}
If the local fibre dimension, \(d\), \(R_H\), \(R_Q\), and the growth of
\(S^{[R_{\rm W}]}\) are bounded independently of the total graph volume,
then evaluation of (77.10) has volume-independent matrix size.
\end{corollary}

\begin{proof}
The compressed carrier dimension is
\(\sum_{x\in\Lambda_{\rm W}}\dim{\cal H}_x\), which is controlled by the
listed local stocks and not by \(|V|\).
\end{proof}

\subsection{Localization of the V66 used-image integrand}

Assume along the reflected interpolation that:
\[
\begin{array}{c|cccccc}
\hbox{operator}
 &H_s^R&Q_s^R&\Omega_s^R&\Delta H&\Delta Q&\Delta\Omega\\ \hline
\hbox{range}
 &R_H&R_Q&0&R_{\Delta H}&R_{\Delta Q}&0\\
\hbox{support}
 &V&V&S&V&V&S .
\end{array}                                                   \tag{77.12}
\]
Here range zero for \(\Omega_s^R,\Delta\Omega\) means site diagonal; a
finite internal matrix at each site is allowed.  Put
\[
\boxed{
\begin{split}
 R_{\rm used}(d)=\max\{&
 (d-2)R_H+R_{\Delta H}+R_Q,\\
 &(d-1)R_H+R_{\Delta Q},\\
 &(d-1)R_H+R_Q\}.
\end{split}}                                                \tag{77.13}
\]
Terms involving \(d<2\) omit the first row rather than using a negative
coefficient.

\begin{theorem}[Exact used-image collar]
\label{thm:v18v67-used-collar}
For a degree-\(d\) polynomial, every trace in the compiled V66 integrand
\({\cal I}_p(s)\) is unchanged when all operators are compressed to
\[
 \Lambda_{\rm used}=S^{[R_{\rm used}(d)]}.                  \tag{77.14}
\]
Consequently
\[
 \boxed{
 \Delta_p^{\rm used}
 =
 \left|\int_0^1{\cal I}_{p,\Lambda_{\rm used}}(s)\,ds\right|.}          \tag{77.15}
\]
\end{theorem}

\begin{proof}
By Lemma \ref{lem:v18v66-second-derivative}, a degree-\(k\) term of
\[
 D^2P_{p,H_s^R}[\Delta H,Q_s^R]\Omega_s^R                  \tag{77.16}
\]
contains \(k-2\) copies of \(H_s^R\), one \(\Delta H\), one \(Q_s^R\),
and the supported generator.  Its total range is bounded by the first row
of (77.13).  A term of
\[
 DP_{p,H_s^R}[\Delta Q]\Omega_s^R                           \tag{77.17}
\]
has the second radius, and a term of
\[
 DP_{p,H_s^R}[Q_s^R]\Delta\Omega                            \tag{77.18}
\]
has the third.  Lemma \ref{lem:v18v67-anchored-walk} applies to every word.
Finite signed summation and integration give (77.15).
\end{proof}

\begin{assessment}[The reflected defect need not be globally small]
\label{ass:v18v67-local-defect}
The norms of \(\Delta H\) or \(\Delta Q\) outside
\(\Lambda_{\rm used}\) do not enter (77.15) at all.  A global reflection
residual can therefore be large while the degree-\(d\) used-image residual
at a local Ward insertion is exactly zero.  This is the spatial version of
the V66 trace-blind generator example.
\end{assessment}

\subsection{A local Hilbert--Schmidt sign tail}

The polynomial collar is exact.  The difference between the exact sign
projector derivative and its polynomial surrogate is controlled by V60.
Suppose
\[
 Q=\Pi_TQ\Pi_T,\qquad \Omega=\Pi_S\Omega\Pi_S,               \tag{77.19}
\]
and let
\[
 r_Q=\operatorname{rank}\Pi_T,\qquad
 r_\Omega=\operatorname{rank}\Pi_S.                         \tag{77.20}
\]

\begin{proposition}[Support-rank sign-tail bound]
\label{prop:v18v67-local-sign-tail}
Under the V60 spectral enclosure \(K_{g,M}\),
\[
\boxed{
\begin{split}
 &\left|
 -i\operatorname{Tr}
 \bigl((DP_{-,H}[Q]-DP_{m,H}[Q])\Omega\bigr)
 \right|\\
 &\quad\le
 \frac{C_\kappa}{2g}\rho_\kappa^{-m}
 \sqrt{r_Qr_\Omega}\,
 \|Q\|_{\rm op}\|\Omega\|_{\rm op}.
\end{split}}                                                \tag{77.21}
\]
\end{proposition}

\begin{proof}
Cauchy--Schwarz in Hilbert--Schmidt norm and (70.9) give
\[
 |\operatorname{Tr}((DP_--DP_m)\Omega)|
 \le\frac{C_\kappa}{2g}\rho_\kappa^{-m}
 \|Q\|_{\rm HS}\|\Omega\|_{\rm HS}.                         \tag{77.22}
\]
The support assumptions imply
\(\|Q\|_{\rm HS}\le\sqrt{r_Q}\|Q\|_{\rm op}\) and the analogous inequality
for \(\Omega\).
\end{proof}

\begin{corollary}[Volume-free local exact-sign certificate]
\label{cor:v18v67-volume-free}
If the support ranks in (77.20), the local operator norms, and the spectral
constants are volume independent, then the sign tail (77.21) has no
full-volume dimension factor.  Combining it at the original and reflected
points with (77.15) gives
\[
 \Delta_{\rm sign}^{\rm rev}
 \le\Delta_{p_m,\Lambda_{\rm used}}^{\rm used}
 +S_{m,{\rm loc}}+S_{m,{\rm loc}}^R.                        \tag{77.23}
\]
\end{corollary}

\begin{proof}
Apply Proposition \ref{prop:v18v67-local-sign-tail} at the two endpoints
and use the V66 exact-sign insertion argument.
\end{proof}

\subsection{Local card sufficiency}

\begin{theorem}[Finite collar replacement of the V65 tuple]
\label{thm:v18v67-local-card}
For a degree-\(d\) sign polynomial and local tangent/generator supports,
the following data determine a rigorous reversal enclosure:
\begin{enumerate}[leftmargin=2.2em]
\item the compressed original and reflected matrices
      \(H,Q,\Omega,\Delta H,\Delta Q,\Delta\Omega\) on
      \(\Lambda_{\rm used}\);
\item the polynomial coefficients and antiunitary carrier map on that
      collar;
\item global scalar spectral bounds \(0<g\le|H|\le M\);
\item local support ranks and operator norms for (77.21).
\end{enumerate}
No full-volume matrix entries are required.
\end{theorem}

\begin{proof}
Items one and two compute (77.15) exactly or by intervals.  Items three and
four compute the two local sign tails in (77.23).  The sum is the required
reversal enclosure.
\end{proof}

\begin{firewall}[A global gap is still global information]
\label{fire:v18v67-global-gap}
The collar cannot certify the spectral gap of the full Wilson kernel.
Admissibility or a separate global spectral theorem must supply \(g\) and
\(M\).  A gapped compressed matrix does not imply a gapped full matrix,
because exterior modes and boundary coupling can create small eigenvalues.
\end{firewall}

\subsection{Cofinal collar capsule}

\begin{theorem}[Support-local parity handoff]
\label{thm:v18v67-cofinal-collar}
Suppose at cutoff \(n\) the local card of Theorem
\ref{thm:v18v67-local-card} is populated and
\[
 \Delta_{p_{m_n},\Lambda_n}^{\rm used}
 +S_{m_n,{\rm loc},n}+S_{m_n,{\rm loc},n}^R\longrightarrow0.           \tag{77.24}
\]
Assume also that the physical collar radii
\[
 a_nR_{{\rm used},n}\longrightarrow0                        \tag{77.25}
\]
and the V63 incidence, V57 determinant-line, and V58--V61 transport rows
hold.  Then the local reversal leakage vanishes in the continuum vertical
Ward residual with shrinking physical support.
\end{theorem}

\begin{proof}
Equation (77.23) and (77.24) make the exact-sign reversal residual vanish.
Equation (77.25) makes the determining card local at continuum scale.  The
V63 conditions control the remaining canonical odd bank, and the cited
global rows perform the gauge-line and continuum handoff.
\end{proof}

\begin{firewall}[Literal collar still not emitted]
\label{fire:v18v67-unpopulated}
The source materials do not emit the compressed Wilson, tangent,
generator, and reflection-defect matrices on any selected physical collar.
V67 proves that this finite local packet is sufficient, but does not
populate it or its global gap.
\end{firewall}

\subsection{Acquisition and next bottleneck}

V67 removes the need for a full-volume matrix read.  The next upstream
target is now sharply finite: select one supported gauge generator and
emit the Wilson hopping blocks on the radius (77.13), together with an
admissibility gap.  If no physical history is specified, the programme can
still develop a basis-independent interval schema which accepts such a
collar and returns (77.23) without silently fitting the missing data.

\begin{assessment}[V67 acquisition and boundary]
\label{ass:v18v67}
V67 proves exact anchored-walk compression for the polynomial Ward trace
and V66 used-image residual, and replaces the full dimension loss in the
sign tail by local support ranks.  It reduces the missing literal input to
a finite collar plus two global spectral scalars.
\end{assessment}

\begin{record}[V67 append record]
\label{rec:v18v67}
V67 is appended to the validated V66 whose installed SHA--256 was
\texttt{0BAD4062A794D77FC615BECFF0BBEE1B90C7DC6C46A3250C6ABFBE0ABB860B50}.
After installation only V67 is the active numeric SM note; the frozen
\texttt{V100\_f2} through \texttt{V100\_f17} archives are unchanged.
\end{record}


\section{V68: an exact Gaussian-rational collar-card compiler}
\label{sec:v18v68}

V67 reduced the polynomial part of the local reversal problem to a finite
matrix collar, but it left the phrase ``populate the collar'' informal.  This
note turns that phrase into a machine-readable contract.  The contract is
deliberately narrower than a physical Standard-Model certificate: it accepts
finite matrices already extracted from a proposed collar, verifies their
algebraic types, constructs the ideal antiunitary image, and returns the exact
compiled polynomial Ward residual.  All accepted arithmetic lies in the
Gaussian rationals, so the output contains no floating-point tolerance or
fitted entry.

\subsection{The finite certificate contract}

Write
\[
 \mathbb Q(i)=\{a+ib:a,b\in\mathbb Q\}.                     \tag{78.1}
\]
For every sector \(a\), a collar card supplies a sign
\(\sigma_a\in\{-1,1\}\), a dimension \(N_a\), and matrices
\[
 U_a,H_a,Q_a,\Omega_a,H_a^R,Q_a^R,\Omega_a^R
 \in M_{N_a}(\mathbb Q(i)).                                 \tag{78.2}
\]
It also supplies common parities
\(\epsilon_Q,\epsilon_\Omega\in\{-1,1\}\) and a real rational polynomial
\[
 p(x)=\sum_{k=0}^d b_kx^k,
 \qquad b_k\in\mathbb Q.                                    \tag{78.3}
\]
The algebraic admissibility rows are
\[
\begin{split}
 U_aU_a^*&=I,\\
 H_a^*=H_a,\quad Q_a^*&=Q_a,\quad \Omega_a^*&=-\Omega_a,\\
 (H_a^R)^*&=H_a^R,\quad (Q_a^R)^*=Q_a^R,
 \quad(\Omega_a^R)^*=-\Omega_a^R.                           \tag{78.4}
\end{split}
\]
Define the antiunitary matrix action and the ideal reflected tuple by
\[
 \mathfrak C_a(A)=U_a\overline A,U_a^*,                    \tag{78.5}
\]
\[
 H_a^0=\mathfrak C_a(H_a),\qquad
 Q_a^0=\epsilon_Q\mathfrak C_a(Q_a),\qquad
 \Omega_a^0=\epsilon_\Omega\mathfrak C_a(\Omega_a).         \tag{78.6}
\]
The measured defects are not thresholded; they are the exact matrices
\[
 \Delta H_a=H_a^R-H_a^0,\quad
 \Delta Q_a=Q_a^R-Q_a^0,\quad
 \Delta\Omega_a=\Omega_a^R-\Omega_a^0.                      \tag{78.7}
\]

For any admissible \((H,Q,\Omega)\), put
\[
 DP_{p,H}[Q]
 =-\frac12\sum_{k=1}^d b_k
   \sum_{j=0}^{k-1}H^jQH^{k-1-j},                           \tag{78.8}
\]
and
\[
 {\cal W}_p(H,Q,\Omega)
 =-i\operatorname{Tr}\bigl(DP_{p,H}[Q]\Omega\bigr).        \tag{78.9}
\]

\begin{lemma}[Reality of an accepted Ward value]
\label{lem:v18v68-reality}
Under (78.3)--(78.4), the quantity in (78.9) is rational and real.
\end{lemma}

\begin{proof}
Every coefficient and entry in (78.8) belongs to \(\mathbb Q(i)\).  Since
\(p\) is real and \(H,Q\) are Hermitian, differentiating
\(p(H+tQ)^*=p(H+tQ)\) at \(t=0\) gives
\[
 DP_{p,H}[Q]^*=DP_{p,H}[Q].                                 \tag{78.10}
\]
Let \(D=DP_{p,H}[Q]\) and \(z=\operatorname{Tr}(D\Omega)\).  Cyclicity
and \(\Omega^*=-\Omega\) give
\[
 \overline z=\operatorname{Tr}(\Omega^*D^*)
 =-\operatorname{Tr}(\Omega D)=-z.                         \tag{78.11}
\]
Thus \(z\in i\mathbb Q\), and \(-iz\in\mathbb Q\).
\end{proof}

\subsection{Exact antiunitary soundness}

\begin{lemma}[Polynomial derivative transport]
\label{lem:v18v68-derivative-transport}
For the map (78.5),
\[
 DP_{p,\mathfrak C_a(H)}[\mathfrak C_a(Q)]
 =\mathfrak C_a\bigl(DP_{p,H}[Q]\bigr).                    \tag{78.12}
\]
Moreover,
\[
 \operatorname{Tr}(\mathfrak C_a(A))
 =\overline{\operatorname{Tr}(A)}.                         \tag{78.13}
\]
\end{lemma}

\begin{proof}
Unitary conjugation and entrywise conjugation preserve products in their
original order, so
\[
 \mathfrak C_a(H^jQH^{k-1-j})
 =\mathfrak C_a(H)^j\mathfrak C_a(Q)
  \mathfrak C_a(H)^{k-1-j}.                                \tag{78.14}
\]
The \(b_k\) are real.  Applying \(\mathfrak C_a\) term by term to (78.8)
proves (78.12).  Unitary invariance of trace proves (78.13).
\end{proof}

\begin{theorem}[Sectorwise ideal parity identity]
\label{thm:v18v68-ideal-parity}
Every algebraically admissible sector obeys
\[
\boxed{
 {\cal W}_p(H_a^0,Q_a^0,\Omega_a^0)
 +\epsilon_Q\epsilon_\Omega
  {\cal W}_p(H_a,Q_a,\Omega_a)=0.}                          \tag{78.15}
\]
\end{theorem}

\begin{proof}
Linearity in the tangent and generator, followed by Lemma
\ref{lem:v18v68-derivative-transport}, gives
\[
\begin{split}
 &\operatorname{Tr}
 \bigl(DP_{p,H_a^0}[Q_a^0]\Omega_a^0\bigr)\\
 &\qquad=\epsilon_Q\epsilon_\Omega
 \overline{
 \operatorname{Tr}
 \bigl(DP_{p,H_a}[Q_a]\Omega_a\bigr)}.                    \tag{78.16}
\end{split}
\]
The traced number on the right is purely imaginary by (78.11); conjugation
therefore changes its sign.  Multiplication by \(-i\) proves (78.15).
\end{proof}

Define the three compiled readings
\[
\begin{split}
 W_p^{\rm F}&=\sum_a\sigma_a
 {\cal W}_p(H_a,Q_a,\Omega_a),\\
 W_p^{0,R}&=\sum_a\sigma_a
 {\cal W}_p(H_a^0,Q_a^0,\Omega_a^0),\\
 W_p^{R}&=\sum_a\sigma_a
 {\cal W}_p(H_a^R,Q_a^R,\Omega_a^R).                       \tag{78.17}
\end{split}
\]
The signed measured reversal residual is
\[
 R_p^{\rm card}=W_p^R+epsilon_Q\epsilon_\Omega W_p^{\rm F}.
                                                                  \tag{78.18}
\]

\begin{corollary}[Measured-minus-ideal form]
\label{cor:v18v68-measured-minus-ideal}
For every accepted card,
\[
\boxed{
 R_p^{\rm card}=W_p^R-W_p^{0,R}
 =\sum_a\sigma_a\left(
 {\cal W}_p(H_a^R,Q_a^R,\Omega_a^R)
 -{\cal W}_p(H_a^0,Q_a^0,\Omega_a^0)\right).}              \tag{78.19}
\]
Thus cancellations between sectors occur before an absolute value is taken.
\end{corollary}

\begin{proof}
Multiply (78.15) by \(\sigma_a\), sum over sectors, and substitute into
(78.18).
\end{proof}

\begin{theorem}[Exactness of the Gaussian-rational certificate]
\label{thm:v18v68-exact-certificate}
Suppose a finite card passes (78.4).  Then every matrix in (78.6)--(78.8),
every sector Ward value, and \(R_p^{\rm card}\) is exactly computable in
\(\mathbb Q(i)\).  If the emitted result states
\(R_p^{\rm card}=0\), the polynomial reversal identity is an exact equality,
not a conclusion from numerical closeness.
\end{theorem}

\begin{proof}
Gaussian rationals are closed under conjugation, addition, and
multiplication.  The formulas use only those operations, rational scalar
multiplication, finite sums, and trace.  Fractions are kept as reduced
integer numerator--denominator pairs.  No limit, eigenvalue routine, decimal
conversion, or tolerance occurs.  The returned equality is therefore
equality of rational numerators after exact reduction.
\end{proof}

\subsection{Executable realization}

The verifier
\[
 \texttt{scripts/verify\_sm\_v68\_exact\_collar\_card.py}              \tag{78.20}
\]
implements (78.2)--(78.19).  It rejects a nonsquare carrier, inconsistent
sector dimensions, nonunitary \(U_a\), incorrectly typed Hermitian or
anti-Hermitian matrices, invalid signs, a nonreal Ward trace, or failure of
the independently forced identity (78.15).  On success it emits canonical
JSON with lexicographically sorted keys and compact separators.

\begin{proposition}[Verifier acceptance invariant]
\label{prop:v18v68-verifier-invariant}
If the verifier terminates successfully on a card, then:
\begin{enumerate}[leftmargin=2.2em]
\item all algebraic conditions in (78.4) hold exactly;
\item its displayed \(\Delta H_a,\Delta Q_a,\Delta\Omega_a\) equal (78.7);
\item every displayed sector value equals (78.9);
\item every ideal parity residual equals zero by an exact assertion; and
\item the displayed compiled residual equals (78.18) and (78.19).
\end{enumerate}
\end{proposition}

\begin{proof}
The parser maps every pair of rational strings to an element of
\(\mathbb Q(i)\).  The program then performs the tests in item one, forms
(78.6)--(78.7), evaluates the finite double sum (78.8), and applies (78.9).
Before sector accumulation it asserts (78.15).  Its final accumulation is
the signed sum (78.17), followed by (78.18).  These are direct evaluations
of the displayed formulas in exact arithmetic.
\end{proof}

This proposition is a transparent code-to-formula audit, not a formal
verification of the Python interpreter.  The mathematical soundness of the
identity itself is Theorem \ref{thm:v18v68-ideal-parity}.

\subsection{Executed trace-blind card}

The sample input is
\[
 \texttt{artifacts/sm\_v68\_sample\_collar\_card.json}.                 \tag{78.21}
\]
It uses one sector, \(U=I\),
\(\epsilon_Q=\epsilon_\Omega=-1\), \(p(x)=x/5\), and
\[
 H=\begin{pmatrix}
 -3&(-12+16i)/5\\
 (-12-16i)/5&3
 \end{pmatrix},
 \quad
 Q=\begin{pmatrix}
 0&(16+12i)/5\\
 (16-12i)/5&0
 \end{pmatrix},                                             \tag{78.22}
\]
\[
 \Omega=\begin{pmatrix}0&i\\ i&0\end{pmatrix}.
                                                                  \tag{78.23}
\]
The measured reflected \(H^R,Q^R\) equal their ideal values, whereas
\[
 \Omega^R=\Omega^0+\Delta\Omega,
 \qquad
 \Delta\Omega=\frac73
 \begin{pmatrix}i&0\\0&-i\end{pmatrix}.                    \tag{78.24}
\]

\begin{theorem}[Executed exact benchmark]
\label{thm:v18v68-executed-benchmark}
The verifier accepts the card (78.21) and emits
\[
 W_p^{\rm F}=-\frac{16}{25},\qquad
 W_p^{0,R}=\frac{16}{25},\qquad
 W_p^R=\frac{16}{25},\qquad
 R_p^{\rm card}=0.                                         \tag{78.25}
\]
It also emits
\[
 \Delta H=0,qquad \Delta Q=0,qquad
 \Delta\Omega\ne0.                                        \tag{78.26}
\]
\end{theorem}

\begin{proof}
For \(p(x)=x/5\), formula (78.8) reduces to
\[
 DP_{p,H}[Q]=-\frac1{10}Q.                                 \tag{78.27}
\]
Direct multiplication of (78.22)--(78.23) gives the forward value in
(78.25).  Theorem \ref{thm:v18v68-ideal-parity} gives the ideal reflected
value.  At the measured endpoint, \(DP_{p,H^R}[Q^R]\) is off diagonal while
\(\Delta\Omega\) in (78.24) is diagonal, hence
\[
 \operatorname{Tr}
 \bigl(DP_{p,H^R}[Q^R]\Delta\Omega\bigr)=0.                \tag{78.28}
\]
This proves the measured value and residual.  Exact execution independently
reproduces these rational operations and the zero/nonzero defect flags.
\end{proof}

The canonical result is stored at
\[
 \texttt{artifacts/sm\_v68\_sample\_collar\_result.json}.               \tag{78.29}
\]
The immutable SHA--256 ledger is
\[
\begin{array}{c|l}
\hbox{object}&\hbox{SHA--256}\\ \hline
\hbox{verifier}&
\texttt{336C164C3DE6DBD8C5EBAB58903E0ABFBDF6A06D57692D3E7C63465F24832FC2}\\
\hbox{input card}&
\texttt{533F01FFD61F46FD5B8D839323CA104CCDB908A9EEFC97D86A98A6FA00822330}\\
\hbox{canonical result}&
\texttt{8B1B54EF0D3225EF5E834A890873480CDEC72D144D83B5D39201D226A2926DD2}.
\end{array}                                                   \tag{78.30}
\]

\subsection{Handoff to the exact-sign collar bound}

For a physical cutoff card let \(R_{p_m,n}^{\rm card}\) be (78.18), and
let \(S_{m,n}^{\rm F},S_{m,n}^{R}\) be the two V67 support-rank sign tails.

\begin{theorem}[Executable local reversal enclosure]
\label{thm:v18v68-executable-enclosure}
If the input matrices are the literal V67 collars of the forward and
reflected Wilson histories, and the supplied global gap data justify the two
sign tails, then the exact sign-projector reversal residual obeys
\[
\boxed{
 \Delta_{{\rm sign},n}^{\rm rev}
 \le |R_{p_m,n}^{\rm card}|
      +S_{m,n}^{\rm F}+S_{m,n}^{R}.}                        \tag{78.31}
\]
If the right-hand side tends to zero along a cofinal sequence satisfying
the V67 shrinking-collar and V63 incidence hypotheses, the reversal leakage
vanishes in the continuum vertical Ward residual.
\end{theorem}

\begin{proof}
Corollary \ref{cor:v18v68-measured-minus-ideal} computes the signed
polynomial principal residual with all sector cancellations retained.  Add
and subtract the polynomial projector derivative at the forward and
reflected endpoints.  The two remaining trace differences are bounded by
the V67 local sign tails, which proves (78.31).  The cofinal conclusion is
Theorem \ref{thm:v18v67-cofinal-collar} with the principal term now supplied
by the exact card compiler.
\end{proof}

\begin{firewall}[What an accepted card does not certify]
\label{fire:v18v68-card-boundary}
Acceptance certifies the algebra and arithmetic of the matrices that were
supplied.  It does not prove that they came from the selected physical
Wilson operator, that the collar radius satisfies (77.13), that the full
operator has gap \(g\), or that the Standard-Model sector and path list is
complete.  The executed card is the V65--V66 benchmark, not a literal
Standard-Model overlap card.  It therefore proves that the acquisition
interface is executable and that coarse defect norms can overcharge the
used trace; it does not close the physical-source row.
\end{firewall}

\subsection{Acquisition and next bottleneck}

V68 replaces an informal request for local matrices by a checked data
contract and an exact compiled output.  The nearest upstream obstruction is
now data provenance rather than trace algebra.  A usable physical producer
must emit each hopping block together with its site, internal-index, sector,
path-orientation, and reflection metadata.  Since physical input may be
rational only through certified enclosures, the next note should extend the
contract from points in \(\mathbb Q(i)\) to complex rational rectangles and
prove a width bound for the compiled Ward residual.  That extension must
retain signed sector cancellation before taking the final enclosure radius.

\begin{assessment}[V68 acquisition and boundary]
\label{ass:v18v68}
V68 proves the finite-card soundness theorem, supplies an executable exact
Gaussian-rational compiler, and reproduces the V66 trace-blind example from
a serialized card with immutable hashes.  The remaining physical inputs are
the literal Standard-Model collar, a completeness ledger, and a global gap
certificate.
\end{assessment}

\begin{record}[V68 append record]
\label{rec:v18v68}
V68 is appended to the validated V67 whose installed SHA--256 was
\texttt{C6FAA091EAB1A58635701200853C88506C0839B42D316144209724C0627C240C}.
After installation only V68 is the active numeric SM note; the frozen
\texttt{V100\_f2} through \texttt{V100\_f17} archives are unchanged.
\end{record}


\section{V69: rational-rectangle certification for uncertain physical collars}
\label{sec:v18v69}

The exact V68 card is suitable for symbolic benchmarks, but a physical
collar producer may obtain hopping coefficients through certified numerical
evaluation.  Its honest output is then a family of rational rectangles, not
a guessed rational point.  V69 extends the compiler to those rectangles and
proves that the final real interval encloses the physical signed reversal
residual.  Ideal reflection is formed from the same forward rectangles, and
measured-minus-ideal sector contributions are summed before the absolute
bound is taken.

\subsection{Complex rational rectangles}

For rational endpoints \(a_-\le a_+\) and \(b_-\le b_+\), write
\[
 \mathbf z=[a_-,a_+]+i[b_-,b_+]
 =\{x+iy:a_-\le x\le a_+,\ b_-\le y\le b_+\}.              \tag{79.1}
\]
For real intervals \(I,J\), use their inclusion operations
\[
 I+J=\{x+y:x\in I,y\in J\},
 \qquad
 I\cdot J=operatorname{hull}\{xy:x\in I,y\in J\}.         \tag{79.2}
\]
The second hull has endpoints equal to the minimum and maximum of the four
endpoint products.  Complex-rectangle operations are
\[
\begin{split}
 (I+iJ)+(K+iL)&=(I+K)+i(J+L),\\
 (I+iJ)(K+iL)&=(IK-JL)+i(IL+JK),\\
 \overline{I+iJ}&=I-iJ.                                    \tag{79.3}
\end{split}
\]
All endpoints remain rational.

An exact center \(z^c=x^c+iy^c\in\mathbb Q(i)\) and nonnegative rational
radii \(r^{\rm re},r^{\rm im}\) generate
\[
 \mathbf z(z^c;r^{\rm re},r^{\rm im})
 =[x^c-r^{\rm re},x^c+r^{\rm re}]
 +i[y^c-r^{\rm im},y^c+r^{\rm im}].                        \tag{79.4}
\]
A matrix rectangle is an entrywise array of such sets.  Matrix addition,
multiplication, adjoint, powers, and trace use (79.2)--(79.3) entrywise.

\begin{lemma}[Natural inclusion]
\label{lem:v18v69-natural-inclusion}
Let \({\bf A}_1,\ldots,{\bf A}_r\) be complex matrix rectangles, and choose
arbitrary point matrices \(A_j\in{\bf A}_j\) entrywise.  If \(F\) is any
finite matrix expression made from addition, subtraction, multiplication,
adjoint, rational scalar multiplication, and trace, then
\[
 F(A_1,\ldots,A_r)\in
 \mathbf F({\bf A}_1,\ldots,{\bf A}_r),                    \tag{79.5}
\]
where \(\mathbf F\) is obtained by replacing every scalar operation by
(79.2)--(79.3).
\end{lemma}

\begin{proof}
For scalar addition, multiplication, and conjugation the claim follows
directly from the set definitions in (79.2)--(79.3).  Matrix multiplication
is a finite sum of enclosed scalar products, and trace is a finite sum of
enclosed diagonal entries.  Structural induction on the expression proves
(79.5).  Repeated occurrences of the same input may enlarge the natural
interval, but cannot invalidate inclusion.
\end{proof}

\subsection{The interval Ward compiler}

Let the exact V68 center card pass all algebraic tests.  Replace each of
\(H_a,Q_a,\Omega_a,H_a^R,Q_a^R,\Omega_a^R\) by a centered matrix rectangle
containing the corresponding unknown physical matrix.  The carrier
\(U_a\), signs, and polynomial coefficients remain exact.  Define
\[
 \mathbf H_a^0
 =U_a\overline{\mathbf H_a}\,U_a^*,\qquad
 \mathbf Q_a^0
 =\epsilon_QU_a\overline{\mathbf Q_a}\,U_a^*,
 \qquad
 \boldsymbol\Omega_a^0
 =\epsilon_\Omega U_a\overline{\boldsymbol\Omega_a}\,U_a^*.
                                                                  \tag{79.6}
\]
Apply the natural extension of (78.8)--(78.9) to obtain the complex
rectangle
\[
 \mathbf W_p(\mathbf H,\mathbf Q,\boldsymbol\Omega)
 =-i\operatorname{Tr}
 \bigl(\mathbf {DP}_{p,\mathbf H}[\mathbf Q]
       \boldsymbol\Omega\bigr).                             \tag{79.7}
\]
The compiled rectangle is formed in measured-minus-ideal form:
\[
\boxed{
 \mathbf R_p
 =\sum_a\sigma_a\left(
 \mathbf W_p(\mathbf H_a^R,\mathbf Q_a^R,
              \boldsymbol\Omega_a^R)
 -\mathbf W_p(\mathbf H_a^0,\mathbf Q_a^0,
              \boldsymbol\Omega_a^0)
 \right).}                                                 \tag{79.8}
\]

\begin{theorem}[Rational-rectangle reversal enclosure]
\label{thm:v18v69-rectangle-enclosure}
Suppose the literal physical matrices lie entrywise in the supplied
rectangles and satisfy the V68 Hermitian and anti-Hermitian type conditions.
Then their polynomial reversal residual satisfies
\[
 R_p^{\rm phys}\in \mathbf R_p\cap\mathbb R.               \tag{79.9}
\]
Writing the real projection of \(\mathbf R_p\) as
\([R_-,R_+]\), one has the rigorous bound
\[
\boxed{
 |R_p^{\rm phys}|\le
 B_p^{\rm rect}:=\max\{|R_-|,|R_+|\}.}                    \tag{79.10}
\]
\end{theorem}

\begin{proof}
The physical ideal reflected tuple lies in (79.6) by Lemma
\ref{lem:v18v69-natural-inclusion}.  The same lemma applied to (78.8)--(78.9)
encloses every measured and ideal Ward value by its rectangle in (79.7).
Signed interval summation encloses their physical sum in (79.8).  By Lemma
\ref{lem:v18v68-reality}, the physical sum is real, proving (79.9).  Every
real number in \([R_-,R_+]\) has modulus at most (79.10).
\end{proof}

\begin{remark}[The imaginary interval is a compatibility diagnostic]
\label{rem:v18v69-imaginary-diagnostic}
Independent entry rectangles generally contain matrices that violate the
adjoint correlations, so \(\mathbf R_p\) may have nonzero imaginary width.
This causes no loss of validity: the physical typed point lies on the real
slice.  If the imaginary interval excludes zero, however, the rectangles
cannot contain any typed physical realization satisfying the asserted
formula.  The V69 verifier rejects that case.
\end{remark}

\begin{proposition}[Exact-center reduction]
\label{prop:v18v69-exact-reduction}
If every radius in (79.4) is zero, the compiler returns the singleton
\[
 \mathbf R_p=\{R_p^{\rm card}\},                            \tag{79.11}
\]
where \(R_p^{\rm card}\) is the V68 exact result.
\end{proposition}

\begin{proof}
Every input rectangle is then a singleton.  Each operation in
(79.2)--(79.3) agrees with the corresponding Gaussian-rational operation,
so (79.8) is (78.19).
\end{proof}

\subsection{An analytic width control}

The natural interval is executable and often sharper than a norm estimate.
For cofinal reasoning it is also useful to have a closed width bound.  Put
\[
 L_1(M)=\sum_{k=1}^d k|b_k|M^{k-1},qquad
 L_2(M)=\sum_{k=2}^d k(k-1)|b_k|M^{k-2}.                   \tag{79.12}
\]

\begin{lemma}[Three-input Ward stability]
\label{lem:v18v69-ward-stability}
Let \(X=(H,Q,\Omega)\) and
\(X^c=(H^c,Q^c,\Omega^c)\) have dimension \(N\), with
\[
 \|H\|_{\rm op},\|H^c\|_{\rm op}\le M,
 \quad
 \|Q\|_{\rm op},\|Q^c\|_{\rm op}\le Q_*,
 \quad
 \|\Omega\|_{\rm op},\|\Omega^c\|_{\rm op}\le\Omega_*.
                                                                  \tag{79.13}
\]
If the three operator-norm deviations are
\(\delta_H,\delta_Q,\delta_\Omega\), then
\[
\boxed{
 |{\cal W}_p(X)-{\cal W}_p(X^c)|
 \le\frac N2\left[
 \Omega_*\bigl(L_2(M)Q_*\delta_H+L_1(M)\delta_Q\bigr)
 +L_1(M)Q_*\delta_\Omega\right].}                          \tag{79.14}
\]
\end{lemma}

\begin{proof}
In a degree-\(k\) derivative word there are \(k-1\) copies of \(H\) and one
copy of \(Q\).  Telescoping the \(H\) factors across the \(k\) possible
positions of \(Q\) gives at most \(k(k-1)\) terms, each bounded by
\(M^{k-2}Q_*\delta_H\).  Changing \(Q\) gives \(k\) terms bounded by
\(M^{k-1}\delta_Q\).  Therefore
\[
 \|DP_{p,H}[Q]-DP_{p,H^c}[Q^c]\|_{\rm op}
 \le\frac12\left(L_2Q_*\delta_H+L_1\delta_Q\right).        \tag{79.15}
\]
Also \(\|DP_{p,H^c}[Q^c]\|_{\rm op}\le L_1Q_*/2\).  Split the trace
difference by adding
\(DP_{p,H^c}[Q^c]\Omega\), use
\(|\operatorname{Tr}A|\le N\|A\|_{\rm op}\), and obtain (79.14).
\end{proof}

For an entry rectangle centered at \(A^c\), define the rational stock
\[
 \widehat\delta_A
 =\sum_{r,s}\left(r_{rs}^{\rm re}+r_{rs}^{\rm im}\right).   \tag{79.16}
\]

\begin{corollary}[Rational radius-to-operator bound]
\label{cor:v18v69-entry-radius}
Every point matrix \(A\) in the rectangle satisfies
\[
 \|A-A^c\|_{\rm op}\le\widehat\delta_A.                   \tag{79.17}
\]
Consequently (79.14), applied to the measured and ideal triples, gives a
fully rational a priori radius for (79.8).  The ideal deviations equal the
corresponding forward deviations because antiunitary conjugation preserves
operator norm.
\end{corollary}

\begin{proof}
Each entry deviation obeys
\(|z_{rs}-z_{rs}^c|\le r_{rs}^{\rm re}+r_{rs}^{\rm im}\).  Hence
\[
 \|A-A^c\|_{\rm op}
 \le\|A-A^c\|_{\rm HS}
 \le\sum_{r,s}|z_{rs}-z_{rs}^c|
 \le\widehat\delta_A.                                     \tag{79.18}
\]
The last statement follows from (78.5).
\end{proof}

\begin{corollary}[Cofinal interval collapse]
\label{cor:v18v69-cofinal-collapse}
Assume the sector dimensions, polynomial degree, \(M,Q_*,\Omega_*\), and
coefficient stocks in (79.12) are uniformly bounded.  If the exact center
residuals tend to zero and every entry-radius stock (79.16) tends to zero
after signed sector compilation, then
\[
 B_{p,n}^{\rm rect}\longrightarrow0.                       \tag{79.19}
\]
If the two V67 sign tails also vanish, the exact sign-projector reversal
residual tends to zero.
\end{corollary}

\begin{proof}
Lemma \ref{lem:v18v69-ward-stability} and Corollary
\ref{cor:v18v69-entry-radius} make the deviation of the physical compiled
residual from its exact center tend to zero.  The center tends to zero by
hypothesis, proving (79.19).  Add the two sign tails as in (78.31).
\end{proof}

\subsection{Executed interval card}

The uncertainty file
\[
 \texttt{artifacts/sm\_v69\_sample\_collar\_uncertainty.json}           \tag{79.20}
\]
uses the exact V68 center.  All radii vanish except the real and imaginary
radii of the two off-diagonal entries of \(\Omega^R\), each of which is
\(1/1000\).  The interval compiler is
\[
 \texttt{scripts/verify\_sm\_v69\_interval\_collar\_card.py}.          \tag{79.21}
\]

\begin{theorem}[Executed rational enclosure]
\label{thm:v18v69-executed-enclosure}
Exact execution returns
\[
 \mathbf W_p^{0,R}=\left\{\frac{16}{25}\right\},          \tag{79.22}
\]
and the real projection
\[
 \operatorname{Re}\mathbf W_p^R
 =\left[\frac{3993}{6250},\frac{4007}{6250}\right].         \tag{79.23}
\]
Therefore
\[
\boxed{
 R_p^{\rm phys}\in
 \left[-\frac7{6250},\frac7{6250}\right],
 \qquad
 |R_p^{\rm phys}|\le\frac7{6250}.}                         \tag{79.24}
\]
The imaginary diagnostic is the same symmetric interval and contains zero.
\end{theorem}

\begin{proof}
The ideal inputs have zero radius, so Proposition
\ref{prop:v18v69-exact-reduction} and (78.25) give (79.22).  With
\(DP_{p,H^R}[Q^R]=-Q^R/10\), interval multiplication of the two uncertain
off-diagonal generator entries gives (79.23).  Subtracting the singleton
(79.22) gives (79.24).  The executable evaluates precisely these rational
endpoint operations.
\end{proof}

The canonical result is
\[
 \texttt{artifacts/sm\_v69\_sample\_collar\_result.json}.                \tag{79.25}
\]
Its immutable ledger is
\[
\begin{array}{c|l}
\hbox{object}&\hbox{SHA--256}\\ \hline
\hbox{V68 center card}&
\texttt{533F01FFD61F46FD5B8D839323CA104CCDB908A9EEFC97D86A98A6FA00822330}\\
\hbox{uncertainty card}&
\texttt{894AF011F923371D1B33C2DA0C420559A56E77DB068041366ABC21F248D41A53}\\
\hbox{interval compiler}&
\texttt{C3FBF1151A1D4E17A61927AED88B23B90BBE4B1DD3F59913498460812C2BDD57}\\
\hbox{canonical result}&
\texttt{9914CF69955AF20CB39B5DD884D3380E9F8647CF4DB5A9027E0EB56B45821F7F}.
\end{array}                                                   \tag{79.26}
\]

\begin{firewall}[Interval enclosure does not create provenance]
\label{fire:v18v69-provenance}
The compiler proves that any compatible physical matrices inside the
rectangles satisfy (79.10).  It does not prove that the rectangles contain
the Wilson hopping blocks of the intended Standard-Model history.  It also
does not validate collar completeness, the global gap, or path and sector
completeness.  Independent-entry rectangles discard correlations and may be
wide; this affects sharpness, not soundness.
\end{firewall}

\subsection{Acquisition and next bottleneck}

The numerical-representation bottleneck is now closed at the local
polynomial level: exact points and certified rational rectangles share one
sound measured-minus-ideal interface.  The next upstream task is a
provenance ledger that maps every emitted collar block to a Wilson stencil,
oriented link, representation, chirality, and reflection partner, and proves
that the emitted list is complete for the V67 radius.  Once that finite
ledger exists, its interval output can feed (79.8) directly; the independent
global task remains an admissibility-to-gap theorem.

\begin{assessment}[V69 acquisition and boundary]
\label{ass:v18v69}
V69 proves sound rational-rectangle evaluation, gives an analytic width
bound, and executes a nonzero-width card enclosing the exact V68 result by
\(7/6250\).  This removes floating-point exactness as an excuse for leaving
the physical collar unpopulated.  Provenance, completeness, and the global
gap remain open.
\end{assessment}

\begin{record}[V69 append record]
\label{rec:v18v69}
V69 is appended to the validated V68 whose installed SHA--256 was
\texttt{891E4A153686B6C0DD1262411870C5101AD7CF6C3FF54FD12257D96EF9C60B02}.
After installation only V69 is the active numeric SM note; the frozen
\texttt{V100\_f2} through \texttt{V100\_f17} archives are unchanged.
\end{record}


\section{V70: a finite Standard-Model collar provenance ledger}
\label{sec:v18v70}

V69 can rigorously evaluate interval matrices, but an interval with unknown
origin is not physical evidence.  The first provenance question is finite:
given a declared Wilson--Yukawa stencil on the V67 collar, has the producer
emitted every allowed block exactly once, with the correct carrier shape and
an explicit reflected partner?  V70 gives an independent row enumerator and
an exact equality test between that required set and an emitted ledger.  It
uses the V62 five-sector charge packet and includes three generations inside
each sector block.

\subsection{Declared minimal collar stencil}

Let \(X\) be a finite set of collar sites, \(E\) a set of undirected
nearest-neighbor edges, and
\[
 \rho_X:X\longrightarrow X,
 \qquad \rho_X^2=1,                                         \tag{80.1}
\]
a reflection preserving \(E\).  Let \(E^{\to}\) contain both orientations
of every edge.  Choose a reflection-closed tangent set
\(T\subseteq E^{\to}\) and a reflection-closed generator support
\(S\subseteq X\).

The charged left-chiral sector set is
\[
 \mathfrak F=\{Q,U^c,D^c,L,E^c\},                          \tag{80.2}
\]
with V62 multiplicities and integer hypercharges
\[
\begin{array}{c|ccccc}
f&Q&U^c&D^c&L&E^c\\ \hline
m_f&6&3&3&2&1\\
y_f&1&-4&2&-3&6.
\end{array}                                                  \tag{80.3}
\]
The representation and chirality labels are respectively
\[
 (\mathbf3,\mathbf2),
 (\overline{\mathbf3},\mathbf1),
 (\overline{\mathbf3},\mathbf1),
 (\mathbf1,\mathbf2),
 (\mathbf1,\mathbf1),
 \qquad \hbox{all left chiral}.                             \tag{80.4}
\]
For \(g\) generations, each site-sector carrier has dimension \(gm_f\).
Generation mixing is contained inside the emitted block; it is not
represented by duplicating ledger keys.

The minimal charged-fermion Yukawa incidence set is
\[
 \mathfrak Y=\bigl\{\{Q,U^c\},\{Q,D^c\},\{L,E^c\}\bigr\}.  \tag{80.5}
\]
Each unordered pair contributes both ordered matrix blocks.

\begin{definition}[Required provenance row]
\label{def:v18v70-required-row}
A row key is the seven-tuple
\[
 k=(\eta,A,f_1,f_2,x_1,x_2,\tau),                           \tag{80.6}
\]
where \(\eta\in\{F,R\}\) is the history side, \(A\) is the operator,
\(f_1,f_2\in\mathfrak F\), \(x_1,x_2\in X\), and \(\tau\) is the stencil
type.  The required set \(\mathcal K\) contains, on each side:
\begin{enumerate}[leftmargin=2.2em]
\item \((H,f,f,x,x,\mathrm{onsite})\) for every \(f,x\);
\item \((H,f,f,x,y,\mathrm{wilson\_hop})\) for every
      \(f,(x,y)\in E^{\to}\);
\item \((Q,f,f,x,y,\mathrm{link\_tangent})\) for every
      \(f,(x,y)\in T\);
\item \((\Omega,f,f,x,x,\mathrm{gauge\_generator})\) for every
      \(f,x\in S\); and
\item \((H,f_1,f_2,x,x,\mathrm{yukawa\_onsite})\) for every \(x\) and
      both orderings of each pair in \(\mathfrak Y\).
\end{enumerate}
The suppressed first component \(\eta\) takes both values.  A row has block
shape
\[
 \operatorname{shape}(k)=(gm_{f_1})\times(gm_{f_2}).         \tag{80.7}
\]
\end{definition}

\begin{remark}[Internal indices are payload data]
\label{rem:v18v70-internal-payload}
Spin, color, weak-isospin, and generation entries occur inside the block of
shape (80.7).  The ledger certifies where the block belongs.  Its numerical
contents and internal ordering must be supplied by the referenced payload
schema.
\end{remark}

\subsection{Reflection pairing and finite counts}

Let \(\bar F=R\) and \(\bar R=F\).  Define
\[
 \mathscr R(\eta,A,f_1,f_2,x_1,x_2,\tau)
 =(\bar\eta,A,f_1,f_2,\rho_Xx_1,\rho_Xx_2,\tau).             \tag{80.8}
\]

\begin{lemma}[Row reflection is a fixed-point-free involution]
\label{lem:v18v70-row-reflection}
If \(E,T,S\) are reflection closed, then
\[
 \mathscr R(\mathcal K)=\mathcal K,
 \qquad \mathscr R^2=1,                                    \tag{80.9}
\]
and no sided row is fixed by \(\mathscr R\).
\end{lemma}

\begin{proof}
Reflection closure sends every site, directed edge, tangent edge, and
supported site in Definition \ref{def:v18v70-required-row} to an object of
the same family.  It leaves sector and stencil labels unchanged and toggles
the side, so the image remains in \(\mathcal K\).  Applying (80.8) twice
uses \(\rho_X^2=1\) and toggles the side twice.  A fixed point is impossible
because \(F\ne R\).
\end{proof}

Let \(v=|X|\), \(e=|E|\), \(t=|T|\), \(s=|S|\),
\(F=|\mathfrak F|=5\), and \(Y=|\mathfrak Y|=3\).

\begin{proposition}[Required-row count]
\label{prop:v18v70-row-count}
The exact number of required sided rows is
\[
\boxed{
 |\mathcal K|
 =2\bigl(Fv+2Fe+Ft+Fs+2Yv\bigr).}                          \tag{80.10}
\]
The number of reflection orbits is \(|\mathcal K|/2\).
\end{proposition}

\begin{proof}
On one side the five families in Definition
\ref{def:v18v70-required-row} have respectively
\(Fv,2Fe,Ft,Fs,2Yv\) rows.  The forward and reflected sets are disjoint and
have the same size.  Lemma \ref{lem:v18v70-row-reflection} partitions their
union into two-element orbits.
\end{proof}

\subsection{Exact completeness criterion}

An emitted row contains its seven-tuple key, the two dimensions in (80.7),
a 64-hexadecimal-character payload digest, and the canonical string of its
declared reflection partner.  Let \(L\) be the finite emitted row list and
\(K(L)\) its set of keys.

\begin{theorem}[Finite ledger completeness]
\label{thm:v18v70-ledger-completeness}
Assume the verifier establishes all of the following:
\[
\begin{split}
 &\text{no two rows in }L\text{ have the same key},\\
 &K(L)=\mathcal K,\\
 &\operatorname{shape}_L(k)=\operatorname{shape}(k)
   \quad(k\in\mathcal K),\\
 &\operatorname{partner}_L(k)=\mathscr R(k)
   \quad(k\in\mathcal K).                                  \tag{80.11}
\end{split}
\]
Then every block allowed by the declared stencil occurs exactly once, no
forbidden block occurs, every block has the required carrier shape, and the
declared reflection pairing is a complete involution.
\end{theorem}

\begin{proof}
Set equality gives at least one emitted row for every allowed key and rules
out every forbidden key.  Duplicate rejection gives at most one row for
each key, hence exactly one.  The last two lines of (80.11) give the shapes
and partners.  Lemma \ref{lem:v18v70-row-reflection} proves that those
partners exhaust the list in disjoint two-element orbits.
\end{proof}

\begin{corollary}[Unique block assembly relative to the stencil]
\label{cor:v18v70-unique-assembly}
Fix a canonical order of sites and sector internal indices.  If every
ledger digest resolves to one block of the declared shape, a complete
ledger determines unique block matrices
\[
 (H,Q,\Omega),\qquad(H^R,Q^R,\Omega^R),                    \tag{80.12}
\]
by placing each payload in its keyed block and setting all blocks outside
\(\mathcal K\) to zero.
\end{corollary}

\begin{proof}
Theorem \ref{thm:v18v70-ledger-completeness} assigns exactly one correctly
shaped payload to every allowed block location.  Distinct locations in the
canonical block decomposition do not overlap.  Placement therefore exists
and is unique, while the declared support rule determines every remaining
block as zero.
\end{proof}

\begin{remark}[Assembly and adjoint validation are separate]
\label{rem:v18v70-adjoint-separate}
Completeness does not force the two orientations of a Wilson hop to be
adjoints, or the two Yukawa blocks to be adjoints.  After payload assembly,
the V68 algebraic gate must still verify that \(H,Q\) are Hermitian and
\(\Omega\) is anti-Hermitian.  This separation prevents a complete but
algebraically inconsistent ledger from entering the Ward compiler.
\end{remark}

\subsection{Executable verifier and negative controls}

The specification
\[
 \texttt{artifacts/sm\_v70\_sample\_provenance\_spec.json}              \tag{80.13}
\]
declares two sites exchanged by reflection, one undirected edge, both
directed tangent edges, full two-site generator support, three generations,
the five V62 sectors, and the three Yukawa pairs (80.5).  The emitted
specimen is
\[
 \texttt{artifacts/sm\_v70\_sample\_provenance\_ledger.json}.          \tag{80.14}
\]
The verifier is
\[
 \texttt{scripts/verify\_sm\_v70\_provenance\_ledger.py}.              \tag{80.15}
\]
It constructs \(\mathcal K\) from the geometry and hard-checked SM profile,
not from the emitted row list.

\begin{theorem}[Executed V70 completeness certificate]
\label{thm:v18v70-executed}
For the specimen (80.13)--(80.14),
\[
 |\mathcal K|=|K(L)|=104,
 \qquad
 |\mathcal K/\mathscr R|=52,                               \tag{80.16}
\]
with family counts
\[
\begin{array}{c|ccccc}
\tau&\mathrm{onsite}&\mathrm{wilson\_hop}&
\mathrm{link\_tangent}&\mathrm{gauge\_generator}&
\mathrm{yukawa\_onsite}\\ \hline
\#&20&20&20&20&24.
\end{array}                                                  \tag{80.17}
\]
There are no missing, extra, or duplicate rows; every shape and reflection
partner passes.  Four destructive controls are rejected: deletion of one
row, duplication of one row, alteration of one shape, and alteration of one
reflection partner.
\end{theorem}

\begin{proof}
Here \(v=2,e=1,t=2,s=2,F=5,Y=3\).  Formula (80.10) gives
\[
 2(5\cdot2+2\cdot5\cdot1+5\cdot2+5\cdot2+2\cdot3\cdot2)
 =104.                                                       \tag{80.18}
\]
The family summands give (80.17).  Exact set comparison returns empty
missing and extra sets after duplicate rejection.  Formula (80.7) checks
each shape and (80.8) checks each partner.  The verifier then mutates a copy
of the accepted ledger in each of the four stated ways and requires every
mutation to raise a rejection before emitting the certificate.
\end{proof}

The canonical result is
\[
 \texttt{artifacts/sm\_v70\_sample\_provenance\_result.json}.           \tag{80.19}
\]
The immutable ledger is
\[
\begin{array}{c|l}
\hbox{object}&\hbox{SHA--256}\\ \hline
\hbox{specification}&
\texttt{75B24A6AA170E7D373F18DCC6EDD3170D72A730E32295C4E88363080B227E523}\\
\hbox{emitted row ledger}&
\texttt{B011CA9494D94A921505C42947C3F1B4363E0B9861A69FADAB6B0DC9EBFA55F7}\\
\hbox{verifier}&
\texttt{EE380C9F64AF17884D29A92348CCDA40C6F21DA8BFFEB32283DD70DEA6152B4C}\\
\hbox{canonical result}&
\texttt{2BF5BE0948DDE96C3F2175FA9D07793D9CA5B9AFE5C1D69A866B784F6154C7E3}.
\end{array}                                                   \tag{80.20}
\]

\subsection{Composition with the interval Ward certificate}

\begin{theorem}[Finite provenance-to-Ward handoff]
\label{thm:v18v70-provenance-handoff}
At cutoff \(n\), suppose:
\begin{enumerate}[leftmargin=2.2em]
\item the declared stencil covers the full V67 collar compression of the
      selected physical Wilson--Yukawa kernel, tangent, and generator;
\item an emitted ledger passes Theorem
      \ref{thm:v18v70-ledger-completeness};
\item every digest is recomputed from a rational-rectangle payload with a
      canonical internal-index order;
\item assembly passes the V68 algebraic gate; and
\item the assembled rectangles give the V69 bound
      \(B_{p,n}^{\rm rect}\).
\end{enumerate}
Then the polynomial principal reversal residual of that selected physical
collar obeys
\[
 |R_{p,n}^{\rm phys}|\le B_{p,n}^{\rm rect}.                \tag{80.21}
\]
Adding the two V67 sign tails yields a valid exact-sign reversal enclosure.
\end{theorem}

\begin{proof}
The first two assumptions and Corollary
\ref{cor:v18v70-unique-assembly} show that the payload set covers every and
only every block of the physical collar.  Digest recomputation binds those
blocks to the checked bytes.  The algebraic gate makes the Ward functional
well typed.  Theorem \ref{thm:v18v69-rectangle-enclosure} then proves
(80.21), and (78.31) adds the sign tails.
\end{proof}

\begin{firewall}[Relative completeness and unresolved byte binding]
\label{fire:v18v70-relative-completeness}
The certificate proves completeness relative to the declared minimal
nearest-neighbor Wilson plus onsite charged-Yukawa stencil.  It does not
prove that this stencil is the kernel selected by the attached physical
programme; improved, clover, Majorana, right-neutrino, or other couplings
would require added row families.  The specimen payload digests are
deterministic placeholders derived from row keys.  The current verifier
checks their syntax but does not resolve files and recompute their contents.
Therefore Theorem \ref{thm:v18v70-provenance-handoff} is not yet instantiated
by the specimen.
\end{firewall}

\subsection{Acquisition and checkpoint boundary}

V70 supplies the missing finite coverage logic and makes omission of one
block mechanically visible.  The immediate upstream tasks are now to bind
digests to canonical interval-block bytes and to derive the required row
families from a selected literal Wilson--Yukawa formula rather than the
minimal declared specimen.  This is a checkpoint version: before choosing
between that byte-binding task and the independent global gap row, the next
step will inspect the active YM and NS notes for a reusable exact-image or
coercivity mechanism, then continue in V71.

\begin{assessment}[V70 acquisition and boundary]
\label{ass:v18v70}
V70 proves exact finite completeness, reflection-orbit closure, shape
correctness, and unique block assembly relative to a declared stencil.  The
104-row three-generation specimen and four negative controls execute those
claims.  A literal kernel-derived stencil, content-addressed block payloads,
and the global gap remain open.
\end{assessment}

\begin{record}[V70 append record]
\label{rec:v18v70}
V70 is appended to the validated V69 whose installed SHA--256 was
\texttt{FE50CA37199CD79388A3380348FAA70C3006DA9CC10FD39DB54B363080608BE3}.
After installation only V70 is the active numeric SM note; the frozen
\texttt{V100\_f2} through \texttt{V100\_f17} archives are unchanged.
\end{record}


\section{V71: content-addressed hybrid collar blocks after the V70 checkpoint}
\label{sec:v18v71}

V70 proved that a row ledger can be complete while leaving its digests as
unresolved strings.  The checkpoint audit suggests how to close that gap
without imposing a wasteful representation.  The current YM note certifies
exact coefficient ranks and keeps a hybrid of low-rank factors and dense
exceptions.  The current NS note likewise shows that a norm improvement may
exist only on the structured class actually used by the equation.  V71
therefore binds each provenance row to canonical block bytes, reconstructs
each center over \(\mathbb Q(i)\), certifies its exact rank, and permits
factor storage only when it is strictly smaller than dense storage.

\subsection{Compulsory companion audit}

At the V70 checkpoint the active companion notes were inspected read-only.
The YM task advanced during the inspection, so the audit used its final
active file
\[
 \texttt{Renewal\_Yang\_Mills\_Mass\_Gap\_Research\_Notes\_V12.tex},   \tag{81.1}
\]
whose SHA--256 was
\[
 \texttt{1A1D41A289DA206CCB88023306DD33743C46058F731345900765A724BA1DDCBE}.
                                                                  \tag{81.2}
\]
YM V12 proves an exact modular--structural rank sandwich, uses exact
Gaussian elimination in the cancellation cases, and selects factors for
344 of 436 coefficient images while retaining 92 dense images.  Its cyclic
trace formula contracts complete low-rank paths before scalar coefficient
accumulation.  The transferable rule is that rank compression must follow
an exact rank certificate and must admit dense exceptions.

The active NS file was
\[
 \texttt{Renewal\_NS\_Stability\_Research\_Notes\_V26.tex},            \tag{81.3}
\]
with SHA--256
\[
 \texttt{82C887F8C2C69E4F28FAD7C3DC8DE5B9099CB5A4BF431EBA1FFF3E91935978AD}.
                                                                  \tag{81.4}
\]
NS V26 computes norms on the actual rank-one tensor class: the first
derivative gains nothing, while the second derivative gains up to about
eleven percent on part of the radial range.  No NS theorem is imported into
the SM programme.  Its transferable rule is to measure a block on its used
input class and to accept a zero gain when the structured extremizer already
realizes the full norm.

\begin{decision}[V71 checkpoint decision]
\label{dec:v18v71-audit}
Bind the V70 rows to actual canonical payload objects and use exact-rank
hybrid storage.  Do not postulate low rank for every Standard-Model block.
Decode every accepted representation to the same dense rational-rectangle
semantics before applying V69.
\end{decision}

\subsection{Canonical payload objects}

Let \(k\in\mathcal K\) be a V70 row with shape \(m\times n\).  A V71
payload object contains:
\begin{enumerate}[leftmargin=2.2em]
\item the shape \((m,n)\);
\item either a dense center \(C\in M_{m\times n}(\mathbb Q(i))\), or
      factors \(U\in M_{m\times r}(\mathbb Q(i))\) and
      \(V\in M_{r\times n}(\mathbb Q(i))\);
\item the declared exact rank \(r\); and
\item nonnegative rational real and imaginary uncertainty radii, uniformly
      or entrywise specified.
\end{enumerate}
Sorted-key compact JSON followed by one newline is denoted
\(\operatorname{Can}(P)\).  Its content identifier is
\[
 d(P)=\operatorname{SHA256}(\operatorname{Can}(P)).          \tag{81.5}
\]
The block bundle is a finite dictionary whose key \(d\) points to a payload
\(P_d\).  A ledger row resolves precisely when
\[
 d_k\in\operatorname{dom}(\mathcal B),
 \qquad
 d(P_{d_k})=d_k.                                             \tag{81.6}
\]

\begin{lemma}[Unambiguous canonical decoding]
\label{lem:v18v71-canonical-decoding}
For a fixed JSON parser and the V71 schema, a payload byte string in
canonical form determines one shape, center encoding, rational uncertainty
rectangle, and declared rank.  Re-encoding the parsed object by
\(\operatorname{Can}\) reproduces the bytes used in (81.5).
\end{lemma}

\begin{proof}
Object keys are sorted, separators and the terminal newline are fixed, and
every rational is a JSON string parsed as one reduced fraction.  Arrays
retain order.  The schema makes the dense and factor encodings disjoint.
Thus parsing and canonical re-encoding are deterministic.
\end{proof}

\begin{remark}[Meaning of the digest]
\label{rem:v18v71-digest-meaning}
The verifier does not infer byte equality from a hash alone.  It retrieves
the concrete bundle object selected by the ledger key, canonicalizes that
object, and recomputes its digest.  Standard collision resistance is needed
only if the digest is used as a hostile external identity across different
stores.  The certificate's interval semantics always refer to the actual
object decoded in the checked bundle.
\end{remark}

\subsection{Exact hybrid representation}

All ranks below are ranks over \(\mathbb Q(i)\).

\begin{lemma}[Sound rank-factor center]
\label{lem:v18v71-rank-factor}
Let \(U\in M_{m\times r}(\mathbb Q(i))\) have column rank \(r\) and
\(V\in M_{r\times n}(\mathbb Q(i))\) have row rank \(r\).  Then
\[
 C=UV\quad\hbox{has rank }r.                                \tag{81.7}
\]
\end{lemma}

\begin{proof}
As a linear map \(V:\mathbb Q(i)^n\to\mathbb Q(i)^r\) is surjective, while
\(U:\mathbb Q(i)^r\to\mathbb Q(i)^m\) is injective.  Hence
\(\operatorname{im}(UV)=U(\mathbb Q(i)^r)\) has dimension \(r\).
\end{proof}

The center-entry costs are
\[
 c_{\rm fac}=r(m+n),
 \qquad
 c_{\rm den}=mn.                                            \tag{81.8}
\]

\begin{definition}[Hybrid selection gate]
\label{def:v18v71-hybrid-gate}
A factor payload is accepted only if exact elimination proves
\(\operatorname{rank}U=\operatorname{rank}V=r\) and
\[
 r(m+n)<mn.                                                  \tag{81.9}
\]
A dense payload is accepted only if exact elimination proves its declared
rank and
\[
 r(m+n)\ge mn.                                               \tag{81.10}
\]
\end{definition}

\begin{theorem}[Representation-independent rectangle semantics]
\label{thm:v18v71-representation-semantics}
Every payload passing Definition \ref{def:v18v71-hybrid-gate} decodes to a
unique center \(C\in M_{m\times n}(\mathbb Q(i))\) and an entrywise complex
rational rectangle
\[
 \mathbf C_{ij}
 =[\Re C_{ij}-r_{ij}^{\rm re},\Re C_{ij}+r_{ij}^{\rm re}]
 +i[\Im C_{ij}-r_{ij}^{\rm im},\Im C_{ij}+r_{ij}^{\rm im}]. \tag{81.11}
\]
The result is independent of whether the accepted center was stored densely
or as factors.
\end{theorem}

\begin{proof}
For a dense payload take the displayed center.  For a factor payload,
finite Gaussian-rational multiplication gives \(C=UV\), and Lemma
\ref{lem:v18v71-rank-factor} validates its declared rank.  Nonnegative
rational radii then determine (81.11).  Both branches end in exactly the
same dense-center-plus-radii object used by the V69 semantics.
\end{proof}

\begin{assessment}[Why the hybrid gate is conservative]
\label{ass:v18v71-hybrid-conservative}
The gate optimizes only stored center entries.  It makes no claim that a
factor representation reduces later block assembly or Ward multiplication;
those operations may densify.  Dense storage is retained whenever factors
do not strictly save entries.  This is the direct SM use of the YM V12
mixed-rank finding and the NS V26 actual-input warning.
\end{assessment}

\subsection{Content-addressed completeness}

\begin{theorem}[Resolved provenance bundle]
\label{thm:v18v71-resolved-bundle}
Suppose a V70 ledger passes Theorem
\ref{thm:v18v70-ledger-completeness}, and a V71 bundle satisfies:
\begin{enumerate}[leftmargin=2.2em]
\item every ledger digest resolves by (81.6);
\item every resolved payload has the ledger shape and passes Definition
      \ref{def:v18v71-hybrid-gate}; and
\item every bundle object is referenced by at least one ledger row.
\end{enumerate}
Then every required stencil row has exactly one checked payload reference,
every reference decodes to a rational rectangle of the correct shape, and
there are no unresolved or extraneous payload objects.
\end{theorem}

\begin{proof}
V70 set equality and duplicate rejection give exactly one ledger row per
required key.  The first two hypotheses resolve its digest to an accepted
object and Theorem \ref{thm:v18v71-representation-semantics} supplies the
rectangle.  The third hypothesis rules out objects that have no provenance
row.  These statements hold for every key in the finite required set.
\end{proof}

\begin{corollary}[Byte-bound interval Ward handoff]
\label{cor:v18v71-byte-bound-handoff}
Under the hypotheses of Theorem \ref{thm:v18v70-provenance-handoff}, replace
its unresolved payload assumption by a resolved V71 bundle.  Decoding and
block assembly then produce the exact V69 matrix rectangles, so
\[
 |R_{p,n}^{\rm phys}|\le B_{p,n}^{\rm rect}                 \tag{81.12}
\]
is bound to the canonical payload objects in that bundle.
\end{corollary}

\begin{proof}
Theorem \ref{thm:v18v71-resolved-bundle} supplies every correctly shaped
block rectangle.  Corollary \ref{cor:v18v70-unique-assembly} places those
blocks, the V68 gate checks the assembled algebraic types, and Theorem
\ref{thm:v18v69-rectangle-enclosure} proves (81.12).
\end{proof}

\subsection{Executed hybrid specimen}

The V71 verifier is
\[
 \texttt{scripts/verify\_sm\_v71\_content\_addressed\_blocks.py}.       \tag{81.13}
\]
It uses the V70 two-site, three-generation specification and emits a new
ledger and bundle:
\[
\begin{split}
 &\texttt{artifacts/sm\_v71\_bound\_provenance\_ledger.json},\\
 &\texttt{artifacts/sm\_v71\_block\_bundle.json}.                       \tag{81.14}
\end{split}
\]
The specimen uses exact rank-one factors for every profitable row.  The
four \(E^c\) onsite rows use dense full-rank \(3\times3\) centers because
factor storage would cost \(18\) entries rather than \(9\).  Gauge-tangent
payloads also carry uniform rational real and imaginary radii \(1/10000\);
other specimen radii vanish.

\begin{theorem}[Executed V71 content certificate]
\label{thm:v18v71-executed}
Exact execution gives
\[
\begin{array}{c|r}
\hbox{complete and resolved rows}&104\\
\hbox{unique canonical content objects}&20\\
\hbox{rank-factor rows}&100\\
\hbox{dense rows}&4\\
\hbox{dense center-entry baseline}&11232\\
\hbox{hybrid center entries}&1956\\
\hbox{saved center entries}&9276.
\end{array}                                                   \tag{81.15}
\]
Thus the row-multiplicity storage ratio is
\[
 \frac{1956}{11232}=\frac{163}{936}.                        \tag{81.16}
\]
Every digest, exact rank, shape, storage decision, and rational uncertainty
radius passes.  Four negative controls are rejected: a tampered payload, a
missing content object, an unresolved ledger digest, and a rank-deficient
factor.
\end{theorem}

\begin{proof}
The V70 enumerator again produces 104 required keys.  For each row the
generator constructs a canonical payload and replaces the placeholder
digest by (81.5).  The verifier independently recomputes every digest,
checks exact Gaussian elimination on the center or both factors, applies
(81.9)--(81.10), and sums the costs (81.8) with row multiplicity.  The
displayed counts and integer sums are its exact output.  It then mutates
copies in the four stated ways and requires rejection before emitting the
certificate.
\end{proof}

The canonical result is
\[
 \texttt{artifacts/sm\_v71\_content\_addressed\_result.json}.           \tag{81.17}
\]
The immutable digest ledger is
\[
\begin{array}{c|l}
\hbox{object}&\hbox{SHA--256}\\ \hline
\hbox{V70 specification}&
\texttt{75B24A6AA170E7D373F18DCC6EDD3170D72A730E32295C4E88363080B227E523}\\
\hbox{bound provenance ledger}&
\texttt{405C0FDF3F40D5525B60B651A3B755D9B5EFABFA1726DF058F30FF21F4D489BF}\\
\hbox{block bundle}&
\texttt{6ADAA757D5F6EF2695F59C3B3D3FB95AE6D1FBA5CB53F02228E5A7465897E80F}\\
\hbox{verifier}&
\texttt{6BF64A58856E2C2EB2ECCA76E68DB00E33805C872ED093F7E3967BC6FBDB6AD1}\\
\hbox{canonical result}&
\texttt{194B882251FD2FD0EA21B13AC12881CA90D701008FCEEF74F8694690CA584042}.
\end{array}                                                   \tag{81.18}
\]

\begin{firewall}[The bound specimen is still synthetic]
\label{fire:v18v71-synthetic}
The V71 payload bytes are generated specimens: their centers do not come
from the attached Wilson--Yukawa formula, and repeated rows deliberately
share content objects.  Content addressing proves which bytes were checked;
it does not make those bytes physical.  The bundle also lacks a canonical
internal-index label list, cross-row adjoint constraints, and an assembled
V68 result.  The global full-kernel gap remains independent and unresolved.
\end{firewall}

\subsection{Acquisition and next bottleneck}

V71 closes the placeholder-digest defect at the schema level and shows that
structure-aware storage composes with rational rectangles.  The next
upstream step is semantic assembly: give every block axis a canonical tuple
of site, generation, color, weak-isospin, and spin labels; verify that paired
Wilson and Yukawa blocks satisfy the required adjoint relations; then build
the six global collar matrices and run the V68--V69 gate on the decoded
bundle.  If the literal source still does not select its kernel formula,
that construction must remain a parameterized producer theorem rather than
silently choosing a discretization.

\begin{assessment}[V71 acquisition and boundary]
\label{ass:v18v71}
V71 records the compulsory YM/NS audit, proves exact hybrid payload
semantics and content resolution, and executes a fully resolved 104-row
specimen with four dense exceptions and four tamper controls.  Canonical
axis labels, physical block formulas, assembled adjoint checks, and the
global gap remain open.
\end{assessment}

\begin{record}[V71 append record]
\label{rec:v18v71}
V71 is appended to the validated V70 whose installed SHA--256 was
\texttt{9404A7B2CEF7DFCA90E525223E2F74D539AE7E0CA1430B33EF3611F579929EFC}.
After installation only V71 is the active numeric SM note; the frozen
\texttt{V100\_f2} through \texttt{V100\_f17} archives are unchanged.
\end{record}


\section{V72: canonical-axis assembly and an end-to-end exact collar gate}
\label{sec:v18v72}

V71 resolved every provenance row to checked payload bytes, but stopped
before the blocks were placed in global matrices.  V72 performs that
placement on a three-generation typed specimen.  It fixes a canonical axis
order, derives every reflected block from its forward partner by the stated
antiunitary parity, assembles all six matrices, and sends the resulting
90-dimensional card through the unmodified V68 exact gate.

\subsection{Canonical collar axes}

For a site \(x\), sector \(f\in\mathfrak F\), generation
\(a\in\{1,\ldots,g\}\), and representation component
\(c\in\{1,\ldots,m_f\}\), define the axis label
\[
 \iota=(x,f,a,c).                                           \tag{82.1}
\]
The finite carrier is
\[
 \mathcal I
 =\coprod_{x\in X}\coprod_{f\in\mathfrak F}
   \{x\}\times\{f\}\times\{1,\ldots,g\}
   \times\{1,\ldots,m_f\}.                                 \tag{82.2}
\]
V72 orders it by the declared site order, the sector order in the
specification, generation, and then representation component.  Consequently
\[
 |\mathcal I|=|X|g\sum_{f\in\mathfrak F}m_f.                \tag{82.3}
\]
For the two-site, three-generation V70 specimen,
\[
 |\mathcal I|=2\cdot3\cdot(6+3+3+2+1)=90.                 \tag{82.4}
\]
Each pair \((x,f)\) occupies one contiguous range of length \(gm_f\),
matching the V70 row-shape rule.

\begin{lemma}[Canonical block placement]
\label{lem:v18v72-block-placement}
Let a complete V70 ledger resolve to V71 payloads.  Placing the decoded
block for
\[
 (\eta,A,f_1,f_2,x_1,x_2,\tau)                              \tag{82.5}
\]
in the rectangle of rows \((x_1,f_1)\) and columns \((x_2,f_2)\) is
well-defined.  Within each fixed pair \((\eta,A)\), two distinct required
row keys never occupy the same block rectangle.
\end{lemma}

\begin{proof}
Equation (80.7) equals the lengths of the two canonical ranges, so the
decoded shape fits.  For \(H\), an onsite row has equal sectors and equal
sites, a Wilson hop has equal sectors and distinct adjacent sites, and a
Yukawa row has distinct sectors and equal sites.  Those cases are disjoint.
Within a case the sector and site labels determine the key uniquely.  The
operators \(Q\) and \(\Omega\) are stored in different global matrices; a
\(Q\) row has distinct adjacent sites and an \(\Omega\) row has equal sites.
Thus no two keys for the same global matrix select one block rectangle.
\end{proof}

\begin{theorem}[Unique typed center assembly]
\label{thm:v18v72-unique-typed-assembly}
Fix the canonical axis order (82.1)--(82.2).  A resolved V71 ledger uniquely
determines center matrices
\[
 H,Q,\Omega,H^R,Q^R,\Omega^R\in
 M_{|\mathcal I|}(\mathbb Q(i))                             \tag{82.6}
\]
by block placement and zero extension outside the declared stencil.
\end{theorem}

\begin{proof}
Theorem \ref{thm:v18v71-resolved-bundle} supplies one decoded center for
every required row.  Lemma \ref{lem:v18v72-block-placement} places those
centers in disjoint rectangles.  The V70 completeness theorem says that no
required rectangle is omitted, while the stencil declares every other
rectangle zero.  Hence existence and uniqueness follow.
\end{proof}

\subsection{Blockwise antiunitary reflection}

Define the real permutation matrix \(U_\rho\) on the canonical carrier by
\[
 U_\rho e_{(x,f,a,c)}=e_{(\rho_Xx,f,a,c)}.                  \tag{82.7}
\]
Then
\[
 U_\rho U_\rho^*=I,
 \qquad U_\rho^2=I.                                        \tag{82.8}
\]
Let
\[
 \epsilon_H=1,qquad\epsilon_Q=-1,
 \qquad\epsilon_\Omega=-1.                                 \tag{82.9}
\]
For every forward row key \(k\), V72 constructs the payload at its V70
partner \(\mathscr Rk\) by conjugating the center and multiplying it by
\(\epsilon_A\); the rational radii are retained.

\begin{theorem}[Blockwise-to-global reflection]
\label{thm:v18v72-global-reflection}
If the reflected partner of every forward block obeys the V72 payload rule,
then the assembled centers satisfy
\[
\boxed{
 H^R=U_\rho\overline H,U_\rho^*,\qquad
 Q^R=-U_\rho\overline Q,U_\rho^*,\qquad
 \Omega^R=-U_\rho\overline\Omega,U_\rho^*.}              \tag{82.10}
\]
\end{theorem}

\begin{proof}
Conjugation by \(U_\rho\) sends the block rectangle indexed by
\((x_1,f_1;x_2,f_2)\) to the rectangle indexed by
\((\rho_Xx_1,f_1;\rho_Xx_2,f_2)\).  This is exactly the row-key map (80.8).
The block value there is the conjugated forward center with the parity
(82.9).  Equality therefore holds on every required block.  Reflection
closure of the row set and common zero extension make it hold on every
remaining block as well.
\end{proof}

\begin{corollary}[Zero exact reflection defects]
\label{cor:v18v72-zero-defects}
For the V68 ideal tuple formed with \(\mathfrak C(A)=
U_\rho\overline A U_\rho^*\),
\[
 \Delta H=\Delta Q=\Delta\Omega=0.                         \tag{82.11}
\]
\end{corollary}

\begin{proof}
Compare (82.10) with (78.6)--(78.7).
\end{proof}

\subsection{Adjoint consistency of the specimen}

The forward specimen chooses real rank-one onsite blocks, equal real blocks
on the two orientations of each Wilson edge, equal transposed real blocks on
the two orientations of every Yukawa pair, real equal blocks on the two
orientations of \(Q\), and imaginary diagonal blocks for \(\Omega\).

\begin{lemma}[Assembled algebraic types]
\label{lem:v18v72-assembled-types}
The V72 centers obey
\[
 H^*=H,qquad Q^*=Q,qquad\Omega^*=-\Omega,                 \tag{82.12}
\]
and the reflected centers obey the corresponding identities.
\end{lemma}

\begin{proof}
Every onsite \(H\) block is Hermitian.  Oppositely oriented Wilson blocks
are real transposes of one another, as are the two ordered Yukawa blocks.
This proves the first identity blockwise.  The two real tangent orientations
are transposes, proving the second.  The generator is site diagonal and
purely imaginary with Hermitian real coefficient, proving the third.  The
relations (82.10) transport adjoints to the reflected centers.
\end{proof}

The point statement does not assert that independently chosen points from
all uncertainty rectangles satisfy the same adjoint relations.  The actual
physical producer must correlate partner uncertainties; the entrywise V69
box remains a sound over-enclosure once it contains that typed point.

\subsection{An explicit nonzero Ward reading}

Use the V68 polynomial
\[
 p(x)=\frac{x^2}{25},
 \qquad
 DP_{p,H}[Q]=-\frac1{50}(QH+HQ).                             \tag{82.13}
\]
Only the first representation component of each sector is active in the
specimen Ward trace.  On its two site coordinates, the sector-diagonal
blocks are
\[
 H_f=\begin{pmatrix}1&-1\\-1&1\end{pmatrix},qquad
 Q_f=\begin{pmatrix}0&2\\2&0\end{pmatrix},qquad
 \Omega_f=iI_2.                                             \tag{82.14}
\]
The off-sector Yukawa part of \(H\) does not contribute to this trace:
after multiplication by sector-diagonal \(Q\) it remains off sector, and
multiplication by sector-diagonal \(\Omega\) cannot create a diagonal
sector block.

\begin{proposition}[Exact forward value]
\label{prop:v18v72-forward-value}
For the assembled 90-dimensional center,
\[
 {\cal W}_p(H,Q,\Omega)=\frac45.                            \tag{82.15}
\]
\end{proposition}

\begin{proof}
From (82.14),
\[
 Q_fH_f+H_fQ_f
 =\begin{pmatrix}-4&4\\4&-4\end{pmatrix},
 \quad
 DP_{p,H_f}[Q_f]
 =\begin{pmatrix}2/25&-2/25\\-2/25&2/25\end{pmatrix}.      \tag{82.16}
\]
Its trace against \(iI_2\) is \(4i/25\), so the Ward value is \(4/25\)
per sector.  There are five sectors and every other axis gives zero, yielding
(82.15).
\end{proof}

\begin{theorem}[End-to-end V68 parity gate]
\label{thm:v18v72-end-to-end}
The assembled exact card passes the V68 gate and returns
\[
 W_p^{\rm F}=\frac45,qquad
 W_p^{0,R}=W_p^R=-\frac45,qquad
 R_p^{\rm card}=0.                                         \tag{82.17}
\]
All three exact reflected defects are zero.
\end{theorem}

\begin{proof}
Lemma \ref{lem:v18v72-assembled-types} and (82.8) pass the V68 algebraic
checks.  Proposition \ref{prop:v18v72-forward-value} gives the forward
reading.  Since \(\epsilon_Q\epsilon_\Omega=1\), Theorem
\ref{thm:v18v68-ideal-parity} gives the ideal reflected value \(-4/5\).
Corollary \ref{cor:v18v72-zero-defects} makes the measured and ideal tuples
identical, proving the remaining statements.
\end{proof}

\subsection{Executable artifacts}

The assembly verifier is
\[
 \texttt{scripts/verify\_sm\_v72\_typed\_assembly.py}.                  \tag{82.18}
\]
It generates and checks
\[
\begin{split}
 &\texttt{artifacts/sm\_v72\_typed\_provenance\_ledger.json},\\
 &\texttt{artifacts/sm\_v72\_typed\_block\_bundle.json},\\
 &\texttt{artifacts/sm\_v72\_axis\_manifest.json},\\
 &\texttt{artifacts/sm\_v72\_assembled\_exact\_card.json},\\
 &\texttt{artifacts/sm\_v72\_typed\_assembly\_result.json}.            \tag{82.19}
\end{split}
\]
The executed certificate reports 90 axes, 104 complete provenance rows,
24 resolved content objects, the three adjoint-type checks, unitarity of the
reflection carrier, the three zero defects, and (82.17).

The immutable SHA--256 ledger is
\[
\begin{array}{c|l}
\hbox{object}&\hbox{SHA--256}\\ \hline
\hbox{V70 specification}&
\texttt{75B24A6AA170E7D373F18DCC6EDD3170D72A730E32295C4E88363080B227E523}\\
\hbox{typed ledger}&
\texttt{41A244891862A934FFC689606FBDDCE1B0E14A818D70969987D8EE054F60A69C}\\
\hbox{typed bundle}&
\texttt{BFC79D8DFC3C0B7659700CEE815487E98D54D453A2BF558F15E5F12E177430A5}\\
\hbox{axis manifest}&
\texttt{F8151EE099BF33AF727557C2E439CBA09AF3D970048436BC6E9771B414B9F634}\\
\hbox{assembled exact card}&
\texttt{90B3E4EA5192A79C71F4B52D081D6721E95D1277E746A4DDA8F20845467F7D3D}\\
\hbox{verifier}&
\texttt{15ACFFB9DF173566D26BCB8B4DB3924DE0D900B10ADE33BF1C0912424043ED39}\\
\hbox{canonical result}&
\texttt{616B6CD6B29FDBB5E216EB280D56360D670789FD3C99D70D97429E89BF0272B9}.
\end{array}                                                   \tag{82.20}
\]

\begin{firewall}[Typed execution is not a physical Wilson card]
\label{fire:v18v72-synthetic-card}
The coefficients \(1,-1,2,i,3\) in the specimen are synthetic and were
chosen to exercise every row family and produce a nonzero Ward reading.
The generic component label \(c\) does not yet decompose into canonical
color, weak-isospin, and spin labels.  The attached programme still does not
select a literal Wilson mass, lattice spacing, spin convention, link
history, Higgs history, or reflection card.  V72 validates the complete
assembly interface, not those missing physical data.  Its rational
uncertainty radii are not yet assembled into a 90-dimensional V69 interval
result, and no global spectral gap is proved.
\end{firewall}

\subsection{Acquisition and next bottleneck}

The finite algebraic chain from row enumeration to exact Ward parity is now
executable end to end.  The next local step is to lift the same assembly from
centers to correlated rational rectangles and obtain a 90-dimensional V69
residual enclosure without expanding invalid independent adjoint choices.
Upstream of that lies the actual physical producer: a selected Wilson--Yukawa
formula must assign the component axes and block coefficients.  The global
gap remains a separate theorem and cannot be inferred from this collar.

\begin{assessment}[V72 acquisition and boundary]
\label{ass:v18v72}
V72 proves canonical placement, blockwise-to-global reflection, and typed
assembly, then executes the full V70--V71--V68 chain on a 90-dimensional
three-generation specimen with nonzero Ward value.  The local interface is
now tested as a whole; physical coefficients, correlated interval assembly,
and the global gap remain open.
\end{assessment}

\begin{record}[V72 append record]
\label{rec:v18v72}
V72 is appended to the validated V71 whose installed SHA--256 was
\texttt{18AF673F1C376111524A8C6962E86853B00C4A3BB551C1B8ECA2867C67E225C9}.
After installation only V72 is the active numeric SM note; the frozen
\texttt{V100\_f2} through \texttt{V100\_f17} archives are unchanged.
\end{record}


\section{V73: correlated parameter compilation before interval evaluation}
\label{sec:v18v73}

V72 assembled exact centers, while the V69 rectangle compiler treated each
matrix entry as independently variable.  Physical forward and reflected
blocks are not independent: a calibration or quadrature error can occur once
upstream and then be transported to both histories by the same antiunitary
map.  Splitting that one variable into two interval copies destroys an exact
cancellation.  V73 keeps shared variables symbolic in a sparse polynomial
ring, compiles the signed reversal polynomial coefficientwise, and evaluates
an interval only afterward.

\subsection{Sparse polynomial matrix semantics}

Let
\[
 \mathscr P_s=\mathbb Q(i)[z_1,\ldots,z_s].                 \tag{83.1}
\]
A sparse scalar polynomial is a finite dictionary
\(\alpha\mapsto c_\alpha\), \(\alpha\in\mathbb N^s\), with zero
coefficients deleted.  A sparse matrix over \(\mathscr P_s\) is in turn a
finite dictionary from matrix positions to sparse scalar polynomials.  Its
addition and multiplication are exact dictionary addition and convolution.

For \(z\in\mathbb R^s\), evaluation is
\[
 \operatorname{ev}_z\left(\sum_\alpha c_\alpha z^\alpha\right)
 =\sum_\alpha c_\alpha z^\alpha.                            \tag{83.2}
\]

\begin{lemma}[Evaluation homomorphism]
\label{lem:v18v73-evaluation}
Entrywise evaluation is a unital algebra homomorphism
\[
 \operatorname{ev}_z:M_N(\mathscr P_s)
 \longrightarrow M_N(\mathbb C).                           \tag{83.3}
\]
It commutes with trace and with coefficientwise conjugation when the
parameters \(z_j\) are real.
\end{lemma}

\begin{proof}
Finite coefficient convolution is precisely the coefficient rule for a
polynomial product.  Therefore evaluation preserves sums, products, and the
identity.  Trace is a finite sum of diagonal entries.  Real parameters are
fixed by conjugation, so only the Gaussian-rational coefficients are
conjugated.
\end{proof}

For a real rational polynomial \(p\), use the same finite formula (78.8) in
\(M_N(\mathscr P_s)\) and define
\[
 \mathfrak W_p(\mathbf H,\mathbf Q,\boldsymbol\Omega)
 =-i\operatorname{Tr}
 \bigl(DP_{p,\mathbf H}[\mathbf Q]\boldsymbol\Omega\bigr)
 \in\mathscr P_s.                                          \tag{83.4}
\]

\begin{theorem}[Symbolic Ward correctness]
\label{thm:v18v73-symbolic-correctness}
For every real parameter point \(z\),
\[
 \operatorname{ev}_z\mathfrak W_p
 (\mathbf H,\mathbf Q,\boldsymbol\Omega)
 ={cal W}_p
 (\operatorname{ev}_z\mathbf H,
  \operatorname{ev}_z\mathbf Q,
  \operatorname{ev}_z\boldsymbol\Omega).                  \tag{83.5}
\]
Consequently coefficientwise subtraction of sector, history, and reflection
partners before any interval operation preserves the exact pointwise Ward
residual.
\end{theorem}

\begin{proof}
Formula (78.8) is a finite sum of matrix products and rational multiples.
Apply Lemma \ref{lem:v18v73-evaluation} to it, then to the product and trace
in (83.4).  Signed finite sums are also preserved.
\end{proof}

\subsection{Shared antiunitary parameters}

Let \(t\) be a shared real producer uncertainty.  Suppose a forward family
is
\[
 H(t),\qquad Q(t),\qquad\Omega(t),                           \tag{83.6}
\]
and its ideal reflected family uses the same value of \(t\):
\[
 H^0(t)=\mathfrak C(H(t)),\quad
 Q^0(t)=\epsilon_Q\mathfrak C(Q(t)),\quad
 \Omega^0(t)=\epsilon_\Omega\mathfrak C(\Omega(t)).         \tag{83.7}
\]

\begin{theorem}[Coefficientwise ideal parity]
\label{thm:v18v73-coefficientwise-parity}
In \(\mathbb Q(i)[t]\),
\[
\boxed{
 \mathfrak W_p(H^0(t),Q^0(t),\Omega^0(t))
 +\epsilon_Q\epsilon_\Omega
  \mathfrak W_p(H(t),Q(t),\Omega(t))=0.}                   \tag{83.8}
\]
Thus every coefficient of the shared parameter cancels separately.
\end{theorem}

\begin{proof}
The proof of Theorem \ref{thm:v18v68-ideal-parity} uses only finite
polynomial operations, conjugation of coefficients, trace cyclicity, and
the fact that the parameters are real.  It therefore holds in the polynomial
ring.  Alternatively evaluate the left side at every real \(t\) and use
Theorem \ref{thm:v18v68-ideal-parity}; a complex polynomial vanishing on all
real \(t\) is the zero polynomial.
\end{proof}

Now let \(e\) be an independent reflected defect parameter and write
\[
 (H^R,Q^R,\Omega^R)(t,e)
 =(H^0,Q^0,\Omega^0)(t)+\mathcal D(t,e).                    \tag{83.9}
\]
The compiled polynomial residual is
\[
 \mathfrak R_p(t,e)
 =\mathfrak W_p(H^R,Q^R,\Omega^R)(t,e)
 +\epsilon_Q\epsilon_\Omega
  \mathfrak W_p(H,Q,\Omega)(t).                             \tag{83.10}
\]
Theorem \ref{thm:v18v73-coefficientwise-parity} guarantees that any
coefficient surviving in (83.10) comes from the measured defect interaction,
not from duplicating the ideal shared uncertainty.

\subsection{Interval evaluation after compilation}

For rational parameter intervals \(I_j\), define the natural monomial
enclosure
\[
 \mathbf E_I\left(\sum_\alpha c_\alpha z^\alpha\right)
 =\sum_\alpha c_\alpha I_1^{\alpha_1}\cdots I_s^{\alpha_s}, \tag{83.11}
\]
using rational interval operations.

\begin{proposition}[Correlated parameter enclosure]
\label{prop:v18v73-parameter-enclosure}
If \(z_j\in I_j\), then
\[
 \mathfrak R_p(z)\in\mathbf E_I(\mathfrak R_p).             \tag{83.12}
\]
When \(\mathfrak R_p\) is affine in each independent parameter, (83.11) is
the exact range on a rectangular parameter box.
\end{proposition}

\begin{proof}
Each power and product interval contains the corresponding monomial value;
finite interval addition proves (83.12).  An affine function attains its
extrema at box vertices, and interval scalar multiplication and addition
produce those same extrema.
\end{proof}

\begin{assessment}[Correlation is algebra, not a smaller guessed radius]
\label{ass:v18v73-correlation-algebra}
The improvement occurs because the same symbol appears in the forward and
reflected polynomials and cancels before interval evaluation.  Assigning two
independent copies enlarges the admissible uncertainty model.  A producer
may use the correlated bound only when its provenance certifies that both
entries derive from the same upstream variable and the declared
antiunitary transformation.
\end{assessment}

\subsection{Executed V72 parameter family}

Let \(T=Q/2\) for the V72 exact center, so that on each active two-site
sector block \(T\) has off-diagonal entries one.  Let
\[
 T^R=-U_\rho\overline T,U_\rho^*.                          \tag{83.13}
\]
Keep \(H,\Omega,H^R,\Omega^R\) fixed and set
\[
 Q^{F}(t)=Q+tT,
 \qquad
 Q^{R}(t,e)=Q^R+tT^R+eT^R.                                 \tag{83.14}
\]
The variable \(t\) is shared ideal transport; \(e\) is the measured
reflected mismatch.

\begin{theorem}[Exact compiled two-parameter polynomial]
\label{thm:v18v73-executed-polynomial}
For \(p(x)=x^2/25\), exact sparse compilation on the 90-dimensional V72
card gives
\[
\begin{split}
 W_p^F(t)&=\frac45+\frac25t,\\
 W_p^R(t,e)&=-\frac45-\frac25t-\frac25e,\\
 \mathfrak R_p(t,e)&=-\frac25e.                            \tag{83.15}
\end{split}
\]
In particular, the constant and shared-\(t\) coefficients vanish exactly in
the residual dictionary.
\end{theorem}

\begin{proof}
Replacing the off-diagonal value \(2\) in (82.14) by \(2+t\), the calculation
in (82.16) gives \(2(2+t)/25\) per sector and therefore
\(2(2+t)/5\) over five sectors.  This is the first line.  Ideal reflection
negates it by (83.8).  The additional reflected coefficient \(eT^R\) has
the same linear response as \(tT^R\), giving the second line.  Subtraction
proves the third.  The sparse executable independently convolves the exact
matrix-polynomial dictionaries and verifies these three dictionaries.
\end{proof}

Take
\[
 t\in\left[-\frac1{10000},\frac1{10000}\right],
 \qquad
 e\in\left[-\frac1{100000},\frac1{100000}\right].          \tag{83.16}
\]

\begin{corollary}[Twenty-one-fold correlation gain]
\label{cor:v18v73-correlation-gain}
The compiled correlated enclosure is
\[
 \mathfrak R_p(t,e)
 \in\left[-\frac1{250000},\frac1{250000}\right].           \tag{83.17}
\]
If the forward and reflected occurrences of \(t\) are instead bounded as
independent variables, the triangle bound is
\[
 \frac25\left(\frac1{10000}+\frac1{10000}
                    +\frac1{100000}\right)
 =\frac{21}{250000}.                                       \tag{83.18}
\]
Thus exact correlation improves this executed principal bound by a factor
of \(21\).
\end{corollary}

\begin{proof}
Evaluate the last line of (83.15) on (83.16) to obtain (83.17).  Without the
shared symbol, the two \(t\)-coefficients and the \(e\)-coefficient are
charged separately, giving (83.18).
\end{proof}

\subsection{Executable provenance}

The exact sparse compiler
\[
 \texttt{scripts/verify\_sm\_v73\_correlated\_parameter\_card.py}       \tag{83.19}
\]
reads the V72 assembled card.  Its active sparse matrix counts are
\[
 |\operatorname{supp}H|=36,
 \quad |\operatorname{supp}Q|=10,
 \quad |\operatorname{supp}\Omega|=10,
 \quad |\operatorname{supp}T|=10.                          \tag{83.20}
\]
It emits
\[
 \texttt{artifacts/sm\_v73\_correlated\_parameter\_result.json}.       \tag{83.21}
\]
The immutable hashes are
\[
\begin{array}{c|l}
\hbox{object}&\hbox{SHA--256}\\ \hline
\hbox{V72 assembled card}&
\texttt{90B3E4EA5192A79C71F4B52D081D6721E95D1277E746A4DDA8F20845467F7D3D}\\
\hbox{compiler}&
\texttt{844654B1661286144E6E716136AE0C9B39D8759E0F449FE3941664FE0AC6AF97}\\
\hbox{canonical result}&
\texttt{14EAF5D6E6A6CC1FF5C146828C4A0888F934F2A36855728D46824E9E67D804C9}.
\end{array}                                                   \tag{83.22}
\]

\begin{theorem}[Correlated exact-sign enclosure]
\label{thm:v18v73-correlated-sign}
For a physical parameter card satisfying the V73 provenance assumptions,
let \(B_{p,n}^{\rm corr}\) be the maximum modulus of the compiled interval
(83.11).  Then
\[
 \Delta_{{\rm sign},n}^{\rm rev}
 \le B_{p,n}^{\rm corr}
 +S_{m,n}^{F}+S_{m,n}^{R}.                                  \tag{83.23}
\]
If all three terms tend to zero along the V67 cofinal collars, the exact-sign
reversal leakage vanishes.
\end{theorem}

\begin{proof}
Theorem \ref{thm:v18v73-symbolic-correctness} and Proposition
\ref{prop:v18v73-parameter-enclosure} bound the polynomial principal term.
Add the two V67 sign-functional tails exactly as in (78.31).  The limiting
statement follows by addition.
\end{proof}

\begin{firewall}[Synthetic correlation and unresolved sign tails]
\label{fire:v18v73-synthetic-correlation}
The variables \(t,e\) in (83.14) are an executed synthetic uncertainty
model.  A physical producer must identify its actual primitive uncertain
quantities and prove every reuse map.  Coefficientwise cancellation cannot
be assumed merely because two marginal intervals coincide.  V73 controls
the polynomial principal trace; it does not supply the full-kernel gap or
the V67 sign tails.
\end{firewall}

\subsection{Acquisition and next bottleneck}

V73 replaces independent-entry overcharging by a sparse, exact correlation
language.  The next local producer theorem should derive its primitive
variables from link and Higgs interval data and compile every block as a
polynomial or rational enclosure in those variables.  The principal global
bottleneck is now sharper: to turn the polynomial certificate into an exact
overlap certificate one must prove a uniform admissibility-to-gap bound for
the selected full Wilson kernel.  Going upstream to that spectral row is
preferable if the physical block formula remains unselected.

\begin{assessment}[V73 acquisition and boundary]
\label{ass:v18v73}
V73 proves exact sparse parameter semantics and coefficientwise
antiunitary cancellation, then executes a two-parameter 90-dimensional card
where correct correlation improves the residual bound by a factor of 21.
Physical primitive-variable provenance and the global gap remain open.
\end{assessment}

\begin{record}[V73 append record]
\label{rec:v18v73}
V73 is appended to the validated V72 whose installed SHA--256 was
\texttt{277EDDFF7E5D565D15BC845147180FCCAC929DCD0B1CD26096D01904894B3D60}.
After installation only V73 is the active numeric SM note; the frozen
\texttt{V100\_f2} through \texttt{V100\_f17} archives are unchanged.
\end{record}


\section{V74: an admissibility-to-gap route for the Hermitian Wilson kernel}
\label{sec:v18v74}

V73 leaves the sign-functional tails behind one global scalar: a uniform
gap for the full Hermitian Wilson kernel.  The Grand Tensor shell assumes
such a window but does not derive it, and the preceding notes use it only as
a hypothesis.  V74 goes upstream.  It proves the free flat-connection gap in
the physical Wilson mass interval, a representation-uniform perturbation
bound in link distance from a flat reference, and a finite tree-gauge lemma
which converts plaquette error into link error whenever explicit filling
words are available.

\subsection{Dimensionless Wilson convention}

Let \(X\) be a finite \(d\)-dimensional periodic lattice, let
\(\gamma_1,\ldots,\gamma_d\) be Hermitian Euclidean gamma matrices, and let
\(T_\mu^U\) be the unitary covariant forward shift in a fixed unitary gauge
representation.  Use Wilson parameter one and define
\[
 D_W(U)
 =dI-\frac12\sum_{\mu=1}^d
 \left((I-\gamma_\mu)T_\mu^U
 +(I+\gamma_\mu)(T_\mu^U)^*\right),                         \tag{84.1}
\]
\[
 H_m(U)=\gamma_{d+1}(D_W(U)-mI).                            \tag{84.2}
\]
The gamma relations make \(H_m(U)\) Hermitian.  The normalization is
dimensionless; a physical factor \(a^{-1}\) multiplies both gap and norm.

Assume first that \(U^0\) is gauge equivalent to a translation-invariant
flat connection whose represented torus holonomies commute.  Simultaneous
diagonalization of those holonomies and Fourier transformation reduce each
color channel to shifted momenta \(p\).  Then
\[
 H_m(U^0;p)^2
 =\left[
 \sum_{\mu=1}^d\sin^2p_\mu
 +\left(\sum_{\mu=1}^d(1-\cos p_\mu)-m\right)^2
 \right]I.                                                  \tag{84.3}
\]

\begin{theorem}[Exact flat Wilson gap]
\label{thm:v18v74-flat-gap}
For every momentum and every such flat connection,
\[
 |H_m(U^0;p)|
 \ge g_{\rm flat}(m)
 :=\min_{0\le k\le d}|2k-m|.                               \tag{84.4}
\]
In particular, if \(0<m<2\),
\[
\boxed{
 g_{\rm flat}(m)=\min\{m,2-m\}>0.}                         \tag{84.5}
\]
One also has
\[
 \|H_m(U^0)\|_{\rm op}\le
 M_{\rm flat}(m):=\max_{0\le k\le d}|2k-m|.               \tag{84.6}
\]
\end{theorem}

\begin{proof}
Put \(t_\mu=1-\cos p_\mu\in[0,2]\).  Since
\(\sin^2p_\mu=2t_\mu-t_\mu^2\), the scalar in (84.3) is
\[
 F_m(t)=m^2+2(1-m)\sum_\mu t_\mu
       +2\sum_{\mu<\nu}t_\mu t_\nu.                       \tag{84.7}
\]
This polynomial is affine in each coordinate separately.  Its minimum and
maximum on the cube \([0,2]^d\) are therefore attained at vertices.  A
vertex with exactly \(k\) coordinates equal to two has value
\[
 F_m=(2k-m)^2.                                               \tag{84.8}
\]
Taking square roots proves (84.4) and (84.6).  For \(0<m<2\), the two
nearest even integers are zero and two, proving (84.5).
\end{proof}

\begin{remark}[Twisted flat holonomies do not reduce the bound]
\label{rem:v18v74-flat-twists}
The holonomies only shift the allowed momenta.  The proof bounds every point
of the continuous momentum cube, so it is uniform in lattice size, boundary
twist, and the commuting flat holonomy eigenvalues.
\end{remark}

\subsection{Link-distance perturbation}

For two gauge fields in the same representation, put
\[
 \delta(U,U^0)=max_{x,\mu}
 \|\rho(U_{x,\mu})-\rho(U^0_{x,\mu})\|_{\rm op}.             \tag{84.9}
\]

\begin{lemma}[Wilson hopping perturbation]
\label{lem:v18v74-hopping-perturbation}
For the convention (84.1),
\[
 \|H_m(U)-H_m(U^0)\|_{\rm op}
 \le 2d\,\delta(U,U^0).                                    \tag{84.10}
\]
\end{lemma}

\begin{proof}
The mass and diagonal Wilson terms cancel.  For each \(\mu\), the forward
and backward differences are multiplied by
\((I-\gamma_\mu)/2\) and \((I+\gamma_\mu)/2\), respectively.  Both have
operator norm one.  Covariant shifts and \(\gamma_{d+1}\) are unitary.
There are two hopping differences in each of \(d\) directions, so the
triangle inequality gives (84.10).
\end{proof}

\begin{theorem}[Flat-reference Wilson window]
\label{thm:v18v74-flat-reference-window}
Let \(U^0\) satisfy the flat hypotheses of Theorem
\ref{thm:v18v74-flat-gap}.  If
\[
 2d\,\delta(U,U^0)<g_{\rm flat}(m),                         \tag{84.11}
\]
then
\[
\boxed{
 \operatorname{dist}(0,\operatorname{spec}H_m(U))
 \ge g_{\rm flat}(m)-2d\,\delta(U,U^0)>0,}                 \tag{84.12}
\]
and
\[
 \|H_m(U)\|_{\rm op}
 \le M_{\rm flat}(m)+2d\,\delta(U,U^0).                    \tag{84.13}
\]
\end{theorem}

\begin{proof}
For self-adjoint matrices, the Hausdorff distance between spectra is at most
the operator norm of their difference.  Combine this with (84.4) and
Lemma \ref{lem:v18v74-hopping-perturbation} to obtain (84.12).  The ordinary
operator-norm triangle inequality and (84.6) give (84.13).
\end{proof}

For the complete charged packet define a representation-uniform flat
distance
\[
 \delta_{\rm SM}(U,U^0)
 =\max_{f\in\mathfrak F}\max_{x,\mu}
 \|\rho_f(U_{x,\mu})-\rho_f(U^0_{x,\mu})\|_{\rm op}.         \tag{84.14}
\]

\begin{corollary}[Common five-sector gap]
\label{cor:v18v74-five-sector-gap}
If sector \(f\) has Wilson mass \(m_f\in(0,2)\), put
\[
 \bar g_0=\min_f\min\{m_f,2-m_f\}.                         \tag{84.15}
\]
Then all five Hermitian Wilson kernels have the common lower bound
\[
 \bar g=\bar g_0-2d\,\delta_{\rm SM}(U,U^0),                \tag{84.16}
\]
provided the right side is positive.
\end{corollary}

\begin{proof}
Apply Theorem \ref{thm:v18v74-flat-reference-window} in each unitary
representation and take the minimum.
\end{proof}

\subsection{From plaquettes to links on a fillable complex}

The perturbation theorem needs link distance.  A plaquette hypothesis can
supply it only after topology and gauge fixing are handled explicitly.

\begin{lemma}[Unitary word telescoping]
\label{lem:v18v74-unitary-word}
For unitary matrices \(A_1,\ldots,A_r\),
\[
 \left\|A_1A_2\cdots A_r-I\right\|_{\rm op}
 \le\sum_{j=1}^r\|A_j-I\|_{\rm op}.                        \tag{84.17}
\]
The same estimate holds when any \(A_j\) is replaced by a unitary conjugate
or inverse of a plaquette holonomy.
\end{lemma}

\begin{proof}
Use the exact telescoping identity
\[
 A_1\cdots A_r-I
 =\sum_{j=1}^r A_1\cdots A_{j-1}(A_j-I),                   \tag{84.18}
\]
with the missing product equal to the identity.  Left multiplication by a
unitary preserves norm.  Conjugation preserves
\(\|A-I\|\), and \(\|A^{-1}-I\|=\|A-I\|\).
\end{proof}

Let \(G\) be the one-skeleton of a finite connected cell complex and
\(\mathcal T\) a spanning tree.  Gauge-fix all tree links to the identity.
For every chord \(e\notin\mathcal T\), its tree fundamental loop has one
remaining link, namely the gauged chord.  Assume a declared filling word
expresses that loop holonomy as a product of at most \(A_*\) conjugates of
oriented plaquette holonomies or their inverses.

\begin{theorem}[Finite tree-gauge plaquette conversion]
\label{thm:v18v74-tree-gauge}
If in a unitary representation
\[
 \max_p\|\rho(U_p)-I\|_{\rm op}\le\varepsilon,             \tag{84.19}
\]
then in the stated tree gauge
\[
 \max_e\|\rho(U_e)-I\|_{\rm op}\le A_*\varepsilon.        \tag{84.20}
\]
Consequently, if
\[
 2dA_*\varepsilon<g_{\rm flat}(m),                         \tag{84.21}
\]
then
\[
 \operatorname{dist}(0,\operatorname{spec}H_m(U))
 \ge g_{\rm flat}(m)-2dA_*\varepsilon.                    \tag{84.22}
\]
\end{theorem}

\begin{proof}
Tree links already equal the identity.  For a chord, the tree paths in its
fundamental loop also equal the identity, so its gauged link matrix is the
fundamental-loop holonomy, up to inversion.  Apply Lemma
\ref{lem:v18v74-unitary-word} to its declared filling word to obtain
(84.20).  Theorem \ref{thm:v18v74-flat-reference-window} with
\(U^0=I\) proves (84.22).
\end{proof}

\begin{corollary}[A cofinal gap condition]
\label{cor:v18v74-cofinal-gap}
Suppose \(m_{f,n}\) remain in a compact subinterval of \((0,2)\), with
common free gap \(\bar g_0>0\).  If the five representations share filling
stocks \(A_{*,n}\) and plaquette errors \(\varepsilon_n\) satisfying
\[
 \sup_n 2dA_{*,n}\varepsilon_n
 \le\bar g_0-\bar g                                      \tag{84.23}
\]
for some \(\bar g>0\), then every sector has the uniform full-kernel window
\[
 \bar g I\preceq |H_{m_{f,n}}(U)|
 \preceq \bar M I,                                         \tag{84.24}
\]
where \(\bar M\) follows from (84.13).  This supplies the global scalar
input required by V60, V67, and V73.
\end{corollary}

\begin{proof}
Equation (84.22) and (84.23) give the common lower bound.  Compactness of
the mass interval and the same uniform perturbation stock give the upper
bound through (84.13).
\end{proof}

\begin{assessment}[Why fixed plaquette admissibility may be insufficient]
\label{ass:v18v74-filling-growth}
The filling stock \(A_{*,n}\) can grow with lattice diameter.  A cutoff-
independent local plaquette bound \(\varepsilon\) then does not imply
(84.23); one needs shrinking curvature error, uniformly bounded filling
words, or a sharper gauge-covariant spectral theorem.  On a torus,
noncontractible holonomies are not plaquette filling words.  They should be
absorbed into a flat reference \(U^0\); proving that a general almost-flat
nonabelian field is uniformly close to such a reference is a separate
global topology problem.
\end{assessment}

\subsection{Exact rational flat-link pilot}

Take \(d=4\), \(m=1\), and the rational unitary phase
\[
 u=\frac{1023+64i}{1025},
 \qquad |u|^2=1.                                            \tag{84.25}
\]
Use the same abelian phase on all four constant links.  Every plaquette is
the identity.  Moreover
\[
 |u-1|^2=\frac4{1025}<\frac1{256},
 \qquad |u-1|<\frac1{16}.                                  \tag{84.26}
\]

\begin{theorem}[Executed V74 gap pilot]
\label{thm:v18v74-executed-pilot}
Theorem \ref{thm:v18v74-flat-reference-window}, used with the rational upper
bound \(\delta=1/16\), gives
\[
 \operatorname{gap}H_1(U)
 \ge 1-8\left(\frac1{16}\right)=\frac12.                  \tag{84.27}
\]
For the exact constant twisted mode, direct evaluation of (84.3) gives
\[
 H_1(U)^2
 =\frac{1050673}{1050625}I>I.                              \tag{84.28}
\]
All comparisons are exact rational comparisons.
\end{theorem}

\begin{proof}
The parametrization
\(\cos\theta=(1-s^2)/(1+s^2)\),
\(\sin\theta=2s/(1+s^2)\) with \(s=1/32\) gives (84.25).
Squaring the link difference gives (84.26).  Since
\(g_{\rm flat}(1)=1\), (84.27) follows.  Substitute
\(\cos\theta=1023/1025\) and \(\sin\theta=64/1025\) into (84.3) to obtain
(84.28).
\end{proof}

The executable
\[
 \texttt{scripts/verify\_sm\_v74\_wilson\_gap\_pilot.py}                \tag{84.29}
\]
emits
\[
 \texttt{artifacts/sm\_v74\_wilson\_gap\_pilot.json}.                  \tag{84.30}
\]
Their SHA--256 values are
\[
\begin{array}{c|l}
\hbox{object}&\hbox{SHA--256}\\ \hline
\hbox{verifier}&
\texttt{9D96663D436CF91B085B2F64A9D8CD54692F69FC83DE124C1A90AC24E7331F70}\\
\hbox{canonical result}&
\texttt{9408BEB41FD0BA32EE625E44509F469EF763FCA57C84820E914394DECD42D5D4}.
\end{array}                                                   \tag{84.31}
\]

\begin{firewall}[The physical admissibility row is not yet instantiated]
\label{fire:v18v74-physical-gap}
The theorem fixes a Wilson convention and proves a sufficient gap condition.
The attached programme does not supply sector masses \(m_f\), boundary and
topological sector, a flat-reference selection, or a cofinal bound for
\(A_{*,n}\varepsilon_n\).  The rational pilot is a flat abelian specimen,
not the physical Standard-Model link history.  Plaquette smallness alone is
not silently promoted to (84.23).
\end{firewall}

\subsection{Acquisition and next bottleneck}

V74 no longer treats the Wilson gap as an opaque scalar assumption.  It
reduces one rigorous branch to masses in \((0,2)\), a representation-uniform
distance from a commuting flat connection, and an explicit topology or
filling stock.  The next task is to decide which branch the physical
programme can actually populate: a tree-gauge filling certificate on
finite reflected boxes, an almost-flat-to-flat theorem on periodic
regulators, or a sharper gauge-covariant Wilson inequality that bypasses
linkwise closeness.  The next version is V75 and therefore requires the
companion-note checkpoint before V76 continues.

\begin{assessment}[V74 acquisition and boundary]
\label{ass:v18v74}
V74 proves the exact flat Wilson gap, a common five-sector perturbative
window, a finite plaquette-to-link conversion, and a cofinal sufficient
condition.  The executed rational phase verifies a nontrivial flat-link
instance.  The selected physical topology and cofinal admissibility stock
remain open.
\end{assessment}

\begin{record}[V74 append record]
\label{rec:v18v74}
V74 is appended to the validated V73 whose installed SHA--256 was
\texttt{945DB7C55EEF39924AEA047C93E0D0492568BC4A1A0DBA9552130410B4E38BCB}.
After installation only V74 is the active numeric SM note; the frozen
\texttt{V100\_f2} through \texttt{V100\_f17} archives are unchanged.
\end{record}


\section{V75: companion checkpoint and the curvature-aggregation decision}
\label{sec:v18v75}

V74 reduced one Wilson-gap branch to a flat reference and a link-distance
bound.  Its tree-gauge estimate charges each chord by the length of one
plaquette filling word and then takes the maximum \(A_*\).  Before extending
that route, the required five-version checkpoint was performed.  The audit
shows that literal filling-word enumeration is the wrong next
representation: curvature contributions should be aggregated by a finite
cellular operator before any norm is taken.

\subsection{Read-only companion audit}

The active Yang--Mills note at the time of inspection was
\[
 \texttt{Renewal\_Yang\_Mills\_Mass\_Gap\_Research\_Notes\_V13.tex},   \tag{85.1}
\]
with SHA--256
\[
 \texttt{8FE2899FEB248D8398E1A18B0A73D2C135C035C86B4249753FA12330DDBD9F27}.
                                                                  \tag{85.2}
\]
YM V13 constructs all 344 exact low-rank factors selected by its V12 rank
gate, retains 92 dense exceptions, and proves a cyclic collapse before
coefficient evaluation.  More importantly for the present step, it counts
over \(10^{13}\) Wilson and \(10^{18}\) endpoint complete paths and rejects
literal path streaming.  Its replacement is to aggregate partial paths by
frequency and only then pair the two halves and take norms.

The active Navier--Stokes note remained
\[
 \texttt{Renewal\_NS\_Stability\_Research\_Notes\_V26.tex},            \tag{85.3}
\]
with unchanged SHA--256
\[
 \texttt{82C887F8C2C69E4F28FAD7C3DC8DE5B9099CB5A4BF431EBA1FFF3E91935978AD}.
                                                                  \tag{85.4}
\]
Its rank-one kernel calculation still supplies the relevant methodological
constraint: restrict an operator to the structured image actually used,
and record honestly when that restriction gives no improvement.

\begin{theorem}[V75 audit conclusion]
\label{thm:v18v75-audit-conclusion}
No mathematical theorem from YM V13 or NS V26 is imported as an SM field
theorem.  Their finite-algebra lessons apply to the V74 certificate design:
\begin{enumerate}[leftmargin=2.2em]
\item sum all plaquette contributions through the cellular incidence map
      before applying an absolute value;
\item measure the right inverse only on the exact curvature image allowed by
      the Bianchi and boundary constraints; and
\item retain exceptional harmonic or topological modes explicitly rather
      than forcing them into a local filling representation.
\end{enumerate}
\end{theorem}

\begin{proof}
The cited YM result concerns Laurent-frequency response paths and the NS
result concerns Oseen-kernel inputs, so neither directly proves a Wilson
gap.  In V74, however, a filling word is likewise an intermediate path
representation.  The physical quantity is the aggregate link potential
produced by the curvature data.  Linearity in the abelianized problem allows
that aggregate to be formed first.  Curvatures lie in the image of the
cellular coboundary and obey compatibility conditions, so a norm on the full
plaquette coordinate space can overcharge them.  Noncontractible modes are
not in the local curvature image and must be represented separately, exactly
as dense exceptions are retained in YM V13.
\end{proof}

\subsection{The operator to be constructed}

Let \(C^0,C^1,C^2\) be real cellular cochains with fixed bases and
coboundaries
\[
 d_0:C^0\to C^1,
 \qquad d_1:C^1\to C^2,
 \qquad d_1d_0=0.                                           \tag{85.5}
\]
Choose a gauge slice \(G\subset C^1\) complementary to
\(\operatorname{im}d_0\), and let
\[
 \mathcal R=\operatorname{im}(d_1|_G)\subseteq C^2.         \tag{85.6}
\]
The desired aggregate filling operator is a right inverse
\[
 K:\mathcal R\longrightarrow G,
 \qquad d_1K=I_{\mathcal R}.                                \tag{85.7}
\]
For a norm pair \((\|\cdot\|_{2,*},\|\cdot\|_{1,*})\), its used-image
stock is
\[
 \kappa_{X,*}
 =\sup_{0\ne F\in\mathcal R}
 \frac{\|KF\|_{1,*}}{\|F\|_{2,*}}.                         \tag{85.8}
\]

\begin{proposition}[Why the aggregate stock can replace \(A_*\)]
\label{prop:v18v75-aggregate-replacement}
In an abelian phase chart, suppose the plaquette phase cochain is
\(F=d_1A\in\mathcal R\), the chosen gauge representative is \(A=KF\), and
\(\|F\|_{2,*}\le\varepsilon\).  Then
\[
 \|A\|_{1,*}\le\kappa_{X,*}\varepsilon.                   \tag{85.9}
\]
Consequently the V74 perturbative Wilson window holds with
\(\delta\) bounded by the representation charge times
\(\kappa_{X,*}\varepsilon\), rather than by the maximum length of a
separately chosen filling word.
\end{proposition}

\begin{proof}
Equation (85.9) is the definition of the restricted operator norm in
(85.8).  In an abelian phase chart, the distance of a link phase from its
flat reference is bounded by the absolute phase difference.  A charge-
\(q\) representation multiplies that phase by \(q\).  Insert the resulting
link-distance bound into Theorem \ref{thm:v18v74-flat-reference-window}.
\end{proof}

\begin{remark}[This is not yet the nonabelian theorem]
\label{rem:v18v75-nonabelian}
For a nonabelian connection, logarithms of plaquette holonomies do not add
linearly; conjugations and commutators produce a controlled nonlinear
remainder only inside a sufficiently small logarithm chart.  V76 will first
construct and execute the exact finite abelian operator.  The subsequent
nonabelian step must bound the Baker--Campbell--Hausdorff remainder rather
than identifying group multiplication with cochain addition.
\end{remark}

\subsection{V76 acquisition order}

\begin{decision}[Checkpoint-adjusted V76 direction]
\label{dec:v18v75-v76}
On an explicit reflected square complex:
\begin{enumerate}[leftmargin=2.2em]
\item build the integer matrices \(d_0,d_1\) and verify \(d_1d_0=0\);
\item fix a spanning-tree gauge slice and compute the exact rational right
      inverse \(K\) on \(\mathcal R\);
\item verify \(d_1K=I\) on the curvature image and compute the exact
      \(\ell^\infty\)-to-\(\ell^\infty\) used-image stock;
\item compare that stock with the crude filling-length debit; and
\item insert the result into the V74 Wilson-gap inequality for all integer
      Standard-Model hypercharges.
\end{enumerate}
This continues in V76 rather than ending at the audit.
\end{decision}

\begin{firewall}[Checkpoint scope]
\label{fire:v18v75-scope}
V75 records a verified cross-program decision and proves the abstract
abelian handoff (85.9).  It does not yet construct \(K\), control a
nonabelian logarithm, or establish a cofinal topology stock.  Those are
subsequent mathematical tasks.
\end{firewall}

\begin{assessment}[V75 acquisition and boundary]
\label{ass:v18v75}
The checkpoint replaces maximum filling-word length by a precise
used-image right-inverse target and isolates harmonic modes as explicit
exceptions.  The active YM and NS files were inspected without modification.
\end{assessment}

\begin{record}[V75 append record]
\label{rec:v18v75}
V75 is appended to the validated V74 whose installed SHA--256 was
\texttt{36C95858316E39E5FAD69BBF2EAFD6A01A2588185DAC575794D04457C8DDF336}.
After installation only V75 is the active numeric SM note; the frozen
\texttt{V100\_f2} through \texttt{V100\_f17} archives are unchanged.
\end{record}


\section{V76: exact cellular curvature aggregation and a sharper abelian gap stock}
\label{sec:v18v76}

V75 replaced literal filling-word enumeration by a cellular right-inverse
problem.  V76 constructs that inverse exactly on a reflected \(2\times2\)
square complex.  The coexact representative aggregates all four plaquette
curvatures before measuring the link field.  Its exact
\(\ell^\infty\)-operator norm is four times smaller than the spanning-tree
filling norm, and the resulting integer-hypercharge Wilson bound is inserted
into V74.

\subsection{Finite cellular setup}

Let \(X_2\) be the square disk with vertices
\[
 V=\{0,1,2\}^2,                                             \tag{86.1}
\]
all nearest-neighbor horizontal and vertical edges oriented in the positive
coordinate directions, and four counterclockwise unit plaquettes.  Thus
\[
 |V|=9,qquad |E|=12,qquad |P|=4.                          \tag{86.2}
\]
In the canonical cellular bases, let
\[
 d_0\in M_{12\times9}(\mathbb Z),
 \qquad d_1\in M_{4\times12}(\mathbb Z)                    \tag{86.3}
\]
be the coboundary matrices.  Each edge row of \(d_0\) has \(-1\) at its
tail and \(+1\) at its head.  Each plaquette row of \(d_1\) is its oriented
four-edge boundary.

\begin{lemma}[Exact cellular identity and full face rank]
\label{lem:v18v76-cellular-rank}
For \(X_2\),
\[
 d_1d_0=0,
 \qquad \operatorname{rank}d_1=4.                          \tag{86.4}
\]
Consequently the face Gram matrix
\[
 G=d_1d_1^*                                                \tag{86.5}
\]
is positive definite and invertible over \(\mathbb Q\).
\end{lemma}

\begin{proof}
Every vertex occurs twice with opposite signs in the boundary of a
plaquette, proving \(d_1d_0=0\).  Order the four vertical chord edges by
row and increasing positive \(x\)-coordinate.  The corresponding
\(4\times4\) submatrix of \(d_1\) is block lower triangular with unit
diagonal after the left boundary value is fixed, so it has determinant one.
Thus \(d_1\) has row rank four.  Hence
\(u^*Gu=\|d_1^*u\|_2^2>0\) for \(u\ne0\).
\end{proof}

The exact Gram data are
\[
 G=
 \begin{pmatrix}
 4&-1&-1&0\\
 -1&4&0&-1\\
 -1&0&4&-1\\
 0&-1&-1&4
 \end{pmatrix},                                             \tag{86.6}
\]
\[
 G^{-1}=\frac1{24}
 \begin{pmatrix}
 7&2&2&1\\
 2&7&1&2\\
 2&1&7&2\\
 1&2&2&7
 \end{pmatrix}.                                             \tag{86.7}
\]

\subsection{The coexact and tree right inverses}

Define
\[
 K_{\rm co}=d_1^*G^{-1}:C^2(X_2;\mathbb R)\to C^1(X_2;\mathbb R).
                                                                  \tag{86.8}
\]

\begin{theorem}[Exact coexact reconstruction]
\label{thm:v18v76-coexact}
The map (86.8) satisfies
\[
 d_1K_{\rm co}=I_{C^2},
 \qquad
 \operatorname{ran}K_{\rm co}=\operatorname{ran}d_1^*
 =\ker(d_1)^\perp.                                         \tag{86.9}
\]
For every curvature cochain \(F\), \(K_{\rm co}F\) is the unique
minimum-\(\ell^2\)-norm solution of \(d_1A=F\).
\end{theorem}

\begin{proof}
Equation (86.5) gives
\(d_1K_{\rm co}=d_1d_1^*G^{-1}=I\).  The range statement is immediate from
(86.8), and finite-dimensional orthogonality gives
\(\operatorname{ran}d_1^*=\ker(d_1)^\perp\).  Every solution has the
orthogonal decomposition
\[
 A=K_{\rm co}F+Z,
 \qquad Z\in\ker d_1,                                      \tag{86.10}
\]
so \(\|A\|_2^2=\|K_{\rm co}F\|_2^2+\|Z\|_2^2\), proving uniqueness and
minimality.
\end{proof}

For comparison, fix the spanning tree consisting of all horizontal edges
and the two vertical edges on the left boundary.  Set those tree-link
potentials to zero.  On each plaquette row, reconstruct the two remaining
vertical chord potentials by cumulative sums from left to right.  Denote the
resulting right inverse by \(K_{\rm tr}\).

\begin{proposition}[Exact infinity stocks]
\label{prop:v18v76-infinity-stocks}
Both maps are right inverses, and
\[
\boxed{
 \|K_{\rm co}\|_{\infty\to\infty}=\frac12,
 \qquad
 \|K_{\rm tr}\|_{\infty\to\infty}=2.}                    \tag{86.11}
\]
Hence
\[
 \frac{\|K_{\rm co}\|_{\infty\to\infty}}
      {\|K_{\rm tr}\|_{\infty\to\infty}}
 =\frac14.                                                   \tag{86.12}
\]
\end{proposition}

\begin{proof}
The coexact identity is Theorem \ref{thm:v18v76-coexact}.  The cumulative
tree formula gives, for the two faces in each row,
\[
 A_1=F_0,qquad A_2=F_0+F_1,                                \tag{86.13}
\]
so applying \(d_1\) returns \((F_0,F_1)\) and the largest absolute row sum
is two.  Direct multiplication of (86.7) by \(d_1^*\) gives a
\(12\times4\) rational matrix.  Every one of its twelve absolute row sums is
\[
 \frac{12}{24}=\frac12.                                    \tag{86.14}
\]
The induced infinity norm is the maximum absolute row sum, proving (86.11).
The complete rational matrices are retained in the executable payload.
\end{proof}

\begin{assessment}[The gain is gauge-slice dependent but valid]
\label{ass:v18v76-gauge-slice}
The coexact slice minimizes \(\ell^2\), not \(\ell^\infty\), yet on this
complex it also improves the exact infinity stock.  The comparison does not
claim that coexact gauge is infinity-optimal on every complex.  It proves
that the V74 tree-filling debit is not intrinsic even on the first
nontrivial square disk.
\end{assessment}

\subsection{Reflection equivariance}

Reflect the disk by \(r(x,y)=(2-x,y)\).  Let \(R_1\) be its signed action on
oriented edges and \(R_2\) its signed action on oriented plaquettes.  A
horizontal edge and a plaquette reverse orientation; a vertical edge
retains orientation.  Thus \(R_1,R_2\) are orthogonal involutions.

\begin{theorem}[Reflection-compatible aggregate gauge]
\label{thm:v18v76-reflection}
One has
\[
 d_1R_1=R_2d_1,
 \qquad
 R_1K_{\rm co}=K_{\rm co}R_2.                              \tag{86.15}
\]
Therefore reflected curvature data reconstruct to reflected link
potentials without an additional gauge correction.
\end{theorem}

\begin{proof}
The first identity is naturality of cellular coboundary with the orientation
signs; the executable verifies it entrywise over \(\mathbb Z\).  Taking
adjoints gives \(R_1d_1^*=d_1^*R_2\).  The first identity also implies that
\(G=d_1d_1^*\) commutes with \(R_2\), hence so does \(G^{-1}\).  Therefore
\[
 R_1K_{\rm co}
 =R_1d_1^*G^{-1}
 =d_1^*R_2G^{-1}
 =d_1^*G^{-1}R_2
 =K_{\rm co}R_2.                                           \tag{86.16}
\]
\end{proof}

\subsection{Integer-hypercharge Wilson insertion}

Use the integer convention \(y=6Y\), so a primitive abelian link angle
\(A_e\) acts in sector \(f\) by \(e^{iy_fA_e}\).  For real \(s\),
\[
 |e^{is}-1|\le|s|.                                         \tag{86.17}
\]
Since \(\max_f|y_f|=6\), a curvature bound
\(\|F\|_\infty\le\varepsilon\) gives
\[
 \delta_{U(1)}
 \le6\|K_{\rm co}\|_{\infty\to\infty}\varepsilon
 =3\varepsilon.                                            \tag{86.18}
\]

\begin{theorem}[Executed U(1) Wilson gap comparison]
\label{thm:v18v76-gap-comparison}
Take \(\varepsilon=1/1000\), four-dimensional Wilson hopping, and
\(m=1\).  The coexact certificate gives
\[
 \delta_{U(1)}^{\rm co}\le\frac3{1000},
 \qquad
 \operatorname{gap}H_1
 \ge1-8\frac3{1000}=\frac{122}{125}.                       \tag{86.19}
\]
The tree certificate gives
\[
 \delta_{U(1)}^{\rm tr}\le\frac3{250},
 \qquad
 \operatorname{gap}H_1
 \ge1-8\frac3{250}=\frac{113}{125}.                        \tag{86.20}
\]
All inequalities are simultaneous for the five integer hypercharges in
(80.3).
\end{theorem}

\begin{proof}
Insert the two exact norms from (86.11) and \(\max|y_f|=6\) into (86.17).
Then apply the \(d=4,m=1\) case of Theorem
\ref{thm:v18v74-flat-reference-window}.
\end{proof}

\begin{corollary}[Cofinal abelian stock]
\label{cor:v18v76-cofinal-abelian}
For reflected disk complexes \(X_n\) with full-row-rank \(d_{1,n}\), put
\[
 K_n=d_{1,n}^*(d_{1,n}d_{1,n}^*)^{-1},
 \qquad
 \kappa_n=\|K_n\|_{\infty\to\infty}.                      \tag{86.21}
\]
If the Wilson masses have common free gap \(g_0>0\) and
\[
 \sup_n 2d\,y_{\max}\kappa_n\varepsilon_n
 \le g_0-g                                                     \tag{86.22}
\]
for some \(g>0\), then the U(1) contribution to link perturbation is
compatible with a uniform Wilson lower gap \(g\), before adding the
nonabelian link deviations.
\end{corollary}

\begin{proof}
The right-inverse identity gives
\(\|A_n\|_\infty\le\kappa_n\varepsilon_n\).  Equation (86.17) adds the
integer charge, and V74 supplies the Wilson perturbation factor \(2d\).
\end{proof}

\subsection{Executable certificate}

The verifier
\[
 \texttt{scripts/verify\_sm\_v76\_cellular\_right\_inverse.py}          \tag{86.23}
\]
constructs \(d_0,d_1,G^{-1},K_{\rm co},K_{\rm tr},R_1,R_2\) over exact
rationals.  It verifies (86.4), both right-inverse identities, both
reflection identities, the exact infinity stocks, and the two gap readings.
Its canonical payload is
\[
 \texttt{artifacts/sm\_v76\_cellular\_right\_inverse.json}.             \tag{86.24}
\]
The SHA--256 ledger is
\[
\begin{array}{c|l}
\hbox{object}&\hbox{SHA--256}\\ \hline
\hbox{verifier}&
\texttt{C2C39D9F726127873AD499F2E28C33DCFF2A1FF71EEB8AB98B218669D939D7E6}\\
\hbox{canonical payload}&
\texttt{5FD1A262C14E3B3E8840F1E26E164754FB198E6E15CA524BE6E5B559C8D34F59}.
\end{array}                                                   \tag{86.25}
\]

\begin{firewall}[Abelian disk result only]
\label{fire:v18v76-boundary}
The exact certificate concerns the linear U(1) phase on a simply connected
two-dimensional square disk, used as a finite curvature sheet inside the
four-dimensional Wilson estimate.  It does not control SU(2) or SU(3)
logarithms, nonabelian commutators, four-dimensional two-cycle
compatibilities, periodic harmonic modes, or growth of \(\kappa_n\).  The
synthetic curvature bound \(1/1000\) is not extracted from the attached
physical measure.
\end{firewall}

\subsection{Acquisition and next bottleneck}

V76 executes the checkpoint decision and replaces the first tree-filling
stock by an exact aggregate operator norm.  The next upstream problem has
two branches.  The finite analytic branch is to put SU(2) and SU(3)
plaquettes in a principal logarithm chart and bound the nonlinear
Baker--Campbell--Hausdorff remainder around the coexact linear solution.  The
cofinal branch is to estimate \(\kappa_n\) on reflected four-dimensional
regulators and separate harmonic modes.  Either branch must precede a claim
that V74 supplies the physical global gap.

\begin{assessment}[V76 acquisition and boundary]
\label{ass:v18v76}
V76 proves and executes an exact reflection-equivariant cellular right
inverse with a fourfold infinity-stock gain over tree filling, then inserts
it into all five integer-hypercharge Wilson bounds.  Nonabelian logarithms,
cofinal four-dimensional stocks, and physical curvature probabilities
remain open.
\end{assessment}

\begin{record}[V76 append record]
\label{rec:v18v76}
V76 is appended to the validated V75 whose installed SHA--256 was
\texttt{9B4899CE441E07043DEB24AF034C7366DC850B06B4A743104C2445CA0403A9CB}.
After installation only V76 is the active numeric SM note; the frozen
\texttt{V100\_f2} through \texttt{V100\_f17} archives are unchanged.
\end{record}


\section{V77: a nonabelian logarithm-chart lift of the coexact curvature bound}
\label{sec:v18v77}

V76 solved the abelian curvature equation by an exact cellular right
inverse.  For SU(2) or SU(3), the logarithm of a plaquette product differs
from the signed sum of its four link logarithms.  V77 bounds that difference
without discarding it.  A rational majorant makes the coexact reconstruction
a self-map of a small Lie-algebra ball, yielding both an existence theorem
and an a posteriori bound for any physical connection already placed in the
coexact logarithmic gauge.

\subsection{A four-factor logarithm remainder}

Let \(X_1,\ldots,X_4\) be skew-Hermitian matrices in a fixed unitary
representation, with
\[
 \|X_j\|_{\rm op}\le r,
 \qquad S=4r<\frac12.                                      \tag{87.1}
\]
Put
\[
 P=e^{X_1}e^{X_2}e^{X_3}e^{X_4},
 \qquad Y=P-I.                                              \tag{87.2}
\]

\begin{lemma}[Rational product majorants]
\label{lem:v18v77-product-majorants}
One has
\[
 \|Y\|_{\rm op}
 \le e^S-1
 \le q(S):=\frac{S}{1-S}<1,                                \tag{87.3}
\]
and
\[
 \left\|P-I-\sum_{j=1}^4X_j\right\|_{\rm op}
 \le e^S-1-S
 \le\frac{S^2}{1-S}.                                      \tag{87.4}
\]
\end{lemma}

\begin{proof}
Expand the four exponentials into their absolutely convergent power series.
The sum of operator-norm majorants of all homogeneous terms is \(e^S\).
Removing the constant term gives the first exponential bound; removing the
linear terms gives the second.  For \(0\le S<1\), the inequalities
\[
 e^S\le\sum_{n\ge0}S^n=\frac1{1-S},
 \qquad
 e^S-1-S\le\sum_{n\ge2}S^n=\frac{S^2}{1-S}                 \tag{87.5}
\]
give the rational bounds.  The condition \(S<1/2\) makes \(q(S)<1\).
\end{proof}

Since \(\|Y\|<1\), define
\[
 \log P=\sum_{n=1}^{\infty}\frac{(-1)^{n+1}}nY^n.          \tag{87.6}
\]
For a unitary \(P\) in this ball, this is its principal skew-Hermitian
logarithm.

\begin{theorem}[Explicit nonabelian face remainder]
\label{thm:v18v77-face-remainder}
The nonlinear remainder
\[
 \mathcal N(X_1,\ldots,X_4)
 =\log P-\sum_{j=1}^4X_j                                   \tag{87.7}
\]
satisfies
\[
\boxed{
 \|\mathcal N\|_{\rm op}
 \le B(S):=
 \frac{S^2}{1-S}
 +\frac{S^2}{2(1-S)(1-2S)}.}                               \tag{87.8}
\]
\end{theorem}

\begin{proof}
Split
\[
 \log P-\sum_jX_j
 =\bigl(P-I-\sum_jX_j\bigr)+\bigl(\log(I+Y)-Y\bigr).       \tag{87.9}
\]
Lemma \ref{lem:v18v77-product-majorants} bounds the first term.  For the
second,
\[
 \|\log(I+Y)-Y\|
 \le\sum_{n\ge2}\frac{\|Y\|^n}{n}
 \le\frac{q(S)^2}{2(1-q(S))}
 =\frac{S^2}{2(1-S)(1-2S)}.                                \tag{87.10}
\]
Adding the bounds proves (87.8).
\end{proof}

\begin{remark}[The bound is deliberately elementary]
\label{rem:v18v77-elementary}
Formula (87.8) does not use a truncated formal Baker--Campbell--Hausdorff
series.  It controls the complete convergent logarithm by operator norms and
therefore remains valid for noncommuting matrices in any finite-dimensional
unitary representation.  Sharper commutator-sensitive constants remain
possible.
\end{remark}

\subsection{Nonlinear coexact reconstruction}

Work on the V76 square disk and apply \(K_{\rm co}\) componentwise to
Lie-algebra-valued face cochains.  Its matrix coefficients are real, so it
preserves skew-Hermitian values.  Let
\[
 \kappa=\|K_{\rm co}\|_{\infty\to\infty}=\frac12.          \tag{87.11}
\]
For a link-log cochain \(A\), define the face cochain
\[
 \mathcal F(A)_p
 =\log\left(\prod_{e\in\partial p}^{\to}e^{\epsilon_{p,e}A_e}\right)
 = (d_1A)_p+\mathcal N_p(A).                                \tag{87.12}
\]
The product follows the oriented plaquette boundary; inversely oriented
links use \(-A_e\).  If \(\|A\|_\infty\le r\), Theorem
\ref{thm:v18v77-face-remainder} gives
\[
 \|\mathcal N(A)\|_\infty\le B(4r).                       \tag{87.13}
\]

For prescribed skew-Hermitian face logs \(F\), define on the coexact
link-log space
\[
 \mathcal T_F(A)=K_{\rm co}\bigl(F-\mathcal N(A)\bigr).    \tag{87.14}
\]

\begin{theorem}[Finite nonabelian coexact existence]
\label{thm:v18v77-nonabelian-existence}
Suppose
\[
 \|F\|_\infty\le\varepsilon,
 \qquad 4r<\frac12,
 \qquad
 \kappa\bigl(\varepsilon+B(4r)\bigr)\le r.                \tag{87.15}
\]
Then there exists a coexact skew-Hermitian link-log cochain \(A\) with
\[
 \|A\|_\infty
 \le\kappa\bigl(\varepsilon+B(4r)\bigr)                   \tag{87.16}
\]
such that every principal plaquette logarithm equals the prescribed value:
\[
 \mathcal F(A)=F.                                          \tag{87.17}
\]
\end{theorem}

\begin{proof}
The closed radius-\(r\) ball in the finite-dimensional real vector space of
coexact skew-Hermitian cochains is compact and convex.  The principal
logarithm is continuous on the ball by (87.3), so \(\mathcal T_F\) is
continuous.  Equations (87.11), (87.13), and (87.15) show that it maps the
ball into itself, in fact into the smaller ball (87.16).  Brouwer's theorem
gives a fixed point.  Applying \(d_1\) and using
\(d_1K_{\rm co}=I\) gives
\[
 d_1A=F-\mathcal N(A),                                     \tag{87.18}
\]
which is exactly (87.17) by (87.12).
\end{proof}

\begin{theorem}[A posteriori control in coexact logarithmic gauge]
\label{thm:v18v77-aposteriori}
Let an existing link field admit logarithms \(A_e\) satisfying
\[
 A\in\operatorname{ran}d_1^*,
 \quad \|A\|_\infty\le r,
 \quad \|\mathcal F(A)\|_\infty\le\varepsilon,            \tag{87.19}
\]
with (87.15).  Then it obeys the improved bound (87.16).
\end{theorem}

\begin{proof}
Coexactness and (86.9) give \(A=K_{\rm co}d_1A\).  Equation (87.12) gives
\(d_1A=\mathcal F(A)-\mathcal N(A)\).  Apply \(K_{\rm co}\), (87.11), and
(87.13).
\end{proof}

\subsection{Exact rational chart}

Choose
\[
 \kappa=\frac12,qquad
 \varepsilon=\frac1{1000},
 \qquad r=\frac1{100},qquad S=4r=\frac1{25}.               \tag{87.20}
\]
Then the rational majorants are
\[
 q(S)=\frac1{24},qquad
 \frac{S^2}{1-S}=\frac1{600},qquad
 \frac{S^2}{2(1-S)(1-2S)}=\frac1{1104},                    \tag{87.21}
\]
and hence
\[
 B(S)=\frac{71}{27600}.                                    \tag{87.22}
\]

\begin{theorem}[Executed V77 self-map and gap]
\label{thm:v18v77-executed}
The fixed-point image radius is
\[
 r_{\rm im}
 =\frac12\left(\frac1{1000}+\frac{71}{27600}\right)
 =\frac{493}{276000}<\frac1{100}.                           \tag{87.23}
\]
Thus Theorem \ref{thm:v18v77-nonabelian-existence} applies.  In the same
unitary representation,
\[
 \|e^{A_e}-I\|_{\rm op}\le\|A_e\|_{\rm op}le r_{\rm im}. \tag{87.24}
\]
For \(d=4,m=1\), V74 gives the Wilson lower gap
\[
\boxed{
 \operatorname{gap}H_1
 \ge1-8r_{\rm im}
 =\frac{34007}{34500}.}                                    \tag{87.25}
\]
\end{theorem}

\begin{proof}
The fraction computations in (87.21)--(87.23) are direct.  For
skew-Hermitian \(A\),
\[
 e^A-I=\int_0^1e^{tA}A\,dt,                                \tag{87.26}
\]
and \(e^{tA}\) is unitary, proving (87.24).  Insert it into Theorem
\ref{thm:v18v74-flat-reference-window} to obtain (87.25).
\end{proof}

The verifier
\[
 \texttt{scripts/verify\_sm\_v77\_nonabelian\_log\_chart.py}           \tag{87.27}
\]
checks all rational inequalities and emits
\[
 \texttt{artifacts/sm\_v77\_nonabelian\_log\_chart.json}.              \tag{87.28}
\]
Their SHA--256 values are
\[
\begin{array}{c|l}
\hbox{object}&\hbox{SHA--256}\\ \hline
\hbox{verifier}&
\texttt{48CE6A1D819C72F3A5E66906278ECA3A58B2006BB1C1414F2E75F75AE16D0A02}\\
\hbox{canonical payload}&
\texttt{A8438758366B5201144D11B431B0A49DA7725990DFA8D27B46D04743D9218609}.
\end{array}                                                   \tag{87.29}
\]

\subsection{Reflection and representation scope}

The V76 identity \(R_1K_{\rm co}=K_{\rm co}R_2\) persists for
Lie-algebra-valued cochains.  The plaquette product and principal logarithm
are natural under unitary conjugation.  Therefore, for a reflection action
whose orientation sign and internal conjugation preserve the logarithm
chart, \(\mathcal T_F\) preserves the corresponding reflection-fixed real
subspace.  Brouwer may be applied on that subspace to produce a reflected
solution whenever \(F\) is reflected.

The bound is representationwise.  To treat the complete SM link, one must
apply it to the SU(3) and SU(2) represented logarithms and combine those
deviations with the U(1) estimate from V76.  The norm of the differential of
each representation must be included; it is not automatically one for an
arbitrary normalization.

\begin{firewall}[Existence is not gauge control of every physical field]
\label{fire:v18v77-gauge-control}
Theorem \ref{thm:v18v77-nonabelian-existence} constructs at least one
coexact link-log field for small prescribed face logs on the square disk.
It does not prove that every physical connection with those plaquette logs
is gauge equivalent to that solution.  The a posteriori theorem applies
only after a physical field has been placed in the stated coexact logarithm
gauge with the trial radius.  Uniqueness modulo gauge, four-dimensional
Bianchi constraints, periodic harmonic modes, representation constants,
and cofinal uniformity remain open.
\end{firewall}

\subsection{Acquisition and next bottleneck}

V77 supplies the first complete nonabelian remainder instead of silently
linearizing plaquette holonomy.  The next local theorem must control the
derivative of \(\mathcal N\) sharply enough to make \(\mathcal T_F\) a
contraction; uniqueness would then identify the coexact representative and
support a gauge-fixing theorem.  In parallel, the square-disk calculation
must be lifted to four-dimensional reflected complexes with harmonic modes
split off.  Those steps precede a physical cofinal Wilson gap.

\begin{assessment}[V77 acquisition and boundary]
\label{ass:v18v77}
V77 proves a full noncommutative logarithm remainder, a finite coexact
existence theorem, and an a posteriori link and Wilson-gap bound with exact
rational constants.  It does not yet establish uniqueness or place every
physical admissible field in the coexact chart.
\end{assessment}

\begin{record}[V77 append record]
\label{rec:v18v77}
V77 is appended to the validated V76 whose installed SHA--256 was
\texttt{24C4B9B9F05E30D2FC49904ECA787A2081B0C4CEFB69F898CD18622743FCECC5}.
After installation only V77 is the active numeric SM note; the frozen
\texttt{V100\_f2} through \texttt{V100\_f17} archives are unchanged.
\end{record}


\section{V78: contraction, uniqueness, and stability in the nonabelian coexact chart}
\label{sec:v18v78}

V77 used a complete logarithm remainder to prove existence by Brouwer, but
that argument did not identify a unique coexact representative.  The
remainder has no constant or linear term, so its derivative is small on a
small link-log ball.  V78 proves an explicit Lipschitz bound, verifies that
the V76 right-inverse stock turns the reconstruction into a strict
contraction, and obtains uniqueness, reflection covariance, and stable
dependence on plaquette logarithms.

\subsection{A Lipschitz bound with the linear term removed}

For a four-tuple \(X=(X_1,\ldots,X_4)\) of skew-Hermitian matrices, write
\[
 P(X)=e^{X_1}e^{X_2}e^{X_3}e^{X_4},                        \tag{88.1}
\]
\[
 R_{\rm prod}(X)=P(X)-I-\sum_{j=1}^4X_j,                   \tag{88.2}
\]
and
\[
 R_{\log}(Y)=\log(I+Y)-Y.                                  \tag{88.3}
\]
Then the V77 face remainder is
\[
 \mathcal N(X)=R_{\rm prod}(X)+R_{\log}(P(X)-I).            \tag{88.4}
\]

Let \(X,Z\) satisfy
\[
 \sum_j\|X_j\|_{\rm op}\le S,qquad
 \sum_j\|Z_j\|_{\rm op}\le S,qquad S<\frac12,           \tag{88.5}
\]
and put
\[
 \Delta_1(X,Z)=\sum_{j=1}^4\|X_j-Z_j\|_{\rm op}.           \tag{88.6}
\]

\begin{lemma}[Product-tail Lipschitz majorant]
\label{lem:v18v78-product-tail}
One has
\[
 \|R_{\rm prod}(X)-R_{\rm prod}(Z)\|_{\rm op}
 \le(e^S-1)\Delta_1(X,Z)
 \le q(S)\Delta_1(X,Z),                                   \tag{88.7}
\]
where \(q(S)=S/(1-S)\).
\end{lemma}

\begin{proof}
Join \(Z\) to \(X\) by the affine path of four-tuples.  Convexity of the
operator norm keeps the sum of component norms at most \(S\).  Expand the
product of exponentials as a noncommutative entire power series.  After the
constant and degree-one terms are removed, the norm majorant is
\(e^s-1-s\).  Its directional derivative is bounded by
\((e^S-1)\Delta_1(X,Z)\).  Integrating along the affine path gives the first
inequality.  The rational estimate (87.5) gives the second.
\end{proof}

\begin{lemma}[Logarithm-tail Lipschitz majorant]
\label{lem:v18v78-log-tail}
If \(\|Y\|,\|W\|\le q<1\), then
\[
 \|R_{\log}(Y)-R_{\log}(W)\|_{\rm op}
 \le\frac{q}{1-q}\|Y-W\|_{\rm op}.                         \tag{88.8}
\]
\end{lemma}

\begin{proof}
From the convergent series (87.6),
\[
 R_{\log}(Y)-R_{\log}(W)
 =\sum_{n\ge2}\frac{(-1)^{n+1}}n(Y^n-W^n).                 \tag{88.9}
\]
The noncommutative telescoping identity gives
\[
 \|Y^n-W^n\|
 \le n q^{n-1}\|Y-W\|.                                   \tag{88.10}
\]
Summing \(\sum_{n\ge2}q^{n-1}=q/(1-q)\) proves (88.8).
\end{proof}

For skew-Hermitian \(A,B\), the exponential is operator-norm Lipschitz:
\[
 \|e^A-e^B\|_{\rm op}\le\|A-B\|_{\rm op}.                 \tag{88.11}
\]
Indeed the difference is the integral of a product of two unitary
exponentials and \(A-B\).  Telescoping the four unitary factors therefore
gives
\[
 \|P(X)-P(Z)\|_{\rm op}\le\Delta_1(X,Z).                   \tag{88.12}
\]

\begin{theorem}[Nonlinear face Lipschitz bound]
\label{thm:v18v78-face-lipschitz}
Under (88.5),
\[
\boxed{
 \|\mathcal N(X)-\mathcal N(Z)\|_{\rm op}
 \le \frac{q(S)(2-q(S))}{1-q(S)}\Delta_1(X,Z).}             \tag{88.13}
\]
If every link logarithm has norm at most \(r\), then with
\(S=4r\) and
\(\delta=\max_j\|X_j-Z_j\|\),
\[
 \|\mathcal N(X)-\mathcal N(Z)\|_{\rm op}
 \le L(r)\delta,                                           \tag{88.14}
\]
where
\[
 L(r)=4\left(q(4r)+\frac{q(4r)}{1-q(4r)}\right).            \tag{88.15}
\]
\end{theorem}

\begin{proof}
Apply Lemma \ref{lem:v18v78-product-tail} to the first term in (88.4).
Lemma \ref{lem:v18v78-log-tail}, (87.3), and (88.12) bound the second by
\(q(S)/(1-q(S))\) times \(\Delta_1\).  Adding the two coefficients gives
(88.13).  Since \(\Delta_1\le4\delta\), (88.14)--(88.15) follow.
\end{proof}

\subsection{Banach reconstruction on the square disk}

Retain the V77 map
\[
 \mathcal T_F(A)=K_{\rm co}(F-\mathcal N(A))               \tag{88.16}
\]
on the coexact link-log ball of radius \(r\), and let
\(\kappa=\|K_{\rm co}\|_{\infty\to\infty}\).

\begin{theorem}[Unique coexact nonabelian reconstruction]
\label{thm:v18v78-unique-reconstruction}
Assume the V77 self-map inequalities and
\[
 c:=\kappa L(r)<1.                                         \tag{88.17}
\]
Then for every admitted curvature-log cochain \(F\):
\begin{enumerate}[leftmargin=2.2em]
\item \(\mathcal T_F\) has a unique fixed point \(A_F\) in the radius-
      \(r\) coexact ball;
\item the Picard iterates \(A^{(n+1)}=\mathcal T_F(A^{(n)})\) satisfy
\[
 \|A^{(n)}-A_F\|_\infty
 \le\frac{c^n}{1-c}
 \|A^{(1)}-A^{(0)}\|_\infty;                              \tag{88.18}
\]
\item any existing coexact solution in that ball equals \(A_F\); and
\item for two admitted face logs,
\[
\boxed{
 \|A_F-A_G\|_\infty
 \le\frac{\kappa}{1-c}\|F-G\|_\infty.}                   \tag{88.19}
\]
\end{enumerate}
\end{theorem}

\begin{proof}
Theorem \ref{thm:v18v78-face-lipschitz} applied facewise and the definition
of \(\kappa\) give
\[
 \|\mathcal T_F(A)-\mathcal T_F(B)\|_\infty
 \le c\|A-B\|_\infty.                                    \tag{88.20}
\]
Banach's fixed-point theorem proves uniqueness and (88.18).  Any coexact
solution satisfies the fixed-point equation by the proof of Theorem
\ref{thm:v18v77-aposteriori}, proving item three.  Finally,
\[
 \|A_F-A_G\|_\infty
 \le\kappa\|F-G\|_\infty+c\|A_F-A_G\|_\infty,             \tag{88.21}
\]
which rearranges to (88.19).
\end{proof}

\begin{corollary}[Reflection is forced by uniqueness]
\label{cor:v18v78-reflection-unique}
Suppose the face data are invariant under the V76 reflection action,
including the required internal unitary conjugation.  If the nonlinear
plaquette map is equivariant and preserves the logarithm chart, then the
unique fixed point \(A_F\) is reflection invariant.
\end{corollary}

\begin{proof}
Reflecting \(A_F\) gives another coexact radius-\(r\) solution for the same
face data because \(R_1K_{\rm co}=K_{\rm co}R_2\).  Uniqueness in Theorem
\ref{thm:v18v78-unique-reconstruction} makes the two solutions equal.
\end{proof}

\subsection{Exact contraction constants}

Use the V77 radius \(r=1/100\), so
\[
 S=\frac1{25},qquad q(S)=\frac1{24}.                       \tag{88.22}
\]
Then
\[
 \frac{q}{1-q}=\frac1{23},
 \qquad
 L(r)=4\left(\frac1{24}+\frac1{23}\right)
 =\frac{47}{138}.                                          \tag{88.23}
\]

\begin{theorem}[Executed V78 contraction certificate]
\label{thm:v18v78-executed}
With the V76 stock \(\kappa=1/2\),
\[
\boxed{
 c=\kappa L(r)=\frac{47}{276}<1.}                          \tag{88.24}
\]
The curvature-to-link stability coefficient is
\[
\boxed{
 \frac{\kappa}{1-c}=\frac{138}{229}.}                      \tag{88.25}
\]
Thus the V77 fixed point is unique in the certified ball and Picard
iteration converges at least geometrically with ratio \(47/276\).
\end{theorem}

\begin{proof}
Substitute (88.23) and \(\kappa=1/2\) into (88.17) and (88.19).  Exact
fraction reduction gives (88.24)--(88.25).
\end{proof}

The executable
\[
 \texttt{scripts/verify\_sm\_v78\_nonabelian\_contraction.py}           \tag{88.26}
\]
checks these exact rational identities and emits
\[
 \texttt{artifacts/sm\_v78\_nonabelian\_contraction.json}.              \tag{88.27}
\]
Their hashes are
\[
\begin{array}{c|l}
\hbox{object}&\hbox{SHA--256}\\ \hline
\hbox{verifier}&
\texttt{B318B3B8874EF36D0B4C8EFA3A859010A76E4882A7DF48598511B11E16F3D886}\\
\hbox{canonical payload}&
\texttt{2FDFCDC6C27DEDCF2EB102A59D7830DF853D9F697AE5933F6488F0555F66CB26}.
\end{array}                                                   \tag{88.28}
\]

\subsection{Cofinal stability capsule}

Let \(K_n\) be reflection-compatible coexact right inverses on finite
regulators, with stocks \(\kappa_n\).  Let \(r_n\) be logarithm-chart radii
and define
\[
 c_n=\kappa_n L(r_n).                                       \tag{88.29}
\]

\begin{corollary}[Uniform nonlinear curvature inversion]
\label{cor:v18v78-cofinal-inversion}
If the V77 self-map inequalities hold and
\[
 \sup_n c_n\le c_*<1,
 \qquad
 \sup_n\frac{\kappa_n}{1-c_n}<\infty,                      \tag{88.30}
\]
then the coexact logarithmic reconstruction is unique and uniformly
Lipschitz at every cutoff.  If in addition the resulting link radii make the
V74 five-sector perturbation debit strictly smaller than the common free
gap, those reconstructed connections have a uniform Wilson gap.
\end{corollary}

\begin{proof}
Apply Theorem \ref{thm:v18v78-unique-reconstruction} at each cutoff.  The
second assertion is Corollary \ref{cor:v18v74-five-sector-gap} after
combining the represented factor deviations.
\end{proof}

\begin{firewall}[Uniqueness in a slice is not a global gauge theorem]
\label{fire:v18v78-global-gauge}
V78 proves uniqueness among coexact logarithmic solutions inside one small
ball.  It does not prove that every admissible lattice connection admits
such a gauge, that two arbitrary representatives with equal plaquette logs
are gauge equivalent, or that periodic harmonic sectors enter a common
ball.  The Lipschitz constants are for one represented nonabelian factor;
SU(3), SU(2), and U(1) deviations still must be combined with their actual
representation norms.
\end{firewall}

\subsection{Acquisition and next bottleneck}

The finite square-disk logarithm chart now has existence, uniqueness,
reflection covariance, and stable dependence.  The next upstream step is a
four-dimensional Hodge stock: construct \(K_n\) for a reflected hypercubic
box, identify the Bianchi-compatible face image, and bound
\(\kappa_n\) without expanding all filling surfaces.  That calculation will
decide whether the coexact branch can remain uniform cofinally or whether a
different gauge-covariant Wilson estimate is required.

\begin{assessment}[V78 acquisition and boundary]
\label{ass:v18v78}
V78 closes the local nonuniqueness left by V77 with contraction factor
\(47/276\) and stability coefficient \(138/229\).  A global/cofinal gauge
chart and four-dimensional right-inverse stock remain unresolved.
\end{assessment}

\begin{record}[V78 append record]
\label{rec:v18v78}
V78 is appended to the validated V77 whose installed SHA--256 was
\texttt{6FD20D0B9427FE033E5BACA653AEBD5DBD6356E9501C02CFB68C040559696C99}.
After installation only V78 is the active numeric SM note; the frozen
\texttt{V100\_f2} through \texttt{V100\_f17} archives are unchanged.
\end{record}


\section{V79: exact Hodge inversion on a reflected four-dimensional cube}
\label{sec:v18v79}

V78 made the coexact nonlinear chart quantitative in terms of one finite
stock \(\kappa=\|K\|_{\infty\to\infty}\).  A four-dimensional complex
introduces a new issue absent from the square disk: its plaquette variables
obey Bianchi compatibility conditions, so \(d_1d_1^*\) is singular on the
full face space.  V79 constructs the exact rational Moore--Penrose inverse
on the compatible curvature image of one reflected four-cube, verifies all
four Penrose identities and reflection equivariance, and closes the V77--V78
nonlinear chart with the computed stock.

\subsection{Cell counts and curvature rank}

Let \(Q^4=[0,1]^4\) with its standard cubical cell structure and positive
coordinate orientations.  Its zero-, one-, and two-cell counts are
\[
 |C_0|=16,qquad |C_1|=32,qquad |C_2|=24.                  \tag{89.1}
\]
Let
\[
 d_0:C^0(Q^4)\to C^1(Q^4),
 \qquad
 B:=d_1:C^1(Q^4)\to C^2(Q^4)                              \tag{89.2}
\]
be the integer coboundary matrices in the canonical cell bases.

\begin{proposition}[Four-cube curvature rank]
\label{prop:v18v79-curvature-rank}
The exact matrices satisfy
\[
 Bd_0=0,qquad \operatorname{rank}B=17.                    \tag{89.3}
\]
Thus the compatible curvature space
\[
 \mathcal R=\operatorname{im}B\subset C^2(Q^4)             \tag{89.4}
\]
has dimension \(17\) and codimension \(7\).
\end{proposition}

\begin{proof}
The first identity follows because the oriented boundary of every square
has zero boundary.  The four-cube is contractible, so
\(\ker d_1=\operatorname{im}d_0\).  The vertex-edge incidence matrix on a
connected graph has rank \(|C_0|-1=15\).  Rank-nullity therefore gives
\[
 \operatorname{rank}B
 =|C_1|-\dim\ker B=32-15=17.                                \tag{89.5}
\]
The codimension is \(24-17=7\).  The executable independently obtains the
same rank by exact elimination and records 17 independent face rows.
\end{proof}

\begin{remark}[Bianchi constraints are not optional data]
\label{rem:v18v79-bianchi}
An arbitrary 24-component face-log vector need not lie in \(\mathcal R\).
The seven-dimensional complement is the linear Bianchi obstruction on this
cell.  Any curvature producer must either prove compatibility or project
and charge the incompatible residual; deleting seven coordinates without a
certificate is insufficient.
\end{remark}

\subsection{Exact rational Moore--Penrose construction}

Choose 17 independent rows of \(B\) by exact pivoting and call their matrix
\(R\in M_{17\times32}(\mathbb Z)\).  Every row of \(B\) is a unique linear
combination of the rows of \(R\), giving
\[
 B=LR,qquad
 L\in M_{24\times17}(\mathbb Q),                            \tag{89.6}
\]
where \(L\) has full column rank and \(R\) full row rank.  Define
\[
\boxed{
 K_4
 =R^*(RR^*)^{-1}(L^*L)^{-1}L^*.}                           \tag{89.7}
\]

\begin{theorem}[Exact compatible Hodge inverse]
\label{thm:v18v79-penrose}
The rational matrix \(K_4\in M_{32\times24}(\mathbb Q)\) satisfies
\[
 BK_4B=B,qquad K_4BK_4=K_4,                               \tag{89.8}
\]
\[
 (BK_4)^*=BK_4,qquad(K_4B)^*=K_4B.                        \tag{89.9}
\]
Hence it is the Moore--Penrose inverse \(B^+\).  In particular,
\[
 P_{\mathcal R}:=BK_4                                     \tag{89.10}
\]
is the orthogonal projection onto \(\mathcal R\), and
\[
 BK_4F=F\quad\hbox{for every }F\in\mathcal R.              \tag{89.11}
\]
The reconstructed potential \(K_4F\) is coexact and has minimum
\(\ell^2\) norm among solutions of \(BA=F\).
\end{theorem}

\begin{proof}
Because \(R\) has full row rank,
\(C_R=R^*(RR^*)^{-1}\) satisfies \(RC_R=I_{17}\).  Because \(L\) has full
column rank, \(C_L=(L^*L)^{-1}L^*\) satisfies \(C_LL=I_{17}\).  Thus
\(K_4=C_RC_L\).  Direct substitution into \(B=LR\) gives the first two
Penrose identities.  Moreover
\[
 BK_4=L(L^*L)^{-1}L^*,
 \qquad
 K_4B=R^*(RR^*)^{-1}R,                                    \tag{89.12}
\]
which are symmetric orthogonal projections, proving (89.9).  Equations
(89.10)--(89.11) follow.  The range of \(K_4\) is
\(\operatorname{ran}R^*=\operatorname{ran}B^*\), the orthogonal complement
of \(\ker B\); the same Pythagorean argument as in (86.10) proves minimality.
\end{proof}

\begin{corollary}[Exact compatibility gate]
\label{cor:v18v79-compatibility-gate}
A rational face cochain \(F\) is linearly compatible exactly when
\[
 P_{\mathcal R}F=F.                                       \tag{89.13}
\]
When this holds, \(A=K_4F\) is a canonical rational coexact solution.
\end{corollary}

\begin{proof}
An orthogonal projection fixes exactly its range.  Apply Theorem
\ref{thm:v18v79-penrose}.
\end{proof}

\subsection{Reflection and the infinity stock}

Reflect the four-cube in the first coordinate.  Let \(R_1\) and \(R_2\) be
the signed permutation actions on oriented edges and faces.  An edge or
face reverses orientation exactly when it contains the reflected coordinate.

\begin{theorem}[Four-cube reflection equivariance]
\label{thm:v18v79-reflection}
The exact matrices obey
\[
 BR_1=R_2B,qquad R_1K_4=K_4R_2.                           \tag{89.14}
\]
Thus both the compatible-curvature projection and coexact reconstruction
commute with reflection.
\end{theorem}

\begin{proof}
The first identity is cellular naturality with orientation signs and is
verified entrywise over \(\mathbb Z\).  Orthogonal changes of bases commute
with the Moore--Penrose construction: applying the four Penrose identities
to \(R_1K_4R_2\) shows it is another pseudoinverse of \(B\).  Uniqueness of
the Moore--Penrose inverse proves the second identity.
\end{proof}

\begin{theorem}[Exact four-cube infinity stock]
\label{thm:v18v79-infinity-stock}
The 32-by-24 rational matrix (89.7) satisfies
\[
\boxed{
 \kappa_4^{\rm full}
 :=\|K_4\|_{\infty\to\infty}=\frac34.}                    \tag{89.15}
\]
Therefore every compatible face cochain obeys
\[
 \|K_4F\|_\infty\le\frac34\|F\|_\infty.                 \tag{89.16}
\]
\end{theorem}

\begin{proof}
The induced infinity norm is the maximum of the 32 absolute row sums.  The
executable constructs all entries by exact fraction arithmetic and obtains
maximum \(3/4\).  The full matrix is retained in the hashed payload, so the
claim is reproducible without decimal rounding.  Equation (89.16) is the
induced-norm inequality.
\end{proof}

\begin{assessment}[Full-face stock versus exact used image]
\label{ass:v18v79-full-face-stock}
The value \(3/4\) norms \(K_4\) on all 24 face coordinates, even though
physical linear curvatures lie in the 17-dimensional range \(\mathcal R\).
It is a sound stock, but it need not equal the smaller restricted norm
\(\sup_{F\in\mathcal R}\|K_4F\|_\infty/\|F\|_\infty\).  Computing that
polyhedral used-image norm is a possible further sharpening.
\end{assessment}

\subsection{Nonlinear chart and Wilson insertion}

Use the V77 chart values
\[
 \varepsilon=\frac1{1000},qquad r=\frac1{100},qquad
 B(4r)=\frac{71}{27600},qquad L(r)=\frac{47}{138}.         \tag{89.17}
\]

\begin{theorem}[Executed four-cube nonlinear closure]
\label{thm:v18v79-nonlinear-closure}
With \(\kappa_4^{\rm full}=3/4\), the nonlinear reconstruction has
contraction factor
\[
 c_4=\frac34\frac{47}{138}=\frac{47}{184}<1,              \tag{89.18}
\]
and its self-map image lies in the smaller radius
\[
 r_{4,{\rm im}}
 =\frac34\left(\frac1{1000}+\frac{71}{27600}\right)
 =\frac{493}{184000}<\frac1{100}.                          \tag{89.19}
\]
For one unitary representation and \(d=4,m=1\), the associated Wilson
bound is
\[
\boxed{
 \operatorname{gap}H_1
 \ge1-8r_{4,{\rm im}}
 =\frac{22507}{23000}.}                                    \tag{89.20}
\]
\end{theorem}

\begin{proof}
Equations (89.18)--(89.19) insert (89.15) into the V77 self-map and V78
contraction bounds.  A skew-Hermitian link logarithm of radius
\(r_{4,{\rm im}}\) gives the same link-unitary distance by (87.24).  The V74
mass-one perturbation theorem then gives (89.20).
\end{proof}

\subsection{Executable certificate}

The verifier
\[
 \texttt{scripts/verify\_sm\_v79\_four\_cube\_hodge.py}                 \tag{89.21}
\]
constructs the full cubical incidence matrices, selects the 17 independent
face rows, forms (89.7), and verifies the rank factorization, all four
Penrose identities, both reflection identities, and the nonlinear and gap
fractions.  It emits
\[
 \texttt{artifacts/sm\_v79\_four\_cube\_hodge.json}.                    \tag{89.22}
\]
The hashes are
\[
\begin{array}{c|l}
\hbox{object}&\hbox{SHA--256}\\ \hline
\hbox{verifier}&
\texttt{336AC2C8C95621C21A6C2497894D2EF777A94BC19AC2A77FE9464C5D4DA936BE}\\
\hbox{canonical payload}&
\texttt{75F691E5D7E39FEF918CDB3EAB26FE47D3EC75340A873ACE7643DF63F7D5D126}.
\end{array}                                                   \tag{89.23}
\]

\begin{firewall}[One cell is not a cofinal regulator]
\label{fire:v18v79-one-cell}
The unit four-cube is the first four-dimensional local cell, not a growing
lattice.  Its pseudoinverse does not determine how
\(\|K_n\|_{\infty\to\infty}\) scales with box diameter, lattice spacing,
or boundary conditions.  The nonlinear theorem assumes compatible
principal face logarithms and one represented gauge factor.  It does not
yet combine SU(3), SU(2), and U(1), handle periodic harmonic modes, or prove
that the physical field measure lies in the logarithm chart.
\end{firewall}

\subsection{Acquisition and next bottleneck}

V79 proves that Bianchi compatibility and reflection can be handled exactly
without deleting dependent plaquettes or choosing surface words.  The next
version is V80 and therefore begins with the companion checkpoint.  After
that audit, the substantive continuation should compute a sequence of
larger four-dimensional boxes or derive a discrete Green-kernel estimate
for \(K_n\).  The decisive quantity is the scaling of
\(\kappa_n\varepsilon_n\), not \(\kappa_n\) alone.

\begin{assessment}[V79 acquisition and boundary]
\label{ass:v18v79}
V79 supplies an exact compatible-curvature projection, reflection-
equivariant Moore--Penrose reconstruction, infinity stock \(3/4\), and a
closed nonabelian pilot on the first four-dimensional cell.  Cofinal Hodge
scaling and the full product-gauge combination remain open.
\end{assessment}

\begin{record}[V79 append record]
\label{rec:v18v79}
V79 is appended to the validated V78 whose installed SHA--256 was
\texttt{59B7A9F2B3310807C706340D663BABC9939F484995A50C88C962D8E7D206DAA2}.
After installation only V79 is the active numeric SM note; the frozen
\texttt{V100\_f2} through \texttt{V100\_f17} archives are unchanged.
\end{record}


\section{V80: companion checkpoint and the exact-height gate}
\label{sec:v18v80}

This is the mandated five-version checkpoint following V75.  It reads the
newest active Yang--Mills and Navier--Stokes companion notes before choosing
the next Standard Model--Einstein calculation.  The purpose is methodological:
the companion programmes may supply reusable finite-dimensional estimates,
but no statement is imported unless its hypotheses and conclusion match the
present complex.

\subsection{Authenticated companion inventory}

The inspected Yang--Mills note was
\[
 \texttt{Renewal\_Yang\_Mills\_Mass\_Gap\_Research\_Notes\_V14.tex},
                                                               \tag{90.1}
\]
with SHA--256
\[
 \texttt{E8A71C2B09204402F33E1006F0CE0D11ABF6E4D4D4DD1EFE12CFD71FD228275E}.
                                                               \tag{90.2}
\]
The inspected Navier--Stokes note was
\[
 \texttt{Renewal\_NS\_Stability\_Research\_Notes\_V26.tex},   \tag{90.3}
\]
with SHA--256
\[
 \texttt{82C887F8C2C69E4F28FAD7C3DC8DE5B9099CB5A4BF431EBA1FFF3E91935978AD}.
                                                               \tag{90.4}
\]

\begin{proposition}[Checkpoint transfer]
\label{prop:v18v80-transfer}
The companion inspection yields the following two admissible transfers.
\begin{enumerate}
\item The Yang--Mills V14 coefficient-height argument shows how an exact
integer or rational computation can be reconstructed by a centered Chinese
remainder calculation only after a numerator bound has been proved and the
product of moduli exceeds twice that bound.
\item The Navier--Stokes V26 note continues to support measuring a right
inverse on the structured image actually supplied to it, rather than silently
replacing that domain by an ambient coordinate cube.
\end{enumerate}
Neither companion note proves a new Standard Model mass gap, a uniform
four-dimensional Hodge estimate, or an Einstein continuum limit.
\end{proposition}

\begin{proof}
The Yang--Mills note derives a submultiplicative Laurent coefficient height,
propagates the actual denominators and heights through its response
calculation, and certifies a finite list of 31-bit primes whose product is
large enough for unique centered reconstruction.  Its own firewall separates
that uniqueness statement from the execution of the modular residues and
from a sign theorem.  Thus the transferable content is the conditional
reconstruction lemma, not the response sign.

The Navier--Stokes note has the same authenticated hash as at the preceding
checkpoint and contains no newly appended theorem.  Its relevant lesson is
therefore unchanged.  Finally, the cell complex, operator, and probability
measure in the two companions are different from those of V79.  Logical
specialisation cannot turn their conclusions into any of the three continuum
claims excluded above.
\end{proof}

\subsection{Height certificate required for growing Hodge blocks}

For an integer matrix \(B\in M_{f\times e}(\mathbb Z)\) of rank \(r\), choose
independent rows \(R\in M_{r\times e}(\mathbb Z)\) and write \(B=LR\) as in
(89.6).  The exact pseudoinverse is
\[
 K=R^*(RR^*)^{-1}(L^*L)^{-1}L^*.                         \tag{90.5}
\]
Write each entry in lowest terms, \(K_{ij}=a_{ij}/b_{ij}\), \(b_{ij}>0\), and
define the observed rational height
\[
 h(K)=\max_{i,j}\max\{\operatorname{bitlen}|a_{ij}|,
                           \operatorname{bitlen} b_{ij}\}. \tag{90.6}
\]

\begin{lemma}[Certified centered reconstruction]
\label{lem:v18v80-crt}
Let \(x=a/b\in\mathbb Q\) be in lowest terms, with
\(0<b\le D\) and \(|a|\le A\).  Suppose pairwise coprime primes
\(p_1,\ldots,p_s\) divide neither \(b\) nor a known common denominator
multiple, and put \(M=\prod_kp_k\).  If
\[
 M>2AD,                                                     \tag{90.7}
\]
then \(x\) is the unique rational \(u/v\) with
\(|u|\le A\), \(0<v\le D\), and \(u\equiv xv\pmod M\).
\end{lemma}

\begin{proof}
If \(u/v\) and \(u'/v'\) satisfy the stated bounds and have the same modular
value, then \(uv'-u'v\equiv0\pmod M\).  But
\[
 |uv'-u'v|\le |u|v'+|u'|v\le2AD<M.                        \tag{90.8}
\]
The only multiple of \(M\) in that interval is zero, so \(uv'=u'v\).
Reducedness then gives equality.  Taking one candidate to be \(a/b\) proves
the claim.
\end{proof}

\begin{remark}[What the lemma does not certify]
\label{rem:v18v80-crt-boundary}
A residue that vanishes modulo one prime need not vanish over the rationals.
Lemma \ref{lem:v18v80-crt} applies only after valid numerator and denominator
bounds are available and enough good primes have been accumulated.  The next
pilot will still use direct exact fractions; its reported heights will be data
for deriving such bounds, not a substitute for them.
\end{remark}

\subsection{Decision for the next version}

\begin{decision}[Rectangular four-box pilot]
\label{dec:v18v80-next}
V81 will replace the single cube by the contractible cubical box
\[
 Q_{2,1,1,1}=[0,2]\times[0,1]^3.                          \tag{90.9}
\]
It will construct \(d_0,d_1\) exactly, prove the curvature rank from
contractibility and by elimination, form the Moore--Penrose inverse (90.5),
verify the four Penrose identities and long-axis reflection equivariance,
and report both
\[
 \|K_{2,1,1,1}\|_{\infty\to\infty}
 \quad\hbox{and}\quad h(K_{2,1,1,1}).                    \tag{90.10}
\]
The same V77--V78 chart constants will then decide whether this larger block
still closes the nonlinear fixed-point and Wilson-gap pilot.
\end{decision}

\begin{proof}[Justification of the decision]
The box is the smallest enlargement that tests propagation in one coordinate
while retaining a contractible complex and a reflection symmetry.  Exact
fractions avoid every reconstruction ambiguity at this size.  Comparing its
norm and height with the V79 cube gives the first nontrivial datum for the
cofinal Hodge problem and simultaneously implements both companion lessons in
Proposition \ref{prop:v18v80-transfer}.
\end{proof}

\begin{firewall}[Checkpoint boundary]
\label{fire:v18v80}
This checkpoint proves a reconstruction lemma and fixes an auditable next
calculation.  It does not supply a volume-uniform Hodge bound, a compatible
nonlinear curvature producer, a product-gauge measure, reflection positivity,
or an Einstein continuum construction.  Those remain programme boundaries.
\end{firewall}

\begin{assessment}[V80 acquisition]
\label{ass:v18v80}
V80 converts the companion Yang--Mills height lesson into an explicit gate for
future Hodge computations and prevents modular evidence from being mistaken
for exact rational evidence.  The immediate bottleneck is now the executed
two-cell rectangular pilot specified in Decision \ref{dec:v18v80-next}.
\end{assessment}

\begin{record}[V80 append record]
\label{rec:v18v80}
V80 is appended to the validated V79 whose installed SHA--256 was
\texttt{73590BAE63475DA7F4AD330A1235D7805CF5E0EF283CD9A84E6D9C91AE25BDFC}.
After installation only V80 is the active numeric SM note; the frozen
\texttt{V100\_f2} through \texttt{V100\_f17} archives are unchanged.
\end{record}


\section{V81: exact Hodge inversion on the two-cell four-box}
\label{sec:v18v81}

V79 solved the compatible curvature equation on one four-cube.  V80 then
required the first enlargement to report coefficient height as well as an
operator norm.  This version executes that decision on
\(Q_{2,1,1,1}=[0,2]\times[0,1]^3\).  The calculation is entirely rational:
there is no floating-point rank decision and no modular reconstruction.

\subsection{Cubical counts, exactness, and Bianchi codimension}

Give each positive coordinate edge and each positively ordered coordinate
square its canonical orientation.  With \(L=(2,1,1,1)\), the numbers of
vertices, edges, and faces are
\[
 |C_0|=\prod_{\mu=0}^3(L_\mu+1)=24,                       \tag{91.1}
\]
\[
 |C_1|=\sum_{\mu=0}^3L_\mu\prod_{\nu\ne\mu}(L_\nu+1)
      =16+3\cdot12=52,                                     \tag{91.2}
\]
\[
 |C_2|=\sum_{\mu<\nu}L_\mu L_\nu
                  \prod_{\rho\notin\{\mu,\nu\}}(L_\rho+1)
      =3\cdot8+3\cdot6=42.                                \tag{91.3}
\]
Let
\[
 d_0:C^0(Q_{2,1,1,1})\longrightarrow C^1(Q_{2,1,1,1}),
 \qquad
 B=d_1:C^1(Q_{2,1,1,1})\longrightarrow C^2(Q_{2,1,1,1})  \tag{91.4}
\]
be the signed integer incidence matrices.

\begin{proposition}[Rectangular curvature rank]
\label{prop:v18v81-rank}
The matrices in (91.4) obey
\[
 Bd_0=0,\qquad \operatorname{rank}d_0=23,
 \qquad \operatorname{rank}B=29.                           \tag{91.5}
\]
Consequently the compatible face space
\(\mathcal R_2=\operatorname{im}B\) has dimension \(29\) and codimension
\[
 \dim(C^2/\mathcal R_2)=42-29=13.                          \tag{91.6}
\]
\end{proposition}

\begin{proof}
The oriented boundary of each square has zero boundary, giving \(Bd_0=0\)
entrywise.  The one-skeleton is connected, so its vertex-edge incidence
matrix has rank \(|C_0|-1=23\).  The rectangular box is contractible; hence
its first cellular cohomology vanishes and
\(\ker d_1=\operatorname{im}d_0\).  Rank-nullity gives
\[
 \operatorname{rank}B=52-\dim\ker B=52-23=29.              \tag{91.7}
\]
The exact verifier independently row-reduces both integer matrices and obtains
23 and 29 pivots, so the enumerated incidence convention also checks the
topological count.  Equation (91.6) follows from (91.3).
\end{proof}

\subsection{The exact rectangular pseudoinverse}

Select 29 independent rows of \(B\) by exact pivoting and assemble them as
\(R\in M_{29\times52}(\mathbb Z)\).  Set
\[
 C=R^*(RR^*)^{-1},\qquad L=BC,                             \tag{91.8}
\]
so that \(RC=I_{29}\) and \(B=LR\).  Define
\[
 \boxed{
 K_2=R^*(RR^*)^{-1}(L^*L)^{-1}L^*}
 \in M_{52\times42}(\mathbb Q).                            \tag{91.9}
\]

\begin{theorem}[Exact two-cell compatible Hodge inverse]
\label{thm:v18v81-penrose}
The matrix \(K_2\) satisfies all four Penrose identities
\[
 BK_2B=B,\qquad K_2BK_2=K_2,                               \tag{91.10}
\]
\[
 (BK_2)^*=BK_2,\qquad(K_2B)^*=K_2B.                       \tag{91.11}
\]
Thus \(K_2=B^+\), \(BK_2\) is the orthogonal projection onto
\(\mathcal R_2\), and, for every compatible curvature \(F\),
\[
 BK_2F=F,qquad K_2F\in\operatorname{ran}B^*,              \tag{91.12}
\]
with \(K_2F\) the unique minimum-
\(\ell^2\)-norm solution of \(BA=F\).
\end{theorem}

\begin{proof}
The rows of \(R\) are independent, so \(RR^*\) is invertible and \(RC=I\).
It follows that \(L=BC\) has full column rank and \(B=LR\).  Equation (91.9)
is the full-rank-factorisation formula
\[
 K_2=C(L^*L)^{-1}L^*.                                      \tag{91.13}
\]
Substitution gives (91.10), while
\[
 BK_2=L(L^*L)^{-1}L^*,qquad
 K_2B=R^*(RR^*)^{-1}R                                      \tag{91.14}
\]
are symmetric orthogonal projections, proving (91.11).  The first assertion
of (91.12) is the range-projection identity.  The second follows from
\(\operatorname{ran}K_2=\operatorname{ran}R^*
=\operatorname{ran}B^*=(\ker B)^\perp\).  Every other solution is
\(K_2F+z\) with \(z\in\ker B\), so Pythagoras proves the final assertion.
All inverses and products used here are carried out over \(\mathbb Q\) in
the certificate.
\end{proof}

\subsection{Long-axis reflection}

Let \(s(x_0,x_1,x_2,x_3)=(2-x_0,x_1,x_2,x_3)\).  Its signed permutation
matrices on edges and faces are denoted by \(S_1\) and \(S_2\).  A cell
containing the zeroth coordinate receives sign (-1); a transverse cell
receives sign (+1).

\begin{theorem}[Reflection-equivariant reconstruction]
\label{thm:v18v81-reflection}
The signed actions are involutions and satisfy
\[
 S_1^2=I_{52},\qquad S_2^2=I_{42},
 \qquad BS_1=S_2B,qquad S_1K_2=K_2S_2.                   \tag{91.15}
\]
\end{theorem}

\begin{proof}
Reflection sends a longitudinal edge based at \(x_0=i\) to the oppositely
oriented edge based at \(x_0=1-i\), and a transverse edge based at \(i\) to
the equally oriented edge based at (2-i).  The same rule for faces uses
base (1-i) exactly when the face contains the long coordinate.  Applying
the rule twice proves the first two identities.  Comparing the four oriented
edges of each reflected face proves \(BS_1=S_2B\).  Since signed permutation
matrices are orthogonal, \(S_1K_2S_2\) satisfies the four Penrose identities
for \(B\).  Uniqueness of the Moore--Penrose inverse gives the last identity.
The certificate also verifies all four matrix equalities entrywise.
\end{proof}

\subsection{Infinity stock and rational height}

\begin{theorem}[Executed two-cell stock]
\label{thm:v18v81-stock}
The exact 52-by-42 matrix \(K_2\) has
\[
 \boxed{\|K_2\|_{\infty\to\infty}=1}.                    \tag{91.16}
\]
Among its reduced rational entries, the largest absolute numerator and
largest denominator are
\[
 \max_{i,j}|\operatorname{num}(K_{2,ij})|=1333,
 \qquad
 \max_{i,j}\operatorname{den}(K_{2,ij})=7560.             \tag{91.17}
\]
Consequently the observed numerator, denominator, and total bit heights are
\[
 h_{\rm num}=11,qquad h_{\rm den}=13,qquad h(K_2)=13.    \tag{91.18}
\]
\end{theorem}

\begin{proof}
For a finite matrix the induced infinity norm is the maximum absolute row
sum.  Exact summation of the 52 rows gives maximum (1), establishing
(91.16).  Reducing every entry of (91.9) and scanning its numerator and
positive denominator gives (91.17).  Since
\(2^{10}\le1333<2^{11}\) and
\(2^{12}\le7560<2^{13}\), (91.18) follows.  The full rational matrix is
retained in the canonical payload, making each maximum independently
checkable.
\end{proof}

\begin{assessment}[First scaling datum]
\label{ass:v18v81-scaling-datum}
The full-face stock rises from (3/4) on the V79 one-cube box to (1) on
the two-cell box.  This disproves neither boundedness nor growth: two data
points do not select a cofinal law.  It does show that substituting the
one-cube constant unchanged at larger volume would be invalid.
\end{assessment}

\subsection{Nonlinear chart and represented Wilson gap}

Retain the proved V77 chart constants
\[
 \varepsilon=\frac1{1000},\qquad r=\frac1{100},
 \qquad B(4r)=\frac{71}{27600},
 \qquad L(r)=\frac{47}{138}.                              \tag{91.19}
\]

\begin{theorem}[Executed rectangular nonlinear closure]
\label{thm:v18v81-nonlinear}
With the exact stock (91.16), the fixed-point map has contraction factor
\[
 c_2=\frac{47}{138}<1                                     \tag{91.20}
\]
and image radius
\[
 r_{2,{\rm im}}
 =\frac1{1000}+\frac{71}{27600}
 =\frac{493}{138000}<\frac1{100}.                         \tag{91.21}
\]
For one unitary representation with \(d=4,m=1\), the V74 perturbation theorem
gives
\[
 \boxed{
 \operatorname{gap}H_1\ge1-8r_{2,{\rm im}}
 =\frac{16757}{17250}.}                                    \tag{91.22}
\]
\end{theorem}

\begin{proof}
Multiplying the Lipschitz constant and the self-map bound in (91.19) by
\(\kappa_2=1\) gives (91.20)--(91.21).  The two strict inequalities close the
Banach self-map argument.  The skew-Hermitian exponential estimate then
bounds the represented link perturbation by \(r_{2,{\rm im}}\); insertion in
the mass-one \(d=4\) estimate gives
\(1-8(493/138000)=16757/17250\).
\end{proof}

\subsection{Executable certificate and next bottleneck}

The verifier
\[
 \texttt{scripts/verify\_sm\_v81\_rectangular\_four\_box\_hodge.py}      \tag{91.23}
\]
constructs the cell bases and both incidence matrices, performs exact
elimination and rational matrix inversion, checks the rank factorisation,
all four Penrose identities, the two reflection involutions and both
reflection intertwining identities, and evaluates (91.16)--(91.22).  It emits
\[
 \texttt{artifacts/sm\_v81\_rectangular\_four\_box\_hodge.json}.        \tag{91.24}
\]
The authenticated hashes are
\[
\begin{array}{c|l}
\hbox{object}&\hbox{SHA--256}\\ \hline
\hbox{verifier}&
\texttt{0C0C7965A578460015082929A2901546D1F7404F5DCA3443AAF1331FF234DE86}\\
\hbox{canonical payload}&
\texttt{B233DA162E45298F12BD4B57FFDFDC578869A1C93B6C8F3C6DCAFEDD630336B4}.
\end{array}                                                   \tag{91.25}
\]

\begin{firewall}[Two cells are not a cofinal estimate]
\label{fire:v18v81}
Theorem \ref{thm:v18v81-stock} concerns the ambient face-coordinate norm on
one finite, open-boundary box.  It does not prove a bound uniform in length,
a sharp norm on the compatible image, control of periodic harmonic modes, or
scaling with lattice spacing.  The nonlinear input still assumes compatible
principal face logarithms and one represented gauge factor.  No product
SU(3)--SU(2)--U(1) measure, chiral determinant, Einstein measure, reflection-
positive continuum limit, or nonzero physical mass gap follows from this
pilot.
\end{firewall}

\begin{decision}[Go upstream from repeated inversion]
\label{dec:v18v81-next}
V82 should compute at least the three-cell long box to distinguish immediate
saturation from growth, while also rewriting the long direction as a block
operator.  If coefficient height begins to dominate exact inversion, the
programme should derive a determinant/minor height bound and invoke Lemma
\ref{lem:v18v80-crt}; it should not infer exact entries from an uncertified
modular sample.
\end{decision}

\begin{assessment}[V81 acquisition and boundary]
\label{ass:v18v81}
V81 supplies the first exact volume-enlargement datum: rank (29), thirteen
linear Bianchi constraints, a reflection-equivariant rational pseudoinverse,
stock (1), height (13), and a closed local nonlinear pilot.  Cofinal Hodge
scaling is the active mathematical bottleneck.
\end{assessment}

\begin{record}[V81 append record]
\label{rec:v18v81}
V81 is appended to the validated V80 whose installed SHA--256 was
\texttt{83ED2527EAF4D7943D850843BDC914B22B6CC7589C9A780E0070B2B98E97F587}.
After installation only V81 is the active numeric SM note; the frozen
\texttt{V100\_f2} through \texttt{V100\_f17} archives are unchanged.
\end{record}


\section{V82: the three-cell Hodge datum and the end of immediate saturation}
\label{sec:v18v82}

V81 found full-face Hodge stock (1) on the two-cell long box.  That single
value left open the possibility that the stocks saturate immediately.  This
version executes the next exact rational case,
\[
 Q_{3,1,1,1}=[0,3]\times[0,1]^3,                           \tag{92.1}
\]
and measures both norm growth and arithmetic height.  It reuses the proved
rank-factor pseudoinverse construction of Theorem
\ref{thm:v18v81-penrose}; it does not introduce a different gauge choice.

\subsection{Exact cellular dimensions}

For \(L=(3,1,1,1)\), direct enumeration gives
\[
 |C_0|=4\cdot2^3=32,                                       \tag{92.2}
\]
\[
 |C_1|=3\cdot2^3+3(4\cdot2^2)=24+48=72,                   \tag{92.3}
\]
\[
 |C_2|=3(3\cdot2^2)+3(4\cdot2)=36+24=60.                 \tag{92.4}
\]
Let \(B_3=d_1:C^1\to C^2\) denote the resulting signed face-edge incidence
matrix.

\begin{proposition}[Three-cell curvature rank]
\label{prop:v18v82-rank}
The exact complex satisfies
\[
 d_1d_0=0,qquad \operatorname{rank}d_0=31,qquad
 \operatorname{rank}B_3=41.                               \tag{92.5}
\]
The compatible curvature image therefore has dimension (41), and the
linear Bianchi codimension is
\[
 60-41=19.                                                  \tag{92.6}
\]
\end{proposition}

\begin{proof}
The first identity is the cellular boundary-of-a-boundary identity.  The
one-skeleton is connected, so the incidence rank is (32-1=31).  The box is
contractible, hence
\(\ker d_1=\operatorname{im}d_0\).  Rank-nullity then yields
\(72-31=41\).  The verifier also obtains 31 and 41 pivots by exact rational
elimination, independently checking the enumerated complex.  Subtracting
from (92.4) proves (92.6).
\end{proof}

\subsection{Pseudoinverse, symmetry, norm, and height}

Choose 41 independent rows \(R_3\) of \(B_3\), put
\[
 C_3=R_3^*(R_3R_3^*)^{-1},qquad L_3=B_3C_3,               \tag{92.7}
\]
and define
\[
 K_3=C_3(L_3^*L_3)^{-1}L_3^*.                              \tag{92.8}
\]
The long-axis reflection is \(s_3(x_0,x')=(3-x_0,x')\), with signed edge
and face matrices \(S_{1,3}\) and \(S_{2,3}\).

\begin{theorem}[Executed three-cell Hodge certificate]
\label{thm:v18v82-certificate}
All entries of \(K_3\) are rational, and the exact matrices satisfy
\[
 B_3K_3B_3=B_3,qquad K_3B_3K_3=K_3,                      \tag{92.9}
\]
\[
 (B_3K_3)^*=B_3K_3,qquad(K_3B_3)^*=K_3B_3,              \tag{92.10}
\]
\[
 S_{1,3}^2=I,qquad S_{2,3}^2=I,qquad
 B_3S_{1,3}=S_{2,3}B_3,qquad S_{1,3}K_3=K_3S_{2,3}.      \tag{92.11}
\]
Its full-face infinity stock is
\[
 \boxed{\kappa_3:=\|K_3\|_{\infty\to\infty}=\frac97}.  \tag{92.12}
\]
The coordinate maxima and bit heights of its reduced entries are
\[
 \max|\operatorname{num}K_{3,ij}|=130681,qquad
 \max\operatorname{den}K_{3,ij}=708288,                   \tag{92.13}
\]
\[
 h_{\rm num}(K_3)=17,qquad h_{\rm den}(K_3)=20,qquad
 h(K_3)=20.                                                 \tag{92.14}
\]
\end{theorem}

\begin{proof}
The factorisation in (92.7) has the same full row and column ranks as in
(91.8).  The algebraic proof of Theorem \ref{thm:v18v81-penrose}, applied
with rank (41), proves (92.9)--(92.10).  Cellular naturality gives the third
identity in (92.11), and orthogonal invariance plus uniqueness of the
Moore--Penrose inverse gives the fourth; the first two follow by applying the
reflection twice.

The executable forms every matrix over \(\mathbb Q\), verifies all identities
entrywise, and sums the absolute values in each of the 72 rows of \(K_3\).
The largest exact row sum is (9/7).  A scan of the retained 72-by-60 reduced
fraction matrix gives (92.13).  The inequalities
\[
 2^{16}\le130681<2^{17},qquad
 2^{19}\le708288<2^{20}                                   \tag{92.15}
\]
then prove (92.14).
\end{proof}

\begin{corollary}[Immediate saturation is false]
\label{cor:v18v82-no-saturation}
The full-face Hodge stock on the family \(Q_{n,1,1,1}\) is not constant for
all \(n\ge2\).  Every uniform bound for that family, if one exists, must be at
least (9/7).
\end{corollary}

\begin{proof}
V81 gives \(\kappa_2=1\), whereas (92.12) gives
\(\kappa_3=9/7\ne1\).  The lower bound on any common upper bound is immediate.
\end{proof}

\subsection{Nonlinear closure at length three}

\begin{theorem}[Three-cell nonlinear pilot]
\label{thm:v18v82-nonlinear}
For the V77 data
\(\varepsilon=1/1000\), \(r=1/100\),
\(B(4r)=71/27600\), and \(L(r)=47/138\), the stock (92.12) gives
\[
 c_3=\frac97\frac{47}{138}=\frac{141}{322}<1,             \tag{92.16}
\]
\[
 r_{3,{\rm im}}
 =\frac97\left(\frac1{1000}+\frac{71}{27600}\right)
 =\frac{1479}{322000}<\frac1{100}.                         \tag{92.17}
\]
For one unitary representation with \(d=4,m=1\), the resulting mass-one
Wilson estimate is
\[
 \boxed{
 \operatorname{gap}H_1\ge1-8r_{3,{\rm im}}
 =\frac{38771}{40250}.}                                    \tag{92.18}
\]
\end{theorem}

\begin{proof}
Equations (92.16)--(92.17) are exact multiplication by (9/7); both strict
inequalities close the contraction and self-map conditions.  Applying the
V74 perturbation estimate gives
\(1-8(1479/322000)=38771/40250\), proving (92.18).
\end{proof}

\subsection{Certificate and the upstream turn}

The verifier
\[
 \texttt{scripts/verify\_sm\_v82\_three\_cell\_four\_box\_hodge.py}     \tag{92.19}
\]
emits the full rational matrix and all check data in
\[
 \texttt{artifacts/sm\_v82\_three\_cell\_four\_box\_hodge.json}.       \tag{92.20}
\]
Their hashes are
\[
\begin{array}{c|l}
\hbox{object}&\hbox{SHA--256}\\ \hline
\hbox{verifier}&
\texttt{7F40E02A52D970D6B225155F000984A5F0DC49EA9411B8F9F952011931D368D0}\\
\hbox{canonical payload}&
\texttt{72C7639DCA2A6044C84FB6EC500FD32633C5252002E376722ACECE67A4C3A558}.
\end{array}                                                   \tag{92.21}
\]

\begin{assessment}[Arithmetic growth]
\label{ass:v18v82-height}
The observed rational height rises from (13) bits at length two to (20)
bits at length three.  This remains easy for direct fractions, but it warns
against treating repeated exact inversion as a cofinal proof.  The height
values are measurements, not determinant bounds; Lemma
\ref{lem:v18v80-crt} cannot yet be invoked for uncomputed sizes.
\end{assessment}

\begin{decision}[Block operator before blind enumeration]
\label{dec:v18v82-next}
V83 should expose the translation-invariant interior and finite transverse
modes of \(B_n\), derive an analytic representation of the minimum-norm right
inverse, and use the next finite length only as a check.  If that derivation
stalls, one more exact length may diagnose the correct recurrence, but finite
data will not be promoted to a uniform theorem.
\end{decision}

\begin{firewall}[Three-cell boundary]
\label{fire:v18v82}
The value (9/7) neither proves divergence nor supplies a volume-uniform
bound.  The norm is still taken on all face coordinates rather than sharply
on the compatible image.  Open boundaries avoid periodic harmonic sectors.
The nonlinear pilot retains the compatible logarithm and single-
representation hypotheses, and proves no full Standard Model--Einstein
continuum construction.
\end{firewall}

\begin{assessment}[V82 acquisition]
\label{ass:v18v82}
V82 rules out immediate norm saturation, extends exact reflection-equivariant
Hodge inversion to three cells, and quantifies the simultaneous increase in
operator stock and arithmetic height.  An analytic long-direction reduction
is now preferable to an unstructured sequence of rational inversions.
\end{assessment}

\begin{record}[V82 append record]
\label{rec:v18v82}
V82 is appended to the validated V81 whose installed SHA--256 was
\texttt{8DB8DED37E545F4E012F4B51B0D665090A5DB71B16C4C24A1F64088E83AADFBF}.
After installation only V82 is the active numeric SM note; the frozen
\texttt{V100\_f2} through \texttt{V100\_f17} archives are unchanged.
\end{record}


\section{V83: interval--three-cube block reduction and the fourth stock}
\label{sec:v18v83}

The exact stocks at long lengths one, two, and three are respectively
(3/4,1,9/7).  V82 therefore directed the programme upstream from repeated
matrix inversion.  This version proves an exact tensor-product formula for
the curvature operator at every long length.  A length-four rational
calculation is retained as a diagnostic and a certificate for the reduction.

\subsection{Cochain product decomposition}

Let \(I_n=[0,n]\) have \(n+1\) vertices and \(n\) positively oriented edges,
and let \(X=[0,1]^3\) be the transverse three-cube.  Write
\[
 D_n:C^0(I_n)\longrightarrow C^1(I_n),qquad
 E_0:C^0(X)\longrightarrow C^1(X),qquad
 E_1:C^1(X)\longrightarrow C^2(X)                         \tag{93.1}
\]
for their integer coboundaries.  The transverse dimensions are
\[
 \dim C^0(X)=8,qquad \dim C^1(X)=12,qquad
 \dim C^2(X)=6.                                            \tag{93.2}
\]
Under the canonical product orientations,
\[
 C^1(I_n\times X)
 =\bigl(C^1(I_n)\otimes C^0(X)\bigr)
  \oplus\bigl(C^0(I_n)\otimes C^1(X)\bigr),               \tag{93.3}
\]
\[
 C^2(I_n\times X)
 =\bigl(C^1(I_n)\otimes C^1(X)\bigr)
  \oplus\bigl(C^0(I_n)\otimes C^2(X)\bigr).               \tag{93.4}
\]

\begin{theorem}[Exact long-box block operator]
\label{thm:v18v83-block}
After the cell permutations specified by (93.3)--(93.4), the face-edge
coboundary \(B_n=d_1\) is
\[
\boxed{
 B_n=
 \begin{pmatrix}
  -I_{C^1(I_n)}\otimes E_0 & D_n\otimes I_{C^1(X)}\\
  0                         & I_{C^0(I_n)}\otimes E_1
 \end{pmatrix}.}                                             \tag{93.5}
\]
Its block dimensions are
\[
 (12n+6)\times(20n+12)
 =\bigl(12n+6\bigr)\times\bigl(8n+12(n+1)\bigr).          \tag{93.6}
\]
\end{theorem}

\begin{proof}
For homogeneous cochains the product coboundary is
\[
 d(\alpha\otimes\beta)
 =d_I\alpha\otimes\beta+(-1)^{\deg\alpha}
  \alpha\otimes d_X\beta.                                  \tag{93.7}
\]
A longitudinal one-cochain lies in
\(C^1(I_n)\otimes C^0(X)\); (93.7) sends it to the mixed face
space by (-I\otimes E_0) and has no transverse-face component.  A transverse
one-cochain lies in \(C^0(I_n)\otimes C^1(X)\); its mixed component is
\(D_n\otimes I\), and its transverse component is \(I\otimes E_1\).
This proves (93.5).  The four block sizes are
\[
 12n\times8n,\quad12n\times12(n+1),\quad
 6(n+1)\times8n,\quad6(n+1)\times12(n+1),                 \tag{93.8}
\]
which prove (93.6).  The executable also constructs both sides from their
integer cell incidences and checks equality entrywise at \(n=4\).
\end{proof}

\begin{corollary}[General finite ranks]
\label{cor:v18v83-ranks}
For every \(n\ge1\), the long box has
\[
 |C_0|=8(n+1),\quad |C_1|=20n+12,\quad |C_2|=18n+6,        \tag{93.9}
\]
\[
 \operatorname{rank}B_n=12n+5,qquad
 \operatorname{codim}\operatorname{im}B_n=6n+1.           \tag{93.10}
\]
\end{corollary}

\begin{proof}
Equation (93.9) follows either from the product decomposition or direct cell
counting.  The box is connected and contractible, so
\(\dim\ker B_n=\operatorname{rank}d_0=|C_0|-1=8n+7\).
Rank-nullity gives
\[
 (20n+12)-(8n+7)=12n+5,                                   \tag{93.11}
\]
and subtraction from (|C_2|) gives (6n+1).
\end{proof}

\subsection{Exact length-four certificate}

At \(n=4\), (93.9)--(93.10) become
\[
 |C_0|=40,qquad |C_1|=92,qquad |C_2|=78,qquad
 \operatorname{rank}B_4=53,qquad\operatorname{codim}=25. \tag{93.12}
\]
Let \(\mathcal K_n=B_n^+\) denote the coexact Moore--Penrose reconstruction;
this length subscript is distinct from the dimension notation used for the
single four-cube in V79.

\begin{theorem}[Fourth exact Hodge stock]
\label{thm:v18v83-stock}
The rational 92-by-78 matrix \(\mathcal K_4\) satisfies all four Penrose
identities and intertwines reflection \(x_0\mapsto4-x_0\).  Its ambient
face-coordinate stock is
\[
 \boxed{\|\mathcal K_4\|_{\infty\to\infty}=\frac{15}{11}}. \tag{93.13}
\]
The coordinate maxima and bit heights are
\[
 \max|\operatorname{num}\mathcal K_{4,ij}|=19645877,qquad
 \max\operatorname{den}\mathcal K_{4,ij}=111433080,        \tag{93.14}
\]
\[
 h_{\rm num}=25,qquad h_{\rm den}=27,qquad h(\mathcal K_4)=27. \tag{93.15}
\]
\end{theorem}

\begin{proof}
Apply the full-rank factorisation construction (91.8)--(91.14) with rank 53.
The executable performs the calculation over \(\mathbb Q\), checks the four
Penrose identities and both reflection intertwining identities entrywise,
and takes the maximum of the 92 exact absolute row sums, obtaining (15/11).
It retains every reduced entry.  Scanning those entries gives (93.14), and
\[
 2^{24}\le19645877<2^{25},qquad
 2^{26}\le111433080<2^{27}                                \tag{93.16}
\]
proves (93.15).
\end{proof}

\begin{assessment}[A candidate ceiling, not a theorem]
\label{ass:v18v83-ceiling}
The four exact values are
\[
 \|\mathcal K_1\|_\infty=\frac34,\quad
 \|\mathcal K_2\|_\infty=1,\quad
 \|\mathcal K_3\|_\infty=\frac97,\quad
 \|\mathcal K_4\|_\infty=\frac{15}{11}<\frac32.          \tag{93.17}
\]
They suggest a limiting stock (3/2), but finite agreement below a candidate
ceiling is not a uniform estimate.  Formula (93.5), rather than extrapolation
from (93.17), is the appropriate starting point for a proof.
\end{assessment}

\subsection{Nonlinear consequence at length four}

\begin{theorem}[Fourth nonlinear pilot]
\label{thm:v18v83-nonlinear}
For the V77 chart constants, the stock (93.13) gives
\[
 c_4=\frac{15}{11}\frac{47}{138}=\frac{235}{506}<1,       \tag{93.18}
\]
\[
 r_{4,{\rm im}}
 =\frac{15}{11}\left(\frac1{1000}+\frac{71}{27600}\right)
 =\frac{493}{101200}<\frac1{100}.                          \tag{93.19}
\]
For one unitary representation at \(d=4,m=1\),
\[
 \boxed{
 \operatorname{gap}H_1\ge1-8r_{4,{\rm im}}
 =\frac{12157}{12650}.}                                    \tag{93.20}
\]
\end{theorem}

\begin{proof}
Equations (93.18)--(93.19) insert (15/11) into the V77 contraction and
self-map bounds.  The two strict inequalities close that finite-volume chart.
The V74 mass-one perturbation inequality gives
\(1-8(493/101200)=12157/12650\).
\end{proof}

\subsection{Executable record and next reduction}

The verifier and canonical payload are
\[
 \texttt{scripts/verify\_sm\_v83\_four\_cell\_four\_box\_hodge.py},
 \qquad
 \texttt{artifacts/sm\_v83\_four\_cell\_four\_box\_hodge.json}.       \tag{93.21}
\]
Their hashes are
\[
\begin{array}{c|l}
\hbox{object}&\hbox{SHA--256}\\ \hline
\hbox{verifier}&
\texttt{028F3E1B066DE106D10E134FB8186AD4A69E14C6015936BF9DB9C29EE8DA3A93}\\
\hbox{canonical payload}&
\texttt{C046A64ED568F37267B83586DD75157D23A243256930FBC99AC586E9F023736D}.
\end{array}                                                   \tag{93.22}
\]

\begin{decision}[Transverse mode reduction]
\label{dec:v18v83-next}
V84 should diagonalise the fixed three-cube Hodge data in (93.5).  The target
is a finite list of scalar tridiagonal Green recurrences along \(I_n\), with
special attention to the row that attains the infinity norm.  A closed
recurrence may prove or refute the candidate (3/2) ceiling without forming
the full rational pseudoinverse at every length.
\end{decision}

\begin{firewall}[Block form is not yet an infinity estimate]
\label{fire:v18v83}
Orthogonal transverse diagonalisation preserves the Euclidean norm but need
not preserve the coordinate infinity norm.  Thus (93.5) alone does not prove
\(\|\mathcal K_n\|_\infty\le3/2\).  The long boxes retain open boundaries and
the finite nonlinear pilot retains its compatibility and single-
representation hypotheses.  No full gauge product or Einstein continuum
measure has yet been constructed.
\end{firewall}

\begin{assessment}[V83 acquisition]
\label{ass:v18v83}
V83 replaces an unstructured family of incidence matrices by the exact
interval--three-cube operator (93.5), while the fourth exact computation
supplies stock (15/11) and height 27 as a check.  The remaining Hodge task is
now a finite transverse-mode and one-dimensional Green-kernel problem.
\end{assessment}

\begin{record}[V83 append record]
\label{rec:v18v83}
V83 is appended to the validated V82 whose installed SHA--256 was
\texttt{F57C20294E5052F8AF7FB3006F613A28A56901961459BFE150A8DBFACB3CDDD8}.
After installation only V83 is the active numeric SM note; the frozen
\texttt{V100\_f2} through \texttt{V100\_f17} archives are unchanged.
\end{record}


\section{V84: exact longitudinal Green law and a uniform sector bound}
\label{sec:v18v84}

V83 reduced the long-box curvature map to the interval--three-cube block
operator (93.5).  This version solves its longitudinal-output sector for every
length.  The proof keeps the coordinate infinity norm: a transverse Walsh
expansion is used to identify signs, after which the absolute row sum collapses
to one positive interval Green kernel.  The result proves the candidate
\(3/2\) ceiling for all longitudinal edge rows.  Transverse edge rows are
explicitly left for the next substantive step.

\subsection{Hodge-Laplacian representation}

Let \(D_n,E_0,E_1\) be as in (93.1), and put
\[
 L_n=D_nD_n^*,\qquad L_X=E_0^*E_0.                         \tag{94.1}
\]
Thus \(L_n\) is the \(n\)-by-\(n\) path-edge Laplacian with diagonal 2 and
nearest-neighbour entries \(-1\), while \(L_X\) is the graph Laplacian of the
three-cube.

\begin{lemma}[Coexact Hodge formula]
\label{lem:v18v84-hodge-formula}
On the contractible long box, the one-cochain Hodge Laplacian
\[
 \Delta_{1,n}=d_0d_0^*+B_n^*B_n                            \tag{94.2}
\]
is invertible and
\[
 \mathcal K_n=B_n^+=\Delta_{1,n}^{-1}B_n^*.                \tag{94.3}
\]
In the decomposition (93.3), \(\Delta_{1,n}\) is block diagonal.  Its
longitudinal block is
\[
 \Delta_{\parallel,n}=L_n\otimes I_8+I_n\otimes L_X.       \tag{94.4}
\]
Consequently the longitudinal-output part of the pseudoinverse is
\[
 (\mathcal K_n)_{\parallel}
 =-\Delta_{\parallel,n}^{-1}(I_n\otimes E_0^*)             \tag{94.5}
\]
on mixed face inputs and is zero on transverse face inputs.
\end{lemma}

\begin{proof}
The kernel of the finite-dimensional Hodge Laplacian represents first
cohomology.  The box is contractible, so that kernel is zero.  The range of
\(B_n^*\) is orthogonal to \(\operatorname{im}d_0\), and (94.2) restricts to
\(B_n^*B_n\) on that coexact subspace.  Hence the right side of (94.3) is the
minimum-norm inverse on \(\operatorname{im}B_n\) and vanishes on its
orthogonal complement; equivalently it satisfies the four Penrose identities.

The product Hodge identity cancels the off-diagonal terms produced separately
by \(d_0d_0^*\) and \(d_1^*d_1\).  On
\(C^1(I_n)\otimes C^0(X)\) it gives (94.4).  Taking the adjoint of the first
block row in (93.5) then gives (94.5).
\end{proof}

\subsection{Walsh resolution and exact row sums}

Index transverse vertices by \(x\in\{0,1\}^3\).  For
\(S\subseteq\{1,2,3\}\), define
\[
 \chi_S(x)=(-1)^{\sum_{r\in S}x_r},\qquad
 P_S=\frac18\chi_S\chi_S^*.                                \tag{94.6}
\]
The Walsh characters diagonalise the cube graph Laplacian:
\[
 L_XP_S=2|S|P_S.                                            \tag{94.7}
\]
For \(a>0\), set
\[
 G_a^{(n)}=(L_n+aI_n)^{-1}.                                 \tag{94.8}
\]
Only nonempty \(S\) occur after \(E_0^*\), so the relevant shifts are
\(a=2,4,6\).

\begin{lemma}[Coordinate Walsh coefficients]
\label{lem:v18v84-walsh-coefficients}
Fix a longitudinal edge position \(i\), a transverse vertex \(x\), a mixed
face position \(j\), and a transverse edge with direction \(r\) and tail
\(z\), where \(z_r=0\).  If \(s,t\) are the other two transverse directions
and
\[
 \alpha=(-1)^{x_s+z_s},\qquad
 \beta=(-1)^{x_t+z_t},                                     \tag{94.9}
\]
then, up to the common orientation sign \((-1)^{x_r}\), the corresponding
entry of (94.5) is one quarter of
\[
 \begin{cases}
  G_2(i,j)+2G_4(i,j)+G_6(i,j),&(\alpha,\beta)=(1,1),\\
  G_2(i,j)-G_6(i,j),&\alpha\ne\beta,\\
  G_2(i,j)-2G_4(i,j)+G_6(i,j),&(\alpha,\beta)=(-1,-1).
 \end{cases}                                                \tag{94.10}
\]
Here \(G_a=G_a^{(n)}\).
\end{lemma}

\begin{proof}
By (94.7), the inverse in (94.5) is
\[
 \Delta_{\parallel,n}^{-1}
 =\sum_{S\subseteq\{1,2,3\}}G_{2|S|}^{(n)}\otimes P_S,     \tag{94.11}
\]
with the \(S=\varnothing\) term annihilated by \(E_0^*\).  For an oriented
edge \(e=(r,z)\),
\[
 (P_SE_0^*)_{x,e}
 =\frac18\chi_S(x)\bigl(\chi_S(z+e_r)-\chi_S(z)\bigr).     \tag{94.12}
\]
This is zero unless \(r\in S\), and otherwise equals
\(-\chi_S(x)\chi_S(z)/4\).  The minus sign cancels the sign in (94.5).
Summing the four subsets that contain \(r\), and grouping them by the two
signs in (94.9), gives (94.10).
\end{proof}

\begin{lemma}[No hidden absolute-value cancellation]
\label{lem:v18v84-positivity}
Every entry of the three matrices
\[
 A=G_2+2G_4+G_6,\qquad
 C=G_2-G_6,\qquad
 F=G_2-2G_4+G_6                                             \tag{94.13}
\]
is nonnegative.
\end{lemma}

\begin{proof}
The matrix \(-L_n\) is Metzler, so
\(H_n(t)=\exp(-tL_n)\) is entrywise nonnegative for \(t\ge0\).  The resolvent
formula gives
\[
 G_a=\int_0^\infty e^{-at}H_n(t)\,dt.                      \tag{94.14}
\]
The scalar weights for \(C\) and \(F\) are respectively
\[
 e^{-2t}-e^{-6t}\ge0,qquad
 e^{-2t}(1-e^{-2t})^2\ge0.                                 \tag{94.15}
\]
The weight for \(A\) is plainly positive.  Entrywise integration proves the
claim.
\end{proof}

\begin{proposition}[Longitudinal row-sum collapse]
\label{prop:v18v84-row-collapse}
For every transverse vertex and every interval edge \(i\), the absolute row
sum of the longitudinal row of \(\mathcal K_n\) is
\[
 R_{n,i}=3\sum_{j=1}^n G_2^{(n)}(i,j).                     \tag{94.16}
\]
\end{proposition}

\begin{proof}
For fixed direction \(r\) and fixed \(j\), the four transverse edges in that
direction realise the sign pairs ((1,1),(1,-1),(-1,1),(-1,-1)) exactly once.
By Lemma \ref{lem:v18v84-positivity}, their absolute contributions from
(94.10) sum to
\[
 \frac14(A+2C+F)(i,j)=G_2(i,j).                            \tag{94.17}
\]
There are three transverse directions.  Summing over \(j\) proves (94.16).
The zero transverse-face block in (94.5) adds nothing.
\end{proof}

\subsection{Closed recurrence and the uniform ceiling}

Let \(u=G_2^{(n)}\mathbf1\).  Its components solve
\[
 4u_i-u_{i-1}-u_{i+1}=1\quad(1\le i\le n),                \tag{94.18}
\]
where the boundary convention is \(u_0=u_{n+1}=0\).  Put
\(\rho=2+\sqrt3\), so \(\rho^{-1}=2-\sqrt3\).

\begin{theorem}[Exact longitudinal stock at every length]
\label{thm:v18v84-longitudinal-stock}
Define positive integers \(q_n\) by
\[
 q_1=4,\quad q_2=6,\quad q_3=14,\quad q_4=22,
 \qquad q_n=4q_{n-2}-q_{n-4}\quad(n\ge5).                 \tag{94.19}
\]
Then the maximum longitudinal row norm is
\[
\boxed{
 \kappa_{\parallel,n}
 :=\max_{\text{longitudinal rows}}
      \sum_f|\mathcal K_n(e,f)|
 =\frac32-\frac3{q_n}<\frac32.}                           \tag{94.20}
\]
The maximum occurs at the central interval edge, or at the two central edges
when \(n\) is even, and
\[
 \lim_{n\to\infty}\kappa_{\parallel,n}=\frac32.           \tag{94.21}
\]
More explicitly, for \(n=2m-1\) and \(n=2m\), respectively,
\[
 q_{2m-1}=\rho^m+\rho^{-m},\qquad
 q_{2m}=\frac{2(\rho^{m+1}+\rho^{-m})}{1+\rho}.             \tag{94.22}
\]
\end{theorem}

\begin{proof}
Set \(v_i=1/2-u_i\).  Equation (94.18) becomes
\[
 v_{i+1}=4v_i-v_{i-1},\qquad v_0=v_{n+1}=\frac12.          \tag{94.23}
\]
Solving this second-order recurrence gives
\[
 v_i=\frac{\rho^i+\rho^{n+1-i}}{2(1+\rho^{n+1})}.          \tag{94.24}
\]
It is symmetric and is smallest at the central index or indices.  At those
indices, (94.24) is \(1/q_n\) with \(q_n\) as in (94.22).  Since
\(\rho+\rho^{-1}=4\), both parity subsequences in (94.22) satisfy (94.19);
the displayed initial values show in particular that all \(q_n\) are
integers.  Proposition \ref{prop:v18v84-row-collapse} now gives
\(3\max_i u_i=3(1/2-1/q_n)\), proving (94.20).
The recurrence or (94.22) shows \(q_n\to\infty\), which proves (94.21).
\end{proof}

\begin{corollary}[Uniform longitudinal chart margin]
\label{cor:v18v84-sector-chart}
Restricted to a source and output for which only the longitudinal block
(94.5) is used, the V77 constants obey uniformly in \(n\)
\[
 c_{\parallel,n}<\frac32\frac{47}{138}=\frac{47}{92}<1,   \tag{94.25}
\]
\[
 r_{\parallel,n}<\frac32\frac{493}{138000}
 =\frac{493}{92000}<\frac1{100}.                           \tag{94.26}
\]
The corresponding sectoral mass-one Wilson estimate is strictly larger than
\[
 1-8\frac{493}{92000}=\frac{11007}{11500}.                 \tag{94.27}
\]
\end{corollary}

\begin{proof}
Insert the strict bound (94.20) into the V77--V78 contraction, self-map, and
V74 perturbation inequalities.  The restriction in the statement is needed
because no transverse-row bound has yet been proved.
\end{proof}

\subsection{Certificate and remaining Hodge gate}

The verifier
\[
 \texttt{scripts/verify\_sm\_v84\_longitudinal\_green\_formula.py}     \tag{94.28}
\]
checks the exact \(G_2,G_4,G_6\) sign combinations, row-sum collapse, central
maximisers, recurrence values through length 12, and agreement with the four
retained full pseudoinverses.  It emits
\[
 \texttt{artifacts/sm\_v84\_longitudinal\_green\_formula.json}.        \tag{94.29}
\]
The hashes are
\[
\begin{array}{c|l}
\hbox{object}&\hbox{SHA--256}\\ \hline
\hbox{verifier}&
\texttt{8933D89A28219DA29EBD3EB350FE41C44FC4DD7952BE8193E7AD183E2F6520F8}\\
\hbox{canonical payload}&
\texttt{8940957C86F71644515FBC52C5530625DAD665BDF759BF67D098C53B294A6D00}.
\end{array}                                                   \tag{94.30}
\]

\begin{firewall}[Sector theorem versus full Hodge theorem]
\label{fire:v18v84}
Theorem \ref{thm:v18v84-longitudinal-stock} is uniform and exact, but it
controls only longitudinal output rows.  It does not bound transverse rows of
\(\mathcal K_n\), so Corollary \ref{cor:v18v84-sector-chart} cannot yet be
applied to an unrestricted compatible curvature source.  Open boundaries,
principal logarithms, the product gauge measure, chiral fermions, and the
Einstein continuum limit remain outside the theorem.
\end{firewall}

\begin{decision}[Checkpoint, then transverse Green rows]
\label{dec:v18v84-next}
V85 is a companion checkpoint.  The substantive V86 continuation should use
the transverse block of the product Hodge Laplacian to obtain a coordinate
row-sum formula analogous to (94.16), or else isolate the precise transverse
mode whose sign pattern prevents it.
\end{decision}

\begin{assessment}[V84 acquisition]
\label{ass:v18v84}
V84 proves the exact formula \(3/2-3/q_n\) and a strict uniform \(3/2\) ceiling
for every longitudinal row of the long-box Moore--Penrose reconstruction.
The cofinal Hodge bottleneck has been reduced to transverse output rows.
\end{assessment}

\begin{record}[V84 append record]
\label{rec:v18v84}
V84 is appended to the validated V83 whose installed SHA--256 was
\texttt{EC050535C8F4DAC6ED17F7F07C89FECAF5380C8FEA5AC8ABC5BB6923231E61A4}.
After installation only V84 is the active numeric SM note; the frozen
\texttt{V100\_f2} through \texttt{V100\_f17} archives are unchanged.
\end{record}


\section{V85: companion checkpoint before the transverse-row proof}
\label{sec:v18v85}

This is the mandated checkpoint five versions after V80.  It was performed
after installing the uniform longitudinal theorem V84 and before beginning the
transverse-output calculation.  The companion documents are treated as
mathematical sources only; their programme directions do not override the
present Standard Model--Einstein workflow.

\subsection{Authenticated companion state}

The newest active Yang--Mills note is
\[
 \texttt{Renewal\_Yang\_Mills\_Mass\_Gap\_Research\_Notes\_V17.tex},
                                                               \tag{95.1}
\]
with SHA--256
\[
 \texttt{C0D6AC6B4F4F527220A10F2590DB9254B792C0539B1ACAB104C3271CF4983260}.
                                                               \tag{95.2}
\]
The newest active Navier--Stokes note remains
\[
 \texttt{Renewal\_NS\_Stability\_Research\_Notes\_V26.tex},           \tag{95.3}
\]
with unchanged SHA--256
\[
 \texttt{82C887F8C2C69E4F28FAD7C3DC8DE5B9099CB5A4BF431EBA1FFF3E91935978AD}.
                                                               \tag{95.4}
\]

\begin{proposition}[V85 audit result]
\label{prop:v18v85-audit}
The Yang--Mills V15--V17 additions supply three admissible computational
principles for the present Hodge problem:
\begin{enumerate}
\item contributions with a common output index must be summed before a zero,
rank, sign, or compression decision is made;
\item a low-rank or factored route may replace a dense multiplication only
when its product is algebraically identical and independently checked; and
\item a one-prime cancellation remains a modular statement until a certified
CRT lift is completed.
\end{enumerate}
The unchanged Navier--Stokes note supplies no new theorem at this checkpoint.
Neither companion proves the missing transverse Hodge estimate or a Standard
Model--Einstein continuum limit.
\end{proposition}

\begin{proof}
Yang--Mills V15 forms complete frequency coefficients of \(E^2\) before
computing their ranks and explicitly distinguishes two whole-matrix modular
cancellations from rational zeros.  V16 constructs canonical factors, verifies
them by multiplication, routes the \(E^3\) products through those factors,
and then recomputes the ranks of the completed aggregate coefficients.  V17
does the same for \(E^4\) and the finite Neumann sum, freezing the fourth-power
object before in-place addition.  These facts prove the three stated
principles.  They are rules about exact finite algebra, so they transfer to an
exact Hodge calculation.  The companion operators and continuum claims differ,
so no physical theorem transfers.  The NS hash equality with (90.4) proves
that no new NS appendix has appeared since V80.
\end{proof}

\subsection{Application to the open transverse gate}

The transverse output of
\(\mathcal K_n=\Delta_{1,n}^{-1}B_n^*\) contains both mixed-face and
transverse-face contributions.  A transverse spectral decomposition may be
used to evaluate it, but the coordinate coefficient must have the form
\[
 k(e,f)=\sum_{\sigma}k_\sigma(e,f)                         \tag{95.5}
\]
before its row norm is formed:
\[
 \sum_f\left|\sum_\sigma k_\sigma(e,f)\right|.             \tag{95.6}
\]
In general, replacing (95.6) by
\(\sum_{f,\sigma}|k_\sigma(e,f)|\) destroys cancellations and proves only a
possibly much larger upper bound.

\begin{lemma}[Aggregate-before-norm rule]
\label{lem:v18v85-aggregate-norm}
For any finite family in (95.5),
\[
 \sum_f|k(e,f)|
 \le\sum_{f,\sigma}|k_\sigma(e,f)|,                        \tag{95.7}
\]
and equality holds exactly when, for every fixed \(f\), all nonzero real
summands \(k_\sigma(e,f)\) have a common sign.
\end{lemma}

\begin{proof}
Apply the scalar triangle inequality at each \(f\) and sum.  Equality in the
real scalar triangle inequality holds precisely when the nonzero summands
have a common sign.  Enforcing that condition independently for each \(f\)
gives the equality criterion.
\end{proof}

\begin{decision}[V86 transverse aggregation protocol]
\label{dec:v18v85-next}
V86 will derive the transverse block of the product Hodge Laplacian, sum all
transverse spectral contributions in the physical edge and face bases, and
only then take absolute row sums.  The exact full pseudoinverses at lengths one
through four will serve as regression checks.  If a uniform coordinate sign
rule emerges, it will be proved using positive interval resolvents; otherwise
the note will retain the rigorous aggregate triangle bound and identify the
specific mode interference that prevents equality.
\end{decision}

\begin{firewall}[Checkpoint boundary]
\label{fire:v18v85}
This audit changes the order of proof obligations, not the physical claim
boundary.  It supplies no full Hodge stock, no cofinal nonlinear chart, no
product-gauge measure, no fermion determinant, and no Einstein continuum
measure.  The Yang--Mills one-prime response programme and the SM long-box
Hodge programme remain logically separate.
\end{firewall}

\begin{assessment}[V85 acquisition]
\label{ass:v18v85}
V85 imports the exact aggregate-before-decision discipline needed for the
transverse coordinate norm and fixes a concrete V86 protocol.  The mandatory
companion inspection is complete, and the next version is substantive.
\end{assessment}

\begin{record}[V85 append record]
\label{rec:v18v85}
V85 is appended to the validated V84 whose installed SHA--256 was
\texttt{7B04BE84500C16D548D83A7D3BBCBA60C273A203A0F710B3BE2597D9D5915B3A}.
After installation only V85 is the active numeric SM note; the frozen
\texttt{V100\_f2} through \texttt{V100\_f17} archives are unchanged.
\end{record}


\section{V86: transverse Green variation and a uniform full-strip chart}
\label{sec:v18v86}

V84 proved a strict (3/2) bound for every longitudinal output row of the
long-box Moore--Penrose inverse.  V85 required the transverse spectral
contributions to be aggregated in the coordinate basis before absolute values
were taken.  This version performs that aggregation exactly.  Each transverse
orientation becomes a scalar massive Green problem on a three-dimensional
prism, and its coordinate row norm becomes the total variation of one positive
Green row.  A degree--mass estimate then gives a length-independent bound.

\subsection{Transverse orientation blocks}

Fix one transverse direction \(r\in\{1,2,3\}\), and write \(s,t\) for the
other two.  An \(r\)-oriented edge of the transverse cube is determined by the
two coordinates \((x_s,x_t)\in\{0,1\}^2\).  Thus the \(r\)-orientation
component of transverse one-cochains on \(I_n\times[0,1]^3\) is naturally
\[
 C^0\(P_{n+1}\)\otimes C^0\(Q_2\),                             \tag{96.1}
\]
where \(P_{n+1}\) is the path on interval vertices \(0,\ldots,n\) and
\(Q_2\) is the four-vertex square.

Let \(A_n=D_n^*D_n\) be the vertex graph Laplacian of \(P_{n+1}\), and let
\(L_{Q_2}\) be the vertex graph Laplacian of the square.

\begin{lemma}[Scalar transverse Hodge block]
\label{lem:v18v86-transverse-block}
The transverse part of the one-cochain Hodge Laplacian preserves each of the
three edge-orientation components.  On the \(r\)-component it is
\[
 H_n=A_n\otimes I_4+I_{n+1}\otimes L_{Q_2}+2I
     =L_{Y_n}+2I,                                          \tag{96.2}
\]
where
\[
 Y_n=P_{n+1}\mathbin{\square}Q_2                           \tag{96.3}
\]
is the Cartesian product graph.  There is no coupling between different
transverse edge orientations.
\end{lemma}

\begin{proof}
Apply the product Hodge identity used in Lemma
\ref{lem:v18v84-hodge-formula} once along the long interval and once in each
of the three transverse intervals.  A one-cochain oriented in direction \(r\)
receives the edge Hodge Laplacian \(D_1D_1^*=[2]\) in that unit direction and
the vertex graph Laplacian in each of the other three directions.  The two
remaining transverse vertex factors combine as \(L_{Q_2}\), and the long
vertex factor is \(A_n\).  Product-Hodge cross terms cancel, so orientation
components do not mix.  This gives (96.2); the standard Cartesian graph
identity gives (96.3).
\end{proof}

\subsection{Faces as prism gradients}

Let \(J_n=H_n^{-1}\).  The mixed faces of orientation ((0,r)) correspond
to the long-direction edges of \(Y_n\).  The transverse faces of orientations
\((r,s)\) and \((r,t)\) correspond to the two square directions of \(Y_n\).
Faces not containing \(r\) have no component in the \(r\)-orientation block.
Choose any orientation of the graph edges of \(Y_n\), and let
\(\partial_{Y_n}:C^0\(Y_n\)\to C^1\(Y_n\)\) be its incidence matrix.

\begin{proposition}[Aggregate transverse row identity]
\label{prop:v18v86-tv-identity}
Up to signed permutations of face columns, the \(r\)-orientation output block
of the pseudoinverse is
\[
 (\mathcal K_n)_{\perp,r}=J_n\partial_{Y_n}^*.             \tag{96.4}
\]
Consequently, for an output edge represented by a prism vertex \(x\), its
full coordinate row norm is
\[
 R_{\perp,n}\(x\)
 =\sum_{\{u,v\}\in E(Y_n)}|J_n(x,u)-J_n(x,v)|.            \tag{96.5}
\]
The sum in (96.5) includes both mixed and transverse face columns and is taken
only after their physical coordinate coefficients have been formed.
\end{proposition}

\begin{proof}
By Lemma \ref{lem:v18v86-transverse-block}, the transverse output of
\(\Delta_{1,n}^{-1}B_n^*\) is obtained by applying \(J_n\) within the fixed
orientation component.  On a mixed face, \(B_n^*\) is the adjoint interval
difference \(D_n^*\); on a face containing \(r\) and one of \(s,t\), it is the
adjoint difference in that square direction, with a sign fixed by the product
orientation.  Together these three families form the full prism incidence
adjoint.  Changing face conventions only signs or permutes columns, proving
(96.4).  A column associated with an oriented prism edge \(u\to v\) has row
entry \(J_n(x,v)-J_n(x,u)\).  Summing its absolute value over all prism edges
gives (96.5).  Other face orientations contribute zero because the Hodge
block does not mix edge orientations.
\end{proof}

\subsection{Positive Green mass and the uniform bound}

\begin{lemma}[Massive prism Green facts]
\label{lem:v18v86-green-facts}
The matrix \(J_n=(L_{Y_n}+2I)^{-1}\) is entrywise nonnegative, every row has
sum \(1/2\), and every vertex of \(Y_n\) has degree at most four.
\end{lemma}

\begin{proof}
The matrix \(-L_{Y_n}\) is Metzler, hence
\(\exp(-tL_{Y_n})\) is entrywise nonnegative.  Therefore
\[
 J_n=\int_0^\infty e^{-2t}\exp(-tL_{Y_n})\,dt              \tag{96.6}
\]
is entrywise nonnegative.  Since \(L_{Y_n}\mathbf1=0\),
\[
 (L_{Y_n}+2I)\frac12\mathbf1=\mathbf1,                    \tag{96.7}
\]
so \(J_n\mathbf1=\mathbf1/2\); symmetry turns this into the asserted row-sum
statement.  A prism vertex has at most two long neighbours and exactly two
square neighbours, hence degree at most four.
\end{proof}

\begin{theorem}[Uniform transverse and full Hodge stocks]
\label{thm:v18v86-uniform-stock}
For every \(n\ge1\),
\[
 \kappa_{\perp,n}:=\max_xR_{\perp,n}\(x\)\le2.              \tag{96.8}
\]
Together with Theorem \ref{thm:v18v84-longitudinal-stock}, this gives the
complete ambient face-coordinate estimate
\[
\boxed{
 \|\mathcal K_n\|_{\infty\to\infty}\le2
 \qquad\text{for every }n\ge1.}                            \tag{96.9}
\]
\end{theorem}

\begin{proof}
Fix \(x\) and write \(g(u)=J_n(x,u)\ge0\).  Proposition
\ref{prop:v18v86-tv-identity} and the scalar triangle inequality give
\[
\begin{split}
 R_{\perp,n}\(x\)
 &\le\sum_{\{u,v\}\in E(Y_n)}(g(u)+g(v))\\
 &=\sum_{u\in V(Y_n)}\deg(u)g(u)
 \le4\sum_ug(u)=2,                                        \tag{96.10}
\end{split}
\]
where the last equality uses Lemma \ref{lem:v18v86-green-facts}.  This proves
(96.8).  Every output row is either longitudinal or one of the three
transverse orientations.  The former has norm strictly below (3/2) by V84;
the latter has norm at most two by (96.8).  Taking the maximum proves (96.9).
\end{proof}

\begin{remark}[Exact diagnostic values]
\label{rem:v18v86-values}
The transverse Green total-variation stocks at lengths one through four are
\[
 \frac34,\qquad\frac9{10},\qquad\frac{27}{28},\qquad
 \frac{39}{38}.                                            \tag{96.11}
\]
They agree exactly with the maximum transverse rows of the four retained full
pseudoinverses.  The verifier also evaluates the scalar prism through length
12.  These finite values check Proposition \ref{prop:v18v86-tv-identity}; the
uniform theorem itself follows from (96.10), independently of that sample.
\end{remark}

\subsection{Uniform nonlinear strip closure}

\begin{theorem}[Cofinal long-strip chart]
\label{thm:v18v86-strip-chart}
With the V77 data
\(\varepsilon=1/1000\), \(r=1/100\),
\(B(4r)=71/27600\), and \(L(r)=47/138\), the full Hodge stock (96.9) gives,
uniformly in the strip length,
\[
 c_n\le2\frac{47}{138}=\frac{47}{69}<1,                   \tag{96.12}
\]
\[
 r_{n,{\rm im}}
 \le2\left(\frac1{1000}+\frac{71}{27600}\right)
 =\frac{493}{69000}<\frac1{100}.                           \tag{96.13}
\]
For one unitary representation with \(d=4,m=1\), the corresponding finite
strip mass-one estimate is
\[
 \boxed{
 \operatorname{gap}H_1\ge1-8\frac{493}{69000}
 =\frac{8132}{8625}.}                                      \tag{96.14}
\]
\end{theorem}

\begin{proof}
Equations (96.12)--(96.13) insert the uniform bound (96.9) into the V77--V78
contraction and self-map estimates.  Their strict inequalities close the
fixed-point chart at every strip length with common constants.  The V74
mass-one perturbation estimate gives (96.14) after exact reduction.
\end{proof}

\subsection{Executable certificate}

The exact verifier
\[
 \texttt{scripts/verify\_sm\_v86\_transverse\_green\_total\_variation.py}
                                                               \tag{96.15}
\]
constructs the prism Laplacian over \(\mathbb Q\), checks Green positivity,
row mass, degree and total-variation bounds through length 12, and compares
(96.5) with the transverse rows of every retained full pseudoinverse.  It emits
\[
 \texttt{artifacts/sm\_v86\_transverse\_green\_total\_variation.json}.
                                                               \tag{96.16}
\]
The authenticated hashes are
\[
\begin{array}{c|l}
\hbox{object}&\hbox{SHA--256}\\ \hline
\hbox{verifier}&
\texttt{D0569906971DB034EBB19C822A8523433A520AD539D174A951AE5BF254E21C64}\\
\hbox{canonical payload}&
\texttt{910FC7726F38CB176F2282D4F44DACF0865F2F268DA47F5F0B4C19FDD2517728}.
\end{array}                                                   \tag{96.17}
\]

\begin{firewall}[Strip cofinality is not four-dimensional cofinality]
\label{fire:v18v86}
Theorem \ref{thm:v18v86-strip-chart} is uniform only while the three
transverse side lengths remain one.  When every side grows, the edge Hodge
factor in a chosen orientation loses the fixed mass two used in (96.2), and
the present proof does not give a volume-independent inverse.  The theorem
also assumes compatible principal logarithms and one represented gauge
factor.  It does not construct the SU(3)--SU(2)--U(1) measure, chiral fermion
weight, Einstein measure, reflection-positive limit, or physical mass gap.
\end{firewall}

\begin{decision}[Pass from strips to growing four-boxes]
\label{dec:v18v86-next}
V87 should formulate the orientation-block Green representation on a general
box \(\prod_{\mu=0}^3[0,L_\mu]\).  It should identify the exact dependence on
the oriented side length, derive a diameter-dependent coordinate bound, and
state the curvature smallness scaling needed for the nonlinear chart.  This
moves upstream to the compensating product \(\kappa_L\varepsilon_L\) rather
than seeking a false volume-independent constant.
\end{decision}

\begin{assessment}[V86 acquisition]
\label{ass:v18v86}
V86 closes the complete open-strip Hodge and nonlinear chart uniformly in its
long length.  The decisive next problem is the simultaneous growth of all four
side lengths and the scaling of curvature smallness against the resulting
Green variation.
\end{assessment}

\begin{record}[V86 append record]
\label{rec:v18v86}
V86 is appended to the validated V85 whose installed SHA--256 was
\texttt{E18DA73B5E740337952C44EC41C00E2D296AA60B518CF785547D9D6578DD0EBF}.
After installation only V86 is the active numeric SM note; the frozen
\texttt{V100\_f2} through \texttt{V100\_f17} archives are unchanged.
\end{record}


\section{V87: general rectangular Green blocks and an explicit cofinal chart}
\label{sec:v18v87}

V86 closed boxes whose three transverse side lengths remain one.  This version
allows all four integer side lengths to vary.  The one-cochain Hodge Laplacian
still separates by edge orientation, but its fixed mass is lost when the
oriented side grows.  A positive Green row gives a rigorous quadratic side-
length bound for the ambient Moore--Penrose inverse.  The V77 rational
majorants then give an explicit, though deliberately conservative, curvature
scaling that closes every finite rectangular chart.

\subsection{Orientation-separated Hodge Laplacian}

Let
\[
 Q_L=\prod_{\nu=0}^3[0,L_\nu],\qquad L_\nu\in\mathbb N,   \tag{97.1}
\]
with open cubical boundary.  For an interval of length \(m\), write
\[
 A_m=D_m^*D_m\quad\hbox{on vertices},\qquad
 E_m=D_mD_m^*\quad\hbox{on edges}.                         \tag{97.2}
\]
An edge of orientation \(\mu\) is indexed by
\[
 \Omega_{\mu,L}
 =\{0,\ldots,L_\mu-1\}
   \times\prod_{\nu\ne\mu}\{0,\ldots,L_\nu\}.           \tag{97.3}
\]

\begin{theorem}[General orientation Hodge block]
\label{thm:v18v87-orientation-block}
The one-cochain Hodge Laplacian is the orthogonal direct sum of four
orientation blocks.  On the \(\mu\)-oriented block it is
\[
\boxed{
 H_{\mu,L}
 =E_{L_\mu}\otimes\bigotimes_{\nu\ne\mu}I_{L_\nu+1}
 +\sum_{\nu\ne\mu}
   I_{L_\mu}\otimes A_{L_\nu}
   \otimes\bigotimes_{\rho\notin\{\mu,\nu\}}I_{L_\rho+1}.} \tag{97.4}
\]
Here each tensor factor is placed in coordinate order.  Different edge
orientations have zero off-diagonal Hodge blocks.
\end{theorem}

\begin{proof}
Apply the tensor-product Hodge identity to the four interval complexes.  A
one-cochain in the \(\mu\) direction has degree one in that interval factor and
degree zero in the other three.  It therefore receives \(E_{L_\mu}\) in the
oriented factor and \(A_{L_\nu}\) in every vertex factor.  The cross terms from
\(d_0d_0^*\) and \(d_1^*d_1\) cancel with their Koszul signs.  This proves
(97.4) and the absence of orientation coupling.  The executable independently
constructs the full integer Hodge matrix and verifies both statements
entrywise in three nonisomorphic boxes.
\end{proof}

\subsection{Green variation and row mass}

Put \(G_{\mu,L}=H_{\mu,L}^{-1}\).  On \(\Omega_{\mu,L}\), join two states when
they differ by one in a coordinate \(\nu\ne\mu\).  Denote this complementary
edge set by \(E_{\mu,L}^{\perp}\).

\begin{proposition}[General coordinate row identity]
\label{prop:v18v87-row-identity}
The \(\mu\)-orientation row of the Moore--Penrose curvature inverse satisfies
\[
 R_{\mu,L}\(x\)
 =\sum_{\{u,v\}\in E_{\mu,L}^{\perp}}
       |G_{\mu,L}(x,u)-G_{\mu,L}(x,v)|.                   \tag{97.5}
\]
The Green matrix is entrywise nonnegative.  If the oriented edge coordinate
of \(x\) is \(j-1\), where \(1\le j\le L_\mu\), then its row mass is
\[
 \sum_{u\in\Omega_{\mu,L}}G_{\mu,L}(x,u)
 =\frac{j(L_\mu+1-j)}2.                                   \tag{97.6}
\]
\end{proposition}

\begin{proof}
Lemma \ref{lem:v18v84-hodge-formula} gives
\(B_L^+=\Delta_{1,L}^{-1}B_L^*\).  A face contributes to a \(\mu\)-oriented
edge only when its orientation contains \(\mu\); the second face direction is
some \(\nu\ne\mu\), and \(B_L^*\) is the signed adjoint difference in that
coordinate.  Thus, up to signed column permutations, the row coefficients are
the complementary differences of a row of \(G_{\mu,L}\).  This proves (97.5)
after all coordinate contributions have been aggregated.

The negative of \(H_{\mu,L}\) has nonnegative off-diagonal entries, and
\(H_{\mu,L}\) is positive definite, so the heat-integral argument of (96.6)
proves Green nonnegativity.  Let
\[
 w_j=\frac{j(L_\mu+1-j)}2,qquad 1\le j\le L_\mu.          \tag{97.7}
\]
Direct substitution gives \(E_{L_\mu}w=\mathbf1\).  The other three vertex
Laplacians annihilate constants.  Hence (97.4) sends the function
\(w_j\), constant in the complementary coordinates, to \(\mathbf1\).
Uniqueness of the inverse proves (97.6).
\end{proof}

\begin{theorem}[Quadratic ambient rectangular bound]
\label{thm:v18v87-quadratic-bound}
For every orientation,
\[
 \max_xR_{\mu,L}\(x\)
 \le3\left\lfloor\frac{(L_\mu+1)^2}{4}\right\rfloor.      \tag{97.8}
\]
With \(L_{\max}=\max_\mu L_\mu\), the full pseudoinverse obeys
\[
\boxed{
 \|B_L^+\|_{\infty\to\infty}
 \le C_L:=3\left\lfloor\frac{(L_{\max}+1)^2}{4}\right\rfloor.} \tag{97.9}
\]
\end{theorem}

\begin{proof}
Every state in \(\Omega_{\mu,L}\) has at most two neighbours in each of the
three complementary directions, so its complementary degree is at most six.
For a fixed nonnegative Green row \(g\), (97.5) gives
\[
 R_{\mu,L}\(x\)
 \le\sum_{\{u,v\}\in E_{\mu,L}^{\perp}}(g(u)+g(v))
 \le6\sum_ug(u).                                           \tag{97.10}
\]
The maximum of \(j(L_\mu+1-j)/2\) is
\(\frac12\lfloor(L_\mu+1)^2/4\rfloor\).  Equations
(97.6) and (97.10) prove (97.8).  Taking the maximum over the four orientation
blocks proves (97.9).
\end{proof}

\begin{remark}[The estimate deliberately charges cancelled modes]
\label{rem:v18v87-crude}
The bound (97.10) replaces a difference by a sum and therefore charges the
complement-constant Green mode even though the face difference annihilates
it.  This explains why (97.9) is much larger than the exact strip bounds.  It
is a valid cofinal stock, not a sharp one.
\end{remark}

\subsection{An explicit size-dependent nonlinear chart}

Recall, for \(S=4r<1/2\),
\[
 B(S)=\frac{S^2}{1-S}+
      \frac{S^2}{2(1-S)(1-2S)},                            \tag{97.11}
\]
\[
 \Lambda(r):=4\left(\frac{S}{1-S}+\frac{S}{1-2S}\right).  \tag{97.12}
\]
The notation \(\Lambda\) distinguishes this Lipschitz majorant from the side
vector \(L\).  Let
\[
 r_L=\frac1{100C_L},\qquad
 \varepsilon_L=\frac1{1000C_L^2}.                         \tag{97.13}
\]

\begin{lemma}[Scaled rational majorants]
\label{lem:v18v87-scaled-majorants}
For every \(C\ge1\), with \(r=1/(100C)\),
\[
 C\Lambda(r)\le\frac{47}{138}<1,
 \qquad
 B(4r)\le\frac{71}{27600C^2}.                             \tag{97.14}
\]
\end{lemma}

\begin{proof}
Here \(S=1/(25C)\le1/25\).  Therefore
\[
 \frac1{1-S}\le\frac{25}{24},\qquad
 \frac1{1-2S}\le\frac{25}{23}.                           \tag{97.15}
\]
Substitution in (97.12) yields
\[
 C\Lambda(r)
 \le C\,4S\left(\frac{25}{24}+\frac{25}{23}\right)
 =\frac{47}{138}.                                          \tag{97.16}
\]
Similarly, (97.11) and (97.15) give
\[
 B(S)\le S^2\left(\frac{25}{24}
            +\frac12\frac{25}{24}\frac{25}{23}\right)
 =\frac{71}{27600C^2}.                                     \tag{97.17}
\]
\end{proof}

\begin{theorem}[General rectangular fixed-point chart]
\label{thm:v18v87-rectangular-chart}
Suppose a compatible skew-Hermitian principal face-log cochain on \(Q_L\)
satisfies
\[
 \|F_L\|_\infty\le\varepsilon_L.                         \tag{97.18}
\]
Then the V77--V78 map formed with \(B_L^+\) is a contraction on the radius-
\(r_L\) coexact ball.  It has a unique fixed point and
\[
 c_L\le\frac{47}{138},                                    \tag{97.19}
\]
\[
 \|A_L\|_\infty
 \le\frac{493}{138000C_L}<r_L.                            \tag{97.20}
\]
For one mass-one unitary representation in four dimensions, the associated
finite-box Wilson estimate obeys
\[
 \operatorname{gap}H_{1,L}
 \ge1-\frac{3944}{138000C_L}
 \ge\frac{16757}{17250}.                                  \tag{97.21}
\]
\end{theorem}

\begin{proof}
Use (97.9) and Lemma \ref{lem:v18v87-scaled-majorants}.  The contraction
constant is at most \(C_L\Lambda(r_L)\), proving (97.19).  The self-map image
radius is at most
\[
 C_L\left(\frac1{1000C_L^2}
          +\frac{71}{27600C_L^2}\right)
 =\frac{493}{138000C_L}<\frac1{100C_L}.                    \tag{97.22}
\]
The Banach argument from V78 gives existence and uniqueness.  The V74
perturbation estimate and \(C_L\ge1\) give (97.21).
\end{proof}

\subsection{Executable certificate and scaling boundary}

The verifier
\[
 \texttt{scripts/verify\_sm\_v87\_general\_box\_orientation\_green.py} \tag{97.23}
\]
checks the orientation block, zero cross-orientation coupling, Green row mass,
rowwise total-variation identity, and bound in the boxes
((1,1,1,1)), ((2,1,1,1)), and ((2,2,1,1)).  It emits
\[
 \texttt{artifacts/sm\_v87\_general\_box\_orientation\_green.json}.   \tag{97.24}
\]
Their hashes are
\[
\begin{array}{c|l}
\hbox{object}&\hbox{SHA--256}\\ \hline
\hbox{verifier}&
\texttt{2FE39EC6103BD767764BCFE595D7475839FC79C0011E5EE66739FA93F9EC58A4}\\
\hbox{canonical payload}&
\texttt{1F19EF18205DA10956F141B9BDE8E7D02BBF03E5F8CEE1089CFBB965B6163C24}.
\end{array}                                                   \tag{97.25}
\]

\begin{firewall}[A mathematically cofinal but physically excessive scaling]
\label{fire:v18v87}
For isotropic side length \(N\), the certified stock is \(O(N^2)\) and
(97.13) requires \(\|F_N\|_\infty=O(N^{-4})\).  A smooth continuum curvature
sampled at mesh \(a\asymp N^{-1}\) naturally has plaquette logarithm of order
\(a^2=O(N^{-2})\), so the present sufficient condition is two powers too
strong for that route.  The gap estimate also remains a represented finite-
box perturbation bound, not an interacting Standard Model--Einstein continuum
mass theorem.
\end{firewall}

\begin{decision}[Use the compatible image before sharpening Moore--Penrose]
\label{dec:v18v87-next}
V88 should construct a cubical contracting homotopy on closed face cochains.
A path-integration right inverse on the compatible image is expected to have
\(O(L_{\max})\), rather than \(O(L_{\max}^2)\), infinity stock.  With the
scaling argument above, that would reduce the required curvature decay to the
natural \(O(L_{\max}^{-2})\) level.  Reflection symmetry can then be restored
by averaging the right inverse over box symmetries.
\end{decision}

\begin{assessment}[V87 acquisition]
\label{ass:v18v87}
V87 proves a full arbitrary-rectangle Hodge bound and an explicit cofinal
nonlinear chart.  Its \(N^{-4}\) curvature gate is rigorous but unsuitable for
the intended smooth \(a^2\) continuum scaling, so the next step moves from the
ambient pseudoinverse to a compatible-image contracting homotopy.
\end{assessment}

\begin{record}[V87 append record]
\label{rec:v18v87}
V87 is appended to the validated V86 whose installed SHA--256 was
\texttt{0ABBD654D47D2EABFE1DBB3F67B2531F32704701D81392D5E89EA8C602B288DB}.
After installation only V87 is the active numeric SM note; the frozen
\texttt{V100\_f2} through \texttt{V100\_f17} archives are unchanged.
\end{record}


\section{V88: reflection-averaged cubical homotopy and the nonlinear Bianchi defect}
\label{sec:v18v88}

V87 obtained an ambient Moore--Penrose bound of quadratic size and therefore
required face logarithms of order \(L_{\max}^{-4}\).  This version instead
uses the compatible curvature image.  A cubical contracting homotopy gives a
linear-size right inverse on closed face cochains, and averaging it over all
coordinate reflections preserves the right-inverse identity while halving the
axial infinity stock.  The same homotopy also exposes a compatibility term
that must vanish before any nonlinear fixed point can be called an exact
prescribed-curvature reconstruction.

\subsection{The tensor contraction}

On the interval \(I_m=[0,m]\), let \(P_m\) send a vertex cochain to the
constant cochain equal to its value at vertex zero and send every positive-
degree cochain to zero.  Define
\[
 (h_m f)(x)=\sum_{i=0}^{x-1}f(i),\qquad f\in C^1(I_m),     \tag{98.1}
\]
and let \(h_m\) vanish on \(C^0(I_m)\).  Directly,
\[
 d h_m+h_m d=I-P_m.                                        \tag{98.2}
\]
Relabel the four coordinates so their lengths satisfy
\[
 1\le\ell_0\le\ell_1\le\ell_2\le\ell_3=L_{\max}.        \tag{98.3}
\]
On their graded tensor product define
\[
 H=\sum_{k=0}^3
 P_{\ell_0}\otimes\cdots\otimes P_{\ell_{k-1}}
 \otimes h_{\ell_k}\otimes I\otimes\cdots\otimes I.     \tag{98.4}
\]
Write \(H_p:C^p(Q_L)\to C^{p-1}(Q_L)\) for its degree-\(p\) component.

\begin{theorem}[Cubical cochain contraction]
\label{thm:v18v88-contraction}
For every positive degree \(p\le4\),
\[
 d_{p-1}H_p+H_{p+1}d_p=I_{C^p(Q_L)},                       \tag{98.5}
\]
where \(H_5d_4\) is zero.  In particular, if a face cochain \(F\) is closed,
then
\[
 d_1H_2F=F.                                                 \tag{98.6}
\]
Moreover,
\[
 \|H_2\|_{\infty\to\infty}
 \le\ell_0+\ell_1+\ell_2.                                 \tag{98.7}
\]
\end{theorem}

\begin{proof}
Equation (98.5) is the tensor contraction identity obtained by inserting
(98.2) successively in the four factors; adjacent graded cross terms cancel.
It can also be seen cellwise.  An input cell is integrated in its first
oriented coordinate, while every preceding vertex coordinate is evaluated at
zero.  Applying \(dH+Hd\) telescopes that prefix sum and returns the original
positive-degree cochain; the residual tensor projection is zero in positive
degree.  Equation (98.6) follows from (98.5) at degree two and \(d_2F=0\).

For an output edge in coordinate \(\nu\), the only contributing input faces
have orientations \((k,\nu)\) with \(k<\nu\).  The \(k\)-th prefix sum has at
most \(\ell_k\) unit coefficients.  Hence that row sum is at most
\(\sum_{k<\nu}\ell_k\), whose maximum is the right side of (98.7).
\end{proof}

\subsection{Reflection averaging}

Let \(\mathscr R\cong(\mathbb Z/2)^4\) be the group generated by the four
coordinate reflections, acting by signed permutation isometries \(R_p(g)\) on
\(p\)-cochains.  Define
\[
 \widehat H_p
 =\frac1{16}\sum_{g\in\mathscr R}
 R_{p-1}(g)H_pR_p(g)^{-1}.                                 \tag{98.8}
\]

\begin{theorem}[Reflection-equivariant compatible right inverse]
\label{thm:v18v88-reflection-average}
The averaged family satisfies
\[
 d_{p-1}\widehat H_p+\widehat H_{p+1}d_p=I               \tag{98.9}
\]
in every positive degree and
\[
 R_{p-1}(g)\widehat H_p=\widehat H_pR_p(g)                \tag{98.10}
\]
for all \(g\in\mathscr R\).  Its face-to-edge and three-to-two stocks obey
\[
\boxed{
 \|\widehat H_2\|_{\infty\to\infty}
 \le C_L^{(2)}:=\frac{\ell_0+\ell_1+\ell_2}{2},}          \tag{98.11}
\]
\[
 \|\widehat H_3\|_{\infty\to\infty}
 \le C_L^{(3)}:=\frac{\ell_0+\ell_1}{2}.                 \tag{98.12}
\]
Consequently \(d_1\widehat H_2F=F\) for every closed face cochain.
\end{theorem}

\begin{proof}
Cellular coboundaries commute with every \(R_p(g)\).  Conjugating (98.5) and
averaging proves (98.9).  Left multiplication of the group average by a fixed
group element merely permutes its summands, proving (98.10).

The reflection average factorises coordinate by coordinate in each summand of
(98.4).  A basepoint projection becomes the average of evaluation at the two
endpoints and has infinity norm one.  The prefix operator becomes
\[
 (\widehat h_m f)(x)
 =\frac12\left(\sum_{i=0}^{x-1}f(i)
              -\sum_{i=x}^{m-1}f(i)\right),               \tag{98.13}
\]
up to the fixed orientation convention.  Every row of \(\widehat h_m\) has
\(m\) coefficients of absolute value \(1/2\), so its norm is \(m/2\).
Repeating the row argument from (98.7) proves (98.11).  For an output
2-cell, only contracted coordinates preceding its first oriented coordinate
contribute; at most the first two lengths occur.  This proves (98.12).  The
last assertion is (98.9) at degree two for \(d_2F=0\).
\end{proof}

\begin{corollary}[Natural compatible-image scaling]
\label{cor:v18v88-natural-scaling}
For an isotropic box of side length \(N\),
\[
 C_N^{(2)}=\frac{3N}{2},\qquad C_N^{(3)}=N.                \tag{98.14}
\]
Thus the reflection-equivariant compatible right inverse has linear rather
than quadratic stock.
\end{corollary}

\subsection{The nonlinear compatibility identity}

Let
\[
 \mathcal F(A)=d_1A+\mathcal N(A)                          \tag{98.15}
\]
be the principal plaquette-log map in the V77 chart, and define the homotopy
fixed-point map
\[
 \mathcal T_F(A)=\widehat H_2\bigl(F-\mathcal N(A)\bigr). \tag{98.16}
\]

\begin{proposition}[Exact nonlinear Bianchi defect]
\label{prop:v18v88-defect}
Every fixed point \(A=\mathcal T_F(A)\) satisfies the identity
\[
\boxed{
 \mathcal F(A)-F
 =-\widehat H_3d_2\bigl(F-\mathcal N(A)\bigr).}            \tag{98.17}
\]
In particular, it reconstructs the prescribed curvature exactly if
\[
 d_2\bigl(F-\mathcal N(A)\bigr)=0.                         \tag{98.18}
\]
If \(d_2F=0\), then
\[
 \mathcal F(A)-F=\widehat H_3d_2\mathcal N(A).            \tag{98.19}
\]
\end{proposition}

\begin{proof}
Apply (98.9) in degree two to \(F-\mathcal N(A)\):
\[
 d_1\widehat H_2(F-\mathcal N(A))
 =F-\mathcal N(A)-\widehat H_3d_2(F-\mathcal N(A)).       \tag{98.20}
\]
At a fixed point the left side is \(d_1A\).  Add \(\mathcal N(A)-F\) to both
sides to obtain (98.17).  Conditions (98.18) and \(d_2F=0\) give the two
specialisations.
\end{proof}

\begin{assessment}[Correction of the finite four-dimensional chart language]
\label{ass:v18v88-correction}
The Banach constants in V79, V81--V83, V86, and V87 rigorously control their
stated finite maps.  In a four-dimensional complex with 3-cells, a fixed point
of a map built from a pseudoinverse or homotopy is an exact prescribed-
curvature solution only when its nonlinear source \(F-\mathcal N(A)\) lies in
\(\operatorname{im}d_1\), equivalently when the appropriate compatibility
projection vanishes; (98.18) is sufficient on the contractible boxes here.
The V76--V78 two-dimensional disk has no 3-cell condition, so this additional
gate is absent there.  Earlier four-dimensional ``closure'' wording must be
read with this compatibility condition.
\end{assessment}

\subsection{Scaled projected chart and quantified residual}

Set
\[
 r_L=\frac1{100C_L^{(2)}},\qquad
 \varepsilon_L=\frac1{1000(C_L^{(2)})^2}.                 \tag{98.21}
\]
Since integer side lengths give \(C_L^{(2)}\ge3/2\), Lemma
\ref{lem:v18v87-scaled-majorants} applies.

\begin{theorem}[Compatible-scale homotopy fixed point]
\label{thm:v18v88-scaled-map}
If \(\|F\|_\infty\le\varepsilon_L\), the map (98.16) is a contraction of the
radius-\(r_L\) ball into itself, with
\[
 c_L\le\frac{47}{138},\qquad
 \|A\|_\infty\le\frac{493}{138000C_L^{(2)}}<r_L.          \tag{98.22}
\]
It therefore has a unique fixed point in that ball.  If \(d_2F=0\), its
curvature mismatch obeys
\[
\boxed{
 \|\mathcal F(A)-F\|_\infty
 \le\frac{71}{4600}
       \frac{C_L^{(3)}}{(C_L^{(2)})^2}.}                   \tag{98.23}
\]
If the compatibility gate (98.18) holds, the mismatch is zero.
\end{theorem}

\begin{proof}
Use the norm (98.11) in the proof of Theorem
\ref{thm:v18v87-rectangular-chart}.  The scaled majorants give the contraction
and self-map bounds (98.22), hence Banach existence and uniqueness.  When
\(d_2F=0\), Proposition \ref{prop:v18v88-defect}, (98.12), and the fact that
each oriented 3-cell boundary has six faces give
\[
 \|\mathcal F(A)-F\|_\infty
 \le6C_L^{(3)}B(4r_L)
 \le6C_L^{(3)}\frac{71}{27600(C_L^{(2)})^2},              \tag{98.24}
\]
which reduces to (98.23).  Equation (98.17) proves the final assertion.
\end{proof}

\begin{corollary}[Isotropic rates]
\label{cor:v18v88-isotropic-rates}
For side length \(N\), the choices and defect bound become
\[
 r_N=\frac1{150N},\qquad
 \varepsilon_N=\frac1{2250N^2},                           \tag{98.25}
\]
\[
 \|A_N\|_\infty\le\frac{493}{207000N},qquad
 \|\mathcal F(A_N)-F_N\|_\infty\le\frac{71}{10350N}.    \tag{98.26}
\]
The face-log hypothesis now has the natural \(N^{-2}\) order.  Exact equality
still requires (98.18); the last bound only proves that the displayed
homotopy residual tends to zero.
\end{corollary}

\begin{proof}
Insert (98.14) into (98.21)--(98.23) and reduce the fractions.
\end{proof}

\subsection{Executable certificate}

The verifier
\[
 \texttt{scripts/verify\_sm\_v88\_cubical\_contracting\_homotopy.py}  \tag{98.27}
\]
constructs all cochain groups of the box \((1,1,2,2)\), verifies (98.5) in
degrees one through four, averages every \(H_p\) over all 16 reflections,
verifies (98.9)--(98.10), and checks the degree-two defect identity.  In this
box it obtains
\[
 \|H_2\|_\infty=4,\quad\|\widehat H_2\|_\infty=2,\quad
 \|H_3\|_\infty=2,\quad\|\widehat H_3\|_\infty=1.        \tag{98.28}
\]
It emits
\[
 \texttt{artifacts/sm\_v88\_cubical\_contracting\_homotopy.json}.    \tag{98.29}
\]
The hashes are
\[
\begin{array}{c|l}
\hbox{object}&\hbox{SHA--256}\\ \hline
\hbox{verifier}&
\texttt{6930EAA0735958D728937C58EF807205C009C1657825FD243BFD7DB6A6B9D3A0}\\
\hbox{canonical payload}&
\texttt{83430F50413E8890B86B2CCDD42CDC52768A713B1FA656E75BFDC8C131CF5177}.
\end{array}                                                   \tag{98.30}
\]

\begin{firewall}[The exact nonlinear Bianchi gate remains open]
\label{fire:v18v88}
The reflection-averaged homotopy proves an exact right inverse on linearly
closed face cochains and a unique small fixed point of the projected map.  It
does not prove that this fixed point satisfies (98.18).  The vanishing residual
limit in (98.26) is not an exact finite-regulator curvature solution and does
not by itself define a consistent projective gauge measure.  Product gauge
factors, chiral fermions, Einstein coupling, reflection positivity, and
continuum reconstruction remain open.
\end{firewall}

\begin{decision}[Solve the compatibility equation rather than suppress it]
\label{dec:v18v88-next}
V89 should treat
\[
 \mathscr B_F(A):=d_2\bigl(F-\mathcal N(A)\bigr)           \tag{98.31}
\]
as an explicit second equation.  The first task is to derive its Lipschitz and
quadratic size bounds and determine whether a face correction in
\(\operatorname{ran}H_3\) can be coupled to (98.16) by a contraction or an
implicit-function argument.  No further Wilson-gap conclusion should be drawn
from the four-dimensional chart until this compatibility residual is closed.
\end{decision}

\begin{assessment}[V88 acquisition]
\label{ass:v18v88}
V88 obtains the natural \(N^{-2}\) compatible-image scaling and a reflection-
equivariant linear right inverse, while isolating the exact nonlinear Bianchi
defect.  The continuum programme now has a sharper and more honest bottleneck:
solve (98.31) jointly with the link reconstruction.
\end{assessment}

\begin{record}[V88 append record]
\label{rec:v18v88}
V88 is appended to the validated V87 whose installed SHA--256 was
\texttt{C40C119F40FF55A814F18220B6E8328FF1E023F5B9C863B6A40A0B79E9EC7741}.
After installation only V88 is the active numeric SM note; the frozen
\texttt{V100\_f2} through \texttt{V100\_f17} archives are unchanged.
\end{record}


\section{V89: the exact nonlinear curvature graph over compatible coordinates}
\label{sec:v18v89}

V88 isolated the residual that appears when one asks an arbitrary linearly
closed face cochain to be the full nonabelian plaquette logarithm.  The correct
local object is not the entire ambient face space.  This version uses the
reflection-averaged homotopy to construct complementary projections and proves
that the actual small-curvature image is an exact nonlinear graph over the
linear compatible space.  The graph coordinate is recovered by a fixed linear
projection, so no nonlinear Bianchi equation is discarded.

\subsection{Compatible and Bianchi projections}

Define on face cochains
\[
 P_L=d_1\widehat H_2,\qquad
 Q_L=I-P_L=\widehat H_3d_2.                                \tag{99.1}
\]
The second equality is the degree-two homotopy identity (98.9).  Let
\[
 \mathcal E_L=\operatorname{im}d_1=\ker d_2.               \tag{99.2}
\]
The equality holds because the rectangular box is contractible.

\begin{proposition}[Exact complementary projections]
\label{prop:v18v89-projections}
The operators in (99.1) satisfy
\[
 P_L^2=P_L,qquad Q_L^2=Q_L,qquad
 P_LQ_L=Q_LP_L=0,                                         \tag{99.3}
\]
\[
 \operatorname{ran}P_L=\mathcal E_L,qquad
 d_2P_L=0,qquad P_Ld_1=d_1.                              \tag{99.4}
\]
Both commute with every coordinate reflection.  Thus every face cochain has
the unique splitting
\[
 F=P_LF+Q_LF,qquad P_LF\in\mathcal E_L,\quad
 Q_LF\in\ker P_L.                                         \tag{99.5}
\]
\end{proposition}

\begin{proof}
The range of \(P_L=d_1\widehat H_2\) lies in \(\operatorname{im}d_1\).  If
\(F\in\mathcal E_L\), then \(d_2F=0\), and (98.9) gives \(P_LF=F\).
Therefore \(P_L\) is a projection onto \(\mathcal E_L\).  Standard algebra for
\(Q_L=I-P_L\) proves (99.3) and the splitting.  The last two relations in
(99.4) follow from \(d_2d_1=0\) and the fact that \(P_L\) is the identity on
\(\operatorname{im}d_1\).  Reflection equivariance follows from (98.10) and
cellular naturality of the coboundaries.
\end{proof}

\subsection{Small compatible coordinates}

Retain the constants
\[
 C_L^{(2)}=\frac{\ell_0+\ell_1+\ell_2}{2},qquad
 C_L^{(3)}=\frac{\ell_0+\ell_1}{2},                        \tag{99.6}
\]
with the ordered side lengths (98.3), and set
\[
 r_L=\frac1{100C_L^{(2)}},qquad
 \varepsilon_L=\frac1{1000(C_L^{(2)})^2}.                 \tag{99.7}
\]
For \(G\in\mathcal E_L\) with \(\|G\|_\infty\le\varepsilon_L\), define
\[
 \mathcal T_G(A)
 =\widehat H_2\bigl(G-\mathcal N(A)\bigr).                \tag{99.8}
\]

\begin{theorem}[Unique link chart over the compatible coordinate]
\label{thm:v18v89-link-chart}
For every admitted \(G\), the map (99.8) has a unique fixed point \(A_G\) in
the radius-\(r_L\) ball.  Uniformly in the box,
\[
 \operatorname{Lip}(\mathcal T_G)\le c_L\le\frac{47}{138}, \tag{99.9}
\]
\[
 \|A_G\|_\infty
 \le\frac{493}{138000C_L^{(2)}}<r_L.                      \tag{99.10}
\]
For two compatible coordinates in the admitted ball,
\[
 \|A_G-A_{G'}\|_\infty
 \le\frac{C_L^{(2)}}{1-c_L}\|G-G'\|_\infty
 \le\frac{138}{91}C_L^{(2)}\|G-G'\|_\infty.             \tag{99.11}
\]
\end{theorem}

\begin{proof}
The self-map and contraction estimates are exactly those in Theorem
\ref{thm:v18v88-scaled-map}, with \(F\) there replaced by \(G\).  Banach's
theorem gives the unique fixed point and (99.9)--(99.10).  Subtract the two
fixed-point equations and use (98.11):
\[
 \|A_G-A_{G'}\|_\infty
 \le C_L^{(2)}\|G-G'\|_\infty+c_L\|A_G-A_{G'}\|_\infty.  \tag{99.12}
\]
Rearrangement and \(1-c_L\ge91/138\) give (99.11).
\end{proof}

\subsection{The nonlinear image as a graph}

Define
\[
 \Gamma_L(G)=Q_L\mathcal N(A_G),qquad
 \Psi_L(G)=G+\Gamma_L(G).                                 \tag{99.13}
\]

\begin{theorem}[Exact nonlinear curvature graph]
\label{thm:v18v89-curvature-graph}
The full principal plaquette logarithm of the fixed-point link field is
exactly
\[
\boxed{
 \mathcal F(A_G)=\Psi_L(G)=G+Q_L\mathcal N(A_G).}          \tag{99.14}
\]
Moreover,
\[
 P_L\Psi_L(G)=G,qquad P_L\Gamma_L(G)=0.                  \tag{99.15}
\]
Consequently \(\Psi_L\) is injective, and \(P_L\) is its inverse coordinate
on the graph.  An admitted small face cochain \(F\) belongs to this constructed
nonlinear curvature image exactly when
\[
 F=\Psi_L(P_LF)                                            \tag{99.16}
\]
and \(P_LF\) lies in the compatible coordinate ball.
\end{theorem}

\begin{proof}
At the fixed point,
\[
 A_G=\widehat H_2(G-\mathcal N(A_G)).                      \tag{99.17}
\]
Apply \(d_1\) and use \(P_LG=G\):
\[
 d_1A_G=G-P_L\mathcal N(A_G).                             \tag{99.18}
\]
Adding \(\mathcal N(A_G)\) gives (99.14).  Since
\(P_LQ_L=0\), (99.15) follows.  If \(\Psi_L(G)=\Psi_L(G')\), applying \(P_L\)
gives \(G=G'\), proving injectivity.  Finally, every graph value obeys (99.16)
by (99.15), while (99.16) itself supplies \(F=\Psi_L(G)\) with \(G=P_LF\),
proving the equivalence.
\end{proof}

\begin{corollary}[Nonlinear Bianchi condition is built into the graph]
\label{cor:v18v89-bianchi-graph}
For every admitted \(G\), the pair \((A_G,\Psi_L(G))\) is an exact solution of
\[
 \mathcal F(A)=F.                                          \tag{99.19}
\]
The compatibility residual from V88 is precisely the graph height:
\[
 \Psi_L(G)-G=Q_L\mathcal N(A_G)
 =\widehat H_3d_2\mathcal N(A_G).                          \tag{99.20}
\]
Thus one need not impose \(d_2(G-\mathcal N(A_G))=0\) when the full face log is
defined by (99.14); that condition is required only if one insists that the
full face log equal the linear coordinate \(G\) itself.
\end{corollary}

\begin{proof}
Equation (99.19) is (99.14) with \(F=\Psi_L(G)\).  Equation (99.20) follows
from (99.1) and (99.13).  The last assertion is the defect identity (98.17).
\end{proof}

\subsection{Graph height, stability, and symmetry}

\begin{theorem}[Quantitative graph bounds]
\label{thm:v18v89-graph-bounds}
The graph height satisfies
\[
 \|\Gamma_L(G)\|_\infty
 \le\frac{71}{4600}\frac{C_L^{(3)}}{(C_L^{(2)})^2}.       \tag{99.21}
\]
For two admitted coordinates,
\[
 \|\Gamma_L(G)-\Gamma_L(G')\|_\infty
 \le\frac{282}{91}C_L^{(3)}\|G-G'\|_\infty.             \tag{99.22}
\]
The maps \(G\mapsto A_G\), \(\Gamma_L\), and \(\Psi_L\) intertwine every
coordinate reflection.
\end{theorem}

\begin{proof}
Using \(Q_L=\widehat H_3d_2\), the six-face incidence bound, (98.12), and
(97.17) gives
\[
 \|\Gamma_L(G)\|_\infty
 \le6C_L^{(3)}B(4r_L)
 \le\frac{71}{4600}\frac{C_L^{(3)}}{(C_L^{(2)})^2},      \tag{99.23}
\]
which is (99.21).  Similarly, the V78 face-remainder Lipschitz bound gives
\[
 \|\Gamma_L(G)-\Gamma_L(G')\|_\infty
 \le6C_L^{(3)}\Lambda(r_L)\|A_G-A_{G'}\|_\infty.        \tag{99.24}
\]
Insert \(C_L^{(2)}\Lambda(r_L)\le47/138\) and (99.11) to obtain
\(6\cdot47/91=282/91\), proving (99.22).

Reflections commute with \(P_L,Q_L,\widehat H_2\) and with the natural
plaquette-log remainder inside the principal chart.  Conjugating the fixed-
point equation for \(G\) therefore gives the equation for the reflected
coordinate.  Uniqueness identifies the reflected fixed point, and (99.13)
then proves equivariance of the two graph maps.
\end{proof}

\begin{remark}[Isotropic size of the graph correction]
\label{rem:v18v89-isotropic}
For side length \(N\), (99.21) becomes
\[
 \|\Gamma_N(G)\|_\infty\le\frac{71}{10350N}.             \tag{99.25}
\]
This tends to zero, but the bound is only \(O(N^{-1})\), while the coordinate
ball has radius \(O(N^{-2})\).  The estimate therefore does not yet show that
the full graph has the smooth plaquette-log scale.  A structured bound on
\(Q_L\mathcal N(A_G)\), rather than the ambient \(\|Q_L\|\) bound, is needed.
\end{remark}

\subsection{Executable certificate}

The verifier
\[
 \texttt{scripts/verify\_sm\_v89\_nonlinear\_curvature\_graph.py}     \tag{99.26}
\]
constructs \(P_L,Q_L\) from the V88 averaged homotopy on the box
\((1,1,2,2)\), verifies the projection, rank, reflection, and exact-complex
identities, and checks (99.14)--(99.15) on a deterministic rational probe.  It
obtains
\[
 \operatorname{rank}P_L=49,qquad \operatorname{rank}Q_L=24,          \tag{99.27}
\]
and emits
\[
 \texttt{artifacts/sm\_v89\_nonlinear\_curvature\_graph.json}.      \tag{99.28}
\]
The hashes are
\[
\begin{array}{c|l}
\hbox{object}&\hbox{SHA--256}\\ \hline
\hbox{verifier}&
\texttt{5BA13BAD32C4CB0790944AF4A05D558F09C6DE38AF3188F90C253E4134C6FE07}\\
\hbox{canonical payload}&
\texttt{A2BB71FFBA395AC9BD222D26EB00D21F2F46BCAAB47CF5254175934949215860}.
\end{array}                                                   \tag{99.29}
\]

\begin{firewall}[A local graph is not a cofinal measure]
\label{fire:v18v89}
Theorem \ref{thm:v18v89-curvature-graph} resolves finite-chart compatibility
by parameterising the nonlinear image.  It does not provide a volume-uniform
Jacobian, a probability density on the graph, consistency under restriction of
boxes, reflection positivity of a pushed measure, or the product-gauge,
fermion, and Einstein sectors.  The current graph-height estimate is too weak
to identify the full curvature with a smooth \(N^{-2}\) field uniformly.
\end{firewall}

\begin{decision}[Checkpoint, then structured graph height]
\label{dec:v18v89-next}
V90 is the mandatory companion checkpoint.  The substantive V91 continuation
should estimate \(Q_L\mathcal N(A_G)\) on the used fixed-point image.  It should
exploit the local commutator structure and the discrete derivatives in
\(d_2\mathcal N(A_G)\), rather than charging six unrelated face remainders and
the full \(H_3\) path length.
\end{decision}

\begin{assessment}[V89 acquisition]
\label{ass:v18v89}
V89 replaces an unsolved arbitrary-curvature compatibility demand by an exact,
injective nonlinear curvature graph over the compatible linear coordinate.
The graph coordinate is recovered by \(P_L\), and exact link reconstruction is
proved throughout the small chart.  The next analytic bottleneck is the
cofinal height of the graph correction.
\end{assessment}

\begin{record}[V89 append record]
\label{rec:v18v89}
V89 is appended to the validated V88 whose installed SHA--256 was
\texttt{C7E326B3B61B11908F8535D7BF5CA9506954AC3CECE3771BFEB5762CA1CFAA53}.
After installation only V89 is the active numeric SM note; the frozen
\texttt{V100\_f2} through \texttt{V100\_f17} archives are unchanged.
\end{record}


\section{V90: mandatory companion audit and composite-first direction}
\label{sec:v18v90}

V89 identified the exact nonlinear curvature image as the graph
\(\Psi_L(G)=G+Q_L\mathcal N(A_G)\), but its available graph-height estimate
first bounded \(d_2\mathcal N\) by six unrelated face terms and then charged
the full norm of \(\widehat H_3\).  The standing five-version rule now
requires an audit of the active Yang--Mills and Navier--Stokes notes.  This
section records the sources, separates transferable methods from
problem-specific claims, and fixes the next SM acquisition.

\subsection{Audited sources}

The active companion files at the time of the audit were
\[
\begin{array}{c|l}
\hbox{programme and active file}&\hbox{SHA--256}\\ \hline
\hbox{Yang--Mills,
 \texttt{Renewal\_Yang\_Mills\_Mass\_Gap\_Research\_Notes\_V21.tex}}&
\texttt{B020585A4AB1BD97B9EBF22C4AFE801D65B040D1DAFA71802D2109996968EFF6}\\
\hbox{Navier--Stokes,
 \texttt{Renewal\_NS\_Stability\_Research\_Notes\_V26.tex}}&
\texttt{82C887F8C2C69E4F28FAD7C3DC8DE5B9099CB5A4BF431EBA1FFF3E91935978AD}.
\end{array}                                                   \tag{100.1}
\]
The Navier--Stokes file is unchanged from the V85 audit.  The Yang--Mills
file has advanced from V17 to V21, so its V18--V21 conclusions are new input
to this checkpoint.

\subsection{What transfers and what does not}

\begin{proposition}[Composite-first transfer from Yang--Mills V21]
\label{prop:v18v90-ym-transfer}
The following methodological inference is valid for the present graph-height
problem.

Suppose a target map is a composition \(TS\).  Sparsity or locality of \(S\)
alone gives no useful size conclusion when \(T\) is long-range or dense.
Before multiplying the separate bounds for \(T\) and \(S\), one should
derive an exact formula for \(TS\) and test whether algebraic cancellation or
restricted input structure survives in the composition.

Applied here, the object to estimate is
\[
 \Gamma_L(G)=Q_L\mathcal N(A_G)
 =\widehat H_3d_2\mathcal N(A_G),                           \tag{100.2}
\]
with \(A_G\) constrained by the fixed-point equation.  Thus V91 should study
the composite in (100.2), rather than the three factors
\(\widehat H_3\), \(d_2\), and \(\mathcal N\) independently.
\end{proposition}

\begin{proof}
Yang--Mills V21 gives an exact example of the premise.  Its raw coefficient
banks are sparse, but its bridge \(K=ZA\) has full rank \(45\) and \(1937\)
nonzero entries out of \(2025\); a common coordinate-block reduction is
therefore absent.  Nevertheless an exact push-through identity
\[
 {\cal S}(I-ACZ)A=A{\cal S}(I-CZA)                           \tag{100.3}
\]
moves the computation to the raw side.  A complete declared-operation ledger,
rather than sparsity in isolation, then selects that architecture.  The proof
of the transfer is elementary: a separate estimate gives
\(\|TSx\|\le\|T\|\,\|Sx\|\), whereas an exact identity for \(TSx\) may expose
cancellation already destroyed by the two absolute values.  Substitution of
\(T=\widehat H_3\), \(S=d_2\mathcal N\), and \(x=A_G\) gives (100.2).
No Yang--Mills spectral or mass-gap assertion is imported.
\end{proof}

\begin{proposition}[Restricted-input transfer from Navier--Stokes V26]
\label{prop:v18v90-ns-transfer}
The fixed-point range
\[
 {\cal A}_L
 =\left\{\widehat H_2\bigl(G-\mathcal N(A_G)\bigr):
 G\in\mathcal E_L,\ \|G\|_\infty\le\varepsilon_L\right\}    \tag{100.4}
\]
is a strict structured subset of the ambient link-cochain ball.  Therefore a
uniform estimate for \(d_2\mathcal N(A)\) on all link cochains may be
strictly weaker than an estimate on \(A\in{\cal A}_L\).  Any claimed gain
must be proved on this actual range and then propagated through
\(\widehat H_3\).
\end{proposition}

\begin{proof}
The inclusion in the ambient radius-\(r_L\) ball is Theorem
\ref{thm:v18v89-link-chart}.  Navier--Stokes V26 supplies a methodological
analogue: contracting the kernel derivatives only against the rank-one
inputs actually present in its nonlinear term produces an exact
one-variable optimization, but the improvement is small and even vanishes
for the first derivative.  Hence restriction to the used input class is
legitimate, while an improvement cannot be assumed.  Formula (100.4)
identifies the class that must be used here.  No Navier--Stokes regularity or
kernel inequality is imported.
\end{proof}

\begin{corollary}[Audited change of direction]
\label{cor:v18v90-direction}
A V91 estimate is admissible only after one of the following has been proved:
\begin{enumerate}
\item an exact local expansion of \(d_2\mathcal N(A)\) whose cancellations are
      retained before absolute values;
\item a kernel formula for \(\widehat H_3d_2\mathcal N(A)\) on
      \(A\in{\cal A}_L\);
\item a counterexample showing that the desired \(O(N^{-2})\) graph height
      fails under the current scaling.
\end{enumerate}
Merely sharpening the constants in
\(\|\widehat H_3\|\,\|d_2\|\,\|\mathcal N(A)\|\) cannot close the scale gap.
\end{corollary}

\begin{proof}
The existing factorwise estimate is \(O(N^{-1})\) by (99.25), while the
coordinate ball is \(O(N^{-2})\).  A constant improvement leaves the power of
\(N\) unchanged.  Each listed route either proves the missing cancellation,
evaluates the relevant composite on its used range, or disproves the intended
scale.  These exhaust the honest immediate outcomes of a composite-first
test.
\end{proof}

\begin{firewall}[Companion results are methodological evidence]
\label{fire:v18v90}
The audit proves no Standard Model field construction, Einstein coupling,
Yang--Mills mass gap, or Navier--Stokes regularity theorem.  It neither
transfers a companion theorem across subjects nor treats exact arithmetic
certificates as analytic proofs.  It changes the order of the next
calculation and records why factorwise norm sharpening is insufficient.
\end{firewall}

\begin{decision}[V91 acquisition]
\label{dec:v18v90-next}
V91 will work upstream at the plaquette logarithm.  It will derive the
homogeneous quadratic BCH term of \(\mathcal N(A)\), apply \(d_2\) before
taking norms, and determine exactly which quadratic terms cancel on a
cubical boundary.  Only after that identity is established will a
volume-scaling estimate be attempted.
\end{decision}

\begin{assessment}[V90 acquisition]
\label{ass:v18v90}
The mandatory companion audit is complete.  Its actionable result is the
composite-first rule (100.2) and the restriction to the actual fixed-point
range (100.4).  The \(O(N^{-1})\) graph-height bottleneck remains open.
\end{assessment}

\begin{record}[V90 append record]
\label{rec:v18v90}
V90 is appended to the validated V89 whose installed SHA--256 was
\texttt{870CBC2DA65FFF0A5BC0BB7F3E5C43AC11C7114B82F4E5E1FB28E74DE7A6D464}.
After installation only V90 is the active numeric SM note; the frozen
\texttt{V100\_f2} through \texttt{V100\_f17} archives and all unrelated
notes are unchanged.
\end{record}


\section{V91: exact lattice Bianchi cancellation and cofinal graph height}
\label{sec:v18v91}

The V89 estimate treated the six face remainders on a cube independently.
V90 required the composite to be simplified first.  This section does so.
An ordered product of six transported plaquette holonomies freely reduces to
the identity.  Its logarithm converts the apparent \(A^2\) cube source into
an \(A\,F+F^2\) source, where \(F\) is the full face logarithm.  On the V89
fixed-point graph this closes a continuity barrier and proves the missing
\(O(N^{-2})\) isotropic graph-height bound.

\subsection{The exact transported cube word}

For three coordinate directions \(\mu<\nu<\rho\), write
\[
 U_\alpha(x)=e^{A_\alpha(x)}
\]
and define the positively oriented plaquette holonomy
\[
 P_{\alpha\beta}(x)
 =U_\alpha(x)U_\beta(x+e_\alpha)
  U_\alpha(x+e_\beta)^{-1}U_\beta(x)^{-1},
 \qquad \alpha<\beta.                                     \tag{101.1}
\]
Let \(F_{\alpha\beta}(x)=\log P_{\alpha\beta}(x)\) in the principal chart.
The cubical coboundary convention is
\[
\begin{split}
 (d_2F)_{\mu\nu\rho}(x)
 ={}&F_{\nu\rho}(x+e_\mu)-F_{\nu\rho}(x)\\
 &-F_{\mu\rho}(x+e_\nu)+F_{\mu\rho}(x)\\
 &+F_{\mu\nu}(x+e_\rho)-F_{\mu\nu}(x).                    \tag{101.2}
\end{split}
\]

\begin{lemma}[Ordered nonabelian cube identity]
\label{lem:v18v91-cube-word}
With \(\operatorname{Ad}_U(V)=UVU^{-1}\), the six based face holonomies
satisfy
\[
\begin{split}
&I=\operatorname{Ad}_{U_\mu(x)}P_{\nu\rho}(x+e_\mu)\,
 P_{\mu\rho}(x)\,
 \operatorname{Ad}_{U_\rho(x)}P_{\mu\nu}(x+e_\rho)\\
&\hspace{12mm}\cdot P_{\nu\rho}(x)^{-1}\,
 \operatorname{Ad}_{U_\nu(x)}P_{\mu\rho}(x+e_\nu)^{-1}\,
 P_{\mu\nu}(x)^{-1}.                                      \tag{101.3}
\end{split}
\]
\end{lemma}

\begin{proof}
Put \(a=U_\mu(x)\), \(b=U_\nu(x)\), and \(c=U_\rho(x)\), and append
subscripts to denote translation in the other coordinate directions.  The
six factors of (101.3) expand in the free edge group as
\[
\begin{array}{lll}
 a\,b_\mu c_{\mu\nu}b_{\mu\rho}^{-1}c_\mu^{-1}a^{-1},
&
 a\,c_\mu a_\rho^{-1}c^{-1},
&
 c\,a_\rho b_{\mu\rho}a_{\nu\rho}^{-1}b_\rho^{-1}c^{-1},
\\
 c\,b_\rho c_\nu^{-1}b^{-1},
&
 b\,c_\nu a_{\nu\rho}c_{\mu\nu}^{-1}a_\nu^{-1}b^{-1},
&
 b\,a_\nu b_\mu^{-1}a^{-1}.
\end{array}                                                \tag{101.4}
\]
Multiplication in the displayed order cancels adjacent inverse edge pairs:
after the first three factors the reduced word is
\(a b_\mu c_{\mu\nu}a_{\nu\rho}^{-1}b_\rho^{-1}c^{-1}\);
after the fourth it is
\(a b_\mu c_{\mu\nu}a_{\nu\rho}^{-1}c_\nu^{-1}b^{-1}\);
the fifth reduces it to \(a b_\mu a_\nu^{-1}b^{-1}\), and the sixth
reduces it to the empty word.  This proves (101.3) without commutativity.
\end{proof}

Define six Lie-algebra elements in the same order:
\[
\begin{array}{lll}
Y_1=\operatorname{Ad}_{U_\mu(x)}F_{\nu\rho}(x+e_\mu),
&
Y_2=F_{\mu\rho}(x),
&
Y_3=\operatorname{Ad}_{U_\rho(x)}F_{\mu\nu}(x+e_\rho),
\\
Y_4=-F_{\nu\rho}(x),
&
Y_5=-\operatorname{Ad}_{U_\nu(x)}F_{\mu\rho}(x+e_\nu),
&
Y_6=-F_{\mu\nu}(x).
\end{array}                                                \tag{101.5}
\]
Then (101.3) is \(\prod_{j=1}^6e^{Y_j}=I\).

\begin{lemma}[Six-factor BCH remainder]
\label{lem:v18v91-six-bch}
Let \(S=\sum_{j=1}^6\|Y_j\|_{\rm op}<1/2\), and put
\[
 {\cal R}_6(Y)
 =\log\left(\prod_{j=1}^6e^{Y_j}\right)-\sum_{j=1}^6Y_j.
                                                                  \tag{101.6}
\]
Then
\[
 \|{\cal R}_6(Y)\|_{\rm op}\le B(S)
 =\frac{S^2}{1-S}+\frac{S^2}{2(1-S)(1-2S)},                \tag{101.7}
\]
and its homogeneous quadratic term is
\[
 {\cal R}_6^{(2)}(Y)
 =\frac12\sum_{1\le i<j\le6}[Y_i,Y_j].                    \tag{101.8}
\]
\end{lemma}

\begin{proof}
The proof of Theorem \ref{thm:v18v77-face-remainder} uses only the sum of
the factor norms, not the number four.  Applying its product and logarithm
series argument to six factors gives (101.7).  For (101.8), multiply
\(e^{Y_j}=I+Y_j+\frac12Y_j^2+O_3\).  If the product is \(I+Z\), then
\[
 Z=\sum_jY_j+\frac12\sum_jY_j^2+\sum_{i<j}Y_iY_j+O_3.
\]
Substitution in \(\log(I+Z)=Z-\frac12Z^2+O_3\) leaves
\(\frac12\sum_{i<j}(Y_iY_j-Y_jY_i)\) beyond the linear term.
\end{proof}

\subsection{Covariant rather than facewise cancellation}

\begin{theorem}[Exact lattice Bianchi estimate]
\label{thm:v18v91-bianchi-estimate}
Suppose every link logarithm on the cube has norm at most \(r\), every face
logarithm has norm at most \(f\), and \(6f<1/2\).  Then
\[
\boxed{
 \|d_2F\|_\infty
 \le 3\bigl(e^{2r}-1\bigr)f+B(6f).}                        \tag{101.9}
\]
More precisely, on each oriented cube,
\[
 d_2F=-{\cal C}(A,F)-{\cal R}_6(Y),                        \tag{101.10}
\]
where
\[
\begin{split}
 {\cal C}(A,F)={}&
 (\operatorname{Ad}_{U_\mu(x)}-I)F_{\nu\rho}(x+e_\mu)\\
 &-(\operatorname{Ad}_{U_\nu(x)}-I)F_{\mu\rho}(x+e_\nu)\\
 &+(\operatorname{Ad}_{U_\rho(x)}-I)F_{\mu\nu}(x+e_\rho).
                                                                  \tag{101.11}
\end{split}
\]
\end{theorem}

\begin{proof}
Since the product in (101.3) is the identity, (101.6) gives
\(\sum_jY_j=-{\cal R}_6(Y)\).  Comparing (101.5) with (101.2) gives
\[
 \sum_jY_j=d_2F+{\cal C}(A,F),
\]
which is (101.10).  Conjugation preserves the operator norm, so
\(S\le6f\).  Moreover
\[
 \|(\operatorname{Ad}_{e^A}-I)X\|_{\rm op}
 \le\sum_{k\ge1}\frac{\|\operatorname{ad}_A\|^k}{k!}\|X\|_{\rm op}
 \le(e^{2\|A\|_{\rm op}}-1)\|X\|_{\rm op},                 \tag{101.12}
\]
because \(\|\operatorname{ad}_A\|\le2\|A\|_{\rm op}\).  Apply (101.7)
to (101.10), use (101.12) on its three transport terms, and take the
maximum over cubes.
\end{proof}

\begin{corollary}[The explicit quadratic cancellation]
\label{cor:v18v91-quadratic}
Let \(F^{(1)}=d_1A\), and let \(\mathcal N^{(2)}(A)\) be the homogeneous
quadratic part of the plaquette-log remainder.  Write \(y_j\) for the six
signed face values in (101.5) with every transport removed and \(F\)
replaced by \(F^{(1)}\).  Then
\[
\begin{split}
 d_2\mathcal N^{(2)}(A)
={}&-[A_\mu(x),F^{(1)}_{\nu\rho}(x+e_\mu)]
   +[A_\nu(x),F^{(1)}_{\mu\rho}(x+e_\nu)]\\
 &-[A_\rho(x),F^{(1)}_{\mu\nu}(x+e_\rho)]
   -\frac12\sum_{i<j}[y_i,y_j].                            \tag{101.13}
\end{split}
\]
Thus the quadratic cube source contains one link logarithm times one linear
face curvature, or two linear face curvatures; it contains no uncancelled
sum of six arbitrary \(A^2\) face terms.
\end{corollary}

\begin{proof}
Since \(F=d_1A+\mathcal N(A)\) and \(d_2d_1=0\), one has
\(d_2F=d_2\mathcal N(A)\).  Take the homogeneous quadratic part of
(101.10).  The quadratic part of
\((\operatorname{Ad}_{e^A}-I)F\) is \([A,F^{(1)}]\), while (101.8) gives
the last term.  This is (101.13).
\end{proof}

\subsection{The isotropic continuity barrier}

\begin{theorem}[Cofinal graph-height scale]
\label{thm:v18v91-graph-height}
Let the four-dimensional box be isotropic of integer side length \(N\ge1\).
For
\[
 r_N=\frac1{150N},\qquad
 \varepsilon_N=\frac1{2250N^2},                            \tag{101.14}
\]
and every \(G\in\mathcal E_N\) with
\(\|G\|_\infty\le\varepsilon_N\), the V89 fixed point satisfies
\[
\boxed{
 \|\Gamma_N(G)\|_\infty<\frac{\varepsilon_N}{10}
 =\frac1{22500N^2}.}                                      \tag{101.15}
\]
Consequently its full plaquette logarithm obeys
\[
 \|\Psi_N(G)\|_\infty
 \le\frac{11}{10}\varepsilon_N
 =\frac{11}{24750N^2}.                                    \tag{101.16}
\]
\end{theorem}

\begin{proof}
For \(0\le t\le1\), let \(A_t=A_{tG}\),
\(F_t=\Psi_N(tG)\), and \(\Gamma_t=F_t-tG\).  The stability estimate
(99.11) and continuity of \(\mathcal N\) make \(t\mapsto\Gamma_t\)
continuous.  At \(t=0\), zero is a fixed point; uniqueness gives
\(A_0=0\) and \(\Gamma_0=0\).

Suppose, for contradiction, that there is a first \(t_*\) such that
\(\|\Gamma_{t_*}\|_\infty=\varepsilon_N/10\).  Put
\(f=\|F_{t_*}\|_\infty\).  Then
\[
 f\le\frac{11}{10}\varepsilon_N,\qquad
 6f\le\frac{11}{3750N^2}<\frac1{300}.                     \tag{101.17}
\]
The link bound from V89 gives \(\|A_{t_*}\|_\infty<r_N\).  Because
\(d_2F_t=d_2\mathcal N(A_t)\), (99.20), (98.12), and Theorem
\ref{thm:v18v91-bianchi-estimate} give
\[
 \|\Gamma_{t_*}\|_\infty
 \le N\left[3(e^{2r_N}-1)f+B(6f)\right].                  \tag{101.18}
\]
For \(x<1\), the exponential series gives \(e^x\le(1-x)^{-1}\).
Hence
\[
 3N(e^{2r_N}-1)
 \le\frac{3N}{75N-1}\le\frac3{74}.                        \tag{101.19}
\]
For \(0\le s\le1/300\), direct substitution in (101.7) gives
\[
 \frac{B(s)}{s^2}
 \le\frac{300}{299}\left(1+\frac{75}{149}\right)
 =\frac{67200}{44551}<\frac85.                            \tag{101.20}
\]
Divide (101.18) by \(\varepsilon_N\), and use
(101.17)--(101.20):
\[
\begin{split}
 \frac{\|\Gamma_{t_*}\|_\infty}{\varepsilon_N}
 &\le\frac3{74}\frac{11}{10}
 +\frac{288}{5}N\left(\frac{11}{10}\right)^2\varepsilon_N\\
 &\le\frac{33}{740}+\frac{484}{15625}
 =\frac{174757}{2312500}<\frac1{10}.                      \tag{101.21}
\end{split}
\]
This contradicts the definition of \(t_*\).  Therefore the barrier is never
crossed, proving (101.15).  Equation (101.16) follows from
\(F_1=G+\Gamma_1\).
\end{proof}

\begin{remark}[What closed the lost power of \(N\)]
\label{rem:v18v91-power}
The facewise estimate used \(B(4r_N)=O(N^{-2})\), then paid
\(\|\widehat H_3\|=O(N)\).  Equation (101.9) instead gives
\(d_2F=O(r_Nf+f^2)=O(N^{-3})\) once the continuity barrier keeps
\(f=O(N^{-2})\).  The same homotopy factor then yields
\(\Gamma_N=O(N^{-2})\).  The improvement comes from the exact transported
cube word, not from a smaller constant in the old factorwise bound.
\end{remark}

\subsection{Executable certificate}

The verifier
\[
 \texttt{scripts/verify\_sm\_v91\_lattice\_bianchi\_graph\_height.py}
                                                                  \tag{101.22}
\]
freely reduces the ordered word (101.3), checks (101.8) through degree two
in the free associative algebra on six generators, and verifies the exact
rational inequalities (101.20)--(101.21).  It emits
\[
 \texttt{artifacts/sm\_v91\_lattice\_bianchi\_graph\_height.json}. \tag{101.23}
\]
The SHA--256 values are
\[
\begin{array}{c|l}
\hbox{object}&\hbox{SHA--256}\\ \hline
\hbox{verifier}&
\texttt{47CC58C85C0AB1A3AD17B541CA36DA6B4723D2262FD919067875F37C716EE73D}\\
\hbox{canonical payload}&
\texttt{8B2FF41F9A872EEE2C3BC4A76308F38F24AFA90AB5C28614F4A3189B6C288831}.
\end{array}                                                   \tag{101.24}
\]

\begin{firewall}[Uniform chart scale is not yet a continuum measure]
\label{fire:v18v91}
Theorem \ref{thm:v18v91-graph-height} proves a volume-uniform scale for the
finite-box nonlinear curvature graph.  It does not prove consistency of
these graphs under restriction, a volume-uniform Jacobian for their
coordinates, reflection positivity of a measure supported on them,
gauge-orbit control, a Standard Model fermion determinant, or an Einstein
continuum limit.
\end{firewall}

\begin{decision}[Next bottleneck]
\label{dec:v18v91-next}
The graph-height power loss is closed.  V92 should compute the derivative
of \(\Psi_N\) and its deviation from the compatible inclusion on the same
\(N^{-2}\) chart.  The immediate target is a volume-uniform invertibility
and determinant-density statement for the finite-dimensional graph
coordinate, with the determinant claim withheld unless a local or
trace-class factorisation is actually proved.
\end{decision}

\begin{assessment}[V91 acquisition]
\label{ass:v18v91}
V91 proves the exact transported lattice Bianchi identity, identifies the
quadratic cancellation, and upgrades the isotropic nonlinear graph height
from \(O(N^{-1})\) to the natural \(O(N^{-2})\) plaquette scale.  This is a
finite-chart cofinal estimate, not a construction of the full SM--Einstein
continuum.
\end{assessment}

\begin{record}[V91 append record]
\label{rec:v18v91}
V91 is appended to the validated V90 whose installed SHA--256 was
\texttt{7CD74400A4BA5830860F45660822D3D7BEFFF47D06494C6DC7E89D680D608351}.
After installation only V91 is the active numeric SM note; the frozen
archives and unrelated notes are unchanged.
\end{record}


\section{V92: a volume-uniform differential chart for the curvature graph}
\label{sec:v18v92}

V91 closed the graph-height power loss.  The next question is whether the
graph can fold or develop an unbounded slope as the box grows.  This section
differentiates the fixed-point and exact cube identities, proves a uniform
small-slope estimate in the infinity norm, and separates that statement from
the still-open induced-volume determinant problem.

\begin{remark}[Correction of the denominator in (101.16)]
\label{corr:v18v92-denominator}
The last equality in (101.16) contains a typographical denominator:
\[
 \frac{11}{10}\varepsilon_N
 =\frac{11}{22500N^2},                                    \tag{102.1}
\]
not \(11/(24750N^2)\).  The barrier proof (101.17)--(101.21) uses
\(\varepsilon_N=1/(2250N^2)\) and is unchanged.  In this section set
\[
 \bar f_N:=\frac{11}{22500N^2}.                            \tag{102.2}
\]
\end{remark}

\subsection{Differentiating the fixed point}

\begin{lemma}[Link-coordinate differential]
\label{lem:v18v92-link-differential}
On the admitted compatible ball, \(G\mapsto A_G\) is analytic.  For a
tangent vector \(\dot G\in\mathcal E_N\),
\[
 \bigl(I+\widehat H_2D\mathcal N(A_G)\bigr)\dot A
 =\widehat H_2\dot G,\qquad \dot A=DA_G[\dot G],           \tag{102.3}
\]
and
\[
 \|\dot A\|_\infty
 \le\frac{138}{91}C_N^{(2)}\|\dot G\|_\infty
 =\frac{207}{91}N\|\dot G\|_\infty.                       \tag{102.4}
\]
\end{lemma}

\begin{proof}
The matrix exponential and principal logarithm are analytic on the V77
chart, so \(\mathcal N\) is analytic there.  Differentiating
\(A_G=\widehat H_2(G-\mathcal N(A_G))\) gives (102.3).
The contraction estimate (99.9) says
\[
 \|\widehat H_2D\mathcal N(A_G)\|_{\infty\to\infty}
 \le\frac{47}{138}.
\]
Its negative has a convergent Neumann series, so the implicit solution is
analytic and the inverse in (102.3) has norm at most
\((1-47/138)^{-1}=138/91\).  Insert
\(C_N^{(2)}=3N/2\).
\end{proof}

\subsection{A differentiated cube estimate}

\begin{lemma}[Derivative of the six-factor remainder]
\label{lem:v18v92-bch-derivative}
Under the hypotheses of Lemma \ref{lem:v18v91-six-bch}, a variation
\(\dot Y=(\dot Y_1,\ldots,\dot Y_6)\) satisfies
\[
 \|D{\cal R}_6(Y)[\dot Y]\|_{\rm op}
 \le B'(S)\sum_{j=1}^6\|\dot Y_j\|_{\rm op}.               \tag{102.5}
\]
For \(0\le S\le1/300\),
\[
 B'(S)\le\frac{16}{5}S.                                   \tag{102.6}
\]
\end{lemma}

\begin{proof}
The proof of (101.7) majorises the absolutely convergent product and
logarithm series coefficient by coefficient with nonnegative scalar series
\(B(S)\).  Termwise directional differentiation marks one factor in every
monomial.  Summing over the marked positions gives exactly the scalar
derivative \(B'(S)\) times
\(\sum_j\|\dot Y_j\|\), proving (102.5).

The power series of \(B\) has nonnegative coefficients from degree two
onward, so \(B'(S)/S\) is increasing.  Direct differentiation at
\(S=1/300\) gives
\[
 \frac{B'(1/300)}{1/300}
 =\frac{6004394700}{1984791601}<\frac{16}{5},              \tag{102.7}
\]
which proves (102.6).
\end{proof}

Let \(A,F\) vary through link fields and their full plaquette logs.  Put
\[
 a=\|\dot A\|_\infty,\qquad
 h=\|\dot F\|_\infty,\qquad
 r=\|A\|_\infty,\qquad f=\|F\|_\infty.
\]

\begin{theorem}[Differential covariant Bianchi bound]
\label{thm:v18v92-differential-bianchi}
If \(6f\le1/300\), then
\[
\begin{split}
 \|d_2\dot F\|_\infty
 \le{}&
 \left[3(e^{2r}-1)+\frac{576}{5}f\right]h\\
 &+\left[6e^{2r}f+\frac{576}{5}e^{2r}f^2\right]a.
                                                                  \tag{102.8}
\end{split}
\]
\end{theorem}

\begin{proof}
Differentiate (101.10).  For one transported face,
the exponential derivative formula on the operator
\(e^{\operatorname{ad}_A}\) gives
\[
 \|D_A\operatorname{Ad}_{e^A}[\dot A]F\|
 \le2e^{2r}f\,a.                                           \tag{102.9}
\]
Therefore the derivative of the three terms in \({\cal C}\) is bounded by
\[
 3(e^{2r}-1)h+6e^{2r}f\,a.                                \tag{102.10}
\]
Among the six \(Y_j\), three are transported and three are not.  Hence
\[
 \sum_{j=1}^6\|\dot Y_j\|
 \le6h+6e^{2r}f\,a.                                       \tag{102.11}
\]
Since \(S\le6f\), (102.5)--(102.6) bound the differentiated BCH remainder
by
\[
 \frac{576}{5}f\,h+\frac{576}{5}e^{2r}f^2a.               \tag{102.12}
\]
Adding (102.10) and (102.12) proves (102.8).
\end{proof}

\subsection{Uniform small slope}

\begin{theorem}[Uniform differential graph bound]
\label{thm:v18v92-graph-differential}
For every \(N\ge1\), every admitted \(G\), and every
\(\dot G\in\mathcal E_N\),
\[
\boxed{
 \|D\Gamma_N(G)[\dot G]\|_\infty
 \le\frac3{25}\|\dot G\|_\infty.}                          \tag{102.13}
\]
The sharper certified coefficient is
\[
 \frac{109097867}{950272050}<\frac3{25}.                  \tag{102.14}
\]
Consequently, for all admitted \(G,G'\),
\[
 \frac{22}{25}\|G-G'\|_\infty
 \le\|\Psi_N(G)-\Psi_N(G')\|_\infty
 \le\frac{28}{25}\|G-G'\|_\infty.                         \tag{102.15}
\]
\end{theorem}

\begin{proof}
Write \(g=\|\dot G\|_\infty\) and
\(z=\|D\Gamma_N(G)[\dot G]\|_\infty\).  Since
\(F=\Psi_N(G)=G+\Gamma_N(G)\), one has \(h\le g+z\).
Theorems \ref{thm:v18v91-graph-height} and
Remark \ref{corr:v18v92-denominator} give \(f\le\bar f_N\), while
\(r\le r_N\).  Moreover
\[
 D\Gamma_N(G)[\dot G]
 =\widehat H_3d_2\dot F,\qquad
 \|\widehat H_3\|_{\infty\to\infty}\le N.                 \tag{102.16}
\]
Apply (102.8), (102.4), and the estimates
\(e^{2r_N}\le75/74\) and
\(e^{2r_N}-1\le(75N-1)^{-1}\).  The coefficient of \(h\) after
multiplication by \(N\) is at most
\[
 u_0:=\frac3{74}+\frac{576}{5}\frac{11}{22500}
 =\frac{22399}{231250}.                                   \tag{102.17}
\]
The remaining coefficient of \(g\), including (102.4), is at most
\[
 v_0:=\frac{7182417}{1052187500}.                         \tag{102.18}
\]
Thus
\[
 z\le u_0(g+z)+v_0g.
\]
Since \(u_0<1\), absorption gives
\[
 z\le\frac{u_0+v_0}{1-u_0}g
 =\frac{109097867}{950272050}g<\frac3{25}g,               \tag{102.19}
\]
proving (102.13).  Integrate this derivative on the line segment from
\(G\) to \(G'\); the compatible ball is convex.  Then use
\(\Psi_N(G)-\Psi_N(G')=(G-G')+(\Gamma_N(G)-\Gamma_N(G'))\)
and the two triangle inequalities to obtain (102.15).
\end{proof}

\begin{corollary}[No folding in the cofinal chart]
\label{cor:v18v92-no-folding}
The curvature graph is a uniformly bi-Lipschitz analytic embedding of the
compatible ball, and
\[
 P_ND\Psi_N(G)=I_{\mathcal E_N}.                           \tag{102.20}
\]
In particular, its coordinate never loses rank as \(N\) grows.
\end{corollary}

\begin{proof}
The analytic statement follows from Lemma
\ref{lem:v18v92-link-differential}.  The lower bound in (102.15) gives an
embedding and excludes folding.  Differentiating the exact identity
\(P_N\Psi_N(G)=G\) from (99.15) gives (102.20).
\end{proof}

\subsection{Executable certificate}

The verifier
\[
 \texttt{scripts/verify\_sm\_v92\_graph\_differential.py} \tag{102.21}
\]
checks the rational derivative majorant and every constant in
(102.17)--(102.19), and emits
\[
 \texttt{artifacts/sm\_v92\_graph\_differential.json}.     \tag{102.22}
\]
Its hashes are
\[
\begin{array}{c|l}
\hbox{object}&\hbox{SHA--256}\\ \hline
\hbox{verifier}&
\texttt{D1C689B8FB8B8806BC2AEC81803619D1D86BF04E852F619525C4A337D49C1199}\\
\hbox{canonical payload}&
\texttt{292D763B02B6A93BB2DA38894147AEB297A5BC93F6C020B229DE12D30E4AB407}.
\end{array}                                                   \tag{102.23}
\]

\begin{firewall}[A small graph slope is not a determinant-density theorem]
\label{fire:v18v92}
The infinity-norm bound (102.13) is uniform in \(N\), but the dimension of
\(\mathcal E_N\) grows with the volume.  It does not by itself bound the
Jacobian of ambient Hausdorff measure relative to \(G\)-coordinate measure:
a fixed singular-value distortion can accumulate over all dimensions.
No volume-uniform determinant, Gibbs density, reflection-positive measure,
or continuum field has yet been constructed.
\end{firewall}

\begin{decision}[Next bottleneck]
\label{dec:v18v92-next}
V93 should use \(P_N\Psi_N=I\) to define the finite curvature law directly
in the compatible coordinate \(G\), so no ambient graph Jacobian is inserted
by convention.  It must then test whether restriction from a larger box
commutes with the reflection-averaged homotopy chart.  Failure of that
commutation would require a local specification or boundary-kernel
correction before any projective continuum claim.
\end{decision}

\begin{assessment}[V92 acquisition]
\label{ass:v18v92}
V92 corrects the denominator in (101.16), differentiates the exact lattice
Bianchi cancellation, and proves a uniform graph slope below \(3/25\).
The small-curvature graph is now a cofinally nonfolding analytic chart.
The measure-consistency problem remains open.
\end{assessment}

\begin{record}[V92 append record]
\label{rec:v18v92}
V92 is appended to the validated V91 whose installed SHA--256 was
\texttt{E914D25B5D1FCF291E4057E11D74A961CFD837CFD4F6DB948A0EC90B261E327F}.
After installation only V92 is the active numeric SM note; the frozen
archives and unrelated notes are unchanged.
\end{record}


\section{V93: restriction failure of the averaged chart and the rooted repair}
\label{sec:v18v93}

V92 made the reflection-averaged curvature graph uniformly nonfolding.  A
continuum construction also needs compatibility between nested boxes.  This
section tests that requirement exactly.  Reflection averaging introduces
dependence on the far endpoint and therefore fails anchored restriction.
The failure on compatible linear curvatures is an exact additive gauge, while
the original corner-rooted homotopy commutes with restriction in every
degree.  This identifies a projective chart candidate for the next version.

\subsection{Restriction maps}

Let
\[
 Q_{\mathbf m}\subset Q_{\mathbf M},\qquad
 1\le m_\alpha\le M_\alpha,                               \tag{103.1}
\]
be cubical boxes with the same corner at the origin.  Let
\[
 R_p^{\mathbf M,\mathbf m}:
 C^p(Q_{\mathbf M})\longrightarrow C^p(Q_{\mathbf m})     \tag{103.2}
\]
be cellular restriction.  It commutes with coboundaries:
\[
 R_{p+1}^{\mathbf M,\mathbf m}d_p^{\mathbf M}
 =d_p^{\mathbf m}R_p^{\mathbf M,\mathbf m}.                \tag{103.3}
\]

\begin{theorem}[Rooted homotopy is exactly projective]
\label{thm:v18v93-rooted-projective}
For the corner-rooted tensor homotopy \(H_p\) of (98.4),
\[
\boxed{
 R_{p-1}^{\mathbf M,\mathbf m}H_p^{\mathbf M}
 =H_p^{\mathbf m}R_p^{\mathbf M,\mathbf m}}
 \qquad(1\le p\le4).                                      \tag{103.4}
\]
\end{theorem}

\begin{proof}
Evaluate both sides on an output \((p-1)\)-cell contained in
\(Q_{\mathbf m}\).  In the \(k\)-th summand of (98.4), every preceding
vertex coordinate is evaluated at the common origin and the \(k\)-th
coordinate is summed only over the prefix from zero to the output
coordinate.  Every input cell used by that row therefore also lies in
\(Q_{\mathbf m}\), and its coefficient is independent of the larger endpoint
\(\mathbf M\).  The two rows agree cell by cell.  This proves (103.4).
\end{proof}

\subsection{Why reflection averaging breaks restriction}

On a single interval, (98.13) reads
\[
 (\widehat h_Mf)(x)
 =\frac12\left(\sum_{i=0}^{x-1}f(i)
              -\sum_{i=x}^{M-1}f(i)\right).               \tag{103.5}
\]
Let \(m<M\) and let \(f\) be supported on the first exterior edge \(i=m\).
Then \(R_1^{M,m}f=0\), whereas for \(0\le x\le m\),
\[
 R_0^{M,m}\widehat h_Mf(x)=-\frac12.                       \tag{103.6}
\]
Thus \(R_0\widehat h_M\ne\widehat h_mR_1\).

\begin{proposition}[Failure on a closed four-dimensional face cochain]
\label{prop:v18v93-averaged-failure}
There are nested four-dimensional boxes and a closed face cochain
\(F\in\ker d_2^{\mathbf M}\) such that
\[
 R_2^{\mathbf M,\mathbf m}F=0,\qquad
 R_1^{\mathbf M,\mathbf m}\widehat H_2^{\mathbf M}F\ne0.
                                                                  \tag{103.7}
\]
For \(\mathbf M=(3,2,2,2)\) and \(\mathbf m=(2,2,2,2)\), one may take
\(F\) equal to one on every \((0,1)\)-face with first base coordinate
\(2\), constant in the other transverse vertex coordinates, and zero
elsewhere.  The resulting restriction defect has infinity norm \(1/2\).
\end{proposition}

\begin{proof}
The probe is constant in coordinates \(2\) and \(3\).  Those are the only
transverse differences of a \((0,1)\)-face that enter \(d_2\), so it is
closed.  Its support lies outside \(Q_{\mathbf m}\), proving the first
relation in (103.7).  On output edges in direction \(1\), the degree-two
tensor homotopy contracts coordinate \(0\).  Formula (103.6), with the
remaining coordinates held fixed, gives a coefficient of magnitude
\(1/2\) on inner edges.  Averaging reflections in the other coordinates
does not alter the constant transverse probe.  This proves the second
relation and the norm assertion.
\end{proof}

\subsection{The linear defect is gauge}

\begin{theorem}[Restriction mismatch on compatible inputs is exact]
\label{thm:v18v93-linear-gauge}
For every closed \(F\in C^2(Q_{\mathbf M})\), define
\[
 D(F):=
 R_1^{\mathbf M,\mathbf m}\widehat H_2^{\mathbf M}F
 -\widehat H_2^{\mathbf m}R_2^{\mathbf M,\mathbf m}F.      \tag{103.8}
\]
Then
\[
 d_1^{\mathbf m}D(F)=0,\qquad
 D(F)=d_0^{\mathbf m}\widehat H_1^{\mathbf m}D(F).         \tag{103.9}
\]
Thus the failure of the reflection-averaged linear reconstruction to
restrict is an exact additive gauge cochain.
\end{theorem}

\begin{proof}
Use (103.3), closedness, and the right-inverse part of (98.9):
\[
\begin{split}
 d_1^{\mathbf m}D(F)
 &=R_2d_1^{\mathbf M}\widehat H_2^{\mathbf M}F
   -d_1^{\mathbf m}\widehat H_2^{\mathbf m}R_2F\\
 &=R_2F-R_2F=0.                                           \tag{103.10}
\end{split}
\]
Apply (98.9) in degree one on \(Q_{\mathbf m}\):
\[
 d_0\widehat H_1D+\widehat H_2d_1D=D.
\]
The second term vanishes by (103.10), proving (103.9).
\end{proof}

\begin{remark}[What the gauge statement does not yet prove]
\label{rem:v18v93-nonlinear}
Equation (103.9) concerns the linearised gauge action
\(A\mapsto A+d_0\phi\).  A finite lattice gauge transformation acts on link
holonomies by endpoint multiplication, and its logarithm contains BCH
terms.  Hence (103.9) alone does not prove that the nonlinear averaged
fixed-point graphs commute with restriction modulo gauge.
\end{remark}

\subsection{Executable certificate}

The verifier
\[
 \texttt{scripts/verify\_sm\_v93\_homotopy\_restriction.py} \tag{103.11}
\]
builds the rooted and reflection-averaged homotopies over the rationals for
the boxes in Proposition \ref{prop:v18v93-averaged-failure}.  It verifies
(103.4) in degrees \(1,2,3,4\), finds failure of the averaged commutator in
all four degrees, constructs the closed face slab, and checks its exact
gauge primitive.  It emits
\[
 \texttt{artifacts/sm\_v93\_homotopy\_restriction.json}.    \tag{103.12}
\]
The hashes are
\[
\begin{array}{c|l}
\hbox{object}&\hbox{SHA--256}\\ \hline
\hbox{verifier}&
\texttt{0BA84F7D7B008486FD161234943F7BCC7258FFB12653C7582471A8C58600D452}\\
\hbox{canonical payload}&
\texttt{7DD7E6BF74FBCBBD4C29FE504CC2899E458EB7D4B1C39718E1AEED089718F33A}.
\end{array}                                                   \tag{103.13}
\]

\begin{firewall}[Projective homotopy is not yet a projective field law]
\label{fire:v18v93}
Theorem \ref{thm:v18v93-rooted-projective} is a linear operator identity.
It does not yet establish existence and uniqueness for a rooted nonlinear
chart at volume-dependent radii, compatibility of those nonlinear fixed
points, reflection positivity, or a limiting probability measure.
Reflection equivariance and anchored projectivity are presently obtained by
different homotopies.
\end{firewall}

\begin{decision}[V94 acquisition]
\label{dec:v18v93-next}
V94 should build the nonlinear fixed-point graph with the rooted homotopy.
Its radii must be rescaled using
\(\|H_2\|_{\infty\to\infty}\le3N\) and
\(\|H_3\|_{\infty\to\infty}\le2N\).  It should prove exact restriction
compatibility of fixed points by applying (103.4) and locality of the
plaquette remainder, and then repeat the V91 barrier with the new constants.
\end{decision}

\begin{assessment}[V93 acquisition]
\label{ass:v18v93}
V93 identifies the first projective-consistency obstruction.  The
reflection-averaged chart is nonprojective, although its linear mismatch on
closed curvatures is pure gauge.  The corner-rooted homotopy is exactly
projective and supplies a concrete nonlinear repair route.
\end{assessment}

\begin{record}[V93 append record]
\label{rec:v18v93}
V93 is appended to the validated V92 whose installed SHA--256 was
\texttt{8CEB5FBA6101639891EB421A5F801EF78DF3F34721B7D48B98789CC84124D4FB}.
After installation only V93 is the active numeric SM note; the frozen
archives and unrelated notes are unchanged.
\end{record}


\section{V94: a rooted projective graph and the shrinking-radius obstruction}
\label{sec:v18v94}

V93 showed that the corner-rooted homotopy commutes with restriction.  This
section carries the nonlinear construction over to that homotopy.  The
resulting fixed points and full curvature graphs are exactly compatible on
nested anchored boxes.  The same calculation exposes a decisive limitation:
if the \(N^{-2}\) chart radius is tied to the total side length and ordinary
cell restriction is used, every infinite projective family in those balls is
zero.

\subsection{Rooted graph constants}

On the isotropic box \(Q_N=[0,N]^4\), retain the unaveraged homotopy \(H_p\)
of (98.4).  Its row-stock bounds are
\[
 \|H_2\|_{\infty\to\infty}\le3N,\qquad
 \|H_3\|_{\infty\to\infty}\le2N.                          \tag{104.1}
\]
Define
\[
 P_N^\circ=d_1H_2,\qquad Q_N^\circ=H_3d_2,\qquad
 r_N^\circ=\frac1{300N},\qquad
 \varepsilon_N^\circ=\frac1{9000N^2}.                    \tag{104.2}
\]
The homotopy identity gives
\[
 P_N^\circ+Q_N^\circ=I,\qquad
 \operatorname{ran}P_N^\circ=\mathcal E_N.                \tag{104.3}
\]

For \(G\in\mathcal E_N\) with
\(\|G\|_\infty\le\varepsilon_N^\circ\), set
\[
 {\cal T}_G^\circ(A)=H_2(G-\mathcal N(A)).                 \tag{104.4}
\]

\begin{theorem}[Rooted nonlinear curvature graph]
\label{thm:v18v94-rooted-graph}
The map (104.4) has a unique fixed point \(A_G^\circ\) in the
radius-\(r_N^\circ\) ball.  Define
\[
 \Gamma_N^\circ(G)=Q_N^\circ\mathcal N(A_G^\circ),\qquad
 \Psi_N^\circ(G)=G+\Gamma_N^\circ(G).                     \tag{104.5}
\]
Then
\[
 \mathcal F(A_G^\circ)=\Psi_N^\circ(G),\qquad
 P_N^\circ\Psi_N^\circ(G)=G,                              \tag{104.6}
\]
and
\[
 \|\Gamma_N^\circ(G)\|_\infty<\frac1{90000N^2},\qquad
 \|\Psi_N^\circ(G)\|_\infty\le\frac{11}{90000N^2}.        \tag{104.7}
\]
The graph differential satisfies
\[
 \|D\Gamma_N^\circ(G)\|_{\infty\to\infty}
 \le\frac{79618373}{986810825}<\frac1{12}.                \tag{104.8}
\]
\end{theorem}

\begin{proof}
With \(C_2=3N\), the choices in (104.2) are exactly
\(r=(100C_2)^{-1}\) and \(\varepsilon=(1000C_2^2)^{-1}\).
The self-map and derivative estimates of V88 therefore give the same
contraction coefficient \(47/138\).  Banach's theorem proves existence and
uniqueness.  Applying \(d_1\) to the fixed-point equation and using
(104.3) proves (104.6), exactly as in V89.

It remains to record the two changed homotopy constants.  At a hypothetical
first crossing
\(\|\Gamma_N^\circ(tG)\|=\varepsilon_N^\circ/10\), the full face norm is at
most \(11/(90000N^2)\).  Equations (101.9) and (104.1) give, after division
by \(\varepsilon_N^\circ\),
\[
 \frac{\|\Gamma_N^\circ(tG)\|}{\varepsilon_N^\circ}
 \le\frac6{149}+\frac{242}{15625}
 =\frac{129808}{2328125}<\frac1{10}.                      \tag{104.9}
\]
Here \(e^{1/(150N)}-1\le(150N-1)^{-1}\), and (101.20) controls
\(B(6f)\).  The continuity argument of V91 proves (104.7).

For the derivative, (102.3) now gives
\[
 \|DA_G^\circ\|\le\frac{414}{91}N.                        \tag{104.10}
\]
Using \(2N\) in place of \(N\), \(r_N^\circ\), and the face bound in
(104.7), the coefficient of \(\|\dot F\|\) in the analogue of (102.16) is
\[
 u_\circ=\frac{31862}{465625},                             \tag{104.11}
\]
and the coefficient contributed through \(\dot A\) is
\[
 v_\circ=\frac{7132323}{1059296875}.                       \tag{104.12}
\]
Thus \(z\le u_\circ(g+z)+v_\circ g\), and absorption gives
\[
 \frac{u_\circ+v_\circ}{1-u_\circ}
 =\frac{79618373}{986810825}<\frac1{12},                  \tag{104.13}
\]
which is (104.8).
\end{proof}

\subsection{Exact nonlinear restriction compatibility}

\begin{theorem}[Rooted fixed points are projective]
\label{thm:v18v94-nonlinear-projective}
Let \(M\ge N\), let \(Q_N\subset Q_M\) share the root corner, and let
\(G_M\in\mathcal E_M\) satisfy
\(\|G_M\|_\infty\le\varepsilon_M^\circ\).  Put
\(G_N=R_2^{M,N}G_M\).  Then
\[
 R_1^{M,N}A_{G_M}^\circ=A_{G_N}^\circ,                    \tag{104.14}
\]
and
\[
 R_2^{M,N}\Psi_M^\circ(G_M)=\Psi_N^\circ(G_N).            \tag{104.15}
\]
\end{theorem}

\begin{proof}
Restriction commutes with \(d_2\), so \(G_N\in\mathcal E_N\), and
\(\|G_N\|\le\varepsilon_M^\circ\le\varepsilon_N^\circ\).
Plaquette products are local:
\[
 R_2^{M,N}\mathcal N(A)=
 \mathcal N(R_1^{M,N}A).                                  \tag{104.16}
\]
Apply \(R_1^{M,N}\) to the large fixed-point equation, then use
Theorem \ref{thm:v18v93-rooted-projective} and (104.16):
\[
 R_1A_{G_M}^\circ
 =H_2^N\bigl(G_N-\mathcal N(R_1A_{G_M}^\circ)\bigr).
                                                                  \tag{104.17}
\]
The restricted field has norm at most
\(r_M^\circ\le r_N^\circ\), so it lies in the small uniqueness ball.
Equation (104.17) therefore proves (104.14).  Locality of the full
plaquette-log map and (104.6) then give (104.15).
\end{proof}

\begin{corollary}[Pushforward preserves any admitted projective coordinate law]
\label{cor:v18v94-pushforward}
If probability laws \(\mu_N\) on the admitted compatible balls satisfy
\((R_2^{M,N})_\#\mu_M=\mu_N\), then
\[
 \nu_N:=(\Psi_N^\circ)_\#\mu_N
\]
satisfy \((R_2^{M,N})_\#\nu_M=\nu_N\).
\end{corollary}

\begin{proof}
This is the functorial pushforward consequence of (104.15).
\end{proof}

\subsection{The projective family inside these balls is trivial}

\begin{theorem}[Shrinking-radius obstruction]
\label{thm:v18v94-trivial-projective}
Let \(G_N\in\mathcal E_N\) be a deterministic family for all \(N\ge1\)
such that
\[
 R_2^{M,N}G_M=G_N\quad(M\ge N),\qquad
 \|G_N\|_\infty\le\frac1{9000N^2}.                        \tag{104.18}
\]
Then \(G_N=0\) for every \(N\).  The same conclusion holds almost surely
for any projective family of laws supported in these balls.  Consequently
Corollary \ref{cor:v18v94-pushforward} yields only the zero field if
ordinary cell restriction and the global \(N^{-2}\) radius are imposed at
every stage.
\end{theorem}

\begin{proof}
Fix a face \(p\) contained in some \(Q_{N_0}\).  For every \(M\ge N_0\),
projectivity gives
\[
 G_{N_0}(p)=G_M(p),
\]
while (104.18) gives
\[
 |G_{N_0}(p)|\le\frac1{9000M^2}.
\]
Letting \(M\to\infty\) proves \(G_{N_0}(p)=0\).  Every face occurs in some
finite box, proving the deterministic assertion.  For projective laws, the
same coordinate random variable is almost surely bounded by
\(1/(9000M^2)\) for every \(M\); its absolute value is therefore zero almost
surely.  Countability of the lattice faces completes the probabilistic
statement.
\end{proof}

\begin{remark}[Two limits were being conflated]
\label{rem:v18v94-two-limits}
Ordinary restriction of nested lattice boxes describes volume exhaustion at
fixed mesh.  The scaling \(G=O(N^{-2})\) describes face holonomy under mesh
refinement when the mesh is \(a\asymp N^{-1}\).  Combining the restriction
map from the first problem with the shrinking amplitude from the second
forces Theorem \ref{thm:v18v94-trivial-projective}.  A nontrivial programme
must either keep the local chart radius independent of the exhausting
volume, or introduce a refinement and blocking map that rescales curvature.
\end{remark}

\subsection{Executable certificate}

The verifier
\[
 \texttt{scripts/verify\_sm\_v94\_rooted\_projective\_chart.py}     \tag{104.19}
\]
checks all rational constants in (104.9)--(104.13) and the strict shrinking
of the declared coordinate radii.  It emits
\[
 \texttt{artifacts/sm\_v94\_rooted\_projective\_chart.json}.        \tag{104.20}
\]
Its hashes are
\[
\begin{array}{c|l}
\hbox{object}&\hbox{SHA--256}\\ \hline
\hbox{verifier}&
\texttt{E9C1B876DB5B9013B157292A20E51A1A4DDAAA92805EC3F74179230E4D202D87}\\
\hbox{canonical payload}&
\texttt{4EE14F1A17DB728378BD27E7361BEA3AF16C6A871403059DC5533B249951B6B6}.
\end{array}                                                   \tag{104.21}
\]

\begin{firewall}[Exact projectivity can still be physically empty]
\label{fire:v18v94}
The rooted graph commutes with anchored restriction, but Theorem
\ref{thm:v18v94-trivial-projective} shows that the present global shrinking
balls support only the zero raw-restriction limit.  No nontrivial continuum
measure follows.  The rooted chart also sacrifices the manifest reflection
equivariance of the averaged chart.
\end{firewall}

\begin{decision}[Mandatory checkpoint and separation of limits]
\label{dec:v18v94-next}
V95 is the mandatory companion audit.  V96 must then separate volume
exhaustion from mesh refinement.  It should define the curvature blocking
map between two mesh scales, including the necessary parallel transports in
the nonabelian case, and test its linearisation before asserting any
projective continuum family.
\end{decision}

\begin{assessment}[V94 acquisition]
\label{ass:v18v94}
V94 constructs a unique rooted nonlinear curvature graph with
\(O(N^{-2})\) height, slope below \(1/12\), and exact anchored restriction
compatibility.  It also proves that combining ordinary restriction with the
global \(N^{-2}\) radius gives only the zero projective family.  The next
programme must use the correct map for each limiting operation.
\end{assessment}

\begin{record}[V94 append record]
\label{rec:v18v94}
V94 is appended to the validated V93 whose installed SHA--256 was
\texttt{2DEB4F413CE2ECD4FB055DE3E36DFEBDF3E7F340A60D542BE06A9DA1BB7ACADA}.
After installation only V94 is the active numeric SM note; the frozen
archives and unrelated notes are unchanged.
\end{record}


\section{V95: mandatory companion audit before the refinement map}
\label{sec:v18v95}

V94 separated anchored volume restriction from mesh refinement and proved
that applying the \(N^{-2}\) chart radius to an ordinary restriction system
forces the zero field.  Before constructing a refinement map, the standing
five-version rule requires a new companion audit.

\subsection{Source record}

The active files read at this checkpoint are
\[
\begin{array}{c|l}
\hbox{programme and active file}&\hbox{SHA--256}\\ \hline
\hbox{Yang--Mills,
 \texttt{Renewal\_Yang\_Mills\_Mass\_Gap\_Research\_Notes\_V23.tex}}&
\texttt{5DA58A569F82DF22D7F070266BCFBA353241F3E1653595E343AC396375551878}\\
\hbox{Navier--Stokes,
 \texttt{Renewal\_NS\_Stability\_Research\_Notes\_V26.tex}}&
\texttt{82C887F8C2C69E4F28FAD7C3DC8DE5B9099CB5A4BF431EBA1FFF3E91935978AD}.
\end{array}                                                   \tag{105.1}
\]
The Navier--Stokes file is unchanged.  The Yang--Mills file has advanced
from V21 to V23.

\begin{proposition}[Transferable conclusions]
\label{prop:v18v95-transfer}
The audit changes the refinement attack in two ways.
\begin{enumerate}
\item The nonabelian blocking map must be derived as an exact local group
      word before seeking a reduced coordinate representation.
\item Its executable test should batch independent coarse cells only after
      the exact word and its linearisation have been certified.
\end{enumerate}
No companion theorem is imported as an SM--Einstein theorem.
\end{proposition}

\begin{proof}
Yang--Mills V22 proves that sixteen rank-one coefficient matrices together
with its bridge generate the full algebra \(M_{45}\), and that every
standard nontrivial Kronecker reshape of the bridge has maximal rank.  Thus
coordinate sparsity, low rank of individual pieces, and a preferred tensor
labelling do not imply a surviving common reduction.  V23 then proves that
the trace-influence mask is also full, but obtains a practical exact kernel
by batching the actual scalar contractions and certifying binary64 limb
exactness.  The relevant lesson is the order of operations: establish the
exact algebraic producer, rule out unsupported reductions, and optimize the
declared output.

For curvature blocking, the producer is the holonomy around a coarse
plaquette.  In a nonabelian theory its fine plaquettes must be transported
to a common base vertex before multiplication.  Those transports depend on
link data and cannot be discarded merely because the linearised curvature
is a face cochain.  This proves item 1.  Once the local word is fixed,
different coarse plaquettes can be evaluated independently, which is the
legitimate analogue of the V23 batching step and proves item 2.

Navier--Stokes V26 still gives the same restricted-input lesson recorded in
V90: a norm may be sharpened on the inputs actually produced, but the gain
has to be computed and can vanish.  It supplies no new refinement identity.
\end{proof}

\begin{firewall}[Audit boundary]
\label{fire:v18v95}
The full-algebra result concerns a finite Yang--Mills response producer, and
the rank-one result concerns Oseen-kernel contractions.  Neither proves a
blocking identity, continuum measure, Standard Model interaction, or
Einstein limit.  Only their verified methodological constraints are used.
\end{firewall}

\begin{decision}[V96 acquisition]
\label{dec:v18v95-next}
V96 will define dyadic coarsening on link holonomies by ordered multiplication
of the two fine links along each coarse edge.  It will freely reduce an
ordered product of four transported fine plaquettes to the coarse plaquette
boundary, prove associativity of repeated link blocking, and show that the
linearised rescaled curvature map is the average over the four fine faces.
\end{decision}

\begin{assessment}[V95 acquisition]
\label{ass:v18v95}
The mandatory audit is complete.  It supports an exact word-first,
output-specific refinement calculation and warns against a face-only
nonabelian shortcut.  The V94 zero-field obstruction remains in force for
ordinary restriction.
\end{assessment}

\begin{record}[V95 append record]
\label{rec:v18v95}
V95 is appended to the validated V94 whose installed SHA--256 was
\texttt{8FA5D0A1AD5D57D48C11F367B5785785435F1516BE15BF1461A512C18B1B3004}.
After installation only V95 is the active numeric SM note; the frozen
archives and unrelated notes are unchanged.
\end{record}


\section{V96: exact dyadic link blocking and the rescaled curvature average}
\label{sec:v18v96}

V94 proved that ordinary cell restriction combined with a global
\(N^{-2}\) face radius has only the zero projective limit.  The correct map
for mesh refinement is coarsening, not restriction.  This section defines
dyadic coarsening at the link level, derives the coarse plaquette as an exact
word in transported fine plaquettes, and proves that its linearised
continuum-scaled curvature is the four-face average.

\subsection{Link blocking}

Let \(\Lambda_{2N}\) be the dyadic refinement of \(\Lambda_N\) on the same
physical box.  For a fine link field \(U\), define the coarse link
\[
 (\mathfrak B_1U)_\mu(x)
 :=U_\mu(2x)U_\mu(2x+e_\mu).                              \tag{106.1}
\]

\begin{proposition}[Associativity and gauge covariance]
\label{prop:v18v96-link-block}
Ordered link blocking is associative under repeated dyadic coarsening.  If
a fine gauge transformation is \(g\), then
\[
 (\mathfrak B_1U^g)_\mu(x)
 =g(2x)(\mathfrak B_1U)_\mu(x)g(2x+2e_\mu)^{-1}.           \tag{106.2}
\]
Thus the induced coarse gauge transformation is the restriction of \(g\)
to the even vertices.
\end{proposition}

\begin{proof}
Two successive blockings multiply four collinear fine links in their path
order.  Direct blocking by factor four gives the same ordered product by
associativity of the group law.  Under a gauge transformation, the endpoint
factor at the middle fine vertex in
\[
 g(2x)U_\mu(2x)g(2x+e_\mu)^{-1}
 g(2x+e_\mu)U_\mu(2x+e_\mu)g(2x+2e_\mu)^{-1}
\]
cancels, giving (106.2).
\end{proof}

\subsection{The four-plaquette word}

Fix one coarse plaquette and label the fine horizontal links by
\(h_{ij}\), from \((i,j)\) to \((i+1,j)\), and the vertical links by
\(v_{ij}\), from \((i,j)\) to \((i,j+1)\), for \(0\le i,j\le2\) whenever
the edge exists.  Put
\[
 p_{ij}=h_{ij}v_{i+1,j}h_{i,j+1}^{-1}v_{ij}^{-1},
 \qquad i,j\in\{0,1\}.                                    \tag{106.3}
\]
The coarse boundary holonomy is
\[
 P^{\rm c}
 =h_{00}h_{10}v_{20}v_{21}
  h_{12}^{-1}h_{02}^{-1}v_{01}^{-1}v_{00}^{-1}.           \tag{106.4}
\]

\begin{theorem}[Exact transported fine-face decomposition]
\label{thm:v18v96-four-face-word}
The coarse plaquette satisfies
\[
\boxed{
 P^{\rm c}
 =\operatorname{Ad}_{h_{00}}(p_{10})\,
  \operatorname{Ad}_{h_{00}v_{10}}(p_{11})\,
  p_{00}\,
  \operatorname{Ad}_{v_{00}}(p_{01}).}                    \tag{106.5}
\]
\end{theorem}

\begin{proof}
Expand the first two factors:
\[
\begin{split}
 &h_{00}p_{10}h_{00}^{-1}
 h_{00}v_{10}p_{11}v_{10}^{-1}h_{00}^{-1}\\
 &\quad=
 h_{00}h_{10}v_{20}v_{21}
 h_{12}^{-1}v_{11}^{-1}v_{10}^{-1}h_{00}^{-1}.            \tag{106.6}
\end{split}
\]
Multiplication by
\(p_{00}=h_{00}v_{10}h_{01}^{-1}v_{00}^{-1}\) cancels the last
\(h_{00}^{-1}h_{00}\) and \(v_{10}^{-1}v_{10}\).  Multiplication by
\[
 \operatorname{Ad}_{v_{00}}(p_{01})
 =v_{00}h_{01}v_{11}h_{02}^{-1}v_{01}^{-1}v_{00}^{-1}
\]
then cancels \(v_{00}^{-1}v_{00}\),
\(h_{01}^{-1}h_{01}\), and \(v_{11}^{-1}v_{11}\).  The reduced word is
exactly (106.4).
\end{proof}

\begin{corollary}[A nonabelian face block needs transport data]
\label{cor:v18v96-transport-data}
The coarse plaquette is determined locally by the four fine plaquettes
together with the three displayed transport paths.  A map of four
unbased plaquette matrices alone is not defined invariantly: changing the
base vertex conjugates the corresponding factor.
\end{corollary}

\begin{proof}
Equation (106.5) supplies the construction.  Under a gauge transformation,
each transported loop is conjugated at the common coarse base vertex, so
their product is covariant.  Omitting its transport leaves a loop based at
a different vertex and destroys that common conjugation law.
\end{proof}

\subsection{Linearised cellular block}

For a coarse \(p\)-cell \(c\), let
\[
 (S_p\omega)(c)
 :=\sum_{\substack{\widetilde c\ {\rm fine}\\
                    \widetilde c\ {\rm tiles}\ c}}
      \omega(\widetilde c),\qquad
 B_p:=2^{-p}S_p.                                           \tag{106.7}
\]
There are \(2^p\) fine \(p\)-cells in the sum.

\begin{theorem}[Subdivision cochain identity]
\label{thm:v18v96-subdivision}
For \(0\le p\le3\),
\[
 d_p^{N}S_p=S_{p+1}d_p^{2N},                              \tag{106.8}
\]
and consequently
\[
 d_p^{N}B_p=2B_{p+1}d_p^{2N}.                             \tag{106.9}
\]
In particular, \(B_2\) sends closed fine face cochains to closed coarse
face cochains and preserves constant face cochains.
\end{theorem}

\begin{proof}
Evaluate \(S_{p+1}d\omega\) on one coarse \((p+1)\)-cell.  It is the sum of
\(d\omega\) over its \(2^{p+1}\) fine subcells.  Every interior fine
\(p\)-face occurs twice with opposite incidence signs and cancels.  The
remaining boundary faces are precisely the \(2^p\) fine tiles of each
coarse boundary \(p\)-face, with the coarse incidence sign.  This is
\(dS_p\), proving (106.8).  Multiply by \(2^{-p}\) to obtain (106.9).
Closedness follows immediately.  A constant cochain is summed \(2^p\)
times and then divided by \(2^p\), so it is fixed.
\end{proof}

\begin{theorem}[Linearised rescaled curvature block]
\label{thm:v18v96-linear-curvature-block}
Let the fine link logarithms be \(tA\), and let \(F_{2N}(t)\) and
\(F_N^{\rm c}(t)\) be the fine and blocked coarse plaquette logarithms.
Then
\[
 \left.\frac{d}{dt}\right|_{t=0}F_N^{\rm c}(t)
 =S_2\left.\frac{d}{dt}\right|_{t=0}F_{2N}(t).             \tag{106.10}
\]
For the continuum-scaled face variables
\[
 G_N=N^2F_N,\qquad G_{2N}=(2N)^2F_{2N},                   \tag{106.11}
\]
the derivative of the refinement map is
\[
\boxed{
 DG_N^{\rm c}(0)=B_2\,DG_{2N}(0)
 =\frac14S_2\,DG_{2N}(0).}                                \tag{106.12}
\]
Thus a constant continuum curvature survives dyadic refinement.
\end{theorem}

\begin{proof}
In (106.5), every transport equals the identity at \(t=0\).
The derivative of a product at the identity is the sum of the four
derivatives, and the derivative of the logarithm at the identity is the
identity map.  This proves (106.10).  Multiply it by \(N^2\), substitute
\(F_{2N}=G_{2N}/(4N^2)\), and obtain (106.12).  The last assertion follows
from Theorem \ref{thm:v18v96-subdivision}.
\end{proof}

\begin{corollary}[Refinement avoids the V94 zero-field argument]
\label{cor:v18v96-nontrivial-linear}
The constant family \(G_N(c)=X\), for a fixed Lie-algebra element \(X\),
satisfies the linearised refinement relation
\[
 G_N=B_2G_{2N}
\]
for every \(N\).  Hence the shrinking-radius argument of Theorem
\ref{thm:v18v94-trivial-projective} does not apply to the rescaled
refinement coordinate.
\end{corollary}

\begin{proof}
Every coarse face averages four copies of \(X\), so \(B_2G_{2N}=X\).
The V94 argument used equality of the same unscaled coordinate under raw
cell restriction; neither condition holds here.
\end{proof}

\subsection{Executable certificate}

The verifier
\[
 \texttt{scripts/verify\_sm\_v96\_dyadic\_curvature\_blocking.py}    \tag{106.13}
\]
freely reduces (106.5) to (106.4), checks associativity on four fine links,
constructs \(S_p\) over the rationals in four dimensions, and verifies
(106.8)--(106.9) in degrees \(0,1,2,3\).  It emits
\[
 \texttt{artifacts/sm\_v96\_dyadic\_curvature\_blocking.json}.       \tag{106.14}
\]
The hashes are
\[
\begin{array}{c|l}
\hbox{object}&\hbox{SHA--256}\\ \hline
\hbox{verifier}&
\texttt{447F704CC676CE1E63270B037624FB0D8290EB45C77219B2577179172CD12C00}\\
\hbox{canonical payload}&
\texttt{6D799F8ACD7B16582E4D30C712971D87C8F2F57547B2703616238A65E3810433}.
\end{array}                                                   \tag{106.15}
\]

\begin{firewall}[A blocking map is not a renormalisation trajectory]
\label{fire:v18v96}
Theorems \ref{thm:v18v96-four-face-word} and
\ref{thm:v18v96-linear-curvature-block} provide the exact kinematic
coarsening map and its derivative.  They do not construct probability laws
invariant under that map, show that a Standard Model action is stable under
blocking, preserve reflection positivity after integrating fine variables,
or control the Einstein and fermion sectors.
\end{firewall}

\begin{decision}[Next bottleneck]
\label{dec:v18v96-next}
V97 should quantify the nonlinear deviation of the rescaled block from
\(B_2\).  The first target is a local bound uniform in the global lattice
size, expressed in terms of the four fine face-log norms and the three
transport-link norms.  Only then can the V94 curvature charts be compared
across mesh scales.
\end{decision}

\begin{assessment}[V96 acquisition]
\label{ass:v18v96}
V96 replaces the incorrect raw-restriction model of refinement by an exact,
associative, gauge-covariant link block.  The coarse face is an exact word
of four transported fine faces, and the continuum-scaled linear curvature
block is their average.  This supplies a nontrivial refinement kinematics;
the interacting measure flow remains open.
\end{assessment}

\begin{record}[V96 append record]
\label{rec:v18v96}
V96 is appended to the validated V95 whose installed SHA--256 was
\texttt{4776AEE9E96771ADAFCE1C957340F4D96BE00BB3B0C5EB4D334EB737146E3A6C}.
After installation only V96 is the active numeric SM note; the frozen
archives and unrelated notes are unchanged.
\end{record}


\section{V97: nonlinear dyadic curvature blocking is asymptotically linear}
\label{sec:v18v97}

V96 identified the exact four-plaquette word and its linearised four-face
average.  This section controls the complete nonlinear logarithm.  The
parallel-transport error is proportional to one link logarithm times one
face logarithm, while the four-factor BCH error is quadratic in the face
logs.  On the rooted chart at mesh \(1/(2N)\), the continuum-rescaled
blocking error is \(O(N^{-1})\) with an explicit constant independent of
the global volume.

\subsection{A local four-face estimate}

Use the notation of (106.3)--(106.5), and let
\[
 f:=\max_{i,j}\|\log p_{ij}\|_{\rm op}.                    \tag{107.1}
\]
Assume that every link logarithm occurring in the three transport paths
\(h_{00}\), \(h_{00}v_{10}\), and \(v_{00}\) has norm at most \(a\).

\begin{theorem}[Nonlinear block versus the raw face sum]
\label{thm:v18v97-local-block}
If \(4f<1/2\) and \(F^{\rm c}=\log P^{\rm c}\) is taken in the principal
chart, then
\[
\boxed{
 \left\|F^{\rm c}-\sum_{i,j=0}^1\log p_{ij}\right\|_{\rm op}
 \le\left[2(e^{2a}-1)+(e^{4a}-1)\right]f+B(4f).}           \tag{107.2}
\]
\end{theorem}

\begin{proof}
Write the four factors in (106.5) as \(e^{Y_1},\ldots,e^{Y_4}\), in their
displayed order.  Unitary conjugation preserves operator norm, so
\(\sum_k\|Y_k\|\le4f\).  The V77 logarithm estimate, which depends only on
the sum of the factor norms, gives
\[
 \left\|\log\prod_{k=1}^4e^{Y_k}-\sum_{k=1}^4Y_k\right\|
 \le B(4f).                                                \tag{107.3}
\]
For a one-link transport \(T=e^A\), (101.12) gives
\[
 \|\operatorname{Ad}_T X-X\|\le(e^{2a}-1)\|X\|.            \tag{107.4}
\]
There are two such transports.  For the two-link transport
\(T=e^Ae^{A'}\),
\[
 \|\operatorname{Ad}_T-I\|
 \le e^{2a}e^{2a}-1=e^{4a}-1,                             \tag{107.5}
\]
obtained by telescoping the two conjugations and using
\(\|\operatorname{Ad}_{e^A}\|\le e^{2a}\).
Apply (107.4)--(107.5) to
\(\sum_kY_k-\sum_{i,j}\log p_{ij}\), then add (107.3).
\end{proof}

\begin{corollary}[The first nonlinear terms have the correct dimensions]
\label{cor:v18v97-dimensions}
For small \(a\) and \(f\), the right side of (107.2) is
\(O(af+f^2)\).  In particular, it has no term linear in \(a\) alone or
linear in \(f\) beyond the four-face sum.
\end{corollary}

\begin{proof}
The transport factors \(e^{2a}-1\) and \(e^{4a}-1\) vanish at \(a=0\),
and \(B(4f)=O(f^2)\) by (87.8).
\end{proof}

\subsection{Uniform estimate on the rooted fine chart}

Let a fine field on \(\Lambda_{2N}\) lie in the rooted chart of V94.  Then
(104.2) and (104.7), with \(N\) replaced by \(2N\), give
\[
 a\le\frac1{600N},\qquad
 f\le\frac{11}{360000N^2}.                                \tag{107.6}
\]
Write
\[
 {\cal G}_{2N}:=(2N)^2F_{2N},\qquad
 {\cal G}_N^{\rm c}:=N^2F_N^{\rm c}.                       \tag{107.7}
\]

\begin{theorem}[Cofinal nonlinear blocking error]
\label{thm:v18v97-rescaled-error}
For every \(N\ge1\),
\[
\boxed{
 \|{\cal G}_N^{\rm c}-B_2{\cal G}_{2N}\|_\infty
 <\frac1{2000000N}.}                                      \tag{107.8}
\]
\end{theorem}

\begin{proof}
The identity
\[
 B_2{\cal G}_{2N}
 =\frac14S_2(2N)^2F_{2N}
 =N^2S_2F_{2N}                                           \tag{107.9}
\]
turns the left side of (107.8) into \(N^2\) times the left side of
(107.2).  From (107.6) and \(e^x-1\le x/(1-x)\),
\[
 e^{2a}-1\le\frac1{299N},\qquad
 e^{4a}-1\le\frac1{149N}.                                 \tag{107.10}
\]
Also \(4f\le11/(90000N^2)<1/300\), so (101.20) gives
\[
 B(4f)\le\frac85(4f)^2=\frac{128}{5}f^2.                  \tag{107.11}
\]
Substitution in (107.2) gives
\[
\begin{split}
 \|{\cal G}_N^{\rm c}-B_2{\cal G}_{2N}\|_\infty
 &\le\frac1N\,
 \frac{11}{360000}\left(\frac2{299}+\frac1{149}\right)
 +\frac1{N^2}\,\frac{128}{5}\frac{121}{360000^2}\\
 &=\frac1N\,\frac{2189}{5346120000}
   +\frac1{N^2}\,\frac{121}{5062500000}.                  \tag{107.12}
\end{split}
\]
Since \(N^{-2}\le N^{-1}\) and
\[
 \frac{2189}{5346120000}+\frac{121}{5062500000}
 =\frac{195478217}{451078875000000}
 <\frac1{2000000},                                        \tag{107.13}
\]
equation (107.8) follows.
\end{proof}

\begin{corollary}[Asymptotic commutation with the linear refinement map]
\label{cor:v18v97-asymptotic}
On the V94 rooted charts,
\[
 {\cal G}_N^{\rm c}-B_2{\cal G}_{2N}\longrightarrow0
\]
uniformly over all admitted fine coordinates as \(N\to\infty\).
\end{corollary}

\begin{proof}
The bound (107.8) is uniform over the admitted fine chart and tends to zero.
\end{proof}

\subsection{Executable certificate}

The verifier
\[
 \texttt{scripts/verify\_sm\_v97\_nonlinear\_dyadic\_block\_bound.py}
                                                                  \tag{107.14}
\]
checks the principal-chart threshold and every rational constant in
(107.10)--(107.13).  It emits
\[
 \texttt{artifacts/sm\_v97\_nonlinear\_dyadic\_block\_bound.json}. \tag{107.15}
\]
The hashes are
\[
\begin{array}{c|l}
\hbox{object}&\hbox{SHA--256}\\ \hline
\hbox{verifier}&
\texttt{21F8CF1980AF913C410C508D9D0DD0954F3854ED1A9069CA4E48E5E44B9B8EFD}\\
\hbox{canonical payload}&
\texttt{C024A83FE3B0B0039DBDC3A1C66D8AC83C095B7369F672D41349AC36C742B260}.
\end{array}                                                   \tag{107.16}
\]

\begin{firewall}[One-step convergence is not a continuum law]
\label{fire:v18v97}
The estimate (107.8) controls a single kinematic dyadic block on small
fields.  It does not show that a blocked field remains inside the next
rooted chart, construct compatible probability laws, control accumulated
errors over an arbitrary refinement tower, or prove stability of an
interacting Standard Model--Einstein action.
\end{firewall}

\begin{decision}[Next bottleneck]
\label{dec:v18v97-next}
V98 should compare the blocked full curvature with the rooted coarse graph
coordinate
\(P_N^\circ F_N^{\rm c}\).  It must determine whether the projector's
\(O(N)\) stock destroys (107.8), or whether the exact closedness and
Bianchi structure retain the vanishing refinement error.  If the raw
projector loses a power, the programme should move upstream to link
blocking, where associativity and gauge covariance are exact.
\end{decision}

\begin{assessment}[V97 acquisition]
\label{ass:v18v97}
V97 proves a local complete nonlinear block estimate and an explicit
uniform \(O(N^{-1})\) error for continuum-scaled curvature.  Thus the exact
nonabelian dyadic block converges to the four-face average on the rooted
small-field charts.  Measure compatibility and the interacting continuum
remain open.
\end{assessment}

\begin{record}[V97 append record]
\label{rec:v18v97}
V97 is appended to the validated V96 whose installed SHA--256 was
\texttt{3D29ADE9AC72EC5D8851508687024CD827BD886486A33C31D1AAE1F2DB5A450D}.
After installation only V97 is the active numeric SM note; the frozen
archives and unrelated notes are unchanged.
\end{record}


\section{V98: rooted homotopy commutes with subdivision, but the projected error plateaus}
\label{sec:v18v98}

V97 left open whether the \(O(N)\) stock of the rooted face projector would
destroy the vanishing one-step block error.  The first part of this section
proves an exact subdivision identity: the rooted homotopy and both of its
degree-two projectors commute with dyadic block summation.  Consequently the
fine graph correction cancels from the comparison.  The remaining nonlinear
four-face error has a uniform projected bound, but the available factorwise
estimate no longer tends to zero.

\subsection{Homotopy and subdivision}

Retain the raw dyadic sums \(S_p:C^p(\Lambda_{2N})\to
C^p(\Lambda_N)\) from (106.7).

\begin{theorem}[Rooted homotopy commutes with dyadic summation]
\label{thm:v18v98-homotopy-block}
For every positive degree \(1\le p\le4\),
\[
\boxed{
 H_p^N S_p=S_{p-1}H_p^{2N}.}                              \tag{108.1}
\]
\end{theorem}

\begin{proof}
Fix a coarse output \((p-1)\)-cell.  In a summand of the rooted tensor
homotopy (98.4), every coordinate before the contracted coordinate is
evaluated at the common origin.  The contracted coordinate is summed over
the prefix ending at the output cell, and the remaining \(p-1\) oriented
coordinates are left unchanged.

On the left of (108.1), \(S_p\) first sums the two fine choices in every
oriented input coordinate.  The subsequent coarse prefix sum groups the
fine contracted edges into consecutive pairs.  On the right,
\(H_p^{2N}\) first performs the full fine prefix sum, while \(S_{p-1}\)
sums the two fine choices in every remaining oriented coordinate.  Both
orders therefore visit exactly the same fine input cells, once each, with
the same cubical sign.  Equality holds row by row.
\end{proof}

Define the rooted degree-two projections
\[
 P_N^\circ=d_1^NH_2^N,\qquad
 Q_N^\circ=I-P_N^\circ=H_3^Nd_2^N.                        \tag{108.2}
\]

\begin{corollary}[Both projections commute with face blocking]
\label{cor:v18v98-projector-block}
One has
\[
 P_N^\circ S_2=S_2P_{2N}^\circ,\qquad
 Q_N^\circ S_2=S_2Q_{2N}^\circ,                           \tag{108.3}
\]
and the same identities with \(S_2\) replaced by \(B_2=S_2/4\).
\end{corollary}

\begin{proof}
Use (108.1) and the subdivision identity (106.8):
\[
 P_N^\circ S_2
 =d_1^NH_2^NS_2
 =d_1^NS_1H_2^{2N}
 =S_2d_1^{2N}H_2^{2N}
 =S_2P_{2N}^\circ.                                       \tag{108.4}
\]
The identity for \(Q^\circ=I-P^\circ\) follows by subtraction, and scalar
normalisation does not affect commutation.
\end{proof}

\subsection{Exact coordinate comparison}

Let a fine rooted-chart field have full and compatible face variables
\[
 {\cal F}_{2N}:=(2N)^2\Psi_{2N}^\circ(G_{2N}),\qquad
 {\cal G}_{2N}:=(2N)^2G_{2N}.                             \tag{108.5}
\]
Block its links by (106.1), let \(F_N^{\rm c}\) be the resulting coarse
plaquette logarithm, and define
\[
 {\cal F}_N^{\rm c}:=N^2F_N^{\rm c},\qquad
 {\cal G}_N^{\rm c}:=P_N^\circ{\cal F}_N^{\rm c}.          \tag{108.6}
\]

\begin{theorem}[Projected block-error identity]
\label{thm:v18v98-projected-identity}
The compatible coarse coordinate obeys the exact identity
\[
\boxed{
 {\cal G}_N^{\rm c}-B_2{\cal G}_{2N}
 =P_N^\circ\bigl({\cal F}_N^{\rm c}-B_2{\cal F}_{2N}\bigr).}
                                                                  \tag{108.7}
\]
In particular, no separate fine graph-height term appears.
\end{theorem}

\begin{proof}
Equation (104.6), multiplied by \((2N)^2\), gives
\(P_{2N}^\circ{\cal F}_{2N}={\cal G}_{2N}\).  Apply
Corollary \ref{cor:v18v98-projector-block}:
\[
\begin{split}
 {\cal G}_N^{\rm c}-B_2{\cal G}_{2N}
 &=P_N^\circ{\cal F}_N^{\rm c}
   -B_2P_{2N}^\circ{\cal F}_{2N}\\
 &=P_N^\circ\bigl({\cal F}_N^{\rm c}-B_2{\cal F}_{2N}\bigr).
\end{split}
\]
This is (108.7).
\end{proof}

\begin{corollary}[Uniform factorwise bound]
\label{cor:v18v98-uniform-bound}
The V97 estimate implies
\[
 \|{\cal G}_N^{\rm c}-B_2{\cal G}_{2N}\|_\infty
 <\frac3{500000}.                                         \tag{108.8}
\]
\end{corollary}

\begin{proof}
Every face row of \(d_1\) has four unit incidences, while (104.1) gives
\(\|H_2^N\|\le3N\).  Hence
\[
 \|P_N^\circ\|_{\infty\to\infty}
 \le12N.                                                   \tag{108.9}
\]
Combine (108.7), (108.9), and (107.8):
\[
 12N\cdot\frac1{2000000N}=\frac3{500000}.
\]
\end{proof}

\begin{assessment}[The remaining plateau]
\label{ass:v18v98-plateau}
The exact diagram (108.7) removes the potentially constant contribution of
the fine graph correction.  The remaining factorwise bound (108.8) is small
and uniform, but it does not tend to zero.  Therefore the present estimates
do not prove that the compatible rooted coordinates commute asymptotically
with refinement.  A vanishing claim would require cancellation inside
\(P_N^\circ({\cal F}_N^{\rm c}-B_2{\cal F}_{2N})\), not another estimate of
the two factors separately.
\end{assessment}

\begin{remark}[Why link space remains the exact projective level]
\label{rem:v18v98-link-level}
Dyadic link blocking is associative and gauge covariant by Proposition
\ref{prop:v18v96-link-block}.  Projection to \(P_N^\circ F\) discards
link transport and gauge-slice data, and the present curvature chart has not
been proved to reconstruct the blocked link field.  Thus exact iteration is
currently justified on link configurations and their gauge classes, while
the compatible curvature coordinate is an analytic local chart used for
bounds.
\end{remark}

\subsection{Executable certificate}

The verifier
\[
 \texttt{scripts/verify\_sm\_v98\_rooted\_block\_projector.py}       \tag{108.10}
\]
constructs \(H_p,S_p,P^\circ,Q^\circ\) over the rationals on a dyadic
four-box, checks (108.1) in every positive degree and (108.3) for both
degree-two projections, and records the bound (108.8).  It emits
\[
 \texttt{artifacts/sm\_v98\_rooted\_block\_projector.json}.          \tag{108.11}
\]
The hashes are
\[
\begin{array}{c|l}
\hbox{object}&\hbox{SHA--256}\\ \hline
\hbox{verifier}&
\texttt{DCD1D10D7CB1AC2136F79E6400240A27498A4ED34B34317AC2FBAA5890E59DAE}\\
\hbox{canonical payload}&
\texttt{173BA8D63821D6ECBB4DB7EA57D51BE80B2DB3DA337C72F5016DD99830D025BA}.
\end{array}                                                   \tag{108.12}
\]

\begin{firewall}[A uniform plateau is not asymptotic consistency]
\label{fire:v18v98}
The constant \(3/500000\) is an upper bound on a small-field coordinate
mismatch.  It does not vanish with the mesh and cannot be called a
continuum consistency theorem.  No probability law, interacting
renormalisation trajectory, reflection-positive limit, fermion sector, or
Einstein coupling is supplied.
\end{firewall}

\begin{decision}[Next upstream move]
\label{dec:v18v98-next}
V99 should remain on the exact link-projective level.  It should construct
a nontrivial family of link laws that is exactly consistent under dyadic
multiplication, beginning with convolution-semigroup measures on the compact
gauge group.  The result must be labelled kinematic unless an interacting
plaquette action and reflection positivity are also proved.
\end{decision}

\begin{assessment}[V98 acquisition]
\label{ass:v18v98}
V98 proves exact commutation of the rooted homotopy and both curvature
projections with dyadic subdivision.  The projected comparison reduces to
one nonlinear blocking error and has a uniform bound \(3/500000\), but the
available estimate plateaus.  Exact projectivity therefore moves upstream
to link space.
\end{assessment}

\begin{record}[V98 append record]
\label{rec:v18v98}
V98 is appended to the validated V97 whose installed SHA--256 was
\texttt{D5B2414339CE0DD358CFA116ABA1CB59C3C82E4228B012069382867D22D084E1}.
After installation only V98 is the active numeric SM note; the frozen
archives and unrelated notes are unchanged.
\end{record}


\section{V99: an exact gauge-invariant projective link law}
\label{sec:v18v99}

V98 showed that exact iteration should remain on link configurations.  This
section constructs a nontrivial projective probability law there.  A
convolution semigroup on the compact gauge group makes ordered products of
two fine links have exactly the next coarse law.  Averaging vertex gauges
then gives a gauge-invariant projective family.  The construction is
kinematic: it is not the interacting Yang--Mills--Higgs--fermion measure.

\subsection{Convolution-semigroup link increments}

Let \(K\) be a compact gauge group with normalized Haar measure.  Let
\((\mu_t)_{t\ge0}\) be a central symmetric convolution semigroup:
\[
 \mu_s*\mu_t=\mu_{s+t},\qquad
 \iota_\#\mu_t=\mu_t,\quad \iota(g)=g^{-1}.                \tag{109.1}
\]
The heat-kernel measures on \(K\) are the principal example.  Fix
\(\tau>0\), put
\[
 t_N=\frac{\tau}{N},                                      \tag{109.2}
\]
and on the positively oriented links \(E_N\) define
\[
 \nu_N:=\bigotimes_{e\in E_N}\mu_{t_N}.                   \tag{109.3}
\]

\begin{theorem}[Exact dyadic consistency of the product link law]
\label{thm:v18v99-product-projective}
For the link-blocking map \(\mathfrak B_1\) of (106.1),
\[
 (\mathfrak B_1)_\#\nu_{2N}=\nu_N.                         \tag{109.4}
\]
The family is invariant under lattice coordinate reflections.
\end{theorem}

\begin{proof}
Each coarse link is the ordered product of two fine links with laws
\(\mu_{\tau/(2N)}\).  Their product law is
\[
 \mu_{\tau/(2N)}*\mu_{\tau/(2N)}=\mu_{\tau/N}.             \tag{109.5}
\]
Distinct coarse edges use disjoint pairs of fine edges, so their blocked
values remain independent.  Fine edges unused by the block map integrate
to one.  This proves (109.4).  A coordinate reflection permutes links and
may reverse their orientations.  Product structure and inversion symmetry
in (109.1) preserve \(\nu_N\).
\end{proof}

\subsection{Gauge averaging}

Let \({\cal G}_N=K^{V_N}\) be the vertex gauge group, with product Haar
measure \(dg\).  For \(g\in{\cal G}_N\), write
\[
 (T_gU)_{xy}=g_xU_{xy}g_y^{-1}.                            \tag{109.6}
\]
Define the gauge-averaged law
\[
 \overline\nu_N
 :=\int_{{\cal G}_N}(T_g)_\#\nu_N\,dg.                    \tag{109.7}
\]

\begin{theorem}[Gauge-invariant projective family]
\label{thm:v18v99-gauge-projective}
Each \(\overline\nu_N\) is invariant under the full vertex gauge group, and
\[
 (\mathfrak B_1)_\#\overline\nu_{2N}
 =\overline\nu_N.                                         \tag{109.8}
\]
It is also invariant under coordinate reflections.
\end{theorem}

\begin{proof}
For a fixed \(h\in{\cal G}_N\), left translation \(g\mapsto hg\) preserves
product Haar measure.  Applying this change of variable in (109.7) proves
gauge invariance.

Proposition \ref{prop:v18v96-link-block} gives
\[
 \mathfrak B_1(T_gU)=T_{r(g)}(\mathfrak B_1U),             \tag{109.9}
\]
where \(r(g)\) is the restriction of \(g\) to the even vertices.  The
pushforward of product Haar measure on \({\cal G}_{2N}\) under \(r\) is
product Haar measure on \({\cal G}_N\).  Push (109.7) through
\(\mathfrak B_1\), use (109.9), and then use (109.4); the result is
(109.7) at scale \(N\), proving (109.8).  Reflection invariance follows
from Theorem \ref{thm:v18v99-product-projective} and invariance of product
Haar measure under the induced vertex permutation.
\end{proof}

\begin{corollary}[Existence of a kinematic projective-limit law]
\label{cor:v18v99-limit-law}
Along the dyadic refinement tower, the measures
\(\overline\nu_N\) define a probability measure on the inverse limit of
the compact link-configuration spaces under \(\mathfrak B_1\).
\end{corollary}

\begin{proof}
Every finite link-configuration space is compact Hausdorff, the blocking
maps are continuous, and (109.8) supplies consistent marginals.  The
standard projective-limit theorem for probability measures on compact
spaces gives the asserted measure.
\end{proof}

\subsection{The law is not the Haar product after gauge averaging}

For an irreducible unitary representation \(\pi\), normalize the heat-kernel
time so that
\[
 \int_K\pi(g)\,d\mu_t(g)=e^{-c_\pi t}I,\qquad c_\pi>0
\]
for nontrivial \(\pi\).

\begin{proposition}[A nonzero plaquette character]
\label{prop:v18v99-character}
Let \(P\) be a unit plaquette and
\(W_\pi(P)=d_\pi^{-1}\operatorname{Tr}\pi(P)\).  Then
\[
 \int W_\pi(P)\,d\overline\nu_N
 =e^{-4c_\pi\tau/N}>0.                                    \tag{109.10}
\]
Hence \(\overline\nu_N\) is not the independent Haar link law.
\end{proposition}

\begin{proof}
The Wilson character is gauge invariant, so gauge averaging does not change
its expectation.  Under \(\nu_N\), the four boundary links are independent.
Centrality and inversion symmetry give the same scalar Fourier transform
\(e^{-c_\pi t_N}I\) for each oriented or inversely oriented link.
Successive integration of the four representation matrices therefore gives
\(e^{-4c_\pi t_N}I\); normalized trace proves (109.10).  Under independent
Haar links the expectation of a nontrivial irreducible character is zero,
so the two laws differ.
\end{proof}

\subsection{Executable certificate}

The verifier
\[
 \texttt{scripts/verify\_sm\_v99\_projective\_link\_law.py}          \tag{109.11}
\]
checks the convolution pushforward over \(\mathbb Z/3\mathbb Z\) with exact
rational weights and exhausts all \(3^5=243\) link and endpoint-gauge
assignments for the covariance identity.  It emits
\[
 \texttt{artifacts/sm\_v99\_projective\_link\_law.json}.             \tag{109.12}
\]
Its hashes are
\[
\begin{array}{c|l}
\hbox{object}&\hbox{SHA--256}\\ \hline
\hbox{verifier}&
\texttt{61D4F6F18269472B74167B86BFBE57F7A2B8D270663F149892AD218831848F39}\\
\hbox{canonical payload}&
\texttt{2D2645CA56692F06ED98609D5B1D9398D51621D87962DC1AB47911E090C15D79}.
\end{array}                                                   \tag{109.13}
\]

\begin{firewall}[Kinematic projectivity is not the Standard Model measure]
\label{fire:v18v99}
The measure in Corollary \ref{cor:v18v99-limit-law} contains independent
heat-kernel link increments before gauge averaging.  It has no interacting
plaquette action, Higgs potential, chiral fermion determinant, anomaly
analysis, Einstein metric integration, or proved Osterwalder--Schrader
positivity.  Its link fluctuations are not confined to the V94 smooth
small-field chart.  Reflection invariance proved above is weaker than
reflection positivity.
\end{firewall}

\begin{decision}[V100 rollover checkpoint]
\label{dec:v18v99-next}
V100 must perform the mandatory companion audit and close this volume with
a precise acquisition ledger.  Under the standing rollover rule it must
then be frozen as \texttt{V100\_f18}, and the next active note restarted at
V1 with context and a summary.  The substantive post-rollover direction is
to replace the independent heat-kernel law by a local interacting density
while preserving exact block consistency or a controlled renormalisation
flow.
\end{decision}

\begin{assessment}[V99 acquisition]
\label{ass:v18v99}
V99 constructs the first nontrivial, exactly dyadic-projective,
gauge-invariant link law in this volume and proves a nonzero Wilson
character.  It is a rigorous kinematic continuum measure on the inverse
system, not the SM--Einstein continuum.
\end{assessment}

\begin{record}[V99 append record]
\label{rec:v18v99}
V99 is appended to the validated V98 whose installed SHA--256 was
\texttt{84921AE26EE4A69149585FBFB9F01B26A69095EA97AD34BFFFA73256DFC0191E}.
After installation only V99 is the active numeric SM note; the frozen
archives and unrelated notes are unchanged.
\end{record}


\section{V100: companion audit and close of the eighteenth note cycle}
\label{sec:v18v100}

This is the mandatory five-version audit and the terminal note of the
eighteenth active cycle.  V99 constructed an exact gauge-invariant
projective link law from a convolution semigroup.  The audit asks whether
the newest companion work changes the next interface test before the cycle
is frozen and restarted.

\subsection{Audited files}

The active companion sources were
\[
\begin{array}{c|l}
\hbox{programme and active file}&\hbox{SHA--256}\\ \hline
\hbox{Yang--Mills,
 \texttt{Renewal\_Yang\_Mills\_Mass\_Gap\_Research\_Notes\_V25.tex}}&
\texttt{5D6759B911E076304A17D09955E6777C3ACD6C99D3B47CD42BBA75EB0EA96BE6}\\
\hbox{Navier--Stokes,
 \texttt{Renewal\_NS\_Stability\_Research\_Notes\_V26.tex}}&
\texttt{82C887F8C2C69E4F28FAD7C3DC8DE5B9099CB5A4BF431EBA1FFF3E91935978AD}.
\end{array}                                                   \tag{110.1}
\]
The Navier--Stokes note is unchanged.  Yang--Mills has advanced from V23
to V25.

\begin{proposition}[Interface lesson from Yang--Mills V24]
\label{prop:v18v100-interface}
An exact construction at one stage of the programme cannot be used at the
next stage until the two stages' algebraic requirements are checked
together.  Applied here, exact projectivity of the V99 link law does not
imply that it has appreciable mass in the V94 logarithmic chart.
\end{proposition}

\begin{proof}
Yang--Mills V24 proves a concrete interface failure.  Its earlier fifteen
CRT moduli satisfied the abstract height requirement, but thirteen of them
admitted no element of order \(31\), even in the declared quadratic
extension, and therefore could not execute the later length-\(31\)
transform.  V24 replaced the moduli by fifteen primes satisfying the height,
denominator, and transform-root conditions simultaneously.  The lesson is
logical rather than numerical: separate correctness of two modules does
not prove their composability.

V99 proves convolution-semigroup projectivity on compact-group links.
V94 proves analytic control only when link logarithms are \(O(N^{-1})\)
and face coordinates are \(O(N^{-2})\).  Neither theorem estimates the
V99 probability of that event.  Their compatibility is therefore an open
interface condition, exactly as claimed.
\end{proof}

\begin{proposition}[Other companion conclusions]
\label{prop:v18v100-other}
Yang--Mills V25 retains the SM exact-link-block route for its later
continuum stage but imports no SM estimate into its finite spectral
certificate.  Navier--Stokes V26 supplies no new statement relevant to the
link-law/chart interface.
\end{proposition}

\begin{proof}
These are the explicit audit decisions in the two companion notes.
Yang--Mills V25 separates its finite Laurent sign computation from later
cofinal blocking.  Navier--Stokes V26 concerns rank-one contractions of
Oseen-kernel derivatives.  Neither provides a compact-group small-ball
probability.
\end{proof}

\subsection{Position at the cycle boundary}

\begin{assessment}[Principal acquisitions of V1--V100]
\label{ass:v18v100-cycle}
The eighteenth cycle established the following chain.
\begin{enumerate}
\item It built exact finite cubical right inverses and a
      reflection-averaged contracting homotopy, then corrected the
      four-dimensional prescribed-curvature claim by exposing the nonlinear
      Bianchi compatibility defect.
\item It parameterised the actual small nonabelian curvature image as an
      injective graph over \(\ker d_2\), proved the transported cube Bianchi
      identity, and obtained graph height \(O(N^{-2})\) with a uniform
      differential slope.
\item It proved that reflection averaging fails anchored restriction, that
      the rooted homotopy is projective, and that raw restriction combined
      with a global \(N^{-2}\) radius forces the zero field.
\item It replaced raw restriction by exact associative gauge-covariant
      dyadic link blocking, derived the transported four-plaquette word, and
      proved an \(O(N^{-1})\) nonlinear error for rescaled curvature.
\item It proved exact commutation of rooted homotopies and projectors with
      subdivision, isolated the remaining projected plateau, and constructed
      a nontrivial gauge-invariant kinematic projective link law from compact
      convolution semigroups.
\end{enumerate}
None of these results supplies the interacting Standard Model measure,
fermion determinant, Higgs vacuum control, Einstein-sector continuum,
reflection positivity, or reconstruction of the physical theory.
\end{assessment}

\begin{firewall}[Cycle boundary]
\label{fire:v18v100}
The exact projective law of V99 is kinematic.  Calling it a continuum
Standard Model would ignore its untested small-field concentration, the
absence of the interacting action and matter sectors, and the missing
Osterwalder--Schrader and Einstein limits.
\end{firewall}

\begin{decision}[Post-rollover V1 acquisition]
\label{dec:v18v100-next}
After this file is frozen as
\texttt{Renewal\_SM\_Einstein\_Continuum\_Research\_Notes\_V100\_f18.tex},
the new active V1 will give context and a compact result ledger.  Its first
new theorem will test the V99 heat-kernel increment scale against the V94
\(O(N^{-1})\) link-log chart.  If the projective law concentrates at the
larger \(N^{-1/2}\) scale, the programme must change either the convolution
clock or the chart architecture before adding interactions.
\end{decision}

\begin{record}[V100 append record]
\label{rec:v18v100}
V100 is appended to the validated V99 whose installed SHA--256 was
\texttt{76F817BF8214D186346F8CBCB997127B8199BCA3540825993A6F8C3E40832C45}.
This file closes the eighteenth cycle and is to be preserved as the
\texttt{f18} archive.  Unrelated notes are unchanged.
\end{record}

\end{document}
